
\usepackage{fancybox}
\usepackage{tikz}
\usetikzlibrary{shapes}
\usetikzlibrary{decorations.pathmorphing}
\usepackage{caption,subcaption}
\usepackage{makeidx}
\usepackage{alltt}
\usepackage{amsmath}
\usepackage{amsfonts}
\usepackage{graphicx}
\usepackage{rotating}
\usepackage{booktabs}
\usepackage{pifont}

\usepackage{cmap}

\usepackage[utf8]{inputenc}
\usepackage[math]{anttor}
\usepackage{libris}
\renewcommand*{\ttdefault}{txtt}

\newcommand*{\docdate}{\formatdate{25}{9}{2012}}

\authordetails{1970}{Talbot}{Nicola L.~C.}
\title[LaTeX for Complete Novices]{\LaTeX\ for Complete Novices}
\version{1.4}
\edition{1}
\volume{1}
\series{The Dickimaw \LaTeX\ Series}
\seriesurl{http://www.dickimaw-books.com/latex/}
\affiliation
 {Dickimaw Books}
 {http://www.dickimaw-books.com/}
 {Saxlingham Nethergate}
\keywords{LaTeX,beginners,introduction,typesetting}
\subject{LaTeX}
\date{\docdate}

\html{
\newcommand*{\reportlinebreak}{}
\newcommand*{\reportpagebreak}{}
\newcommand*{\bookpagebreak}{}
\newcommand*{\booklinebreak}{}
\newcommand*{\screenpagebreak}{}
\newcommand*{\screenlinebreak}{}

\newcommand*{\dohtmlcolorbox}[2][0.8\linewidth]{%
 \colorbox{white}{%
   \begin{minipage}{#1}%
   \setlength{\parindent}{1.5em}%
    #2%
   \end{minipage}
 }%
}

\newcommand{\suppmaterial}[2][]{\htmlref{#2}{ch:supp}#1}
\newcommand{\suppref}[3][]{\htmlref{#2}{sec:#3}#1}

\providecommand{\dosidewaysfigure}[3]{}

\newcommand{\mathorit}[1]{\textit{#1}}

\newcommand{\objectref}[2]{\htmlref{#1~}{#2}\ref{#2}}
\newcommand{\figureref}[1]{\objectref{Figure}{#1}}
\newcommand{\xfigureref}[1]{\objectref{Figure}{#1}}
\newcommand{\vfigureref}[1]{\objectref{Figure}{#1}}
\newcommand{\xpageref}[2]{\htmlref{#1}{#2}}

\providecommand*{\vref}[1]{\ref{#1}}

\newcommand{\figuretop}[1]{}
\newcommand{\makeimg}[2]{$\htmlimage{ALT=#1}{}\mbox{#2}$}

\newcommand{\tableref}[1]{\objectref{Table}{#1}}

\newcommand{\exerciseref}[1]{\objectref{Exercise}{#1}}

\newcommand{\modification}[1]{\textbf{\texttt{#1}}}

\newcommand{\comptexttex}{\htmladdnormallink{comp.text.tex}{http://groups.google.com/group/comp.text.tex}}
\newcommand{\texhax}{\htmladdnormallink{texhax}{http://tug.org/pipermail/texhax/}}

\newcommand{\latexthesis}{%
   \htmladdnormallink{Using LaTeX to Write a PhD
Thesis}{\baseurl/latex/thesis/index.html}%
}

\newcommand{\latexpdfdoc}{\htmladdnormallink{Creating
a PDF Document using PDFLaTeX}{\baseurl/latex/pdfdoc/}}

\newenvironment{resultS}[1][]{\begin{htmlresult}{#1}\begin{makeimage}}{\end{makeimage}\end{htmlresult}}

\newenvironment{result}[1][]{\begin{htmlresult}{#1}\begin{makeimage}}{\end{makeimage}\end{htmlresult}}

\newcounter{object}

\newenvironment{inlinedef}{\begin{ttfamily}}{\end{ttfamily}}

\newcommand{\casemedia}[4]{#4}

\newcommand{\download}[1]{%
 \htmladdnormallink{download}{\exerciseurl/#1.tex}%
}

\newcommand{\downloadorview}[1]{%
 \htmladdnormallink{download}{\exerciseurl/#1.tex} or
 \htmladdnormallink{view}{\exerciseurl/#1.html}%
}

\newcommand{\doifbook}[1]{}
\newcommand{\doifnotbook}[1]{#1}
\newcommand{\doifreport}[1]{}

\newcommand{\doublebar}{\ensuremath{\|}}

\newcommand{\istystart}[1]{%
 \isty{#1}%
}

\newcommand{\istyend}[1]{%
 \isty{#1}%
}

\newcommand{\idxpagestyle}[1]{\index{page style!#1}}

\newcommand{\pagenumberingfmt}[1]{\texttt{#1}}

\newcommand{\indexpagenumbering}[1]{%
  \index{page numbering!#1}%
}

\newcommand{\ipagenumbering}[1]{%
  \pagenumberingfmt{#1}\indexpagenumbering{#1}%
}

\newcommand{\iunit}[2]{%
  \texttt{#1}\index{#1 (#2)@\texttt{#1} (#2)}%
  \index{units!\texttt{#1} (#2)}%
}

}

\makeindex
\makeglossaries

\newkeyword[\glstockey={Source Code}]{source}{source code}

\newkeyword[\glstockey={Output File}]{output}{output file}

\newkeyword
 [firstplural={\protect\emph{DVI (DeVice Independent) files}},
  \glstockey={DVI File}]
{dvi}{DVI file}

\newkeyword
 [\glsindexkey={auxiliary file (.aux)@auxiliary file (\texttt{.aux})},
  \glsextrakey={\extraindexaux},
  name={Auxiliary Files},\glstockey={Auxiliary Files}]
 {auxiliary}{auxiliary file}

\newkeyword
 [name={Terminal or Command Prompt},%
  \glsextrakey={\extraindexterminal},
  \glstockey={Terminal or Command Prompt}
 ]
{terminal}{terminal}

\newkeyword
 [name={Commands},\glstockey={Commands}]
 {command}{command}

\newkeyword
 [name={Grouping (or Scope)},%
  \glsextrakey={\extraindexgroup},
  \glstockey={Grouping (or Scope)}
 ]
 {group}{group}%

\newkeyword
 [name={Arguments (also called \dq{Parameters})},%
  \glsextrakey={\extraindexarg},
  \glstockey={Arguments}
 ]
{argument}{argument}%

\newkeyword
 [parent=argument,
  \glsindexkey={argument!mandatory},
  \glsextrakey={\extraindexmandatory},
  first={mandatory (or compulsory) argument},
  firstplural={mandatory (or compulsory) arguments},
  name={Mandatory Arguments},
  \glstockey={Mandatory Arguments}
 ]
{mandatory}{mandatory argument}

\newkeyword
 [parent=argument,
  \glsindexkey={argument!optional},
  \glsextrakey={\extraindexoptional},
  name={Optional Arguments},
  \glstockey={Optional Arguments}
]
{optional}{optional argument}


\newkeyword
 [\glsextrakey={\extraindexfragile},
  \glsindexkey={command!fragile},
  \glstockey={Moving Arguments and Fragile Commands},
   name={Moving Arguments and Fragile Commands}
 ]
{fragile}{fragile command}

\newkeyword
 [name={Robust Commands},
  \glstockey={Robust Commands},
  \glsindexkey={command!robust},
  \glsextrakey={\extraindexrobust}
 ]
{robust}{robust command}%

\newkeyword
 [name={Short and Long Commands},
  \glstockey={Short and Long Commands},
  \glsindexkey={command!short},
  \glsextrakey={\extraindexshortlong}]
{short}{short command}%


\newkeyword
[name={Declarations},\glstockey={Declarations}]
{declaration}{declaration}


\newkeyword
[name={Inter-Sentence Spacing},
 \glstockey={Inter-Sentence Spacing},
 \glsindexkey={spacing!inter-sentence},
 \glsextrakey={\extraindexsentencespacing}
]
{intersentencespacing}{inter-sentence spacing}%

\newkeyword
 [\glsingkey={hyphenating},
  \glsdkey={hyphenated},
  \glstockey={Hyphenation}]
 {hyphenation}{hyphenation}

\newkeyword
[name={Environments},
 \glstockey={Environments}]
{environment}{environment}%

\newkeyword
[name={The Preamble},
 \glstockey={The Preamble}]
{preamble}{preamble}%


\newkeyword
[name={Lengths},\glstockey={Lengths}]
{length}{length}

\newkeyword[\glstockey={Class File}]{cls}{class file}%

\newkeyword
[\glsindexkey={TeX@\protect\TeX},\glstockey={\texorpdfstring{\protect\TeX}{TeX}}]
{tex}{\TeX}

\newkeyword{perl}{Perl}

\glsaddall[types=keywords]
\glsunsetall[keywords]

\renewcommand{\summarypreamble}{%
Commands or environments defined in the \LaTeX\ kernel are always
available.%
}

\defgactivechar
 {backslashchar}
 {\backslashsym}
 {}
 {\LaTeX\ Kernel}
 {Escape character (indicates a command).}
 {}

\defgactivechar
 {dbbackslashchar}
 {\dbbackslashsym}
 {\oarg{\meta{height}}}
 {\LaTeX\ Kernel}
 {\nopostdesc}
 {%
   \BeginArgList
    \csentryargitem{height} Extra vertical space.
   \EndArgList
 }

\defgchildactivechar
 {newline.dbbackslashchar}
 {dbbackslashchar}
 {\dbbackslashsym}
 {Breaks a line without justification (\protect\htmlref{starred form}{itm:starredcommand}
   forbids a page break)}

\defgchildactivechar
 {tab.dbbackslashchar}
 {dbbackslashchar}
 {\dbbackslashsym}
 {Starts a new row in tabular-style environments}

\defgcs{tabularnewline}%
 {}%
 {\LaTeX\ Kernel}%
 {Behaves like \nxglsi{tab.dbbackslashchar} in a
   \nxglsni{env-tabular}-like environment but helps to disambiguate a
    \nxglslink{newline.dbbackslashchar}{line break} in a paragraph
    cell from a \nxglslink{tab.dbbackslashchar}{row separator}.}%
 {}

\defgidxactivechar
 {leftbracechar}
 {\leftbracesym}
 {}
 {\LaTeX\ Kernel}
 {Marks the beginning of a \nxglsi{group}.}
 {}

\defgidxactivechar
 {rightbracechar}
 {\rightbracesym}
 {}
 {\LaTeX\ Kernel}
 {Marks the end of a \nxglsi{group}.}
 {}

\defgidxactivecharcs
 {leftbrace}
 {\leftbracesym}
 {}
 {\LaTeX\ Kernel}
 {Left brace \{ character. In math mode may be used as a delimiter.}
 {}

\defgidxactivecharcs
 {rightbrace}
 {\rightbracesym}
 {}
 {\LaTeX\ Kernel}
 {Right brace \{ character. In math mode may be used as a delimiter.}
 {}

\defgidxactivechar
 {emdash}%
 {\emDashcs}%
 {}%
 {\LaTeX\ Kernel}%
 {Em-dash \textemdash\ symbol. (Normally used to
  indicate omissions or interruptions or to highlight a parenthetical element.)
  See also \nxglsni{textemdash}.}%
 {}

\defgidxactivechar
 {endash}%
 {\enDashcs}%
 {}%
 {\LaTeX\ Kernel}%
 {En-dash \textendash\ symbol.  (Normally used for number ranges.)
  See also \nxglsni{textendash}.}%
 {}

\defgidxactivechar
 {questiondown}%
 {\questiondowncs}%
 {}%
 {\LaTeX\ Kernel}%
 {Upside-down question mark \textquestiondown\ symbol. See also \nxglsni{textquestiondown}.}%
 {}

\defgidxactivechar
 [!Y]
 {exclamdown}%
 {\exclamdowncs}%
 {}%
 {\LaTeX\ Kernel}%
 {Upside-down exclamation mark \textexclamdown\ symbol. See also \nxglsni{textexclamdown}.}%
 {}

\defgchar
 {quoteleft}%
 {\quoteleftcs}%
 {}%
 {\LaTeX\ Kernel}%
 {Open quote \textquoteleft\ symbol. See also \nxglsni{textquoteleft}.}%
 {}

\defgchar
 {quotedblleft}%
 {\quotedblleftcs}%
 {}%
 {\LaTeX\ Kernel}%
 {Open double quote \textquotedblleft\ symbol. See also \nxglsni{textquotedblleft}.}%
 {}

\defgchar
 {quoteright}%
 {\quoterightcs}%
 {}%
 {\LaTeX\ Kernel}%
 {Closing quote or apostrophe \textquoteright\ symbol in text mode 
  or prime symbol \ensuremath{'} in math mode. See also \nxglsni{textquoteright}.}%
 {}

\defgchar
 {quotedblright}%
 {\quotedblrightcs}%
 {}%
 {\LaTeX\ Kernel}%
 {Closing double quote \textquotedblright\ symbol in text mode
  or double prime \ensuremath{''} in math mode. See also \nxglsni{textquotedblright}.}%
 {}

\defgchar
 {dash}%
 {\dashcs}%
 {}%
 {\LaTeX\ Kernel}%
 {Hyphen - in text mode or minus sign $-$ in math mode.}%
 {}

\defgidxactivecharcs
 {exclam}
 {\exclamsym}
 {}
 {\LaTeX\ Kernel (Math Mode)}
 {Negative thin space.}
 {}

\defgidxactivechar
 {exclamchar}
 {\exclamsym}
 {}
 {}
 {\nopostdesc}
 {}

\defgchildidxactivechar
 {sentence.exclamchar}
 {exclamchar}
 {\exclamsym}
 {Exclamation symbol (end of sentence marker)}
 

\defgchildidxactivechar
 {makeindex.exclamchar}
 {exclamchar}
 {\exclamsym}
 {\nxiappname{makeindex} sublevel special character}

\defgchildidxactivechar
 {resizebox.exclamchar}
 {exclamchar}
 {\exclamsym}
 {Used in \nxglsi{resizebox} to maintain aspect ratio}

\defgidxactivecharcs
 {csbar}
 {\vbarsym}
 {}
 {\LaTeX\ Kernel (Math Mode)}
 {Double vertical bar \doublebar{} delimiter}
 {}

\defgidxactivechar
 {barchar}
 {\vbarsym}
 {}
 {\LaTeX\ Kernel}
 {\nopostdesc}
 {}

\defgchildidxactivechar
 {delim.barchar}
 {barchar}
 {\vbarsym}
 {Delimiter. (Math mode only. Use \nxglsni{textbar} in text mode.)}

\defgchildidxactivechar
 {array.barchar}
 {barchar}
 {\vbarsym}
 {Vertical rule specifier (\nxglsi{env-tabular} or
   \nxglsi{env-array})}

\defgchar
 {slashchar}
 {\slashsym}
 {}
 {}
 {\nopostdesc}
 {}

\defgchildchar
 {text.slash}
 {slashchar}
 {\slashsym}
 {Forward slash symbol (see also \nxglsi{slash})}

\defgxchildchar
 {dir.slash}
 {slashchar}
 {\slashsym}
 {Directory divider}
 {directory divider}

\defgchildchar
 {delim.slash}
 {slashchar}
 {\slashsym}
 {Forward slash delimiter (math mode)}

\defgidxactivechar
 [!Z]
 {visiblespace}
 {\textvisiblespace}
 {}
 {}
 {A visual indication of a space in the code. When you type up 
   the code, replace all instances of this symbol with a space via the space bar on your 
   keyboard.}
 {}

\defgidxactivecharcs
 {at}
 {\atsym}
 {}
 {\LaTeX\ Kernel}
 {Used when a sentence ends with a capital letter.
   This command should be placed after the letter and before the
   punctuation mark.}
 {}

\defgidxactivechar
 {atchar}
 {\atsym}
 {\marg{\meta{text}}}
 {\LaTeX\ Kernel}
 {%
   Used in the argument of \nxglsi{env-tabular} or 
   \nxglsi{env-array} like environments to specify text to insert between columns.%
 }
 {%
     \BeginArgList
      \csentryargitem{text} Text to insert between columns
     \EndArgList
 }
  
\defgidxactivechar
 {questionchar}
 {\questionsym}
 {}
 {}
 {Question mark (end of sentence marker).}
 {}

\defgchar
 {ltchar}
 {\lesssym}
 {}
 {\LaTeX\ Kernel (Math Mode)}
 {Less than symbol. (Use \nxglsni{textless} in text mode.)}
 {}

\defgchar
 {gtchar}
 {\greatersym}
 {}
 {\LaTeX\ Kernel (Math Mode)}
 {Greater than symbol. (Use \nxglsni{textgreater} in text mode.)}
 {}

\defgchar
 {gtcol}
 {\greatersym}
 {\marg{\meta{decl}}}
 {\nxisty{array} package}
 {%
   Used in \nxglsi{env-tabular} or \nxglsi{env-array} column specifiers before
   \texttt{l}, \texttt{r}, \texttt{c}, \texttt{p}, \texttt{m} or
   \texttt{b} to insert \meta{decl} directly in front of the entry
   for that column.%
 }
 {%
     \BeginArgList
       \csentryargitem{decl} The code to insert at the start of the
         column.
     \EndArgList
 }

\defgchar
 {ltcol}
 {\lesssym}
 {\marg{\meta{decl}}}
 {\nxisty{array} package}
 {%
   Used in \nxglsi{env-tabular} or \nxglsi{env-array} column specifiers after
   \texttt{l}, \texttt{r}, \texttt{c}, \texttt{p}, \texttt{m} or
   \texttt{b} to insert \meta{decl} directly after the entry
   for that column.%
 }
 {%
     \BeginArgList
       \csentryargitem{decl} The code to insert at the end of the
         column.
     \EndArgList
 }

\defgcs{begin}%
 {\marg{\meta{env-name}}\oarg{\meta{env-option}}\marg{\meta{env-arg-1}}\ldots\marg{\meta{env-arg-n}}}%
 {\LaTeX\ Kernel}%
 {%
   Starts an environment. (Must have a matching \nxglsi{end}.)
 }%
 {%
   \BeginArgList
    \csentryargitem{env-name} The name of the environment. (\emph{No
      backslash.})
    \csentryargitem{env-option} An optional argument that may be 
      passed to the environment. Not all environments have optional
      arguments.
    \csentryargitem{env-arg-1}\ldots\meta{env-arg-n} Any mandatory
      arguments required by the environment. Not all environments
      require arguments.
   \EndArgList
 }

\defgcs{end}%
  {\marg{\meta{env-name}}}%
  {\LaTeX\ Kernel}%
  {Ends an environment. (Must have a matching \nxglsi{begin}.)}
  {%
    \BeginArgList
      \csentryargitem{env-name} The name of the environment.
    \EndArgList
  }

\defgcs{documentclass}%
 {\oarg{\meta{option-list}}\marg{\meta{class-name}}}%
 {\LaTeX\ Kernel}%
 {%
   Loads the document class file, which sets up the type of document
   you wish to write.%
 }%
 {%
   \BeginArgList
    \csentryargitem{option-list} A comma-separated list of options to
     pass to the class file or any packages that will later be
     loaded.
    \csentryargitem{class-name} The name of the document class. This
    corresponds to a file called \meta{class-name}\texttt{.cls},
    which must be installed. 
   \EndArgList
 }

\defgcs{usepackage}%
 {\oarg{\meta{option-list}}\marg{\meta{package-list}}}%
 {\LaTeX\ Kernel}%
 {%
   Loads the named packages.%
 }%
 {%
   \BeginArgList
    \csentryargitem{option-list} A comma-separated list of options to
     pass to the package.
    \csentryargitem{package-list} A comma-separated list of package
     names (without the \texttt{.sty} extension).
   \EndArgList
 }

\defgcs{footnote}%
 {\oarg{\meta{number}}\marg{\meta{text}}}%
 {\LaTeX\ Kernel}%
 {Inserts a footnote.}%
 {%
  \BeginArgList
    \csentryargitem{number} Overrides the default footnote number with
    the specified \meta{number}.
    \csentryargitem{text} The footnote text.
  \EndArgList
 }%

\defgcs{textbackslash}%
 {}%
 {\LaTeX\ Kernel (Text Mode)}%
 {Backlash \textbackslash\ symbol. (Use
  \nxglsi{backslash} for math mode.)}%
 {}

\defgcs{backslash}%
 {}%
 {\LaTeX\ Kernel (Math Mode)}%
 {Backslash \protect\ensuremath{\protect\backslash} symbol, which may be used
  as a delimiter. (Use
 \nxglsi{textbackslash} for text mode.)}%
 {}

\defgactivecharcs
 [\underscoresym]
 {underscore}
 {\textunderscore}
 {}
 {\LaTeX\ Kernel}
 {Underscore \_ symbol (see also \nxglsni{textunderscore}).}
 {}

\defgcs{textunderscore}
 {}
 {\LaTeX\ Kernel}
 {Underscore \textunderscore\ symbol (see also \nxglsni{textunderscore}).}
 {}

\defgactivechar
 [\underscoresym]
 {underscorechar}
 {\_}
 {\marg{\meta{maths}}}
 {\LaTeX\ Kernel (Math Mode)}
 {Displays its argument as a subscript.}
 {%
   \BeginArgList
     \csentryargitem{maths} The subscript.
   \EndArgList
 }

\defgcs{textgreater}%
 {}%
 {\LaTeX\ Kernel (Text Mode)}%
 {Greater than \textgreater\ symbol. (Just use \texttt{\textgreater} in math mode.)}%
 {}

\defgcs{textless}%
 {}%
 {\LaTeX\ Kernel (Text Mode)}%
 {Less than \textless\ symbol. (Just use \texttt{\textless} in math mode.)}%
 {}

\defgcs{textasciicircum}%
 {}%
 {\LaTeX\ Kernel}%
 {Circumflex \textasciicircum\ symbol.}%
 {}%

\defgcs[][cshyphenation]{hyphenation}
 {\marg{\meta{word}}}
 {\LaTeX\ Kernel}
 {Specifies hyphenation points.}
 {%
   \BeginArgList
    \csentryargitem{word} Word with hyphen points indicated with a
     dash (\texttt{-}).
   \EndArgList
 }

\defgsymcs[hyphen]{\dashcs}
 {}
 {\LaTeX\ Kernel}
 {Insert discretionary hyphen.}
 {}

\defgactivecharcs
 {dollar}
 {\dollarsym}
 {}
 {\LaTeX\ Kernel}
 {Dollar \$ symbol.}
 {}

\defgactivechar
 {dollarchar}
 {\dollarsym}
 {}
 {\LaTeX\ Kernel}
 {Switches in and out of in-line math mode.}
 {}

\defgcs{textasciitilde}%
 {}%
 {\LaTeX\ Kernel}%
 {Tilde \textasciitilde\ symbol. (If you are typing an
  URL, use the \nxisty{url} package, which provides \nxglsi{url}\marg{\meta{address}}
  that allows you to directly type \textasciitilde\ in the address.)}%
 {}%

\defgcs{url}%
 {\marg{\meta{address}}}%
 {\nxisty{url} package}%
 {Typesets an URL in a typewriter font and allows you to use
  characters such as \textasciitilde.}%
 {%
   \BeginArgList
    \csentryargitem{address} The web address.
   \EndArgList
 }%

\defgcs{dag}%
 {}%
 {\LaTeX\ Kernel}%
 {Dagger \dag\ symbol.}%
 {}

\defgcs{ddag}%
 {}%
 {\LaTeX\ Kernel}%
 {Double-dagger \ddag\ symbol.}%
 {}

\defgcs{textregistered}%
 {}%
 {\LaTeX\ Kernel}%
 {Registered \textregistered\ symbol.}%
 {}

\defgcs{texttrademark}%
 {}%
 {\LaTeX\ Kernel}%
 {Trademark \texttrademark\ symbol.}%
 {}

\defgcs{copyright}%
 {}%
 {\LaTeX\ Kernel}%
 {Copyright \copyright\ symbol.}%
 {}

\defgcs{textquoteright}%
 {}%
 {\LaTeX\ Kernel}%
 {Closing single quote (or apostrophe) \textquoteright\ symbol.}%
 {}

\defgcs{textquoteleft}%
 {}%
 {\LaTeX\ Kernel}%
 {Opening single quote \textquoteleft\ symbol.}%
 {}

\defgcs{textquotedblright}%
 {}%
 {\LaTeX\ Kernel}%
 {Closing double quote \textquotedblright\ symbol.}%
 {}

\defgcs{textquotedblleft}%
 {}%
 {\LaTeX\ Kernel}%
 {Opening double quote \textquotedblleft\ symbol.}%
 {}

\defgcs{textbullet}%
 {}%
 {\LaTeX\ Kernel (Text Mode)}%
 {Bullet \textbullet\ symbol.}%
 {}

\defgcs{textquestiondown}%
 {}%
 {\LaTeX\ Kernel}%
 {Upside-down question mark \textquestiondown\ symbol.}%
 {}

\defgcs{textexclamdown}%
 {}%
 {\LaTeX\ Kernel}%
 {Upside-down exclamation mark \textexclamdown\ symbol.}%
 {}

\defgcs{textendash}%
 {}%
 {\LaTeX\ Kernel}%
 {En-dash \textendash\ symbol. (Normally used for number ranges.)
  See also \nxglsni{endash}.}%
 {}

\defgcs{textemdash}%
 {}%
 {\LaTeX\ Kernel}%
 {Em-dash \textemdash\ symbol. (Normally used to
  indicate omissions or interruptions or to highlight a parenthetical element.)
  See also \nxglsni{emdash}.}%
 {}

\defgcs{textperiodcentered}%
 {}%
 {\LaTeX\ Kernel (Text Mode)}%
 {Centred period \textperiodcentered\ symbol.}%
 {}

\defgcs{i}%
 {}%
 {\LaTeX\ Kernel}%
 {Dotless i character: \i.}%
 {}

\defgcs{j}%
 {}%
 {\LaTeX\ Kernel}%
 {Dotless j character: \dotlessj.}%
 {}

\defgactivecharcs
 {hash}
 {\hashsym}
 {}
 {\LaTeX\ Kernel}
 {Hash \# symbol.}
 {}

\defgactivechar
 {hashchar}
 {\hashsym}
 {\meta{digit}}
 {\LaTeX\ Kernel}
 {Replacement text for argument \meta{digit}.}
 {}

\defgactivecharcs
 {percent}
 {\percentsym}
 {}
 {\LaTeX\ Kernel}
 {Percent \% symbol}
 {}

\defgactivechar
 {percentchar}
 {\percentsym}
 {}
 {\LaTeX\ Kernel}
 {Comment character used to ignore everything up to and including
  the newline character in the \nxgls{source}.}
 {}

\defgactivecharcs
 {amp}
 {\ampsym}
 {}
 {\LaTeX\ Kernel}
 {Ampersand \& symbol}
 {}

\defgactivechar
 {ampchar}
 {\ampsym}
 {}
 {\LaTeX\ Kernel}
 {Alignment tab.}
 {}

\defgcs{S}%
 {}%
 {\LaTeX\ Kernel}%
 {Sectional \S\ symbol.}%
 {}

\defgcs{P}%
 {}%
 {\LaTeX\ Kernel}%
 {Paragraph \P\ symbol.}%
 {}

\defgcs{slash}%
 {}%
 {\LaTeX\ Kernel}%
 {Forward slash \slash\ symbol.}%
 {}

\defgcs{AE}%
 {}%
 {\LaTeX\ Kernel}%
 {\protect\AE\ ligature.}%
 {}

\defgcs{ae}%
 {}%
 {\LaTeX\ Kernel}%
 {\protect\ae\ ligature.}%
 {}

\defgcs{OE}%
 {}%
 {\LaTeX\ Kernel}%
 {\protect\OE\ ligature.}%
 {}

\defgcs{oe}%
 {}%
 {\LaTeX\ Kernel}%
 {\protect\oe\ ligature.}%
 {}

\defgcs{AA}%
 {}%
 {\LaTeX\ Kernel}%
 {Upper case A-ring \AA\ character.}%
 {}

\defgcs{aa}%
 {}%
 {\LaTeX\ Kernel}%
 {Lower case a-ring \aa\ character.}%
 {}

\defgcs{L}%
 {}%
 {\LaTeX\ Kernel}%
 {Upper case L-bar \L\ character.}%
 {}

\defgcs{l}%
 {}%
 {\LaTeX\ Kernel}%
 {Lower case l-bar \l\ character.}%
 {}

\defgcs{O}%
 {}%
 {\LaTeX\ Kernel}%
 {Upper case slashed-O \O\ character.}%
 {}

\defgcs{o}%
 {}%
 {\LaTeX\ Kernel}%
 {Lower case slashed-o \o\ character.}%
 {}

\defgcs{ss}%
 {}%
 {\LaTeX\ Kernel}%
 {Eszett \ss\ character.}%
 {}

\defgcs{SS}%
 {}%
 {\LaTeX\ Kernel}%
 {SS (upper case \ss).}%
 {}

\defgcs{left}
 {\meta{delimiter}}
 {\LaTeX\ Kernel (Math Mode)}
 {Indicates a left stretchable delimiter. Must have a matching
  \nxglsi{right}.}
 {%
   \BeginArgList
    \csentryargitem{delimiter} A delimiter symbol, such as
    \nxglsni{openparen}, or a delimiter command, such as \nxglsni{langle}.
   \EndArgList
 }

\defgcs{right}
 {\meta{delimiter}}
 {\LaTeX\ Kernel (Math Mode)}
 {Indicates a right stretchable delimiter. Must have a matching
  \nxglsi{left}.}
 {%
   \BeginArgList
    \csentryargitem{delimiter} A delimiter symbol, such as
    \nxglsni{closeparen}, or a delimiter command, such as \nxglsni{rangle}.
   \EndArgList
 }

\defgcs{bigl}%
 {\meta{delimiter}}%
 {\nxisty{amsmath} package (Math Mode)}%
 {Left delimiter sizing.}%
 {%
   \BeginArgList
    \csentryargitem{delimiter} A delimiter symbol, such as
    \nxglsni{openparen}, or a delimiter command, such as \nxglsni{langle}.
   \EndArgList
 }

\defgcs{bigr}%
 {\meta{delimiter}}%
 {\nxisty{amsmath} package (Math Mode)}%
 {Right delimiter sizing.}%
 {%
   \BeginArgList
    \csentryargitem{delimiter} A delimiter symbol, such as
    \nxglsni{closeparen}, or a delimiter command, such as \nxglsni{rangle}.
   \EndArgList
 }

\defgcs{Bigl}%
 {\meta{delimiter}}%
 {\nxisty{amsmath} package (Math Mode)}%
 {Left delimiter sizing.}%
 {%
   \BeginArgList
    \csentryargitem{delimiter} A delimiter symbol, such as
    \nxglsni{openparen}, or a delimiter command, such as \nxglsni{langle}.
   \EndArgList
 }

\defgcs{Bigr}%
 {\meta{delimiter}}%
 {\nxisty{amsmath} package (Math Mode)}%
 {Right delimiter sizing.}%
 {%
   \BeginArgList
    \csentryargitem{delimiter} A delimiter symbol, such as
    \nxglsni{closeparen}, or a delimiter command, such as \nxglsni{rangle}.
   \EndArgList
 }

\defgcs{biggl}%
 {\meta{delimiter}}%
 {\nxisty{amsmath} package (Math Mode)}%
 {Left delimiter sizing.}%
 {%
   \BeginArgList
    \csentryargitem{delimiter} A delimiter symbol, such as
    \nxglsni{openparen}, or a delimiter command, such as \nxglsni{langle}.
   \EndArgList
 }

\defgcs{biggr}%
 {\meta{delimiter}}%
 {\nxisty{amsmath} package (Math Mode)}%
 {Right delimiter sizing.}%
 {%
   \BeginArgList
    \csentryargitem{delimiter} A delimiter symbol, such as
    \nxglsni{closeparen}, or a delimiter command, such as \nxglsni{rangle}.
   \EndArgList
 }

\defgcs{Biggl}%
 {\meta{delimiter}}%
 {\nxisty{amsmath} package (Math Mode)}%
 {Left delimiter sizing.}%
 {%
   \BeginArgList
    \csentryargitem{delimiter} A delimiter symbol, such as
    \nxglsni{openparen}, or a delimiter command, such as \nxglsni{langle}.
   \EndArgList
 }

\defgcs{Biggr}%
 {\meta{delimiter}}%
 {\nxisty{amsmath} package (Math Mode)}%
 {Right delimiter sizing.}%
 {%
   \BeginArgList
    \csentryargitem{delimiter} A delimiter symbol, such as
    \nxglsni{closeparen}, or a delimiter command, such as \nxglsni{rangle}.
   \EndArgList
 }

\defgcs{pounds}%
 {}%
 {\LaTeX\ Kernel}%
 {Pound \pounds\ symbol.}%
 {}%


\defgcs{v}%
 {\marg{\meta{c}}}%
 {\LaTeX\ Kernel}%
 {Caron diacritic over \meta{c}. Example: 
  \cmdname{v}\marg{o} produces \ocaron.}%
 {%
   \BeginArgList
    \csentryargitem{c} The character that requires the caron
      diacritic.
   \EndArgList
 }

\defgcs{u}%
 {\marg{\meta{c}}}%
 {\LaTeX\ Kernel}%
 {Breve diacritic over \meta{c}. Example: 
  \cmdname{u}\marg{o} produces \obreve.%
 }%
 {%
   \BeginArgList
    \csentryargitem{c} The character that requires the breve
      diacritic.
   \EndArgList
 }

\defgcs{H}%
 {\marg{\meta{c}}}%
 {\LaTeX\ Kernel}%
 {Double acute diacritic over \meta{c}. Example: 
  \cmdname{H}\marg{o} produces \odoubleacute.%
 }%
 {%
   \BeginArgList
    \csentryargitem{c} The character that requires the double
      acute diacritic.
   \EndArgList
 }

\defgcs{t}%
 {\marg{\meta{characters}}}%
 {\LaTeX\ Kernel}%
 {Tie over \meta{characters}. Example: 
  \cmdname{t}\marg{xy} produces \xytie.%
 }%
 {%
   \BeginArgList
    \csentryargitem{characters} The characters that require the tie.
   \EndArgList
 }

\defgcs{c}%
 {\marg{\meta{c}}}%
 {\LaTeX\ Kernel}%
 {Cedilla under \meta{c}. Example: 
  \cmdname{c}\marg{o} produces \ocedilla.%
 }%
 {%
   \BeginArgList
    \csentryargitem{c} The character that requires the cedilla.
   \EndArgList
 }

\defgcs{d}%
 {\marg{\meta{c}}}%
 {\LaTeX\ Kernel}%
 {Dot under \meta{c}. Example: 
  \cmdname{d}\marg{o} produces \odotunder.%
 }%
 {%
   \BeginArgList
    \csentryargitem{c} The character that requires the dot
     under it.
   \EndArgList
 }

\defgcs{b}%
 {\marg{\meta{c}}}%
 {\LaTeX\ Kernel}%
 {Bar under \meta{c}. Example: 
  \cmdname{b}\marg{r} produces \rbarunder.%
 }%
 {%
   \BeginArgList
    \csentryargitem{c} The character that requires the bar
     under it.
   \EndArgList
 }

\defgcs{r}%
 {\marg{\meta{c}}}%
 {\LaTeX\ Kernel}%
 {Ring over \meta{c}. Example: 
  \cmdname{r}\marg{u} produces \uring.%
 }%
 {%
   \BeginArgList
    \csentryargitem{c} The character that requires the ring
      over it.
   \EndArgList
 }

\defgsymcs[acute]{\quoterightcs}%
 {\marg{\meta{c}}}%
 {\LaTeX\ Kernel}%
 {Acute accent over \meta{c}. Example:
   \cmdname{'}\marg{o} produces \protect\'{o}.%
 }%
 {%
   \BeginArgList
    \csentryargitem{c} The character that requires an acute
     accent over it.
   \EndArgList
 }

\defgsymcs[grave]{\quoteleftcs}%
 {\marg{\meta{c}}}%
 {\LaTeX\ Kernel}%
 {Grave accent over \meta{c}. Example:
   \cmdname{`}\marg{o} produces \protect\`{o}.}%
 {%
   \BeginArgList
    \csentryargitem{c} The character that requires a grave
     accent over it.
   \EndArgList
 }

\defgsymcs[period]{\periodsym}%
 {\marg{\meta{c}}}%
 {\LaTeX\ Kernel}%
 {Dot over \meta{c}. Example:
   \cmdname{.}\marg{o} produces \odotover.%
 }%
 {%
   \BeginArgList
    \csentryargitem{c} The character that requires a dot
     over it.
   \EndArgList
 }

\defgchar
 {periodchar}
 {\periodsym}
 {}
 {\LaTeX\ Kernel}
 {\nopostdesc}
 {}

\defgchildchar
 {sentence.periodchar}
 {periodchar}
 {\periodsym}
 {period (full stop) or decimal point}

\defgchildchar
 {delimiter.periodchar}
 {periodchar}
 {\periodsym}
 {invisible delimiter}

\defgsymcs[macron]{\equalsym}%
 {\marg{\meta{c}}}%
 {\LaTeX\ Kernel}%
 {Macron accent over \meta{c}. Example:
   \cmdname{=}\marg{o} produces \omacron.%
 }%
 {%
   \BeginArgList
    \csentryargitem{c} The character that requires a macron
     accent over it.
   \EndArgList
 }

\defgactivecharcs
 [\circumsym\space]
 {circum}
 {\textasciicircum}
 {\marg{\meta{c}}}
 {\LaTeX\ Kernel}
 {%
   Circumflex accent over \meta{c}. Example:
   \cmdname{\textasciicircum}\marg{o} produces \ocircum.%
 }
 {%
   \BeginArgList
    \csentryargitem{c} The character that requires a
     circumflex accent over it.
   \EndArgList
 }

\defgactivechar
 [\circumsym]
 {circumchar}
 {\textasciicircum}
 {\marg{\meta{maths}}}
 {\LaTeX\ Kernel (Math Mode)}
 {Displays its argument as a superscript.}
 {%
   \BeginArgList
     \csentryargitem{maths} The superscript.
   \EndArgList
 }

\defgactivecharcs
 [\tildesym\space]
 {tilde}
 {\textasciitilde}
 {\marg{\meta{c}}}
 {\LaTeX\ Kernel}
 {%
   Tilde accent over \meta{c}. Example:
   \cmdname{\textasciitilde}\marg{o} produces \otilde.%
 }
 {%
   \BeginArgList
    \csentryargitem{c} The character that requires a
     tilde accent over it.
   \EndArgList
 }

\defgactivechar
 [\tildesym]
 {tildechar}
 {\textasciitilde}
 {}
 {\LaTeX\ Kernel}
 {Unbreakable space.}
 {}

\defgidxactivecharcs
 {doublequote}
 {\doublequotesym}
 {\marg{\meta{c}}}
 {\LaTeX\ Kernel}%
 {%
   Umlaut over \meta{c}. Example:
   \quotecs\marg{o} produces \oumlaut.%
 }%
 {%
   \BeginArgList
    \csentryargitem{c} The character that requires an umlaut
     over it.
   \EndArgList
 }

\defgcs{item}%
 {\oarg{\meta{marker}}}%
 {\LaTeX\ Kernel}%
 {Specifies the start of an item in a list. (Only allowed inside one of
  the list making environments.)}%
 {%
   \BeginArgList
    \csentryargitem{marker} If specified, the default item marker is
     replaced with \meta{marker}.
   \EndArgList
 }

\defgsymcs[itcorr]{\slashsym}%
 {}%
 {\LaTeX\ Kernel}%
 {Italic correction.}%
 {}

\defgcs{textrm}%
 {\marg{\meta{text}}}%
 {\LaTeX\ Kernel}%
 {Renders \meta{text} in the predefined serif font. (Defaults to
  Computer Modern Roman.)}%
 {%
   \BeginArgList
     \csentryargitem{text} The text on which to apply the font change.
   \EndArgList
 }

\defgcs{textsf}%
 {\marg{\meta{text}}}%
 {\LaTeX\ Kernel}%
 {Renders \meta{text} in the predefined sans-serif font. (Defaults to
  Computer Modern Sans.)}%
 {%
   \BeginArgList
     \csentryargitem{text} The text on which to apply the font change.
   \EndArgList
 }

\defgcs{texttt}%
 {\marg{\meta{text}}}%
 {\LaTeX\ Kernel}%
 {Renders \meta{text} in the predefined monospaced font. (Defaults to
  Computer Modern Typewriter.)}%
 {%
   \BeginArgList
     \csentryargitem{text} The text on which to apply the font change.
   \EndArgList
 }

\defgcs{textmd}%
 {\marg{\meta{text}}}%
 {\LaTeX\ Kernel}%
 {Renders \meta{text} with a medium weight in the current font
  family.}%
 {%
   \BeginArgList
     \csentryargitem{text} The text on which to apply the font change.
   \EndArgList
 }

\defgcs{textbf}%
 {\marg{\meta{text}}}%
 {\LaTeX\ Kernel}%
 {Renders \meta{text} with a bold weight in the current font family, if it exists.}%
 {%
   \BeginArgList
     \csentryargitem{text} The text on which to apply the font change.
   \EndArgList
 }

\defgcs{textup}%
 {\marg{\meta{text}}}%
 {\LaTeX\ Kernel}%
 {Renders \meta{text} with the upright form of the current font family.}%
 {%
   \BeginArgList
     \csentryargitem{text} The text on which to apply the font change.
   \EndArgList
 }

\defgcs{textit}%
 {\marg{\meta{text}}}%
 {\LaTeX\ Kernel}%
 {Renders \meta{text} with the italic form of the current font family, if it exists.}%
 {%
   \BeginArgList
     \csentryargitem{text} The text on which to apply the font change.
   \EndArgList
 }

\defgcs{textsl}%
 {\marg{\meta{text}}}%
 {\LaTeX\ Kernel}%
 {Renders \meta{text} with the slanted form of the current font family, if it exists.}%
 {%
   \BeginArgList
     \csentryargitem{text} The text on which to apply the font change.
   \EndArgList
 }

\defgcs{textsc}%
 {\marg{\meta{text}}}%
 {\LaTeX\ Kernel}%
 {Renders \meta{text} with the small-caps form of the current font family, if it exists.}%
 {%
   \BeginArgList
     \csentryargitem{text} The text on which to apply the font change.
   \EndArgList
 }

\defgcs{emph}%
 {\marg{\meta{text}}}%
 {\LaTeX\ Kernel}%
 {Toggles the upright and italic\slash slanted rendering of \meta{text}.}%
 {%
   \BeginArgList
     \csentryargitem{text} The text on which to apply the font change.
   \EndArgList
 }

\defgcs{textnormal}%
 {\marg{\meta{text}}}%
 {\LaTeX\ Kernel}%
 {Renders \meta{text} in the default font style.}%
 {%
   \BeginArgList
     \csentryargitem{text} The text on which to apply the font change.
   \EndArgList
 }

\defgcs{ttdefault}%
 {}%
 {\LaTeX\ Kernel}%
 {The name of the default typewriter family as used by \nxglsni{ttfamily}. 
  Defaults to \texttt{cmtt} (Computer Modern Typewriter).}%
 {}

\defgcs{rmdefault}%
 {}%
 {\LaTeX\ Kernel}%
 {The name of the default serif family as used by \nxglsni{rmfamily}. 
  Defaults to \texttt{cmr} (Computer Modern Roman).}%
 {}

\defgcs{sfdefault}%
 {}%
 {\LaTeX\ Kernel}%
 {The name of the default sans-serif family as used by \nxglsni{sffamily}. 
  Defaults to \texttt{cmss} (Computer Modern Sans-serif).}%
 {}

\defgcs{familydefault}%
 {}%
 {\LaTeX\ Kernel}%
 {Specifies the default font family. Defaults to
   \nxglsi{rmdefault} but may be redefined by certain classes.}%
 {}

\defgcs{rmfamily}%
 {}%
 {\LaTeX\ Kernel}%
 {Switches to the predefined serif font. (Defaults to
  Computer Modern Roman.)}%
 {}

\defgcs{sffamily}%
 {}%
 {\LaTeX\ Kernel}%
 {Switches to the predefined sans-serif font. (Defaults to
  Computer Modern Sans.)}%
 {}

\defgcs{ttfamily}%
 {}%
 {\LaTeX\ Kernel}%
 {Switches to the predefined monospaced font. (Defaults to
  Computer Modern Typewriter.)}%
 {}

\defgcs{mdseries}%
 {}%
 {\LaTeX\ Kernel}%
 {Switches to the medium weight in the current font
  family.}%
 {}

\defgcs{bfseries}%
 {}%
 {\LaTeX\ Kernel}%
 {Switches to the bold weight in the current font
  family.}%
 {}


\defgcs{upshape}%
 {}%
 {\LaTeX\ Kernel}%
 {Switches to the upright form of the current font family.}%
 {}

\defgcs{itshape}%
 {}%
 {\LaTeX\ Kernel}%
 {Switches to the italic form of the current font family, if it exists.}%
 {}

\defgcs{slshape}%
 {}%
 {\LaTeX\ Kernel}%
 {Switches to the slanted form of the current font family, if it exists.}%
 {}

\defgcs{scshape}%
 {}%
 {\LaTeX\ Kernel}%
 {Switches to the small-caps form of the current font family, if it exists.}%
 {}

\defgcs{em}%
 {}%
 {\LaTeX\ Kernel}%
 {Toggles the upright and italic\slash slanted form of the current font family.}%
 {}

\defgcs{normalfont}%
 {}%
 {\LaTeX\ Kernel}%
 {Switches to the default font style.}%
 {}

\defgcs{tiny}%
 {}%
 {Most document classes}%
 {Switches to tiny sized text.}%
 {}

\defgcs{scriptsize}%
 {}%
 {Most document classes}%
 {Switches to sub- or superscript sized text.}%
 {}

\defgcs{footnotesize}%
 {}%
 {Most document classes}%
 {Switches to footnote sized text.}%
 {}

\defgcs{small}%
 {}%
 {Most document classes}%
 {Switches to small sized text.}%
 {}

\defgcs{normalsize}%
 {}%
 {\LaTeX\ Kernel}%
 {Switches to normal sized text.}%
 {%
 }

\defgcs{large}%
 {}%
 {Most document classes}%
 {Switches to large sized text.}%
 {%
 }

\defgcs{Large}%
 {}%
 {Most document classes}%
 {Switches to extra-large sized text.}%
 {%
 }

\defgcs{LARGE}%
 {}%
 {Most document classes}%
 {Switches to extra-extra-large sized text.}%
 {%
 }

\defgcs{huge}%
 {}%
 {Most document classes}%
 {Switches to huge sized text.}%
 {%
 }

\defgcs{Huge}%
 {}%
 {Most document classes}%
 {Switches to extra-huge sized text.}%
 {%
 }

\defgcs{multicolumn}%
 {\marg{\meta{cols spanned}}\marg{\meta{col specifier}}\marg{\meta{text}}}%
 {\LaTeX\ Kernel}%
 {Spans multiple columns in a tabular-style environment.}%
 {%
   \BeginArgList
    \csentryargitem{cols spanned} The number of columns to span.
    \csentryargitem{col specifier} The alignment of this column-spanning entry.
    \csentryargitem{text} The contents of this column-spanning entry.
   \EndArgList
 }

\defgcs{toprule}%
 {\oarg{\meta{wd}}}%
 {\nxisty{booktabs} package}%
 {Horizontal rule for the top of a \nxglsi{env-tabular} environment.}%
 {%
   \BeginArgList
    \csentryargitem{wd} Thickness of the rule (dimension).
   \EndArgList
 }

\defgcs{bottomrule}%
 {\oarg{\meta{wd}}}%
 {\nxisty{booktabs} package}%
 {Horizontal rule for the bottom of a \nxglsi{env-tabular} environment.}%
 {%
   \BeginArgList
    \csentryargitem{wd} Thickness of the rule (dimension).
   \EndArgList
 }

\defgcs{midrule}%
 {\oarg{\meta{wd}}}%
 {\nxisty{booktabs} package}%
 {Horizontal rule to go below headings row of a \nxglsi{env-tabular} environment.}%
 {%
   \BeginArgList
    \csentryargitem{wd} Thickness of the rule (dimension).
   \EndArgList
 }

\defgcs{cmidrule}%
 {\oarg{\meta{wd}}\parg{\meta{trim}}\marg{\meta{n}-\meta{m}}}%
 {\nxisty{booktabs} package}%
 {Like \nxglsi{midrule} but only spans column \meta{n} to column
  \meta{m}.}%
 {%
   \BeginArgList
    \csentryargitem{wd} Thickness of the rule (dimension).
    \csentryargitem{trim} Trimming commands. Specifications:
     \texttt{r}, \texttt{r}\marg{\meta{wd}}, \texttt{l} and
      \texttt{l}\marg{\meta{wd}}
    \csentryargitem{n} Index of start column.
    \csentryargitem{m} Index of end column.
   \EndArgList
 }

\defgcs{heavyrulewidth}%
 {}%
 {\nxisty{booktabs} package}%
 {\nxGls{length} register specifying the thickness of \nxglsi{toprule} and
 \nxglsi{bottomrule}.}%
 {}

\defgcs{lightrulewidth}%
 {}%
 {\nxisty{booktabs} package}%
 {\nxGls{length} register specifying the thickness of \nxglsi{midrule}.}%
 {}

\defgcs{addlinespace}%
 {\oarg{\meta{wd}}}%
 {\nxisty{booktabs} package}%
 {Adds extra vertical space (\meta{wd} if specified, otherwise
  \nxglsi{defaultaddspace}) in a row of 
  a \nxglsi{env-tabular} environment. (Must be inserted after
  \nxglsi{dbbackslashchar} marker.)}%
 {%
   \BeginArgList
    \csentryargitem{wd} Height of space (dimension). Default is
     \nxglsi{defaultaddspace}.
   \EndArgList
 }

\defgcs{defaultaddspace}%
 {}%
 {\nxisty{booktabs} package}%
 {Default extra space before or after an adjacent rule.}%
 {}

\defgcs{parbox}%
 {\oarg{\meta{pos}}\oarg{\meta{height}}\marg{\meta{width}}\marg{\meta{text}}}%
 {\LaTeX\ Kernel}%
 {Makes a box with line-wrapped contents. (More restrictive than
  \nxglsi{env-minipage}.)}%
 {%
   \BeginArgList
    \csentryargitem{pos} The vertical alignment of the box relative to
     the surrounding text. (Centred if omitted.)
    \csentryargitem{height} The height of the box. (If omitted the
     height is the natural height of the contents of the box.)
    \csentryargitem{width} The width of the box.
    \csentryargitem{text} The contents of the box.
   \EndArgList
 }

\defgcs{author}%
 {\marg{\meta{name}}}%
 {Most classes that have the concept of a title page}%
 {Specifies the document author (or authors). This command doesn't
  display any text so may be used in the preamble, but if it's not
  in the preamble it must be placed before \nxglsi{maketitle}.}%
 {%
   \BeginArgList
    \csentryargitem{name} The name (or names) of the document author
     (or authors).
   \EndArgList
   Note that some classes, such as those supplied by journals or
   conference proceedings, may also define an optional argument
   that can be used to specify an abbreviated author list for the
   page headers.
 }

\defgcs{title}%
 {\marg{\meta{text}}}%
 {Most classes that have the concept of a title page}%
 {Specifies the document title. This command doesn't
  display any text so may be used in the preamble, but if it's not
  in the preamble it must be placed before \nxglsi{maketitle}.}%
 {%
   \BeginArgList
    \csentryargitem{text} The title of the document.
   \EndArgList
   Note that some classes, such as those supplied by journals or
   conference proceedings, may also define an optional argument
   that can be used to specify an abbreviated title for the page
   headers.
 }

\defgcs{date}%
 {\marg{\meta{text}}}%
 {Most classes that have the concept of a title page}%
 {Specifies the document date. This command doesn't
  display any text so may be used in the preamble, but if it's not
  in the preamble it must be placed before \nxglsi{maketitle}. If omitted, most
  classes assume the current date (as provided by \nxglsi{today}).}%
 {%
   \BeginArgList
    \csentryargitem{text} The document date.
   \EndArgList
 }

\defgcs{titlehead}%
 {\marg{\meta{text}}}%
 {\nxicls{scrartcl}, \nxicls{scrreprt}, \nxicls{scrbook} classes}%
 {Specifies the title header (typeset at the top of the title page).}%
 {%
   \BeginArgList
    \csentryargitem{text} The title header text.
   \EndArgList
 }

\defgcs{subtitle}%
 {\marg{\meta{text}}}%
 {\nxicls{scrartcl}, \nxicls{scrreprt}, \nxicls{scrbook} classes}%
 {Specifies the subtitle (typeset just below the title).}%
 {%
   \BeginArgList
    \csentryargitem{text} The subtitle text.
   \EndArgList
 }

\defgcs{subject}%
 {\marg{\meta{text}}}%
 {\nxicls{scrartcl}, \nxicls{scrreprt}, \nxicls{scrbook} classes}%
 {Specifies the subject (typeset just above the title).}%
 {%
   \BeginArgList
    \csentryargitem{text} The subject.
   \EndArgList
 }

\defgcs{publishers}%
 {\marg{\meta{text}}}%
 {\nxicls{scrartcl}, \nxicls{scrreprt}, \nxicls{scrbook} classes}%
 {Specifies the publisher (typeset after all the other titling
   information).}%
 {%
   \BeginArgList
    \csentryargitem{text} The publisher text.
   \EndArgList
 }

\defgcs{maketitle}%
 {}%
 {Most classes that have the concept of a title page}%
 {Generates the title page (or title block). This command is usually
  placed at the beginning of the document environment.}%
 {}

\defgcs{thanks}%
 {\marg{\meta{text}}}%
 {Most classes that have the concept of a title page}%
 {Inserts a special type of footnote in one of the titling fields,
  such as \nxglsi{author} or \nxglsni{title}. Usually used for some form of
  acknowledgement or affiliation.}%
 {}

\defgcs{today}%
 {}%
 {Most of the commonly-used classes}%
 {Inserts into the output file the date when the \LaTeX\ 
  application created it from the source code.}%
 {}

\defgcs{part}%
 {\oarg{\meta{short title}}\marg{\meta{title}}}%
 {Most classes that have the concept of document structure}%
 {Inserts a part sectional unit. \textbf{This command has a \htmlref{moving
  argument}{sec:fragile}.}}%
 {%
   \BeginArgList
     \csentryargitem{short title} An abbreviated form of the title to
       go in the table of contents or the page header.
     \csentryargitem{title} The title.
   \EndArgList
   The starred form of this command doesn't have an optional
   argument and doesn't increment or display the part counter.
 }

\defgcs{chapter}%
 {\oarg{\meta{short title}}\marg{\meta{title}}}%
 {Book-style classes (such as \nxicls{scrbook} or \nxicls{scrreprt}) that have 
  the concept of chapters}%
 {Inserts a chapter heading. \textbf{This command has a \htmlref{moving
  argument}{sec:fragile}.}}%
 {%
   \BeginArgList
     \csentryargitem{short title} An abbreviated form of the title to
       go in the table of contents or the page header.
     \csentryargitem{title} The title.
   \EndArgList
   The starred form of this command doesn't have an optional
   argument and doesn't increment or display the chapter counter.
 }

\defgcs{section}%
 {\oarg{\meta{short title}}\marg{\meta{title}}}%
 {Most classes that have the concept of document structure}%
 {Inserts a section header. \textbf{This command has a \htmlref{moving
  argument}{sec:fragile}.}}%
 {%
   \BeginArgList
     \csentryargitem{short title} An abbreviated form of the title to
       go in the table of contents or the page header.
     \csentryargitem{title} The title.
   \EndArgList
   The starred form of this command doesn't have an optional
   argument and doesn't increment or display the section counter.
 }

\defgcs{subsection}%
 {\oarg{\meta{short title}}\marg{\meta{title}}}%
 {Most classes that have the concept of document structure}%
 {Inserts a subsection header. \textbf{This command has a \htmlref{moving
  argument}{sec:fragile}.}}%
 {%
   \BeginArgList
     \csentryargitem{short title} An abbreviated form of the title to
       go in the table of contents or the page header.
     \csentryargitem{title} The title.
   \EndArgList
   The starred form of this command doesn't have an optional
   argument and doesn't increment or display the subsection counter.
 }

\defgcs{subsubsection}%
 {\oarg{\meta{short title}}\marg{\meta{title}}}%
 {Most classes that have the concept of document structure}%
 {Inserts a subsubsection header. \textbf{This command has a \htmlref{moving
  argument}{sec:fragile}.}}%
 {%
   \BeginArgList
     \csentryargitem{short title} An abbreviated form of the title to
       go in the table of contents or the page header.
     \csentryargitem{title} The title.
   \EndArgList
   The starred form of this command doesn't have an optional
   argument and doesn't increment or display the subsubsection counter.
 }

\defgcs{paragraph}%
 {\oarg{\meta{short title}}\marg{\meta{title}}}%
 {Most classes that have the concept of document structure}%
 {Inserts a subsubsubsection header. Most classes default to an 
  unnumbered running header for
  this sectional unit. \textbf{This command has a \htmlref{moving
  argument}{sec:fragile}.}}%
 {%
   \BeginArgList
     \csentryargitem{short title} An abbreviated form of the title to
       go in the table of contents or the page header.
     \csentryargitem{title} The title.
   \EndArgList
   The starred form of this command doesn't have an optional
   argument and doesn't increment or display the associated counter.
 }

\defgcs{subparagraph}%
 {\oarg{\meta{short title}}\marg{\meta{title}}}%
 {Most classes that have the concept of document structure}%
 {Inserts a subsubsubsubsection header. Most classes default 
  to an unnumbered running header for
  this sectional unit. \textbf{This command has a \htmlref{moving
  argument}{sec:fragile}.}}%
 {%
   \BeginArgList
     \csentryargitem{short title} An abbreviated form of the title to
       go in the table of contents or the page header.
     \csentryargitem{title} The title.
   \EndArgList
   The starred form of this command doesn't have an optional
   argument and doesn't increment or display the associated counter.
 }

\defgcs{minisec}%
 {\marg{\meta{heading}}}%
 {\nxicls{scrartcl}, \nxicls{scrreprt} and
  \nxicls{scrbook} classes}%
 {An unnumbered heading not associated with any structuring level.}%
 {%
   \BeginArgList
    \csentryargitem{heading} The heading text.
   \EndArgList
 }

\defgcs{addtokomafont}%
 {\marg{\meta{element name}}\marg{\meta{commands}}}%
 {\nxicls{scrartcl}, \nxicls{scrreprt} and
  \nxicls{scrbook} classes}%
 {Sets the font characteristics for the given \koma\ element.}%
 {%
   \BeginArgList
     \csentryargitem{element name} The element's name, for example
      \texttt{chapter}. See the \koma\ manual for a full list
      of defined elements.
     \csentryargitem{commands} The font changing commands to apply
      to that element.
   \EndArgList
 }

\defgcs{appendix}%
 {}%
 {Most classes that have the concept of document structure}%
 {Indicates (but doesn't print anything) that the document is
  switching to the appendices. If chapters exist, the chapter
  numbering is reset and switched to a different format
  (usually upper case letters) otherwise the section numbering
  is reset and switched to a different format.}%
 {}

\defgcs{tableofcontents}
 {}%
 {Most classes that have the concept of document structure}%
 {Inserts the table of contents. A second (possibly third) run
  is required to ensure the page numbering is correct.}%
 {}

\defgcs{addcontentsline}%
 {\marg{\meta{toc}}\marg{\meta{section unit}}\marg{\meta{text}}}%
 {\LaTeX\ Kernel}%
 {Adds a sectional unit header to the contents list.}%
 {%
   \BeginArgList
    \csentryargitem{toc} The extension of the required external file
    without the dot (for example, \texttt{toc} for the table of
     contents or \texttt{lof} for the list of figures).
    \csentryargitem{section unit} The name of the sectional unit (for
     example, \texttt{chapter}).
    \csentryargitem{text} The text to be added to the contents list.
   \EndArgList
 }

\defgcs{label}%
 {\marg{\meta{string}}}%
 {\LaTeX\ Kernel}%
 {Assigns a unique textual label linked to the most recently
  incremented cross-referencing counter in the current \nxglslink{group}{scope} (see
  also \nxglsi{ref}).}%
 {%
   \BeginArgList
    \csentryargitem{string} A unique label that can be referenced
     elsewhere in the document with \nxglsi{ref}. (It's best to 
     just use alphanumerics and non-active punctuation characters in
     the label.)
   \EndArgList
 }

\defgcs{ref}%
 {\marg{\meta{string}}}%
 {\LaTeX\ Kernel}%
 {References the value of the counter linked to the given label.
  A second (possibly third) run of \LaTeX\ is required to ensure the cross-references
  are up-to-date.
 }%
 {%
   \BeginArgList
    \csentryargitem{string} The required label that was used in the
     corresponding \nxglsi{label}.
   \EndArgList
 }

\defgcs{pageref}%
 {\marg{\meta{string}}}%
 {\LaTeX\ Kernel}%
 {Similar to \nxglsi{ref} but inserts the page number where the given
  label was defined.
  A second (possibly third) run of \LaTeX\ is required to ensure the cross-references
  are up-to-date.
 }%
 {%
   \BeginArgList
    \csentryargitem{string} The required label that was used in the
     corresponding \nxglsi{label}.
   \EndArgList
 }

\defgcs{vref}%
 {\marg{\meta{string}}}%
 {\nxisty{varioref} package}%
 {Like \nxglsi{ref} but also adds information about the location, such
  as \dq{on page~\meta{n}} or \dq{on the following page}.
 }%
 {%
   \BeginArgList
    \csentryargitem{string} The required label that was used in the
     corresponding \nxglsi{label}.
   \EndArgList
 }

\defgcs{bibitem}%
 {\oarg{\meta{tag}}\marg{\meta{key}}}%
 {\LaTeX\ Kernel}%
 {Indicates the start of a new reference in the bibliography. May
  only be used inside the contents of \nxglsi{env-thebibliography} environment}%
 {%
   \BeginArgList
    \csentryargitem{tag} If present, overrides the marker at the
    start of the reference.
    \csentryargitem{key} A unique key that identifies this reference so
    it can be cited elsewhere in the document using \nxglsi{cite}.
   \EndArgList
 }

\defgcs{cite}%
 {\oarg{\meta{text}}\marg{\meta{key list}}}%
 {\LaTeX\ Kernel}%
 {Inserts the citation markers of each reference identified in the
  key list. A second run is required to ensure the reference is
  correct.}%
 {%
    \BeginArgList
      \csentryargitem{text} Additional text to insert into the citation
       (such as the chapter number, or a particular page or
       page range within the citation).
      \csentryargitem{key list} A comma-separated list of keys used in
      the corresponding \nxglsi{cite}.
    \EndArgList
 }

\defgcs{pagenumbering}%
 {\marg{\meta{style}}}%
 {\LaTeX\ Kernel}%
 {Sets the style of the page numbers.}%
 {%
   \BeginArgList
     \csentryargitem{style} The page numbering style (e.g.\
      \texttt{roman} for lower case Roman numerals).
   \EndArgList
 }

\defgcs{pagestyle}%
 {\marg{\meta{style}}}%
 {\LaTeX\ Kernel}%
 {Sets the style of the headers and footers.}%
 {%
   \BeginArgList
     \csentryargitem{style} The name of the page style. The \LaTeX\
     kernel defines only two styles: \nxipagestyle{empty} and
     \nxipagestyle{plain}. Most of the standard classes also provide
     the \nxipagestyle{headings} style.
   \EndArgList
 }

\defgcs{thispagestyle}%
 {\marg{\meta{style}}}%
 {\LaTeX\ Kernel}%
 {Like \nxglsi{pagestyle} but only affects the current page.}%
 {%
   \BeginArgList
     \csentryargitem{style} The name of the page style.
   \EndArgList
 }

\defgcs{markboth}%
 {\marg{\meta{left head}}\marg{\meta{right head}}}%
 {\LaTeX\ Kernel}%
 {Specifies information for the left and right page headers. Not all
   page styles use this information, in which case the arguments are
   ignored.}%
 {%
   \BeginArgList
     \csentryargitem{left head} Text for the left (even) page header,
     if the document has been identified as a double-sided document
     (usually with the \nxclsopt{twoside} class option).

     \csentryargitem{right head} Text for the right page header,
     if the \nxclsopt{twoside} option has been used. Otherwise used
     for both odd and even page headers.
   \EndArgList
 }

\defgcs{markright}%
 {\marg{\meta{right head}}}%
 {\LaTeX\ Kernel}%
 {Specifies information for the right (odd) page header. Not all
  page styles use this information, in which case the argument is
  ignored.}%
 {%
   \BeginArgList
     \csentryargitem{right head} Text for the right page header,
     if the \nxclsopt{twoside} option has been used. Otherwise used
     for both odd and even page headers.
   \EndArgList
 }

\defgcs{includegraphics}%
 {\oarg{\meta{key vals}}\marg{\meta{filename}}}%
 {\nxisty{graphicx} package}%
 {Inserts a graphics file into the document.}%
 {%
   \BeginArgList
     \csentryargitem{key vals} Comma-separated list of options that can
      be used to change the way the image is displayed.
   \EndArgList
 }

\defgcs{DeclareGraphicsExtensions}%
 {\marg{\meta{ext-list}}}%
 {\nxisty{graphicx} package}%
 {Specify the file extensions to look for if no extension is used in
 \nxglsi{includegraphics}}%
 {%
   \BeginArgList
     \csentryargitem{ext-list} A comma-separated list of extensions.
   \EndArgList
 }

\defgcs{selectlanguage}
 {\marg{\meta{language name}}}%
 {\nxisty{babel} package}%
 {Switches to the named language. Predefined textual elements, such
  as the date given by \nxglsi{today} or prefixes like \dq{Chapter},
  are redefined to those supplied by the given language.}%
 {%
   \BeginArgList
    \csentryargitem{language} The language identifier.
   \EndArgList
 }

\defgcs{foreignlanguage}%
 {\marg{\meta{language name}}\marg{\meta{text}}}%
 {\nxisty{babel} package}%
 {Typesets the given text using any predefined names or date formats supplied by the
  given language.
 }%
 {%
   \BeginArgList
    \csentryargitem{language} The language identifier.
    \csentryargitem{text} The foreign text.
   \EndArgList
 }

\defgcs{iflanguage}%
 {\marg{\meta{language name}}\marg{\meta{true text}}\marg{\meta{false
text}}}%
 {\nxisty{babel} package}%
 {Tests the current language.}%
 {%
   \BeginArgList
    \csentryargitem{language} The language identifier.
    \csentryargitem{true text} The text to typeset if true.
    \csentryargitem{false text} The text to typeset if not true.
   \EndArgList
 }

\defgcs{caption}%
 {\oarg{\meta{short caption}}\marg{\meta{caption text}}}%
 {\LaTeX\ Kernel}%
 {Inserts the caption for a float such as a figure or table. 
  \textbf{This command has a \htmlref{moving argument}{sec:fragile}.}}%
 {%
   \BeginArgList
    \csentryargitem{short caption} If provided, an abbreviated caption
      to go in the list of figures\slash tables etc.
    \csentryargitem{caption text} The caption contents.
   \EndArgList
 }

\defgcs{listoffigures}%
 {}%
 {Most classes that have the concept of document structure}%
 {Inserts the list of figures. A second (possibly third) run
  is required to ensure the page numbering is correct.}%
 {}

\defgcs{listoftables}%
 {}%
 {Most classes that have the concept of document structure}%
 {Inserts the list of tables. A second (possibly third) run
  is required to ensure the page numbering is correct.}%
 {}

\defgcs{hspace}%
 {\marg{\meta{length}}}%
 {\LaTeX\ Kernel}%
 {Inserts a horizontal gap of the given width.}%
 {%
   \BeginArgList
    \csentryargitem{length} The width of the horizontal gap.
   \EndArgList
 }

\defgcs{vspace}%
 {\marg{\meta{length}}}%
 {\LaTeX\ Kernel}%
 {Inserts a vertical gap of the given height.}%
 {%
   \BeginArgList
    \csentryargitem{length} The \protect\htmlref{height}{sec:length} of the vertical gap.
   \EndArgList
 }

\defgcs{baselineskip}%
 {}%
 {\LaTeX\ Kernel}%
 {A \nxglsi{length} register that stores the current interline spacing.
  This is recalculated whenever the font changes.}%
 {}

\defgcs{parskip}%
 {}%
 {\LaTeX\ Kernel}%
 {A \nxglsi{length} register that stores the spacing between paragraphs.
  (If you're using one of the \koma\ classes, use the \nxscrclsopt{parskip} 
   option to set it to full or half line height.)}%
 {}

\defgcs{twocolumn}%
 {\oarg{\meta{header text}}}%
 {\LaTeX\ Kernel}%
 {Starts a new page and switches to two column mode.}%
 {%
   \BeginArgList
    \csentryargitem{header text} If present, this header text will be
    placed at the top of the page spanning both columns. (For
    example, as used by \nxglsi{maketitle} in two column articles.)
   \EndArgList
 }

\defgcs{onecolumn}%
 {}%
 {\LaTeX\ Kernel}%
 {Starts a new page and switches to one column mode.}%
 {}

\defgcs{TeX}%
 {}%
 {\LaTeX\ Kernel}%
 {Typesets the \TeX\ logo.}%
 {}

\defgcs{LaTeX}%
 {}%
 {\LaTeX\ Kernel}%
 {Typesets the \makeimg{LaTeX}{\LaTeX}\ logo.}%
 {}

\defgcs{LaTeXe}%
 {}%
 {\LaTeX\ Kernel}%
 {Typesets the \makeimg{LaTeX2e}{\LaTeXe}\ logo.}%
 {}

\defgcs{framebox}%
 {\oarg{\meta{width}}\oarg{\meta{align}}\marg{\meta{text}}}%
 {\LaTeX\ Kernel}%
 {Puts a frame around its contents, prohibiting a line break in the
  contents.}%
 {%
   \BeginArgList
     \csentryargitem{width} The width of the box. (If omitted, this value is
      the natural width of the box contents.)
     \csentryargitem{align} Horizontal alignment of box contents
      (\texttt{c} if omitted).
     \csentryargitem{text} The contents of the box.
   \EndArgList
 }

\defgcs{fbox}%
 {\marg{\meta{text}}}%
 {\LaTeX\ Kernel}%
 {Puts a frame around its contents, prohibiting a line break in the
  contents.}%
 {%
   \BeginArgList
     \csentryargitem{text} The contents of the box.
   \EndArgList
 }

\defgcs{shadowbox}%
 {\marg{\meta{text}}}%
 {\nxisty{fancybox} package}%
 {Puts a shadow frame around its contents, prohibiting a line break in the
  contents.}%
 {%
   \BeginArgList
     \csentryargitem{text} The contents of the box.
   \EndArgList
 }

\defgcs{doublebox}%
 {\marg{\meta{text}}}%
 {\nxisty{fancybox} package}%
 {Puts a double-lined frame around its contents, prohibiting a line break in the
  contents.}%
 {%
   \BeginArgList
     \csentryargitem{text} The contents of the box.
   \EndArgList
 }

\defgcs{ovalbox}%
 {\marg{\meta{text}}}%
 {\nxisty{fancybox} package}%
 {Puts a thin-lined oval frame around its contents, prohibiting a line break in the
  contents.}%
 {%
   \BeginArgList
     \csentryargitem{text} The contents of the box.
   \EndArgList
 }

\defgcs{Ovalbox}%
 {\marg{\meta{text}}}%
 {\nxisty{fancybox} package}%
 {Puts a thick-lined oval frame around its contents, prohibiting a line break in the
  contents.}%
 {%
   \BeginArgList
     \csentryargitem{text} The contents of the box.
   \EndArgList
 }

\defgcs{protect}%
 {\meta{command}}%
 {\LaTeX\ Kernel}%
 {Used in a moving argument to prevent a fragile command from expanding.}%
 {%
   \BeginArgList
    \csentryargitem{command} The fragile command that needs protecting.
   \EndArgList
 }

\defgcs{par}%
 {}%
 {\LaTeX\ Kernel}%
 {Insert a paragraph break.}%
 {}

\defgcs{linewidth}%
 {}%
 {\LaTeX\ Kernel}%
 {A length containing the desired current line width. This is
  usually the width of the typeblock, but inside a \nxglsi{env-minipage}
  or \nxglsi{parbox} it will be the width the box. Note that
  the actual contents of the line may fall short of the line width
  (underfull hbox) or extend beyond it (overfull hbox).}%
 {}

\defgcs{textwidth}%
 {}%
 {\LaTeX\ Kernel}%
 {A length containing the width of the typeblock. Note that
  the actual contents of the line may fall short of the line width
  (underfull hbox) or extend beyond it (overfull hbox). This width
  does not include the area for marginal notes.}%
 {}

\defgcs{textheight}%
 {}%
 {\LaTeX\ Kernel}%
 {A length containing the height of the typeblock. Note that
  the actual contents of the page may fall short of the text height
  (underfull vbox) or extend beyond it (overfull vbox). This
  measurement does not include the header and footer areas.}%
 {}

\defgcs{and}%
 {}%
 {\LaTeX\ Kernel}%
 {Used to separate authors in \nxglsi{author}}%
 {}

\defgcs{clearpage}%
 {}%
 {\LaTeX\ Kernel}%
 {Inserts a page break and processes any unprocessed floats}%
 {}

\defgcs{rotatebox}%
 {\oarg{\meta{option list}}\marg{\meta{angle}}\marg{\meta{text}}}%
 {\nxisty{graphicx} package}%
 {Rotates the given contents by the given angle.}%
 {%
   \BeginArgList
     \csentryargitem{option list} A comma-separated list of options:
      \protect\begin{itemize}
       \protect\item \nxikeyvalopt{rotatebox}{origin}=\meta{label}
       \protect\item \nxikeyvalopt{rotatebox}{x}=\meta{dimen}
       \protect\item \nxikeyvalopt{rotatebox}{y}=\meta{dimen}
       \protect\item \nxikeyvalopt{rotatebox}{units}=\meta{number}
      \protect\end{itemize}

     \csentryargitem{angle} The angle of rotation.

     \csentryargitem{text} The text to be rotated.
   \EndArgList
 }

\defgcs{scalebox}%
 {\marg{\meta{h scale}}\oarg{\meta{v scale}}\marg{\meta{text}}}%
 {\nxisty{graphicx} package}%
 {Scales the specified contents.}%
 {%
    \BeginArgList
      \csentryargitem{h scale} The horizontal scale factor.
      \csentryargitem{v scale} The vertical scale factor. If omitted,
       the same as \meta{h scale}.
      \csentryargitem{text} The text to be scaled.
    \EndArgList
 }

\defgcs{reflectbox}%
 {\marg{\meta{text}}}%
 {\nxisty{graphicx} package}%
 {Reflects the specified contents in the \makeimg{y}{\ensuremath{y}}-axis.)}%
 {%
    \BeginArgList
      \csentryargitem{text} The text to be reflected.
    \EndArgList
 }

\defgcs{resizebox}%
 {\marg{\meta{h length}}\marg{\meta{v length}}\marg{\meta{text}}}%
 {\nxisty{graphicx} package}%
 {Scales the specified contents to the given dimensions.}%
 {%
   \BeginArgList
      \csentryargitem{h length} The required width or \texttt{!} to keep
      the aspect ratio.
      \csentryargitem{v length} The required height or \texttt{!} to keep
      the aspect ratio.
      \csentryargitem{text} The text to be scaled.
   \EndArgList
 }

\defgcs{centering}%
 {}%
 {\LaTeX\ Kernel}%
 {Switches the paragraph alignment to centred.}%
 {}

\defgcs{newcommand}%
 {\marg{\meta{cmd}}\oarg{\meta{n-args}}\oarg{\meta{default}}\marg{\meta{text}}}%
 {\LaTeX\ Kernel}%
 {Defines a new command.}%
 {%
   \BeginArgList
    \csentryargitem{cmd} The new command name (including initial backslash).
    \csentryargitem{n-args} The number of arguments this new command
     should have.
    \csentryargitem{default} If the first argument should be optional,
     the default value if omitted.
    \csentryargitem{text} What actions the command should perform.
   \EndArgList
 }

\defgcs{renewcommand}%
 {\marg{\meta{cmd}}\oarg{\meta{n-args}}\oarg{\meta{default}}\marg{\meta{text}}}%
 {\LaTeX\ Kernel}%
 {Redefines an existing command.}%
 {%
   \BeginArgList
    \csentryargitem{cmd} The command name (including initial backslash).
    \csentryargitem{n-args} The number of arguments this command
     should have.
    \csentryargitem{default} If the first argument should be optional,
     the default value if omitted.
    \csentryargitem{text} What actions the command should perform.
   \EndArgList
 }

\defgcs{newenvironment}%
 {\marg{\meta{env-name}}\oarg{\meta{n-args}}\oarg{\meta{default}}\marg{\meta{begin-code}}\marg{\meta{end-code}}}%
 {\LaTeX\ Kernel}%
 {Defines a new environment.}%
 {%
   \BeginArgList
    \csentryargitem{env-name} The new environment name (\emph{no} backslash).
    \csentryargitem{n-args} The number of arguments this new
     environment should have.
    \csentryargitem{default} If the first argument should be optional,
     the default value if omitted.
    \csentryargitem{begin-code} Actions to perform at the start of the
      environment.
    \csentryargitem{end-code} Actions to perform at the end of the
      environment.
   \EndArgList
 }

\defgcs{renewenvironment}%
 {\marg{\meta{env-name}}\oarg{\meta{n-args}}\oarg{\meta{default}}\marg{\meta{begin-code}}\marg{end-code}}%
 {\LaTeX\ Kernel}%
 {Redefines an existing environment.}%
 {%
   \BeginArgList
    \csentryargitem{env-name} The environment name (\emph{no} backslash).
    \csentryargitem{n-args} The number of arguments this new
     environment should have.
    \csentryargitem{default} If the first argument should be optional,
     the default value if omitted.
    \csentryargitem{begin-code} Actions to perform at the start of the
      environment.
    \csentryargitem{end-code} Actions to perform at the end of the
      environment.
   \EndArgList
 }

\defgcs{labelitemi}%
 {}%
 {Classes that define the \nxglsni{env-itemize} environment}%
 {The default label for the first level \nxglsni{env-itemize}.}%
 {}

\defgcs{labelitemii}%
 {}%
 {Classes that define the \nxglsni{env-itemize} environment}%
 {The default label for the second level \nxglsni{env-itemize}.}%
 {}

\defgcs{labelitemiii}%
 {}%
 {Classes that define the \nxglsni{env-itemize} environment}%
 {The default label for the third level \nxglsni{env-itemize}.}%
 {}

\defgcs{labelitemiv}%
 {}%
 {Classes that define the \nxglsni{env-itemize} environment}%
 {The default label for the fourth level \nxglsni{env-itemize}.}%
 {}

\defgcs{mbox}%
 {\marg{\meta{text}}}%
 {\LaTeX\ Kernel}%
 {Ensures that the given text doesn't contain a line break.}%
 {}

\defgcs{text}%
 {\marg{\meta{text}}}%
 {\nxisty{amsmath} package (Math Mode)}%
 {Displays its argument in the normal text font (as opposed to the
  current maths font).}%
 {%
   \BeginArgList
     \csentryargitem{text} The text to be displayed in the normal text font.
   \EndArgList
 }

\defgcs{intertext}%
 {\marg{\meta{text}}}%
 {\nxisty{amsmath} package (Math Mode)}%
 {Used for a short interjection in the middle of a multi-line
  displayed maths, such as in an \nxglsi{env-align} environment.
  May only appear right after \nxglsi{dbbackslashchar}.}%
 {%
   \BeginArgList
     \csentryargitem{text} The interjection text.
   \EndArgList
 }

\defgcs{overleftarrow}%
 {\marg{\meta{maths}}}%
 {\LaTeX\ Kernel (Math Mode)}%
 {Puts an extendible left arrow over \meta{maths}}%
 {%
   \BeginArgList
     \csentryargitem{maths} The maths over which to put the arrow.
   \EndArgList
 }

\defgcs{overrightarrow}%
 {\marg{\meta{maths}}}%
 {\LaTeX\ Kernel (Math Mode)}%
 {Puts an extendible right arrow over \meta{maths}}%
 {%
   \BeginArgList
     \csentryargitem{maths} The maths over which to put the arrow.
   \EndArgList
 }

\defgcs{overleftrightarrow}%
 {\marg{\meta{maths}}}%
 {\nxisty{amsmath} package (Math Mode)}%
 {Puts an extendible left-right arrow over \meta{maths}}%
 {%
   \BeginArgList
     \csentryargitem{maths} The maths over which to put the arrow.
   \EndArgList
 }

\defgcs{underleftarrow}%
 {\marg{\meta{maths}}}%
 {\nxisty{amsmath} package (Math Mode)}%
 {Puts an extendible left arrow under \meta{maths}}%
 {%
   \BeginArgList
     \csentryargitem{maths} The maths under which to put the arrow.
   \EndArgList
 }

\defgcs{underrightarrow}%
 {\marg{\meta{maths}}}%
 {\nxisty{amsmath} package (Math Mode)}%
 {Puts an extendible right arrow under \meta{maths}}%
 {%
   \BeginArgList
     \csentryargitem{maths} The maths under which to put the arrow.
   \EndArgList
 }

\defgcs{underleftrightarrow}%
 {\marg{\meta{maths}}}%
 {\nxisty{amsmath} package (Math Mode)}%
 {Puts an extendible left-right arrow under \meta{maths}}%
 {%
   \BeginArgList
     \csentryargitem{maths} The maths under which to put the arrow.
   \EndArgList
 }

\defgcs{xrightarrow}%
 {\oarg{\meta{subscript}}\marg{\meta{superscript}}}%
 {\nxisty{amsmath} package (Math Mode)}%
 {An extendible right arrow with a superscript and optionally a
  subscript.}%
 {%
   \BeginArgList
     \csentryargitem{subscript} The subscript to go under the arrow.
     \csentryargitem{superscript} The superscript to go over the arrow.
   \EndArgList
 }

\defgcs{xleftarrow}%
 {\oarg{\meta{subscript}}\marg{\meta{superscript}}}%
 {\nxisty{amsmath} package (Math Mode)}%
 {An extendible left arrow with a superscript and optionally a
  subscript.}%
 {%
   \BeginArgList
     \csentryargitem{subscript} The subscript to go under the arrow.
     \csentryargitem{superscript} The superscript to go over the arrow.
   \EndArgList
 }

\defgcs{rightarrow}%
 {}%
 {\LaTeX\ Kernel (Math Mode)}%
 {Right arrow \ensuremath{\rightarrow}.}%
 {}

\defgcs{to}%
 {}%
 {\LaTeX\ Kernel (Math Mode)}%
 {Right arrow \ensuremath{\to}.}%
 {}

\defgcs{Rightarrow}%
 {}%
 {\LaTeX\ Kernel (Math Mode)}%
 {Double-lined right arrow \ensuremath{\Rightarrow}.}%
 {}

\defgcs{leftarrow}%
 {}%
 {\LaTeX\ Kernel (Math Mode)}%
 {Left arrow \ensuremath{\leftarrow}.}%
 {}

\defgcs{gets}%
 {}%
 {\LaTeX\ Kernel (Math Mode)}%
 {Left arrow \ensuremath{\gets}.}%
 {}

\defgcs{Leftarrow}%
 {}%
 {\LaTeX\ Kernel (Math Mode)}%
 {Double-lined left arrow \ensuremath{\Leftarrow}.}%
 {}

\defgcs{downarrow}%
 {}%
 {\LaTeX\ Kernel (Math Mode)}%
 {Down arrow \ensuremath{\downarrow}. (May be used as a delimiter.)}%
 {}

\defgcs{Downarrow}%
 {}%
 {\LaTeX\ Kernel (Math Mode)}%
 {Double-lined down arrow \ensuremath{\Downarrow}. (May be used as a delimiter.)}%
 {}

\defgcs{uparrow}%
 {}%
 {\LaTeX\ Kernel (Math Mode)}%
 {Up arrow \ensuremath{\uparrow}. (May be used as a delimiter.)}%
 {}

\defgcs{Uparrow}%
 {}%
 {\LaTeX\ Kernel (Math Mode)}%
 {Double-lined up arrow \ensuremath{\Uparrow}. (May be used as a delimiter.)}%
 {}

\defgcs{updownarrow}%
 {}%
 {\LaTeX\ Kernel (Math Mode)}%
 {Double-ended vertical arrow \ensuremath{\updownarrow}. (May be used as a delimiter.)}%
 {}

\defgcs{Updownarrow}%
 {}%
 {\LaTeX\ Kernel (Math Mode)}%
 {Double-ended double-lined vertical arrow \ensuremath{\Updownarrow}. (May be used as a delimiter.)}%
 {}

\defgcs{hookrightarrow}%
 {}%
 {\LaTeX\ Kernel (Math Mode)}%
 {Hooked right arrow \ensuremath{\hookrightarrow}.}%
 {}

\defgcs{hookleftarrow}%
 {}%
 {\LaTeX\ Kernel (Math Mode)}%
 {Hooked left arrow \ensuremath{\hookleftarrow}.}%
 {}

\defgcs{leftharpoondown}%
 {}%
 {\LaTeX\ Kernel (Math Mode)}%
 {Left down harpoon \ensuremath{\leftharpoondown}.}%
 {}

\defgcs{leftharpoonup}%
 {}%
 {\LaTeX\ Kernel (Math Mode)}%
 {Left up harpoon \ensuremath{\leftharpoonup}.}%
 {}

\defgcs{rightharpoondown}%
 {}%
 {\LaTeX\ Kernel (Math Mode)}%
 {Right down harpoon \ensuremath{\rightharpoondown}.}%
 {}

\defgcs{rightharpoonup}%
 {}%
 {\LaTeX\ Kernel (Math Mode)}%
 {Right up harpoon \ensuremath{\rightharpoonup}.}%
 {}

\defgcs{rightleftharpoons}%
 {}%
 {\LaTeX\ Kernel (Math Mode)}%
 {Right-left harpoons \ensuremath{\rightleftharpoons}.}%
 {}

\defgcs{leftrightarrow}%
 {}%
 {\LaTeX\ Kernel (Math Mode)}%
 {Double-ended horizontal arrow \ensuremath{\leftrightarrow}.}%
 {}

\defgcs{Leftrightarrow}%
 {}%
 {\LaTeX\ Kernel (Math Mode)}%
 {Double-ended double-lined horizontal arrow \ensuremath{\Leftrightarrow}.}%
 {}

\defgcs{longrightarrow}%
 {}%
 {\LaTeX\ Kernel (Math Mode)}%
 {Long right arrow \ensuremath{\longrightarrow}.}%
 {}

\defgcs{longleftarrow}%
 {}%
 {\LaTeX\ Kernel (Math Mode)}%
 {Long left arrow \ensuremath{\longleftarrow}.}%
 {}

\defgcs{longleftrightarrow}%
 {}%
 {\LaTeX\ Kernel (Math Mode)}%
 {Long double-ended horizontal arrow \ensuremath{\longleftrightarrow}.}%
 {}

\defgcs{Longrightarrow}%
 {}%
 {\LaTeX\ Kernel (Math Mode)}%
 {Long double-lined right arrow \ensuremath{\Longrightarrow}.}%
 {}

\defgcs{Longleftarrow}%
 {}%
 {\LaTeX\ Kernel (Math Mode)}%
 {Long double-lined left arrow \ensuremath{\Longleftarrow}.}%
 {}

\defgcs{Longleftrightarrow}%
 {}%
 {\LaTeX\ Kernel (Math Mode)}%
 {Long double-lined double-ended horizontal arrow \ensuremath{\Longleftrightarrow}.}%
 {}

\defgcs{longmapsto}%
 {}%
 {\LaTeX\ Kernel (Math Mode)}%
 {Long mapping arrow \ensuremath{\longmapsto}.}%
 {}

\defgcs{mapsto}%
 {}%
 {\LaTeX\ Kernel (Math Mode)}%
 {Mapping arrow \ensuremath{\mapsto}.}%
 {}

\defgcs{nearrow}%
 {}%
 {\LaTeX\ Kernel (Math Mode)}%
 {North-East arrow \ensuremath{\nearrow}.}%
 {}

\defgcs{nwarrow}%
 {}%
 {\LaTeX\ Kernel (Math Mode)}%
 {North-West arrow \ensuremath{\nwarrow}.}%
 {}

\defgcs{searrow}%
 {}%
 {\LaTeX\ Kernel (Math Mode)}%
 {South-East arrow \ensuremath{\searrow}.}%
 {}

\defgcs{swarrow}%
 {}%
 {\LaTeX\ Kernel (Math Mode)}%
 {South-West arrow \ensuremath{\swarrow}.}%
 {}

\defgcs{makeindex}%
 {}%
 {\LaTeX\ Kernel (Preamble Only)}%
 {Enables \nxglsi{index}.}%
 {}

\defgcs{index}%
 {\marg{\meta{text}}}%
 {\LaTeX\ Kernel}%
 {Adds indexing information to an external index file. The command
   \nxglsi{makeindex} must be used in the preamble to enable
   this command. The external index file must be post-processed with
   an indexing application, such as \nxiappname{makeindex}.}%
 {%
   \BeginArgList
     \csentryargitem{text} The text to go in the index.
   \EndArgList
 }

\defgcs{printindex}%
 {}%
 {\nxisty{makeidx} package}%
 {Prints the index. Must be used with \nxglsi{makeindex} and \nxglsi{index}.
 (The external index file must first be processed by an indexing application.)}%
 {}

\defgcs{alpha}%
 {}%
 {\LaTeX\ Kernel (Math Mode)}%
 {Greek lower case alpha \ensuremath{\alpha}.}%
 {}

\defgcs{beta}%
 {}%
 {\LaTeX\ Kernel (Math Mode)}%
 {Greek lower case beta \ensuremath{\beta}.}%
 {}

\defgcs{gamma}%
 {}%
 {\LaTeX\ Kernel (Math Mode)}%
 {Greek lower case gamma \ensuremath{\gamma}.}%
 {}

\defgcs{delta}%
 {}%
 {\LaTeX\ Kernel (Math Mode)}%
 {Greek lower case delta \ensuremath{\delta}.}%
 {}

\defgcs{epsilon}%
 {}%
 {\LaTeX\ Kernel (Math Mode)}%
 {Greek lower case epsilon \ensuremath{\epsilon}.}%
 {}

\defgcs{varepsilon}%
 {}%
 {\LaTeX\ Kernel (Math Mode)}%
 {Variant Greek lower case alpha \ensuremath{\varepsilon}.}%
 {}

\defgcs{zeta}%
 {}%
 {\LaTeX\ Kernel (Math Mode)}%
 {Greek lower case zeta \ensuremath{\zeta}.}%
 {}

\defgcs{eta}%
 {}%
 {\LaTeX\ Kernel (Math Mode)}%
 {Greek lower case eta \ensuremath{\eta}.}%
 {}

\defgcs{theta}%
 {}%
 {\LaTeX\ Kernel (Math Mode)}%
 {Greek lower case theta \ensuremath{\theta}.}%
 {}

\defgcs{vartheta}%
 {}%
 {\LaTeX\ Kernel (Math Mode)}%
 {A variant Greek lower case theta \ensuremath{\vartheta}.}%
 {}

\defgcs{iota}%
 {}%
 {\LaTeX\ Kernel (Math Mode)}%
 {Greek lower case iota \ensuremath{\iota}.}%
 {}

\defgcs{kappa}%
 {}%
 {\LaTeX\ Kernel (Math Mode)}%
 {Greek lower case kappa \ensuremath{\kappa}.}%
 {}

\defgcs{lambda}%
 {}%
 {\LaTeX\ Kernel (Math Mode)}%
 {Greek lower case lambda \ensuremath{\lambda}.}%
 {}

\defgcs{mu}%
 {}%
 {\LaTeX\ Kernel (Math Mode)}%
 {Greek lower case mu \ensuremath{\mu}.}%
 {}

\defgcs{nu}%
 {}%
 {\LaTeX\ Kernel (Math Mode)}%
 {Greek lower case nu \ensuremath{\nu}.}%
 {}

\defgcs{xi}%
 {}%
 {\LaTeX\ Kernel (Math Mode)}%
 {Greek lower case xi \ensuremath{\xi}.}%
 {}

\defgcs{pi}%
 {}%
 {\LaTeX\ Kernel (Math Mode)}%
 {Greek lower case pi \ensuremath{\pi}.}%
 {}

\defgcs{varpi}%
 {}%
 {\LaTeX\ Kernel (Math Mode)}%
 {Variant Greek lower case pi \ensuremath{\varpi}.}%
 {}

\defgcs{rho}%
 {}%
 {\LaTeX\ Kernel (Math Mode)}%
 {Greek lower case rho \ensuremath{\rho}.}%
 {}

\defgcs{varrho}%
 {}%
 {\LaTeX\ Kernel (Math Mode)}%
 {Variant Greek lower case rho \ensuremath{\varrho}.}%
 {}

\defgcs{sigma}%
 {}%
 {\LaTeX\ Kernel (Math Mode)}%
 {Greek lower case sigma \ensuremath{\sigma}.}%
 {}

\defgcs{varsigma}%
 {}%
 {\LaTeX\ Kernel (Math Mode)}%
 {Variant Greek lower case sigma \ensuremath{\varsigma}.}%
 {}

\defgcs{tau}%
 {}%
 {\LaTeX\ Kernel (Math Mode)}%
 {Greek lower case tau \ensuremath{\tau}.}%
 {}

\defgcs{upsilon}%
 {}%
 {\LaTeX\ Kernel (Math Mode)}%
 {Greek lower case upsilon \ensuremath{\upsilon}.}%
 {}

\defgcs{phi}%
 {}%
 {\LaTeX\ Kernel (Math Mode)}%
 {Greek lower case phi \ensuremath{\phi}.}%
 {}

\defgcs{varphi}%
 {}%
 {\LaTeX\ Kernel (Math Mode)}%
 {Variant Greek lower case phi \ensuremath{\varphi}.}%
 {}

\defgcs{chi}%
 {}%
 {\LaTeX\ Kernel (Math Mode)}%
 {Greek lower case chi \ensuremath{\chi}.}%
 {}

\defgcs{psi}%
 {}%
 {\LaTeX\ Kernel (Math Mode)}%
 {Greek lower case psi \ensuremath{\psi}.}%
 {}

\defgcs{omega}%
 {}%
 {\LaTeX\ Kernel (Math Mode)}%
 {Greek lower case omega \ensuremath{\omega}.}%
 {}

\defgcs{Gamma}%
 {}%
 {\LaTeX\ Kernel (Math Mode)}%
 {Greek upper case gamma \ensuremath{\Gamma}.}%
 {}

\defgcs{Delta}%
 {}%
 {\LaTeX\ Kernel (Math Mode)}%
 {Greek upper case delta \ensuremath{\Delta}.}%
 {}

\defgcs{Theta}%
 {}%
 {\LaTeX\ Kernel (Math Mode)}%
 {Greek upper case theta \ensuremath{\Theta}.}%
 {}

\defgcs{Lambda}%
 {}%
 {\LaTeX\ Kernel (Math Mode)}%
 {Greek upper case lambda \ensuremath{\Lambda}.}%
 {}

\defgcs{Xi}%
 {}%
 {\LaTeX\ Kernel (Math Mode)}%
 {Greek upper case xi \ensuremath{\Xi}.}%
 {}

\defgcs{Pi}%
 {}%
 {\LaTeX\ Kernel (Math Mode)}%
 {Greek upper case pi \ensuremath{\Pi}.}%
 {}

\defgcs{Sigma}%
 {}%
 {\LaTeX\ Kernel (Math Mode)}%
 {Greek upper case sigma \ensuremath{\Sigma}.}%
 {}

\defgcs{Upsilon}%
 {}%
 {\LaTeX\ Kernel (Math Mode)}%
 {Greek upper case upsilon \ensuremath{\Upsilon}.}%
 {}

\defgcs{Phi}%
 {}%
 {\LaTeX\ Kernel (Math Mode)}%
 {Greek upper case phi \ensuremath{\Phi}.}%
 {}

\defgcs{Psi}%
 {}%
 {\LaTeX\ Kernel (Math Mode)}%
 {Greek upper case psi \ensuremath{\Psi}.}%
 {}

\defgcs{Omega}%
 {}%
 {\LaTeX\ Kernel (Math Mode)}%
 {Greek upper case omega \ensuremath{\Omega}.}%
 {}

\defgcs{sb}%
 {\marg{\meta{maths}}}%
 {\LaTeX\ Kernel (Math Mode)}%
 {Displays its argument as a subscript.}%
 {%
   \BeginArgList
     \csentryargitem{maths} The subscript.
   \EndArgList
 }

\defgcs{sp}%
 {\marg{\meta{maths}}}%
 {\LaTeX\ Kernel (Math Mode)}%
 {Displays its argument as a superscript.}%
 {%
   \BeginArgList
     \csentryargitem{maths} The superscript.
   \EndArgList
 }

\defgcs{pmod}%
 {\marg{\meta{maths}}}%
 {\LaTeX\ Kernel (Math Mode)}%
 {Modulo operator with parentheses.}%
 {%
   \BeginArgList
     \csentryargitem{maths} The modulo (placed in parentheses).
   \EndArgList
 }

\defgcs{mod}%
 {\marg{\meta{maths}}}%
 {\nxisty{amsmath} (Math Mode)}%
 {Modulo operator without parentheses.}%
 {%
   \BeginArgList
     \csentryargitem{maths} The modulo (placed in parentheses).
   \EndArgList
 }

\defgcs{pod}%
 {\marg{\meta{maths}}}%
 {\nxisty{amsmath} (Math Mode)}%
 {Modulo operator with parentheses but no \dq{mod}.}%
 {%
   \BeginArgList
     \csentryargitem{maths} The modulo (placed in parentheses).
   \EndArgList
 }

\defgcs{bmod}%
 {}%
 {\LaTeX\ Kernel (Math Mode)}%
 {Modulo operator.}%
 {}

\defgcs{arccos}%
 {}%
 {\LaTeX\ Kernel (Math Mode)}%
 {Typesets \ensuremath{\arccos} function name.}%
 {}

\defgcs{arcsin}%
 {}%
 {\LaTeX\ Kernel (Math Mode)}%
 {Typesets \ensuremath{\arcsin} function name.}%
 {}

\defgcs{arctan}%
 {}%
 {\LaTeX\ Kernel (Math Mode)}%
 {Typesets \ensuremath{\arctan} function name.}%
 {}

\defgcs{arg}%
 {}%
 {\LaTeX\ Kernel (Math Mode)}%
 {Typesets \ensuremath{\arg} function name.}%
 {}

\defgcs{cos}%
 {}%
 {\LaTeX\ Kernel (Math Mode)}%
 {Typesets \ensuremath{\cos} function name.}%
 {}

\defgcs{cosh}%
 {}%
 {\LaTeX\ Kernel (Math Mode)}%
 {Typesets \ensuremath{\cosh} function name.}%
 {}

\defgcs{cot}%
 {}%
 {\LaTeX\ Kernel (Math Mode)}%
 {Typesets \ensuremath{\cot} function name.}%
 {}

\defgcs{coth}%
 {}%
 {\LaTeX\ Kernel (Math Mode)}%
 {Typesets \ensuremath{\coth} function name.}%
 {}

\defgcs{csc}%
 {}%
 {\LaTeX\ Kernel (Math Mode)}%
 {Typesets \ensuremath{\csc} function name.}%
 {}

\defgcs{deg}%
 {}%
 {\LaTeX\ Kernel (Math Mode)}%
 {Typesets \ensuremath{\deg} function name.}%
 {}

\defgcs{det}%
 {}%
 {\LaTeX\ Kernel (Math Mode)}%
 {Typesets \ensuremath{\det} function name (may have limits via \nxglsi{underscorechar} or \nxglsi{circumchar}).}%
 {}

\defgcs{dim}%
 {}%
 {\LaTeX\ Kernel (Math Mode)}%
 {Typesets \ensuremath{\dim} function name.}%
 {}

\defgcs{exp}%
 {}%
 {\LaTeX\ Kernel (Math Mode)}%
 {Typesets \ensuremath{\exp} function name.}%
 {}

\defgcs{gcd}%
 {}%
 {\LaTeX\ Kernel (Math Mode)}%
 {Typesets \ensuremath{\gcd} function name (may have limits via \nxglsi{underscorechar} or \nxglsi{circumchar}).}%
 {}

\defgcs{hom}%
 {}%
 {\LaTeX\ Kernel (Math Mode)}%
 {Typesets \ensuremath{\hom} function name.}%
 {}

\defgcs{inf}%
 {}%
 {\LaTeX\ Kernel (Math Mode)}%
 {Typesets \ensuremath{\inf} function name (may have limits via \nxglsi{underscorechar} or \nxglsi{circumchar}).}%
 {}

\defgcs{injlim}%
 {}%
 {\nxisty{amsmath} (Math Mode)}%
 {Typesets \ensuremath{\injlim} function name (may have limits via \nxglsi{underscorechar} or \nxglsi{circumchar}).}%
 {}

\defgcs{projlim}%
 {}%
 {\nxisty{amsmath} (Math Mode)}%
 {Typesets \ensuremath{\projlim} function name (may have limits via \nxglsi{underscorechar} or \nxglsi{circumchar}).}%
 {}

\defgcs{varlimsup}%
 {}%
 {\nxisty{amsmath} (Math Mode)}%
 {Typesets \ensuremath{\varlimsup} function name (may have limits via \nxglsi{underscorechar} or \nxglsi{circumchar}).}%
 {}

\defgcs{varliminf}%
 {}%
 {\nxisty{amsmath} (Math Mode)}%
 {Typesets \ensuremath{\varliminf} function name (may have limits via \nxglsi{underscorechar} or \nxglsi{circumchar}).}%
 {}

\defgcs{varinjlim}%
 {}%
 {\nxisty{amsmath} (Math Mode)}%
 {Typesets \ensuremath{\varinjlim} function name (may have limits via \nxglsi{underscorechar} or \nxglsi{circumchar}).}%
 {}

\defgcs{varprojlim}%
 {}%
 {\nxisty{amsmath} (Math Mode)}%
 {Typesets \ensuremath{\varprojlim} function name (may have limits via \nxglsi{underscorechar} or \nxglsi{circumchar}).}%
 {}

\defgcs{ker}%
 {}%
 {\LaTeX\ Kernel (Math Mode)}%
 {Typesets \ensuremath{\ker} function name.}%
 {}

\defgcs{lg}%
 {}%
 {\LaTeX\ Kernel (Math Mode)}%
 {Typesets \ensuremath{\lg} function name.}%
 {}

\defgcs{lim}%
 {}%
 {\LaTeX\ Kernel (Math Mode)}%
 {Typesets \ensuremath{\lim} function name (may have limits via \nxglsi{underscorechar} or \nxglsi{circumchar}).}%
 {}

\defgcs{liminf}%
 {}%
 {\LaTeX\ Kernel (Math Mode)}%
 {Typesets \ensuremath{\liminf} function name (may have limits via \nxglsi{underscorechar} or \nxglsi{circumchar}).}%
 {}

\defgcs{limsup}%
 {}%
 {\LaTeX\ Kernel (Math Mode)}%
 {Typesets \ensuremath{\limsup} function name (may have limits via \nxglsi{underscorechar} or \nxglsi{circumchar}).}%
 {}

\defgcs{ln}%
 {}%
 {\LaTeX\ Kernel (Math Mode)}%
 {Typesets \ensuremath{\ln} function name.}%
 {}

\defgcs{log}%
 {}%
 {\LaTeX\ Kernel (Math Mode)}%
 {Typesets \ensuremath{\log} function name.}%
 {}

\defgcs{max}%
 {}%
 {\LaTeX\ Kernel (Math Mode)}%
 {Typesets \ensuremath{\max} function name (may have limits via \nxglsi{underscorechar} or \nxglsi{circumchar}).}%
 {}

\defgcs{min}%
 {}%
 {\LaTeX\ Kernel (Math Mode)}%
 {Typesets \ensuremath{\min} function name (may have limits via \nxglsi{underscorechar} or \nxglsi{circumchar}).}%
 {}

\defgcs{Pr}%
 {}%
 {\LaTeX\ Kernel (Math Mode)}%
 {Typesets \ensuremath{\Pr} function name (may have limits via \nxglsi{underscorechar} or \nxglsi{circumchar}).}%
 {}

\defgcs{sec}%
 {}%
 {\LaTeX\ Kernel (Math Mode)}%
 {Typesets \ensuremath{\sec} function name.}%
 {}

\defgcs{sin}%
 {}%
 {\LaTeX\ Kernel (Math Mode)}%
 {Typesets \ensuremath{\sin} function name.}%
 {}

\defgcs{sinh}%
 {}%
 {\LaTeX\ Kernel (Math Mode)}%
 {Typesets \ensuremath{\sinh} function name.}%
 {}

\defgcs{sup}%
 {}%
 {\LaTeX\ Kernel (Math Mode)}%
 {Typesets \ensuremath{\sup} function name (may have limits via \nxglsi{underscorechar} or \nxglsi{circumchar}).}%
 {}

\defgcs{tan}%
 {}%
 {\LaTeX\ Kernel (Math Mode)}%
 {Typesets \ensuremath{\tan} function name.}%
 {}

\defgcs{tanh}%
 {}%
 {\LaTeX\ Kernel (Math Mode)}%
 {Typesets \ensuremath{\tanh} function name.}%
 {}

\defgcs{infty}%
 {}%
 {\LaTeX\ Kernel (Math Mode)}%
 {Infinity \ensuremath{\infty} symbol.}%
 {}

\defgcs{partial}%
 {}%
 {\LaTeX\ Kernel (Math Mode)}%
 {Partial \ensuremath{\partial} symbol.}%
 {}

\defgcs{subset}%
 {}%
 {\LaTeX\ Kernel (Math Mode)}%
 {Subset \ensuremath{\subset} symbol.}%
 {}

\defgcs{approx}%
 {}%
 {\LaTeX\ Kernel (Math Mode)}%
 {Relational \ensuremath{\approx} symbol.}%
 {}

\defgcs{asymp}%
 {}%
 {\LaTeX\ Kernel (Math Mode)}%
 {Relational \ensuremath{\asymp} symbol.}%
 {}

\defgcs{bowtie}%
 {}%
 {\LaTeX\ Kernel (Math Mode)}%
 {Relational \ensuremath{\bowtie} symbol.}%
 {}

\defgcs{cong}%
 {}%
 {\LaTeX\ Kernel (Math Mode)}%
 {Relational \ensuremath{\cong} symbol.}%
 {}

\defgcs{dashv}%
 {}%
 {\LaTeX\ Kernel (Math Mode)}%
 {Relational \ensuremath{\dashv} symbol.}%
 {}

\defgcs{doteq}%
 {}%
 {\LaTeX\ Kernel (Math Mode)}%
 {Relational \makeimg{equals symbol with dot on top}{\ensuremath{\doteq}} symbol.}%
 {}

\defgcs{equiv}%
 {}%
 {\LaTeX\ Kernel (Math Mode)}%
 {Relational \makeimg{equivalent}{\ensuremath{\equiv}} symbol.}%
 {}

\defgcs{frown}%
 {}%
 {\LaTeX\ Kernel (Math Mode)}%
 {Relational \ensuremath{\frown} symbol.}%
 {}

\defgcs{ge}%
 {}%
 {\LaTeX\ Kernel (Math Mode)}%
 {Relational \makeimg{greater than or equal to}{\ensuremath{\ge}} symbol.}%
 {}

\defgcs{geq}%
 {}%
 {\LaTeX\ Kernel (Math Mode)}%
 {Relational \makeimg{greater than or equal to}{\ensuremath{\geq}} symbol.}%
 {}

\defgcs{gg}%
 {}%
 {\LaTeX\ Kernel (Math Mode)}%
 {Relational \makeimg{much greater than}{\ensuremath{\gg}} symbol.}%
 {}

\defgcs{in}%
 {}%
 {\LaTeX\ Kernel (Math Mode)}%
 {Relational \ensuremath{\in} symbol.}%
 {}

\defgcs{le}%
 {}%
 {\LaTeX\ Kernel (Math Mode)}%
 {Relational \makeimg{less than or equal to}{\ensuremath{\le}} symbol.}%
 {}

\defgcs{leq}%
 {}%
 {\LaTeX\ Kernel (Math Mode)}%
 {Relational \makeimg{less than or equal to}{\ensuremath{\leq}} symbol.}%
 {}

\defgcs{ll}%
 {}%
 {\LaTeX\ Kernel (Math Mode)}%
 {Relational \makeimg{much less than}{\ensuremath{\ll}} symbol.}%
 {}

\defgcs{mid}%
 {}%
 {\LaTeX\ Kernel (Math Mode)}%
 {Relational \ensuremath{\mid} symbol.}%
 {}

\defgcs{models}%
 {}%
 {\LaTeX\ Kernel (Math Mode)}%
 {Relational \ensuremath{\models} symbol.}%
 {}

\defgcs{neq}%
 {}%
 {\LaTeX\ Kernel (Math Mode)}%
 {Relational \makeimg{not equal to}{\ensuremath{\neq}} symbol.}%
 {}

\defgcs{ni}%
 {}%
 {\LaTeX\ Kernel (Math Mode)}%
 {Relational \makeimg{contains element}{\ensuremath{\ni}} symbol.}%
 {}

\defgcs{notin}%
 {}%
 {\LaTeX\ Kernel (Math Mode)}%
 {Relational \makeimg{not element of}{\ensuremath{\notin}} symbol.}%
 {}

\defgcs{parallel}%
 {}%
 {\LaTeX\ Kernel (Math Mode)}%
 {Relational \ensuremath{\parallel} symbol.}%
 {}

\defgcs{prec}%
 {}%
 {\LaTeX\ Kernel (Math Mode)}%
 {Relational \ensuremath{\prec} symbol.}%
 {}

\defgcs{preceq}%
 {}%
 {\LaTeX\ Kernel (Math Mode)}%
 {Relational \ensuremath{\preceq} symbol.}%
 {}

\defgcs{perp}%
 {}%
 {\LaTeX\ Kernel (Math Mode)}%
 {Relational \makeimg{perpendicular to}{\ensuremath{\perp}} symbol.}%
 {}

\defgcs{propto}%
 {}%
 {\LaTeX\ Kernel (Math Mode)}%
 {Relational \makeimg{proportional to}{\ensuremath{\propto}} symbol.}%
 {}

\defgcs{sim}%
 {}%
 {\LaTeX\ Kernel (Math Mode)}%
 {Relational \ensuremath{\sim} symbol.}%
 {}

\defgcs{simeq}%
 {}%
 {\LaTeX\ Kernel (Math Mode)}%
 {Relational \ensuremath{\simeq} symbol.}%
 {}

\defgcs{smile}%
 {}%
 {\LaTeX\ Kernel (Math Mode)}%
 {Relational \ensuremath{\smile} symbol.}%
 {}

\defgcs{sqsubseteq}%
 {}%
 {\LaTeX\ Kernel (Math Mode)}%
 {Relational \ensuremath{\sqsubseteq} symbol.}%
 {}

\defgcs{sqsupseteq}%
 {}%
 {\LaTeX\ Kernel (Math Mode)}%
 {Relational \ensuremath{\sqsupseteq} symbol.}%
 {}

\defgcs{subseteq}%
 {}%
 {\LaTeX\ Kernel (Math Mode)}%
 {Relational \ensuremath{\subseteq} symbol.}%
 {}

\defgcs{succ}%
 {}%
 {\LaTeX\ Kernel (Math Mode)}%
 {Relational \ensuremath{\succ} symbol.}%
 {}

\defgcs{succeq}%
 {}%
 {\LaTeX\ Kernel (Math Mode)}%
 {Relational \ensuremath{\succeq} symbol.}%
 {}

\defgcs{supset}%
 {}%
 {\LaTeX\ Kernel (Math Mode)}%
 {Relational \ensuremath{\supset} symbol.}%
 {}

\defgcs{supseteq}%
 {}%
 {\LaTeX\ Kernel (Math Mode)}%
 {Relational \ensuremath{\supseteq} symbol.}%
 {}

\defgcs{vdash}%
 {}%
 {\LaTeX\ Kernel (Math Mode)}%
 {Relational \ensuremath{\vdash} symbol.}%
 {}

\defgcs{amalg}%
 {}%
 {\LaTeX\ Kernel (Math Mode)}%
 {Binary operator \ensuremath{\amalg} symbol.}%
 {}

\defgcs{ast}%
 {}%
 {\LaTeX\ Kernel (Math Mode)}%
 {Binary operator \makeimg{asterisk}{\ensuremath{\ast}} symbol.}%
 {}

\defgcs{bullet}%
 {}%
 {\LaTeX\ Kernel (Math Mode)}%
 {Binary operator \makeimg{bullet point}{\ensuremath{\bullet}} symbol.}%
 {}

\defgcs{bigcirc}%
 {}%
 {\LaTeX\ Kernel (Math Mode)}%
 {Binary operator \ensuremath{\bigcirc} symbol.}%
 {}

\defgcs{bigtriangledown}%
 {}%
 {\LaTeX\ Kernel (Math Mode)}%
 {Binary operator \ensuremath{\bigtriangledown} symbol.}%
 {}

\defgcs{bigtriangleup}%
 {}%
 {\LaTeX\ Kernel (Math Mode)}%
 {Binary operator \ensuremath{\bigtriangleup} symbol.}%
 {}

\defgcs{cap}%
 {}%
 {\LaTeX\ Kernel (Math Mode)}%
 {Binary operator \ensuremath{\cap} symbol.}%
 {}

\defgcs{cdot}%
 {}%
 {\LaTeX\ Kernel (Math Mode)}%
 {Centred dot \ensuremath{\cdot} symbol.}%
 {}

\defgcs{circ}%
 {}%
 {\LaTeX\ Kernel (Math Mode)}%
 {Circle \ensuremath{\circ} symbol.}%
 {}

\defgcs{cup}%
 {}%
 {\LaTeX\ Kernel (Math Mode)}%
 {Operator \ensuremath{\cup} symbol.}%
 {}

\defgcs{dagger}%
 {}%
 {\LaTeX\ Kernel (Math Mode)}%
 {Binary operator \ensuremath{\dagger} symbol.}%
 {}

\defgcs{ddagger}%
 {}%
 {\LaTeX\ Kernel (Math Mode)}%
 {Binary operator \ensuremath{\ddagger} symbol.}%
 {}

\defgcs{diamond}%
 {}%
 {\LaTeX\ Kernel (Math Mode)}%
 {Binary operator \ensuremath{\diamond} symbol.}%
 {}

\defgcs{div}%
 {}%
 {\LaTeX\ Kernel (Math Mode)}%
 {Division operator \ensuremath{\div} symbol.}%
 {}

\defgcs{mp}%
 {}%
 {\LaTeX\ Kernel (Math Mode)}%
 {Minus or plus operator \makeimg{minus or plus}{\ensuremath{\mp}} symbol.}%
 {}

\defgcs{odot}%
 {}%
 {\LaTeX\ Kernel (Math Mode)}%
 {Operator \ensuremath{\odot} symbol.}%
 {}

\defgcs{ominus}%
 {}%
 {\LaTeX\ Kernel (Math Mode)}%
 {Operator \ensuremath{\ominus} symbol.}%
 {}

\defgcs{oplus}%
 {}%
 {\LaTeX\ Kernel (Math Mode)}%
 {Operator \ensuremath{\oplus} symbol.}%
 {}

\defgcs{oslash}%
 {}%
 {\LaTeX\ Kernel (Math Mode)}%
 {Operator \ensuremath{\oslash} symbol.}%
 {}

\defgcs{otimes}%
 {}%
 {\LaTeX\ Kernel (Math Mode)}%
 {Operator \ensuremath{\otimes} symbol.}%
 {}

\defgcs{pm}%
 {}%
 {\LaTeX\ Kernel (Math Mode)}%
 {Operator \makeimg{plus or minus}{\ensuremath{\pm}} symbol.}%
 {}

\defgcs{setminus}%
 {}%
 {\LaTeX\ Kernel (Math Mode)}%
 {Operator \ensuremath{\setminus} symbol.}%
 {}

\defgcs{sqcap}%
 {}%
 {\LaTeX\ Kernel (Math Mode)}%
 {Operator \ensuremath{\sqcap} symbol.}%
 {}

\defgcs{sqcup}%
 {}%
 {\LaTeX\ Kernel (Math Mode)}%
 {Operator \ensuremath{\sqcup} symbol.}%
 {}

\defgcs{star}%
 {}%
 {\LaTeX\ Kernel (Math Mode)}%
 {Operator \ensuremath{\star} symbol.}%
 {}

\defgcs{times}%
 {}%
 {\LaTeX\ Kernel (Math Mode)}%
 {Operator \ensuremath{\times} symbol.}%
 {}

\defgcs{triangleleft}%
 {}%
 {\LaTeX\ Kernel (Math Mode)}%
 {Binary operator \ensuremath{\triangleleft} symbol.}%
 {}

\defgcs{triangleright}%
 {}%
 {\LaTeX\ Kernel (Math Mode)}%
 {Binary operator \ensuremath{\triangleright} symbol.}%
 {}

\defgcs{uplus}%
 {}%
 {\LaTeX\ Kernel (Math Mode)}%
 {Operator \ensuremath{\uplus} symbol.}%
 {}

\defgcs{vee}%
 {}%
 {\LaTeX\ Kernel (Math Mode)}%
 {Operator \ensuremath{\vee} symbol.}%
 {}

\defgcs{wedge}%
 {}%
 {\LaTeX\ Kernel (Math Mode)}%
 {Operator \ensuremath{\wedge} symbol.}%
 {}

\defgcs{wr}%
 {}%
 {\LaTeX\ Kernel (Math Mode)}%
 {Operator \ensuremath{\wr} symbol.}%
 {}

\defgcs{sum}%
 {}%
 {\LaTeX\ Kernel (Math Mode)}%
 {Summation \ensuremath{\sum} symbol (may take limits).}%
 {}

\defgcs{int}%
 {}%
 {\LaTeX\ Kernel (Math Mode)}%
 {Integral \ensuremath{\int} symbol (may take limits).}%
 {}

\defgcs{oint}%
 {}%
 {\LaTeX\ Kernel (Math Mode)}%
 {Closed path integral \ensuremath{\oint} symbol (may take limits).}%
 {}

\defgcs{prod}%
 {}%
 {\LaTeX\ Kernel (Math Mode)}%
 {Product \ensuremath{\prod} symbol (may take limits).}%
 {}

\defgcs{coprod}%
 {}%
 {\LaTeX\ Kernel (Math Mode)}%
 {Co-product \ensuremath{\coprod} symbol (may take limits).}%
 {}

\defgcs{bigcap}%
 {}%
 {\LaTeX\ Kernel (Math Mode)}%
 {Collection intersection \ensuremath{\bigcap} symbol (may take limits).}%
 {}

\defgcs{bigcup}%
 {}%
 {\LaTeX\ Kernel (Math Mode)}%
 {Collection union \ensuremath{\bigcup} symbol (may take limits).}%
 {}

\defgcs{bigsqcup}%
 {}%
 {\LaTeX\ Kernel (Math Mode)}%
 {Big operator \ensuremath{\bigsqcup} (may take limits).}%
 {}

\defgcs{bigvee}%
 {}%
 {\LaTeX\ Kernel (Math Mode)}%
 {Big operator \ensuremath{\bigvee} (may take limits).}%
 {}

\defgcs{bigwedge}%
 {}%
 {\LaTeX\ Kernel (Math Mode)}%
 {Big operator \ensuremath{\bigwedge} (may take limits).}%
 {}

\defgcs{bigodot}%
 {}%
 {\LaTeX\ Kernel (Math Mode)}%
 {Big operator \ensuremath{\bigodot} (may take limits).}%
 {}

\defgcs{bigotimes}%
 {}%
 {\LaTeX\ Kernel (Math Mode)}%
 {Big operator \ensuremath{\bigotimes} (may take limits).}%
 {}

\defgcs{bigoplus}%
 {}%
 {\LaTeX\ Kernel (Math Mode)}%
 {Big operator \ensuremath{\bigoplus} (may take limits).}%
 {}

\defgcs{biguplus}%
 {}%
 {\LaTeX\ Kernel (Math Mode)}%
 {Big operator \ensuremath{\biguplus} (may take limits).}%
 {}

\defgcs{not}%
 {\meta{symbol command}}%
 {\LaTeX\ Kernel (Math Mode)}%
 {Negates the following symbol. Example: \cmdname{not}\nxglsi{subset} produces
   \ensuremath{\not\subset}.}%
 {%
   \BeginArgList
    \csentryargitem{symbol command} The symbol to be negated.
   \EndArgList
 }

\defgcs{ldots}%
 {}%
 {\LaTeX\ Kernel}%
 {Ellipses \ldots\ symbol.}%
 {}

\defgcs{cdots}%
 {}%
 {\LaTeX\ Kernel (Math Mode)}%
 {Centred ellipses \ensuremath{\cdots} symbol.}%
 {}

\defgcs{vdots}%
 {}%
 {\LaTeX\ Kernel (Math Mode)}%
 {Vertical ellipses \ensuremath{\vdots} symbol.}%
 {}

\defgcs{ddots}%
 {}%
 {\LaTeX\ Kernel (Math Mode)}%
 {Diagonal ellipses \ensuremath{\ddots} symbol.}%
 {}

\defgcs{dotsc}%
 {}%
 {\nxisty{amsmath} (Math Mode)}%
 {Ellipses \ensuremath{\dotsc} for dots with commas.}%
 {}

\defgcs{dotsb}%
 {}%
 {\nxisty{amsmath} (Math Mode)}%
 {Ellipses \ensuremath{\dotsb} for dots with binary
  operators\slash relations.}%
 {}

\defgcs{dotsm}%
 {}%
 {\nxisty{amsmath} (Math Mode)}%
 {Ellipses \ensuremath{\dotsm} for dots with multiplications.}%
 {}

\defgcs{dotsi}%
 {}%
 {\nxisty{amsmath} (Math Mode)}%
 {Ellipses \ensuremath{\dotsi} for dots with integrals.}%
 {}

\defgcs{dotso}%
 {}%
 {\nxisty{amsmath} (Math Mode)}%
 {Ellipses \ensuremath{\dotso} for general dots.}%
 {}

\defgcs{langle}%
 {}%
 {\LaTeX\ Kernel (Math Mode)}%
 {Left-angled \ensuremath{\langle} delimiter.}%
 {}

\defgcs{rangle}%
 {}%
 {\LaTeX\ Kernel (Math Mode)}%
 {Right-angled \ensuremath{\rangle} delimiter.}%
 {}

\defgcs{lfloor}%
 {}%
 {\LaTeX\ Kernel (Math Mode)}%
 {Left floor \ensuremath{\lfloor} delimiter.}%
 {}

\defgcs{rfloor}%
 {}%
 {\LaTeX\ Kernel (Math Mode)}%
 {Right floor \ensuremath{\rfloor} delimiter.}%
 {}

\defgcs{lceil}%
 {}%
 {\LaTeX\ Kernel (Math Mode)}%
 {Left ceil \ensuremath{\lceil} delimiter.}%
 {}

\defgcs{rceil}%
 {}%
 {\LaTeX\ Kernel (Math Mode)}%
 {Right ceil \ensuremath{\rceil} delimiter.}%
 {}

\defgcs{lvert}%
 {}%
 {\nxisty{amsmath} (Math Mode)}%
 {Left vertical bar \ensuremath{\lvert} delimiter.}%
 {}

\defgcs{rvert}%
 {}%
 {\nxisty{amsmath} (Math Mode)}%
 {Right vertical bar \ensuremath{\rvert} delimiter.}%
 {}

\defgcs{lVert}%
 {}%
 {\nxisty{amsmath} (Math Mode)}%
 {Left double vertical bar \ensuremath{\lVert} delimiter.}%
 {}

\defgcs{rVert}%
 {}%
 {\nxisty{amsmath} (Math Mode)}%
 {Right double vertical bar \ensuremath{\rVert} delimiter.}%
 {}

\defgcs{vec}%
 {\marg{\meta{c}}}%
 {\LaTeX\ Kernel (Math Mode)}%
 {Typesets its argument as a vector (defaults to a right arrow
   accent).}%
 {%
   \BeginArgList
    \csentryargitem{c} The character or symbol that represents a vector.
   \EndArgList
 }

\defgcs{forall}%
 {}%
 {\LaTeX\ Kernel (Math Mode)}%
 {\dq{For all} \ensuremath{\forall} symbol.}%
 {}


\defgcs{noindent}%
 {}%
 {\LaTeX\ Kernel}%
 {Suppress the indentation that would usually occur at the start of
  the next paragraph.}%
 {}

\defgcs{refstepcounter}%
 {\marg{\meta{counter}}}%
 {\LaTeX\ Kernel}%
 {Increments the value of the given counter by one and allows the
  counter to be cross-referenced using \nxglsni{ref} and \nxglsi{label}.}%
 {%
   \BeginArgList
    \csentryargitem{counter} The name of the counter that needs
     incrementing.
   \EndArgList
 }

\defgcs{stepcounter}%
 {\marg{\meta{counter}}}%
 {\LaTeX\ Kernel}%
 {Increments the value of the given counter by one.}%
 {%
   \BeginArgList
    \csentryargitem{counter} The name of the counter that needs
     incrementing.
   \EndArgList
 }

\defgcs{setcounter}%
 {\marg{\meta{counter}}\marg{\meta{number}}}%
 {\LaTeX\ Kernel}%
 {Sets the value of a counter.}%
 {%
   \BeginArgList
    \csentryargitem{counter} The name of the counter that needs
     changing.
    \csentryargitem{number} The new value. (Must be an integer.)
   \EndArgList
 }

\defgcs{addtocounter}%
 {\marg{\meta{counter}}\marg{\meta{increment}}}%
 {\LaTeX\ Kernel}%
 {Increments the value of a counter by the given amount.}%
 {%
   \BeginArgList
    \csentryargitem{counter} The name of the counter that needs
     changing.
    \csentryargitem{number} The amount by which to increment the counter. (Must be an integer.)
   \EndArgList
 }

\defgcs{newcounter}%
 {\marg{\meta{counter}}\oarg{\meta{outer counter}}}%
 {\LaTeX\ Kernel}%
 {Defines a new counter.}%
 {%
   \BeginArgList
    \csentryargitem{counter} The name of the new counter.
    \csentryargitem{outer counter} The name of the parent counter.
   \EndArgList
 }


\defgcs{value}%
 {\marg{\meta{counter}}}%
 {\LaTeX\ Kernel}%
 {References the value of the given counter where a number rather
   than a counter name is required.}%
 {%
   \BeginArgList
    \csentryargitem{counter} The name of the counter.
   \EndArgList
 }

\defgcs{parindent}%
 {}%
 {\LaTeX\ Kernel}%
 {A \nxglsi{length} register that stores the indentation at the start of paragraphs.}%
 {}


\defgcs{setlength}%
 {\marg{\meta{register}}\marg{\meta{dimension}}}%
 {\LaTeX\ Kernel}%
 {Sets the value of a length register.}%
 {%
   \BeginArgList
     \csentryargitem{register} The name of the length register (for
       example \nxglsi{parindent}).
     \csentryargitem{dimension} The new dimension.
   \EndArgList
 }

\defgcs{addtolength}%
 {\marg{\meta{register}}\marg{\meta{dimension}}}%
 {\LaTeX\ Kernel}%
 {Adds \meta{dimension} to the value of the given length register.}%
 {%
   \BeginArgList
     \csentryargitem{register} The name of the length register (for
       example \nxglsi{parindent}).
     \csentryargitem{dimension} The dimension to add to the length.
   \EndArgList
 }

\defgcs{currenttime}%
 {}%
 {\nxisty{datetime} package}%
 {Inserts into the output file the time when the \LaTeX\ 
  application created it from the source code.}%
 {}

\defgcs{ddmmyyyydate}%
 {}%
 {\nxisty{datetime} package}%
 {Changes the format of \nxglsi{today} so that it displays the date in
  the form {\ddmmyyyydate\today} (day\slash month\slash year in digits).}%
 {}

\defgcs{subref}%
 {\marg{\meta{label}}}%
 {\nxisty{subcaption} package}%
 {Analogous to \nxglsi{ref} but only references the subfigure
  or subtable caption.}%
 {%
   \BeginArgList
    \csentryargitem{label} The required label that was used in the
     corresponding \nxglsi{label}.
   \EndArgList
 }

\defgcs{eqref}%
 {\marg{\meta{label}}}%
 {\nxisty{amsmath} package}%
 {Short cut for \texttt{(\nxglsi{ref}\marg{\meta{label}})} for
  referencing equations.}%
 {%
   \BeginArgList
    \csentryargitem{label} The required label that was used in the
     corresponding \nxglsi{label}.
   \EndArgList
 }

\defgcs{frenchspacing}%
 {}%
 {\LaTeX\ Kernel}%
 {Switch to French spacing.}%
 {}

\defgcs{nonfrenchspacing}%
 {}%
 {\LaTeX\ Kernel}%
 {Switch to English spacing.}%
 {}

\defgcs{newline}%
 {}%
 {\LaTeX\ Kernel}%
 {Forces a line break.}%
 {}

\defgcs{linebreak}%
 {\oarg{\meta{n}}}%
 {\LaTeX\ Kernel}%
 {Requests a line break, ensuring the paragraph remains justified.
  This may cause excess white space in the paragraph.}%
 {%
   \BeginArgList
    \csentryargitem{n} An integer from~0 to~4 indicating how strongly
    you want a line break to occur. The higher the number, the
    stronger the request to break the line.
   \EndArgList
 }

\defgcs{displaybreak}%
 {\oarg{\meta{n}}}%
 {\nxisty{amsmath} package}%
 {Allows a page break in multi-lined maths environments, such as
  \nxglsi{env-align}.}%
 {%
   \BeginArgList
    \csentryargitem{n} An integer from~0 to~4 indicating how strongly
    you want a page break to occur. The higher the number, the
    stronger the request.
   \EndArgList
 }

\defgsymcs[beginmath]{\openparensym}%
 {}%
 {\LaTeX\ Kernel}%
 {Equivalent to \nxglsni{begin}\marg{math}.}%
 {}

\defgsymcs[endmath]{\closeparensym}%
 {}%
 {\LaTeX\ Kernel}%
 {Equivalent to \nxglsni{end}\marg{math}.}%
 {}

\defgchar
 {openparen}
 {\openparensym}
 {}
 {\LaTeX\ Kernel}
 {%
   Opening parenthesis in text mode or left round bracket
   delimiter in math mode.%
 }
 {}

\defgchar
 {closeparen}
 {\closeparensym}
 {}
 {\LaTeX\ Kernel}
 {%
   Closing parenthesis in text mode or right round bracket
   delimiter in math mode.%
 }
 {}

\defgsymcs[begindispmath]{\opensqsym}%
 {}%
 {\LaTeX\ Kernel (inconsistency corrected in \nxisty{amsmath})}%
 {Starts an unnumbered single-line of displayed maths.}%
 {}

\defgsymcs[enddispmath]{\closesqsym}%
 {}%
 {\LaTeX\ Kernel (inconsistency corrected in \nxisty{amsmath})}%
 {Ends an unnumbered single-line of displayed maths.}%
 {}

\defgchar
 {opensq}
 {\opensqsym}
 {}
 {\LaTeX\ Kernel}
 {\nopostdesc}
 {}

\defgchildchar
 {opt.opensq}
 {opensq}
 {\opensqsym}
 {Open delimiter of an \nxglsi{optional}}

\defgchildchar
 {text.opensq}
 {opensq}
 {\opensqsym}
 {Open square bracket in text mode}

\defgchildchar
 {delimiter.opensq}
 {opensq}
 {\opensqsym}
 {Left square bracket delimiter in math mode}

\defgchar
 {closesq}
 {\closesqsym}
 {}
 {\LaTeX\ Kernel}
 {\nopostdesc}
 {}

\defgchildchar
 {opt.closesq}
 {closesq}
 {\closesqsym}
 {Closing delimiter of an \nxglsi{optional}}

\defgchildchar
 {text.closesq}
 {closesq}
 {\closesqsym}
 {Closing square bracket in text mode}

\defgchildchar
 {delimiter.closesq}
 {closesq}
 {\closesqsym}
 {Right square bracket delimiter in math mode}

\defgsymcs[comma]{\commasym}%
 {}%
 {\LaTeX\ Kernel}%
 {Thin space.}%
 {}

\defgcs{thinspace}%
 {}%
 {\LaTeX\ Kernel}%
 {Thin space.}%
 {}

\defgsymcs[colon]{\colonsym}%
 {}%
 {\LaTeX\ Kernel (Math Mode)}%
 {Medium space.}%
 {}

\defgcs{medspace}%
 {}%
 {\nxisty{amsmath} package}%
 {Medium space.}%
 {}

\defgsymcs[semicolon]{\semicolonsym}%
 {}%
 {\LaTeX\ Kernel (Math Mode)}%
 {Thick space.}%
 {}

\defgcs{thickspace}%
 {}%
 {\nxisty{amsmath} package}%
 {Thick space.}%
 {}

\defgcs{qquad}%
 {}%
 {\LaTeX\ Kernel}%
 {Horizontal spacing command (twice as wide as \nxglsni{quad}).}%
 {}

\defgsymcs[spacesym]{\textvisiblespace}%
 {}%
 {\LaTeX\ Kernel}%
 {(Backslash followed by space character.) Horizontal spacing command.}%
 {}

\defgcs{enspace}%
 {}%
 {\LaTeX\ Kernel}%
 {Horizontal spacing command (half as wide as \nxglsi{quad}).}%
 {}

\defgcs{quad}%
 {}%
 {\LaTeX\ Kernel}%
 {Horizontal spacing command equal to the current font's \protect\htmlref{em}{sec:length} value.}%
 {}

\defgcs{negthinspace}%
 {}%
 {\LaTeX\ Kernel}%
 {Negative thin space.}%
 {}

\defgcs{negmedspace}%
 {}%
 {\nxisty{amsmath} package}%
 {Negative medium space.}%
 {}

\defgcs{negthickspace}%
 {}%
 {\nxisty{amsmath} package}%
 {Negative thick space.}%
 {}

\defgcs{contentsname}%
 {}%
 {Classes or packages that define a table of contents}%
 {Text used for table of contents heading.}%
 {}

\defgcs{listfigurename}%
 {}%
 {Classes or packages that define a list of figures}%
 {Text used for list of figures heading.}%
 {}

\defgcs{listtablename}%
 {}%
 {Classes or packages that define a list of tables}%
 {Text used for list of tables heading.}%
 {}

\defgcs{bibname}%
 {}%
 {Report or book style classes that define a bibliography chapter}%
 {Text used for bibliography chapter heading. (See also
  \nxglsi{refname}.)}%
 {}

\defgcs{refname}%
 {}%
 {Article style classes that define a bibliography section}%
 {Text used for bibliography section heading. (See also
   \nxglsi{bibname}.)}%
 {}

\defgcs{indexname}%
 {}%
 {Classes or packages that define an index section}%
 {Text used for index heading.}%
 {}

\defgcs{figurename}%
 {}%
 {Classes or packages that define figures}%
 {Number prefix used in figure captions.}%
 {}

\defgcs{tablename}%
 {}%
 {Classes or packages that define tables}%
 {Number prefix used in table captions.}%
 {}

\defgcs{partname}%
 {}%
 {Classes or packages that define parts with a number prefix}%
 {Number prefix used in part headings.}%
 {}

\defgcs{chaptername}%
 {}%
 {Classes or packages that define chapters}%
 {Number prefix used in chapter headings.}%
 {}

\defgcs{appendixname}%
 {}%
 {Classes or packages that define chapters}%
 {Number prefix used in appendix headings.}%
 {}

\defgcs{abstractname}%
 {}%
 {Classes or packages that define an \nxglsni{env-abstract}
  environment}%
 {Text used in abstract heading.}%
 {}

\defgcs{addto}%
 {\marg{\meta{command}}\marg{\meta{code}}}
 {\nxisty{babel} package}
 {Adds \meta{code} to the definition of \meta{command}. (See also \nxglsni{appto}.)}
 {%
   \BeginArgList
     \csentryargitem{command} The command that needs patching.
     \csentryargitem{code} The additional code to add to the command
definition.
   \EndArgList
 }

\defgcs{appto}%
 {\marg{\meta{command}}\marg{\meta{code}}}
 {\nxisty{etoolbox} package}
 {Adds \meta{code} to the definition of \meta{command}.}
 {%
   \BeginArgList
     \csentryargitem{command} The command that needs patching.
     \csentryargitem{code} The additional code to add to the command definition.
   \EndArgList
 }

\defgcs{mathrm}%
 {\marg{\meta{maths}}}%
 {\LaTeX\ Kernel (Math Mode)}%
 {Renders \meta{maths} in the predefined maths serif font.}%
 {%
   \BeginArgList
     \csentryargitem{maths} The text on which to apply the font change.
   \EndArgList
 }

\defgcs{mathsf}%
 {\marg{\meta{maths}}}%
 {\LaTeX\ Kernel (Math Mode)}%
 {Renders \meta{maths} in the predefined maths sans-serif font.}%
 {%
   \BeginArgList
     \csentryargitem{maths} The text on which to apply the font change.
   \EndArgList
 }

\defgcs{mathtt}%
 {\marg{\meta{maths}}}%
 {\LaTeX\ Kernel (Math Mode)}%
 {Renders \meta{maths} in the predefined maths typewriter font.}%
 {%
   \BeginArgList
     \csentryargitem{maths} The text on which to apply the font change.
   \EndArgList
 }

\defgcs{mathit}%
 {\marg{\meta{maths}}}%
 {\LaTeX\ Kernel (Math Mode)}%
 {Renders \meta{maths} in the predefined maths italic font.}%
 {%
   \BeginArgList
     \csentryargitem{maths} The text on which to apply the font change.
   \EndArgList
 }

\defgcs{mathbf}%
 {\marg{\meta{maths}}}%
 {\LaTeX\ Kernel (Math Mode)}%
 {Renders \meta{maths} in the predefined maths bold font. (Doesn't
  work with numbers and nonalphabetical symbols.) See also \nxglsi{boldsymbol}.}%
 {%
   \BeginArgList
     \csentryargitem{maths} The text on which to apply the font change.
   \EndArgList
 }

\defgcs{boldsymbol}%
 {\marg{\meta{symbol}}}%
 {\nxisty{amsmath} package (Math Mode)}%
 {Like \nxglsi{mathbf} but also works for numbers and
   many nonalphabetical symbols. (See also \nxglsi{pmb}.)}%
 {%
   \BeginArgList
     \csentryargitem{symbol} The character or symbol on which to apply the font change.
   \EndArgList
 }

\defgcs{pmb}%
 {\marg{\meta{symbol}}}%
 {\nxisty{amsmath} package (Math Mode)}%
 {\dq{Poor man's bold.} Overlays multiple copies of the symbol to
   produce a bold effect for symbols that don't work with
   \nxglsi{boldsymbol}.}%
 {%
   \BeginArgList
     \csentryargitem{symbol} The symbol on which to apply the bold
      effect.
   \EndArgList
 }

\defgcs{mathcal}%
 {\marg{\meta{maths}}}%
 {\LaTeX\ Kernel (Math Mode)}%
 {Typesets its argument in the maths calligraphic font.
  Example: \nxglsni{beginmath}\protect\cmdname{mathcal}\marg{S}\nxglsni{endmath} produces
 \ensuremath{\mathcal{S}}.}%
 {%
   \BeginArgList
    \csentryargitem{maths} The maths to be displayed in a calligraphic
      font.
   \EndArgList
 }

\defgcs{mathbb}%
 {\marg{\meta{maths}}}%
 {\nxisty{amsfonts} package (Math Mode)}%
 {Typesets its argument in the blackboard bold font.
  Example: \nxglsni{beginmath}\protect\cmdname{mathbb}\marg{R}\nxglsni{endmath} produces
  \ensuremath{\mathbb{R}}.}%
 {%
   \BeginArgList
    \csentryargitem{maths} The maths to be displayed in a blackboard
      bold font.
   \EndArgList
 }

\defgcs{mathfrak}%
 {\marg{\meta{maths}}}%
 {\nxisty{amsfonts} package (Math Mode)}%
 {Typesets its argument in Euler Fraktur letters.
  Example: \nxglsni{beginmath}\protect\cmdname{mathfrak}\marg{U}\nxglsni{endmath} produces
  \ensuremath{\mathfrak{U}}.}%
 {%
   \BeginArgList
    \csentryargitem{maths} The maths to be displayed in Euler Fraktur
      letters.
   \EndArgList
 }

\defgcs{DeclareMathOperator}%
 {\marg{\meta{cmd}}\marg{\meta{operator-name}}}%
 {\nxisty{amsmath} package (Preamble Only)}%
 {Defines a new maths operator. The starred version allows limits.}%
 {%
   \BeginArgList
    \csentryargitem{cmd} The name of the new operator command (must
      begin with a backslash).
    \csentryargitem{operator name} The name of the maths operator.
   \EndArgList
 }

\defgcs{frac}%
 {\marg{\meta{numerator}}\marg{\meta{denominator}}}%
 {\LaTeX\ Kernel (Math Mode)}%
 {Displays a fraction.}%
 {%
    \BeginArgList
     \csentryargitem{numerator} The numerator (above the line).
     \csentryargitem{denominator} The denominator (below the line).
    \EndArgList
 }

\defgcs{cfrac}%
 {\oarg{\meta{pos}}\marg{\meta{numerator}}\marg{\meta{denominator}}}%
 {\nxisty{amsmath} (Math Mode)}%
 {Displays a continued fraction.}%
 {%
    \BeginArgList
     \csentryargitem{pos} May be \texttt{l} (left) or \texttt{r}
      (right).
     \csentryargitem{numerator} The numerator (above the line).
     \csentryargitem{denominator} The denominator (below the line).
    \EndArgList
 }

\defgcs{sqrt}%
 {\oarg{\meta{order}}\marg{\meta{operand}}}%
 {\LaTeX\ Kernel (Math Mode)}%
 {Displays a root.}%
 {%
    \BeginArgList
     \csentryargitem{order} The order of the root. (If omitted, a
       square root).
     \csentryargitem{operand} The operand.
    \EndArgList
 }

\defgcs{arabic}%
 {\marg{\meta{counter}}}%
 {\LaTeX\ Kernel}%
 {Displays counter value as an Arabic number. (1, 2, 3, \ldots)}%
 {%
    \BeginArgList
     \csentryargitem{counter} The name of the counter.
    \EndArgList
 }

\defgcs{roman}%
 {\marg{\meta{counter}}}%
 {\LaTeX\ Kernel}%
 {Displays counter value as a lower case Roman number. (i, ii, iii, \ldots)}%
 {%
    \BeginArgList
     \csentryargitem{counter} The name of the counter.
    \EndArgList
 }

\defgcs{Roman}%
 {\marg{\meta{counter}}}%
 {\LaTeX\ Kernel}%
 {Displays counter value as an upper case Roman number. (I, II, III, \ldots)}%
 {%
    \BeginArgList
     \csentryargitem{counter} The name of the counter.
    \EndArgList
 }

\defgcs{alph}%
 {\marg{\meta{counter}}}%
 {\LaTeX\ Kernel}%
 {Displays counter value as a lower case letter. (a, b, c, \ldots, z)}%
 {%
    \BeginArgList
     \csentryargitem{counter} The name of the counter.
    \EndArgList
 }

\defgcs{Alph}%
 {\marg{\meta{counter}}}%
 {\LaTeX\ Kernel}%
 {Displays counter value as an upper case letter. (A, B, C, \ldots, Z)}%
 {%
    \BeginArgList
     \csentryargitem{counter} The name of the counter.
    \EndArgList
 }

\defgcs{fnsymbol}%
 {\marg{\meta{counter}}}%
 {\LaTeX\ Kernel}%
 {Displays counter value as footnote symbol. (\makeimg{asterisk
  dagger double dagger section mark
  paragraph mark double bar double asterisk two single daggers
  two double daggers}{\footnotesymbols})}%
 {%
    \BeginArgList
     \csentryargitem{counter} The name of the counter.
    \EndArgList
 }

\defgenv{bfseries}%
 {}%
 {\LaTeX\ Kernel}%
 {Typesets the environment contents in a bold font.}%
 {}

\defgenv{itshape}%
 {}%
 {\LaTeX\ Kernel}%
 {Typesets the environment contents in an italic font.}%
 {}

\defgenv{em}%
 {}%
 {\LaTeX\ Kernel}%
 {Typesets the environment contents in an emphasized font. (Switches
  to italic\slash slanted if the surrounding font is upright, or switches to
  upright if the surrounding font is italic\slash slanted.)}%
 {}

\defgenv{itemize}%
 {}%
 {\LaTeX\ Kernel}%
 {Unordered list.}%
 {}

\defgenv{enumerate}%
 {}%
 {\LaTeX\ Kernel}%
 {Ordered list.}%
 {}

\defgenv{description}%
 {}%
 {Most class files}%
 {Labelled list.}%
 {}

\defgenv{document}%
 {}%
 {\LaTeX\ Kernel}%
 {The body of the document.}%
 {}

\defgenv{tabular}%
 {\oarg{\meta{v-pos}}\marg{\meta{column specifiers}}}%
 {\LaTeX\ Kernel (Text Mode)}%
 {%
   Environment for lining things up in rows and columns.
   Use \nxglsi{env-array} for math mode.
 }%
 {%
    \BeginArgList
     \csentryargitem{v-pos} Vertical alignment of the entire
      environment with respect to the surrounding baseline. May be
      one of \texttt{t} (top), \texttt{b} (bottom) or \texttt{c}
      (centered).
     \csentryargitem{column specifiers} Indicates how to align the
     columns.
    \EndArgList
 }

\defgenv{array}%
 {\oarg{\meta{v-pos}}\marg{\meta{column specifiers}}}%
 {\LaTeX\ Kernel (Math Mode)}%
 {%
   Environment for lining things up in rows and columns.
   Use \nxglsi{env-tabular} for text mode.
 }%
 {%
    \BeginArgList
     \csentryargitem{v-pos} Vertical alignment of the entire
      environment with respect to the surrounding baseline. May be
      one of \texttt{t} (top), \texttt{b} (bottom) or \texttt{c}
      (centered).
     \csentryargitem{column specifiers} Indicates how to align the
     columns.
    \EndArgList
 }

\defgenv{align}
 {}
 {\nxisty{amsmath} package}
 {Used for numbered aligned equations.}
 {}

\defgenv{align*}
 {}
 {\nxisty{amsmath} package}
 {Used for unnumbered aligned equations.}
 {}

\defgenv{minipage}%
 {\oarg{\meta{pos}}\oarg{\meta{height}}\marg{\meta{width}}}%
 {\LaTeX\ Kernel}%
 {Makes a box with line-wrapped contents. (See also \nxglsi{parbox}.)}%
 {%
   \BeginArgList
    \csentryargitem{pos} The vertical alignment of the box relative to
     the surrounding text. (Centred if omitted.)
    \csentryargitem{height} The height of the box. (If omitted the
     height is the natural height of the contents of the box.)
    \csentryargitem{width} The width of the box.
   \EndArgList
 }

\defgenv{abstract}%
 {}%
 {Most article- or report-style classes, such as \nxicls{scrartcl} or
  \nxicls{scrreprt}. Not usually defined in book-style classes, such
  as \nxicls{scrbook}, but is defined in \nxicls{memoir}}%
 {Displays its contents as an abstract.}%
 {}

\defgenv{thebibliography}%
 {\marg{\meta{widest entry label}}}%
 {Most classes that define sectioning commands}%
 {Bibliographic list. (See also \nxglsi{bibitem} and
  \nxglsi{cite}).}%
 {%
   \BeginArgList
    \csentryargitem{widest entry label} The widest label in the
     bibliography list.
   \EndArgList
 }

\defgenv{otherlanguage}
 {\marg{\meta{language name}}}%
 {\nxisty{babel} package}%
 {Within the environment contents, predefined textual elements, such
  as the date given by \nxglsi{today} or prefixes like \dq{Chapter},
  are set to those supplied by the given language.}%
 {%
   \BeginArgList
    \csentryargitem{language} The language identifier.
   \EndArgList
 }

\defgenv{figure}
 {\oarg{\meta{placement}}}%
 {Most classes that define sectioning commands}%
 {%
   Floats the contents to the nearest location according to the
   preferred placement options, if possible. Within the environment,
   \nxglsi{caption} may be used one or more times, as required.
   The caption will usually include the prefix given by
   \nxglsni{figurename}.
 }%
 {%
   \BeginArgList
    \csentryargitem{placement} The preferred placement.
   \EndArgList
 }

\defgenv{sidewaysfigure}%
 {}%
 {\nxisty{rotating} package}%
 {Like the \nxglsi{env-figure} environment but rotates the entire
  figure (including caption) sideways.}%
 {}

\defgenv{table}
 {\oarg{\meta{placement}}}%
 {Most classes that define sectioning commands}%
 {%
   Floats the contents to the nearest location according to the
   preferred placement options, if possible. Within the environment,
   \nxglsi{caption} may be used one or more times, as required.
   The caption will usually include the prefix given by
   \nxglsni{tablename}.
 }%
 {%
   \BeginArgList
    \csentryargitem{placement} The preferred placement.
   \EndArgList
 }

\defgenv{sidewaystable}%
 {}%
 {\nxisty{rotating} package}%
 {Like the \nxglsi{env-table} environment but rotates the entire
  table (including caption) sideways.}%
 {}

\defgenv{subfigure}%
 {\oarg{\meta{pos}}\marg{\meta{width}}}%
 {\nxisty{subcaption} package}%
 {Used to form a subfigure within a \nxglsi{env-figure}
  environment. The \nxglsi{caption} command may be used in this
  environment to produce a subcaption.}%
 {%
   \BeginArgList
    \csentryargitem{pos} The vertical alignment of the box relative to
     the surrounding text. (Centred if omitted.)
    \csentryargitem{width} The width of the box.
   \EndArgList
 }

\defgenv{subtable}%
 {\oarg{\meta{pos}}\marg{\meta{width}}}%
 {\nxisty{subcaption} package}%
 {Used to form a subtable within a \nxglsni{env-table}
  environment. The \nxglsi{caption} command may be used in this
  environment to produce a subcaption.}%
 {%
   \BeginArgList
    \csentryargitem{pos} The vertical alignment of the box relative to
     the surrounding text. (Centred if omitted.)
    \csentryargitem{width} The width of the box.
   \EndArgList
 }

\defgenv{equation}%
 {}%
 {\LaTeX\ Kernel}%
 {Displays its contents as a single-lined numbered equation.}%
 {}

\defgenv{displaymath}%
 {}%
 {\LaTeX\ Kernel}%
 {Displays its contents as a single-lined unnumbered equation.}%
 {}

\defgenv{math}%
 {}%
 {\LaTeX\ Kernel}%
 {Sets its contents in in-line math mode.}%
 {}

\defgenv{cases}%
 {}%
 {\nxisty{amsmath} package (Math Mode)}%
 {%
   Like the \nxglsi{env-array} environment, but adds a left
   brace start delimiter and an invisible end delimiter.
 }%
 {}

\defgenv{matrix}%
 {}%
 {\nxisty{amsmath} package (Math Mode)}%
 {%
   Like the \nxglsi{env-array} environment, but doesn't have an
   argument.
 }%
 {}

\defgenv{pmatrix}%
 {}%
 {\nxisty{amsmath} package (Math Mode)}%
 {%
   Like the \nxglsi{env-array} environment, but doesn't have an
   argument and adds round bracket delimiters.
 }%
 {}

\defgenv{smallmatrix}%
 {}%
 {\nxisty{amsmath} package (Math Mode)}%
 {%
   Like the \nxglsi{env-array} environment but doesn't have an
   argument and is designed for in-line maths.
 }%
 {}

\defgenv{bmatrix}%
 {}%
 {\nxisty{amsmath} package (Math Mode)}%
 {%
   Like the \nxglsi{env-array} environment, but doesn't have an
   argument and adds square bracket delimiters.
 }%
 {}

\defgenv{Bmatrix}%
 {}%
 {\nxisty{amsmath} package (Math Mode)}%
 {%
   Like the \nxglsi{env-array} environment, but doesn't have an
   argument and adds curly brace delimiters.
 }%
 {}

\defgenv{vmatrix}%
 {}%
 {\nxisty{amsmath} package (Math Mode)}%
 {%
   Like the \nxglsi{env-array} environment, but doesn't have an
   argument and adds single vertical bar delimiters.
 }%
 {}

\defgenv{Vmatrix}%
 {}%
 {\nxisty{amsmath} package (Math Mode)}%
 {%
   Like the \nxglsi{env-array} environment, but doesn't have an
   argument and adds double vertical bar delimiters.
 }%
 {}

\defgcs{raggedright}%
 {}%
 {\LaTeX\ Kernel}%
 {Ragged-right paragraph justification.}%
 {}

\defgcs{raggedleft}%
 {}%
 {\LaTeX\ Kernel}%
 {Ragged-left paragraph justification.}%
 {}

\defgcs{frontmatter}
 {}%
 {Most book-style classes, such as \nxicls{scrbook}}%
 {Switches to lower case Roman numeral page numbering. Also suppresses
   chapter and section numbering, but still adds unstarred sectional
   units to the table of contents. (See also \nxglsi{mainmatter}
   and \nxglsi{backmatter}.)}%
 {}

\defgcs{mainmatter}
 {}%
 {Most book-style classes, such as \nxicls{scrbook}}%
 {Switches to Arabic page numbering and enables
   chapter and section numbering. (See also
   \nxglsi{frontmatter} and \nxglsi{backmatter}.)}%
 {}

\defgcs{backmatter}
 {}%
 {Most book-style classes, such as \nxicls{scrbook}}%
 {Suppresses chapter and section numbering, but still adds unstarred
  sectional units to the table of contents. (See also \nxglsi{frontmatter}
  and \nxglsi{mainmatter}.)}%
 {}

\defgenv{dinglist}
 {\marg{\meta{number}}}
 {\nxisty{pifont} package}
 {A list where the item marker is given by character \meta{number}
  in the Zapf Dingbats font.}
 {%
   \BeginArgList
    \csentryargitem{number} The character code for the item marker.
   \EndArgList
 }

\defgcs{notag}%
 {}%
 {\nxisty{amsmath} package}%
 {Suppresses equation numbering for the current row in environments such as
   \nxglsi{env-align}.}%
 {}

\defgcs{tag}%
 {\marg{\meta{tag}}}%
 {\nxisty{amsmath} package}%
 {Overrides equation numbering for the current row in environments such as
   \nxglsi{env-align}.}%
 {%
    \BeginArgList
     \csentryargitem{tag} The replacement for the equation number.
    \EndArgList
 }

\defgcs{ding}%
 {\marg{\meta{n}}}%
 {\nxisty{pifont} package}%
 {Inserts PostScript ZapfDingbats character with code \meta{n},
   which must be an integer.}%
 {%
   \BeginArgList
    \csentryargitem{n} The decimal code of the required character.
   \EndArgList
 }

\defgcs{color}%
 {\oarg{\meta{model}}\marg{\meta{specs}}}%
 {\nxisty{color} and \nxisty{xcolor} packages}%
 {A declaration that switches the current foreground colour to the given specification.}%
 {%
    \BeginArgList
      \csentryargitem{model} The colour model, for example
      \texttt{rgb} or \texttt{named}.
      \csentryargitem{specs} The colour specification for the
      given model. For example, if the \texttt{rgb} model is chosen,
      \meta{specs} must be a comma-separated list of three numbers
      each between~0 and~1.
    \EndArgList
 }

\defgcs{textcolor}%
 {\oarg{\meta{model}}\marg{\meta{specs}}\marg{\meta{text}}}%
 {\nxisty{color} and \nxisty{xcolor} packages}%
 {Sets \meta{text} with the foreground colour according to the given \meta{specs}.}%
 {%
    \BeginArgList
      \csentryargitem{\meta{model}} The colour model, for example
      \texttt{rgb} or \texttt{named}.
      \csentryargitem{\meta{specs}} The colour specification for the
      given model. For example, if the \texttt{rgb} model is chosen,
      \meta{specs} must be a comma-separated list of three numbers
      each between~0 and~1.
      \csentryargitem{\meta{text}} The text to be displayed in the
      given colour.
    \EndArgList
 }

\defgcs{tabcolsep}
 {}%
 {\LaTeX\ Kernel}%
 {\nxGls{length} register specifying half the gap between
   columns in a \nxglsni{env-tabular} environment.}%
 {}

\defgcs{arraycolsep}
 {}%
 {\LaTeX\ Kernel}%
 {\nxGls{length} register specifying half the gap between
   columns in an \nxglsni{env-array} environment.}%
 {}


\defgcs{labelformat}%
 {\marg{\meta{ctr}}\marg{\meta{defn}}}%
 {\nxisty{fncylab} package}%
 {Defines how the label for the counter \meta{ctr} should be
 formatted. The definition \meta{defn} should use \texttt{\nxglsi{hashchar}1} to
 indicate the label value.}%
 {%
   \BeginArgList
     \csentryargitem{ctr} The name of the counter.
     \csentryargitem{defn} How to display the value of the counter.
   \EndArgList
 }

\defgcs{thefootnote}%
 {}%
 {\LaTeX\ Kernel}%
 {Displays the current value of the \nxicounter{footnote} counter}%
 {}

\defgcs{thechapter}%
 {}%
 {\LaTeX\ Kernel}%
 {Displays the current value of the \nxicounter{chapter} counter}%
 {}

\defgcs{thesection}%
 {}%
 {\LaTeX\ Kernel}%
 {Displays the current value of the \nxicounter{section} counter}%
 {}

\defgcs{thepage}%
 {}%
 {\LaTeX\ Kernel}%
 {Displays the current value of the \nxicounter{page} counter}%
 {}

\defgcs{thefigure}%
 {}%
 {\LaTeX\ Kernel}%
 {Displays the current value of the \nxicounter{figure} counter}%
 {}

\defgcs{thetable}%
 {}%
 {\LaTeX\ Kernel}%
 {Displays the current value of the \nxicounter{table} counter}%
 {}

\defgcs{ignorespaces}
 {}
 {\LaTeX\ Kernel}
 {Used in begin environment code to suppress any spaces occurring 
  at the start of the environment (see also \nxglsni{ignorespacesafterend}).}
 {}

\defgcs{ignorespacesafterend}%
 {}%
 {\LaTeX\ Kernel}%
 {Used in end environment code to suppress any spaces following the
  end of the environment.}%
 {}

\defgcs{textbar}%
 {}%
 {\LaTeX\ Kernel}%
 {Vertical bar \textbar\ symbol.}%
 {}

\defgcs{the}%
 {\meta{register}}%
 {\LaTeX\ Kernel}%
 {Displays the value of the given register (such as a \nxglsi{length} register).
  Not to be confused with \cmdname{the}\meta{ctr} commands, such as \nxglsni{thefigure}.}%
 {%
   \BeginArgList
     \csentryargitem{register} The name of the register.
   \EndArgList
 }

\defgcs{substack}%
 {\marg{\meta{maths}}}%
 {\nxisty{amsmath} package}%
 {Can be used to produce a multiline subscript or superscript. Lines
  are separated using \nxglsi{tab.dbbackslashchar}.}%
 {%
   \BeginArgList
    \csentryargitem{maths} The subscript or superscript maths with rows
      ended using \nxglsi{tab.dbbackslashchar}
   \EndArgList
 }

\defgcs{DeclareCaptionLabelFormat}
 {\marg{\meta{name}}\marg{\meta{code}}}
 {\nxisty{caption}}
 {Used to defined your own caption label formats.}
 {%
   \BeginArgList
     \csentryargitem{name} The name used to identify this new format
     \csentryargitem{code} The code that defines this format
     (\nxglsi{hashchar}1 gets replaced with the name and \nxglsni{hashchar}2
     gets replaced with the reference number.)
   \EndArgList
 }

\defgcs{captionsetup}
 {\oarg{\meta{float type}}\marg{\meta{options}}}
 {\nxisty{caption} package}
 {Used to set up the options affecting float captions.}
 {%
   \BeginArgList
     \csentryargitem{float type} The float type.
     \csentryargitem{options} A \meta{key}=\meta{value}
     comma-separated list of options.
   \EndArgList
 }



\newacr{ctan}{CTAN}
 {the Comprehensive \TeX\ Archive Network}
 {http://mirror.ctan.org/}

\newacr{tug}{TUG}%
 {\TeX\ User Group}%
 {http://tug.org/}

\newacr{uktug}{UK TUG}%
 {UK \TeX\ User Group}%
 {http://uk.tug.org/}

\newacr{ukfaq}{UK FAQ}%
 {UK List of \TeX\ Frequently Asked Questions}%
 {http://www.tex.ac.uk/faq}

\newacrx{tds}{TDS}{\TeX\ Directory Structure}

\newacrnocite{gui}{GUI}{graphical user interface}

\definecolor{midgray}{gray}{0.4}

\indexpreamble{\latexhtml{Page numbers}{Locations} in \textit{italic} indicate the primary
reference. \latexhtml{Page numbers}{Locations} in \textbf{bold} indicate the entry
definition in the summary.}

\newcommand{\xtableref}[1]{\objectref{Table}{#1}}

% \figconts[graphics opts]{image}{caption}{label}
\newcommand{\figconts}[4][]{%
  \begin{makeimage}\end{makeimage}\figuretop{#4}
  \centering
  \incPgfOrGraphics[#1]{#2}%
  #3% caption
  \label{#4}%
}

\newcommand{\latexbook}{\emph{\LaTeX: A Document Preparation System}~\cite{lamport94}}
\newcommand{\latexcomp}{\emph{The \LaTeX\ Companion}~\cite{goossens94}}
\newcommand{\latexguide}{\emph{A Guide to \LaTeX}~\cite{kopka95}}
\newcommand{\latexgraphic}{\emph{The \LaTeX\ Graphics Companion}~\cite{goossens97}}
\newcommand{\latexweb}{\emph{The \LaTeX\ Web Companion}~\cite{goossens99}}

\newcommand{\baseurl}{http://www.dickimaw-books.com}
\newcommand{\packageurl}{\baseurl/latex/packages}
\newcommand{\exerciseurl}{\baseurl/latex/novices/html/exercises}

\newcommand{\pderiv}[2]{\frac{\partial #1}{\partial #2}}
\newcommand{\e}{\mathrm{e}}

\newcommand{\koma}{KOMA-Script}

\newcommand*{\backcovertext}{If you choose to buy a copy of this book, 
Dickimaw Books asks for your support through buying the Dickimaw
Books edition to help cover costs.}

\copyrighttext{%
  Copyright \textcopyright\ 2004  Nicola L.~C. Talbot

  Permission is granted to copy, distribute and/or modify this
  document under the terms of the GNU Free Documentation License,
  Version 1.2 or any later version published by the Free Software
  Foundation; with no Invariant Sections, no Front-Cover Texts, and
  one Back-Cover Text: \dq{\backcovertext} A copy of the license is included in
  the section entitled \htmlref{\dq{GNU Free Documentation
  License}}{sec:fdl}.

  \doifnotbook
  {%
     The base URL for this document is: \url{\baseurl/latex/novices/}
  }%
}

\renewcommand{\glossaryname}{Summary}

\begin{document}
\maketitle

\frontmatter
\setcounter{tocdepth}{2}
\setnode{contents}
\tableofcontents
\setnode{listoffigures}
\listoffigures
\setnode{listoftables}
\listoftables

\setnode{listofexercises}
\listofexercises

\mainmatter

%%%%%%%%%%%%%%%% INTRODUCTION %%%%%%%%%%%%%%%%%%%%%%%%%%%%

\setnode{introduction}
\chapter{Introduction}
\label{ch:intro}

The aim of this book is to introduce \LaTeX\ to a non-technical
person.  \LaTeX\ is excellent for producing professional looking
documents, however it is a \emph{language} not a word 
processor\faq{Why is TeX not a WYSIWYG system?}{notWYSIWYG}, so
it can take a bit of getting used to, particularly if you have never
had any experience using programming languages.  

\LaTeX\ does take a while to learn, so why should you use it? Here are
a few reasons but it is not an exhaustive list:

\LaTeX\ is far better at typesetting mathematical equations
than word processors. I~wrote my Ph.D. thesis back in the days of
\LaTeX2.09 (the old version of \LaTeX) and given the high quantity of
mathematics that I~had to typeset, it would have taken me considerably
longer to write it in a word processor, and the resulting document
wouldn't have looked nearly as good. Even Microsoft have
acknowledged \glsterm{tex}'s high-quality mathematical
typography~\cite{sargentiii}.

\xminisec{Example:}
Here's an equation taken from some kernel survival analysis:
\begin{result}[latexintromaths.html]
\newcommand{\me}{\mathrm{e}}
\newcommand{\pfrac}[2]{\frac{\partial #1}{\partial #2}}
\[
\pfrac{^2\mathcal{L}}{{z_i^\rho}^2} =
-\pfrac{\rho_i}{z_i^\rho}\left(\pfrac{v_i}{\rho_i}\frac{\me^{v_i}}{1-\me^{v_i}}
+v_i\frac{\me^{v_i}\pfrac{v_i}{\rho_i}(1-\me^{v_i})+\me^{2v_i}\pfrac{v_i}{\rho_i}}{(1-\me^{v_i})^2}\right)
\]
\end{result}
(You can find out how to create this equation \latexhtml{on
page~\pageref{introeq} in \sectionref{sec:delimiters}}{\htmlref{later}{introeq}}.)

That's all very well and good if you want to typeset some equations,
but if your work doesn't involve maths, does that mean that \LaTeX\
is not for you? Although I~am a mathematician, I~have written plenty
of documents with no maths in at all, including prose, poetry,
newsletters, posters and brochures, but I~still opt for \LaTeX\
because using \LaTeX\ ensures consistent formatting, and the style
of the document can be completely changed by simply using a
different \glsterm{cls}, or loading additional \htmlref{packages}{sec:packages}. This
means that I~can concentrate on writing the document, rather than
worrying about how it will look. It also means that if, after having
written a 200 page document, I~then find that I~need to change all
the figure captions so that they are labelled \dq{Fig} instead of
\dq{Figure}, all I~need to do is edit a single line, rather than
going through 200 pages to individually edit every single figure
caption.\footnote{Sure, you could use a search and replace function,
but a sweeping replace-all can have unexpected side effects. For
example, your document may include the sentence, \dq{Figures from
the last quarter showed improvement}, which would get changed to,
\dq{Figs from the last quarter showed improvement}.}

Serious fiction writers are taught never to remind the reader that
they're reading a book. Poor formatting is just as much a reminder
of this as authorial intrusion.

\LaTeX\ makes it very easy to cross-reference chapters, sections,
equations, figures, tables etc, and it also makes it very easy to
generate a table of contents, list of figures, list of tables,
index, glossary\footnote{Glossaries are covered in \latexthesis.}\ and 
bibliography. You don't need to worry about numbering anything,
as this is done automatically, which means that you can insert new
sections or swap sections around without having to worry about
updating all the section numbering etc. Furthermore, if you use
\BiBTeX\footnote{Automating bibliographies is covered in
\latexthesis.}\ in combination with \LaTeX, and you have, say, 100
or more citations, it doesn't matter if you are then told that the
citations have to be re-ordered (say, in order of citation rather
than alphabetically). All that is required is a minor edit to change
the appropriate style file rather than ploughing through the entire
document changing all the citations by hand.

When you are editing a document using a word processor, the
word processor has to work out how to reformat the document every time
you type something.  If you have a large document with a great many
inserted objects (such as figures and equations), the response to
keyboard input can become very slow.  You may find that after typing
a few words you will have to wait until the computer catches up
before you can see what you have typed.  With \LaTeX\ you type in your
code using an ordinary text editor.  The document 
doesn't get formatted until you pass it to \LaTeX, which means that 
you are not slowed down by constant reformatting.

Lastly, there's the fact that \LaTeX\ follows certain typographical
rules, so you can leave most of the typesetting to \LaTeX. You
rarely need to worry about minor things such as
\glstext{intersentencespacing}. The default is English spacing, but
if you have a publisher who disapproves of this, you can switch if
off with a single command. (See \sectionref{sec:intersentencespacing}.)

\LaTeX\ will also automatically deal with
f-ligatures.\footnote{Ligatures can be suppressed using the
\isty{microtype} package if necessary} That is, if any of the
following combination of letters are found: \texttt{fl},
\texttt{ffl}, \texttt{ff}, \texttt{fi}, \texttt{ffi}, they will
automatically be converted into the corresponding ligatures:
% Using CMR as the ligatures are more noticeable than with anttor.
% make latex2html convert this text into an image
\makeimg{fl ligature}{\cmr{fl}},
\makeimg{ffl ligature}{\cmr{ffl}},
\makeimg{ff ligature}{\cmr{ff}},
\makeimg{fi ligature}{\cmr{fi}},
\makeimg{ffi ligature}{\cmr{ffi}}.
  Note the difference between
\makeimg{fl-ligature u ffi-ligature e r}{\cmr{fluffier}} (2
ligatures) and \makeimg{fluffier}{\cmr{f{}l{}uf{}f{}i{}er}} (no ligatures).  These
points may seem minor but they all contribute towards the impact of
the entire document.  When writing technical documents, the
presentation as well as the content is important.  All too often
examiners or referees are put off reading a document because it is
badly formatted.  This provokes an immediate negative reaction and
provides little desire to look favourably upon your work.

To give you an idea of what you can do with \LaTeX, this book was
written in \LaTeX.\footnote{The source code is available at
\url{\baseurl/latex/novices/}, but it really is
\emph{not} the place to start if you are a beginner, as it contains
\LaTeX\ and \glsterm{perl} code beyond the scope of this tutorial.}  The PDF
versions (including \ifbookorother{this}{the} paperback version) were generated using PDF\LaTeX\ and \appname{makeindex} and
the HTML version was generated using the
\latexhtml
 {\LaTeX2HTML\footnote{\url{http://www.latex2html.org/}}}%
 {\htmladdnormallink{\LaTeX2HTML}{http://www.latex2html.org/}}%
\faq{Conversion from (La)TeX to HTML}{LaTeX2HTML} converter. 

For more reasons as to why you might want to use \LaTeX\ instead
of a word processor, have a look at
\ifbookorother
  {\url{http://www.ctan.org/what_is_tex.html\#whytex}.}%
  {\htmladdnormallink{Why TeX?}{http://www.ctan.org/what_is_tex.html\#whytex}}

\setnode{texdoc}
\section{Class and Package Documentation}
\label{sec:texdoc}

There are hundreds of \htmlref{classes}{sec:cls} and
\htmlref{packages}{sec:packages} available on \gls{ctan}. These are
made available by many volunteers. Some provide detailed
documentation to accompany their contribution, while others only
provide a few notes in a README file or comments in the source
files. This book only provides an introductory
look at a small selection of these contributions. If you want
further details on how to use a particular class or package you
should check the documentation that accompanies it. You can use the
\iappname{texdoc} application to search for the documentation. This
is a command line application, which means you need a terminal or
command prompt (see \sectionref{sec:terminal}).

To use \appname{texdoc}, you need to
type (at the command prompt) \appname{texdoc} followed by a space followed by the name of
the class or package you want information about. 
For example, to read the \icls{memoir}
documentation, type the following at the \htmlref{command prompt}{sec:terminal} (press the
return/enter key \enter\ at the end of the line):
\begin{verbatim}
texdoc memoir
\end{verbatim}
Some packages come with more than one set of documentation. For
example, the \isty{glossaries} package comes with the main user
manual, a short guide for beginners and the documented code for
advanced users. Just doing
\begin{verbatim}
texdoc glossaries
\end{verbatim}
will display the advanced documented code. To list all available
documentation for a package, use the \texttt{-l} option:
\begin{verbatim}
texdoc -l glossaries
\end{verbatim}
Then type the number corresponding to the file you want to view. If
you can remember the file name (for example \texttt{glossaries-user}) 
you can type that next time you want to view it:
\begin{verbatim}
texdoc glossaries-user
\end{verbatim}
There is also a \glsterm{perl}/Tk-based \gls{gui} called \iappname{texdoctk},
which is distributed with \reportlinebreak\texdistro{TeX~Live},
that you can use instead of \appname{texdoc} if you can't work out
how to use a \glsterm{terminal} or prefer a \gls*{gui} approach.

Failing that, you can also check on \gls{ctan} using the URL
\texttt{ctan.org\slash pkg\slash}\meta{name}, where \meta{name} is the name of
the package or class. For example, if you want to look up the
documentation for the \icls{memoir} package, you can find it at
\url{http://ctan.org/pkg/memoir} or go to
\url{http://mirror.ctan.org/} and search for the package or class.

Another alternative recently made available is to use the URL
\texttt{texdoc.net\slash pkg\slash}\meta{name}. For example, 
\url{http://texdoc.net/pkg/memoir} will 
fetch the documentation for the \icls{memoir} class.

However, it's better to use \appname{texdoc} or \appname{texdoctk} to 
read the documentation installed with the
class or package on your computer to ensure it
matches the installed class or package version.

Note that it is important to remember that the \glsterm{tex} world is mostly
supported by volunteers. \gls{ctan} itself is maintained by a very small
group (currently two people). It's not like a commercial company
with 24/7 support and hundreds of paid employees constantly updating
the software. At its core, \TeX\ is a community effort. While some
volunteers actively maintain and update their classes or packages,
some people move on to other things and stop maintaining their work.
Occasionally, if the class or package is popular, someone else
might take over maintenance. There is no dedicated helpdesk to go
to, but there are many ways of getting help, see
\appendixref{ch:help}

\setnode{overview}
\section{Overview}

This document is structured as follows:

\begin{description}

\item[\chapterref{ch:def}]
defines terms that will be used throughout this document.  I
strongly suggest that you look through this
\latexhtml{chapter}{section} before you start so that you understand
the terminology used in this document. At the very least, you should
read the first part that details how corresponding input and output
is displayed in this document\dash you need to understand the
difference between \dq{input} (source code) and \dq{output} (how the
source code will appear in the typeset document).\html{ Note that it
is not always possible to reproduce an exact replica of the output
in the HTML version, so if in doubt, look at one of the PDF versions
instead, which can be downloaded from this 
book's \htmladdnormallink{home page}{../index.html}.}

\item[\chapterref{ch:tex2pdf}] details the software that you will
need to use \LaTeX\ and describes how to use the software.

\item[\chapterref{ch:simpledoc}] shows you how to create a very
basic document.  

\item[\chapterref{ch:sections}] 
shows you how to create chapters and other sectional units so that
you end up with a fully structured document.

\item[\chapterref{ch:graphicx}]
shows you how to include external image files and how to scale and
rotate text.

\item[\chapterref{ch:floats}]
describes how to create figures and tables.

\item[\chapterref{ch:newcom}]
describes how to define your own commands, and redefine existing
commands.

\item[\chapterref{ch:maths}] describes how
to typeset mathematics.  

\item[\chapterref{ch:newenv}] describes how to define new
environments.  

\item[\chapterref{ch:counters}]
discusses how numbers are stored in counters, how to change their
values, and how to define your own counter.

\item[\appendixref{ch:installsty}] shows you how
to download and install additional
packages that weren't installed with your \TeX\ distribution.

\item[\appendixref{ch:errors}]
documents possible errors you may encounter, and gives advice on how
to fix them.

\item[\appendixref{ch:help}] gives pointers on where to go for help.

\end{description}

Throughout this document there are pointers to related topics
in the \gls{ukfaq}.
\latexhtml{These are displayed in the margin in
square brackets, as illustrated on the 
\ifbookorother{\ifthispageodd{right}{left}}{right}.}{%
These are displayed in the text like this:}%
\faq{What is LaTeX?}{latex} You may find these resources useful in
answering related questions that are not covered in this 
book.
\begin{latexonly}
\doifbook{%
  To find the resources, go to \url{http://www.tex.ac.uk/faq} and
  either look for the question title in the list, or enter a keyword
  in the search field.%
}

This book and associated files, including  
solutions to the exercises, are available on-line at:
\url{\baseurl/latex/novices/}. 
\doifnotbook
{%
  The links in this document are colour-coded: internal links are 
  blue, external links are magenta.
}%
\end{latexonly}

\setnode{reading}
\section{Recommended Reading}

This document is designed as an introductory text, not a
comprehensive guide.  For further reading try some of the
following\faq{Books on LaTeX}{latex-books}:

\latexbook\ is the user guide and reference manual for \LaTeX, and
is a good basic text for anyone starting out, however it doesn't
cover AMS\TeX\faq{What are the AMS packages?}{AMSpkg}, so anyone who
needs to typeset more than basic mathematics may prefer either
\latexguide\ or \latexcomp.  Both these books cover AMS\TeX,
\iBiBTeX\ and \appname{makeindex}.

In the same series as \emph{The \LaTeX\ Companion}, there is also 
\latexgraphic\ which details how to illustrate
documents with \LaTeX\ and PostScript, including a chapter on colour
(coloured text, background, tables and slides).  This is recommended
to anyone who is contemplating heavy use of graphics, but you do need
a basic knowledge of \LaTeX\ before delving into it.  

The final book in the \dq{Companion} series which you may find useful
is \latexweb. This is recommended for those interested in creating
documents for the web, either as HTML or PDF.  It details how
to convert \LaTeX\ documents into HTML using various applications
such as \appname{LaTeX2HTML} and \appname{TeX4ht}, and how to create
PDF documents using \iPDFLaTeX\faq{What is PDFTeX?}{whatpdftex},
including how to create active links within your document using the
\isty{hyperref} package.

There are two new \LaTeX\ books that I haven't read but have
been recommended to me: \emph{\LaTeX\ Beginner's
Guide}~\cite{kottwitz} and \emph{\LaTeX\ and
Friends}~\cite{vandongen}.

Note that the \gls{uktug} has a 25\% book discount scheme
for members. See \url{http://uk.tug.org/membership} for more details
of that and other associated benefits. If you're not in the
UK, have a look at \url{http://www.tug.org/usergroups.html} to see
if there is a local user group in your area.

There is also a wealth of \LaTeX-related information on the
Internet\faq{How to get help}{gethelp}. \reportlinebreak\Gls{ctan} is a good place
to start. You can check the on-line catalogue~\cite{texcat} for
information about available software and, as mentioned earlier, there is also the list of
\htmladdnormallink{frequently asked
questions}{http://www.tex.ac.uk/faq} which I~recommend you try if
you have any queries. See also \appendixref{ch:help}

%%%%%%%%%%%%%%%%%%%%%%%%% DEFINITIONS %%%%%%%%%%%%%%%%%%%%%%%%%%%%%%

\setnode{definitions}
\printkeywords
\html{\documentclass[html]{novices}


\usepackage{fancybox}
\usepackage{tikz}
\usetikzlibrary{shapes}
\usetikzlibrary{decorations.pathmorphing}
\usepackage{caption,subcaption}
\usepackage{makeidx}
\usepackage{alltt}
\usepackage{amsmath}
\usepackage{amsfonts}
\usepackage{graphicx}
\usepackage{rotating}
\usepackage{booktabs}
\usepackage{pifont}

\usepackage{cmap}

\usepackage[utf8]{inputenc}
\usepackage[math]{anttor}
\usepackage{libris}
\renewcommand*{\ttdefault}{txtt}

\newcommand*{\docdate}{\formatdate{25}{9}{2012}}

\authordetails{1970}{Talbot}{Nicola L.~C.}
\title[LaTeX for Complete Novices]{\LaTeX\ for Complete Novices}
\version{1.4}
\edition{1}
\volume{1}
\series{The Dickimaw \LaTeX\ Series}
\seriesurl{http://www.dickimaw-books.com/latex/}
\affiliation
 {Dickimaw Books}
 {http://www.dickimaw-books.com/}
 {Saxlingham Nethergate}
\keywords{LaTeX,beginners,introduction,typesetting}
\subject{LaTeX}
\date{\docdate}

\html{
\newcommand*{\reportlinebreak}{}
\newcommand*{\reportpagebreak}{}
\newcommand*{\bookpagebreak}{}
\newcommand*{\booklinebreak}{}
\newcommand*{\screenpagebreak}{}
\newcommand*{\screenlinebreak}{}

\newcommand*{\dohtmlcolorbox}[2][0.8\linewidth]{%
 \colorbox{white}{%
   \begin{minipage}{#1}%
   \setlength{\parindent}{1.5em}%
    #2%
   \end{minipage}
 }%
}

\newcommand{\suppmaterial}[2][]{\htmlref{#2}{ch:supp}#1}
\newcommand{\suppref}[3][]{\htmlref{#2}{sec:#3}#1}

\providecommand{\dosidewaysfigure}[3]{}

\newcommand{\mathorit}[1]{\textit{#1}}

\newcommand{\objectref}[2]{\htmlref{#1~}{#2}\ref{#2}}
\newcommand{\figureref}[1]{\objectref{Figure}{#1}}
\newcommand{\xfigureref}[1]{\objectref{Figure}{#1}}
\newcommand{\vfigureref}[1]{\objectref{Figure}{#1}}
\newcommand{\xpageref}[2]{\htmlref{#1}{#2}}

\providecommand*{\vref}[1]{\ref{#1}}

\newcommand{\figuretop}[1]{}
\newcommand{\makeimg}[2]{$\htmlimage{ALT=#1}{}\mbox{#2}$}

\newcommand{\tableref}[1]{\objectref{Table}{#1}}

\newcommand{\exerciseref}[1]{\objectref{Exercise}{#1}}

\newcommand{\modification}[1]{\textbf{\texttt{#1}}}

\newcommand{\comptexttex}{\htmladdnormallink{comp.text.tex}{http://groups.google.com/group/comp.text.tex}}
\newcommand{\texhax}{\htmladdnormallink{texhax}{http://tug.org/pipermail/texhax/}}

\newcommand{\latexthesis}{%
   \htmladdnormallink{Using LaTeX to Write a PhD
Thesis}{\baseurl/latex/thesis/index.html}%
}

\newcommand{\latexpdfdoc}{\htmladdnormallink{Creating
a PDF Document using PDFLaTeX}{\baseurl/latex/pdfdoc/}}

\newenvironment{resultS}[1][]{\begin{htmlresult}{#1}\begin{makeimage}}{\end{makeimage}\end{htmlresult}}

\newenvironment{result}[1][]{\begin{htmlresult}{#1}\begin{makeimage}}{\end{makeimage}\end{htmlresult}}

\newcounter{object}

\newenvironment{inlinedef}{\begin{ttfamily}}{\end{ttfamily}}

\newcommand{\casemedia}[4]{#4}

\newcommand{\download}[1]{%
 \htmladdnormallink{download}{\exerciseurl/#1.tex}%
}

\newcommand{\downloadorview}[1]{%
 \htmladdnormallink{download}{\exerciseurl/#1.tex} or
 \htmladdnormallink{view}{\exerciseurl/#1.html}%
}

\newcommand{\doifbook}[1]{}
\newcommand{\doifnotbook}[1]{#1}
\newcommand{\doifreport}[1]{}

\newcommand{\doublebar}{\ensuremath{\|}}

\newcommand{\istystart}[1]{%
 \isty{#1}%
}

\newcommand{\istyend}[1]{%
 \isty{#1}%
}

\newcommand{\idxpagestyle}[1]{\index{page style!#1}}

\newcommand{\pagenumberingfmt}[1]{\texttt{#1}}

\newcommand{\indexpagenumbering}[1]{%
  \index{page numbering!#1}%
}

\newcommand{\ipagenumbering}[1]{%
  \pagenumberingfmt{#1}\indexpagenumbering{#1}%
}

\newcommand{\iunit}[2]{%
  \texttt{#1}\index{#1 (#2)@\texttt{#1} (#2)}%
  \index{units!\texttt{#1} (#2)}%
}

}

\makeindex
\makeglossaries

\newkeyword[\glstockey={Source Code}]{source}{source code}

\newkeyword[\glstockey={Output File}]{output}{output file}

\newkeyword
 [firstplural={\protect\emph{DVI (DeVice Independent) files}},
  \glstockey={DVI File}]
{dvi}{DVI file}

\newkeyword
 [\glsindexkey={auxiliary file (.aux)@auxiliary file (\texttt{.aux})},
  \glsextrakey={\extraindexaux},
  name={Auxiliary Files},\glstockey={Auxiliary Files}]
 {auxiliary}{auxiliary file}

\newkeyword
 [name={Terminal or Command Prompt},%
  \glsextrakey={\extraindexterminal},
  \glstockey={Terminal or Command Prompt}
 ]
{terminal}{terminal}

\newkeyword
 [name={Commands},\glstockey={Commands}]
 {command}{command}

\newkeyword
 [name={Grouping (or Scope)},%
  \glsextrakey={\extraindexgroup},
  \glstockey={Grouping (or Scope)}
 ]
 {group}{group}%

\newkeyword
 [name={Arguments (also called \dq{Parameters})},%
  \glsextrakey={\extraindexarg},
  \glstockey={Arguments}
 ]
{argument}{argument}%

\newkeyword
 [parent=argument,
  \glsindexkey={argument!mandatory},
  \glsextrakey={\extraindexmandatory},
  first={mandatory (or compulsory) argument},
  firstplural={mandatory (or compulsory) arguments},
  name={Mandatory Arguments},
  \glstockey={Mandatory Arguments}
 ]
{mandatory}{mandatory argument}

\newkeyword
 [parent=argument,
  \glsindexkey={argument!optional},
  \glsextrakey={\extraindexoptional},
  name={Optional Arguments},
  \glstockey={Optional Arguments}
]
{optional}{optional argument}


\newkeyword
 [\glsextrakey={\extraindexfragile},
  \glsindexkey={command!fragile},
  \glstockey={Moving Arguments and Fragile Commands},
   name={Moving Arguments and Fragile Commands}
 ]
{fragile}{fragile command}

\newkeyword
 [name={Robust Commands},
  \glstockey={Robust Commands},
  \glsindexkey={command!robust},
  \glsextrakey={\extraindexrobust}
 ]
{robust}{robust command}%

\newkeyword
 [name={Short and Long Commands},
  \glstockey={Short and Long Commands},
  \glsindexkey={command!short},
  \glsextrakey={\extraindexshortlong}]
{short}{short command}%


\newkeyword
[name={Declarations},\glstockey={Declarations}]
{declaration}{declaration}


\newkeyword
[name={Inter-Sentence Spacing},
 \glstockey={Inter-Sentence Spacing},
 \glsindexkey={spacing!inter-sentence},
 \glsextrakey={\extraindexsentencespacing}
]
{intersentencespacing}{inter-sentence spacing}%

\newkeyword
 [\glsingkey={hyphenating},
  \glsdkey={hyphenated},
  \glstockey={Hyphenation}]
 {hyphenation}{hyphenation}

\newkeyword
[name={Environments},
 \glstockey={Environments}]
{environment}{environment}%

\newkeyword
[name={The Preamble},
 \glstockey={The Preamble}]
{preamble}{preamble}%


\newkeyword
[name={Lengths},\glstockey={Lengths}]
{length}{length}

\newkeyword[\glstockey={Class File}]{cls}{class file}%

\newkeyword
[\glsindexkey={TeX@\protect\TeX},\glstockey={\texorpdfstring{\protect\TeX}{TeX}}]
{tex}{\TeX}

\newkeyword{perl}{Perl}

\glsaddall[types=keywords]
\glsunsetall[keywords]

\renewcommand{\summarypreamble}{%
Commands or environments defined in the \LaTeX\ kernel are always
available.%
}

\defgactivechar
 {backslashchar}
 {\backslashsym}
 {}
 {\LaTeX\ Kernel}
 {Escape character (indicates a command).}
 {}

\defgactivechar
 {dbbackslashchar}
 {\dbbackslashsym}
 {\oarg{\meta{height}}}
 {\LaTeX\ Kernel}
 {\nopostdesc}
 {%
   \BeginArgList
    \csentryargitem{height} Extra vertical space.
   \EndArgList
 }

\defgchildactivechar
 {newline.dbbackslashchar}
 {dbbackslashchar}
 {\dbbackslashsym}
 {Breaks a line without justification (\protect\htmlref{starred form}{itm:starredcommand}
   forbids a page break)}

\defgchildactivechar
 {tab.dbbackslashchar}
 {dbbackslashchar}
 {\dbbackslashsym}
 {Starts a new row in tabular-style environments}

\defgcs{tabularnewline}%
 {}%
 {\LaTeX\ Kernel}%
 {Behaves like \nxglsi{tab.dbbackslashchar} in a
   \nxglsni{env-tabular}-like environment but helps to disambiguate a
    \nxglslink{newline.dbbackslashchar}{line break} in a paragraph
    cell from a \nxglslink{tab.dbbackslashchar}{row separator}.}%
 {}

\defgidxactivechar
 {leftbracechar}
 {\leftbracesym}
 {}
 {\LaTeX\ Kernel}
 {Marks the beginning of a \nxglsi{group}.}
 {}

\defgidxactivechar
 {rightbracechar}
 {\rightbracesym}
 {}
 {\LaTeX\ Kernel}
 {Marks the end of a \nxglsi{group}.}
 {}

\defgidxactivecharcs
 {leftbrace}
 {\leftbracesym}
 {}
 {\LaTeX\ Kernel}
 {Left brace \{ character. In math mode may be used as a delimiter.}
 {}

\defgidxactivecharcs
 {rightbrace}
 {\rightbracesym}
 {}
 {\LaTeX\ Kernel}
 {Right brace \{ character. In math mode may be used as a delimiter.}
 {}

\defgidxactivechar
 {emdash}%
 {\emDashcs}%
 {}%
 {\LaTeX\ Kernel}%
 {Em-dash \textemdash\ symbol. (Normally used to
  indicate omissions or interruptions or to highlight a parenthetical element.)
  See also \nxglsni{textemdash}.}%
 {}

\defgidxactivechar
 {endash}%
 {\enDashcs}%
 {}%
 {\LaTeX\ Kernel}%
 {En-dash \textendash\ symbol.  (Normally used for number ranges.)
  See also \nxglsni{textendash}.}%
 {}

\defgidxactivechar
 {questiondown}%
 {\questiondowncs}%
 {}%
 {\LaTeX\ Kernel}%
 {Upside-down question mark \textquestiondown\ symbol. See also \nxglsni{textquestiondown}.}%
 {}

\defgidxactivechar
 [!Y]
 {exclamdown}%
 {\exclamdowncs}%
 {}%
 {\LaTeX\ Kernel}%
 {Upside-down exclamation mark \textexclamdown\ symbol. See also \nxglsni{textexclamdown}.}%
 {}

\defgchar
 {quoteleft}%
 {\quoteleftcs}%
 {}%
 {\LaTeX\ Kernel}%
 {Open quote \textquoteleft\ symbol. See also \nxglsni{textquoteleft}.}%
 {}

\defgchar
 {quotedblleft}%
 {\quotedblleftcs}%
 {}%
 {\LaTeX\ Kernel}%
 {Open double quote \textquotedblleft\ symbol. See also \nxglsni{textquotedblleft}.}%
 {}

\defgchar
 {quoteright}%
 {\quoterightcs}%
 {}%
 {\LaTeX\ Kernel}%
 {Closing quote or apostrophe \textquoteright\ symbol in text mode 
  or prime symbol \ensuremath{'} in math mode. See also \nxglsni{textquoteright}.}%
 {}

\defgchar
 {quotedblright}%
 {\quotedblrightcs}%
 {}%
 {\LaTeX\ Kernel}%
 {Closing double quote \textquotedblright\ symbol in text mode
  or double prime \ensuremath{''} in math mode. See also \nxglsni{textquotedblright}.}%
 {}

\defgchar
 {dash}%
 {\dashcs}%
 {}%
 {\LaTeX\ Kernel}%
 {Hyphen - in text mode or minus sign $-$ in math mode.}%
 {}

\defgidxactivecharcs
 {exclam}
 {\exclamsym}
 {}
 {\LaTeX\ Kernel (Math Mode)}
 {Negative thin space.}
 {}

\defgidxactivechar
 {exclamchar}
 {\exclamsym}
 {}
 {}
 {\nopostdesc}
 {}

\defgchildidxactivechar
 {sentence.exclamchar}
 {exclamchar}
 {\exclamsym}
 {Exclamation symbol (end of sentence marker)}
 

\defgchildidxactivechar
 {makeindex.exclamchar}
 {exclamchar}
 {\exclamsym}
 {\nxiappname{makeindex} sublevel special character}

\defgchildidxactivechar
 {resizebox.exclamchar}
 {exclamchar}
 {\exclamsym}
 {Used in \nxglsi{resizebox} to maintain aspect ratio}

\defgidxactivecharcs
 {csbar}
 {\vbarsym}
 {}
 {\LaTeX\ Kernel (Math Mode)}
 {Double vertical bar \doublebar{} delimiter}
 {}

\defgidxactivechar
 {barchar}
 {\vbarsym}
 {}
 {\LaTeX\ Kernel}
 {\nopostdesc}
 {}

\defgchildidxactivechar
 {delim.barchar}
 {barchar}
 {\vbarsym}
 {Delimiter. (Math mode only. Use \nxglsni{textbar} in text mode.)}

\defgchildidxactivechar
 {array.barchar}
 {barchar}
 {\vbarsym}
 {Vertical rule specifier (\nxglsi{env-tabular} or
   \nxglsi{env-array})}

\defgchar
 {slashchar}
 {\slashsym}
 {}
 {}
 {\nopostdesc}
 {}

\defgchildchar
 {text.slash}
 {slashchar}
 {\slashsym}
 {Forward slash symbol (see also \nxglsi{slash})}

\defgxchildchar
 {dir.slash}
 {slashchar}
 {\slashsym}
 {Directory divider}
 {directory divider}

\defgchildchar
 {delim.slash}
 {slashchar}
 {\slashsym}
 {Forward slash delimiter (math mode)}

\defgidxactivechar
 [!Z]
 {visiblespace}
 {\textvisiblespace}
 {}
 {}
 {A visual indication of a space in the code. When you type up 
   the code, replace all instances of this symbol with a space via the space bar on your 
   keyboard.}
 {}

\defgidxactivecharcs
 {at}
 {\atsym}
 {}
 {\LaTeX\ Kernel}
 {Used when a sentence ends with a capital letter.
   This command should be placed after the letter and before the
   punctuation mark.}
 {}

\defgidxactivechar
 {atchar}
 {\atsym}
 {\marg{\meta{text}}}
 {\LaTeX\ Kernel}
 {%
   Used in the argument of \nxglsi{env-tabular} or 
   \nxglsi{env-array} like environments to specify text to insert between columns.%
 }
 {%
     \BeginArgList
      \csentryargitem{text} Text to insert between columns
     \EndArgList
 }
  
\defgidxactivechar
 {questionchar}
 {\questionsym}
 {}
 {}
 {Question mark (end of sentence marker).}
 {}

\defgchar
 {ltchar}
 {\lesssym}
 {}
 {\LaTeX\ Kernel (Math Mode)}
 {Less than symbol. (Use \nxglsni{textless} in text mode.)}
 {}

\defgchar
 {gtchar}
 {\greatersym}
 {}
 {\LaTeX\ Kernel (Math Mode)}
 {Greater than symbol. (Use \nxglsni{textgreater} in text mode.)}
 {}

\defgchar
 {gtcol}
 {\greatersym}
 {\marg{\meta{decl}}}
 {\nxisty{array} package}
 {%
   Used in \nxglsi{env-tabular} or \nxglsi{env-array} column specifiers before
   \texttt{l}, \texttt{r}, \texttt{c}, \texttt{p}, \texttt{m} or
   \texttt{b} to insert \meta{decl} directly in front of the entry
   for that column.%
 }
 {%
     \BeginArgList
       \csentryargitem{decl} The code to insert at the start of the
         column.
     \EndArgList
 }

\defgchar
 {ltcol}
 {\lesssym}
 {\marg{\meta{decl}}}
 {\nxisty{array} package}
 {%
   Used in \nxglsi{env-tabular} or \nxglsi{env-array} column specifiers after
   \texttt{l}, \texttt{r}, \texttt{c}, \texttt{p}, \texttt{m} or
   \texttt{b} to insert \meta{decl} directly after the entry
   for that column.%
 }
 {%
     \BeginArgList
       \csentryargitem{decl} The code to insert at the end of the
         column.
     \EndArgList
 }

\defgcs{begin}%
 {\marg{\meta{env-name}}\oarg{\meta{env-option}}\marg{\meta{env-arg-1}}\ldots\marg{\meta{env-arg-n}}}%
 {\LaTeX\ Kernel}%
 {%
   Starts an environment. (Must have a matching \nxglsi{end}.)
 }%
 {%
   \BeginArgList
    \csentryargitem{env-name} The name of the environment. (\emph{No
      backslash.})
    \csentryargitem{env-option} An optional argument that may be 
      passed to the environment. Not all environments have optional
      arguments.
    \csentryargitem{env-arg-1}\ldots\meta{env-arg-n} Any mandatory
      arguments required by the environment. Not all environments
      require arguments.
   \EndArgList
 }

\defgcs{end}%
  {\marg{\meta{env-name}}}%
  {\LaTeX\ Kernel}%
  {Ends an environment. (Must have a matching \nxglsi{begin}.)}
  {%
    \BeginArgList
      \csentryargitem{env-name} The name of the environment.
    \EndArgList
  }

\defgcs{documentclass}%
 {\oarg{\meta{option-list}}\marg{\meta{class-name}}}%
 {\LaTeX\ Kernel}%
 {%
   Loads the document class file, which sets up the type of document
   you wish to write.%
 }%
 {%
   \BeginArgList
    \csentryargitem{option-list} A comma-separated list of options to
     pass to the class file or any packages that will later be
     loaded.
    \csentryargitem{class-name} The name of the document class. This
    corresponds to a file called \meta{class-name}\texttt{.cls},
    which must be installed. 
   \EndArgList
 }

\defgcs{usepackage}%
 {\oarg{\meta{option-list}}\marg{\meta{package-list}}}%
 {\LaTeX\ Kernel}%
 {%
   Loads the named packages.%
 }%
 {%
   \BeginArgList
    \csentryargitem{option-list} A comma-separated list of options to
     pass to the package.
    \csentryargitem{package-list} A comma-separated list of package
     names (without the \texttt{.sty} extension).
   \EndArgList
 }

\defgcs{footnote}%
 {\oarg{\meta{number}}\marg{\meta{text}}}%
 {\LaTeX\ Kernel}%
 {Inserts a footnote.}%
 {%
  \BeginArgList
    \csentryargitem{number} Overrides the default footnote number with
    the specified \meta{number}.
    \csentryargitem{text} The footnote text.
  \EndArgList
 }%

\defgcs{textbackslash}%
 {}%
 {\LaTeX\ Kernel (Text Mode)}%
 {Backlash \textbackslash\ symbol. (Use
  \nxglsi{backslash} for math mode.)}%
 {}

\defgcs{backslash}%
 {}%
 {\LaTeX\ Kernel (Math Mode)}%
 {Backslash \protect\ensuremath{\protect\backslash} symbol, which may be used
  as a delimiter. (Use
 \nxglsi{textbackslash} for text mode.)}%
 {}

\defgactivecharcs
 [\underscoresym]
 {underscore}
 {\textunderscore}
 {}
 {\LaTeX\ Kernel}
 {Underscore \_ symbol (see also \nxglsni{textunderscore}).}
 {}

\defgcs{textunderscore}
 {}
 {\LaTeX\ Kernel}
 {Underscore \textunderscore\ symbol (see also \nxglsni{textunderscore}).}
 {}

\defgactivechar
 [\underscoresym]
 {underscorechar}
 {\_}
 {\marg{\meta{maths}}}
 {\LaTeX\ Kernel (Math Mode)}
 {Displays its argument as a subscript.}
 {%
   \BeginArgList
     \csentryargitem{maths} The subscript.
   \EndArgList
 }

\defgcs{textgreater}%
 {}%
 {\LaTeX\ Kernel (Text Mode)}%
 {Greater than \textgreater\ symbol. (Just use \texttt{\textgreater} in math mode.)}%
 {}

\defgcs{textless}%
 {}%
 {\LaTeX\ Kernel (Text Mode)}%
 {Less than \textless\ symbol. (Just use \texttt{\textless} in math mode.)}%
 {}

\defgcs{textasciicircum}%
 {}%
 {\LaTeX\ Kernel}%
 {Circumflex \textasciicircum\ symbol.}%
 {}%

\defgcs[][cshyphenation]{hyphenation}
 {\marg{\meta{word}}}
 {\LaTeX\ Kernel}
 {Specifies hyphenation points.}
 {%
   \BeginArgList
    \csentryargitem{word} Word with hyphen points indicated with a
     dash (\texttt{-}).
   \EndArgList
 }

\defgsymcs[hyphen]{\dashcs}
 {}
 {\LaTeX\ Kernel}
 {Insert discretionary hyphen.}
 {}

\defgactivecharcs
 {dollar}
 {\dollarsym}
 {}
 {\LaTeX\ Kernel}
 {Dollar \$ symbol.}
 {}

\defgactivechar
 {dollarchar}
 {\dollarsym}
 {}
 {\LaTeX\ Kernel}
 {Switches in and out of in-line math mode.}
 {}

\defgcs{textasciitilde}%
 {}%
 {\LaTeX\ Kernel}%
 {Tilde \textasciitilde\ symbol. (If you are typing an
  URL, use the \nxisty{url} package, which provides \nxglsi{url}\marg{\meta{address}}
  that allows you to directly type \textasciitilde\ in the address.)}%
 {}%

\defgcs{url}%
 {\marg{\meta{address}}}%
 {\nxisty{url} package}%
 {Typesets an URL in a typewriter font and allows you to use
  characters such as \textasciitilde.}%
 {%
   \BeginArgList
    \csentryargitem{address} The web address.
   \EndArgList
 }%

\defgcs{dag}%
 {}%
 {\LaTeX\ Kernel}%
 {Dagger \dag\ symbol.}%
 {}

\defgcs{ddag}%
 {}%
 {\LaTeX\ Kernel}%
 {Double-dagger \ddag\ symbol.}%
 {}

\defgcs{textregistered}%
 {}%
 {\LaTeX\ Kernel}%
 {Registered \textregistered\ symbol.}%
 {}

\defgcs{texttrademark}%
 {}%
 {\LaTeX\ Kernel}%
 {Trademark \texttrademark\ symbol.}%
 {}

\defgcs{copyright}%
 {}%
 {\LaTeX\ Kernel}%
 {Copyright \copyright\ symbol.}%
 {}

\defgcs{textquoteright}%
 {}%
 {\LaTeX\ Kernel}%
 {Closing single quote (or apostrophe) \textquoteright\ symbol.}%
 {}

\defgcs{textquoteleft}%
 {}%
 {\LaTeX\ Kernel}%
 {Opening single quote \textquoteleft\ symbol.}%
 {}

\defgcs{textquotedblright}%
 {}%
 {\LaTeX\ Kernel}%
 {Closing double quote \textquotedblright\ symbol.}%
 {}

\defgcs{textquotedblleft}%
 {}%
 {\LaTeX\ Kernel}%
 {Opening double quote \textquotedblleft\ symbol.}%
 {}

\defgcs{textbullet}%
 {}%
 {\LaTeX\ Kernel (Text Mode)}%
 {Bullet \textbullet\ symbol.}%
 {}

\defgcs{textquestiondown}%
 {}%
 {\LaTeX\ Kernel}%
 {Upside-down question mark \textquestiondown\ symbol.}%
 {}

\defgcs{textexclamdown}%
 {}%
 {\LaTeX\ Kernel}%
 {Upside-down exclamation mark \textexclamdown\ symbol.}%
 {}

\defgcs{textendash}%
 {}%
 {\LaTeX\ Kernel}%
 {En-dash \textendash\ symbol. (Normally used for number ranges.)
  See also \nxglsni{endash}.}%
 {}

\defgcs{textemdash}%
 {}%
 {\LaTeX\ Kernel}%
 {Em-dash \textemdash\ symbol. (Normally used to
  indicate omissions or interruptions or to highlight a parenthetical element.)
  See also \nxglsni{emdash}.}%
 {}

\defgcs{textperiodcentered}%
 {}%
 {\LaTeX\ Kernel (Text Mode)}%
 {Centred period \textperiodcentered\ symbol.}%
 {}

\defgcs{i}%
 {}%
 {\LaTeX\ Kernel}%
 {Dotless i character: \i.}%
 {}

\defgcs{j}%
 {}%
 {\LaTeX\ Kernel}%
 {Dotless j character: \dotlessj.}%
 {}

\defgactivecharcs
 {hash}
 {\hashsym}
 {}
 {\LaTeX\ Kernel}
 {Hash \# symbol.}
 {}

\defgactivechar
 {hashchar}
 {\hashsym}
 {\meta{digit}}
 {\LaTeX\ Kernel}
 {Replacement text for argument \meta{digit}.}
 {}

\defgactivecharcs
 {percent}
 {\percentsym}
 {}
 {\LaTeX\ Kernel}
 {Percent \% symbol}
 {}

\defgactivechar
 {percentchar}
 {\percentsym}
 {}
 {\LaTeX\ Kernel}
 {Comment character used to ignore everything up to and including
  the newline character in the \nxgls{source}.}
 {}

\defgactivecharcs
 {amp}
 {\ampsym}
 {}
 {\LaTeX\ Kernel}
 {Ampersand \& symbol}
 {}

\defgactivechar
 {ampchar}
 {\ampsym}
 {}
 {\LaTeX\ Kernel}
 {Alignment tab.}
 {}

\defgcs{S}%
 {}%
 {\LaTeX\ Kernel}%
 {Sectional \S\ symbol.}%
 {}

\defgcs{P}%
 {}%
 {\LaTeX\ Kernel}%
 {Paragraph \P\ symbol.}%
 {}

\defgcs{slash}%
 {}%
 {\LaTeX\ Kernel}%
 {Forward slash \slash\ symbol.}%
 {}

\defgcs{AE}%
 {}%
 {\LaTeX\ Kernel}%
 {\protect\AE\ ligature.}%
 {}

\defgcs{ae}%
 {}%
 {\LaTeX\ Kernel}%
 {\protect\ae\ ligature.}%
 {}

\defgcs{OE}%
 {}%
 {\LaTeX\ Kernel}%
 {\protect\OE\ ligature.}%
 {}

\defgcs{oe}%
 {}%
 {\LaTeX\ Kernel}%
 {\protect\oe\ ligature.}%
 {}

\defgcs{AA}%
 {}%
 {\LaTeX\ Kernel}%
 {Upper case A-ring \AA\ character.}%
 {}

\defgcs{aa}%
 {}%
 {\LaTeX\ Kernel}%
 {Lower case a-ring \aa\ character.}%
 {}

\defgcs{L}%
 {}%
 {\LaTeX\ Kernel}%
 {Upper case L-bar \L\ character.}%
 {}

\defgcs{l}%
 {}%
 {\LaTeX\ Kernel}%
 {Lower case l-bar \l\ character.}%
 {}

\defgcs{O}%
 {}%
 {\LaTeX\ Kernel}%
 {Upper case slashed-O \O\ character.}%
 {}

\defgcs{o}%
 {}%
 {\LaTeX\ Kernel}%
 {Lower case slashed-o \o\ character.}%
 {}

\defgcs{ss}%
 {}%
 {\LaTeX\ Kernel}%
 {Eszett \ss\ character.}%
 {}

\defgcs{SS}%
 {}%
 {\LaTeX\ Kernel}%
 {SS (upper case \ss).}%
 {}

\defgcs{left}
 {\meta{delimiter}}
 {\LaTeX\ Kernel (Math Mode)}
 {Indicates a left stretchable delimiter. Must have a matching
  \nxglsi{right}.}
 {%
   \BeginArgList
    \csentryargitem{delimiter} A delimiter symbol, such as
    \nxglsni{openparen}, or a delimiter command, such as \nxglsni{langle}.
   \EndArgList
 }

\defgcs{right}
 {\meta{delimiter}}
 {\LaTeX\ Kernel (Math Mode)}
 {Indicates a right stretchable delimiter. Must have a matching
  \nxglsi{left}.}
 {%
   \BeginArgList
    \csentryargitem{delimiter} A delimiter symbol, such as
    \nxglsni{closeparen}, or a delimiter command, such as \nxglsni{rangle}.
   \EndArgList
 }

\defgcs{bigl}%
 {\meta{delimiter}}%
 {\nxisty{amsmath} package (Math Mode)}%
 {Left delimiter sizing.}%
 {%
   \BeginArgList
    \csentryargitem{delimiter} A delimiter symbol, such as
    \nxglsni{openparen}, or a delimiter command, such as \nxglsni{langle}.
   \EndArgList
 }

\defgcs{bigr}%
 {\meta{delimiter}}%
 {\nxisty{amsmath} package (Math Mode)}%
 {Right delimiter sizing.}%
 {%
   \BeginArgList
    \csentryargitem{delimiter} A delimiter symbol, such as
    \nxglsni{closeparen}, or a delimiter command, such as \nxglsni{rangle}.
   \EndArgList
 }

\defgcs{Bigl}%
 {\meta{delimiter}}%
 {\nxisty{amsmath} package (Math Mode)}%
 {Left delimiter sizing.}%
 {%
   \BeginArgList
    \csentryargitem{delimiter} A delimiter symbol, such as
    \nxglsni{openparen}, or a delimiter command, such as \nxglsni{langle}.
   \EndArgList
 }

\defgcs{Bigr}%
 {\meta{delimiter}}%
 {\nxisty{amsmath} package (Math Mode)}%
 {Right delimiter sizing.}%
 {%
   \BeginArgList
    \csentryargitem{delimiter} A delimiter symbol, such as
    \nxglsni{closeparen}, or a delimiter command, such as \nxglsni{rangle}.
   \EndArgList
 }

\defgcs{biggl}%
 {\meta{delimiter}}%
 {\nxisty{amsmath} package (Math Mode)}%
 {Left delimiter sizing.}%
 {%
   \BeginArgList
    \csentryargitem{delimiter} A delimiter symbol, such as
    \nxglsni{openparen}, or a delimiter command, such as \nxglsni{langle}.
   \EndArgList
 }

\defgcs{biggr}%
 {\meta{delimiter}}%
 {\nxisty{amsmath} package (Math Mode)}%
 {Right delimiter sizing.}%
 {%
   \BeginArgList
    \csentryargitem{delimiter} A delimiter symbol, such as
    \nxglsni{closeparen}, or a delimiter command, such as \nxglsni{rangle}.
   \EndArgList
 }

\defgcs{Biggl}%
 {\meta{delimiter}}%
 {\nxisty{amsmath} package (Math Mode)}%
 {Left delimiter sizing.}%
 {%
   \BeginArgList
    \csentryargitem{delimiter} A delimiter symbol, such as
    \nxglsni{openparen}, or a delimiter command, such as \nxglsni{langle}.
   \EndArgList
 }

\defgcs{Biggr}%
 {\meta{delimiter}}%
 {\nxisty{amsmath} package (Math Mode)}%
 {Right delimiter sizing.}%
 {%
   \BeginArgList
    \csentryargitem{delimiter} A delimiter symbol, such as
    \nxglsni{closeparen}, or a delimiter command, such as \nxglsni{rangle}.
   \EndArgList
 }

\defgcs{pounds}%
 {}%
 {\LaTeX\ Kernel}%
 {Pound \pounds\ symbol.}%
 {}%


\defgcs{v}%
 {\marg{\meta{c}}}%
 {\LaTeX\ Kernel}%
 {Caron diacritic over \meta{c}. Example: 
  \cmdname{v}\marg{o} produces \ocaron.}%
 {%
   \BeginArgList
    \csentryargitem{c} The character that requires the caron
      diacritic.
   \EndArgList
 }

\defgcs{u}%
 {\marg{\meta{c}}}%
 {\LaTeX\ Kernel}%
 {Breve diacritic over \meta{c}. Example: 
  \cmdname{u}\marg{o} produces \obreve.%
 }%
 {%
   \BeginArgList
    \csentryargitem{c} The character that requires the breve
      diacritic.
   \EndArgList
 }

\defgcs{H}%
 {\marg{\meta{c}}}%
 {\LaTeX\ Kernel}%
 {Double acute diacritic over \meta{c}. Example: 
  \cmdname{H}\marg{o} produces \odoubleacute.%
 }%
 {%
   \BeginArgList
    \csentryargitem{c} The character that requires the double
      acute diacritic.
   \EndArgList
 }

\defgcs{t}%
 {\marg{\meta{characters}}}%
 {\LaTeX\ Kernel}%
 {Tie over \meta{characters}. Example: 
  \cmdname{t}\marg{xy} produces \xytie.%
 }%
 {%
   \BeginArgList
    \csentryargitem{characters} The characters that require the tie.
   \EndArgList
 }

\defgcs{c}%
 {\marg{\meta{c}}}%
 {\LaTeX\ Kernel}%
 {Cedilla under \meta{c}. Example: 
  \cmdname{c}\marg{o} produces \ocedilla.%
 }%
 {%
   \BeginArgList
    \csentryargitem{c} The character that requires the cedilla.
   \EndArgList
 }

\defgcs{d}%
 {\marg{\meta{c}}}%
 {\LaTeX\ Kernel}%
 {Dot under \meta{c}. Example: 
  \cmdname{d}\marg{o} produces \odotunder.%
 }%
 {%
   \BeginArgList
    \csentryargitem{c} The character that requires the dot
     under it.
   \EndArgList
 }

\defgcs{b}%
 {\marg{\meta{c}}}%
 {\LaTeX\ Kernel}%
 {Bar under \meta{c}. Example: 
  \cmdname{b}\marg{r} produces \rbarunder.%
 }%
 {%
   \BeginArgList
    \csentryargitem{c} The character that requires the bar
     under it.
   \EndArgList
 }

\defgcs{r}%
 {\marg{\meta{c}}}%
 {\LaTeX\ Kernel}%
 {Ring over \meta{c}. Example: 
  \cmdname{r}\marg{u} produces \uring.%
 }%
 {%
   \BeginArgList
    \csentryargitem{c} The character that requires the ring
      over it.
   \EndArgList
 }

\defgsymcs[acute]{\quoterightcs}%
 {\marg{\meta{c}}}%
 {\LaTeX\ Kernel}%
 {Acute accent over \meta{c}. Example:
   \cmdname{'}\marg{o} produces \protect\'{o}.%
 }%
 {%
   \BeginArgList
    \csentryargitem{c} The character that requires an acute
     accent over it.
   \EndArgList
 }

\defgsymcs[grave]{\quoteleftcs}%
 {\marg{\meta{c}}}%
 {\LaTeX\ Kernel}%
 {Grave accent over \meta{c}. Example:
   \cmdname{`}\marg{o} produces \protect\`{o}.}%
 {%
   \BeginArgList
    \csentryargitem{c} The character that requires a grave
     accent over it.
   \EndArgList
 }

\defgsymcs[period]{\periodsym}%
 {\marg{\meta{c}}}%
 {\LaTeX\ Kernel}%
 {Dot over \meta{c}. Example:
   \cmdname{.}\marg{o} produces \odotover.%
 }%
 {%
   \BeginArgList
    \csentryargitem{c} The character that requires a dot
     over it.
   \EndArgList
 }

\defgchar
 {periodchar}
 {\periodsym}
 {}
 {\LaTeX\ Kernel}
 {\nopostdesc}
 {}

\defgchildchar
 {sentence.periodchar}
 {periodchar}
 {\periodsym}
 {period (full stop) or decimal point}

\defgchildchar
 {delimiter.periodchar}
 {periodchar}
 {\periodsym}
 {invisible delimiter}

\defgsymcs[macron]{\equalsym}%
 {\marg{\meta{c}}}%
 {\LaTeX\ Kernel}%
 {Macron accent over \meta{c}. Example:
   \cmdname{=}\marg{o} produces \omacron.%
 }%
 {%
   \BeginArgList
    \csentryargitem{c} The character that requires a macron
     accent over it.
   \EndArgList
 }

\defgactivecharcs
 [\circumsym\space]
 {circum}
 {\textasciicircum}
 {\marg{\meta{c}}}
 {\LaTeX\ Kernel}
 {%
   Circumflex accent over \meta{c}. Example:
   \cmdname{\textasciicircum}\marg{o} produces \ocircum.%
 }
 {%
   \BeginArgList
    \csentryargitem{c} The character that requires a
     circumflex accent over it.
   \EndArgList
 }

\defgactivechar
 [\circumsym]
 {circumchar}
 {\textasciicircum}
 {\marg{\meta{maths}}}
 {\LaTeX\ Kernel (Math Mode)}
 {Displays its argument as a superscript.}
 {%
   \BeginArgList
     \csentryargitem{maths} The superscript.
   \EndArgList
 }

\defgactivecharcs
 [\tildesym\space]
 {tilde}
 {\textasciitilde}
 {\marg{\meta{c}}}
 {\LaTeX\ Kernel}
 {%
   Tilde accent over \meta{c}. Example:
   \cmdname{\textasciitilde}\marg{o} produces \otilde.%
 }
 {%
   \BeginArgList
    \csentryargitem{c} The character that requires a
     tilde accent over it.
   \EndArgList
 }

\defgactivechar
 [\tildesym]
 {tildechar}
 {\textasciitilde}
 {}
 {\LaTeX\ Kernel}
 {Unbreakable space.}
 {}

\defgidxactivecharcs
 {doublequote}
 {\doublequotesym}
 {\marg{\meta{c}}}
 {\LaTeX\ Kernel}%
 {%
   Umlaut over \meta{c}. Example:
   \quotecs\marg{o} produces \oumlaut.%
 }%
 {%
   \BeginArgList
    \csentryargitem{c} The character that requires an umlaut
     over it.
   \EndArgList
 }

\defgcs{item}%
 {\oarg{\meta{marker}}}%
 {\LaTeX\ Kernel}%
 {Specifies the start of an item in a list. (Only allowed inside one of
  the list making environments.)}%
 {%
   \BeginArgList
    \csentryargitem{marker} If specified, the default item marker is
     replaced with \meta{marker}.
   \EndArgList
 }

\defgsymcs[itcorr]{\slashsym}%
 {}%
 {\LaTeX\ Kernel}%
 {Italic correction.}%
 {}

\defgcs{textrm}%
 {\marg{\meta{text}}}%
 {\LaTeX\ Kernel}%
 {Renders \meta{text} in the predefined serif font. (Defaults to
  Computer Modern Roman.)}%
 {%
   \BeginArgList
     \csentryargitem{text} The text on which to apply the font change.
   \EndArgList
 }

\defgcs{textsf}%
 {\marg{\meta{text}}}%
 {\LaTeX\ Kernel}%
 {Renders \meta{text} in the predefined sans-serif font. (Defaults to
  Computer Modern Sans.)}%
 {%
   \BeginArgList
     \csentryargitem{text} The text on which to apply the font change.
   \EndArgList
 }

\defgcs{texttt}%
 {\marg{\meta{text}}}%
 {\LaTeX\ Kernel}%
 {Renders \meta{text} in the predefined monospaced font. (Defaults to
  Computer Modern Typewriter.)}%
 {%
   \BeginArgList
     \csentryargitem{text} The text on which to apply the font change.
   \EndArgList
 }

\defgcs{textmd}%
 {\marg{\meta{text}}}%
 {\LaTeX\ Kernel}%
 {Renders \meta{text} with a medium weight in the current font
  family.}%
 {%
   \BeginArgList
     \csentryargitem{text} The text on which to apply the font change.
   \EndArgList
 }

\defgcs{textbf}%
 {\marg{\meta{text}}}%
 {\LaTeX\ Kernel}%
 {Renders \meta{text} with a bold weight in the current font family, if it exists.}%
 {%
   \BeginArgList
     \csentryargitem{text} The text on which to apply the font change.
   \EndArgList
 }

\defgcs{textup}%
 {\marg{\meta{text}}}%
 {\LaTeX\ Kernel}%
 {Renders \meta{text} with the upright form of the current font family.}%
 {%
   \BeginArgList
     \csentryargitem{text} The text on which to apply the font change.
   \EndArgList
 }

\defgcs{textit}%
 {\marg{\meta{text}}}%
 {\LaTeX\ Kernel}%
 {Renders \meta{text} with the italic form of the current font family, if it exists.}%
 {%
   \BeginArgList
     \csentryargitem{text} The text on which to apply the font change.
   \EndArgList
 }

\defgcs{textsl}%
 {\marg{\meta{text}}}%
 {\LaTeX\ Kernel}%
 {Renders \meta{text} with the slanted form of the current font family, if it exists.}%
 {%
   \BeginArgList
     \csentryargitem{text} The text on which to apply the font change.
   \EndArgList
 }

\defgcs{textsc}%
 {\marg{\meta{text}}}%
 {\LaTeX\ Kernel}%
 {Renders \meta{text} with the small-caps form of the current font family, if it exists.}%
 {%
   \BeginArgList
     \csentryargitem{text} The text on which to apply the font change.
   \EndArgList
 }

\defgcs{emph}%
 {\marg{\meta{text}}}%
 {\LaTeX\ Kernel}%
 {Toggles the upright and italic\slash slanted rendering of \meta{text}.}%
 {%
   \BeginArgList
     \csentryargitem{text} The text on which to apply the font change.
   \EndArgList
 }

\defgcs{textnormal}%
 {\marg{\meta{text}}}%
 {\LaTeX\ Kernel}%
 {Renders \meta{text} in the default font style.}%
 {%
   \BeginArgList
     \csentryargitem{text} The text on which to apply the font change.
   \EndArgList
 }

\defgcs{ttdefault}%
 {}%
 {\LaTeX\ Kernel}%
 {The name of the default typewriter family as used by \nxglsni{ttfamily}. 
  Defaults to \texttt{cmtt} (Computer Modern Typewriter).}%
 {}

\defgcs{rmdefault}%
 {}%
 {\LaTeX\ Kernel}%
 {The name of the default serif family as used by \nxglsni{rmfamily}. 
  Defaults to \texttt{cmr} (Computer Modern Roman).}%
 {}

\defgcs{sfdefault}%
 {}%
 {\LaTeX\ Kernel}%
 {The name of the default sans-serif family as used by \nxglsni{sffamily}. 
  Defaults to \texttt{cmss} (Computer Modern Sans-serif).}%
 {}

\defgcs{familydefault}%
 {}%
 {\LaTeX\ Kernel}%
 {Specifies the default font family. Defaults to
   \nxglsi{rmdefault} but may be redefined by certain classes.}%
 {}

\defgcs{rmfamily}%
 {}%
 {\LaTeX\ Kernel}%
 {Switches to the predefined serif font. (Defaults to
  Computer Modern Roman.)}%
 {}

\defgcs{sffamily}%
 {}%
 {\LaTeX\ Kernel}%
 {Switches to the predefined sans-serif font. (Defaults to
  Computer Modern Sans.)}%
 {}

\defgcs{ttfamily}%
 {}%
 {\LaTeX\ Kernel}%
 {Switches to the predefined monospaced font. (Defaults to
  Computer Modern Typewriter.)}%
 {}

\defgcs{mdseries}%
 {}%
 {\LaTeX\ Kernel}%
 {Switches to the medium weight in the current font
  family.}%
 {}

\defgcs{bfseries}%
 {}%
 {\LaTeX\ Kernel}%
 {Switches to the bold weight in the current font
  family.}%
 {}


\defgcs{upshape}%
 {}%
 {\LaTeX\ Kernel}%
 {Switches to the upright form of the current font family.}%
 {}

\defgcs{itshape}%
 {}%
 {\LaTeX\ Kernel}%
 {Switches to the italic form of the current font family, if it exists.}%
 {}

\defgcs{slshape}%
 {}%
 {\LaTeX\ Kernel}%
 {Switches to the slanted form of the current font family, if it exists.}%
 {}

\defgcs{scshape}%
 {}%
 {\LaTeX\ Kernel}%
 {Switches to the small-caps form of the current font family, if it exists.}%
 {}

\defgcs{em}%
 {}%
 {\LaTeX\ Kernel}%
 {Toggles the upright and italic\slash slanted form of the current font family.}%
 {}

\defgcs{normalfont}%
 {}%
 {\LaTeX\ Kernel}%
 {Switches to the default font style.}%
 {}

\defgcs{tiny}%
 {}%
 {Most document classes}%
 {Switches to tiny sized text.}%
 {}

\defgcs{scriptsize}%
 {}%
 {Most document classes}%
 {Switches to sub- or superscript sized text.}%
 {}

\defgcs{footnotesize}%
 {}%
 {Most document classes}%
 {Switches to footnote sized text.}%
 {}

\defgcs{small}%
 {}%
 {Most document classes}%
 {Switches to small sized text.}%
 {}

\defgcs{normalsize}%
 {}%
 {\LaTeX\ Kernel}%
 {Switches to normal sized text.}%
 {%
 }

\defgcs{large}%
 {}%
 {Most document classes}%
 {Switches to large sized text.}%
 {%
 }

\defgcs{Large}%
 {}%
 {Most document classes}%
 {Switches to extra-large sized text.}%
 {%
 }

\defgcs{LARGE}%
 {}%
 {Most document classes}%
 {Switches to extra-extra-large sized text.}%
 {%
 }

\defgcs{huge}%
 {}%
 {Most document classes}%
 {Switches to huge sized text.}%
 {%
 }

\defgcs{Huge}%
 {}%
 {Most document classes}%
 {Switches to extra-huge sized text.}%
 {%
 }

\defgcs{multicolumn}%
 {\marg{\meta{cols spanned}}\marg{\meta{col specifier}}\marg{\meta{text}}}%
 {\LaTeX\ Kernel}%
 {Spans multiple columns in a tabular-style environment.}%
 {%
   \BeginArgList
    \csentryargitem{cols spanned} The number of columns to span.
    \csentryargitem{col specifier} The alignment of this column-spanning entry.
    \csentryargitem{text} The contents of this column-spanning entry.
   \EndArgList
 }

\defgcs{toprule}%
 {\oarg{\meta{wd}}}%
 {\nxisty{booktabs} package}%
 {Horizontal rule for the top of a \nxglsi{env-tabular} environment.}%
 {%
   \BeginArgList
    \csentryargitem{wd} Thickness of the rule (dimension).
   \EndArgList
 }

\defgcs{bottomrule}%
 {\oarg{\meta{wd}}}%
 {\nxisty{booktabs} package}%
 {Horizontal rule for the bottom of a \nxglsi{env-tabular} environment.}%
 {%
   \BeginArgList
    \csentryargitem{wd} Thickness of the rule (dimension).
   \EndArgList
 }

\defgcs{midrule}%
 {\oarg{\meta{wd}}}%
 {\nxisty{booktabs} package}%
 {Horizontal rule to go below headings row of a \nxglsi{env-tabular} environment.}%
 {%
   \BeginArgList
    \csentryargitem{wd} Thickness of the rule (dimension).
   \EndArgList
 }

\defgcs{cmidrule}%
 {\oarg{\meta{wd}}\parg{\meta{trim}}\marg{\meta{n}-\meta{m}}}%
 {\nxisty{booktabs} package}%
 {Like \nxglsi{midrule} but only spans column \meta{n} to column
  \meta{m}.}%
 {%
   \BeginArgList
    \csentryargitem{wd} Thickness of the rule (dimension).
    \csentryargitem{trim} Trimming commands. Specifications:
     \texttt{r}, \texttt{r}\marg{\meta{wd}}, \texttt{l} and
      \texttt{l}\marg{\meta{wd}}
    \csentryargitem{n} Index of start column.
    \csentryargitem{m} Index of end column.
   \EndArgList
 }

\defgcs{heavyrulewidth}%
 {}%
 {\nxisty{booktabs} package}%
 {\nxGls{length} register specifying the thickness of \nxglsi{toprule} and
 \nxglsi{bottomrule}.}%
 {}

\defgcs{lightrulewidth}%
 {}%
 {\nxisty{booktabs} package}%
 {\nxGls{length} register specifying the thickness of \nxglsi{midrule}.}%
 {}

\defgcs{addlinespace}%
 {\oarg{\meta{wd}}}%
 {\nxisty{booktabs} package}%
 {Adds extra vertical space (\meta{wd} if specified, otherwise
  \nxglsi{defaultaddspace}) in a row of 
  a \nxglsi{env-tabular} environment. (Must be inserted after
  \nxglsi{dbbackslashchar} marker.)}%
 {%
   \BeginArgList
    \csentryargitem{wd} Height of space (dimension). Default is
     \nxglsi{defaultaddspace}.
   \EndArgList
 }

\defgcs{defaultaddspace}%
 {}%
 {\nxisty{booktabs} package}%
 {Default extra space before or after an adjacent rule.}%
 {}

\defgcs{parbox}%
 {\oarg{\meta{pos}}\oarg{\meta{height}}\marg{\meta{width}}\marg{\meta{text}}}%
 {\LaTeX\ Kernel}%
 {Makes a box with line-wrapped contents. (More restrictive than
  \nxglsi{env-minipage}.)}%
 {%
   \BeginArgList
    \csentryargitem{pos} The vertical alignment of the box relative to
     the surrounding text. (Centred if omitted.)
    \csentryargitem{height} The height of the box. (If omitted the
     height is the natural height of the contents of the box.)
    \csentryargitem{width} The width of the box.
    \csentryargitem{text} The contents of the box.
   \EndArgList
 }

\defgcs{author}%
 {\marg{\meta{name}}}%
 {Most classes that have the concept of a title page}%
 {Specifies the document author (or authors). This command doesn't
  display any text so may be used in the preamble, but if it's not
  in the preamble it must be placed before \nxglsi{maketitle}.}%
 {%
   \BeginArgList
    \csentryargitem{name} The name (or names) of the document author
     (or authors).
   \EndArgList
   Note that some classes, such as those supplied by journals or
   conference proceedings, may also define an optional argument
   that can be used to specify an abbreviated author list for the
   page headers.
 }

\defgcs{title}%
 {\marg{\meta{text}}}%
 {Most classes that have the concept of a title page}%
 {Specifies the document title. This command doesn't
  display any text so may be used in the preamble, but if it's not
  in the preamble it must be placed before \nxglsi{maketitle}.}%
 {%
   \BeginArgList
    \csentryargitem{text} The title of the document.
   \EndArgList
   Note that some classes, such as those supplied by journals or
   conference proceedings, may also define an optional argument
   that can be used to specify an abbreviated title for the page
   headers.
 }

\defgcs{date}%
 {\marg{\meta{text}}}%
 {Most classes that have the concept of a title page}%
 {Specifies the document date. This command doesn't
  display any text so may be used in the preamble, but if it's not
  in the preamble it must be placed before \nxglsi{maketitle}. If omitted, most
  classes assume the current date (as provided by \nxglsi{today}).}%
 {%
   \BeginArgList
    \csentryargitem{text} The document date.
   \EndArgList
 }

\defgcs{titlehead}%
 {\marg{\meta{text}}}%
 {\nxicls{scrartcl}, \nxicls{scrreprt}, \nxicls{scrbook} classes}%
 {Specifies the title header (typeset at the top of the title page).}%
 {%
   \BeginArgList
    \csentryargitem{text} The title header text.
   \EndArgList
 }

\defgcs{subtitle}%
 {\marg{\meta{text}}}%
 {\nxicls{scrartcl}, \nxicls{scrreprt}, \nxicls{scrbook} classes}%
 {Specifies the subtitle (typeset just below the title).}%
 {%
   \BeginArgList
    \csentryargitem{text} The subtitle text.
   \EndArgList
 }

\defgcs{subject}%
 {\marg{\meta{text}}}%
 {\nxicls{scrartcl}, \nxicls{scrreprt}, \nxicls{scrbook} classes}%
 {Specifies the subject (typeset just above the title).}%
 {%
   \BeginArgList
    \csentryargitem{text} The subject.
   \EndArgList
 }

\defgcs{publishers}%
 {\marg{\meta{text}}}%
 {\nxicls{scrartcl}, \nxicls{scrreprt}, \nxicls{scrbook} classes}%
 {Specifies the publisher (typeset after all the other titling
   information).}%
 {%
   \BeginArgList
    \csentryargitem{text} The publisher text.
   \EndArgList
 }

\defgcs{maketitle}%
 {}%
 {Most classes that have the concept of a title page}%
 {Generates the title page (or title block). This command is usually
  placed at the beginning of the document environment.}%
 {}

\defgcs{thanks}%
 {\marg{\meta{text}}}%
 {Most classes that have the concept of a title page}%
 {Inserts a special type of footnote in one of the titling fields,
  such as \nxglsi{author} or \nxglsni{title}. Usually used for some form of
  acknowledgement or affiliation.}%
 {}

\defgcs{today}%
 {}%
 {Most of the commonly-used classes}%
 {Inserts into the output file the date when the \LaTeX\ 
  application created it from the source code.}%
 {}

\defgcs{part}%
 {\oarg{\meta{short title}}\marg{\meta{title}}}%
 {Most classes that have the concept of document structure}%
 {Inserts a part sectional unit. \textbf{This command has a \htmlref{moving
  argument}{sec:fragile}.}}%
 {%
   \BeginArgList
     \csentryargitem{short title} An abbreviated form of the title to
       go in the table of contents or the page header.
     \csentryargitem{title} The title.
   \EndArgList
   The starred form of this command doesn't have an optional
   argument and doesn't increment or display the part counter.
 }

\defgcs{chapter}%
 {\oarg{\meta{short title}}\marg{\meta{title}}}%
 {Book-style classes (such as \nxicls{scrbook} or \nxicls{scrreprt}) that have 
  the concept of chapters}%
 {Inserts a chapter heading. \textbf{This command has a \htmlref{moving
  argument}{sec:fragile}.}}%
 {%
   \BeginArgList
     \csentryargitem{short title} An abbreviated form of the title to
       go in the table of contents or the page header.
     \csentryargitem{title} The title.
   \EndArgList
   The starred form of this command doesn't have an optional
   argument and doesn't increment or display the chapter counter.
 }

\defgcs{section}%
 {\oarg{\meta{short title}}\marg{\meta{title}}}%
 {Most classes that have the concept of document structure}%
 {Inserts a section header. \textbf{This command has a \htmlref{moving
  argument}{sec:fragile}.}}%
 {%
   \BeginArgList
     \csentryargitem{short title} An abbreviated form of the title to
       go in the table of contents or the page header.
     \csentryargitem{title} The title.
   \EndArgList
   The starred form of this command doesn't have an optional
   argument and doesn't increment or display the section counter.
 }

\defgcs{subsection}%
 {\oarg{\meta{short title}}\marg{\meta{title}}}%
 {Most classes that have the concept of document structure}%
 {Inserts a subsection header. \textbf{This command has a \htmlref{moving
  argument}{sec:fragile}.}}%
 {%
   \BeginArgList
     \csentryargitem{short title} An abbreviated form of the title to
       go in the table of contents or the page header.
     \csentryargitem{title} The title.
   \EndArgList
   The starred form of this command doesn't have an optional
   argument and doesn't increment or display the subsection counter.
 }

\defgcs{subsubsection}%
 {\oarg{\meta{short title}}\marg{\meta{title}}}%
 {Most classes that have the concept of document structure}%
 {Inserts a subsubsection header. \textbf{This command has a \htmlref{moving
  argument}{sec:fragile}.}}%
 {%
   \BeginArgList
     \csentryargitem{short title} An abbreviated form of the title to
       go in the table of contents or the page header.
     \csentryargitem{title} The title.
   \EndArgList
   The starred form of this command doesn't have an optional
   argument and doesn't increment or display the subsubsection counter.
 }

\defgcs{paragraph}%
 {\oarg{\meta{short title}}\marg{\meta{title}}}%
 {Most classes that have the concept of document structure}%
 {Inserts a subsubsubsection header. Most classes default to an 
  unnumbered running header for
  this sectional unit. \textbf{This command has a \htmlref{moving
  argument}{sec:fragile}.}}%
 {%
   \BeginArgList
     \csentryargitem{short title} An abbreviated form of the title to
       go in the table of contents or the page header.
     \csentryargitem{title} The title.
   \EndArgList
   The starred form of this command doesn't have an optional
   argument and doesn't increment or display the associated counter.
 }

\defgcs{subparagraph}%
 {\oarg{\meta{short title}}\marg{\meta{title}}}%
 {Most classes that have the concept of document structure}%
 {Inserts a subsubsubsubsection header. Most classes default 
  to an unnumbered running header for
  this sectional unit. \textbf{This command has a \htmlref{moving
  argument}{sec:fragile}.}}%
 {%
   \BeginArgList
     \csentryargitem{short title} An abbreviated form of the title to
       go in the table of contents or the page header.
     \csentryargitem{title} The title.
   \EndArgList
   The starred form of this command doesn't have an optional
   argument and doesn't increment or display the associated counter.
 }

\defgcs{minisec}%
 {\marg{\meta{heading}}}%
 {\nxicls{scrartcl}, \nxicls{scrreprt} and
  \nxicls{scrbook} classes}%
 {An unnumbered heading not associated with any structuring level.}%
 {%
   \BeginArgList
    \csentryargitem{heading} The heading text.
   \EndArgList
 }

\defgcs{addtokomafont}%
 {\marg{\meta{element name}}\marg{\meta{commands}}}%
 {\nxicls{scrartcl}, \nxicls{scrreprt} and
  \nxicls{scrbook} classes}%
 {Sets the font characteristics for the given \koma\ element.}%
 {%
   \BeginArgList
     \csentryargitem{element name} The element's name, for example
      \texttt{chapter}. See the \koma\ manual for a full list
      of defined elements.
     \csentryargitem{commands} The font changing commands to apply
      to that element.
   \EndArgList
 }

\defgcs{appendix}%
 {}%
 {Most classes that have the concept of document structure}%
 {Indicates (but doesn't print anything) that the document is
  switching to the appendices. If chapters exist, the chapter
  numbering is reset and switched to a different format
  (usually upper case letters) otherwise the section numbering
  is reset and switched to a different format.}%
 {}

\defgcs{tableofcontents}
 {}%
 {Most classes that have the concept of document structure}%
 {Inserts the table of contents. A second (possibly third) run
  is required to ensure the page numbering is correct.}%
 {}

\defgcs{addcontentsline}%
 {\marg{\meta{toc}}\marg{\meta{section unit}}\marg{\meta{text}}}%
 {\LaTeX\ Kernel}%
 {Adds a sectional unit header to the contents list.}%
 {%
   \BeginArgList
    \csentryargitem{toc} The extension of the required external file
    without the dot (for example, \texttt{toc} for the table of
     contents or \texttt{lof} for the list of figures).
    \csentryargitem{section unit} The name of the sectional unit (for
     example, \texttt{chapter}).
    \csentryargitem{text} The text to be added to the contents list.
   \EndArgList
 }

\defgcs{label}%
 {\marg{\meta{string}}}%
 {\LaTeX\ Kernel}%
 {Assigns a unique textual label linked to the most recently
  incremented cross-referencing counter in the current \nxglslink{group}{scope} (see
  also \nxglsi{ref}).}%
 {%
   \BeginArgList
    \csentryargitem{string} A unique label that can be referenced
     elsewhere in the document with \nxglsi{ref}. (It's best to 
     just use alphanumerics and non-active punctuation characters in
     the label.)
   \EndArgList
 }

\defgcs{ref}%
 {\marg{\meta{string}}}%
 {\LaTeX\ Kernel}%
 {References the value of the counter linked to the given label.
  A second (possibly third) run of \LaTeX\ is required to ensure the cross-references
  are up-to-date.
 }%
 {%
   \BeginArgList
    \csentryargitem{string} The required label that was used in the
     corresponding \nxglsi{label}.
   \EndArgList
 }

\defgcs{pageref}%
 {\marg{\meta{string}}}%
 {\LaTeX\ Kernel}%
 {Similar to \nxglsi{ref} but inserts the page number where the given
  label was defined.
  A second (possibly third) run of \LaTeX\ is required to ensure the cross-references
  are up-to-date.
 }%
 {%
   \BeginArgList
    \csentryargitem{string} The required label that was used in the
     corresponding \nxglsi{label}.
   \EndArgList
 }

\defgcs{vref}%
 {\marg{\meta{string}}}%
 {\nxisty{varioref} package}%
 {Like \nxglsi{ref} but also adds information about the location, such
  as \dq{on page~\meta{n}} or \dq{on the following page}.
 }%
 {%
   \BeginArgList
    \csentryargitem{string} The required label that was used in the
     corresponding \nxglsi{label}.
   \EndArgList
 }

\defgcs{bibitem}%
 {\oarg{\meta{tag}}\marg{\meta{key}}}%
 {\LaTeX\ Kernel}%
 {Indicates the start of a new reference in the bibliography. May
  only be used inside the contents of \nxglsi{env-thebibliography} environment}%
 {%
   \BeginArgList
    \csentryargitem{tag} If present, overrides the marker at the
    start of the reference.
    \csentryargitem{key} A unique key that identifies this reference so
    it can be cited elsewhere in the document using \nxglsi{cite}.
   \EndArgList
 }

\defgcs{cite}%
 {\oarg{\meta{text}}\marg{\meta{key list}}}%
 {\LaTeX\ Kernel}%
 {Inserts the citation markers of each reference identified in the
  key list. A second run is required to ensure the reference is
  correct.}%
 {%
    \BeginArgList
      \csentryargitem{text} Additional text to insert into the citation
       (such as the chapter number, or a particular page or
       page range within the citation).
      \csentryargitem{key list} A comma-separated list of keys used in
      the corresponding \nxglsi{cite}.
    \EndArgList
 }

\defgcs{pagenumbering}%
 {\marg{\meta{style}}}%
 {\LaTeX\ Kernel}%
 {Sets the style of the page numbers.}%
 {%
   \BeginArgList
     \csentryargitem{style} The page numbering style (e.g.\
      \texttt{roman} for lower case Roman numerals).
   \EndArgList
 }

\defgcs{pagestyle}%
 {\marg{\meta{style}}}%
 {\LaTeX\ Kernel}%
 {Sets the style of the headers and footers.}%
 {%
   \BeginArgList
     \csentryargitem{style} The name of the page style. The \LaTeX\
     kernel defines only two styles: \nxipagestyle{empty} and
     \nxipagestyle{plain}. Most of the standard classes also provide
     the \nxipagestyle{headings} style.
   \EndArgList
 }

\defgcs{thispagestyle}%
 {\marg{\meta{style}}}%
 {\LaTeX\ Kernel}%
 {Like \nxglsi{pagestyle} but only affects the current page.}%
 {%
   \BeginArgList
     \csentryargitem{style} The name of the page style.
   \EndArgList
 }

\defgcs{markboth}%
 {\marg{\meta{left head}}\marg{\meta{right head}}}%
 {\LaTeX\ Kernel}%
 {Specifies information for the left and right page headers. Not all
   page styles use this information, in which case the arguments are
   ignored.}%
 {%
   \BeginArgList
     \csentryargitem{left head} Text for the left (even) page header,
     if the document has been identified as a double-sided document
     (usually with the \nxclsopt{twoside} class option).

     \csentryargitem{right head} Text for the right page header,
     if the \nxclsopt{twoside} option has been used. Otherwise used
     for both odd and even page headers.
   \EndArgList
 }

\defgcs{markright}%
 {\marg{\meta{right head}}}%
 {\LaTeX\ Kernel}%
 {Specifies information for the right (odd) page header. Not all
  page styles use this information, in which case the argument is
  ignored.}%
 {%
   \BeginArgList
     \csentryargitem{right head} Text for the right page header,
     if the \nxclsopt{twoside} option has been used. Otherwise used
     for both odd and even page headers.
   \EndArgList
 }

\defgcs{includegraphics}%
 {\oarg{\meta{key vals}}\marg{\meta{filename}}}%
 {\nxisty{graphicx} package}%
 {Inserts a graphics file into the document.}%
 {%
   \BeginArgList
     \csentryargitem{key vals} Comma-separated list of options that can
      be used to change the way the image is displayed.
   \EndArgList
 }

\defgcs{DeclareGraphicsExtensions}%
 {\marg{\meta{ext-list}}}%
 {\nxisty{graphicx} package}%
 {Specify the file extensions to look for if no extension is used in
 \nxglsi{includegraphics}}%
 {%
   \BeginArgList
     \csentryargitem{ext-list} A comma-separated list of extensions.
   \EndArgList
 }

\defgcs{selectlanguage}
 {\marg{\meta{language name}}}%
 {\nxisty{babel} package}%
 {Switches to the named language. Predefined textual elements, such
  as the date given by \nxglsi{today} or prefixes like \dq{Chapter},
  are redefined to those supplied by the given language.}%
 {%
   \BeginArgList
    \csentryargitem{language} The language identifier.
   \EndArgList
 }

\defgcs{foreignlanguage}%
 {\marg{\meta{language name}}\marg{\meta{text}}}%
 {\nxisty{babel} package}%
 {Typesets the given text using any predefined names or date formats supplied by the
  given language.
 }%
 {%
   \BeginArgList
    \csentryargitem{language} The language identifier.
    \csentryargitem{text} The foreign text.
   \EndArgList
 }

\defgcs{iflanguage}%
 {\marg{\meta{language name}}\marg{\meta{true text}}\marg{\meta{false
text}}}%
 {\nxisty{babel} package}%
 {Tests the current language.}%
 {%
   \BeginArgList
    \csentryargitem{language} The language identifier.
    \csentryargitem{true text} The text to typeset if true.
    \csentryargitem{false text} The text to typeset if not true.
   \EndArgList
 }

\defgcs{caption}%
 {\oarg{\meta{short caption}}\marg{\meta{caption text}}}%
 {\LaTeX\ Kernel}%
 {Inserts the caption for a float such as a figure or table. 
  \textbf{This command has a \htmlref{moving argument}{sec:fragile}.}}%
 {%
   \BeginArgList
    \csentryargitem{short caption} If provided, an abbreviated caption
      to go in the list of figures\slash tables etc.
    \csentryargitem{caption text} The caption contents.
   \EndArgList
 }

\defgcs{listoffigures}%
 {}%
 {Most classes that have the concept of document structure}%
 {Inserts the list of figures. A second (possibly third) run
  is required to ensure the page numbering is correct.}%
 {}

\defgcs{listoftables}%
 {}%
 {Most classes that have the concept of document structure}%
 {Inserts the list of tables. A second (possibly third) run
  is required to ensure the page numbering is correct.}%
 {}

\defgcs{hspace}%
 {\marg{\meta{length}}}%
 {\LaTeX\ Kernel}%
 {Inserts a horizontal gap of the given width.}%
 {%
   \BeginArgList
    \csentryargitem{length} The width of the horizontal gap.
   \EndArgList
 }

\defgcs{vspace}%
 {\marg{\meta{length}}}%
 {\LaTeX\ Kernel}%
 {Inserts a vertical gap of the given height.}%
 {%
   \BeginArgList
    \csentryargitem{length} The \protect\htmlref{height}{sec:length} of the vertical gap.
   \EndArgList
 }

\defgcs{baselineskip}%
 {}%
 {\LaTeX\ Kernel}%
 {A \nxglsi{length} register that stores the current interline spacing.
  This is recalculated whenever the font changes.}%
 {}

\defgcs{parskip}%
 {}%
 {\LaTeX\ Kernel}%
 {A \nxglsi{length} register that stores the spacing between paragraphs.
  (If you're using one of the \koma\ classes, use the \nxscrclsopt{parskip} 
   option to set it to full or half line height.)}%
 {}

\defgcs{twocolumn}%
 {\oarg{\meta{header text}}}%
 {\LaTeX\ Kernel}%
 {Starts a new page and switches to two column mode.}%
 {%
   \BeginArgList
    \csentryargitem{header text} If present, this header text will be
    placed at the top of the page spanning both columns. (For
    example, as used by \nxglsi{maketitle} in two column articles.)
   \EndArgList
 }

\defgcs{onecolumn}%
 {}%
 {\LaTeX\ Kernel}%
 {Starts a new page and switches to one column mode.}%
 {}

\defgcs{TeX}%
 {}%
 {\LaTeX\ Kernel}%
 {Typesets the \TeX\ logo.}%
 {}

\defgcs{LaTeX}%
 {}%
 {\LaTeX\ Kernel}%
 {Typesets the \makeimg{LaTeX}{\LaTeX}\ logo.}%
 {}

\defgcs{LaTeXe}%
 {}%
 {\LaTeX\ Kernel}%
 {Typesets the \makeimg{LaTeX2e}{\LaTeXe}\ logo.}%
 {}

\defgcs{framebox}%
 {\oarg{\meta{width}}\oarg{\meta{align}}\marg{\meta{text}}}%
 {\LaTeX\ Kernel}%
 {Puts a frame around its contents, prohibiting a line break in the
  contents.}%
 {%
   \BeginArgList
     \csentryargitem{width} The width of the box. (If omitted, this value is
      the natural width of the box contents.)
     \csentryargitem{align} Horizontal alignment of box contents
      (\texttt{c} if omitted).
     \csentryargitem{text} The contents of the box.
   \EndArgList
 }

\defgcs{fbox}%
 {\marg{\meta{text}}}%
 {\LaTeX\ Kernel}%
 {Puts a frame around its contents, prohibiting a line break in the
  contents.}%
 {%
   \BeginArgList
     \csentryargitem{text} The contents of the box.
   \EndArgList
 }

\defgcs{shadowbox}%
 {\marg{\meta{text}}}%
 {\nxisty{fancybox} package}%
 {Puts a shadow frame around its contents, prohibiting a line break in the
  contents.}%
 {%
   \BeginArgList
     \csentryargitem{text} The contents of the box.
   \EndArgList
 }

\defgcs{doublebox}%
 {\marg{\meta{text}}}%
 {\nxisty{fancybox} package}%
 {Puts a double-lined frame around its contents, prohibiting a line break in the
  contents.}%
 {%
   \BeginArgList
     \csentryargitem{text} The contents of the box.
   \EndArgList
 }

\defgcs{ovalbox}%
 {\marg{\meta{text}}}%
 {\nxisty{fancybox} package}%
 {Puts a thin-lined oval frame around its contents, prohibiting a line break in the
  contents.}%
 {%
   \BeginArgList
     \csentryargitem{text} The contents of the box.
   \EndArgList
 }

\defgcs{Ovalbox}%
 {\marg{\meta{text}}}%
 {\nxisty{fancybox} package}%
 {Puts a thick-lined oval frame around its contents, prohibiting a line break in the
  contents.}%
 {%
   \BeginArgList
     \csentryargitem{text} The contents of the box.
   \EndArgList
 }

\defgcs{protect}%
 {\meta{command}}%
 {\LaTeX\ Kernel}%
 {Used in a moving argument to prevent a fragile command from expanding.}%
 {%
   \BeginArgList
    \csentryargitem{command} The fragile command that needs protecting.
   \EndArgList
 }

\defgcs{par}%
 {}%
 {\LaTeX\ Kernel}%
 {Insert a paragraph break.}%
 {}

\defgcs{linewidth}%
 {}%
 {\LaTeX\ Kernel}%
 {A length containing the desired current line width. This is
  usually the width of the typeblock, but inside a \nxglsi{env-minipage}
  or \nxglsi{parbox} it will be the width the box. Note that
  the actual contents of the line may fall short of the line width
  (underfull hbox) or extend beyond it (overfull hbox).}%
 {}

\defgcs{textwidth}%
 {}%
 {\LaTeX\ Kernel}%
 {A length containing the width of the typeblock. Note that
  the actual contents of the line may fall short of the line width
  (underfull hbox) or extend beyond it (overfull hbox). This width
  does not include the area for marginal notes.}%
 {}

\defgcs{textheight}%
 {}%
 {\LaTeX\ Kernel}%
 {A length containing the height of the typeblock. Note that
  the actual contents of the page may fall short of the text height
  (underfull vbox) or extend beyond it (overfull vbox). This
  measurement does not include the header and footer areas.}%
 {}

\defgcs{and}%
 {}%
 {\LaTeX\ Kernel}%
 {Used to separate authors in \nxglsi{author}}%
 {}

\defgcs{clearpage}%
 {}%
 {\LaTeX\ Kernel}%
 {Inserts a page break and processes any unprocessed floats}%
 {}

\defgcs{rotatebox}%
 {\oarg{\meta{option list}}\marg{\meta{angle}}\marg{\meta{text}}}%
 {\nxisty{graphicx} package}%
 {Rotates the given contents by the given angle.}%
 {%
   \BeginArgList
     \csentryargitem{option list} A comma-separated list of options:
      \protect\begin{itemize}
       \protect\item \nxikeyvalopt{rotatebox}{origin}=\meta{label}
       \protect\item \nxikeyvalopt{rotatebox}{x}=\meta{dimen}
       \protect\item \nxikeyvalopt{rotatebox}{y}=\meta{dimen}
       \protect\item \nxikeyvalopt{rotatebox}{units}=\meta{number}
      \protect\end{itemize}

     \csentryargitem{angle} The angle of rotation.

     \csentryargitem{text} The text to be rotated.
   \EndArgList
 }

\defgcs{scalebox}%
 {\marg{\meta{h scale}}\oarg{\meta{v scale}}\marg{\meta{text}}}%
 {\nxisty{graphicx} package}%
 {Scales the specified contents.}%
 {%
    \BeginArgList
      \csentryargitem{h scale} The horizontal scale factor.
      \csentryargitem{v scale} The vertical scale factor. If omitted,
       the same as \meta{h scale}.
      \csentryargitem{text} The text to be scaled.
    \EndArgList
 }

\defgcs{reflectbox}%
 {\marg{\meta{text}}}%
 {\nxisty{graphicx} package}%
 {Reflects the specified contents in the \makeimg{y}{\ensuremath{y}}-axis.)}%
 {%
    \BeginArgList
      \csentryargitem{text} The text to be reflected.
    \EndArgList
 }

\defgcs{resizebox}%
 {\marg{\meta{h length}}\marg{\meta{v length}}\marg{\meta{text}}}%
 {\nxisty{graphicx} package}%
 {Scales the specified contents to the given dimensions.}%
 {%
   \BeginArgList
      \csentryargitem{h length} The required width or \texttt{!} to keep
      the aspect ratio.
      \csentryargitem{v length} The required height or \texttt{!} to keep
      the aspect ratio.
      \csentryargitem{text} The text to be scaled.
   \EndArgList
 }

\defgcs{centering}%
 {}%
 {\LaTeX\ Kernel}%
 {Switches the paragraph alignment to centred.}%
 {}

\defgcs{newcommand}%
 {\marg{\meta{cmd}}\oarg{\meta{n-args}}\oarg{\meta{default}}\marg{\meta{text}}}%
 {\LaTeX\ Kernel}%
 {Defines a new command.}%
 {%
   \BeginArgList
    \csentryargitem{cmd} The new command name (including initial backslash).
    \csentryargitem{n-args} The number of arguments this new command
     should have.
    \csentryargitem{default} If the first argument should be optional,
     the default value if omitted.
    \csentryargitem{text} What actions the command should perform.
   \EndArgList
 }

\defgcs{renewcommand}%
 {\marg{\meta{cmd}}\oarg{\meta{n-args}}\oarg{\meta{default}}\marg{\meta{text}}}%
 {\LaTeX\ Kernel}%
 {Redefines an existing command.}%
 {%
   \BeginArgList
    \csentryargitem{cmd} The command name (including initial backslash).
    \csentryargitem{n-args} The number of arguments this command
     should have.
    \csentryargitem{default} If the first argument should be optional,
     the default value if omitted.
    \csentryargitem{text} What actions the command should perform.
   \EndArgList
 }

\defgcs{newenvironment}%
 {\marg{\meta{env-name}}\oarg{\meta{n-args}}\oarg{\meta{default}}\marg{\meta{begin-code}}\marg{\meta{end-code}}}%
 {\LaTeX\ Kernel}%
 {Defines a new environment.}%
 {%
   \BeginArgList
    \csentryargitem{env-name} The new environment name (\emph{no} backslash).
    \csentryargitem{n-args} The number of arguments this new
     environment should have.
    \csentryargitem{default} If the first argument should be optional,
     the default value if omitted.
    \csentryargitem{begin-code} Actions to perform at the start of the
      environment.
    \csentryargitem{end-code} Actions to perform at the end of the
      environment.
   \EndArgList
 }

\defgcs{renewenvironment}%
 {\marg{\meta{env-name}}\oarg{\meta{n-args}}\oarg{\meta{default}}\marg{\meta{begin-code}}\marg{end-code}}%
 {\LaTeX\ Kernel}%
 {Redefines an existing environment.}%
 {%
   \BeginArgList
    \csentryargitem{env-name} The environment name (\emph{no} backslash).
    \csentryargitem{n-args} The number of arguments this new
     environment should have.
    \csentryargitem{default} If the first argument should be optional,
     the default value if omitted.
    \csentryargitem{begin-code} Actions to perform at the start of the
      environment.
    \csentryargitem{end-code} Actions to perform at the end of the
      environment.
   \EndArgList
 }

\defgcs{labelitemi}%
 {}%
 {Classes that define the \nxglsni{env-itemize} environment}%
 {The default label for the first level \nxglsni{env-itemize}.}%
 {}

\defgcs{labelitemii}%
 {}%
 {Classes that define the \nxglsni{env-itemize} environment}%
 {The default label for the second level \nxglsni{env-itemize}.}%
 {}

\defgcs{labelitemiii}%
 {}%
 {Classes that define the \nxglsni{env-itemize} environment}%
 {The default label for the third level \nxglsni{env-itemize}.}%
 {}

\defgcs{labelitemiv}%
 {}%
 {Classes that define the \nxglsni{env-itemize} environment}%
 {The default label for the fourth level \nxglsni{env-itemize}.}%
 {}

\defgcs{mbox}%
 {\marg{\meta{text}}}%
 {\LaTeX\ Kernel}%
 {Ensures that the given text doesn't contain a line break.}%
 {}

\defgcs{text}%
 {\marg{\meta{text}}}%
 {\nxisty{amsmath} package (Math Mode)}%
 {Displays its argument in the normal text font (as opposed to the
  current maths font).}%
 {%
   \BeginArgList
     \csentryargitem{text} The text to be displayed in the normal text font.
   \EndArgList
 }

\defgcs{intertext}%
 {\marg{\meta{text}}}%
 {\nxisty{amsmath} package (Math Mode)}%
 {Used for a short interjection in the middle of a multi-line
  displayed maths, such as in an \nxglsi{env-align} environment.
  May only appear right after \nxglsi{dbbackslashchar}.}%
 {%
   \BeginArgList
     \csentryargitem{text} The interjection text.
   \EndArgList
 }

\defgcs{overleftarrow}%
 {\marg{\meta{maths}}}%
 {\LaTeX\ Kernel (Math Mode)}%
 {Puts an extendible left arrow over \meta{maths}}%
 {%
   \BeginArgList
     \csentryargitem{maths} The maths over which to put the arrow.
   \EndArgList
 }

\defgcs{overrightarrow}%
 {\marg{\meta{maths}}}%
 {\LaTeX\ Kernel (Math Mode)}%
 {Puts an extendible right arrow over \meta{maths}}%
 {%
   \BeginArgList
     \csentryargitem{maths} The maths over which to put the arrow.
   \EndArgList
 }

\defgcs{overleftrightarrow}%
 {\marg{\meta{maths}}}%
 {\nxisty{amsmath} package (Math Mode)}%
 {Puts an extendible left-right arrow over \meta{maths}}%
 {%
   \BeginArgList
     \csentryargitem{maths} The maths over which to put the arrow.
   \EndArgList
 }

\defgcs{underleftarrow}%
 {\marg{\meta{maths}}}%
 {\nxisty{amsmath} package (Math Mode)}%
 {Puts an extendible left arrow under \meta{maths}}%
 {%
   \BeginArgList
     \csentryargitem{maths} The maths under which to put the arrow.
   \EndArgList
 }

\defgcs{underrightarrow}%
 {\marg{\meta{maths}}}%
 {\nxisty{amsmath} package (Math Mode)}%
 {Puts an extendible right arrow under \meta{maths}}%
 {%
   \BeginArgList
     \csentryargitem{maths} The maths under which to put the arrow.
   \EndArgList
 }

\defgcs{underleftrightarrow}%
 {\marg{\meta{maths}}}%
 {\nxisty{amsmath} package (Math Mode)}%
 {Puts an extendible left-right arrow under \meta{maths}}%
 {%
   \BeginArgList
     \csentryargitem{maths} The maths under which to put the arrow.
   \EndArgList
 }

\defgcs{xrightarrow}%
 {\oarg{\meta{subscript}}\marg{\meta{superscript}}}%
 {\nxisty{amsmath} package (Math Mode)}%
 {An extendible right arrow with a superscript and optionally a
  subscript.}%
 {%
   \BeginArgList
     \csentryargitem{subscript} The subscript to go under the arrow.
     \csentryargitem{superscript} The superscript to go over the arrow.
   \EndArgList
 }

\defgcs{xleftarrow}%
 {\oarg{\meta{subscript}}\marg{\meta{superscript}}}%
 {\nxisty{amsmath} package (Math Mode)}%
 {An extendible left arrow with a superscript and optionally a
  subscript.}%
 {%
   \BeginArgList
     \csentryargitem{subscript} The subscript to go under the arrow.
     \csentryargitem{superscript} The superscript to go over the arrow.
   \EndArgList
 }

\defgcs{rightarrow}%
 {}%
 {\LaTeX\ Kernel (Math Mode)}%
 {Right arrow \ensuremath{\rightarrow}.}%
 {}

\defgcs{to}%
 {}%
 {\LaTeX\ Kernel (Math Mode)}%
 {Right arrow \ensuremath{\to}.}%
 {}

\defgcs{Rightarrow}%
 {}%
 {\LaTeX\ Kernel (Math Mode)}%
 {Double-lined right arrow \ensuremath{\Rightarrow}.}%
 {}

\defgcs{leftarrow}%
 {}%
 {\LaTeX\ Kernel (Math Mode)}%
 {Left arrow \ensuremath{\leftarrow}.}%
 {}

\defgcs{gets}%
 {}%
 {\LaTeX\ Kernel (Math Mode)}%
 {Left arrow \ensuremath{\gets}.}%
 {}

\defgcs{Leftarrow}%
 {}%
 {\LaTeX\ Kernel (Math Mode)}%
 {Double-lined left arrow \ensuremath{\Leftarrow}.}%
 {}

\defgcs{downarrow}%
 {}%
 {\LaTeX\ Kernel (Math Mode)}%
 {Down arrow \ensuremath{\downarrow}. (May be used as a delimiter.)}%
 {}

\defgcs{Downarrow}%
 {}%
 {\LaTeX\ Kernel (Math Mode)}%
 {Double-lined down arrow \ensuremath{\Downarrow}. (May be used as a delimiter.)}%
 {}

\defgcs{uparrow}%
 {}%
 {\LaTeX\ Kernel (Math Mode)}%
 {Up arrow \ensuremath{\uparrow}. (May be used as a delimiter.)}%
 {}

\defgcs{Uparrow}%
 {}%
 {\LaTeX\ Kernel (Math Mode)}%
 {Double-lined up arrow \ensuremath{\Uparrow}. (May be used as a delimiter.)}%
 {}

\defgcs{updownarrow}%
 {}%
 {\LaTeX\ Kernel (Math Mode)}%
 {Double-ended vertical arrow \ensuremath{\updownarrow}. (May be used as a delimiter.)}%
 {}

\defgcs{Updownarrow}%
 {}%
 {\LaTeX\ Kernel (Math Mode)}%
 {Double-ended double-lined vertical arrow \ensuremath{\Updownarrow}. (May be used as a delimiter.)}%
 {}

\defgcs{hookrightarrow}%
 {}%
 {\LaTeX\ Kernel (Math Mode)}%
 {Hooked right arrow \ensuremath{\hookrightarrow}.}%
 {}

\defgcs{hookleftarrow}%
 {}%
 {\LaTeX\ Kernel (Math Mode)}%
 {Hooked left arrow \ensuremath{\hookleftarrow}.}%
 {}

\defgcs{leftharpoondown}%
 {}%
 {\LaTeX\ Kernel (Math Mode)}%
 {Left down harpoon \ensuremath{\leftharpoondown}.}%
 {}

\defgcs{leftharpoonup}%
 {}%
 {\LaTeX\ Kernel (Math Mode)}%
 {Left up harpoon \ensuremath{\leftharpoonup}.}%
 {}

\defgcs{rightharpoondown}%
 {}%
 {\LaTeX\ Kernel (Math Mode)}%
 {Right down harpoon \ensuremath{\rightharpoondown}.}%
 {}

\defgcs{rightharpoonup}%
 {}%
 {\LaTeX\ Kernel (Math Mode)}%
 {Right up harpoon \ensuremath{\rightharpoonup}.}%
 {}

\defgcs{rightleftharpoons}%
 {}%
 {\LaTeX\ Kernel (Math Mode)}%
 {Right-left harpoons \ensuremath{\rightleftharpoons}.}%
 {}

\defgcs{leftrightarrow}%
 {}%
 {\LaTeX\ Kernel (Math Mode)}%
 {Double-ended horizontal arrow \ensuremath{\leftrightarrow}.}%
 {}

\defgcs{Leftrightarrow}%
 {}%
 {\LaTeX\ Kernel (Math Mode)}%
 {Double-ended double-lined horizontal arrow \ensuremath{\Leftrightarrow}.}%
 {}

\defgcs{longrightarrow}%
 {}%
 {\LaTeX\ Kernel (Math Mode)}%
 {Long right arrow \ensuremath{\longrightarrow}.}%
 {}

\defgcs{longleftarrow}%
 {}%
 {\LaTeX\ Kernel (Math Mode)}%
 {Long left arrow \ensuremath{\longleftarrow}.}%
 {}

\defgcs{longleftrightarrow}%
 {}%
 {\LaTeX\ Kernel (Math Mode)}%
 {Long double-ended horizontal arrow \ensuremath{\longleftrightarrow}.}%
 {}

\defgcs{Longrightarrow}%
 {}%
 {\LaTeX\ Kernel (Math Mode)}%
 {Long double-lined right arrow \ensuremath{\Longrightarrow}.}%
 {}

\defgcs{Longleftarrow}%
 {}%
 {\LaTeX\ Kernel (Math Mode)}%
 {Long double-lined left arrow \ensuremath{\Longleftarrow}.}%
 {}

\defgcs{Longleftrightarrow}%
 {}%
 {\LaTeX\ Kernel (Math Mode)}%
 {Long double-lined double-ended horizontal arrow \ensuremath{\Longleftrightarrow}.}%
 {}

\defgcs{longmapsto}%
 {}%
 {\LaTeX\ Kernel (Math Mode)}%
 {Long mapping arrow \ensuremath{\longmapsto}.}%
 {}

\defgcs{mapsto}%
 {}%
 {\LaTeX\ Kernel (Math Mode)}%
 {Mapping arrow \ensuremath{\mapsto}.}%
 {}

\defgcs{nearrow}%
 {}%
 {\LaTeX\ Kernel (Math Mode)}%
 {North-East arrow \ensuremath{\nearrow}.}%
 {}

\defgcs{nwarrow}%
 {}%
 {\LaTeX\ Kernel (Math Mode)}%
 {North-West arrow \ensuremath{\nwarrow}.}%
 {}

\defgcs{searrow}%
 {}%
 {\LaTeX\ Kernel (Math Mode)}%
 {South-East arrow \ensuremath{\searrow}.}%
 {}

\defgcs{swarrow}%
 {}%
 {\LaTeX\ Kernel (Math Mode)}%
 {South-West arrow \ensuremath{\swarrow}.}%
 {}

\defgcs{makeindex}%
 {}%
 {\LaTeX\ Kernel (Preamble Only)}%
 {Enables \nxglsi{index}.}%
 {}

\defgcs{index}%
 {\marg{\meta{text}}}%
 {\LaTeX\ Kernel}%
 {Adds indexing information to an external index file. The command
   \nxglsi{makeindex} must be used in the preamble to enable
   this command. The external index file must be post-processed with
   an indexing application, such as \nxiappname{makeindex}.}%
 {%
   \BeginArgList
     \csentryargitem{text} The text to go in the index.
   \EndArgList
 }

\defgcs{printindex}%
 {}%
 {\nxisty{makeidx} package}%
 {Prints the index. Must be used with \nxglsi{makeindex} and \nxglsi{index}.
 (The external index file must first be processed by an indexing application.)}%
 {}

\defgcs{alpha}%
 {}%
 {\LaTeX\ Kernel (Math Mode)}%
 {Greek lower case alpha \ensuremath{\alpha}.}%
 {}

\defgcs{beta}%
 {}%
 {\LaTeX\ Kernel (Math Mode)}%
 {Greek lower case beta \ensuremath{\beta}.}%
 {}

\defgcs{gamma}%
 {}%
 {\LaTeX\ Kernel (Math Mode)}%
 {Greek lower case gamma \ensuremath{\gamma}.}%
 {}

\defgcs{delta}%
 {}%
 {\LaTeX\ Kernel (Math Mode)}%
 {Greek lower case delta \ensuremath{\delta}.}%
 {}

\defgcs{epsilon}%
 {}%
 {\LaTeX\ Kernel (Math Mode)}%
 {Greek lower case epsilon \ensuremath{\epsilon}.}%
 {}

\defgcs{varepsilon}%
 {}%
 {\LaTeX\ Kernel (Math Mode)}%
 {Variant Greek lower case alpha \ensuremath{\varepsilon}.}%
 {}

\defgcs{zeta}%
 {}%
 {\LaTeX\ Kernel (Math Mode)}%
 {Greek lower case zeta \ensuremath{\zeta}.}%
 {}

\defgcs{eta}%
 {}%
 {\LaTeX\ Kernel (Math Mode)}%
 {Greek lower case eta \ensuremath{\eta}.}%
 {}

\defgcs{theta}%
 {}%
 {\LaTeX\ Kernel (Math Mode)}%
 {Greek lower case theta \ensuremath{\theta}.}%
 {}

\defgcs{vartheta}%
 {}%
 {\LaTeX\ Kernel (Math Mode)}%
 {A variant Greek lower case theta \ensuremath{\vartheta}.}%
 {}

\defgcs{iota}%
 {}%
 {\LaTeX\ Kernel (Math Mode)}%
 {Greek lower case iota \ensuremath{\iota}.}%
 {}

\defgcs{kappa}%
 {}%
 {\LaTeX\ Kernel (Math Mode)}%
 {Greek lower case kappa \ensuremath{\kappa}.}%
 {}

\defgcs{lambda}%
 {}%
 {\LaTeX\ Kernel (Math Mode)}%
 {Greek lower case lambda \ensuremath{\lambda}.}%
 {}

\defgcs{mu}%
 {}%
 {\LaTeX\ Kernel (Math Mode)}%
 {Greek lower case mu \ensuremath{\mu}.}%
 {}

\defgcs{nu}%
 {}%
 {\LaTeX\ Kernel (Math Mode)}%
 {Greek lower case nu \ensuremath{\nu}.}%
 {}

\defgcs{xi}%
 {}%
 {\LaTeX\ Kernel (Math Mode)}%
 {Greek lower case xi \ensuremath{\xi}.}%
 {}

\defgcs{pi}%
 {}%
 {\LaTeX\ Kernel (Math Mode)}%
 {Greek lower case pi \ensuremath{\pi}.}%
 {}

\defgcs{varpi}%
 {}%
 {\LaTeX\ Kernel (Math Mode)}%
 {Variant Greek lower case pi \ensuremath{\varpi}.}%
 {}

\defgcs{rho}%
 {}%
 {\LaTeX\ Kernel (Math Mode)}%
 {Greek lower case rho \ensuremath{\rho}.}%
 {}

\defgcs{varrho}%
 {}%
 {\LaTeX\ Kernel (Math Mode)}%
 {Variant Greek lower case rho \ensuremath{\varrho}.}%
 {}

\defgcs{sigma}%
 {}%
 {\LaTeX\ Kernel (Math Mode)}%
 {Greek lower case sigma \ensuremath{\sigma}.}%
 {}

\defgcs{varsigma}%
 {}%
 {\LaTeX\ Kernel (Math Mode)}%
 {Variant Greek lower case sigma \ensuremath{\varsigma}.}%
 {}

\defgcs{tau}%
 {}%
 {\LaTeX\ Kernel (Math Mode)}%
 {Greek lower case tau \ensuremath{\tau}.}%
 {}

\defgcs{upsilon}%
 {}%
 {\LaTeX\ Kernel (Math Mode)}%
 {Greek lower case upsilon \ensuremath{\upsilon}.}%
 {}

\defgcs{phi}%
 {}%
 {\LaTeX\ Kernel (Math Mode)}%
 {Greek lower case phi \ensuremath{\phi}.}%
 {}

\defgcs{varphi}%
 {}%
 {\LaTeX\ Kernel (Math Mode)}%
 {Variant Greek lower case phi \ensuremath{\varphi}.}%
 {}

\defgcs{chi}%
 {}%
 {\LaTeX\ Kernel (Math Mode)}%
 {Greek lower case chi \ensuremath{\chi}.}%
 {}

\defgcs{psi}%
 {}%
 {\LaTeX\ Kernel (Math Mode)}%
 {Greek lower case psi \ensuremath{\psi}.}%
 {}

\defgcs{omega}%
 {}%
 {\LaTeX\ Kernel (Math Mode)}%
 {Greek lower case omega \ensuremath{\omega}.}%
 {}

\defgcs{Gamma}%
 {}%
 {\LaTeX\ Kernel (Math Mode)}%
 {Greek upper case gamma \ensuremath{\Gamma}.}%
 {}

\defgcs{Delta}%
 {}%
 {\LaTeX\ Kernel (Math Mode)}%
 {Greek upper case delta \ensuremath{\Delta}.}%
 {}

\defgcs{Theta}%
 {}%
 {\LaTeX\ Kernel (Math Mode)}%
 {Greek upper case theta \ensuremath{\Theta}.}%
 {}

\defgcs{Lambda}%
 {}%
 {\LaTeX\ Kernel (Math Mode)}%
 {Greek upper case lambda \ensuremath{\Lambda}.}%
 {}

\defgcs{Xi}%
 {}%
 {\LaTeX\ Kernel (Math Mode)}%
 {Greek upper case xi \ensuremath{\Xi}.}%
 {}

\defgcs{Pi}%
 {}%
 {\LaTeX\ Kernel (Math Mode)}%
 {Greek upper case pi \ensuremath{\Pi}.}%
 {}

\defgcs{Sigma}%
 {}%
 {\LaTeX\ Kernel (Math Mode)}%
 {Greek upper case sigma \ensuremath{\Sigma}.}%
 {}

\defgcs{Upsilon}%
 {}%
 {\LaTeX\ Kernel (Math Mode)}%
 {Greek upper case upsilon \ensuremath{\Upsilon}.}%
 {}

\defgcs{Phi}%
 {}%
 {\LaTeX\ Kernel (Math Mode)}%
 {Greek upper case phi \ensuremath{\Phi}.}%
 {}

\defgcs{Psi}%
 {}%
 {\LaTeX\ Kernel (Math Mode)}%
 {Greek upper case psi \ensuremath{\Psi}.}%
 {}

\defgcs{Omega}%
 {}%
 {\LaTeX\ Kernel (Math Mode)}%
 {Greek upper case omega \ensuremath{\Omega}.}%
 {}

\defgcs{sb}%
 {\marg{\meta{maths}}}%
 {\LaTeX\ Kernel (Math Mode)}%
 {Displays its argument as a subscript.}%
 {%
   \BeginArgList
     \csentryargitem{maths} The subscript.
   \EndArgList
 }

\defgcs{sp}%
 {\marg{\meta{maths}}}%
 {\LaTeX\ Kernel (Math Mode)}%
 {Displays its argument as a superscript.}%
 {%
   \BeginArgList
     \csentryargitem{maths} The superscript.
   \EndArgList
 }

\defgcs{pmod}%
 {\marg{\meta{maths}}}%
 {\LaTeX\ Kernel (Math Mode)}%
 {Modulo operator with parentheses.}%
 {%
   \BeginArgList
     \csentryargitem{maths} The modulo (placed in parentheses).
   \EndArgList
 }

\defgcs{mod}%
 {\marg{\meta{maths}}}%
 {\nxisty{amsmath} (Math Mode)}%
 {Modulo operator without parentheses.}%
 {%
   \BeginArgList
     \csentryargitem{maths} The modulo (placed in parentheses).
   \EndArgList
 }

\defgcs{pod}%
 {\marg{\meta{maths}}}%
 {\nxisty{amsmath} (Math Mode)}%
 {Modulo operator with parentheses but no \dq{mod}.}%
 {%
   \BeginArgList
     \csentryargitem{maths} The modulo (placed in parentheses).
   \EndArgList
 }

\defgcs{bmod}%
 {}%
 {\LaTeX\ Kernel (Math Mode)}%
 {Modulo operator.}%
 {}

\defgcs{arccos}%
 {}%
 {\LaTeX\ Kernel (Math Mode)}%
 {Typesets \ensuremath{\arccos} function name.}%
 {}

\defgcs{arcsin}%
 {}%
 {\LaTeX\ Kernel (Math Mode)}%
 {Typesets \ensuremath{\arcsin} function name.}%
 {}

\defgcs{arctan}%
 {}%
 {\LaTeX\ Kernel (Math Mode)}%
 {Typesets \ensuremath{\arctan} function name.}%
 {}

\defgcs{arg}%
 {}%
 {\LaTeX\ Kernel (Math Mode)}%
 {Typesets \ensuremath{\arg} function name.}%
 {}

\defgcs{cos}%
 {}%
 {\LaTeX\ Kernel (Math Mode)}%
 {Typesets \ensuremath{\cos} function name.}%
 {}

\defgcs{cosh}%
 {}%
 {\LaTeX\ Kernel (Math Mode)}%
 {Typesets \ensuremath{\cosh} function name.}%
 {}

\defgcs{cot}%
 {}%
 {\LaTeX\ Kernel (Math Mode)}%
 {Typesets \ensuremath{\cot} function name.}%
 {}

\defgcs{coth}%
 {}%
 {\LaTeX\ Kernel (Math Mode)}%
 {Typesets \ensuremath{\coth} function name.}%
 {}

\defgcs{csc}%
 {}%
 {\LaTeX\ Kernel (Math Mode)}%
 {Typesets \ensuremath{\csc} function name.}%
 {}

\defgcs{deg}%
 {}%
 {\LaTeX\ Kernel (Math Mode)}%
 {Typesets \ensuremath{\deg} function name.}%
 {}

\defgcs{det}%
 {}%
 {\LaTeX\ Kernel (Math Mode)}%
 {Typesets \ensuremath{\det} function name (may have limits via \nxglsi{underscorechar} or \nxglsi{circumchar}).}%
 {}

\defgcs{dim}%
 {}%
 {\LaTeX\ Kernel (Math Mode)}%
 {Typesets \ensuremath{\dim} function name.}%
 {}

\defgcs{exp}%
 {}%
 {\LaTeX\ Kernel (Math Mode)}%
 {Typesets \ensuremath{\exp} function name.}%
 {}

\defgcs{gcd}%
 {}%
 {\LaTeX\ Kernel (Math Mode)}%
 {Typesets \ensuremath{\gcd} function name (may have limits via \nxglsi{underscorechar} or \nxglsi{circumchar}).}%
 {}

\defgcs{hom}%
 {}%
 {\LaTeX\ Kernel (Math Mode)}%
 {Typesets \ensuremath{\hom} function name.}%
 {}

\defgcs{inf}%
 {}%
 {\LaTeX\ Kernel (Math Mode)}%
 {Typesets \ensuremath{\inf} function name (may have limits via \nxglsi{underscorechar} or \nxglsi{circumchar}).}%
 {}

\defgcs{injlim}%
 {}%
 {\nxisty{amsmath} (Math Mode)}%
 {Typesets \ensuremath{\injlim} function name (may have limits via \nxglsi{underscorechar} or \nxglsi{circumchar}).}%
 {}

\defgcs{projlim}%
 {}%
 {\nxisty{amsmath} (Math Mode)}%
 {Typesets \ensuremath{\projlim} function name (may have limits via \nxglsi{underscorechar} or \nxglsi{circumchar}).}%
 {}

\defgcs{varlimsup}%
 {}%
 {\nxisty{amsmath} (Math Mode)}%
 {Typesets \ensuremath{\varlimsup} function name (may have limits via \nxglsi{underscorechar} or \nxglsi{circumchar}).}%
 {}

\defgcs{varliminf}%
 {}%
 {\nxisty{amsmath} (Math Mode)}%
 {Typesets \ensuremath{\varliminf} function name (may have limits via \nxglsi{underscorechar} or \nxglsi{circumchar}).}%
 {}

\defgcs{varinjlim}%
 {}%
 {\nxisty{amsmath} (Math Mode)}%
 {Typesets \ensuremath{\varinjlim} function name (may have limits via \nxglsi{underscorechar} or \nxglsi{circumchar}).}%
 {}

\defgcs{varprojlim}%
 {}%
 {\nxisty{amsmath} (Math Mode)}%
 {Typesets \ensuremath{\varprojlim} function name (may have limits via \nxglsi{underscorechar} or \nxglsi{circumchar}).}%
 {}

\defgcs{ker}%
 {}%
 {\LaTeX\ Kernel (Math Mode)}%
 {Typesets \ensuremath{\ker} function name.}%
 {}

\defgcs{lg}%
 {}%
 {\LaTeX\ Kernel (Math Mode)}%
 {Typesets \ensuremath{\lg} function name.}%
 {}

\defgcs{lim}%
 {}%
 {\LaTeX\ Kernel (Math Mode)}%
 {Typesets \ensuremath{\lim} function name (may have limits via \nxglsi{underscorechar} or \nxglsi{circumchar}).}%
 {}

\defgcs{liminf}%
 {}%
 {\LaTeX\ Kernel (Math Mode)}%
 {Typesets \ensuremath{\liminf} function name (may have limits via \nxglsi{underscorechar} or \nxglsi{circumchar}).}%
 {}

\defgcs{limsup}%
 {}%
 {\LaTeX\ Kernel (Math Mode)}%
 {Typesets \ensuremath{\limsup} function name (may have limits via \nxglsi{underscorechar} or \nxglsi{circumchar}).}%
 {}

\defgcs{ln}%
 {}%
 {\LaTeX\ Kernel (Math Mode)}%
 {Typesets \ensuremath{\ln} function name.}%
 {}

\defgcs{log}%
 {}%
 {\LaTeX\ Kernel (Math Mode)}%
 {Typesets \ensuremath{\log} function name.}%
 {}

\defgcs{max}%
 {}%
 {\LaTeX\ Kernel (Math Mode)}%
 {Typesets \ensuremath{\max} function name (may have limits via \nxglsi{underscorechar} or \nxglsi{circumchar}).}%
 {}

\defgcs{min}%
 {}%
 {\LaTeX\ Kernel (Math Mode)}%
 {Typesets \ensuremath{\min} function name (may have limits via \nxglsi{underscorechar} or \nxglsi{circumchar}).}%
 {}

\defgcs{Pr}%
 {}%
 {\LaTeX\ Kernel (Math Mode)}%
 {Typesets \ensuremath{\Pr} function name (may have limits via \nxglsi{underscorechar} or \nxglsi{circumchar}).}%
 {}

\defgcs{sec}%
 {}%
 {\LaTeX\ Kernel (Math Mode)}%
 {Typesets \ensuremath{\sec} function name.}%
 {}

\defgcs{sin}%
 {}%
 {\LaTeX\ Kernel (Math Mode)}%
 {Typesets \ensuremath{\sin} function name.}%
 {}

\defgcs{sinh}%
 {}%
 {\LaTeX\ Kernel (Math Mode)}%
 {Typesets \ensuremath{\sinh} function name.}%
 {}

\defgcs{sup}%
 {}%
 {\LaTeX\ Kernel (Math Mode)}%
 {Typesets \ensuremath{\sup} function name (may have limits via \nxglsi{underscorechar} or \nxglsi{circumchar}).}%
 {}

\defgcs{tan}%
 {}%
 {\LaTeX\ Kernel (Math Mode)}%
 {Typesets \ensuremath{\tan} function name.}%
 {}

\defgcs{tanh}%
 {}%
 {\LaTeX\ Kernel (Math Mode)}%
 {Typesets \ensuremath{\tanh} function name.}%
 {}

\defgcs{infty}%
 {}%
 {\LaTeX\ Kernel (Math Mode)}%
 {Infinity \ensuremath{\infty} symbol.}%
 {}

\defgcs{partial}%
 {}%
 {\LaTeX\ Kernel (Math Mode)}%
 {Partial \ensuremath{\partial} symbol.}%
 {}

\defgcs{subset}%
 {}%
 {\LaTeX\ Kernel (Math Mode)}%
 {Subset \ensuremath{\subset} symbol.}%
 {}

\defgcs{approx}%
 {}%
 {\LaTeX\ Kernel (Math Mode)}%
 {Relational \ensuremath{\approx} symbol.}%
 {}

\defgcs{asymp}%
 {}%
 {\LaTeX\ Kernel (Math Mode)}%
 {Relational \ensuremath{\asymp} symbol.}%
 {}

\defgcs{bowtie}%
 {}%
 {\LaTeX\ Kernel (Math Mode)}%
 {Relational \ensuremath{\bowtie} symbol.}%
 {}

\defgcs{cong}%
 {}%
 {\LaTeX\ Kernel (Math Mode)}%
 {Relational \ensuremath{\cong} symbol.}%
 {}

\defgcs{dashv}%
 {}%
 {\LaTeX\ Kernel (Math Mode)}%
 {Relational \ensuremath{\dashv} symbol.}%
 {}

\defgcs{doteq}%
 {}%
 {\LaTeX\ Kernel (Math Mode)}%
 {Relational \makeimg{equals symbol with dot on top}{\ensuremath{\doteq}} symbol.}%
 {}

\defgcs{equiv}%
 {}%
 {\LaTeX\ Kernel (Math Mode)}%
 {Relational \makeimg{equivalent}{\ensuremath{\equiv}} symbol.}%
 {}

\defgcs{frown}%
 {}%
 {\LaTeX\ Kernel (Math Mode)}%
 {Relational \ensuremath{\frown} symbol.}%
 {}

\defgcs{ge}%
 {}%
 {\LaTeX\ Kernel (Math Mode)}%
 {Relational \makeimg{greater than or equal to}{\ensuremath{\ge}} symbol.}%
 {}

\defgcs{geq}%
 {}%
 {\LaTeX\ Kernel (Math Mode)}%
 {Relational \makeimg{greater than or equal to}{\ensuremath{\geq}} symbol.}%
 {}

\defgcs{gg}%
 {}%
 {\LaTeX\ Kernel (Math Mode)}%
 {Relational \makeimg{much greater than}{\ensuremath{\gg}} symbol.}%
 {}

\defgcs{in}%
 {}%
 {\LaTeX\ Kernel (Math Mode)}%
 {Relational \ensuremath{\in} symbol.}%
 {}

\defgcs{le}%
 {}%
 {\LaTeX\ Kernel (Math Mode)}%
 {Relational \makeimg{less than or equal to}{\ensuremath{\le}} symbol.}%
 {}

\defgcs{leq}%
 {}%
 {\LaTeX\ Kernel (Math Mode)}%
 {Relational \makeimg{less than or equal to}{\ensuremath{\leq}} symbol.}%
 {}

\defgcs{ll}%
 {}%
 {\LaTeX\ Kernel (Math Mode)}%
 {Relational \makeimg{much less than}{\ensuremath{\ll}} symbol.}%
 {}

\defgcs{mid}%
 {}%
 {\LaTeX\ Kernel (Math Mode)}%
 {Relational \ensuremath{\mid} symbol.}%
 {}

\defgcs{models}%
 {}%
 {\LaTeX\ Kernel (Math Mode)}%
 {Relational \ensuremath{\models} symbol.}%
 {}

\defgcs{neq}%
 {}%
 {\LaTeX\ Kernel (Math Mode)}%
 {Relational \makeimg{not equal to}{\ensuremath{\neq}} symbol.}%
 {}

\defgcs{ni}%
 {}%
 {\LaTeX\ Kernel (Math Mode)}%
 {Relational \makeimg{contains element}{\ensuremath{\ni}} symbol.}%
 {}

\defgcs{notin}%
 {}%
 {\LaTeX\ Kernel (Math Mode)}%
 {Relational \makeimg{not element of}{\ensuremath{\notin}} symbol.}%
 {}

\defgcs{parallel}%
 {}%
 {\LaTeX\ Kernel (Math Mode)}%
 {Relational \ensuremath{\parallel} symbol.}%
 {}

\defgcs{prec}%
 {}%
 {\LaTeX\ Kernel (Math Mode)}%
 {Relational \ensuremath{\prec} symbol.}%
 {}

\defgcs{preceq}%
 {}%
 {\LaTeX\ Kernel (Math Mode)}%
 {Relational \ensuremath{\preceq} symbol.}%
 {}

\defgcs{perp}%
 {}%
 {\LaTeX\ Kernel (Math Mode)}%
 {Relational \makeimg{perpendicular to}{\ensuremath{\perp}} symbol.}%
 {}

\defgcs{propto}%
 {}%
 {\LaTeX\ Kernel (Math Mode)}%
 {Relational \makeimg{proportional to}{\ensuremath{\propto}} symbol.}%
 {}

\defgcs{sim}%
 {}%
 {\LaTeX\ Kernel (Math Mode)}%
 {Relational \ensuremath{\sim} symbol.}%
 {}

\defgcs{simeq}%
 {}%
 {\LaTeX\ Kernel (Math Mode)}%
 {Relational \ensuremath{\simeq} symbol.}%
 {}

\defgcs{smile}%
 {}%
 {\LaTeX\ Kernel (Math Mode)}%
 {Relational \ensuremath{\smile} symbol.}%
 {}

\defgcs{sqsubseteq}%
 {}%
 {\LaTeX\ Kernel (Math Mode)}%
 {Relational \ensuremath{\sqsubseteq} symbol.}%
 {}

\defgcs{sqsupseteq}%
 {}%
 {\LaTeX\ Kernel (Math Mode)}%
 {Relational \ensuremath{\sqsupseteq} symbol.}%
 {}

\defgcs{subseteq}%
 {}%
 {\LaTeX\ Kernel (Math Mode)}%
 {Relational \ensuremath{\subseteq} symbol.}%
 {}

\defgcs{succ}%
 {}%
 {\LaTeX\ Kernel (Math Mode)}%
 {Relational \ensuremath{\succ} symbol.}%
 {}

\defgcs{succeq}%
 {}%
 {\LaTeX\ Kernel (Math Mode)}%
 {Relational \ensuremath{\succeq} symbol.}%
 {}

\defgcs{supset}%
 {}%
 {\LaTeX\ Kernel (Math Mode)}%
 {Relational \ensuremath{\supset} symbol.}%
 {}

\defgcs{supseteq}%
 {}%
 {\LaTeX\ Kernel (Math Mode)}%
 {Relational \ensuremath{\supseteq} symbol.}%
 {}

\defgcs{vdash}%
 {}%
 {\LaTeX\ Kernel (Math Mode)}%
 {Relational \ensuremath{\vdash} symbol.}%
 {}

\defgcs{amalg}%
 {}%
 {\LaTeX\ Kernel (Math Mode)}%
 {Binary operator \ensuremath{\amalg} symbol.}%
 {}

\defgcs{ast}%
 {}%
 {\LaTeX\ Kernel (Math Mode)}%
 {Binary operator \makeimg{asterisk}{\ensuremath{\ast}} symbol.}%
 {}

\defgcs{bullet}%
 {}%
 {\LaTeX\ Kernel (Math Mode)}%
 {Binary operator \makeimg{bullet point}{\ensuremath{\bullet}} symbol.}%
 {}

\defgcs{bigcirc}%
 {}%
 {\LaTeX\ Kernel (Math Mode)}%
 {Binary operator \ensuremath{\bigcirc} symbol.}%
 {}

\defgcs{bigtriangledown}%
 {}%
 {\LaTeX\ Kernel (Math Mode)}%
 {Binary operator \ensuremath{\bigtriangledown} symbol.}%
 {}

\defgcs{bigtriangleup}%
 {}%
 {\LaTeX\ Kernel (Math Mode)}%
 {Binary operator \ensuremath{\bigtriangleup} symbol.}%
 {}

\defgcs{cap}%
 {}%
 {\LaTeX\ Kernel (Math Mode)}%
 {Binary operator \ensuremath{\cap} symbol.}%
 {}

\defgcs{cdot}%
 {}%
 {\LaTeX\ Kernel (Math Mode)}%
 {Centred dot \ensuremath{\cdot} symbol.}%
 {}

\defgcs{circ}%
 {}%
 {\LaTeX\ Kernel (Math Mode)}%
 {Circle \ensuremath{\circ} symbol.}%
 {}

\defgcs{cup}%
 {}%
 {\LaTeX\ Kernel (Math Mode)}%
 {Operator \ensuremath{\cup} symbol.}%
 {}

\defgcs{dagger}%
 {}%
 {\LaTeX\ Kernel (Math Mode)}%
 {Binary operator \ensuremath{\dagger} symbol.}%
 {}

\defgcs{ddagger}%
 {}%
 {\LaTeX\ Kernel (Math Mode)}%
 {Binary operator \ensuremath{\ddagger} symbol.}%
 {}

\defgcs{diamond}%
 {}%
 {\LaTeX\ Kernel (Math Mode)}%
 {Binary operator \ensuremath{\diamond} symbol.}%
 {}

\defgcs{div}%
 {}%
 {\LaTeX\ Kernel (Math Mode)}%
 {Division operator \ensuremath{\div} symbol.}%
 {}

\defgcs{mp}%
 {}%
 {\LaTeX\ Kernel (Math Mode)}%
 {Minus or plus operator \makeimg{minus or plus}{\ensuremath{\mp}} symbol.}%
 {}

\defgcs{odot}%
 {}%
 {\LaTeX\ Kernel (Math Mode)}%
 {Operator \ensuremath{\odot} symbol.}%
 {}

\defgcs{ominus}%
 {}%
 {\LaTeX\ Kernel (Math Mode)}%
 {Operator \ensuremath{\ominus} symbol.}%
 {}

\defgcs{oplus}%
 {}%
 {\LaTeX\ Kernel (Math Mode)}%
 {Operator \ensuremath{\oplus} symbol.}%
 {}

\defgcs{oslash}%
 {}%
 {\LaTeX\ Kernel (Math Mode)}%
 {Operator \ensuremath{\oslash} symbol.}%
 {}

\defgcs{otimes}%
 {}%
 {\LaTeX\ Kernel (Math Mode)}%
 {Operator \ensuremath{\otimes} symbol.}%
 {}

\defgcs{pm}%
 {}%
 {\LaTeX\ Kernel (Math Mode)}%
 {Operator \makeimg{plus or minus}{\ensuremath{\pm}} symbol.}%
 {}

\defgcs{setminus}%
 {}%
 {\LaTeX\ Kernel (Math Mode)}%
 {Operator \ensuremath{\setminus} symbol.}%
 {}

\defgcs{sqcap}%
 {}%
 {\LaTeX\ Kernel (Math Mode)}%
 {Operator \ensuremath{\sqcap} symbol.}%
 {}

\defgcs{sqcup}%
 {}%
 {\LaTeX\ Kernel (Math Mode)}%
 {Operator \ensuremath{\sqcup} symbol.}%
 {}

\defgcs{star}%
 {}%
 {\LaTeX\ Kernel (Math Mode)}%
 {Operator \ensuremath{\star} symbol.}%
 {}

\defgcs{times}%
 {}%
 {\LaTeX\ Kernel (Math Mode)}%
 {Operator \ensuremath{\times} symbol.}%
 {}

\defgcs{triangleleft}%
 {}%
 {\LaTeX\ Kernel (Math Mode)}%
 {Binary operator \ensuremath{\triangleleft} symbol.}%
 {}

\defgcs{triangleright}%
 {}%
 {\LaTeX\ Kernel (Math Mode)}%
 {Binary operator \ensuremath{\triangleright} symbol.}%
 {}

\defgcs{uplus}%
 {}%
 {\LaTeX\ Kernel (Math Mode)}%
 {Operator \ensuremath{\uplus} symbol.}%
 {}

\defgcs{vee}%
 {}%
 {\LaTeX\ Kernel (Math Mode)}%
 {Operator \ensuremath{\vee} symbol.}%
 {}

\defgcs{wedge}%
 {}%
 {\LaTeX\ Kernel (Math Mode)}%
 {Operator \ensuremath{\wedge} symbol.}%
 {}

\defgcs{wr}%
 {}%
 {\LaTeX\ Kernel (Math Mode)}%
 {Operator \ensuremath{\wr} symbol.}%
 {}

\defgcs{sum}%
 {}%
 {\LaTeX\ Kernel (Math Mode)}%
 {Summation \ensuremath{\sum} symbol (may take limits).}%
 {}

\defgcs{int}%
 {}%
 {\LaTeX\ Kernel (Math Mode)}%
 {Integral \ensuremath{\int} symbol (may take limits).}%
 {}

\defgcs{oint}%
 {}%
 {\LaTeX\ Kernel (Math Mode)}%
 {Closed path integral \ensuremath{\oint} symbol (may take limits).}%
 {}

\defgcs{prod}%
 {}%
 {\LaTeX\ Kernel (Math Mode)}%
 {Product \ensuremath{\prod} symbol (may take limits).}%
 {}

\defgcs{coprod}%
 {}%
 {\LaTeX\ Kernel (Math Mode)}%
 {Co-product \ensuremath{\coprod} symbol (may take limits).}%
 {}

\defgcs{bigcap}%
 {}%
 {\LaTeX\ Kernel (Math Mode)}%
 {Collection intersection \ensuremath{\bigcap} symbol (may take limits).}%
 {}

\defgcs{bigcup}%
 {}%
 {\LaTeX\ Kernel (Math Mode)}%
 {Collection union \ensuremath{\bigcup} symbol (may take limits).}%
 {}

\defgcs{bigsqcup}%
 {}%
 {\LaTeX\ Kernel (Math Mode)}%
 {Big operator \ensuremath{\bigsqcup} (may take limits).}%
 {}

\defgcs{bigvee}%
 {}%
 {\LaTeX\ Kernel (Math Mode)}%
 {Big operator \ensuremath{\bigvee} (may take limits).}%
 {}

\defgcs{bigwedge}%
 {}%
 {\LaTeX\ Kernel (Math Mode)}%
 {Big operator \ensuremath{\bigwedge} (may take limits).}%
 {}

\defgcs{bigodot}%
 {}%
 {\LaTeX\ Kernel (Math Mode)}%
 {Big operator \ensuremath{\bigodot} (may take limits).}%
 {}

\defgcs{bigotimes}%
 {}%
 {\LaTeX\ Kernel (Math Mode)}%
 {Big operator \ensuremath{\bigotimes} (may take limits).}%
 {}

\defgcs{bigoplus}%
 {}%
 {\LaTeX\ Kernel (Math Mode)}%
 {Big operator \ensuremath{\bigoplus} (may take limits).}%
 {}

\defgcs{biguplus}%
 {}%
 {\LaTeX\ Kernel (Math Mode)}%
 {Big operator \ensuremath{\biguplus} (may take limits).}%
 {}

\defgcs{not}%
 {\meta{symbol command}}%
 {\LaTeX\ Kernel (Math Mode)}%
 {Negates the following symbol. Example: \cmdname{not}\nxglsi{subset} produces
   \ensuremath{\not\subset}.}%
 {%
   \BeginArgList
    \csentryargitem{symbol command} The symbol to be negated.
   \EndArgList
 }

\defgcs{ldots}%
 {}%
 {\LaTeX\ Kernel}%
 {Ellipses \ldots\ symbol.}%
 {}

\defgcs{cdots}%
 {}%
 {\LaTeX\ Kernel (Math Mode)}%
 {Centred ellipses \ensuremath{\cdots} symbol.}%
 {}

\defgcs{vdots}%
 {}%
 {\LaTeX\ Kernel (Math Mode)}%
 {Vertical ellipses \ensuremath{\vdots} symbol.}%
 {}

\defgcs{ddots}%
 {}%
 {\LaTeX\ Kernel (Math Mode)}%
 {Diagonal ellipses \ensuremath{\ddots} symbol.}%
 {}

\defgcs{dotsc}%
 {}%
 {\nxisty{amsmath} (Math Mode)}%
 {Ellipses \ensuremath{\dotsc} for dots with commas.}%
 {}

\defgcs{dotsb}%
 {}%
 {\nxisty{amsmath} (Math Mode)}%
 {Ellipses \ensuremath{\dotsb} for dots with binary
  operators\slash relations.}%
 {}

\defgcs{dotsm}%
 {}%
 {\nxisty{amsmath} (Math Mode)}%
 {Ellipses \ensuremath{\dotsm} for dots with multiplications.}%
 {}

\defgcs{dotsi}%
 {}%
 {\nxisty{amsmath} (Math Mode)}%
 {Ellipses \ensuremath{\dotsi} for dots with integrals.}%
 {}

\defgcs{dotso}%
 {}%
 {\nxisty{amsmath} (Math Mode)}%
 {Ellipses \ensuremath{\dotso} for general dots.}%
 {}

\defgcs{langle}%
 {}%
 {\LaTeX\ Kernel (Math Mode)}%
 {Left-angled \ensuremath{\langle} delimiter.}%
 {}

\defgcs{rangle}%
 {}%
 {\LaTeX\ Kernel (Math Mode)}%
 {Right-angled \ensuremath{\rangle} delimiter.}%
 {}

\defgcs{lfloor}%
 {}%
 {\LaTeX\ Kernel (Math Mode)}%
 {Left floor \ensuremath{\lfloor} delimiter.}%
 {}

\defgcs{rfloor}%
 {}%
 {\LaTeX\ Kernel (Math Mode)}%
 {Right floor \ensuremath{\rfloor} delimiter.}%
 {}

\defgcs{lceil}%
 {}%
 {\LaTeX\ Kernel (Math Mode)}%
 {Left ceil \ensuremath{\lceil} delimiter.}%
 {}

\defgcs{rceil}%
 {}%
 {\LaTeX\ Kernel (Math Mode)}%
 {Right ceil \ensuremath{\rceil} delimiter.}%
 {}

\defgcs{lvert}%
 {}%
 {\nxisty{amsmath} (Math Mode)}%
 {Left vertical bar \ensuremath{\lvert} delimiter.}%
 {}

\defgcs{rvert}%
 {}%
 {\nxisty{amsmath} (Math Mode)}%
 {Right vertical bar \ensuremath{\rvert} delimiter.}%
 {}

\defgcs{lVert}%
 {}%
 {\nxisty{amsmath} (Math Mode)}%
 {Left double vertical bar \ensuremath{\lVert} delimiter.}%
 {}

\defgcs{rVert}%
 {}%
 {\nxisty{amsmath} (Math Mode)}%
 {Right double vertical bar \ensuremath{\rVert} delimiter.}%
 {}

\defgcs{vec}%
 {\marg{\meta{c}}}%
 {\LaTeX\ Kernel (Math Mode)}%
 {Typesets its argument as a vector (defaults to a right arrow
   accent).}%
 {%
   \BeginArgList
    \csentryargitem{c} The character or symbol that represents a vector.
   \EndArgList
 }

\defgcs{forall}%
 {}%
 {\LaTeX\ Kernel (Math Mode)}%
 {\dq{For all} \ensuremath{\forall} symbol.}%
 {}


\defgcs{noindent}%
 {}%
 {\LaTeX\ Kernel}%
 {Suppress the indentation that would usually occur at the start of
  the next paragraph.}%
 {}

\defgcs{refstepcounter}%
 {\marg{\meta{counter}}}%
 {\LaTeX\ Kernel}%
 {Increments the value of the given counter by one and allows the
  counter to be cross-referenced using \nxglsni{ref} and \nxglsi{label}.}%
 {%
   \BeginArgList
    \csentryargitem{counter} The name of the counter that needs
     incrementing.
   \EndArgList
 }

\defgcs{stepcounter}%
 {\marg{\meta{counter}}}%
 {\LaTeX\ Kernel}%
 {Increments the value of the given counter by one.}%
 {%
   \BeginArgList
    \csentryargitem{counter} The name of the counter that needs
     incrementing.
   \EndArgList
 }

\defgcs{setcounter}%
 {\marg{\meta{counter}}\marg{\meta{number}}}%
 {\LaTeX\ Kernel}%
 {Sets the value of a counter.}%
 {%
   \BeginArgList
    \csentryargitem{counter} The name of the counter that needs
     changing.
    \csentryargitem{number} The new value. (Must be an integer.)
   \EndArgList
 }

\defgcs{addtocounter}%
 {\marg{\meta{counter}}\marg{\meta{increment}}}%
 {\LaTeX\ Kernel}%
 {Increments the value of a counter by the given amount.}%
 {%
   \BeginArgList
    \csentryargitem{counter} The name of the counter that needs
     changing.
    \csentryargitem{number} The amount by which to increment the counter. (Must be an integer.)
   \EndArgList
 }

\defgcs{newcounter}%
 {\marg{\meta{counter}}\oarg{\meta{outer counter}}}%
 {\LaTeX\ Kernel}%
 {Defines a new counter.}%
 {%
   \BeginArgList
    \csentryargitem{counter} The name of the new counter.
    \csentryargitem{outer counter} The name of the parent counter.
   \EndArgList
 }


\defgcs{value}%
 {\marg{\meta{counter}}}%
 {\LaTeX\ Kernel}%
 {References the value of the given counter where a number rather
   than a counter name is required.}%
 {%
   \BeginArgList
    \csentryargitem{counter} The name of the counter.
   \EndArgList
 }

\defgcs{parindent}%
 {}%
 {\LaTeX\ Kernel}%
 {A \nxglsi{length} register that stores the indentation at the start of paragraphs.}%
 {}


\defgcs{setlength}%
 {\marg{\meta{register}}\marg{\meta{dimension}}}%
 {\LaTeX\ Kernel}%
 {Sets the value of a length register.}%
 {%
   \BeginArgList
     \csentryargitem{register} The name of the length register (for
       example \nxglsi{parindent}).
     \csentryargitem{dimension} The new dimension.
   \EndArgList
 }

\defgcs{addtolength}%
 {\marg{\meta{register}}\marg{\meta{dimension}}}%
 {\LaTeX\ Kernel}%
 {Adds \meta{dimension} to the value of the given length register.}%
 {%
   \BeginArgList
     \csentryargitem{register} The name of the length register (for
       example \nxglsi{parindent}).
     \csentryargitem{dimension} The dimension to add to the length.
   \EndArgList
 }

\defgcs{currenttime}%
 {}%
 {\nxisty{datetime} package}%
 {Inserts into the output file the time when the \LaTeX\ 
  application created it from the source code.}%
 {}

\defgcs{ddmmyyyydate}%
 {}%
 {\nxisty{datetime} package}%
 {Changes the format of \nxglsi{today} so that it displays the date in
  the form {\ddmmyyyydate\today} (day\slash month\slash year in digits).}%
 {}

\defgcs{subref}%
 {\marg{\meta{label}}}%
 {\nxisty{subcaption} package}%
 {Analogous to \nxglsi{ref} but only references the subfigure
  or subtable caption.}%
 {%
   \BeginArgList
    \csentryargitem{label} The required label that was used in the
     corresponding \nxglsi{label}.
   \EndArgList
 }

\defgcs{eqref}%
 {\marg{\meta{label}}}%
 {\nxisty{amsmath} package}%
 {Short cut for \texttt{(\nxglsi{ref}\marg{\meta{label}})} for
  referencing equations.}%
 {%
   \BeginArgList
    \csentryargitem{label} The required label that was used in the
     corresponding \nxglsi{label}.
   \EndArgList
 }

\defgcs{frenchspacing}%
 {}%
 {\LaTeX\ Kernel}%
 {Switch to French spacing.}%
 {}

\defgcs{nonfrenchspacing}%
 {}%
 {\LaTeX\ Kernel}%
 {Switch to English spacing.}%
 {}

\defgcs{newline}%
 {}%
 {\LaTeX\ Kernel}%
 {Forces a line break.}%
 {}

\defgcs{linebreak}%
 {\oarg{\meta{n}}}%
 {\LaTeX\ Kernel}%
 {Requests a line break, ensuring the paragraph remains justified.
  This may cause excess white space in the paragraph.}%
 {%
   \BeginArgList
    \csentryargitem{n} An integer from~0 to~4 indicating how strongly
    you want a line break to occur. The higher the number, the
    stronger the request to break the line.
   \EndArgList
 }

\defgcs{displaybreak}%
 {\oarg{\meta{n}}}%
 {\nxisty{amsmath} package}%
 {Allows a page break in multi-lined maths environments, such as
  \nxglsi{env-align}.}%
 {%
   \BeginArgList
    \csentryargitem{n} An integer from~0 to~4 indicating how strongly
    you want a page break to occur. The higher the number, the
    stronger the request.
   \EndArgList
 }

\defgsymcs[beginmath]{\openparensym}%
 {}%
 {\LaTeX\ Kernel}%
 {Equivalent to \nxglsni{begin}\marg{math}.}%
 {}

\defgsymcs[endmath]{\closeparensym}%
 {}%
 {\LaTeX\ Kernel}%
 {Equivalent to \nxglsni{end}\marg{math}.}%
 {}

\defgchar
 {openparen}
 {\openparensym}
 {}
 {\LaTeX\ Kernel}
 {%
   Opening parenthesis in text mode or left round bracket
   delimiter in math mode.%
 }
 {}

\defgchar
 {closeparen}
 {\closeparensym}
 {}
 {\LaTeX\ Kernel}
 {%
   Closing parenthesis in text mode or right round bracket
   delimiter in math mode.%
 }
 {}

\defgsymcs[begindispmath]{\opensqsym}%
 {}%
 {\LaTeX\ Kernel (inconsistency corrected in \nxisty{amsmath})}%
 {Starts an unnumbered single-line of displayed maths.}%
 {}

\defgsymcs[enddispmath]{\closesqsym}%
 {}%
 {\LaTeX\ Kernel (inconsistency corrected in \nxisty{amsmath})}%
 {Ends an unnumbered single-line of displayed maths.}%
 {}

\defgchar
 {opensq}
 {\opensqsym}
 {}
 {\LaTeX\ Kernel}
 {\nopostdesc}
 {}

\defgchildchar
 {opt.opensq}
 {opensq}
 {\opensqsym}
 {Open delimiter of an \nxglsi{optional}}

\defgchildchar
 {text.opensq}
 {opensq}
 {\opensqsym}
 {Open square bracket in text mode}

\defgchildchar
 {delimiter.opensq}
 {opensq}
 {\opensqsym}
 {Left square bracket delimiter in math mode}

\defgchar
 {closesq}
 {\closesqsym}
 {}
 {\LaTeX\ Kernel}
 {\nopostdesc}
 {}

\defgchildchar
 {opt.closesq}
 {closesq}
 {\closesqsym}
 {Closing delimiter of an \nxglsi{optional}}

\defgchildchar
 {text.closesq}
 {closesq}
 {\closesqsym}
 {Closing square bracket in text mode}

\defgchildchar
 {delimiter.closesq}
 {closesq}
 {\closesqsym}
 {Right square bracket delimiter in math mode}

\defgsymcs[comma]{\commasym}%
 {}%
 {\LaTeX\ Kernel}%
 {Thin space.}%
 {}

\defgcs{thinspace}%
 {}%
 {\LaTeX\ Kernel}%
 {Thin space.}%
 {}

\defgsymcs[colon]{\colonsym}%
 {}%
 {\LaTeX\ Kernel (Math Mode)}%
 {Medium space.}%
 {}

\defgcs{medspace}%
 {}%
 {\nxisty{amsmath} package}%
 {Medium space.}%
 {}

\defgsymcs[semicolon]{\semicolonsym}%
 {}%
 {\LaTeX\ Kernel (Math Mode)}%
 {Thick space.}%
 {}

\defgcs{thickspace}%
 {}%
 {\nxisty{amsmath} package}%
 {Thick space.}%
 {}

\defgcs{qquad}%
 {}%
 {\LaTeX\ Kernel}%
 {Horizontal spacing command (twice as wide as \nxglsni{quad}).}%
 {}

\defgsymcs[spacesym]{\textvisiblespace}%
 {}%
 {\LaTeX\ Kernel}%
 {(Backslash followed by space character.) Horizontal spacing command.}%
 {}

\defgcs{enspace}%
 {}%
 {\LaTeX\ Kernel}%
 {Horizontal spacing command (half as wide as \nxglsi{quad}).}%
 {}

\defgcs{quad}%
 {}%
 {\LaTeX\ Kernel}%
 {Horizontal spacing command equal to the current font's \protect\htmlref{em}{sec:length} value.}%
 {}

\defgcs{negthinspace}%
 {}%
 {\LaTeX\ Kernel}%
 {Negative thin space.}%
 {}

\defgcs{negmedspace}%
 {}%
 {\nxisty{amsmath} package}%
 {Negative medium space.}%
 {}

\defgcs{negthickspace}%
 {}%
 {\nxisty{amsmath} package}%
 {Negative thick space.}%
 {}

\defgcs{contentsname}%
 {}%
 {Classes or packages that define a table of contents}%
 {Text used for table of contents heading.}%
 {}

\defgcs{listfigurename}%
 {}%
 {Classes or packages that define a list of figures}%
 {Text used for list of figures heading.}%
 {}

\defgcs{listtablename}%
 {}%
 {Classes or packages that define a list of tables}%
 {Text used for list of tables heading.}%
 {}

\defgcs{bibname}%
 {}%
 {Report or book style classes that define a bibliography chapter}%
 {Text used for bibliography chapter heading. (See also
  \nxglsi{refname}.)}%
 {}

\defgcs{refname}%
 {}%
 {Article style classes that define a bibliography section}%
 {Text used for bibliography section heading. (See also
   \nxglsi{bibname}.)}%
 {}

\defgcs{indexname}%
 {}%
 {Classes or packages that define an index section}%
 {Text used for index heading.}%
 {}

\defgcs{figurename}%
 {}%
 {Classes or packages that define figures}%
 {Number prefix used in figure captions.}%
 {}

\defgcs{tablename}%
 {}%
 {Classes or packages that define tables}%
 {Number prefix used in table captions.}%
 {}

\defgcs{partname}%
 {}%
 {Classes or packages that define parts with a number prefix}%
 {Number prefix used in part headings.}%
 {}

\defgcs{chaptername}%
 {}%
 {Classes or packages that define chapters}%
 {Number prefix used in chapter headings.}%
 {}

\defgcs{appendixname}%
 {}%
 {Classes or packages that define chapters}%
 {Number prefix used in appendix headings.}%
 {}

\defgcs{abstractname}%
 {}%
 {Classes or packages that define an \nxglsni{env-abstract}
  environment}%
 {Text used in abstract heading.}%
 {}

\defgcs{addto}%
 {\marg{\meta{command}}\marg{\meta{code}}}
 {\nxisty{babel} package}
 {Adds \meta{code} to the definition of \meta{command}. (See also \nxglsni{appto}.)}
 {%
   \BeginArgList
     \csentryargitem{command} The command that needs patching.
     \csentryargitem{code} The additional code to add to the command
definition.
   \EndArgList
 }

\defgcs{appto}%
 {\marg{\meta{command}}\marg{\meta{code}}}
 {\nxisty{etoolbox} package}
 {Adds \meta{code} to the definition of \meta{command}.}
 {%
   \BeginArgList
     \csentryargitem{command} The command that needs patching.
     \csentryargitem{code} The additional code to add to the command definition.
   \EndArgList
 }

\defgcs{mathrm}%
 {\marg{\meta{maths}}}%
 {\LaTeX\ Kernel (Math Mode)}%
 {Renders \meta{maths} in the predefined maths serif font.}%
 {%
   \BeginArgList
     \csentryargitem{maths} The text on which to apply the font change.
   \EndArgList
 }

\defgcs{mathsf}%
 {\marg{\meta{maths}}}%
 {\LaTeX\ Kernel (Math Mode)}%
 {Renders \meta{maths} in the predefined maths sans-serif font.}%
 {%
   \BeginArgList
     \csentryargitem{maths} The text on which to apply the font change.
   \EndArgList
 }

\defgcs{mathtt}%
 {\marg{\meta{maths}}}%
 {\LaTeX\ Kernel (Math Mode)}%
 {Renders \meta{maths} in the predefined maths typewriter font.}%
 {%
   \BeginArgList
     \csentryargitem{maths} The text on which to apply the font change.
   \EndArgList
 }

\defgcs{mathit}%
 {\marg{\meta{maths}}}%
 {\LaTeX\ Kernel (Math Mode)}%
 {Renders \meta{maths} in the predefined maths italic font.}%
 {%
   \BeginArgList
     \csentryargitem{maths} The text on which to apply the font change.
   \EndArgList
 }

\defgcs{mathbf}%
 {\marg{\meta{maths}}}%
 {\LaTeX\ Kernel (Math Mode)}%
 {Renders \meta{maths} in the predefined maths bold font. (Doesn't
  work with numbers and nonalphabetical symbols.) See also \nxglsi{boldsymbol}.}%
 {%
   \BeginArgList
     \csentryargitem{maths} The text on which to apply the font change.
   \EndArgList
 }

\defgcs{boldsymbol}%
 {\marg{\meta{symbol}}}%
 {\nxisty{amsmath} package (Math Mode)}%
 {Like \nxglsi{mathbf} but also works for numbers and
   many nonalphabetical symbols. (See also \nxglsi{pmb}.)}%
 {%
   \BeginArgList
     \csentryargitem{symbol} The character or symbol on which to apply the font change.
   \EndArgList
 }

\defgcs{pmb}%
 {\marg{\meta{symbol}}}%
 {\nxisty{amsmath} package (Math Mode)}%
 {\dq{Poor man's bold.} Overlays multiple copies of the symbol to
   produce a bold effect for symbols that don't work with
   \nxglsi{boldsymbol}.}%
 {%
   \BeginArgList
     \csentryargitem{symbol} The symbol on which to apply the bold
      effect.
   \EndArgList
 }

\defgcs{mathcal}%
 {\marg{\meta{maths}}}%
 {\LaTeX\ Kernel (Math Mode)}%
 {Typesets its argument in the maths calligraphic font.
  Example: \nxglsni{beginmath}\protect\cmdname{mathcal}\marg{S}\nxglsni{endmath} produces
 \ensuremath{\mathcal{S}}.}%
 {%
   \BeginArgList
    \csentryargitem{maths} The maths to be displayed in a calligraphic
      font.
   \EndArgList
 }

\defgcs{mathbb}%
 {\marg{\meta{maths}}}%
 {\nxisty{amsfonts} package (Math Mode)}%
 {Typesets its argument in the blackboard bold font.
  Example: \nxglsni{beginmath}\protect\cmdname{mathbb}\marg{R}\nxglsni{endmath} produces
  \ensuremath{\mathbb{R}}.}%
 {%
   \BeginArgList
    \csentryargitem{maths} The maths to be displayed in a blackboard
      bold font.
   \EndArgList
 }

\defgcs{mathfrak}%
 {\marg{\meta{maths}}}%
 {\nxisty{amsfonts} package (Math Mode)}%
 {Typesets its argument in Euler Fraktur letters.
  Example: \nxglsni{beginmath}\protect\cmdname{mathfrak}\marg{U}\nxglsni{endmath} produces
  \ensuremath{\mathfrak{U}}.}%
 {%
   \BeginArgList
    \csentryargitem{maths} The maths to be displayed in Euler Fraktur
      letters.
   \EndArgList
 }

\defgcs{DeclareMathOperator}%
 {\marg{\meta{cmd}}\marg{\meta{operator-name}}}%
 {\nxisty{amsmath} package (Preamble Only)}%
 {Defines a new maths operator. The starred version allows limits.}%
 {%
   \BeginArgList
    \csentryargitem{cmd} The name of the new operator command (must
      begin with a backslash).
    \csentryargitem{operator name} The name of the maths operator.
   \EndArgList
 }

\defgcs{frac}%
 {\marg{\meta{numerator}}\marg{\meta{denominator}}}%
 {\LaTeX\ Kernel (Math Mode)}%
 {Displays a fraction.}%
 {%
    \BeginArgList
     \csentryargitem{numerator} The numerator (above the line).
     \csentryargitem{denominator} The denominator (below the line).
    \EndArgList
 }

\defgcs{cfrac}%
 {\oarg{\meta{pos}}\marg{\meta{numerator}}\marg{\meta{denominator}}}%
 {\nxisty{amsmath} (Math Mode)}%
 {Displays a continued fraction.}%
 {%
    \BeginArgList
     \csentryargitem{pos} May be \texttt{l} (left) or \texttt{r}
      (right).
     \csentryargitem{numerator} The numerator (above the line).
     \csentryargitem{denominator} The denominator (below the line).
    \EndArgList
 }

\defgcs{sqrt}%
 {\oarg{\meta{order}}\marg{\meta{operand}}}%
 {\LaTeX\ Kernel (Math Mode)}%
 {Displays a root.}%
 {%
    \BeginArgList
     \csentryargitem{order} The order of the root. (If omitted, a
       square root).
     \csentryargitem{operand} The operand.
    \EndArgList
 }

\defgcs{arabic}%
 {\marg{\meta{counter}}}%
 {\LaTeX\ Kernel}%
 {Displays counter value as an Arabic number. (1, 2, 3, \ldots)}%
 {%
    \BeginArgList
     \csentryargitem{counter} The name of the counter.
    \EndArgList
 }

\defgcs{roman}%
 {\marg{\meta{counter}}}%
 {\LaTeX\ Kernel}%
 {Displays counter value as a lower case Roman number. (i, ii, iii, \ldots)}%
 {%
    \BeginArgList
     \csentryargitem{counter} The name of the counter.
    \EndArgList
 }

\defgcs{Roman}%
 {\marg{\meta{counter}}}%
 {\LaTeX\ Kernel}%
 {Displays counter value as an upper case Roman number. (I, II, III, \ldots)}%
 {%
    \BeginArgList
     \csentryargitem{counter} The name of the counter.
    \EndArgList
 }

\defgcs{alph}%
 {\marg{\meta{counter}}}%
 {\LaTeX\ Kernel}%
 {Displays counter value as a lower case letter. (a, b, c, \ldots, z)}%
 {%
    \BeginArgList
     \csentryargitem{counter} The name of the counter.
    \EndArgList
 }

\defgcs{Alph}%
 {\marg{\meta{counter}}}%
 {\LaTeX\ Kernel}%
 {Displays counter value as an upper case letter. (A, B, C, \ldots, Z)}%
 {%
    \BeginArgList
     \csentryargitem{counter} The name of the counter.
    \EndArgList
 }

\defgcs{fnsymbol}%
 {\marg{\meta{counter}}}%
 {\LaTeX\ Kernel}%
 {Displays counter value as footnote symbol. (\makeimg{asterisk
  dagger double dagger section mark
  paragraph mark double bar double asterisk two single daggers
  two double daggers}{\footnotesymbols})}%
 {%
    \BeginArgList
     \csentryargitem{counter} The name of the counter.
    \EndArgList
 }

\defgenv{bfseries}%
 {}%
 {\LaTeX\ Kernel}%
 {Typesets the environment contents in a bold font.}%
 {}

\defgenv{itshape}%
 {}%
 {\LaTeX\ Kernel}%
 {Typesets the environment contents in an italic font.}%
 {}

\defgenv{em}%
 {}%
 {\LaTeX\ Kernel}%
 {Typesets the environment contents in an emphasized font. (Switches
  to italic\slash slanted if the surrounding font is upright, or switches to
  upright if the surrounding font is italic\slash slanted.)}%
 {}

\defgenv{itemize}%
 {}%
 {\LaTeX\ Kernel}%
 {Unordered list.}%
 {}

\defgenv{enumerate}%
 {}%
 {\LaTeX\ Kernel}%
 {Ordered list.}%
 {}

\defgenv{description}%
 {}%
 {Most class files}%
 {Labelled list.}%
 {}

\defgenv{document}%
 {}%
 {\LaTeX\ Kernel}%
 {The body of the document.}%
 {}

\defgenv{tabular}%
 {\oarg{\meta{v-pos}}\marg{\meta{column specifiers}}}%
 {\LaTeX\ Kernel (Text Mode)}%
 {%
   Environment for lining things up in rows and columns.
   Use \nxglsi{env-array} for math mode.
 }%
 {%
    \BeginArgList
     \csentryargitem{v-pos} Vertical alignment of the entire
      environment with respect to the surrounding baseline. May be
      one of \texttt{t} (top), \texttt{b} (bottom) or \texttt{c}
      (centered).
     \csentryargitem{column specifiers} Indicates how to align the
     columns.
    \EndArgList
 }

\defgenv{array}%
 {\oarg{\meta{v-pos}}\marg{\meta{column specifiers}}}%
 {\LaTeX\ Kernel (Math Mode)}%
 {%
   Environment for lining things up in rows and columns.
   Use \nxglsi{env-tabular} for text mode.
 }%
 {%
    \BeginArgList
     \csentryargitem{v-pos} Vertical alignment of the entire
      environment with respect to the surrounding baseline. May be
      one of \texttt{t} (top), \texttt{b} (bottom) or \texttt{c}
      (centered).
     \csentryargitem{column specifiers} Indicates how to align the
     columns.
    \EndArgList
 }

\defgenv{align}
 {}
 {\nxisty{amsmath} package}
 {Used for numbered aligned equations.}
 {}

\defgenv{align*}
 {}
 {\nxisty{amsmath} package}
 {Used for unnumbered aligned equations.}
 {}

\defgenv{minipage}%
 {\oarg{\meta{pos}}\oarg{\meta{height}}\marg{\meta{width}}}%
 {\LaTeX\ Kernel}%
 {Makes a box with line-wrapped contents. (See also \nxglsi{parbox}.)}%
 {%
   \BeginArgList
    \csentryargitem{pos} The vertical alignment of the box relative to
     the surrounding text. (Centred if omitted.)
    \csentryargitem{height} The height of the box. (If omitted the
     height is the natural height of the contents of the box.)
    \csentryargitem{width} The width of the box.
   \EndArgList
 }

\defgenv{abstract}%
 {}%
 {Most article- or report-style classes, such as \nxicls{scrartcl} or
  \nxicls{scrreprt}. Not usually defined in book-style classes, such
  as \nxicls{scrbook}, but is defined in \nxicls{memoir}}%
 {Displays its contents as an abstract.}%
 {}

\defgenv{thebibliography}%
 {\marg{\meta{widest entry label}}}%
 {Most classes that define sectioning commands}%
 {Bibliographic list. (See also \nxglsi{bibitem} and
  \nxglsi{cite}).}%
 {%
   \BeginArgList
    \csentryargitem{widest entry label} The widest label in the
     bibliography list.
   \EndArgList
 }

\defgenv{otherlanguage}
 {\marg{\meta{language name}}}%
 {\nxisty{babel} package}%
 {Within the environment contents, predefined textual elements, such
  as the date given by \nxglsi{today} or prefixes like \dq{Chapter},
  are set to those supplied by the given language.}%
 {%
   \BeginArgList
    \csentryargitem{language} The language identifier.
   \EndArgList
 }

\defgenv{figure}
 {\oarg{\meta{placement}}}%
 {Most classes that define sectioning commands}%
 {%
   Floats the contents to the nearest location according to the
   preferred placement options, if possible. Within the environment,
   \nxglsi{caption} may be used one or more times, as required.
   The caption will usually include the prefix given by
   \nxglsni{figurename}.
 }%
 {%
   \BeginArgList
    \csentryargitem{placement} The preferred placement.
   \EndArgList
 }

\defgenv{sidewaysfigure}%
 {}%
 {\nxisty{rotating} package}%
 {Like the \nxglsi{env-figure} environment but rotates the entire
  figure (including caption) sideways.}%
 {}

\defgenv{table}
 {\oarg{\meta{placement}}}%
 {Most classes that define sectioning commands}%
 {%
   Floats the contents to the nearest location according to the
   preferred placement options, if possible. Within the environment,
   \nxglsi{caption} may be used one or more times, as required.
   The caption will usually include the prefix given by
   \nxglsni{tablename}.
 }%
 {%
   \BeginArgList
    \csentryargitem{placement} The preferred placement.
   \EndArgList
 }

\defgenv{sidewaystable}%
 {}%
 {\nxisty{rotating} package}%
 {Like the \nxglsi{env-table} environment but rotates the entire
  table (including caption) sideways.}%
 {}

\defgenv{subfigure}%
 {\oarg{\meta{pos}}\marg{\meta{width}}}%
 {\nxisty{subcaption} package}%
 {Used to form a subfigure within a \nxglsi{env-figure}
  environment. The \nxglsi{caption} command may be used in this
  environment to produce a subcaption.}%
 {%
   \BeginArgList
    \csentryargitem{pos} The vertical alignment of the box relative to
     the surrounding text. (Centred if omitted.)
    \csentryargitem{width} The width of the box.
   \EndArgList
 }

\defgenv{subtable}%
 {\oarg{\meta{pos}}\marg{\meta{width}}}%
 {\nxisty{subcaption} package}%
 {Used to form a subtable within a \nxglsni{env-table}
  environment. The \nxglsi{caption} command may be used in this
  environment to produce a subcaption.}%
 {%
   \BeginArgList
    \csentryargitem{pos} The vertical alignment of the box relative to
     the surrounding text. (Centred if omitted.)
    \csentryargitem{width} The width of the box.
   \EndArgList
 }

\defgenv{equation}%
 {}%
 {\LaTeX\ Kernel}%
 {Displays its contents as a single-lined numbered equation.}%
 {}

\defgenv{displaymath}%
 {}%
 {\LaTeX\ Kernel}%
 {Displays its contents as a single-lined unnumbered equation.}%
 {}

\defgenv{math}%
 {}%
 {\LaTeX\ Kernel}%
 {Sets its contents in in-line math mode.}%
 {}

\defgenv{cases}%
 {}%
 {\nxisty{amsmath} package (Math Mode)}%
 {%
   Like the \nxglsi{env-array} environment, but adds a left
   brace start delimiter and an invisible end delimiter.
 }%
 {}

\defgenv{matrix}%
 {}%
 {\nxisty{amsmath} package (Math Mode)}%
 {%
   Like the \nxglsi{env-array} environment, but doesn't have an
   argument.
 }%
 {}

\defgenv{pmatrix}%
 {}%
 {\nxisty{amsmath} package (Math Mode)}%
 {%
   Like the \nxglsi{env-array} environment, but doesn't have an
   argument and adds round bracket delimiters.
 }%
 {}

\defgenv{smallmatrix}%
 {}%
 {\nxisty{amsmath} package (Math Mode)}%
 {%
   Like the \nxglsi{env-array} environment but doesn't have an
   argument and is designed for in-line maths.
 }%
 {}

\defgenv{bmatrix}%
 {}%
 {\nxisty{amsmath} package (Math Mode)}%
 {%
   Like the \nxglsi{env-array} environment, but doesn't have an
   argument and adds square bracket delimiters.
 }%
 {}

\defgenv{Bmatrix}%
 {}%
 {\nxisty{amsmath} package (Math Mode)}%
 {%
   Like the \nxglsi{env-array} environment, but doesn't have an
   argument and adds curly brace delimiters.
 }%
 {}

\defgenv{vmatrix}%
 {}%
 {\nxisty{amsmath} package (Math Mode)}%
 {%
   Like the \nxglsi{env-array} environment, but doesn't have an
   argument and adds single vertical bar delimiters.
 }%
 {}

\defgenv{Vmatrix}%
 {}%
 {\nxisty{amsmath} package (Math Mode)}%
 {%
   Like the \nxglsi{env-array} environment, but doesn't have an
   argument and adds double vertical bar delimiters.
 }%
 {}

\defgcs{raggedright}%
 {}%
 {\LaTeX\ Kernel}%
 {Ragged-right paragraph justification.}%
 {}

\defgcs{raggedleft}%
 {}%
 {\LaTeX\ Kernel}%
 {Ragged-left paragraph justification.}%
 {}

\defgcs{frontmatter}
 {}%
 {Most book-style classes, such as \nxicls{scrbook}}%
 {Switches to lower case Roman numeral page numbering. Also suppresses
   chapter and section numbering, but still adds unstarred sectional
   units to the table of contents. (See also \nxglsi{mainmatter}
   and \nxglsi{backmatter}.)}%
 {}

\defgcs{mainmatter}
 {}%
 {Most book-style classes, such as \nxicls{scrbook}}%
 {Switches to Arabic page numbering and enables
   chapter and section numbering. (See also
   \nxglsi{frontmatter} and \nxglsi{backmatter}.)}%
 {}

\defgcs{backmatter}
 {}%
 {Most book-style classes, such as \nxicls{scrbook}}%
 {Suppresses chapter and section numbering, but still adds unstarred
  sectional units to the table of contents. (See also \nxglsi{frontmatter}
  and \nxglsi{mainmatter}.)}%
 {}

\defgenv{dinglist}
 {\marg{\meta{number}}}
 {\nxisty{pifont} package}
 {A list where the item marker is given by character \meta{number}
  in the Zapf Dingbats font.}
 {%
   \BeginArgList
    \csentryargitem{number} The character code for the item marker.
   \EndArgList
 }

\defgcs{notag}%
 {}%
 {\nxisty{amsmath} package}%
 {Suppresses equation numbering for the current row in environments such as
   \nxglsi{env-align}.}%
 {}

\defgcs{tag}%
 {\marg{\meta{tag}}}%
 {\nxisty{amsmath} package}%
 {Overrides equation numbering for the current row in environments such as
   \nxglsi{env-align}.}%
 {%
    \BeginArgList
     \csentryargitem{tag} The replacement for the equation number.
    \EndArgList
 }

\defgcs{ding}%
 {\marg{\meta{n}}}%
 {\nxisty{pifont} package}%
 {Inserts PostScript ZapfDingbats character with code \meta{n},
   which must be an integer.}%
 {%
   \BeginArgList
    \csentryargitem{n} The decimal code of the required character.
   \EndArgList
 }

\defgcs{color}%
 {\oarg{\meta{model}}\marg{\meta{specs}}}%
 {\nxisty{color} and \nxisty{xcolor} packages}%
 {A declaration that switches the current foreground colour to the given specification.}%
 {%
    \BeginArgList
      \csentryargitem{model} The colour model, for example
      \texttt{rgb} or \texttt{named}.
      \csentryargitem{specs} The colour specification for the
      given model. For example, if the \texttt{rgb} model is chosen,
      \meta{specs} must be a comma-separated list of three numbers
      each between~0 and~1.
    \EndArgList
 }

\defgcs{textcolor}%
 {\oarg{\meta{model}}\marg{\meta{specs}}\marg{\meta{text}}}%
 {\nxisty{color} and \nxisty{xcolor} packages}%
 {Sets \meta{text} with the foreground colour according to the given \meta{specs}.}%
 {%
    \BeginArgList
      \csentryargitem{\meta{model}} The colour model, for example
      \texttt{rgb} or \texttt{named}.
      \csentryargitem{\meta{specs}} The colour specification for the
      given model. For example, if the \texttt{rgb} model is chosen,
      \meta{specs} must be a comma-separated list of three numbers
      each between~0 and~1.
      \csentryargitem{\meta{text}} The text to be displayed in the
      given colour.
    \EndArgList
 }

\defgcs{tabcolsep}
 {}%
 {\LaTeX\ Kernel}%
 {\nxGls{length} register specifying half the gap between
   columns in a \nxglsni{env-tabular} environment.}%
 {}

\defgcs{arraycolsep}
 {}%
 {\LaTeX\ Kernel}%
 {\nxGls{length} register specifying half the gap between
   columns in an \nxglsni{env-array} environment.}%
 {}


\defgcs{labelformat}%
 {\marg{\meta{ctr}}\marg{\meta{defn}}}%
 {\nxisty{fncylab} package}%
 {Defines how the label for the counter \meta{ctr} should be
 formatted. The definition \meta{defn} should use \texttt{\nxglsi{hashchar}1} to
 indicate the label value.}%
 {%
   \BeginArgList
     \csentryargitem{ctr} The name of the counter.
     \csentryargitem{defn} How to display the value of the counter.
   \EndArgList
 }

\defgcs{thefootnote}%
 {}%
 {\LaTeX\ Kernel}%
 {Displays the current value of the \nxicounter{footnote} counter}%
 {}

\defgcs{thechapter}%
 {}%
 {\LaTeX\ Kernel}%
 {Displays the current value of the \nxicounter{chapter} counter}%
 {}

\defgcs{thesection}%
 {}%
 {\LaTeX\ Kernel}%
 {Displays the current value of the \nxicounter{section} counter}%
 {}

\defgcs{thepage}%
 {}%
 {\LaTeX\ Kernel}%
 {Displays the current value of the \nxicounter{page} counter}%
 {}

\defgcs{thefigure}%
 {}%
 {\LaTeX\ Kernel}%
 {Displays the current value of the \nxicounter{figure} counter}%
 {}

\defgcs{thetable}%
 {}%
 {\LaTeX\ Kernel}%
 {Displays the current value of the \nxicounter{table} counter}%
 {}

\defgcs{ignorespaces}
 {}
 {\LaTeX\ Kernel}
 {Used in begin environment code to suppress any spaces occurring 
  at the start of the environment (see also \nxglsni{ignorespacesafterend}).}
 {}

\defgcs{ignorespacesafterend}%
 {}%
 {\LaTeX\ Kernel}%
 {Used in end environment code to suppress any spaces following the
  end of the environment.}%
 {}

\defgcs{textbar}%
 {}%
 {\LaTeX\ Kernel}%
 {Vertical bar \textbar\ symbol.}%
 {}

\defgcs{the}%
 {\meta{register}}%
 {\LaTeX\ Kernel}%
 {Displays the value of the given register (such as a \nxglsi{length} register).
  Not to be confused with \cmdname{the}\meta{ctr} commands, such as \nxglsni{thefigure}.}%
 {%
   \BeginArgList
     \csentryargitem{register} The name of the register.
   \EndArgList
 }

\defgcs{substack}%
 {\marg{\meta{maths}}}%
 {\nxisty{amsmath} package}%
 {Can be used to produce a multiline subscript or superscript. Lines
  are separated using \nxglsi{tab.dbbackslashchar}.}%
 {%
   \BeginArgList
    \csentryargitem{maths} The subscript or superscript maths with rows
      ended using \nxglsi{tab.dbbackslashchar}
   \EndArgList
 }

\defgcs{DeclareCaptionLabelFormat}
 {\marg{\meta{name}}\marg{\meta{code}}}
 {\nxisty{caption}}
 {Used to defined your own caption label formats.}
 {%
   \BeginArgList
     \csentryargitem{name} The name used to identify this new format
     \csentryargitem{code} The code that defines this format
     (\nxglsi{hashchar}1 gets replaced with the name and \nxglsni{hashchar}2
     gets replaced with the reference number.)
   \EndArgList
 }

\defgcs{captionsetup}
 {\oarg{\meta{float type}}\marg{\meta{options}}}
 {\nxisty{caption} package}
 {Used to set up the options affecting float captions.}
 {%
   \BeginArgList
     \csentryargitem{float type} The float type.
     \csentryargitem{options} A \meta{key}=\meta{value}
     comma-separated list of options.
   \EndArgList
 }



\newacr{ctan}{CTAN}
 {the Comprehensive \TeX\ Archive Network}
 {http://mirror.ctan.org/}

\newacr{tug}{TUG}%
 {\TeX\ User Group}%
 {http://tug.org/}

\newacr{uktug}{UK TUG}%
 {UK \TeX\ User Group}%
 {http://uk.tug.org/}

\newacr{ukfaq}{UK FAQ}%
 {UK List of \TeX\ Frequently Asked Questions}%
 {http://www.tex.ac.uk/faq}

\newacrx{tds}{TDS}{\TeX\ Directory Structure}

\newacrnocite{gui}{GUI}{graphical user interface}

\definecolor{midgray}{gray}{0.4}

\indexpreamble{\latexhtml{Page numbers}{Locations} in \textit{italic} indicate the primary
reference. \latexhtml{Page numbers}{Locations} in \textbf{bold} indicate the entry
definition in the summary.}

\newcommand{\xtableref}[1]{\objectref{Table}{#1}}

% \figconts[graphics opts]{image}{caption}{label}
\newcommand{\figconts}[4][]{%
  \begin{makeimage}\end{makeimage}\figuretop{#4}
  \centering
  \incPgfOrGraphics[#1]{#2}%
  #3% caption
  \label{#4}%
}

\newcommand{\latexbook}{\emph{\LaTeX: A Document Preparation System}~\cite{lamport94}}
\newcommand{\latexcomp}{\emph{The \LaTeX\ Companion}~\cite{goossens94}}
\newcommand{\latexguide}{\emph{A Guide to \LaTeX}~\cite{kopka95}}
\newcommand{\latexgraphic}{\emph{The \LaTeX\ Graphics Companion}~\cite{goossens97}}
\newcommand{\latexweb}{\emph{The \LaTeX\ Web Companion}~\cite{goossens99}}

\newcommand{\baseurl}{http://www.dickimaw-books.com}
\newcommand{\packageurl}{\baseurl/latex/packages}
\newcommand{\exerciseurl}{\baseurl/latex/novices/html/exercises}

\newcommand{\pderiv}[2]{\frac{\partial #1}{\partial #2}}
\newcommand{\e}{\mathrm{e}}

\newcommand{\koma}{KOMA-Script}

\newcommand*{\backcovertext}{If you choose to buy a copy of this book, 
Dickimaw Books asks for your support through buying the Dickimaw
Books edition to help cover costs.}

\copyrighttext{%
  Copyright \textcopyright\ 2004  Nicola L.~C. Talbot

  Permission is granted to copy, distribute and/or modify this
  document under the terms of the GNU Free Documentation License,
  Version 1.2 or any later version published by the Free Software
  Foundation; with no Invariant Sections, no Front-Cover Texts, and
  one Back-Cover Text: \dq{\backcovertext} A copy of the license is included in
  the section entitled \htmlref{\dq{GNU Free Documentation
  License}}{sec:fdl}.

  \doifnotbook
  {%
     The base URL for this document is: \url{\baseurl/latex/novices/}
  }%
}

\renewcommand{\glossaryname}{Summary}

\begin{document}
\maketitle

\frontmatter
\setcounter{tocdepth}{2}
\setnode{contents}
\tableofcontents
\setnode{listoffigures}
\listoffigures
\setnode{listoftables}
\listoftables

\setnode{listofexercises}
\listofexercises

\mainmatter

%%%%%%%%%%%%%%%% INTRODUCTION %%%%%%%%%%%%%%%%%%%%%%%%%%%%

\setnode{introduction}
\chapter{Introduction}
\label{ch:intro}

The aim of this book is to introduce \LaTeX\ to a non-technical
person.  \LaTeX\ is excellent for producing professional looking
documents, however it is a \emph{language} not a word 
processor\faq{Why is TeX not a WYSIWYG system?}{notWYSIWYG}, so
it can take a bit of getting used to, particularly if you have never
had any experience using programming languages.  

\LaTeX\ does take a while to learn, so why should you use it? Here are
a few reasons but it is not an exhaustive list:

\LaTeX\ is far better at typesetting mathematical equations
than word processors. I~wrote my Ph.D. thesis back in the days of
\LaTeX2.09 (the old version of \LaTeX) and given the high quantity of
mathematics that I~had to typeset, it would have taken me considerably
longer to write it in a word processor, and the resulting document
wouldn't have looked nearly as good. Even Microsoft have
acknowledged \glsterm{tex}'s high-quality mathematical
typography~\cite{sargentiii}.

\xminisec{Example:}
Here's an equation taken from some kernel survival analysis:
\begin{result}[latexintromaths.html]
\newcommand{\me}{\mathrm{e}}
\newcommand{\pfrac}[2]{\frac{\partial #1}{\partial #2}}
\[
\pfrac{^2\mathcal{L}}{{z_i^\rho}^2} =
-\pfrac{\rho_i}{z_i^\rho}\left(\pfrac{v_i}{\rho_i}\frac{\me^{v_i}}{1-\me^{v_i}}
+v_i\frac{\me^{v_i}\pfrac{v_i}{\rho_i}(1-\me^{v_i})+\me^{2v_i}\pfrac{v_i}{\rho_i}}{(1-\me^{v_i})^2}\right)
\]
\end{result}
(You can find out how to create this equation \latexhtml{on
page~\pageref{introeq} in \sectionref{sec:delimiters}}{\htmlref{later}{introeq}}.)

That's all very well and good if you want to typeset some equations,
but if your work doesn't involve maths, does that mean that \LaTeX\
is not for you? Although I~am a mathematician, I~have written plenty
of documents with no maths in at all, including prose, poetry,
newsletters, posters and brochures, but I~still opt for \LaTeX\
because using \LaTeX\ ensures consistent formatting, and the style
of the document can be completely changed by simply using a
different \glsterm{cls}, or loading additional \htmlref{packages}{sec:packages}. This
means that I~can concentrate on writing the document, rather than
worrying about how it will look. It also means that if, after having
written a 200 page document, I~then find that I~need to change all
the figure captions so that they are labelled \dq{Fig} instead of
\dq{Figure}, all I~need to do is edit a single line, rather than
going through 200 pages to individually edit every single figure
caption.\footnote{Sure, you could use a search and replace function,
but a sweeping replace-all can have unexpected side effects. For
example, your document may include the sentence, \dq{Figures from
the last quarter showed improvement}, which would get changed to,
\dq{Figs from the last quarter showed improvement}.}

Serious fiction writers are taught never to remind the reader that
they're reading a book. Poor formatting is just as much a reminder
of this as authorial intrusion.

\LaTeX\ makes it very easy to cross-reference chapters, sections,
equations, figures, tables etc, and it also makes it very easy to
generate a table of contents, list of figures, list of tables,
index, glossary\footnote{Glossaries are covered in \latexthesis.}\ and 
bibliography. You don't need to worry about numbering anything,
as this is done automatically, which means that you can insert new
sections or swap sections around without having to worry about
updating all the section numbering etc. Furthermore, if you use
\BiBTeX\footnote{Automating bibliographies is covered in
\latexthesis.}\ in combination with \LaTeX, and you have, say, 100
or more citations, it doesn't matter if you are then told that the
citations have to be re-ordered (say, in order of citation rather
than alphabetically). All that is required is a minor edit to change
the appropriate style file rather than ploughing through the entire
document changing all the citations by hand.

When you are editing a document using a word processor, the
word processor has to work out how to reformat the document every time
you type something.  If you have a large document with a great many
inserted objects (such as figures and equations), the response to
keyboard input can become very slow.  You may find that after typing
a few words you will have to wait until the computer catches up
before you can see what you have typed.  With \LaTeX\ you type in your
code using an ordinary text editor.  The document 
doesn't get formatted until you pass it to \LaTeX, which means that 
you are not slowed down by constant reformatting.

Lastly, there's the fact that \LaTeX\ follows certain typographical
rules, so you can leave most of the typesetting to \LaTeX. You
rarely need to worry about minor things such as
\glstext{intersentencespacing}. The default is English spacing, but
if you have a publisher who disapproves of this, you can switch if
off with a single command. (See \sectionref{sec:intersentencespacing}.)

\LaTeX\ will also automatically deal with
f-ligatures.\footnote{Ligatures can be suppressed using the
\isty{microtype} package if necessary} That is, if any of the
following combination of letters are found: \texttt{fl},
\texttt{ffl}, \texttt{ff}, \texttt{fi}, \texttt{ffi}, they will
automatically be converted into the corresponding ligatures:
% Using CMR as the ligatures are more noticeable than with anttor.
% make latex2html convert this text into an image
\makeimg{fl ligature}{\cmr{fl}},
\makeimg{ffl ligature}{\cmr{ffl}},
\makeimg{ff ligature}{\cmr{ff}},
\makeimg{fi ligature}{\cmr{fi}},
\makeimg{ffi ligature}{\cmr{ffi}}.
  Note the difference between
\makeimg{fl-ligature u ffi-ligature e r}{\cmr{fluffier}} (2
ligatures) and \makeimg{fluffier}{\cmr{f{}l{}uf{}f{}i{}er}} (no ligatures).  These
points may seem minor but they all contribute towards the impact of
the entire document.  When writing technical documents, the
presentation as well as the content is important.  All too often
examiners or referees are put off reading a document because it is
badly formatted.  This provokes an immediate negative reaction and
provides little desire to look favourably upon your work.

To give you an idea of what you can do with \LaTeX, this book was
written in \LaTeX.\footnote{The source code is available at
\url{\baseurl/latex/novices/}, but it really is
\emph{not} the place to start if you are a beginner, as it contains
\LaTeX\ and \glsterm{perl} code beyond the scope of this tutorial.}  The PDF
versions (including \ifbookorother{this}{the} paperback version) were generated using PDF\LaTeX\ and \appname{makeindex} and
the HTML version was generated using the
\latexhtml
 {\LaTeX2HTML\footnote{\url{http://www.latex2html.org/}}}%
 {\htmladdnormallink{\LaTeX2HTML}{http://www.latex2html.org/}}%
\faq{Conversion from (La)TeX to HTML}{LaTeX2HTML} converter. 

For more reasons as to why you might want to use \LaTeX\ instead
of a word processor, have a look at
\ifbookorother
  {\url{http://www.ctan.org/what_is_tex.html\#whytex}.}%
  {\htmladdnormallink{Why TeX?}{http://www.ctan.org/what_is_tex.html\#whytex}}

\setnode{texdoc}
\section{Class and Package Documentation}
\label{sec:texdoc}

There are hundreds of \htmlref{classes}{sec:cls} and
\htmlref{packages}{sec:packages} available on \gls{ctan}. These are
made available by many volunteers. Some provide detailed
documentation to accompany their contribution, while others only
provide a few notes in a README file or comments in the source
files. This book only provides an introductory
look at a small selection of these contributions. If you want
further details on how to use a particular class or package you
should check the documentation that accompanies it. You can use the
\iappname{texdoc} application to search for the documentation. This
is a command line application, which means you need a terminal or
command prompt (see \sectionref{sec:terminal}).

To use \appname{texdoc}, you need to
type (at the command prompt) \appname{texdoc} followed by a space followed by the name of
the class or package you want information about. 
For example, to read the \icls{memoir}
documentation, type the following at the \htmlref{command prompt}{sec:terminal} (press the
return/enter key \enter\ at the end of the line):
\begin{verbatim}
texdoc memoir
\end{verbatim}
Some packages come with more than one set of documentation. For
example, the \isty{glossaries} package comes with the main user
manual, a short guide for beginners and the documented code for
advanced users. Just doing
\begin{verbatim}
texdoc glossaries
\end{verbatim}
will display the advanced documented code. To list all available
documentation for a package, use the \texttt{-l} option:
\begin{verbatim}
texdoc -l glossaries
\end{verbatim}
Then type the number corresponding to the file you want to view. If
you can remember the file name (for example \texttt{glossaries-user}) 
you can type that next time you want to view it:
\begin{verbatim}
texdoc glossaries-user
\end{verbatim}
There is also a \glsterm{perl}/Tk-based \gls{gui} called \iappname{texdoctk},
which is distributed with \reportlinebreak\texdistro{TeX~Live},
that you can use instead of \appname{texdoc} if you can't work out
how to use a \glsterm{terminal} or prefer a \gls*{gui} approach.

Failing that, you can also check on \gls{ctan} using the URL
\texttt{ctan.org\slash pkg\slash}\meta{name}, where \meta{name} is the name of
the package or class. For example, if you want to look up the
documentation for the \icls{memoir} package, you can find it at
\url{http://ctan.org/pkg/memoir} or go to
\url{http://mirror.ctan.org/} and search for the package or class.

Another alternative recently made available is to use the URL
\texttt{texdoc.net\slash pkg\slash}\meta{name}. For example, 
\url{http://texdoc.net/pkg/memoir} will 
fetch the documentation for the \icls{memoir} class.

However, it's better to use \appname{texdoc} or \appname{texdoctk} to 
read the documentation installed with the
class or package on your computer to ensure it
matches the installed class or package version.

Note that it is important to remember that the \glsterm{tex} world is mostly
supported by volunteers. \gls{ctan} itself is maintained by a very small
group (currently two people). It's not like a commercial company
with 24/7 support and hundreds of paid employees constantly updating
the software. At its core, \TeX\ is a community effort. While some
volunteers actively maintain and update their classes or packages,
some people move on to other things and stop maintaining their work.
Occasionally, if the class or package is popular, someone else
might take over maintenance. There is no dedicated helpdesk to go
to, but there are many ways of getting help, see
\appendixref{ch:help}

\setnode{overview}
\section{Overview}

This document is structured as follows:

\begin{description}

\item[\chapterref{ch:def}]
defines terms that will be used throughout this document.  I
strongly suggest that you look through this
\latexhtml{chapter}{section} before you start so that you understand
the terminology used in this document. At the very least, you should
read the first part that details how corresponding input and output
is displayed in this document\dash you need to understand the
difference between \dq{input} (source code) and \dq{output} (how the
source code will appear in the typeset document).\html{ Note that it
is not always possible to reproduce an exact replica of the output
in the HTML version, so if in doubt, look at one of the PDF versions
instead, which can be downloaded from this 
book's \htmladdnormallink{home page}{../index.html}.}

\item[\chapterref{ch:tex2pdf}] details the software that you will
need to use \LaTeX\ and describes how to use the software.

\item[\chapterref{ch:simpledoc}] shows you how to create a very
basic document.  

\item[\chapterref{ch:sections}] 
shows you how to create chapters and other sectional units so that
you end up with a fully structured document.

\item[\chapterref{ch:graphicx}]
shows you how to include external image files and how to scale and
rotate text.

\item[\chapterref{ch:floats}]
describes how to create figures and tables.

\item[\chapterref{ch:newcom}]
describes how to define your own commands, and redefine existing
commands.

\item[\chapterref{ch:maths}] describes how
to typeset mathematics.  

\item[\chapterref{ch:newenv}] describes how to define new
environments.  

\item[\chapterref{ch:counters}]
discusses how numbers are stored in counters, how to change their
values, and how to define your own counter.

\item[\appendixref{ch:installsty}] shows you how
to download and install additional
packages that weren't installed with your \TeX\ distribution.

\item[\appendixref{ch:errors}]
documents possible errors you may encounter, and gives advice on how
to fix them.

\item[\appendixref{ch:help}] gives pointers on where to go for help.

\end{description}

Throughout this document there are pointers to related topics
in the \gls{ukfaq}.
\latexhtml{These are displayed in the margin in
square brackets, as illustrated on the 
\ifbookorother{\ifthispageodd{right}{left}}{right}.}{%
These are displayed in the text like this:}%
\faq{What is LaTeX?}{latex} You may find these resources useful in
answering related questions that are not covered in this 
book.
\begin{latexonly}
\doifbook{%
  To find the resources, go to \url{http://www.tex.ac.uk/faq} and
  either look for the question title in the list, or enter a keyword
  in the search field.%
}

This book and associated files, including  
solutions to the exercises, are available on-line at:
\url{\baseurl/latex/novices/}. 
\doifnotbook
{%
  The links in this document are colour-coded: internal links are 
  blue, external links are magenta.
}%
\end{latexonly}

\setnode{reading}
\section{Recommended Reading}

This document is designed as an introductory text, not a
comprehensive guide.  For further reading try some of the
following\faq{Books on LaTeX}{latex-books}:

\latexbook\ is the user guide and reference manual for \LaTeX, and
is a good basic text for anyone starting out, however it doesn't
cover AMS\TeX\faq{What are the AMS packages?}{AMSpkg}, so anyone who
needs to typeset more than basic mathematics may prefer either
\latexguide\ or \latexcomp.  Both these books cover AMS\TeX,
\iBiBTeX\ and \appname{makeindex}.

In the same series as \emph{The \LaTeX\ Companion}, there is also 
\latexgraphic\ which details how to illustrate
documents with \LaTeX\ and PostScript, including a chapter on colour
(coloured text, background, tables and slides).  This is recommended
to anyone who is contemplating heavy use of graphics, but you do need
a basic knowledge of \LaTeX\ before delving into it.  

The final book in the \dq{Companion} series which you may find useful
is \latexweb. This is recommended for those interested in creating
documents for the web, either as HTML or PDF.  It details how
to convert \LaTeX\ documents into HTML using various applications
such as \appname{LaTeX2HTML} and \appname{TeX4ht}, and how to create
PDF documents using \iPDFLaTeX\faq{What is PDFTeX?}{whatpdftex},
including how to create active links within your document using the
\isty{hyperref} package.

There are two new \LaTeX\ books that I haven't read but have
been recommended to me: \emph{\LaTeX\ Beginner's
Guide}~\cite{kottwitz} and \emph{\LaTeX\ and
Friends}~\cite{vandongen}.

Note that the \gls{uktug} has a 25\% book discount scheme
for members. See \url{http://uk.tug.org/membership} for more details
of that and other associated benefits. If you're not in the
UK, have a look at \url{http://www.tug.org/usergroups.html} to see
if there is a local user group in your area.

There is also a wealth of \LaTeX-related information on the
Internet\faq{How to get help}{gethelp}. \reportlinebreak\Gls{ctan} is a good place
to start. You can check the on-line catalogue~\cite{texcat} for
information about available software and, as mentioned earlier, there is also the list of
\htmladdnormallink{frequently asked
questions}{http://www.tex.ac.uk/faq} which I~recommend you try if
you have any queries. See also \appendixref{ch:help}

%%%%%%%%%%%%%%%%%%%%%%%%% DEFINITIONS %%%%%%%%%%%%%%%%%%%%%%%%%%%%%%

\setnode{definitions}
\printkeywords
\html{\documentclass[html]{novices}

\input{novices-main}
}


\glsaddall[types=keywords]

%%%%%%%%%%%%%%%%%%%%%% SOURCE TO PDF %%%%%%%%%%%%%%%%%%%%%%%%%%

\setnode{fromsource2output}
\chapter{From Source Code to Typeset Output}
\label{ch:tex2pdf}

Every time you want to create or edit a \LaTeX\ document, there are three
basic steps you will always need to follow:

\begin{enumerate}
\item\label{itm:step1} Write or edit the \gls{source}.

\item\label{itm:step2} Pass the source code to the \iappname{latex} or \iappname{pdflatex} application (\dq{\LaTeX\ the document}).

\begin{itemize}
\item If there are any error messages, return to \objectref{Step}{itm:step1}.
\item If there are no error messages, a 
\glslink{output}{PDF file}\indexPDF\ is created.
\end{itemize}

\item\label{itm:step3} View the \glslink{output}{PDF file} to check
the result.  If you need to modify your document, go back to
\objectref{Step}{itm:step1}. 
\end{enumerate}

You will therefore need:

\begin{enumerate}
\item A text editor  (to perform \objectref{Step}{itm:step1}). For
example \iappnamelink[,]{Vim}{http://www.vim.org/} 
\iappnamelink{Emacs}{http://www.gnu.org/software/emacs/} or
\iappnamelink[.]{Gedit}{http://projects.gnome.org/gedit/}

\refstepcounter{object}\label{obj:MiKTeX}%
\item The \gls{tex} software\faq{(La)TeX for different
machines}{TeXsystems} (to perform \objectref{Step}{itm:step2}). If
you don't already have \TeX\ on your machine, you will need to
install it. The most convenient way to do this is
to install from the \TeX\ Collection DVD ROM, which is distributed
to all \gls{tug} members, but you can also download and install free
\TeX\ distributions, such as \itexdistro{TeX~Live},
\itexdistro{MiKTeX} or \itexdistro{MacTeX}, from the Internet (see
\xpageref{below}{obj:install}).  There is also \itexdistro{proTeXt},
an enhancement of \texdistro{MiKTeX} that aims to be an
easy-to-install \TeX\ Distribution.  For more information including
up-to-date links, go to \url{http://www.ctan.org/starter.html}.

\item A PDF viewer (to perform
\objectref{Step}{itm:step3}). For example 
\iappnamelink[,]{Adobe Reader}{http://www.adobe.com/} 
\iappnamelink[,]{Sumatra}{http://blog.kowalczyk.info/software/sumatrapdf/free-pdf-reader.html}
\iappnamelink{Evince}{http://projects.gnome.org/evince/} or
\iappnamelink[.]{Okular}{http://okular.kde.org/}

\end{enumerate}

This can be rather complicated for a beginner, especially for those
with no experience writing computer code. Fortunately, there are
some all-in-one applications (often called a \keyword{front-end})
that provide a text editor (for \objectref{Step}{itm:step1}),
buttons or menu items to run the \appname{latex} or
\appname{pdflatex} \glslink{terminal}{command-line application} (for
\objectref{Step}{itm:step2}) and, in some cases, a viewer to perform
\objectref{Step}{itm:step3}.

\sectionref[The next section]{sec:texworks} describes one such front-end called
\iappnamelink[.]{TeXWorks}{http://www.tug.org/texworks/} I~have chosen
to describe TeXWorks because it is a free, cross-platform
application. Binaries are available for Microsoft Windows, Mac~OS~X
and GNU/Linux. The screen shots of TeXWorks in this
book were taken from the Linux version running
under Fedora. If you run TeXWorks on other operating systems, it may
have a slightly different look, but it has the same functionality.

New versions of \texdistro{TeX~Live} and \texdistro{MiKTeX} include TeXWorks for MS~Windows,
and new versions of \texdistro{MacTeX} include TeXWorks for Mac~OS~X users.
GNU/Linux users can use their Add\slash Remove Software utility to 
install TeXWorks. Alternatively, you can download TeXWorks by
following the links provided at \url{http://www.tug.org/texworks/}.

\refstepcounter{object}\label{obj:install}%
If you're confused by all the options, let's keep
things as simple as possible:

\begin{itemize}

\item MS~Windows:

 You have a choice between \texdistro{MiKTeX} (or \texdistro{proTeXt}) and 
 \texdistro{TeX~Live}. \texdistro{MiKTeX} provides a
 smaller and quicker installation, but the downside is that you may
 not have the classes or packages you want to use. \texdistro{MiKTeX} can
 install these whenever you try to \LaTeX\ a document that uses them, 
 but you need an Internet connection while it does this. 
 \texdistro{TeX~Live} installs everything, so it takes longer and needs 
 more space, but you should have the majority of packages and classes that you need.

 \begin{description}

   \item[TeX~Live:]\mbox{}

   \begin{enumerate}
    \item Fetch and unpack
     \url{http://mirror.ctan.org/systems/texlive/tlnet/install-tl.zip}

    \item Run \iappname{install-tl} and follow the instructions.
     This can take an hour or more.
   \end{enumerate}

   \item[proTeXt:]\mbox{}

   \begin{enumerate}
    \item Go to \url{http://tug.org/protext/}

    \item Click on the \dq{download the self-extract protext.exe} link
     to download and run the executable.
   \end{enumerate}

   \item[MiKTeX:]\mbox{}

   \begin{enumerate}
    \item Go to \url{http://www.miktex.org/}

    \item In the left-hand panel, there is a link to the download page
     for the latest version. At time of writing, it is \texdistro{MiKTeX}~2.9.
     Click on that link.

    \item Scroll down to the section \dq{Installing a basic
      \texdistro{MiKTeX} system}.

    \item If you're happy with the selected mirror location, click on the 
      \dq{Download} button.

    \item Run the executable.
   \end{enumerate}
 \end{description}

\item Mac~OS~X:

 \begin{enumerate}
  \item Go to \url{http://tug.org/mactex/}

  \item Follow the instructions on that page.
 \end{enumerate}

\item GNU/Linux:

 \begin{enumerate}
  \item Fetch and unpack
    \url{http://mirror.ctan.org/systems/texlive/tlnet/install-tl-unx.tar.gz}

  \item Follow the instructions at
    \url{http://tug.org/texlive/quickinstall.html}

  \item Once \texdistro{TeX~Live} has finished installing, run your
   system's \dq{Add\slash Remove Software} tool.

  \item Find \dq{texworks}, select the newest version and install.
 \end{enumerate}

\end{itemize}

If you run into problems, there are mailing lists at
\url{http://tug.org/mailman/listinfo/tex-live} and
\url{http://docs.miktex.org/manual/lists.html} for
\itexdistro{TeX~Live} and \itexdistro{MiKTeX}, and
\itexdistro{MacTeX} help at \url{http://www.tug.org/mactex/help/}.
There is also a list of places where you can ask for help in
\appendixref{ch:help}.

%%%%%%%%%%%%%%%%%%% TeXWorks %%%%%%%%%%%%%%%%%%%%%

\setnode{texworks}
\section{TeXWorks}
\label{sec:texworks}

Hopefully you've managed to successfully install \gls{tex} and TeXWorks as
described \latexhtml{above}{in the \htmlref{previous
section}{ch:tex2pdf}}, so let's test it out.

First run TeXWorks. On Windows, you can access it via the
\menu{Start} menu. On GNU/Linux, it's probably located in
\menu{Applications}\menuto\menu{Office}, or you can type
\iappname{texworks} in a \gls{terminal}. You should now see the TeXWorks
window. The button marked
with a grey triangle in a green circle is the build or typeset function. It
runs the application in the drop-down list next to it. This is set to
pdfLaTeX, which is what we want for now.

It's a good idea to switch on the syntax highlighting, if it isn't
already on. This is done via the \menu{Format}\menuto\menu{Syntax
Coloring} sub-menu. Make sure the \menu{LaTeX} item is selected.

Next, type in the following sample source code, as shown in \figureref{fig:texworks2}
(the commands used here will be described in more detail in
\chapterref{ch:simpledoc}):
\begin{code}
\begin{verbatim}
\documentclass{scrartcl}

\begin{document}

This is an example document.

\end{document}
\end{verbatim}
\end{code}
\textbf{Pay close attention to the backslashes at the start of each
command name.} If you find the font is a bit too small for you,
you can make it larger using the \menu{Format}\menuto\menu{Font} menu item.
\emph{This doesn't affect the font size in your PDF file, just the
font size of your code.} This displays the \dq{Select Font} dialog
box. Set the font size as appropriate.

Then save the document, using the
\menu{File}\menuto\menu{Save As} menu item. I~called my document
\texttt{example1.tex} (remember the \texttt{.tex} extension and
stick to file names that only consist of alphabetical characters,
digits and hyphens\dash don't uses spaces or underscores).

\begin{figure}[htbp]
\figconts
  {pictures/texworks2}
  {%
    \caption[Source Code for an Example Document]{Source Code for an 
     Example Document. (Syntax highlighting switched on.)}%
  }%
  {fig:texworks2}
\end{figure}

Now that you have saved the file, you can run pdfLaTeX. Make sure
the drop-down list next to the build button has \dq{pdfLaTeX}
selected and click on the build button. If all goes well, a new
window should open displaying the typeset document
(\figureref{fig:texworks3}).

\begin{figure}[htbp]
\figconts
  {pictures/texworks3}
  {%
    \caption{The Typeset Document}
  }%
  {fig:texworks3}
\end{figure}

\begin{figure}[htbp]
\figconts
  {pictures/texworks4}
  {%
    \caption{The Source Code Has a Misspelt Command}
  }%
  {fig:texworks4}
\end{figure}

Now let's see what happens if there is an error in the source code.
In \figureref{fig:texworks4} I~have misspelt the
\glsni{documentclass} command. This time, when I~click on the
build button, I~get the error message:
\begin{verbatim}
! Undefined control sequence.
l.1 \documentclas
                 {scrartcl}
?
\end{verbatim}
(Shown in \figureref{fig:texworks5}.)

\begin{figure}[htbp]
\figconts
  {pictures/texworks5-annote}
  {%
     \caption{An Error Message is Displayed}%
  }%
  {fig:texworks5}
\end{figure}


Here \dq{Undefined control sequence} means an unrecognised command,
and below that message, \dq{\texttt{l.1}} means the error was encountered on
line~1. An input line at the bottom of the window has appeared with a
cursor. \LaTeX\ is in interactive mode
and is awaiting a response. There are several responses, but I'm
only going to mention two of them:
\begin{enumerate}
\item Type \texttt{h} and press the Return/Enter key \enter. This displays
a short help message and awaits a new response (see
\figureref{fig:texworks6}).

\item Type \texttt{x} and press Return/Enter. This aborts the
\LaTeX\ run.
\end{enumerate}

Notice that the green circle button with the grey triangle has turned
into a red stop button. This button can be used to abort the
process instead of typing \texttt{x}.

\begin{figure}[htbp]
\figconts
  {pictures/texworks6}
  {%
    \caption{A Short Help Message}%
  }%
  {fig:texworks6}
\end{figure}

Now, there is a second tab at the bottom of the TeXWork's window
(\figureref{fig:texworks7}). This lists the error message and
provides a link to the line where the error occurred. Clicking on
this link highlights line~1. Next I~need to fix the error by
correcting the spelling of the command. Once it's fixed, I~can click
on the build button.

\begin{figure}[htbp]
\figconts
  {pictures/texworks7}
  {%
    \caption{Error Tab}%
  }%
  {fig:texworks7}
\end{figure}

Here's another error you might encounter: I'm now going to misspell
the class name. It should be \texttt{scrartcl}, but in
\figureref{fig:texworks8} it has been misspelt. This time, when I
click on the build button, I~get the error:
\begin{verbatim}
! LaTeX Error: File `scrartc.cls' not found.
\end{verbatim}
I~have two choices: type in the correct name on the line below
\dq{Enter file name:} or I~can abort the process using the red abort
button. In either case I~need to go back and fix the error in my
code.

\begin{figure}[htbp]
\figconts
  {pictures/texworks8}
  {%
   \caption{Misspelt Class File}%
  }%
  {fig:texworks8}
\end{figure}


%%%%%%%%%%%%% CREATING A SIMPLE DOCUMENT %%%%%%%%%%%%%%%%%%

\setnode{simpledoc}
\chapter{Creating a Simple Document}
\label{ch:simpledoc}

Having installed and tested the software, let's now look at how to
actually write the \gls{source}.  The very first line of any document that you create must have
the \gls{command}:
\begin{definition}
\gls{documentclass}\oarg{\meta{option-list}}\marg{\meta{class-name}}
\end{definition}%
This tells \LaTeX\ what type of document you want to create
(for example an article, a technical report or correspondence).  The
\glsni{documentclass} command takes one \gls{mandatory}, \meta{class-name},
that specifies the \gls{cls}.

There are many class files available, and some publishers,
institutions or journals provide their own custom classes (for
example, the \icls{jmlr} class for the Journal of Machine Learning
Research). Popular classes include \icls{memoir} (for books and
reports) and those supplied in the \koma\ bundle (for books,
reports, articles and correspondence). There's also \icls{beamer}
(for presentations) as well as classes for typesetting exams,
flashcards, concert programmes etc. For simplicity, this
book will concentrate on three of the
\koma\ classes \icls{scrartcl} (for articles), \icls{scrreprt}
(for technical reports, theses etc) and \icls{scrbook} (for books).

We'll start with a very simple document, so let's use the
\icls{scrartcl} class file.  In this case the very first line of the
\gls{source} should be:
\begin{codeS}
\glsni{documentclass}\marg{scrartcl}
\end{codeS}
The \glsni{documentclass} command also takes an 
\gls{optional}, \meta{option-list},
which should be a comma separated list of options to be passed to the
class file.  This allows you to override the class file defaults.
For example, the \icls{scrartcl} class file by default uses A4 
paper, but if you are in the USA you will probably want to use
letter paper.  This can be achieved
using the option \clsopt{letterpaper}.  So you would need to edit the
above line to:
\begin{codeS}
\glsni{documentclass}\oarg{letterpaper}\marg{scrartcl}
\end{codeS}
\refstepcounter{object}\label{obj:sizeopt}%
Let's change another option.  The normal font size is \clsopt{11pt}
by default, but we have the option to change it, so let's use
\clsopt{12pt}:
\begin{codeS}
\glsni{documentclass}\oarg{letterpaper,12pt}\marg{scrartcl}
\end{codeS}
You can also change your document so that it is in a two-column
format using the \clsopt{twocolumn} option:
\begin{codeS}
\glsni{documentclass}\oarg{letterpaper,12pt,twocolumn}\marg{scrartcl}
\end{codeS}

After deciding what type of document you want, you now need to specify
the contents of the document.  This is done inside the \gls{env-document}
\gls{environment}.  The document is started with the
command:
\begin{codeS}
\glsni{begin}\marg{document}
\end{codeS}
and ended with
\begin{codeS}
\glsni{end}\marg{document}
\end{codeS}
(\LaTeX\ stops reading the file when it reaches the above line, so
anything occurring after it is ignored.)

My \gls{source} now looks like:
\begin{code}
\glsni{documentclass}[12pt]\marg{scrartcl}\newline
\mbox{}\newline
\glsni{begin}\marg{document}\newline
\mbox{}\newline
\glsni{end}\marg{document}
\end{code}
Every document you create must have this form.  You can't simply
start typing the document text.  You must first specify
your class file\indexCLS, and then place the contents of the document
inside the \glsni{env-document} environment.\screenpagebreak

So far so good, but at the moment we have an empty document, so we
won't get any output.  Let's now put some text into our document:
\begin{code}
\begin{alltt}
\glsni{documentclass}\oarg{12pt}\marg{scrartcl}

\glsni{begin}\marg{document}

This is a simple document.
Here is the first paragraph.

Here is the second paragraph. As you
can       see it's 
a rather 
short paragraph, but not as short as the previous one.

\glsni{end}\marg{document}
\end{alltt}
\end{code}

\xminisec{Top Five Mistakes Made by New Users}
I first started teaching \LaTeX\ in 1998, and these are the most
common errors I've seen when people start learning \LaTeX:
\begin{enumerate}
\item Missing out the backslash \glsni{backslashchar} at the start
of one or more of the \glspl{command}.

\item Using a forward slash \glsni{slashchar} instead of a backslash
\glsni{backslashchar}.

\item Forgetting \texttt{\glsni{end}\marg{document}}.

\item Misspelling \dq{document} (in
\texttt{\glsni{begin}\marg{document}} and
\texttt{\glsni{end}\marg{document}}).

\item Missing a closing brace \glsni{rightbracechar}.

\end{enumerate}
If you encounter any problems when you start out, go through that check list first.
Then check \appendixref{ch:errors}.

Whenever you start a new document, always type in the
\glsni{documentclass}, \texttt{\glsni{begin}\marg{document}} and
\texttt{\glsni{end}\marg{document}} commands first
(\figureref{fig:newdoc-1}). Then move your cursor between the
\glsni{begin} and \glsni{end} lines and type the document text
(\figureref{fig:newdoc-2}).

\begin{figure}[htbp]
\figconts
  {pictures/newdoc-1}
  {%
    \caption[Starting a New Document (1)]{Starting a New Document: 
    always type these three lines first.}%
  }%
  {fig:newdoc-1}

\ifscreenorother
{%
  \end{figure}
  \begin{figure}[htbp]
}%
{%
  \strut
}

\figconts
  {pictures/newdoc-2}
  {%
    \caption[Starting a New Document (2)]{Starting a New Document: 
     move the cursor inside the \glsni{env-document} environment and 
     start typing the document text.}%
  }%
  {fig:newdoc-2}
\end{figure}

\begin{exercise}{Simple Document}{ex:simpledoc}

Try typing the code in the above example into TeXWorks or the editor
of your choice (see \chapterref[the previous section]{ch:tex2pdf} if
you can't remember what to do.) You can also \download{simpledoc} a
copy of this file, but I~recommend that you try typing it in to give
yourself some practice.

Things to note while you are typing: firstly, when you press the
return character at the end of a line this end-of-line character is
converted into a space in the \gls{output}.  So the
fact that I~have some very ragged lines in my \gls{source} has
no effect on the final result. (Note that some front-ends will 
reformat your lines as you type.) Whereas a completely blank line
will be converted into a paragraph break (\gls{par} has the same
effect).

Secondly, multiple \Index{spaces} are converted into a single space,
so the large gap between the words \dq{\texttt{can}} and \dq{\texttt{see}} is
no different from having a single space.

Once you have typed up your \gls{source}, save your file 
(called, for example, \texttt{exercise1.tex}), and run PDF\LaTeX\ as described in 
\sectionref{sec:texworks}. If all goes well, TeXWorks should display
the resulting PDF file in a new window, usually alongside the window
containing the source code.

\xminisec{Notes:}
\begin{enumerate}

\item Each paragraph automatically starts with an indentation in the
PDF.

\item There is no blank line between the paragraphs in the PDF
document. (See what happens if you add the \koma\ class option
\scrclsopt[full]{parskip}:
\begin{codeS}
\glsni{documentclass}\oarg{12pt,parskip=full}\marg{scrartcl}
\end{codeS}
and rebuild the PDF.)

\item Move the mouse over one of paragraphs in the PDF viewer and
pop-up the context menu (usually a right mouse click). Select
\menu{Jump to Source}. The window containing the source code should
now gain the focus and the line of code matching the typeset line you clicked on
in the PDF should now be highlighted.

\end{enumerate}

\end{exercise}

%%%%%%%%%%%%%%%%%% Simple Commands %%%%%%%%%%%%%%%%%%

\setnode{simplecommands}
\section{Using Simple Commands}
\label{sec:simple}

\faq{Typesetting all those TeX-related logos}{logos}Now let's try
adding a few simple \glspl{command} to our
document. The command \gls{LaTeX} produces the
\makeimg{ALT=LaTeX}{\LaTeX} logo and the command \gls{today} prints
the current date. \LaTeX\ always ignores any spaces that follow a
command name that consists of letters, as it uses the space to
indicate the end of the command name. This means that if we want a
space to occur immediately after the command, we need to explicitly
say so using the command \gls{spacesym} (recall from
\ifscreenorother{\htmlref{earlier}{obj:visiblespace}}{page~\pageref{obj:visiblespace}}
\gls{visiblespace} indicates a space character). So, for example:
\begin{codeS}
\glsni{LaTeX}\glsni{spacesym}logo
\end{codeS}%
produces the output:
\begin{resultS}[LaTeX logo]
\LaTeX\ logo
\end{resultS}%
Some people when starting out can get a bit confused by this and
read it as the entity \dq{\cmdname{LaTeX}\cmdname{}} whereas it is
in fact two commands: \dq{\cmdname{LaTeX}} (print the \LaTeX\ logo)
followed \dq{\cmdname{\ }} (print a space.)

Let's also try using a command that takes an 
\gls{argument}.  The command 
\begin{definition}
\gls{footnote}\refstepcounter{object}\label{cmd:footnote}\marg{\meta{text}}
\end{definition}%
takes one argument that specifies the text that should
appear in the footnote.  This command must be placed where you want
the footnote marker to appear.

\begin{exercise}{Using Simple Commands}{ex:simple}

Edit the document you created in \exerciseref{ex:simpledoc},
so that it looks like the following: (You can \download{simple} it if
you like, but again it is better if you try typing it in yourself.)
\begin{bcode}
\glsni{documentclass}\oarg{12pt}\marg{scrartcl}\newline
\mbox{}\newline
\glsni{begin}\marg{document}\newline
\mbox{}\newline
This is a simple \glsni{LaTeX}\glsni{spacesym}document.
Here is the first paragraph.\newline
\mbox{}\newline
Here is the second paragraph\glsni{footnote}\marg{with a footnote}. 
As you can see it's a rather short paragraph, but not 
as short as the previous one. This document was 
created on: \glsni{today}.\newline
\glsni{end}\marg{document}
\end{bcode}

Now \htmlref{\LaTeX\ your document}{ch:tex2pdf} and view the result.
(Remember to check \appendixref{ch:errors}
if you have a problem.) You should see the \LaTeX\ logo, the footnote
marker and the current date.  If you scroll down to the bottom of the
page, you should see the footnote.
\end{exercise}

%%%%%%%%%%%%%%%%%%%%%%%%%% PACKAGES %%%%%%%%%%%%%%%%%%%%%%%%

\setnode{packages}
\section{Packages}
\label{sec:packages}

Packages\indexSTY\ are files with the extension \texttt{.sty} that
either define new \glspl{command} or redefine
existing commands\faq{What are LaTeX classes and packages?}{clsvpkg}.
They're like a type of add-on or plug-in. Most of the commonly used
packages should have been installed when you installed your \gls{tex}
distribution (see \chapterref{ch:tex2pdf}).
\appendixref{ch:installsty} covers how to install new packages.
Most packages come with documentation that can be accessed using the
\iappname{texdoc} application described in \sectionref{sec:texdoc}.

Packages are loaded in the \gls{preamble} (after
\gls{documentclass} and before \glsni{begin}\marg{document}) using
\begin{definition}
\gls{usepackage}\oarg{\meta{option list}}\marg{\meta{package}}
\end{definition}
where \meta{package} is the name of the package and \meta{option
list} is a list of comma-separated options. For example, to load the
package \isty{graphicx} with the \istyopt{graphicx}{draft} option:
\begin{codeS}
\glsni{usepackage}\oarg{draft}\marg{graphicx}
\end{codeS}
Any applicable class options are also passed to packages, so in
\begin{code}
\gls{documentclass}\oarg{draft}\marg{scrartcl}\newline
\glsni{usepackage}\marg{graphicx}
\end{code}
the \clsopt{draft} option is set for both the \icls{scrartcl} class
and the \isty{graphicx} package.

You can specify more than one package in the argument of
\glsni{usepackage}, where each package name is separated by a
comma. For example:
\begin{codeS}
\glsni{usepackage}\marg{amsmath,amsfonts}
\end{codeS}
The \isty{graphicx} package is covered in \chapterref{ch:graphicx}
and the \isty{amsmath} package is covered in \chapterref{ch:maths},
so let's start out with a~relatively simple example.

\setnode{datetime}
\subsection{Changing the Format of
\latexhtml{\texorpdfstring{\protect\cmdname{today}}{\protect\textbackslash
today}}{\protect\cmdname{today}}}
\label{sec:today}

In the \htmlref{previous exercise}{ex:simple}, we used the
\gls{today} command to produce the current date. By default, this
command displays the date in US format. To illustrate how to use
packages, this section will look at how to use the \isty{datetime}
package to change the way that \gls{today} displays the date.

The \isty{datetime} package has various options that can be used to
change the format of \gls{today}.  For example, by default the
\isty{datetime} package redefines \gls{today} to display the
date in the form: {\showdowtrue\longdate
\renewcommand{\fmtord}[1]{\textsuperscript{#1}}\docdate}. The option
\istyopt{datetime}{short} will produce an abbreviated form, (for
example {\shortdate\docdate}) and the option \istyopt{datetime}{nodayofweek}
won't display the day of the week (for example
{\showdowfalse\docdate}).
For those who don't like the raised ordinal, there is the
\istyopt{datetime}{level} option.
These can be passed as a comma separated list in the
\gls{optional} to the
\gls{usepackage} command.  It is also possible to use a
declaration instead. For example, to redefine \gls{today} to
display the date in the form {\ddmmyyyydate\docdate}, you can either do
\begin{codeS}
\glsni{usepackage}\oarg{ddmmyyyy}\marg{datetime}
\end{codeS}%
or
\begin{code}
\begin{alltt}
\glsni{usepackage}\marg{datetime}
\gls{ddmmyyyydate}
\end{alltt}
\end{code}

The \isty{datetime} package also defines the command
\begin{definition}
\gls{currenttime}
\end{definition}
which displays the current time, where again
the format can be changed by the package options. So the option
\istyopt{datetime}{12hr} will cause \glsni{currenttime} to display
the date in 12 hour format (for example, \ampmtime) and the option
\istyopt{datetime}{24hr} will cause \glsni{currenttime} to display
the date in 24 hour format (for example, \xxivtime).

\begin{exercise}{Using the datetime Package}{ex:datetime}

Edit your document from \exerciseref{ex:simple} so that it uses the \isty{datetime}
package. Experiment with the different package options, for example
\begin{codeS}
\glsni{usepackage}\oarg{short,nodayofweek,level,12hr}\marg{datetime}
\end{codeS}
and add the current time
\begin{codeS}
This document was created on: \glsni{today}\glsni{spacesym}at \glsni{currenttime}.
\end{codeS}
You can \downloadorview{datetime} an example.
For a full list \isty{datetime} of package options, see the
\sty{datetime} \htmlref{documentation}{sec:texdoc}.
\latex{(Refer to \sectionref{sec:texdoc} on how to find package
documentation.)}

\end{exercise}

%%%%%%%%%%%%%%%%%%%%%% Special Characters %%%%%%%%%%%%%%%%%%%%

\setnode{symbols}
\section{Special Characters and Symbols}
\label{sec:chars}

You can use most of the standard characters that you find on your keyboard, but the 
10 symbols shown in \tableref{tab:specialchars} have a special
meaning.

\begin{table}[htbp]
\caption{Special Characters}
\label{tab:specialchars}
\centering
\begin{makeimage}\end{makeimage}
\glsi{leftbracechar} \glsi{rightbracechar} 
\glsi{percentchar}
\glsi{ampchar} 
\glsi{dollarchar} 
\glsi{hashchar} 
\glsi{underscorechar}
\glsi{circumchar} 
\glsi{tildechar} 
\glsi{backslashchar}
\end{table}

We have already used the curly braces \gls{leftbracechar} and
\gls{rightbracechar}. The percent symbol \gls{percentchar} is a
comment character. Everything from the percent symbol up to the end
of line is ignored by \LaTeX. This means you can have comments in
your \gls{source} to remind you what a particular part of your code
is doing. We have also used the backslash symbol
\glsni{backslashchar} which indicates that we are using a \LaTeX\
command, as in \glstext{LaTeX} or \glstext{today}. The meaning of
the other special characters will be covered later.

So what do you do if you want one of these symbols to actually appear
in your document? \faq{Where can I find the symbol for
\ldots?}{symbols} \tableref{tab:symbols} lists commands that produce
these and other symbols\faq{How to get copyright, trademark,
etc}{tradesyms}. Note that some of the commands have short cuts, such
as \glsni{emdash} instead of \glsni{textemdash} and
\glsni{questiondown} instead of \glsni{textquestiondown}.

The symbol \texttt{`} is the backtick (or grave) symbol, as
opposed to the apostrophe symbol~\texttt{'}.  \latexhtml{The backtick
symbol usually looks like~\texttt{\`{}} on a keyboard, and on}{On}
most UK keyboards it is situated to the left of the \texttt{1}~key.
The opening double quote is created using two adjacent backtick
symbols and the closing double quote with two adjacent
apostrophe symbols. This gives 66 and 99 style quotes, which you
wouldn't get using the double quote symbol on your keyboard.

Note that the symbols \gls{barchar} \gls{ltchar} and \gls{gtchar} have to
be created using \glsni{textbar}, \glsni{textless} and
\glsni{textgreater} when in normal text mode. If you try to enter
them using the corresponding keyboard characters you may get
\makeimg{em dash}{\textemdash} \makeimg{upside down exclamation
mark}{\textexclamdown} and \makeimg{upside down question
mark}{\textquestiondown}. (They do however work if you are in
\htmlref{maths mode}{ch:maths}.\footnote{There are also some text
fonts that will display them correctly, but don't rely on it.}) The
slash character \gls{text.slash} may be used directly, as in
\dq{\texttt{and\glsni{text.slash}or}}, but no line break will be permitted at the slash,
whereas \glsni{slash} (as in \dq{\texttt{and\glsni{slash}\glsni{visiblespace}or}}) will allow a
line break at that point.

\begin{table}[htbp]
\caption{Symbols}
\label{tab:symbols}
\centering
\begin{latexonly}
\begin{tabular}{@{}ll@{\hspace{4\tabcolsep}}ll@{\hspace{4\tabcolsep}}ll@{}}
\begin{inlinedef}\gls{textbackslash}\end{inlinedef} & \textbackslash&
\begin{inlinedef}\gls{slash}\end{inlinedef} & \slash &
\begin{inlinedef}\gls{textgreater}\end{inlinedef} & \textgreater \\
\begin{inlinedef}\gls{textasciicircum}\end{inlinedef} & \textasciicircum & 
\begin{inlinedef}\gls{dollar}\end{inlinedef} & \$ &
\begin{inlinedef}\gls{textbar}\end{inlinedef} & \textbar \\
\begin{inlinedef}\gls{textasciitilde}\end{inlinedef} & \textasciitilde&
\begin{inlinedef}\gls{leftbrace}\end{inlinedef} & \{ &
\begin{inlinedef}\gls{textless}\end{inlinedef} & \textless \\
\begin{inlinedef}\gls{pounds}\end{inlinedef} & \pounds &
\begin{inlinedef}\gls{rightbrace}\end{inlinedef} & \} &
\begin{inlinedef}\gls{dag}\end{inlinedef} & \dag \\
\begin{inlinedef}\gls{textregistered}\end{inlinedef}& 
\textregistered&
\begin{inlinedef}\gls{hash}\end{inlinedef} & \# &
\begin{inlinedef}\gls{ddag}\end{inlinedef} & \ddag \\
\begin{inlinedef}\gls{texttrademark}\end{inlinedef} & \texttrademark &
\begin{inlinedef}\gls{percent}\end{inlinedef} & \% &
\begin{inlinedef}\gls{quoteright}\end{inlinedef}
or \begin{inlinedef}\gls{textquoteright}\end{inlinedef} & \textquoteright\\
\begin{inlinedef}\gls{copyright}\end{inlinedef} & \copyright &
\begin{inlinedef}\gls{amp}\end{inlinedef} & \& &
\begin{inlinedef}\gls{quoteleft}\end{inlinedef} 
or \begin{inlinedef}\gls{textquoteleft}\end{inlinedef} & \textquoteleft \\
\begin{inlinedef}\gls{textbullet}\end{inlinedef} & \textbullet & 
\begin{inlinedef}\gls{i}\end{inlinedef} & \i & 
\begin{inlinedef}\gls{quotedblright}\end{inlinedef}
or \begin{inlinedef}\gls{textquotedblright}\end{inlinedef} & \textquotedblright\\
\begin{inlinedef}\gls{questiondown}\end{inlinedef}
or \begin{inlinedef}\gls{textquestiondown}\end{inlinedef} & ?` &
\begin{inlinedef}\gls{j}\end{inlinedef} & \dotlessj & 
\begin{inlinedef}\gls{quotedblleft}\end{inlinedef} 
or \begin{inlinedef}\gls{textquotedblleft}\end{inlinedef} & \textquotedblleft\\
\begin{inlinedef}\gls{exclamdown}\end{inlinedef}
or \begin{inlinedef}\gls{textexclamdown}\end{inlinedef} & !` &
\begin{inlinedef}\gls{dash}\end{inlinedef} & - &
\begin{inlinedef}\gls{endash}\end{inlinedef}
or \begin{inlinedef}\gls{textendash}\end{inlinedef} & \textendash \\
\begin{inlinedef}\gls{emdash}\end{inlinedef}
or \begin{inlinedef}\gls{textemdash}\end{inlinedef} & \textemdash &
\begin{inlinedef}\gls{S}\end{inlinedef} & \S &
\begin{inlinedef}\gls{textperiodcentered}\end{inlinedef} &
\textperiodcentered\\
\begin{inlinedef}\gls{ldots}\end{inlinedef} & \ldots  &
\begin{inlinedef}\gls{P}\end{inlinedef} & \P &
\begin{inlinedef}\gls{underscore}\end{inlinedef} 
or \begin{inlinedef}\gls{textunderscore}\end{inlinedef}& \_ 
\end{tabular}
\end{latexonly}
\begin{htmlonly}
\begin{tabular}{ll|ll}
\begin{inlinedef}\gls{textbackslash}\end{inlinedef} &
\textbackslash&
\begin{inlinedef}\gls{underscore}\end{inlinedef}\space 
or \begin{inlinedef}\gls{textunderscore}\end{inlinedef}& \_ \\
\begin{inlinedef}\gls{dash}\end{inlinedef} & - &
\begin{inlinedef}\gls{endash}\end{inlinedef}\space
or \begin{inlinedef}\gls{textendash}\end{inlinedef} & \makeimg{en dash character}{\textendash} \\
\begin{inlinedef}\gls{emdash}\end{inlinedef}\space  
or \begin{inlinedef}\gls{textemdash}\end{inlinedef} & \makeimg{em dash character}{\textemdash} &
\begin{inlinedef}\gls{P}\end{inlinedef} & \P \\
\begin{inlinedef}\gls{textasciicircum}\end{inlinedef} & \textasciicircum & 
\begin{inlinedef}\gls{dollar}\end{inlinedef} & \$ \\
\begin{inlinedef}\gls{S}\end{inlinedef} & \S &
\begin{inlinedef}\gls{textasciitilde}\end{inlinedef} & \textasciitilde\\
\begin{inlinedef}\gls{leftbrace}\end{inlinedef} & \{ &
\begin{inlinedef}\gls{ldots}\end{inlinedef} & \ldots \\
\begin{inlinedef}\gls{pounds}\end{inlinedef} & \pounds &
\begin{inlinedef}\gls{rightbrace}\end{inlinedef} & \} \\
\begin{inlinedef}\gls{dag}\end{inlinedef} & 
\makeimg{dagger symbol}{\dag} &
\begin{inlinedef}\gls{questiondown}\end{inlinedef}\space  
or \begin{inlinedef}\gls{textquestiondown}\end{inlinedef} & 
\makeimg{upside-down question mark}{?`} \\
\begin{inlinedef}\gls{textregistered}\end{inlinedef}& 
\makeimg{registered symbol}{\textregistered}&
\begin{inlinedef}\gls{hash}\end{inlinedef} & \# \\
\begin{inlinedef}\gls{ddag}\end{inlinedef} & 
\makeimg{double dagger symbol}{\ddag} &
\begin{inlinedef}\gls{exclamdown}\end{inlinedef}\space  
or \begin{inlinedef}\gls{textexclamdown}\end{inlinedef} & 
\makeimg{upside-down exclamation mark}{!`} \\
\begin{inlinedef}\gls{texttrademark}\end{inlinedef} & \texttrademark &
\begin{inlinedef}\gls{percent}\end{inlinedef} & \% \\
\begin{inlinedef}\gls{quoteright}\end{inlinedef}\space  
or \begin{inlinedef}\gls{textquoteright}\end{inlinedef} &
\makeimg{closing single quote}{\textquoteright}&
\begin{inlinedef}\gls{quotedblright}\end{inlinedef}\space  
or \begin{inlinedef}\gls{textquotedblright}\end{inlinedef} &
\makeimg{closing double quote}{\textquotedblright}\\
\begin{inlinedef}\gls{copyright}\end{inlinedef} & \copyright &
\begin{inlinedef}\gls{amp}\end{inlinedef} & \& \\
\begin{inlinedef}\gls{quoteleft}\end{inlinedef}\space  
or \begin{inlinedef}\gls{textquoteleft}\end{inlinedef} &
\makeimg{opening single quote}{\textquoteleft} &
\begin{inlinedef}\gls{quotedblleft}\end{inlinedef}\space  
or \begin{inlinedef}\gls{textquotedblleft}\end{inlinedef} &
\makeimg{opening double quote}{\textquotedblleft}\\
\begin{inlinedef}\gls{textbullet}\end{inlinedef} & \makeimg{bullet point}{\textbullet} & 
\begin{inlinedef}\gls{i}\end{inlinedef} & \makeimg{dotless i}{\i} \\ 
\begin{inlinedef}\gls{j}\end{inlinedef} & \makeimg{dotless j}{\dotlessj} & 
\begin{inlinedef}\gls{textbar}\end{inlinedef} & \textbar\\
\begin{inlinedef}\gls{textperiodcentered}\end{inlinedef} &
\middot &
\begin{inlinedef}\gls{textless}\end{inlinedef} & \textless\\
\begin{inlinedef}\gls{textgreater}\end{inlinedef} & \textgreater&
\begin{inlinedef}\gls{slash}\end{inlinedef} & / \\
\end{tabular}
\end{htmlonly}

\end{table}

Ligatures and special symbols are shown in \tableref{tab:ligatures}.
(Note that, as mentioned in the \htmlref{introduction}{ch:intro}, the
f-ligatures are automatically converted.)  When using a command in
the middle of a word, take care that the command doesn't run into the
rest of the word.  For example, the British spelling of the word
man\oe{}uvre has an
\latexhtml{\oe}{oe}-ligature in the middle of it.
You will get an error if you try:
\begin{alltt}\wrong
man\cmdname{oe}uvre
\end{alltt}
as \LaTeX\ will interpret it as the command
\verb|\oeuvre| which doesn't exist.

There are several ways to code this in \LaTeX:
\begin{enumerate}

\item Place a space after the command:
\begin{codeS}
man\glsni{oe}\glsni{visiblespace}uvre
\end{codeS}

\item Place an empty brace after the command:
\begin{codeS}
man\glsni{oe}\marg{}uvre
\end{codeS}

\item Group the command:
\begin{codeS}
man\marg{\glsni{oe}}uvre
\end{codeS}
(This can adversely affect the kerning so is best avoided.)

\end{enumerate}

\begin{table}[htbp]
\caption[Ligatures and Special Symbols]{Ligatures and Special Symbols (Computer Modern Font)}
\label{tab:ligatures}
\centering
\begin{tabular}{ll@{\hspace{3\tabcolsep}}ll@{\hspace{3\tabcolsep}}ll@{\hspace{3\tabcolsep}}ll}
\gls{AE} & 
  \makeimg{upper case AE ligature}{\fontfamily{cmr}\selectfont\AE} &
\gls{ae} & 
  \makeimg{lower case ae ligature}{\fontfamily{cmr}\selectfont\ae} &
\gls{OE} & 
  \makeimg{upper case OE ligature}{\fontfamily{cmr}\selectfont\OE} &
\gls{oe} & 
  \makeimg{lower case oe ligature}{\fontfamily{cmr}\selectfont\oe}\\
\Indextt{fi} & 
  \makeimg{fi ligature}{\fontfamily{cmr}\selectfont fi} &
\Indextt{ffi} & 
  \makeimg{ffi ligature}{\fontfamily{cmr}\selectfont ffi} &
\Indextt{fl} & 
  \makeimg{fl ligature}{\fontfamily{cmr}\selectfont fl} &
\Indextt{ffl} & 
  \makeimg{ffl ligature}{\fontfamily{cmr}\selectfont ffl}\\
\gls{AA} & \makeimg{upper case A ring}{\fontfamily{cmr}\selectfont\AA} &
\gls{aa} & \makeimg{lower case a ring}{\fontfamily{cmr}\selectfont\aa} &
\gls{L} & \makeimg{upper case L with a stroke}{\fontfamily{cmr}\selectfont\L} &
\gls{l} & \makeimg{lower case l with a stroke}{\fontfamily{cmr}\selectfont\l}\\
\gls{O} & \makeimg{upper case slashed O}{\fontfamily{cmr}\selectfont\O}&
\gls{o} & \makeimg{lower case slashed o}{\fontfamily{cmr}\selectfont\o} &
\gls{SS} & \makeimg{upper case Eszett}{\fontfamily{cmr}\selectfont\SS} &
\gls{ss} & \makeimg{lower case eszett}{\fontfamily{cmr}\selectfont\ss}
\end{tabular}
\end{table}

\refstepcounter{object}\label{obj:accents}%
English speakers are by and large very lackadaisical when it comes
to accents, but accents affect pronunciation, and so are just as
important as the correct spelling. There is a big difference
between putting your knife into someone's p\^at\'e (meat paste), and
putting your knife into someone's pate (head)!

Accented letters are created by specifying which accent you want, and
the letter on which to put the accent.  The accent commands are
listed in \xtableref{tab:accents}, and each command takes one
\gls{mandatory}.  The command indicates
what accent to use, and the argument indicates the letter on which to put
the accent.

You may have noticed in \tableref{tab:symbols} the commands \gls{i}
and \gls{j} which produce a dotless i and j (\i\ and \dotlessj).
With old versions of \LaTeX\ (or \TeX) an accent over a normal
\dq{i} or \dq{j} left the original dot in, which is incorrect, so a
dotless \dq{\i} or \dq{\dotlessj} were required. With modern
distributions, an accented \dq{i} or \dq{j} is correctly rendered.

\bookpagebreak
\xminisec{Example:}
\begin{code}
It's na\glsni{doublequote}ive to think that eating mouldy 
p\glsni{circum}at\glsni{acute}e won't result in food poisoning.
\end{code}%
Result:
\begin{resultS}[naive.html]
It's na\"{\i}ve to think that eating mouldy p\^at\'e won't result in food poisoning.
\end{resultS}%

\begin{table}[htbp]
\caption{Accent Commands}
\label{tab:accents}
\centering
\begin{tabular}{llc@{\hspace{3\tabcolsep}}llc}
\toprule
  & \multicolumn{2}{c@{\hspace{3\tabcolsep}}}{Example} &
  & \multicolumn{2}{c}{Example} \\
 Definition & Input & Output & Definition & Input & Output\\\midrule
\begin{inlinedef}\gls{acute}\marg{\meta{object}}\end{inlinedef} &
\verb|\'{c}| & \makeimg{c with an acute accent}{\'{c}} &
\begin{inlinedef}\gls{macron}\marg{\meta{object}}\end{inlinedef} &
\verb|\={c}| & \makeimg{c with a macron accent}{\={c}} \\
\begin{inlinedef}\gls{grave}\marg{\meta{object}}\end{inlinedef} &
\verb|\`{c}| & \makeimg{c with a grave accent}{\`{c}} &
\begin{inlinedef}\gls{period}\marg{\meta{object}}\end{inlinedef} &
\verb|\.{c}| & \makeimg{c with a dot over it}{\.{c}} \\
\begin{inlinedef}\gls{circum}\marg{\meta{object}}\end{inlinedef} &
\verb|\^{c}| & \makeimg{c with a circumflex}{\^{c}} &
\begin{inlinedef}\gls{tilde}\marg{\meta{object}}\end{inlinedef} &
\verb|\~{c}| & \makeimg{c with a tilde above}{\~{c}} \\
\begin{inlinedef}\gls{doublequote}\marg{\meta{object}}\end{inlinedef} &
\verb|\"{c}| & \makeimg{c with an umlaut}{\"{c}} &
\begin{inlinedef}\gls{v}\marg{\meta{object}}\end{inlinedef} &
\verb|\v{c}| & \makeimg{c with a caron accent}{\v{c}} \\
\begin{inlinedef}\gls{u}\marg{\meta{object}}\end{inlinedef} &
\verb|\u{c}| & \makeimg{c with a breve accent over it}{\u{c}} &
\begin{inlinedef}\gls{H}\marg{\meta{object}}\end{inlinedef} &
\verb|\H{c}| & \makeimg{c with double acute accent}{\H{c}} \\
\begin{inlinedef}\gls{t}\marg{\meta{object}}\end{inlinedef} &
\verb|\t{xy}| & \makeimg{x and y with a tie over them}{\t{xy}} &
\begin{inlinedef}\gls{c}\marg{\meta{object}}\end{inlinedef} &
\verb|\c{c}| & \makeimg{c with a cedilla}{\c{c}}\\
\begin{inlinedef}\gls{d}\marg{\meta{object}}\end{inlinedef} &
\verb|\d{c}| & \makeimg{c with a dot under it}{\d{c}} &
\begin{inlinedef}\gls{b}\marg{\meta{object}}\end{inlinedef} &
\verb|\b{c}| & \makeimg{c with a bar under it}{\b{c}}\\
\begin{inlinedef}\gls{r}\marg{\meta{object}}\end{inlinedef} &
\verb|\r{c}| & \makeimg{c with a ring above}{\r{c}} & & &\\
\bottomrule
\end{tabular}
\end{table}

This book only covers a very small subset of available symbol
commands. If the command you want isn't here, try Scott~Pakin's
comprehensive symbol list~\cite{pakin09}. Another useful resource is
\iappnamelink[.]{detexify}{http://detexify.kirelabs.org/classify.html}

\setnode{inputenc}
\subsection{The \sty{inputenc} Package}
\label{sec:inputenc}

Instead of using the accent and ligature commands described
\htmlref{above}{sec:chars}, you can use the \isty{inputenc} package
and enter the character directly, but you must ensure you match the
encoding with that used by your text editor. For example, this book
uses UTF8 encoding so I~have loaded the \sty{inputenc} package in
the \glsterm{preamble} with the \istyopt{inputenc}{utf8} option:
\begin{codeS}
\glsni{usepackage}\oarg{utf8}\marg{inputenc}
\end{codeS}
Note that it's a good idea to also use the \isty{fontenc} package as
well. For example, if you want to use Type~1 fonts:
\begin{code}
\begin{alltt}
\glsni{usepackage}\oarg{T1}\marg{fontenc}
\glsni{usepackage}\oarg{utf8}\marg{inputenc}
\end{alltt}
\end{code}

Returning to an earlier example, 
I~can directly enter the Unicode character (U+0153) for the
lower case \oe\ ligature:
\bookpagebreak
% actually I can't do it because latex2html won't like it, but lets
% pretend I have
\begin{codeS}
man\oe uvre
\end{codeS}
Note that if you are collaborating on a document and you want to
use this approach, you must ensure that all your co-authors use the
same input encoding. For example, suppose you decide to use
ISO~Latin~1 encoding (\istyopt{inputenc}{latin1} option):
\begin{codeS}
\glsni{usepackage}\oarg{latin1}\marg{inputenc}
\end{codeS}
but your co-author is using a UTF-8 editor and types:
\begin{codeS}
na\"ive
\end{codeS}
where \texttt{\"i} is the Unicode character U+00EF\@. UTF-8 characters
use one to four 8-bit bytes whereas ISO Latin~1 uses an 8-bit
single-byte character set. So the U+00EF binary sequence is
interpreted by ISO Latin~1 encoding as two characters: \~A\ (0xC3)
and \={ } (0xAF)\@.  Therefore the resulting PDF file will end up
containing the rather odd looking:
\begin{flushleft}\wrong
na\~A\={ }ve
\end{flushleft}\html{\par}%
(If you are using TeXWorks, you can set your preferred encoding
using \menu{Edit}\menuto\menu{Preferences} and select the \dq{Editor} tab 
where there is an \dq{Encoding} setting. Make sure this setting
matches the \isty{inputenc} option you use in your document.)

\begin{exercise}{Using Special Characters}{ex:spchar}
Start a new file in TeXworks, and
see if you can write the source code to create the output below.
\doifnotbook{(Ignore any \gls{hyphenation} that may appear below, \LaTeX\ does
that automatically where necessary, see \sectionref{sec:hyphenation}. Likewise, 
ignore where the line breaks occur, except for the paragraph break.) }Choose 
whether you want to use the \isty{inputenc} package or if you
want to use commands such as \gls{c}, but in either case you need to
be careful of the \htmlref{special characters}{tab:specialchars}\doifbook{\ 
listed in \tableref{tab:specialchars}}.

\begin{result}[spcharex.html]\setlength{\parindent}{1.5em}%
Item \#1: Our travel expenditure came to \$2000.00 
\& our equipment expenditure came
to \pounds 100.00 plus VAT @ 17.5\%.

Chlo\"e collected Zo\"e from the cr\`eche. They stopped to
admire the fa\c{c}ade of a new caf\'e and then went to a matin\'ee.
\end{result}

You can \downloadorview{spchar} the source code if you can't work out
how to do it, and remember to check \appendixref{ch:errors} if you have a problem.

\end{exercise}

%%%%%%%%%%%%%%%%%%%%%%%%% Lists %%%%%%%%%%%%%%%%%%%%%%%%%%

\setnode{lists}
\section{Lists}
\label{sec:lists}

Now you've had a go at using some \glspl{command},
let's use some \glspl{environment}\latex{ (recall
\sectionref{sec:environment})}. A good example of environments are the list
making environments. There are three basic list making environments:
\glsni{env-itemize} (for unordered lists),
\glsni{env-enumerate} (for ordered lists) and
\glsni{env-description} (for lists where you want to specify
your own label.)

In each of these environments, each item in the list must be started
with the command: 
\begin{definition}
\gls{item}\oarg{\meta{marker}}
\end{definition}
The optional argument \meta{marker} can be used to override 
the default marker for that particular item. (For example, to replace
the bullet point for an individual item in an \htmlref{unordered
list}{sec:itemize} to make that item stand out from all the other
items.) We will be looking at how to change the default marker
in \sectionref{sec:renewcom}.

\xminisec{Related \gls{ukfaq} topics:}
\begin{itemize}
\item \faqlink{Perhaps a missing \cmdname{item}?}{errmissitem}
\item \faqlink{Fancy enumeration lists}{enumerate}
\item \faqlink{How to adjust list spacing}{complist}
\item \faqlink{Interrupting enumerated lists}{interruptlist}
\item \faqlink{\dq{Too deeply nested}}{toodeep}
\end{itemize}

%---------------------- Itemize --------------------------

\setnode{itemize}
\subsection{Unordered Lists}
\label{sec:itemize}

Unordered lists are created using the \gls{env-itemize} environment.

\xminisec{Example:}
\begin{code}
\glsni{begin}\marg{itemize}\newline
\glsni{item} Animal\newline
\glsni{item} Vegetable\newline
\glsni{item} Mineral\newline
\glsni{end}\marg{itemize}
\end{code}
\begin{result}[itemize.html]
\begin{itemize}
\item Animal
\item Vegetable
\item Mineral
\end{itemize}
\end{result}

\xminisec{Another Example:}

Changing the default markers is covered in \sectionref{sec:renewcom}, 
but it's also possible to override the default marker for a
particular item, as in this example (recall the double-dagger
symbol command \gls{ddag} from \tableref{tab:symbols}):
\begin{code}
\glsni{begin}\marg{itemize}\newline
\glsni{item} Animal\newline
\glsni{item}\oarg{\glsni{ddag}} Vegetable\newline
\glsni{item} Mineral\newline
\glsni{end}\marg{itemize}
\end{code}%
\begin{result}[itemizedag.html]
\begin{itemize}
\item Animal
\item[\ddag] Vegetable
\item Mineral
\end{itemize}
\end{result}

Be careful about using square brackets
\gls{text.opensq}\gls{text.closesq} inside an optional
argument. Grouping is required, as in:
\begin{code}
\glsni{begin}\marg{itemize}\newline
\glsni{item} Animal\newline
\glsni{item}\oarg{\marg{\gls{text.opensq}X\gls{text.closesq}}} Vegetable\newline
\glsni{item} Mineral\newline
\glsni{end}\marg{itemize}
\end{code}%
\begin{result}[itemizesq.html]
\begin{itemize}
\item Animal
\item[{[X]}] Vegetable
\item Mineral
\end{itemize}
\end{result}
Similarly if the item starts with an open square bracket
\glsni{opensq}, as in:
\begin{code}
\glsni{begin}\marg{itemize}\newline
\glsni{item} Animal\newline
\glsni{item} \marg{\gls{text.opensq}sic\gls{text.closesq}} Wegetable\newline
\glsni{item} Mineral\newline
\glsni{end}\marg{itemize}
\end{code}
\begin{result}[itemizetextsqbr.html]
\begin{itemize}
\item Animal
\item {[sic]} Wegetable
\item Mineral
\end{itemize}
\end{result}

\setnode{nesteditemize}
\subsubsection{Nested Lists}
\label{sec:nesteditemise}

It is also possible to nest \glsni{env-itemize} environments.
The following example has three levels, each using its own default marker.

\begin{code}
\glsni{begin}\marg{itemize}\newline
\glsni{item} Animal\newline
\glsni{begin}\marg{itemize}\newline
\glsni{item} Mammals\newline
\glsni{item} Birds\newline
\glsni{item} Reptiles. For example:\newline
\glsni{begin}\marg{itemize}\newline
\glsni{item} dinosaurs\newline
\glsni{item} crocodiles\newline
\glsni{end}\marg{itemize}\newline
\glsni{end}\marg{itemize}\newline
\glsni{item} Vegetable\newline
\glsni{begin}\marg{itemize}\newline
\glsni{item} Cultivated\newline
\glsni{item} Wild\newline
\glsni{end}\marg{itemize}\newline
\glsni{item} Mineral\newline
\glsni{end}\marg{itemize}
\end{code}
\begin{result}[nesteditemize.html]
\begin{itemize}
\item Animal
\begin{itemize}
\item Mammals
\item Birds
\item Reptiles. For example:
\begin{itemize}
\item dinosaurs
\item crocodiles
\end{itemize}
\end{itemize}
\item Vegetable
\begin{itemize}
\item Cultivated
\item Wild
\end{itemize}
\item Mineral
\end{itemize}
\end{result}

You might have noticed the code in the above example is a little
difficult to read. Each new list item starts a new paragraph, so it
doesn't matter if we have blank lines before each item. Also, recall
from \chapterref{ch:def} that spaces at the start of each line of
code are ignored, so it's possible to make the code more readable
without affecting the final result:

\begin{code}
\begin{alltt}
\glsni{begin}\marg{itemize}

   \glsni{item} Animal

   \glsni{begin}\marg{itemize}

      \glsni{item} Mammals

      \glsni{item} Birds

      \glsni{item} Reptiles. For example:

      \glsni{begin}\marg{itemize}

         \glsni{item} dinosaurs

         \glsni{item} crocodiles

       \glsni{end}\marg{itemize}

   \glsni{end}\marg{itemize}

   \glsni{item} Vegetable

   \glsni{begin}\marg{itemize}

      \glsni{item} Cultivated

      \glsni{item} Wild

   \glsni{end}\marg{itemize}

   \glsni{item} Mineral

\glsni{end}\marg{itemize}
\end{alltt}
\end{code}

It's now a little easier to see which \verb|\begin{itemize}| matches
up with the corresponding \verb|\end{itemize}|.

\xminisec{Example (Four Levels)}

This example has four levels, which is the maximum allowed by most
classes.

\begin{code}
\begin{alltt}
\glsni{begin}\marg{itemize}
 \glsni{item} Animal

 \glsni{begin}\marg{itemize}
   \glsni{item} Mammal

   \glsni{begin}\marg{itemize}
    \glsni{item} Placental

    \glsni{item} Monotreme

    \glsni{begin}\marg{itemize}
      \glsni{item} Platypus
    \glsni{end}\marg{itemize}

    \glsni{item} Marsupial

    \glsni{begin}\marg{itemize}
     \glsni{item} Kangaroo

     \glsni{item} Koala
    \glsni{end}\marg{itemize}
   \glsni{end}\marg{itemize}

   \glsni{item} Reptile
 \glsni{end}\marg{itemize}

 \glsni{item} Vegetable

 \glsni{item} Mineral
\glsni{end}\marg{itemize}
\end{alltt}
\end{code}
\begin{result}[fourlevelitemize.html]
\begin{itemize}
 \item Animal
 \begin{itemize}
   \item Mammal

   \begin{itemize}
    \item Placental

    \item Monotreme
    \begin{itemize}
      \item Platypus
    \end{itemize}


    \item Marsupial

    \begin{itemize}
     \item Kangaroo

     \item Koala
    \end{itemize}
   \end{itemize}

   \item Reptile
 \end{itemize}

 \item Vegetable

 \item Mineral
\end{itemize}
\end{result}

If you try adding a further level, \LaTeX\ will give a \dq{Too
deeply nested} error.

%---------------- Enumerate -------------------------

\setnode{enumerate}
\subsection{Ordered Lists}
\label{sec:enumerate}

Ordered lists are created using the \gls{env-enumerate} 
environment.
It has exactly the same format as the \gls{env-itemize} environment
described \htmlref{above}{sec:itemize}.

We can use the same example as before, only this time use
\glsni{env-enumerate} instead of \gls{env-itemize}.

\begin{code}
\begin{alltt}
\glsni{begin}\marg{enumerate}
  \glsni{item} Animal
  \glsni{item} Vegetable
  \glsni{item} Mineral
\glsni{end}\marg{enumerate}
\end{alltt}
\end{code}%
The above input will produce the following output:
\begin{result}[enumerate.html]
\begin{enumerate}
\item Animal
\item Vegetable
\item Mineral
\end{enumerate}
\end{result}

As before, the marker for a particular item can be overridden:
\begin{code}
\glsni{begin}\marg{enumerate}\newline
\strut~~\glsni{item} Animal\newline
\strut~~\glsni{item}\oarg{\marg{\gls{text.opensq}X\gls{text.closesq}}} Vegetable\newline
\strut~~\glsni{item} Mineral\newline
\glsni{end}\marg{enumerate}
\end{code}%
\begin{result}[enumeratesqbr.html]
\begin{enumerate}
\item Animal
\item[{[X]}] Vegetable
\item Mineral
\end{enumerate}
\end{result}

\xminisec{Example (Nested):}

As with the \gls{env-itemize} environment, most classes allow a
maximum of four nested \gls{env-enumerate} environments.

\begin{code}
\begin{alltt}
\glsni{begin}\marg{enumerate}
 \glsni{item} Animal

 \glsni{begin}\marg{enumerate}
   \glsni{item} Mammal

   \glsni{begin}\marg{enumerate}
    \glsni{item} Placental

    \glsni{item} Monotreme

    \glsni{begin}\marg{enumerate}
      \glsni{item} Platypus
    \glsni{end}\marg{enumerate}

    \glsni{item} Marsupial

    \glsni{begin}\marg{enumerate}
     \glsni{item} Kangaroo

     \glsni{item} Koala
    \glsni{end}\marg{enumerate}
   \glsni{end}\marg{enumerate}

   \glsni{item} Reptile
 \glsni{end}\marg{enumerate}

 \glsni{item} Vegetable

 \glsni{item} Mineral
\glsni{end}\marg{enumerate}
\end{alltt}
\end{code}
\begin{result}[nestedenumerate.html]
\begin{enumerate}
 \item Animal
 \begin{enumerate}
   \item Mammal

   \begin{enumerate}
    \item Placental

    \item Monotreme
    \begin{enumerate}
      \item Platypus
    \end{enumerate}


    \item Marsupial

    \begin{enumerate}
     \item Kangaroo

     \item Koala
    \end{enumerate}
   \end{enumerate}

   \item Reptile
 \end{enumerate}

 \item Vegetable

 \item Mineral
\end{enumerate}
\end{result}

%--------------------- Description -------------------------

\setnode{description}
\subsection{Description Environment}
\label{sec:description}

The \gls{env-description} environment has the same format as
the \gls{env-itemize} environment described in
\sectionref{sec:itemize}, only this time you need to
specify a marker as an \gls{optional} to
the \gls{item} command, since there is no default marker
for this environment. The marker may be a textual label, and most
classes will typeset it in bold. The \koma\ classes, such as
\icls{scrartcl}, default to a bold sans-serif font, as illustrated
in this next example:

\begin{code}
\begin{alltt}
\glsni{begin}\marg{description}
  \glsni{item}\oarg{Animal} Living being

  \glsni{item}\oarg{Vegetable} Plant
 
  \glsni{item}\oarg{Mineral} Natural inorganic substance
\glsni{end}\marg{description}
\end{alltt}
\end{code}
\begin{result}[description.html]
\dodescriptionexample
\end{result}

The \koma\ classes provide a way of changing the font style in
the \glsni{env-description} label markers. (The font changing commands
\glsni{normalfont} and \glsni{scshape} will be covered in 
\sectionref{sec:fonts}, and the \koma\ command
\glsni{addtokomafont} in \sectionref{sec:section}.)

\begin{code}
\begin{alltt}
\glsni{addtokomafont}\marg{descriptionlabel}\marg{\glsni{normalfont}\glsni{scshape}}
\glsni{begin}\marg{description}
  \glsni{item}\oarg{Animal} Living being

  \glsni{item}\oarg{Vegetable} Plant

  \glsni{item}\oarg{Mineral} Natural inorganic substance
\glsni{end}\marg{description}
\end{alltt}
\end{code}
\begin{result}[descriptionkoma.html]
% Make sure LaTeX2HTML doesn't think this should be the same
% image as the previous example
\dodescriptionkomaexample
\end{result}

It is possible to nest all the listing environments, as long as you
don't exceed four \gls{env-itemize} and four \gls{env-enumerate}
environments. The \gls{env-description} environment has no
restriction on the number of times it can be nested. However, just
because you can do something, doesn't mean you should. In general
it's best to avoid an excessively complicated block of text in your
document.

\xminisec{Example (Assorted Nesting):}

This example uses each of the listing environments described above.

\bookpagebreak
\begin{code}
\begin{alltt}
\glsni{begin}\marg{description}

   \glsni{item}\oarg{Animal} Living being

   \glsni{begin}\marg{itemize}

      \glsni{item} Mammals

      \glsni{item} Birds

      \glsni{item} Reptiles. For example:

      \glsni{begin}\marg{enumerate}

         \glsni{item} dinosaurs

         \glsni{item} crocodiles

       \glsni{end}\marg{enumerate}

   \glsni{end}\marg{itemize}

   \glsni{item}\oarg{Vegetable} Plant

   \glsni{begin}\marg{itemize}

      \glsni{item} Cultivated. For example:

      \glsni{begin}\marg{enumerate}

         \glsni{item} Carrots

         \glsni{item} Broccoli

         \glsni{item} Potatoes

      \glsni{end}\marg{enumerate}

      \glsni{item} Wild

   \glsni{end}\marg{itemize}

   \glsni{item}\oarg{Mineral} Natural inorganic substance

\glsni{end}\marg{description}
\end{alltt}
\end{code}

\begin{result}[nestedlists.html]
\begin{description}

   \item[Animal] Living being

   \begin{itemize}

      \item Mammals

      \item Birds

      \item Reptiles. For example:
      \begin{enumerate}

         \item dinosaurs

         \item crocodiles

       \end{enumerate}

   \end{itemize}

   \item[Vegetable] Plant

   \begin{itemize}

      \item Cultivated. For example:

      \begin{enumerate}

         \item Carrots

         \item Broccoli

         \item Potatoes

      \end{enumerate}


      \item Wild

   \end{itemize}

   \item[Mineral] Natural inorganic substance

\end{description}
\end{result}

\begin{exercise}{Lists}{ex:lists}

Try writing the \gls{source} that will create
the output shown below.

\begin{result}[listex.html]
\begin{description}
\item[Village] A small collection of dwelling places.
Examples:
\begin{enumerate}
\item Marlingford
\item Saxlingham Nethergate
\end{enumerate}

\item[Town] A large collection of dwelling places.
Examples:
\begin{enumerate}
\item Great Yarmouth
\item Beccles
\end{enumerate}

\item[City] A large town, usually containing a cathedral.
Examples:
\begin{enumerate}
\item Norwich
\item Birmingham
\item London
\end{enumerate}
\end{description}
\end{result}\bookpagebreak
You can \downloadorview{lists} the answer if you can't work out how to do it.

\end{exercise}

%%%%%%%%%%%%%%%%%% Font Changing Commands %%%%%%%%%%%%%%%%%%%%%%

\setnode{fontcommands}
\section{Fonts}
\label{sec:fonts}

\LaTeX\ uses Donald~Knuth's Computer Modern fonts by
default\faq{Using PostScript fonts with TeX}{usepsfont}. This
supplies three font families: serif, sans-serif and a typewriter (or
monospaced) font (as well as the maths
fonts which are discussed in
\sectionref{sec:mathsfonts}). With each font family, you can
change the shape and weight, as well as the size.

\reportpagebreak
\setnode{fontstyle}
\subsection{Changing the Font Style}
\label{sec:fontstyle}

There are two basic ways of changing fonts: you can either change
the font for a small selection of text, for example, if you want to
\emph{emphasize} a word, or you may wish to change the font \dq{from
this point onwards}. The \glspl{command} shown in
\tableref{tab:fontsI} are of the first type (text-block
commands\index{command!text-block|hyperpage}), whereas those shown
in \tableref{tab:fontsII} are of the second type\dash a
\gls{declaration} (or modal command\index{command!modal|hyperpage}).

\xminisec{Note:}

\warning Don't be tempted to use \cmdname{bf}, \cmdname{md}, \cmdname{it}, 
\cmdname{sl}, \cmdname{sc}, \cmdname{sf}, \cmdname{tt} or \cmdname{rm}. 
These commands are \textbf{obsolete}~\cite{l2tabu}.\faq{What's wrong with \cmdname{bf}, \cmdname{it}
etc.?}{2letterfontcmd}

If you use an italic or slanted font declaration, such as
\glsni{itshape}, you will need to add an \keyword{italic correction}
\gls{itcorr} at the end of the block of text, when the last letter
of the sloping text leans too far over. This isn't necessary for
text-block commands, such as \glsni{textit}, just for the modal
commands. The effect is more noticeable when part of a word is
stressed, particularly with certain fonts.

\xminisec{Example:}

In the code below, the first instance of \dq{repeated} doesn't have an italic
correction but the second does: 
\begin{codeS}
\marg{\glsni{itshape} repeated}ly \marg{\glsni{itshape}
repeated\glsni{itcorr}}ly
\end{codeS}%
Using Computer Modern:
\begin{resultS}[repeatedly repeatedly]\fontfamily{cmr}\selectfont
{\itshape repeated}ly {\itshape repeated\/}ly
\end{resultS}
Using Helvetica:
\begin{resultS}[repeatedly repeatedly]\fontfamily{phv}\selectfont
{\itshape repeated}ly {\itshape repeated\/}ly
\end{resultS}
Using Antykwa Toru\'nska typeface:
\begin{resultS}[repeatedly repeatedly]
{\itshape repeated}ly {\itshape repeated\/}ly
\end{resultS}

\begin{table}[tp]
\caption{Font Changing Text-Block Commands}
\label{tab:fontsI}
\centering
\begin{tabular}{@{}lll@{}}
\toprule
\bfseries Command & \bfseries Example Input & \bfseries Corresponding output\\
& & (Computer Modern)\\
\midrule
\begin{inlinedef}\gls{textrm}\marg{\meta{text}}\end{inlinedef} &
\verb|\textrm{roman} text| & \makeimg{Roman text}{\fontfamily{cmr}\selectfont{\fontfamily{cmr}\selectfont roman} text}\\
\begin{inlinedef}\gls{textsf}\marg{\meta{text}}\end{inlinedef} &
\verb|\textsf{sans serif} text| & \makeimg{sans serif text}{\fontfamily{cmr}\selectfont{\fontfamily{cmss}\selectfont sans serif} text}\\
\begin{inlinedef}\gls{texttt}\marg{\meta{text}}\end{inlinedef} &
\verb|\texttt{typewriter} text| & \makeimg{typewriter text}{\fontfamily{cmr}\selectfont{\fontfamily{cmtt}\selectfont typewriter} text}\\[5pt]
\begin{inlinedef}\gls{textmd}\marg{\meta{text}}\end{inlinedef} &
\verb|\textmd{medium} text| & \makeimg{medium text}{\fontfamily{cmr}\selectfont\textmd{medium} text}\\
\begin{inlinedef}\gls{textbf}\marg{\meta{text}}\end{inlinedef} &
\verb|\textbf{bold} text| & \makeimg{bold text}{\fontfamily{cmr}\selectfont\textbf{bold} text}\\[5pt]
\begin{inlinedef}\gls{textup}\marg{\meta{text}}\end{inlinedef} &
\verb|\textup{upright} text| & \makeimg{upright text}{\fontfamily{cmr}\selectfont\textup{upright} text}\\
\begin{inlinedef}\gls{textit}\marg{\meta{text}}\end{inlinedef} &
\verb|\textit{italic} text| & \makeimg{italic text}{\fontfamily{cmr}\selectfont\textit{italic} text}\\
\begin{inlinedef}\gls{textsl}\marg{\meta{text}}\end{inlinedef} &
\verb|\textsl{slanted} text| & \makeimg{slanted text}{\fontfamily{cmr}\selectfont\textsl{slanted} text}\\[5pt]
\begin{inlinedef}\gls{textsc}\marg{\meta{text}}\end{inlinedef} &
\verb|\textsc{Small Caps} text| &
\makeimg{Small Caps text}{\fontfamily{cmr}\selectfont\textsc{Small Caps} text}\\
\begin{inlinedef}\gls{emph}\marg{\meta{text}}\end{inlinedef} &
\verb|\emph{emphasized} text| &
\makeimg{emphasized text}{\fontfamily{cmr}\selectfont\emph{emphasized} text}\\[5pt]
\begin{inlinedef}\gls{textnormal}\marg{\meta{text}}\end{inlinedef} &
\verb|\textnormal{default} text| &
\makeimg{default text}{\fontfamily{cmr}\selectfont default text}
\\\bottomrule
\end{tabular}
\end{table}

Note that if you want to typeset an URL, rather than using
\glsni{texttt} it is better to use\bookpagebreak
\begin{definition}
\gls{url}\marg{\meta{address}}
\end{definition}
which is defined in the \isty{url} package. For example:
\begin{codeS}%
\begin{alltt}
\glsni{url}\marg{http://theoval.cmp.uea.ac.uk/\glsnl{tildechar}nlct/}
\end{alltt}
\end{codeS}
produces:
\begin{resultS}[http://theoval.cmp.uea.ac.uk/\tildesym nlct/]
% simulate to prevent link, which might confuse the reader into
% thinking links automatically occur.
\texttt{http://theoval.cmp.uea.ac.uk/\tildesym nlct/}
\end{resultS}
(Note there is no need to do anything with the \gls{tildechar}
(tilde) \htmlref{special character}{sec:chars} if you use it in the argument of \glsni{url}.)

\begin{table}[hbtp]
\caption{Font Changing Declarations}
\label{tab:fontsII}
\centering
\begin{tabular}{@{}lll@{}}
\toprule
\bfseries Declaration & \bfseries Example Input & \bfseries Corresponding output\\
& & (Computer Modern)\\
\midrule
\begin{inlinedef}\gls{rmfamily}\end{inlinedef} &
\verb|\rmfamily roman text| &
\makeimg{Roman text}{\fontfamily{cmr}\selectfont roman text}\\
\begin{inlinedef}\gls{sffamily}\end{inlinedef} &
\verb|\sffamily sans serif text| &
\makeimg{sans serif text}{\fontfamily{cmss}\selectfont sans serif text}\\
\begin{inlinedef}\gls{ttfamily}\end{inlinedef} &
\verb|\ttfamily typewriter text| &
\makeimg{typewriter text}{\fontfamily{cmtt}\selectfont typewriter text}\\[5pt]
\begin{inlinedef}\gls{mdseries}\end{inlinedef} &
\verb|\mdseries medium text| &
\makeimg{medium text}{\fontfamily{cmr}\selectfont\mdseries medium text}\\
\begin{inlinedef}\gls{bfseries}\end{inlinedef} &
\verb|\bfseries bold text| &
\makeimg{bold text}{\fontfamily{cmr}\selectfont\bfseries bold text}\\[5pt]
\begin{inlinedef}\gls{upshape}\end{inlinedef} &
\verb|\upshape upright text| &
\makeimg{upright text}{\fontfamily{cmr}\selectfont\upshape upright text}\\
\begin{inlinedef}\gls{itshape}\end{inlinedef} &
\verb|\itshape italic text| &
\makeimg{italic text}{\fontfamily{cmr}\selectfont\itshape italic text}\\
\begin{inlinedef}\gls{slshape}\end{inlinedef} &
\verb|\slshape slanted text| &
\makeimg{slanted text}{\fontfamily{cmr}\selectfont\slshape slanted text}\\
\begin{inlinedef}\gls{scshape}\end{inlinedef} &
\verb|\scshape Small Caps text| &
\makeimg{Small Caps text}{\fontfamily{cmr}\selectfont\scshape Small Caps text}\\[5pt]
\begin{inlinedef}\gls{em}\end{inlinedef} &
\verb|\em emphasized text| & 
\makeimg{emphasized text}{\fontfamily{cmr}\selectfont\em emphasized text}\\[5pt]
\begin{inlinedef}\gls{normalfont}\end{inlinedef} &
\verb|\normalfont default text| &
\makeimg{default text}{\fontfamily{cmr}\selectfont default text}
\\\bottomrule
\end{tabular}

\end{table}

\refstepcounter{object}\label{obj:fontenv}\Glspl{environment} can be used instead.
Each \gls{environment} has the same name as its corresponding
declaration, but \emph{without} the preceding backslash. For example:
\begin{codeS}
\begin{alltt}
\glsni{begin}\marg{sffamily}Some sans-serif text.\glsni{end}\marg{sffamily}
\end{alltt}
\end{codeS}%
yields:
\begin{resultS}[Some sans-serif text.]
\begin{sffamily}Some sans-serif text.\end{sffamily}
\end{resultS}

You can combine a font family with a given shape and weight using
a variety of methods.

\xminisec{Examples:}

\begin{enumerate}
\item Localised declarations:

\begin{codeS}
\marg{\glsni{sffamily}\glsni{slshape} Some slanted sans-serif text.}
\end{codeS}

\item Declarations that later get explicitly reset:

\begin{codeS}
\glsni{sffamily}\glsni{slshape} Some slanted sans-serif
text.\glsni{normalfont}
\end{codeS}

\bookpagebreak
\item Mixing text-block and modal commands:

\begin{codeS}
\glsni{textsf}\marg{\glsni{slshape} Some slanted sans-serif text.}
\end{codeS}

\item Nested commands:

\begin{codeS}
\glsni{textsf}\marg{\glsni{textsl}\marg{Some slanted sans-serif text.}}
\end{codeS}

\item Mixing environments and declarations:

\begin{code}
\glsni{begin}\marg{sffamily}\glsni{slshape} Some slanted sans-serif 
text.\glsni{end}\marg{sffamily}
\end{code}
\end{enumerate}

All of the above produce the same output:
\begin{resultS}[Some slanted sans-serif text.]
\textsf{\slshape Some slanted sans-serif text.}
\end{resultS}

\faq{Warning: \dq{Font shape \ldots\ not
available}}{fontunavail}Note that some combinations are not
available, in which case \LaTeX\ will give a warning message, and
will substitute the font for what it considers to be the closest
available match.

\xminisec{Example:}
\begin{codeS}
\glsni{textsc}\marg{\glsni{bfseries} Text}
\end{codeS}
With the Antykwa Toru\'nska typeface, this appears as:
\begin{resultS}[Text (in bold small caps)]
\textsc{\bfseries Text}
\end{resultS}\bookpagebreak\noindent
whereas with Computer Modern, the result is:
\begin{resultS}[Text (in bold)]\fontfamily{cmr}\selectfont
\textsc{\bfseries Text}
\end{resultS}
This is because Computer Modern doesn't have a bold small-caps
font, so it just uses bold. \LaTeX\ gives the following warning:
\begin{verbatim}
LaTeX Font Warning: Font shape `T1/cmr/b/sc' undefined
(Font)              using `T1/cmr/b/n' instead on input line 2792.
\end{verbatim}
Most sans-serif fonts don't provide a small-caps variant, so 
\begin{codeS}
\glsni{textsf}\marg{\glsni{scshape} Text}
\end{codeS}
will either appear in regular sans-serif or small-caps serif,
depending on the font in use. Using Libris sans-serif the result is:
\begin{resultS}[Text (in sans-serif)]
\textsf{\scshape Text}
\end{resultS}
whereas using Computer Modern Sans, the result is:
\begin{resultS}[Text (in small caps)]
\fontfamily{cmss}
\scshape Text
\end{resultS}

\setnode{emph}
\subsubsection{Emphasizing Words or Phrases}
\label{sec:emph}

The \gls{command} \gls{emph},
the \gls{declaration} \gls{em} and the
\gls{environment} \gls{env-em} behave slightly
differently to the corresponding \gls{textit} command,
\gls{itshape} declaration and \gls{env-itshape} environment.
The latter group simply use an italic font, whereas the former will toggle
between sloping and upright.  So if the surrounding font is upright
then \gls{emph}, \gls{em} and \glsni{env-em} will use the
sloping font, but if the surrounding font is italic or slanted,
\gls{emph}, \gls{em} and \glsni{env-em} will use an upright
font.  This is particularly useful in abstracts where the
\Index{abstract} font varies between \glspl{cls}.  It is recommended that if your intention
is to emphasize something, you should use \gls{emph} etc.\ rather
than \gls{textit} etc.\screenpagebreak

\xminisec{Examples:}
\begin{enumerate}
\item Emphasized text in upright surrounding:
\begin{codeS}
Some \glsni{emph}\marg{emphasized} text.
\end{codeS}%
yields
\begin{resultS}[Some emphasized text.]
Some \emph{emphasized} text.
\end{resultS}

\item Emphasized text in italic surrounding:
\begin{codeS}
\marg{\glsni{itshape} Some \glsni{emph}\marg{emphasized} text.}
\end{codeS}%
yields
\begin{resultS}[Some emphasized text]
{\itshape Some \emph{emphasized} text.}
\end{resultS}

\bookpagebreak
\item Emphasized text in upright sans-serif surrounding:
\begin{codeS}
\marg{\glsni{sffamily} Some \glsni{emph}\marg{emphasized} text.}
\end{codeS}%
yields
\begin{resultS}[Some emphasized text]
{\sffamily Some \emph{emphasized} text.}
\end{resultS}
\end{enumerate}

\setnode{fontsize}
\subsection{Changing the Font Size}
\label{sec:fontsize}

When you start writing a document, you need to decide what the base
font size should be. The \koma\ classes default to 11pt, but this can
be changed using the class options \scrclsopt{8pt}, \scrclsopt{9pt},
\clsopt{10pt}, \clsopt{12pt}, \scrclsopt{14pt}, \scrclsopt{17pt} or
\scrclsopt{20pt}.

You can then change the font size \emph{relative to} the base size,
using one of the declarations shown in \tableref{tab:fontsize}. That
way, if you later decide to change the normal font size from, say,
11pt to 12pt, all you need do is change the class option (see
\latexhtml{page~\pageref{obj:sizeopt}}{\htmlref{earlier}{obj:sizeopt}})
and re-run \LaTeX. Note that there are no equivalent text-block
commands.

\begin{table}[tbhp]
\caption{Font Size Changing Declarations}
\label{tab:fontsize}
\centering
\begin{tabular}{@{}lll@{}}
\toprule
\bfseries Declaration & \bfseries Example Input & 
\bfseries Corresponding Output\\
\midrule
\begin{inlinedef}\gls{tiny}\end{inlinedef} &
\verb|\tiny tiny text| & \makeimg{tiny text}{\tiny tiny text}\\
\begin{inlinedef}\gls{scriptsize}\end{inlinedef} &
\verb|\scriptsize script size| &
\makeimg{script size}{\scriptsize script size}\\
\begin{inlinedef}\gls{footnotesize}\end{inlinedef} &
\verb|\footnotesize footnote size| &
\makeimg{footnote size}{\footnotesize footnote size}\\
\begin{inlinedef}\gls{small}\end{inlinedef} &
\verb|\small small text| &
\makeimg{small text}{\small small text}\\
\begin{inlinedef}\gls{normalsize}\end{inlinedef} &
\verb|\normalsize normal size| &
\makeimg{normal size}{\normalsize normal size}\\
\begin{inlinedef}\gls{large}\end{inlinedef} &
\verb|\large large text| &
\makeimg{large text}{\large large text}\\
\begin{inlinedef}\gls{Large}\end{inlinedef} &
\verb|\Large even larger| &
\makeimg{even larger}{\Large even larger}\\
\begin{inlinedef}\gls{LARGE}\end{inlinedef} &
\verb|\LARGE larger still| &
\makeimg{larger still}{\LARGE larger still}\\
\begin{inlinedef}\gls{huge}\end{inlinedef} &
\verb|\huge huge| & \makeimg{huge}{\huge huge}\\
\begin{inlinedef}\gls{Huge}\end{inlinedef} &
\verb|\Huge extra huge| &
\makeimg{extra huge}{\Huge extra huge}
\\\bottomrule
\end{tabular}

\end{table}

Again, \glspl{environment} can be used instead, where
each \gls{environment} has the same name as its corresponding
declaration, but \emph{without} the preceding backslash.
Font environments may be nested, for example:%
\indexEnv{itshape}\indexEnv{Large}
\begin{code}
\glsni{begin}\marg{itshape} Some italic text.
\glsni{begin}\marg{Large}This text is large.\glsni{end}\marg{Large}
\glsni{end}\marg{itshape} Back to normal.
\end{code}%
Output:
\begin{resultS}[Some italic text. This text is large. Back to
normal.]
\begin{itshape} Some italic text.
\begin{Large}This text is large.\end{Large}
\end{itshape}
Back to normal.
\end{resultS}

\setnode{documentfonts}
\subsection{Changing Document Fonts}
\label{sec:changingfonts}

\faq{Choice of scalable outline fonts}{psfchoice}What if you don't
want to use the default Computer Modern fonts? Some publishers and
institutions insist on a combination of Times Roman (serif),
Helvetica (sans-serif) and Courier (typewriter). To do this, you can
load the following packages:
\begin{description}
\item[\isty{mathptmx}] (Times) Only affects \gls{rmfamily} and
\gls{textrm}.

\item[\isty{helvet}] (Helvetica) Only affects \gls{sffamily} and
\gls{textsf}.

\item[\isty{courier}] (Courier) Only affects \gls{ttfamily}
and \gls{texttt}.
\end{description}

\xminisec{Notes:}

\begin{enumerate}

\item \warning Don't be tempted to use the \sty{times} package. It's
obsolete~\cite{l2tabu}. Use \isty{mathptmx} instead.

\item Although Times and Helvetica are commonly used together, they
don't match, as illustrated below (temporarily switching from this
book's fonts to Times and Helvetica):
\begin{codeS}
\glsni{rmfamily} xx \glsni{sffamily} xx
\end{codeS}
Results in:
\begin{resultS}[xx xx]
\fontfamily{ptm}\selectfont xx \fontfamily{phv}\selectfont xx
\end{resultS}
The first two x's are in Times Roman and the second two are in
Helvetica, which are somewhat larger. To compensate for this you
need to scale the Helvetica font using the \istyopt{helvet}{scaled}
option:
\begin{alltt}
\glsni{usepackage}\oarg{scaled=0.9}\marg{helvet}
\end{alltt}

\item Loading \isty{helvet} or \isty{courier} doesn't change the
default font family. Consider the following:
\begin{alltt}
\glsni{documentclass}\marg{scrartcl}

\glsni{usepackage}\marg{helvet}

\glsni{begin}\marg{document}
This is a sample document.
\glsni{end}\marg{document}
\end{alltt}
Here, the text \dq{This is a sample document} will be typeset in
Computer Modern Roman. This is because \gls{rmfamily} is the default
font and \isty{helvet} only affects \gls{sffamily}, which hasn't
been used. (See \sectionref{sec:renewcom} to find out how to change
the default font family.)
\end{enumerate}
\latexhtml{This book}{The PDF version of this document} has used the following packages:
\begin{alltt}
\glsni{usepackage}\oarg{T1}\marg{fontenc}
\glsni{usepackage}\oarg{math}\marg{anttor}
\glsni{usepackage}\marg{libris}
\end{alltt}
The \isty{fontenc} package is used to switch to Type~1 font
encoding\faq{Why bother with \sty{inputenc} and \sty{fontenc}?}{why-inp-font},
the \isty{anttor} package is used to set the serif family to Antykwa
Toru\'nska typeface, and the \isty{libris} package is used to set
the sans-serif family to the Libris ADF typeface.

\begin{exercise}{Fonts}{ex:fonts}

Go back to the document you created in \exerciseref{ex:simpledoc} and
change the first paragraph to a large bold font and the second
paragraph to normal size italic.  Emphasize the words \dq{simple} and
\dq{short}.  (Again, you can \downloadorview{fonts} the solution.)

If you like, you can try experimenting with loading
different font packages, such as \isty{mathptmx}, to change the
default typeface. The \LaTeX\ Font Catalogue~\cite{fontcat} provides
a useful list of fonts, although you may not have all of them
installed.
\end{exercise}

%%%%%%%%%%%%%% Aligning Material in Rows and Columns %%%%%%%%%%%%%%%%

\setnode{tabular}
\section{Aligning Material in Rows and Columns}
\label{sec:tabular}

Text can be aligned in rows and columns using the \gls{env-tabular}
environment.
\begin{definition}
\glsni{begin}\marg{tabular}\oarg{\meta{pos}}\marg{\meta{column specifiers}}
\end{definition}%
This \glsterm{environment} has a
\glsterm{mandatory} \meta{column specifiers}
that specifies how to align each column. Within \meta{column
specifiers}, there must be a specifier for each column. The three
basic are: \texttt{r} (right aligned), \texttt{l} (left aligned) and
\texttt{c} (centred).
(Make sure you don't confuse \texttt{l} (the letter \dq{ell}) with \texttt{1} 
(the digit one).)
The optional argument \meta{pos} is covered 
\latexhtml{in \sectionref{sec:boxes}}{\htmlref{later}{sec:boxes}}.

\xminisec{Example:}
Three columns (\underline{l}eft, \underline{c}entred, \underline{c}entred):
\begin{codeS}
\glsni{begin}\marg{tabular}\marg{lcc}
\end{codeS} 

\xminisec{Another Example:}
Four columns (\underline{l}eft, \underline{c}entred, \underline{c}entred, \underline{r}ight):
\begin{codeS}
\glsni{begin}\marg{tabular}\marg{lccr}
\end{codeS} 

The \texttt{r}, \texttt{l} and \texttt{c} specifiers don't allow
line breaks or paragraphs within a cell. It's not a good idea to
have too much text in a cell, but if it's required you can use
\begin{definition}
p\marg{\meta{width}}
\end{definition}
which indicates a paragraph cell of the given width.

\xminisec{Example:}
Three columns (\underline{p}aragraph of width 1in,
\underline{c}entred, \underline{r}ight):
\begin{codeS}
\glsni{begin}\marg{tabular}\marg{p\marg{1in}cr}
\end{codeS}
The paragraph cell will be formatted fully justified, which is often
inappropriate for a narrow block of text. The \isty{array} package
provides
\begin{definition}
\gls{gtcol}\marg{\meta{declaration}}
\end{definition}
which can be used directly in front of the \texttt{l}, \texttt{c},
\texttt{r} or \texttt{p} column specifiers. This inserts
\meta{declaration} in front of the entries for that column, so it
can be used to insert, say, \gls{raggedright}.

\xminisec{Example:}
Three columns, the first \underline{l}eft justified where each entry
is in bold, the second a \underline{p}aragraph column of width 1in set to
ragged right and the third \underline{c}entered:
\begin{codeS}
\glsni{begin}\marg{tabular}\marg{\glsni{gtcol}\marg{\glsni{bfseries}}l\glsni{gtcol}\marg{\glsni{raggedright}}p\marg{1in}c}
\end{codeS}

The \isty{array} package also provides
\begin{definition}
\gls{ltcol}\marg{\meta{declaration}}
\end{definition}
which can be used directly after the \texttt{l}, \texttt{c},
\texttt{r} or \texttt{p} column specifiers. This inserts
\meta{declaration} after the entries for that column.

\xminisec{Inter-Column Gap:}

The gap between columns is given by twice the value of the
\gls{length} register:
\begin{definition}
\gls{tabcolsep}
\end{definition}
A gap of \glsni{tabcolsep} is
also inserted before the first column and after the last column.
This length can be changed using one of the commands described in
\sectionref{sec:length}. For example:
\begin{codeS}
\gls{setlength}\marg{\glsni{tabcolsep}}\marg{4pt}
\end{codeS}
This will put an 8pt gap between columns and a 4pt gap before the
first column and after the last column.

The column specifiers can also include:
\begin{definition}
\gls{atchar}\marg{\meta{inter-column text}}
\end{definition}
This inserts \meta{inter-column text} at that place on each row of
the table, replacing the default inter-column gap.

\xminisec{Example:}

Suppose we want a centred first column, a right justified second
column and a left justified third column with a dot between the
second and third columns:
\begin{codeS}
\begin{alltt}
\glsni{begin}\marg{tabular}\marg{cr\glsni{atchar}\marg{.}l}
\end{alltt}
\end{codeS}
Alternatively, you may want a larger gap between groups of columns,
for example, two groups of three left justified columns:
\begin{codeS}
\begin{alltt}
\glsni{begin}\marg{tabular}\marg{lll\glsni{atchar}\marg{\glsni{hspace}\marg{4\glsni{tabcolsep}}}lll}
\end{alltt}
\end{codeS}
This uses the command:
\begin{definition}
\gls{hspace}\marg{\meta{length}}
\end{definition}
which inserts a horizontal space of a given \gls{length}. In this
case, four times the value of \glsni{tabcolsep}. This makes the gap
between the third and fourth columns twice as wide as the gap
between the other columns.

\setnode{rowscols}
\subsection{Column and Row Separation}
\label{sec:rowscols}

Remember the \latexhtml{special characters mentioned in
\sectionref{sec:chars}}{\htmlref{special characters}{sec:chars}}?
The ampersand character \gls{ampchar} is used to separate column
entries. Rows are separated using:
\begin{definition}
\gls{tab.dbbackslashchar}\oarg{\meta{vertical space}}
\end{definition}
where \meta{vertical space} is extra vertical spacing between rows,
if required. There is also a longer equivalent:
\begin{definition}
\gls{tabularnewline}
\end{definition}
\faq{Alignment tab changed to \cmdname{cr}}{altabcr}%
If you have used something like
\texttt{\gls{gtcol}\marg{\gls{raggedright}}p\marg{\meta{length}}}
as the specifier for your last column, you must use
\glsni{tabularnewline} instead of \glsni{tab.dbbackslashchar}
to indicate the row break otherwise you will get the following
error:
\begin{verbatim}
! Extra alignment tab has been changed to \cr.
<recently read> \endtemplate
\end{verbatim}
\screenpagebreak

\xminisec{Example:} 

Let's have two columns, the first left justified and the second right justified:
\begin{code}
\glsni{begin}\marg{tabular}\marg{lr}\newline
Video \glsni{ampchar}
8.99\glsni{tab.dbbackslashchar}\newline
CD \glsni{ampchar} 9.11\glsni{tab.dbbackslashchar}\newline
DVD \glsni{ampchar}
15.00\glsni{tab.dbbackslashchar}\newline
Total \glsni{ampchar} 33.10\newline
\glsni{end}\marg{tabular}
\end{code}

\begin{result}[tabular1.html]
\begin{tabular}{lr}
Video & 8.99\\
CD & 9.11\\
DVD & 15.00\\
Total & 33.10
\end{tabular}
\end{result}%
Recall from \chapterref{ch:def} that \LaTeX\ ignores spaces at the start of a line and
treats multiple spaces as a single space, so I~could just have
easily done:
\begin{code}
\glsni{begin}\marg{tabular}\marg{lr}\newline
\strut~~Video \glsni{ampchar} 8.99\glsni{tab.dbbackslashchar}\newline
\strut~~CD~~~~\glsni{ampchar} 9.11\glsni{tab.dbbackslashchar}\newline
\strut~~DVD~~~\glsni{ampchar} 15.00\glsni{tab.dbbackslashchar}\newline
\strut~~Total \glsni{ampchar} 33.10\newline
\glsni{end}\marg{tabular}
\end{code}%
and still have got the same result, but now the code is easier to
read.

Entries form implicit \glsing{group}, so
\glspl{declaration} made
within a \gls{env-tabular} environment only have an effect up to the
next \gls{ampchar} or \gls{tab.dbbackslashchar}.  

\xminisec{Example:}
\begin{code}
\glsni{begin}\marg{tabular}\marg{lr}\newline
\strut~~Video~\glsni{ampchar} 8.99\glsni{tab.dbbackslashchar}\newline
\strut~~CD~~~~\glsni{ampchar} 9.11\glsni{tab.dbbackslashchar}\newline
\strut~~DVD~~~\glsni{ampchar}
15.00\glsni{tab.dbbackslashchar}\newline
\strut~~\glsni{bfseries} Total \glsni{ampchar} 33.10\newline
\glsni{end}\marg{tabular}
\end{code}%
Output:\screenpagebreak
\begin{result}[Image: as the previous example except that the word
Total has appeared in bold]
\begin{tabular}{lr}
Video & 8.99\\
CD & 9.11\\
DVD & 15.00\\
\bfseries Total & 33.10
\end{tabular}
\end{result}%
Let's add an extra column and a header row:
\begin{code}\obeyspaces
\glsni{begin}\marg{tabular}\marg{lrr}\newline
\strut~~Item~~\glsni{ampchar} ex VAT \glsni{ampchar} inc
VAT\glsni{tab.dbbackslashchar}\newline
\strut~~Video~\glsni{ampchar}~~~8.99 \glsni{ampchar}~~10.56 \glsni{tab.dbbackslashchar}\newline
\strut~~CD~~~~\glsni{ampchar}~~~9.11 \glsni{ampchar}~~10.70 \glsni{tab.dbbackslashchar}\newline
\strut~~DVD~~~\glsni{ampchar}~~15.00 \glsni{ampchar}~~17.63 \glsni{tab.dbbackslashchar}\newline
\strut~~\glsni{bfseries} Total \glsni{ampchar}~~33.10 \glsni{ampchar}~~39.89\newline
\glsni{end}\marg{tabular}
\end{code}%
Output:
\begin{result}[tabular2.html]
\begin{tabular}{lrr}
Item  & ex VAT & inc VAT\\ 
Video & 8.99 & 10.56\\
CD & 9.11 & 10.70\\
DVD & 15.00 & 17.63\\
\bfseries Total & 33.10 & 39.89
\end{tabular}
\end{result}%

\xminisec{Example (Aligning on a Decimal Point):}

\label{xmp:decpt}If you want to align on the decimal point, it's best to use the
\isty{siunitx} package. That's beyond the scope of this
book, but for simple data this can be achieved
using the \gls{atchar} inter-column specifier. For example:
\begin{code}\obeyspaces
\glsni{begin}\marg{tabular}\marg{lr\glsni{atchar}\marg{.}l}\newline
\strut~~Video \glsni{ampchar} 8 \glsni{ampchar} 99\glsni{tab.dbbackslashchar}\newline
\strut~~CD~~~~\glsni{ampchar} 9 \glsni{ampchar} 11\glsni{tab.dbbackslashchar}\newline
\strut~~DVD~~~\glsni{ampchar} 15 \glsni{ampchar} 00\glsni{tab.dbbackslashchar}\newline
\strut~~\glsni{bfseries} Total \glsni{ampchar} 33 \glsni{ampchar} 10\newline
\glsni{end}\marg{tabular}
\end{code}%
Output:\reportpagebreak
\begin{result}[Image: as the second example but the decimal points
are aligned]
\begin{tabular}{lr@{.}l}
Video & 8 & 99\\
CD & 9 & 11\\
DVD & 15 & 00\\
\bfseries Total & 33 & 10
\end{tabular}
\end{result}%

\setnode{multicolumn}
\subsection{Spanning Columns}
\label{sec:multicolumn}

You may have noticed I~omitted the column headers in the \htmlref{above
example}{xmp:decpt}. \faq{Merging cells in
a column of a table}{multirow}
The problem with rewriting the table using 
\texttt{r\glsni{atchar}\marg{.}l} to align the decimal point is that the header
now needs to span the last two columns.
This can be done using the command:
\begin{definition}
\gls{multicolumn}\marg{\meta{cols spanned}}\marg{\meta{col specifier}}\marg{\meta{text}}
\end{definition}%
The first \gls{mandatory} \meta{cols spanned} is the number of
columns you want to span, the second argument \meta{col specifier}
indicates how to align this column-spanning entry, the third
argument \meta{text} indicates what should go in this entry. Note
that \meta{col specifier} should only have a single column
specifier, such as \texttt{\marg{c}} or \texttt{\marg{r}}.  We can
use \glsni{multicolumn} to modify an earlier example as
follows:\screenpagebreak
\begin{code}
\glsni{begin}\marg{tabular}\marg{lrr}\newline
\strut~~~~~~~~\glsni{ampchar} \glsni{multicolumn}\marg{2}\marg{c}\marg{Price (\glsni{pounds})}\glsni{tab.dbbackslashchar}\newline
\strut~~Item~~\glsni{ampchar} ex VAT \glsni{ampchar} inc
VAT\glsni{tab.dbbackslashchar}\newline
\strut~~Video~\glsni{ampchar}~~~8.99 \glsni{ampchar}~~10.56 \glsni{tab.dbbackslashchar}\newline
\strut~~CD~~~~\glsni{ampchar}~~~9.11 \glsni{ampchar}~~10.70 \glsni{tab.dbbackslashchar}\newline
\strut~~DVD~~~\glsni{ampchar}~~15.00 \glsni{ampchar}~~17.63 \glsni{tab.dbbackslashchar}\newline
\strut~~\glsni{bfseries} Total \glsni{ampchar}~~33.10 \glsni{ampchar}~~39.89\newline
\glsni{end}\marg{tabular}
\end{code}%
Output:
\begin{result}[tabular2multicol.html]
\begin{tabular}{lrr}
 & \multicolumn{2}{c}{Price (\pounds)}\\
Item  & ex VAT & inc VAT\\ 
Video & 8.99 & 10.56\\
CD & 9.99 & 11.74\\
DVD & 15.00 & 17.63\\
\bfseries Total & 33.98 & 39.93
\end{tabular}
\end{result}%
Here we are spanning two columns, so the first argument to \glsni{multicolumn} is \verb|{2}|,
we want the entry centred, so the second argument is \verb|{c}| and the text to go in this entry 
is simply \texttt{\marg{Price (\gls{pounds})}}.

\faq{How to alter the alignment
of tabular cells}{tabcellalign}The \glsni{multicolumn} command can also be used to override
the alignment of individual entries.  Consider
the following example:
\begin{code}
\glsni{begin}\marg{tabular}\marg{lrr}\newline
\strut~~~~~~~~~~\glsni{ampchar} Year1   \glsni{ampchar} Year2
\glsni{tab.dbbackslashchar}\newline
Travel~~~~\glsni{ampchar} 100,000 \glsni{ampchar}
110,000\glsni{tab.dbbackslashchar}\newline
Equipment \glsni{ampchar} 50,000  \glsni{ampchar} 60,000\newline
\glsni{end}\marg{tabular}
\end{code}
Output:
\begin{result}[tabular3.html]
\begin{tabular}{lrr}
 & Year1 & Year2 \\
Travel & 100,000 & 110,000\\
Equipment & 50,000 & 60,000
\end{tabular}
\end{result}%
In this example, the headers \dq{Year1} and \dq{Year2} would look
better centred, but the rest of the entries in the second
and third columns look best right aligned.  We can
use \glsni{multicolumn} to span just one column, and
use the second argument of \glsni{multicolumn} to override
the column specification:
\begin{code}
\glsni{begin}\marg{tabular}\marg{lrr}\newline
\strut~~~~~~~~~~\glsni{ampchar} \glsni{multicolumn}\marg{1}\marg{c}\marg{Year1}\newline
\strut~~~~~~~~~~\glsni{ampchar} \glsni{multicolumn}\marg{1}\marg{c}\marg{Year2}\glsni{tab.dbbackslashchar}\newline
Travel~~~~\glsni{ampchar} 100,000 \glsni{ampchar} 110,000\glsni{tab.dbbackslashchar}\newline
Equipment \glsni{ampchar} 50,000  \glsni{ampchar} 60,000\newline
\glsni{end}\marg{tabular}
\end{code}
Output:
\begin{result}[Image: as the previous example except that the words
'Year1' and 'Year2' have been centred in their columns]
\begin{tabular}{lrr}
 & \multicolumn{1}{c}{Year1} & \multicolumn{1}{c}{Year2} \\
Travel & 100,000 & 110,000\\
Equipment & 50,000 & 60,000
\end{tabular}
\end{result}

\setnode{tabrules}
\subsection{Rules}
\label{sec:rules}

In general, vertical rules are considered superfluous~\cite{oxford}.
Although Turabian~\cite{turabian96} allows for the possibility of
vertical rules for tabulated material containing more than two
columns but still advises against having too many and deprecates
the use of them at either end.

Horizontal rules may be used at the top and bottom of the
tabulated material, but other horizontal rules should be kept to a
minimum. In general, the top and bottom rule should be thicker than
the mid rules.

The \isty{booktabs} package provides:
\begin{definition}
\gls{toprule}\oarg{\meta{wd}}
\end{definition}
for the top horizontal rule,
\begin{definition}
\gls{bottomrule}\oarg{\meta{wd}}
\end{definition}
for the bottom horizontal rule, and
\begin{definition}
\gls{midrule}\oarg{\meta{wd}}
\end{definition}
for horizontal rules in between rows, such as after the header row.

These commands should all go at the start of the appropriate row.
This means that if you want a bottom rule, you need to add
\gls{tab.dbbackslashchar} followed by \glsni{bottomrule} at the end of
the tabulated material.

\xminisec{Example:}
\begin{code}
\glsni{begin}\marg{tabular}\marg{lrr}\newline
\glsni{toprule}\newline
\strut~~~~~~~~~~\glsni{ampchar}
\glsni{multicolumn}\marg{1}\marg{c}\marg{Year1}\newline
\strut~~~~~~~~~~\glsni{ampchar}
\glsni{multicolumn}\marg{1}\marg{c}\marg{Year2}%
\glsni{tab.dbbackslashchar}\newline
\glsni{midrule}\newline
Travel~~~~\glsni{ampchar} 100,000 \glsni{ampchar}
110,000\glsni{tab.dbbackslashchar}\newline
Equipment \glsni{ampchar} 50,000  \glsni{ampchar} 60,000%
\glsni{tab.dbbackslashchar}\newline
\glsni{bottomrule}\newline
\glsni{end}\marg{tabular}
\end{code}
results in:
\begin{result}[Image: as previous example but with horizontal rules
at the start and end and between the first and second rows.]
\begin{tabular}{lrr}
\toprule
          & \multicolumn{1}{c}{Year1} 
          & \multicolumn{1}{c}{Year2} \\
\midrule
Travel    & 100,000 & 110,000\\
Equipment & 50,000  & 60,000
\\\bottomrule
\end{tabular}
\end{result}

The thickness of the top and bottom rule is given by
the \gls{length} register \booklinebreak\gls{heavyrulewidth}, and the thickness of
the mid rule is given by the \gls{length} register
\gls{lightrulewidth}. These rule thicknesses can be overridden using the
optional argument \meta{wd} for \glsni{toprule}, \glsni{midrule}
and \glsni{bottomrule}.

\begin{exercise}{Aligning Material}{ex:tabular}
Go back to the document you created in \exerciseref{ex:simple} (and later modified in \exerciseref{ex:datetime}),
and add the following: 
\begin{result}[tabex.html]
\begin{tabular}{lrr}
 & \multicolumn{2}{c}{\bfseries Expenditure (\pounds)}\\
 & \multicolumn{1}{c}{Year1} & \multicolumn{1}{c}{Year2} \\
\bfseries Travel & 100,000 & 110,000\\
\bfseries Equipment & 50,000 & 60,000
\end{tabular}
\end{result}

Note that the \gls{env-tabular} environment
doesn't create a caption, all it does is arrange its contents in
rows and columns. You can find out how to turn your
\gls{env-tabular} environment into a table 
in \sectionref{sec:tables}.

You can \downloadorview{tabular} the solution to this exercise. (Remember to check
\appendixref{ch:errors} if you encounter an error message.)
\end{exercise}

For more information about using the \gls{env-tabular} environment see \latexbook,
\latexguide\ or \latexcomp.  \emph{The \LaTeX\ Companion} also describes how
to span rows using the \isty{multirow}
\htmlref{package}{sec:packages}.  For information on how to create
coloured tables using the \isty{colortbl}
\htmlref{package}{sec:packages}, see \latexgraphic.

\xminisec{Related \gls{ukfaq} topics:}
\begin{itemize}
\item \faqlink{How to change a whole row of a table}{wholerow}
\item \faqlink{Merging cells in a column of a table}{multirow}
\item \faqlink{Fixed width tables}{fixwidtab}
\item \faqlink{Variable-width columns in tables}{varwidcol}
\item \faqlink{Spacing lines in tables}{struttab}
\end{itemize}

%%%%%%%%%%%%%%%% Boxes and Mini-Pages %%%%%%%%%%%%%%%%%%%

\setnode{boxes}
\section{Boxes and Mini-Pages}
\label{sec:boxes}

\gls{tex} views everything on a page as a form of box. Each box has an
associated width, height and depth, and the boxes are placed together
on the page with \keyword{glue}. This is reminiscent of the days
of manual typesetting, where each letter or symbol was on a wooden
block, and the wooden blocks were glued in place. The simplest form
of box is a single letter. Some letters, such as \dq{a} only have a
height and width, whereas other letters, such as \dq{y} have a
height, width and depth (see \figureref{fig:letterbox}).

\begin{figure}[htbp]
\figconts[alt=Image of the letters y and a illustrating their
height depth and width. The letter a has zero depth as it does not
extend below the baseline.]
  {pictures/letterbox}%
  {%
    \caption{\texorpdfstring{\protect\TeX}{TeX}\protect\space Views Each Letter as a Box}%
  }%
  {fig:letterbox}%
\end{figure}

For example, the phrase \dq{cabbages and peas} is made up of 15 boxes:
\begin{center}
\makeimg{c a b b a g e s a n d p e a s}{\setlength{\fboxsep}{0.5pt}%
\setlength{\fboxrule}{0.1pt}\LARGE
\fbox{c}\fbox{a}\fbox{b}\fbox{b}\fbox{a}%
\fbox{g}\fbox{e}\fbox{s}~\fbox{a}\fbox{n}\fbox{d}~\fbox{p}\fbox{e}%
\fbox{a}\fbox{s}}
\end{center}
whereas the word \dq{cauliflower} consists of 10 boxes:\footnote{The 
fl-ligature is a single character, and so is one box not
two.}
\begin{center}
\makeimg{c a u l i fl-ligature o u r}{\setlength{\fboxsep}{0.5pt}%
\setlength{\fboxrule}{0.1pt}\LARGE
\fbox{c}\fbox{a}\fbox{u}\fbox{l}\fbox{i}%
\fbox{fl}\fbox{o}\fbox{w}\fbox{e}\fbox{r}}
\end{center}

More complicated boxes are made up of smaller boxes. We have already
encountered one of these more complicated boxes: the 
\gls{env-tabular} environment, discussed in the 
\htmlref{previous section}{sec:tabular}. This type of box is called
a \keyword{horizontal box}, which means that it can go in a line
of text. For example:
\begin{code}\obeyspaces
Here is some text.\newline
\glsni{begin}\marg{tabular}\marg{cc}\newline
A \glsni{ampchar} B\glsni{tab.dbbackslashchar}\newline
C \glsni{ampchar} D\newline
\glsni{end}\marg{tabular}\newline
The rest of the line.
\end{code}
produces:
\begin{result}[Image illustrating the typeset output: the tabulated
material is centred vertically within the line of text.]
Here is some text.
\begin{tabular}{cc}
A & B\\
C & D
\end{tabular}
The rest of the line.
\end{result}
Recall from \htmlref{the previous section}{sec:tabular} that the 
\gls{env-tabular} environment had an optional argument
\meta{pos}. This governs the vertical alignment when the
\gls{env-tabular} environment occurs within a line of text.
This can be one of \texttt{c} (centred\dash the default, as illustrated
above), \texttt{t} (top) and \texttt{b} (bottom). For example,
\begin{code}\obeyspaces
Here is some text.\newline
\glsni{begin}\marg{tabular}\oarg{b}\marg{cc}\newline
A \glsni{ampchar} B\glsni{tab.dbbackslashchar}\newline
C \glsni{ampchar} D\newline
\glsni{end}\marg{tabular}\newline
The rest of the line.
\end{code}%
produces:
\begin{result}[Image illustrating the typeset output: the last row
of the tabulated material rests on the baseline of the text outside it.]
Here is some text.
\begin{tabular}[b]{cc}
A & B\\
C & D
\end{tabular}
The rest of the line.
\end{result}

Since a box can't be broken across a line of text, you can use
the box making command:
\begin{definition}
\gls{mbox}\marg{\meta{text}}
\end{definition}
to prevent \meta{text} from spanning a line break.

\xminisec{Example:}
Compare:
\begin{code}
\gls{raggedright}
Some text at the beginning of a paragraph. Some
text in the middle of the paragraph. Some more text.
\glsni{par}
\end{code}
\begin{result}[Image of typeset output: paragraph is left-justified]
\raggedright
Some text at the beginning of a paragraph. Some text in the
middle of the paragraph. Some more text.
\par
\end{result}
with:
\begin{code}
\gls{raggedright}
Some text at the beginning of a paragraph. \glsni{mbox}\marg{Some
text in the middle of the paragraph.} Some more text.
\glsni{par}
\end{code}
\begin{result}[Image of typeset output: line break occurs at the
start of the second sentence leaving a large white space at the end
of the first line.]
\raggedright
Some text at the beginning of a paragraph. \mbox{Some text in the
middle of the paragraph.} Some more text.
\par
\end{result}
(If \gls{raggedright} had not been used, the text in the \glsni{mbox}
would've spilt out over the edge of the page.)

Another type of box which can again be placed in a line of text,
is the \gls{env-minipage} environment.
\begin{definition}
\glsni{begin}\marg{minipage}\oarg{\meta{pos}}\oarg{\meta{height}}%
\marg{\meta{width}}
\end{definition}
As the name suggests, this environment creates a \dq{mini-page}
of the given width.

\xminisec{Example:}
\begin{code}
Some text.\newline
\glsni{begin}\marg{minipage}\marg{2in}\newline
This is a mini-page. The text inside
it is formatted as usual.\newline
\mbox{}\newline
Paragraph breaks can also be used, but 
there is no indentation by default\gls{footnote}\marg{and 
this is how a footnote appears}.\newline
\glsni{end}\marg{minipage}\newline
The rest of the line.
\end{code}%
which produces:\reportpagebreak
\begin{result}[Image: the contents of the minipage including
footnote have been typeset in a rectangular block within the outer
line of text.]
Some text.
\begin{minipage}{2in}
This is a mini-page. The text inside
it is formatted as usual.

Paragraph breaks can also be used, but 
there is no indentation by default\mpexfootnote{and 
this is how a footnote appears}.
\end{minipage}
The rest of the line.
\end{result}

You can optionally specify a height, and how
the mini-page is aligned with the rest of the text. As with
the \gls{env-tabular} environment, the alignment
option \meta{pos} can be one of \texttt{t} (top), \texttt{c} (centred)
or \texttt{b} (bottom). The default is \texttt{c}, which is why
the above example has the mini-page centred vertically. This
can be changed, for example:
\begin{code}
Some text.\newline
\glsni{begin}\marg{minipage}\oarg{t}\marg{2in}\newline
This is a mini-page. The text inside
it is formatted as usual.\newline
\mbox{}\newline
Paragraph breaks can also be used, but 
there is no indentation by default\gls{footnote}\marg{and 
this is how a footnote appears}.\newline
\glsni{end}\marg{minipage}\newline
The rest of the line.
\end{code}
which produces
\begin{result}[Image: as previous example but the first line of the
minipage is aligned with the outer text.]
\begin{makeimage}
Some text.
\begin{minipage}[t]{2in}
This is a mini-page. The text inside
it is formatted as usual.

Paragraph breaks can also be used, but 
there is no indentation by default\mpexfootnote{and this
is how a footnote appears}.
\end{minipage}
The rest of the line.
\end{makeimage}
\end{result}

Note that the width can be specified relative to the current
line width, using the \gls{length} register \gls{linewidth}. For example,
\begin{alltt}
\glsni{begin}\marg{minipage}\marg{0.5\glsni{linewidth}}
\end{alltt}
will start a mini-page that is half the width of the current
line.

There is also a corresponding command
\begin{definition}
\gls{parbox}\oarg{\meta{pos}}\oarg{\meta{height}}%
\marg{\meta{width}}\marg{\meta{text}}
\end{definition}%
which behaves in a similar way. So the above example can be 
rewritten using a \glsni{parbox}:
\begin{code}
Some text.\newline
\glsni{parbox}[t]\marg{2in}\marg{This is a parbox. 
The text inside
it is formatted as usual.\newline
\mbox{}\newline
Paragraph breaks can also be used, but 
there is no indentation by default.}\newline
The rest of the line.
\end{code}
which produces
\begin{result}[Image: as previous example but without a footnote]
Some text.
\parbox[t]{2in}{This is a parbox. 
The text inside
it is formatted as usual.

Paragraph breaks can also be used, but 
there is no indentation by default.}
The rest of the line.
\end{result}
You may have noticed that the \gls{footnote} command has
not been used in the above example.  The \glsni{parbox} command 
is more restricted than the
\gls{env-minipage} environment, so you can't use the 
\gls{footnote} command in it. There are also certain 
environments, such as the 
\htmlref{list-making
environments}{sec:lists}\latex{\ifscreenorother{}{\ 
described in \sectionref{sec:lists}}}, that can be
used in a \gls{env-minipage} but not in a \glsni{parbox}.

\setnode{fbox}
\subsection{Framed Boxes}
\label{sec:fbox}

Recall the \gls{framebox} command described in \sectionref{sec:optional}:
\begin{definition}
\glsni{framebox}\oarg{\meta{width}}\oarg{\meta{align}}\marg{\meta{text}}
\end{definition}
This treats \meta{text} as a box of width \meta{width} and puts a frame around it.
The second optional argument may be one of: \texttt{c} (centred
contents), \texttt{l} (left-aligned contents), \texttt{r}
(right-aligned contents).

\xminisec{Example:}

\begin{codeS}
Some \glsni{framebox}\oarg{2in}\marg{framed} text.
\end{codeS}
\begin{resultS}[Some framed text (the word 'framed' is inside a 2
inch wide box)]
Some \framebox[2in]{framed} text.
\end{resultS}

There is a shorter related command with no optional arguments:
\begin{definition}
\gls{fbox}\marg{\meta{text}}
\end{definition}

The \isty{fancybox} package provides some additional framing
commands:
\begin{definition}
\gls{shadowbox}\marg{\meta{text}}
\end{definition}
Puts a shadow-style frame around its contents:\bookpagebreak
\begin{codeS}
Some \glsni{shadowbox}\marg{framed} text.
\end{codeS}
\begin{resultS}[Some framed text (the word 'framed' is in a 
box with a shadow)]
Some \shadowbox{framed} text.
\end{resultS}

\begin{definition}
\gls{doublebox}\marg{\meta{text}}
\end{definition}
Puts a double-lined frame around its contents:
\begin{codeS}
Some \glsni{doublebox}\marg{framed} text.
\end{codeS}
\begin{resultS}[Some framed text (the word 'framed' is in a 
box with a double lined frame)]
Some \doublebox{framed} text.
\end{resultS}

\begin{definition}
\gls{ovalbox}\marg{\meta{text}}
\end{definition}
Puts a thin-lined oval frame around its contents:
\begin{codeS}
Some \glsni{ovalbox}\marg{framed} text.
\end{codeS}
\begin{resultS}[Some framed text (the word 'framed' is in a 
a box with rounded corners)]
Some \ovalbox{framed} text.
\end{resultS}

\begin{definition}
\gls{Ovalbox}\marg{\meta{text}}
\end{definition}
Puts a thick-lined oval frame around its contents:

\begin{codeS}
Some \glsni{Ovalbox}\marg{framed} text.
\end{codeS}
\begin{resultS}[Some framed text (the word 'framed' is in a 
a thick bordered box with rounded corners)]
Some \Ovalbox{framed} text.
\end{resultS}

If you want a different frame effect, you will need to use a
graphical package, such as \isty{pgf}/\isty{tikz}.

\xminisec{Example:}

This example uses commands beyond the scope of this book,
but gives you an idea of what's possible.

\begin{code}
\begin{verbatim}
\documentclass{scrartcl}

\usepackage{tikz}
\usetikzlibrary{shapes}
\usetikzlibrary{decorations.pathmorphing}

\begin{document}

Some
\begin{tikzpicture}[baseline=(n.base),decoration=bumps]
\node[draw,ellipse,decorate] (n) {framed};
\end{tikzpicture}
text.

\end{document}
\end{verbatim}
\end{code}
\begin{resultS}[Some framed text (the word 'framed' is in a 
decorated with a wiggly border)]
Some
\begin{tikzpicture}[baseline=(n.base),decoration=bumps]
\node[draw,ellipse,decorate] (n) {framed};
\end{tikzpicture}
text.
\end{resultS}

For further details, see the \isty{pgf} \htmlref{documentation}{sec:texdoc}. 

\xminisec{Related \gls{ukfaq} topics:}
\begin{itemize}
\item \faqlink{Automatic sizing of minipage}{varwidth}
\item \faqlink{Float(s) lost}{fllost}
\item \faqlink{Perhaps a missing \cmdname{item}?}{errmissitem}
\end{itemize}

%%%%%%%%%%%%%%%%% CHAPTERS, SECTIONS ETC %%%%%%%%%%%%%%%%%%%%%%%

\setnode{structure}
\chapter{Structuring Your Document}
\label{ch:sections}

Let's go back to the document we modified in \exerciseref{ex:tabular}.
In this \latexhtml{chapter}{section} we shall edit that document
step by step until we have a fully-fledged document with title,
abstract, table of contents, sections etc.

%%%%%%%%%%%%%%%%%% Maketitle %%%%%%%%%%%%%%%%%%%%%%%%%%%%%%%%%

\setnode{title}
\section{Author and Title Information}
\label{sec:title}

The term \keyword{title page} is used to indicate the author, title and
date information that can appear either on the front cover by itself
or along the top of the first page of text.  In order to do this, you
must first specify the information.  Once this information has been
specified it can then be displayed.

The author, title and date are entered using the 
\glspl{command}:
\begin{definition}
\gls{author}\marg{\meta{author names}}
\\%\end{definition}
%\begin{definition}
\gls{title}\marg{\meta{title text}}
\\%\end{definition}
%\begin{definition}
\gls{date}\marg{\meta{document date}}
\end{definition}%
The \koma\ classes also define:
\begin{definition}
\gls{titlehead}\marg{\meta{Title heading}}\\
\gls{subject}\marg{\meta{Subject}}\\
\gls{subtitle}\marg{\meta{Subtitle}}\\
\gls{publishers}\marg{\meta{Publisher}}
\end{definition}

\faq{The style of document titles}{titlsty}All these title-related
commands only \emph{store} information, they don't actually display
anything. These commands can be put in the \gls{preamble}. With most
classes, you will typically need to use at least \glsni{author} and
\glsni{title}.

Once you have used these commands, you can then
display the information using the command:\refstepcounter{object}\label{obj:maketitle}
\begin{definition}
\gls{maketitle}
\end{definition}
This command should be placed where you want the title page to
appear, which is normally at the start of the \gls{env-document}
environment.

Note that if you don't use the \glsni{date} command, the current 
date will be inserted.  If you want no date to appear, you need to 
specify an empty argument:
\begin{codeS}
\glsni{date}\marg{}
\end{codeS}
\bookpagebreak
Multiple authors should be separated by the command \gls{and}.
For example:
\begin{code}
\begin{alltt}
\glsni{author}\marg{A. Jones\glsni{tab.dbbackslashchar}University of Somewhere \glsni{and}
B. Smith\glsni{tab.dbbackslashchar}University of Somewhere Else}
\end{alltt}
\end{code}%
Within these titling fields, you can also use the command:
\begin{definition}
\gls{thanks}\marg{text}
\end{definition}%
which produces a special type of footnote. For example:
\begin{codeS}
\glsni{title}\marg{A Great Project\glsni{thanks}\marg{funded by XYZ}}
\end{codeS}%
Note that the footnote marker produced using \glsni{thanks}
is considered to have zero width, so if it occurs in the middle
of a line, rather than the end, you will need to insert some
extra space using \gls{spacesym} (backslash space).
The argument of \glsni{thanks} is a \glslink{fragile}{moving 
argument}\index{argument!moving|hyperpage}.

\begin{exercise}{Creating Title Pages}{ex:title}

Try editing the document you modified in \exerciseref{ex:tabular} to
include title information.  Modifications are illustrated in bold
\modification{like this}:

\begin{bcode}
\begin{alltt}
\glsni{documentclass}\oarg{12pt}\marg{scrartcl}

\glsni{usepackage}\marg{datetime}

\modification{\glsni{title}\marg{A Simple Document}}
\modification{\glsni{author}\marg{Me}}

\glsni{begin}\marg{document}

\modification{\glsni{maketitle}}

This is a simple \glsni{LaTeX}\glsni{spacesym}document.
Here is the first paragraph.

Here is the second paragraph\glsni{footnote}\marg{with a footnote}.
As you can see it's a rather short paragraph, but not
as short as the previous one. This document was
created on: \glsni{today}\glsni{spacesym}at \glsni{currenttime}.

\glsni{begin}\marg{tabular}\marg{lrr}
 \glsni{ampchar} \glsni{multicolumn}\marg{2}\marg{c}\marg{\glsni{bfseries} Expenditure}\glsni{tab.dbbackslashchar}
 \glsni{ampchar} \glsni{multicolumn}\marg{1}\marg{c}\marg{Year1} \glsni{ampchar} \glsni{multicolumn}\marg{1}\marg{c}\marg{Year2}\glsni{tab.dbbackslashchar}
\glsni{bfseries} Travel \glsni{ampchar} 100,000 \glsni{ampchar} 110,000\glsni{tab.dbbackslashchar}
\glsni{bfseries} Equipment \glsni{ampchar} 50,000 \glsni{ampchar} 60,000
\glsni{end}\marg{tabular}

\glsni{end}\marg{document}
\end{alltt}
\end{bcode}

You can \download{title} this document.
\end{exercise}

%%%%%%%%%%%%%%%%%%%%%%%%% Abstract %%%%%%%%%%%%%%%%%%%%%%%%

\setnode{abstr}
\section{Abstract}
\label{sec:abstract}

The \gls{env-abstract} \glsterm{environment} is used to
create an \Index{abstract} for the document\faq{1-column abstract in
2-column document}{onecolabs}.  The way in which the
abstract is formatted depends on the class file\indexCLS.  The
\icls{scrreprt} class file will put the abstract on a page by itself,
some class files will indent the abstract and some will typeset the
abstract in italic.  Note also that some class files (such as
\icls{scrbook}) don't have an \glsni{env-abstract}
environment.  Abstracts traditionally go at the start of the document
after the title, so the \glsni{env-abstract} environment should go
after the \gls{maketitle} command.

\begin{exercise}{Creating an Abstract}{ex:abstr}

Try editing your document so that it has an abstract:
Modifications are illustrated \modification{like this}:

\begin{bcode}
\begin{alltt}
\glsni{documentclass}\oarg{12pt}\marg{scrartcl}

\glsni{usepackage}\marg{datetime}

\glsni{title}\marg{A Simple Document}
\glsni{author}\marg{Me}

\glsni{begin}\marg{document}

\glsni{maketitle}

\modification{\glsni{begin}\marg{abstract}}
\modification{A brief document to illustrate how to use \glsni{LaTeX}.}
\modification{\glsni{end}\marg{abstract}}

This is a simple \glsni{LaTeX}\glsni{spacesym}document.
Here is the first paragraph.

Here is the second paragraph\glsni{footnote}\marg{with a footnote}.
As you can see it's a rather short paragraph, but not
as short as the previous one. This document was
created on: \glsni{today}\glsni{spacesym}at \glsni{currenttime}.

\glsni{begin}\marg{tabular}\marg{lrr}
 \glsni{ampchar} \glsni{multicolumn}\marg{2}\marg{c}\marg{\glsni{bfseries} Expenditure}\glsni{tab.dbbackslashchar}
 \glsni{ampchar} \glsni{multicolumn}\marg{1}\marg{c}\marg{Year1} \glsni{ampchar} \glsni{multicolumn}\marg{1}\marg{c}\marg{Year2}\glsni{tab.dbbackslashchar}
\glsni{bfseries} Travel \glsni{ampchar} 100,000 \glsni{ampchar} 110,000\glsni{tab.dbbackslashchar}
\glsni{bfseries} Equipment \glsni{ampchar} 50,000 \glsni{ampchar} 60,000
\glsni{end}\marg{tabular}


\glsni{end}\marg{document}
\end{alltt}
\end{bcode}

You can \download{abstr} this document.
\end{exercise}

%%%%%%%%%%%%%%%%% Sections, Subsections etc %%%%%%%%%%%%%%%%%%%%%%%%%%%%

\setnode{sectionunits}
\section{Chapters, Sections, Subsections \ldots}
\label{sec:section}

Chapters, sections, subsections etc can be inserted using the
commands:
\begin{definition}
\gls{part}\oarg{\meta{short title}}\marg{\meta{title}}\\
\gls{chapter}\oarg{\meta{short title}}\marg{\meta{title}}\\
\gls{section}\oarg{\meta{short title}}\marg{\meta{title}}\\
\gls{subsection}\oarg{\meta{short title}}\marg{\meta{title}}\\
\gls{subsubsection}\oarg{\meta{short title}}\marg{\meta{title}}\\
\gls{paragraph}\oarg{\meta{short title}}\marg{\meta{title}}\\
\gls{subparagraph}\oarg{\meta{short title}}\marg{\meta{title}}
\end{definition}%
All these commands have a \latexhtml{moving argument (see
\sectionref{sec:fragile})}{\glslink{fragile}{moving
argument}\index{argument!moving|hyperpage}}, so fragile commands will need to
be protected using \gls{protect}. \faq{How to create a
\cmdname{subsubsub\-section}}{subsubsub}The final two commands in the
above list, \glsni{paragraph} and \glsni{subparagraph}, represent
subsubsubsections and subsubsubsubsections, although most
\glspl{cls} typeset their arguments as unnumbered running titles.

Note that the availability of these commands depends on the
\gls{cls} you are using. For example, the \icls{scrartcl} class
file that we have been using is designed for articles, so the
\glsni{chapter} command is not defined in that class, whereas it is
defined in the \icls{scrreprt} and \icls{scrbook} class files.

Each of the commands above has a \gls{mandatory} \meta{title} and an
\gls{optional} \meta{short title}. The mandatory argument
\meta{title} is simply the title of the chapter\slash section\slash subsection
etc. For
example:
\begin{codeS}
\glsni{section}\marg{Introduction}
\end{codeS}%
If you are using the \icls{scrartcl} class file, the output will look 
like:
\begin{resultS}[Image: The text '1 Introduction' is typeset in a
large bold font]
\textbf{\sffamily\Large 1\quad Introduction}
\end{resultS}%
Note that you don't specify the section number as \LaTeX\ does this
automatically.  This means that you can insert a new section or
chapter or swap sections around or even change a section to a
subsection etc, without having to worry about updating all the
section numbers\faq{The style of section headings}{secthead}.

\renewcommand{\fmtord}[1]{\textsuperscript{#1}}%
If you are using a class file\indexCLS\ that contains chapters as
well as sections, the section number will depend on the chapter.  So,
for example, the current section is the \ordinal{section}\ section
of chapter~\arabic{chapter}, so the section number is \thesection.
\faq{Why are my sections numbered 0.1 ...?}{zerochap}(Note that if
you are using a class file where the section number depends on the
chapter number, you must have a \gls{chapter} command before your
first \gls{section} command, otherwise your section numbers will
come out as 0.1, 0.2 etc.)

Unnumbered chapters/sections etc are produced by placing an 
asterisk \texttt{*} after the command name.  For example:
\begin{codeS}
\glsni{chapter}*\marg{Acknowledgements}
\end{codeS}

You can switch to appendices using the 
command
\begin{definition}
\gls{appendix}
\end{definition}%
then continue using \gls{chapter}, \gls{section} 
etc\faq{Appendixes}{appendix}.  For example (using the \icls{scrreprt} class file):
\begin{code}
\begin{alltt}
\glsni{appendix}
\glsni{chapter}\marg{Derivations}
Some derivations.

\glsni{chapter}\marg{Tables}
Some tables.
\end{alltt}
\end{code}

\xminisec{Note:}
The \koma\ classes have another type of sectioning command:
\begin{definition}
\gls{minisec}\marg{\meta{heading}}
\end{definition}\screenpagebreak
This provides an unnumbered heading not associated with any of the
structuring levels. For example, the above was produce using:
\begin{code}
\begin{alltt}
\glsni{minisec}\marg{Note:}
The \koma\ classes have another type of sectioning command:
\end{alltt}
\end{code}
The next note below was produced using:
\begin{code}
\begin{alltt}
\glsni{minisec}\marg{Important Note:}
If you want to change the font style used by headings, 
\glsni{emph}\marg{\glsni{bfseries} do not} use font declarations in the sectioning
command arguments.
\end{alltt}
\end{code}

\xminisec{Important Note:}
If you want to change the font style used by headings, 
\emph{\bfseries do not} use font declarations in the sectioning command
arguments. Don't do, for example:
\begin{alltt}\wrong
\glsni{chapter}\marg{\glsni{itshape} Introduction}
\end{alltt}
The \koma\ classes provide the command:
\begin{definition}
\gls{addtokomafont}\marg{\meta{element}}\marg{\meta{commands}}
\end{definition}
where \meta{element} is the name of a structuring element (no
backslash) and \meta{commands} is the list of font changing
declarations (see \tableref{tab:fontsII}) to apply to that element style.
For example, \latexhtml{this book}{the PDF version of this document}
uses the commands:
\begin{code}
\begin{alltt}
\glsni{addtokomafont}\marg{section}\marg{\glsni{rmfamily}\glsni{bfseries}}
\glsni{addtokomafont}\marg{minisec}\marg{\glsni{rmfamily}\glsni{bfseries}\glsni{scshape}}
\end{alltt}
\end{code}

\begin{exercise}{Creating Chapters, Sections etc}{ex:section}

Let's try editing our document so that it now has chapters, sections
and an appendix.  Since the \icls{scrartcl} class file doesn't have
chapters, let's change to the \icls{scrreprt} class.  Changes from our
previous document are shown \modification{like this}.

\begin{bcode}%
\begin{alltt}
\glsni{documentclass}\oarg{12pt}\marg{\modification{scrreprt}}

\glsni{usepackage}\marg{datetime}

\glsni{title}\marg{A Simple Document}
\glsni{author}\marg{Me}

\glsni{begin}\marg{document}
\glsni{maketitle}

\glsni{begin}\marg{abstract}
A brief document to illustrate how to use \glsni{LaTeX}.
\glsni{end}\marg{abstract}

\modification{\glsni{chapter}\marg{Introduction}}

\modification{\glsni{section}\marg{The First Section}}
This is a simple \glsni{LaTeX}\glsni{spacesym}document.
Here is the first paragraph.

\modification{\glsni{section}\marg{The Next Section}}
Here is the second paragraph\glsni{footnote}\marg{with a footnote}. 
As you can see it's a rather short paragraph, but not 
as short as the previous one. This document was 
created on: \glsni{today}\glsni{spacesym}at \glsni{currenttime}.

\modification{\glsni{chapter}\marg{Another Chapter}}

\modification{Here\textquoteright{}s another very interesting chapter.}
\modification{We\textquoteright{}re going to put a picture here later.}

\modification{\glsni{chapter}*\marg{Acknowledgements}}

\modification{I would like to acknowledge all those}
\modification{very helpful people who have assisted me in my work.}

\modification{\glsni{appendix}}

\modification{\glsni{chapter}\marg{Tables}}
\modification{We will turn this tabular environment into a table later.}

\glsni{begin}\marg{tabular}\marg{lrr}
 \glsni{ampchar} \glsni{multicolumn}\marg{2}\marg{c}\marg{\glsni{bfseries} Expenditure}\glsni{tab.dbbackslashchar}
 \glsni{ampchar} \glsni{multicolumn}\marg{1}\marg{c}\marg{Year1} \glsni{ampchar} \glsni{multicolumn}\marg{1}\marg{c}\marg{Year2}\glsni{tab.dbbackslashchar}
\glsni{bfseries} Travel \glsni{ampchar} 100,000 \glsni{ampchar} 110,000\glsni{tab.dbbackslashchar}
\glsni{bfseries} Equipment \glsni{ampchar} 50,000 \glsni{ampchar} 60,000
\glsni{end}\marg{tabular}

\glsni{end}\marg{document}
\end{alltt}
\end{bcode}

(You can \download{section} a copy of this file if you like, but
I~recommend that you try editing the file yourself to give you
practice.)
\end{exercise}

%%%%%%%%%%%%%%%%%%% Table of Contents %%%%%%%%%%%%%%%%%%%%%%%%

\setnode{toc}
\section{Creating a Table of Contents}
\label{sec:toc}

Once you have all your \htmlref{sectioning commands}{sec:section},
such as \gls{chapter} and \gls{section},
you can create a table of contents with the command
\begin{definition}
\gls{tableofcontents}
\end{definition}%
This command should go where you want your table of
contents to appear (usually after \gls{maketitle}).\faq{The format of the
Table of Contents, etc}{tocloft} The \koma\ classes provide two
options that govern the format of the table of contents:
\scrclsopt[graduated]{toc} and \scrclsopt[flat]{toc}. The first is the
default and indents the different sectioning levels. The second
doesn't use any indentation.

\xminisec{Example:}

\begin{codeS}
\glsni{documentclass}\oarg{12pt,toc=flat}\marg{scrreprt}
\end{codeS}

You may recall from \htmlref{the previous section}{sec:section} that
the sectioning commands all had an optional argument \meta{short
title}.  If your chapter or section title is particularly long\faq{My
section title is too wide for the page header}{runheadtoobig}, you
can use \meta{short title} to specify a shorter title that should go
in the table of contents.\footnote{\label{ftn:header}and in the page header, depending
on the page style.}  The longer title (given by the other
argument \meta{title}) will still appear in the section heading in
the main part of the document.

\LaTeX\ processes all source code sequentially, so when it first
encounters the \glsni{tableofcontents} command, it doesn't yet
know anything about the chapters, sections etc.  So the first time
the \htmlref{document is \LaTeX{}ed}{ch:tex2pdf} the necessary
information is written to the table of contents 
(\texttt{.toc})\indexTOC\ file (see \sectionref{sec:auxiliary}).  The subsequent pass reads the
information in from the \texttt{.toc}\indexTOC\ file, and generates
the table of contents.  You will therefore need to \LaTeX\ your
document twice to make sure that the table of contents is
up-to-date\faq{Numbers too large in table of contents,
etc}{tocloftwrong}.

\xminisec{Adding Extra Information}

The \htmlref{starred versions}{itm:starredcommand} of the sectional 
commands (such as \glsnl{chapter}\texttt{*})
don't get added to the table of contents. It may be that you want to
add it, in which case you need to use
\begin{definition}
\gls{addcontentsline}\marg{\meta{toc}}\marg{\meta{section
unit}}\marg{\meta{text}}
\end{definition}
after the heading. The first argument \meta{toc} is the file
extension without the dot. As mentioned above, the table of contents
file has the extension \texttt{.toc}, so the first argument should
be \texttt{toc} (\htmlref{later}{ch:floats}\latex{ in \chapterref{ch:floats}}, 
we'll be adding a list of figures and a list of tables, and those
have file extensions \texttt{.lof}\indexLOF\ and \texttt{.lot}\indexLOT, respectively).
The second argument \meta{section unit} is the name of the section
unit. This is just the name of the relevant sectioning command
\emph{without} the backslash. The final argument \meta{text} is the
entry text. For example (using \icls{scrreprt} class):
\begin{code}
\gls{chapter}*\marg{Acknowledgments}\newline
\glsni{addcontentsline}\marg{toc}\marg{chapter}\marg{Acknowledgements}
\end{code}

\begin{exercise}{Creating a Table of Contents}{ex:toc}

Try modifying your document so that it has a table of contents.
Modifications from the previous exercise are illustrated 
\modification{like this}:
\begin{bcode}
\begin{alltt}
\glsni{documentclass}\oarg{12pt}\marg{scrreprt}

\glsni{usepackage}\marg{datetime}

\glsni{title}\marg{A Simple Document}
\glsni{author}\marg{Me}

\glsni{begin}\marg{document}

\glsni{maketitle}

\modification{\glsni{tableofcontents}}

\glsni{begin}\marg{abstract}
A brief document to illustrate how to use \glsni{LaTeX}.
\glsni{end}\marg{abstract}

\glsni{chapter}\marg{Introduction}

\glsni{section}\marg{The First Section}

This is a simple \glsni{LaTeX}\glsni{spacesym}document.  Here is the first paragraph.

\glsni{section}\marg{The Next Section}

Here is the second paragraph\glsni{footnote}\marg{with a footnote}. 
As you can see it's a rather short paragraph, but not 
as short as the previous one. This document was 
created on: \glsni{today}\glsni{spacesym}at \glsni{currenttime}.

\glsni{chapter}\marg{Another Chapter}

Here's another very interesting chapter.
We're going to put a picture here later.

\glsni{chapter}*\marg{Acknowledgements}

I would like to acknowledge all those
very helpful people who have assisted
me in my work.

\glsni{appendix}
\glsni{chapter}\marg{Tables}

We will turn this tabular environment into a table later.

\glsni{begin}\marg{tabular}\marg{lrr}
 \glsni{ampchar} \glsni{multicolumn}\marg{2}\marg{c}\marg{\glsni{bfseries} Expenditure}\glsni{tab.dbbackslashchar}
 \glsni{ampchar} \glsni{multicolumn}\marg{1}\marg{c}\marg{Year1} \glsni{ampchar} \glsni{multicolumn}\marg{1}\marg{c}\marg{Year2}\glsni{tab.dbbackslashchar}
\glsni{bfseries} Travel \glsni{ampchar} 100,000 \glsni{ampchar} 110,000\glsni{tab.dbbackslashchar}
\glsni{bfseries} Equipment \glsni{ampchar} 50,000 \glsni{ampchar} 60,000
\glsni{end}\marg{tabular}


\glsni{end}\marg{document}
\end{alltt}
\end{bcode}

If your table of contents doesn't come out right, try \LaTeX{}ing it
again.  (Again, you can \download{toc} this file.)

You might want to try experimenting with the
\scrclsopt[flat]{toc} class options to
see what difference it makes:
\begin{codeS}
\glsni{documentclass}\oarg{12pt,toc=flat}\marg{scrreprt}
\end{codeS}

\end{exercise}

%%%%%%%%%%%%%%%%% Cross-Referencing %%%%%%%%%%%%%%%%%%%%%%%%%%%%

\setnode{crossref}
\section{Cross-Referencing}
\label{sec:crossref}

We have already seen that \LaTeX\ takes care of all the numbering for
the chapters etc, but what happens if you want to refer to a chapter
or section?  There's no point leaving \LaTeX\ to automatically
generate the section numbers if you have to keep track of them all,
and change all your cross-references every time you add a new
section.  Fortunately \LaTeX\ provides a way to generate the correct
number. All you have to do is label the part of the document you want
to reference, and then refer to this label when you want to
cross-reference it\faq{Referring to labels in other
documents}{extref}.  \LaTeX\ will then determine the correct number
that needs to be inserted at that point.

The first part, labelling the place you want to reference, is done 
using the command:
\begin{definition}
\gls{label}\marg{\meta{string}}
\end{definition}\refstepcounter{object}\label{obj:label}%
The \gls{argument} \meta{string} should be a
unique textual label.  This label can be anything you like as long as
it is unique, but it's a good idea to make it something obvious
so that, firstly, you can remember the label when you want to 
use it, and secondly, when you read through your code at some later
date, it's immediately apparent to you to which part of the
document you are referring.  People tend to have their own conventions
for labelling.  I usually start the label with two or three letters
that signify what type of thing I'm labelling.  For example, if I'm
labelling a chapter I'll start with \texttt{ch}, if I'm labelling a
section I'll start with \texttt{sec}. 

\xminisec{Examples:}

\begin{enumerate}
\item Labelling a chapter:

\begin{code}
\begin{alltt}
\glsni{chapter}\marg{Introduction}
\glsni{label}\marg{ch:intro}
\end{alltt}
\end{code}

\item Labelling a section:

\begin{code}
\begin{alltt}
\glsni{section}\marg{Technical Details}
\glsni{label}\marg{sec:details}
\end{alltt}
\end{code}
\end{enumerate}

Note that the \glsni{label} command doesn't produce any text, it
simply assigns a label.  You can now refer to that object
using the command:
\begin{definition}
\gls{ref}\marg{\meta{string}}
\end{definition}
which will produce the relevant number.

\xminisec{Example:}
\begin{codeS}
See Section \glsni{ref}\marg{sec:results} for an analysis
of the results.
\end{codeS}
\refstepcounter{object}\label{cmd:tilde}%
It is a typographical convention that you should never
start a new line with a number.  For example,
if you have the text \dq{Chapter~1} the \dq{1} must
be on the same line as the \dq{Chapter}.  We can 
do this by using an \keyword{unbreakable space}, which will
put a space but won't allow \LaTeX\ to break the line at
that point.  This is done using the tilde (\gls{tildechar}) 
\htmlref{special character}{sec:chars},
so the example above should actually be:\bookpagebreak
\begin{codeS}
See Section\glsni{tildechar}\glsni{ref}\marg{sec:results} for an analysis
of the results.
\end{codeS}
There is a similar command to reference the page number:
\begin{definition}
\gls{pageref}\marg{\meta{string}}
\end{definition}

\xminisec{Example:}
\begin{code}
See Chapter\glsni{tildechar}\glsni{ref}\marg{ch:def} on
page\glsni{tildechar}\glsni{pageref}\marg{ch:def} for a list of definitions.
\end{code}
The label \texttt{ch:def} obviously needs to
be defined somewhere:
\begin{code}
\begin{alltt}
\glsni{chapter}\marg{Definitions}
\glsni{label}\marg{ch:def}
\end{alltt}
\end{code}
In fact, I~have done this in my source code for \chapterref[the
definitions section]{ch:def} of this document, so the above example
would look like:
\begin{resultS}[See Chapter 2 on page 8 for a list of definitions]
See Chapter~\ref*{ch:def} on
page~\pageref*{ch:def} for a list of definitions.
\end{resultS}

\faq{Referring to things by their name}{nameref}%
It's not just chapters and sections that you can reference, most of
the numbers that \LaTeX\ automatically generates can be
cross-referenced. 

\xminisec{Example:} 
The source code for footnote~\ref{ftn:header} \latexhtml{on
page~\pageref{ftn:header}}{in \sectionref{sec:toc}} is:
\begin{code}
\glsni{footnote}\marg{\glsni{label}\marg{ftn:header}and in the page header, depending
on the page style}
\end{code}
\latexhtml{and it was referenced above using:}{and can be referenced
using:}
\begin{code}
The source code for footnote\glsni{tildechar}\glsni{ref}\marg{ftn:header} on
page\glsni{tildechar}\glsni{pageref}\marg{ftn:header} is:
\end{code}

The \isty{varioref} package provides a more convenient way of doing this using the
command:
\begin{definition}
\gls{vref}\marg{\meta{label}}
\end{definition}
This is like \gls{ref} but also adds information about the location, 
such as \dq{on page~\meta{n}} or \dq{on the following page}, if the
corresponding \gls{label} occurs on a different page, so the above
example can be changed to:
\begin{codeS}
The source code for footnote\glsni{tildechar}\glsni{vref}\marg{ftn:header} is:
\end{codeS}
\begin{latexonly}
\ifscreenorother
{%
which produces
\begin{resultS}
The source code for footnote~\vref{ftn:header} is:
\end{resultS}
}%
{%
\bookpagebreak\noindent which still produces
\begin{resultS}
The source code for footnote~\vref{ftn:header} is:
\end{resultS}
Compare with a reference to one of the labels in the next example:
\begin{codeS}
See step\glsni{tildechar}\cmdname{vref}\marg{itm:edit}.
\end{codeS}
which produces:
\begin{resultS}
See step~\vref{itm:edit}.
\end{resultS}
}%
\end{latexonly}
\begin{htmlonly}
which produces
\begin{resultS}[The source code for footnote 5.1 on page 110 is:]
The source code for footnote~\vref{ftn:header} is:
\end{resultS}
\end{htmlonly}

\xminisec{Caveat:}
You can run into trouble if the \glsni{vref} command occurs on a
page break. When it tries to insert the location information, such
as \dq{on the next page}, the information is no longer correct. This
can cause an \dq{Infinite loop} error. When this happens, either
edit your paragraph so the reference no longer falls on the page
break or use \gls{ref} instead of \glsni{vref} for that instance.

\xminisec{Another Example:}

The \gls{env-enumerate} environment \latex{\ifscreenorother{}{described 
in \sectionref{sec:enumerate} }}%
automatically numbers the items within an ordered list, so it's
possible to label list items. Recall the numbered list of
instructions at the start of
\chapterref{ch:tex2pdf}. Here's the code:

\begin{code}[0.9\linewidth]
\begin{alltt}
\glsni{begin}\marg{enumerate}

   \glsni{item}\glsni{label}\marg{itm:edit} Write or edit the source code.

   \glsni{item} Pass the source code to the \glsni{texttt}\marg{latex} or \glsni{texttt}\marg{pdflatex} 
        application (``\glsni{LaTeX}\glsni{spacesym}the document'').

      \glsni{begin}\marg{itemize}

         \glsni{item} If there are any error messages, 
               return to Step\glsni{tildechar}\glsni{ref}\marg{itm:edit}.

         \glsni{item} If there are no error messages, a PDF file
               is created, go to Step\glsni{tildechar}\glsni{ref}\marg{itm:view}.
      \glsni{end}\marg{itemize}

   \glsni{item}\glsni{label}\marg{itm:view} View the PDF file to check the result.

\glsni{end}\marg{enumerate}
\end{alltt}
\end{code}
Output:
\begin{result}[enumerateref.html]
\begin{enumerate}

   \item\label{itm:edit} Write or edit the source code.

   \item Pass the source code to the \texttt{latex} or \texttt{pdflatex} application 
        (\dq{\LaTeX\ the document}).

      \begin{itemize}

         \item If there are any error messages, 
               return to Step~\ref*{itm:edit}.

         \item If there are no error messages, a PDF file
               is created, go to Step~\ref*{itm:view}.
      \end{itemize}

   \item\label{itm:view} View the PDF file to check the result.

\end{enumerate}
\end{result}

The \gls{ref} and \gls{pageref} commands may come before or
after the corresponding \gls{label} command.  As with the table
of contents, \LaTeX\ first writes out all the cross-referencing
information to another file (the \iauxfile, see \sectionref{sec:auxiliary}) and then reads it in the
next time, so you will need to \LaTeX\ your document twice to get
everything up-to-date.

If the references aren't up-to-date, you will see the following
message at the end of the \LaTeX\ run:
\begin{verbatim}
LaTeX Warning: Label(s) may have changed. 
Rerun to get cross-references right.
\end{verbatim}
The following warning
\begin{verbatim}
LaTeX Warning: There were undefined references.
\end{verbatim}
means that \LaTeX\ found a reference to a label that
does not appear in the auxiliary file\indexAUX.  This could
mean that it's a new label, and the warning
will go away the next time you \LaTeX\ your document,
or it could mean that either you've forgotten to
define your label with the \gls{label} command,
or you've simply misspelt the label. The undefined references
will show up as two question marks \doublequestionmark\ in the \gls{output}.

\faq{\dq{Rerun} messages won't go away}{rerun}Very occasionally, if you have cross-references
and a table of contents, you might have to \LaTeX\ 
your document three times to get everything up to
date.  Just check to see if the \texttt{Label(s) may have changed}
warning appears.

If you find it inconvenient having to remember to click the typeset
button twice, you can use \iappname{latexmk}. This will run \LaTeX\
the required number of times to ensure the document is up-to-date.
To do this in TeXWorks, change the drop-down menu to \dq{LaTeXmk}, as
illustrated in \figureref{fig:latexmk}. Note that \appname{latexmk}
is a \gls{perl} script, so you need to make sure you have
\iappname{perl} installed (see \sectionref{sec:perl}).

\begin{figure}[htbp]
\figconts
 {pictures/texworks-latexmk}
 {%
   \caption{Selecting LaTeXmk in TeXWorks}
 }%
 {fig:latexmk}
\end{figure}

If \appname{latexmk} isn't listed in the drop-down menu, you can add
it via \menu{Edit}\menuto\menu{Preferences}. This opens the dialog
box shown in \figureref{fig:texworks-pref}. You can add a new tool
as follows:

\begin{enumerate}
\item To the right of the box
labelled \dq{Processing Tools} there is a button marked with a plus
(+) sign. Click on it to open the tool configuration dialog, shown
in \figureref{fig:texworks-toolconfig}.

\item Fill in the name \dq{LaTeXmk} in the box labelled \dq{Name} and either type
in the location of \appname{latexmk} in the box labelled \dq{Program} 
or use the \dq{Browse} button to locate it on your filing system.
(See \figureref{fig:texworks-toolconfig2}.) This will vary depending
on your operating system and \gls{tex}-distribution, but it will
probably be in a subdirectory (folder) called \texttt{bin} somewhere
in the \gls*{tex}-distribution tree.

\item There are lots of options that can be passed to
\appname{latexmk}, but if you want to produce PDF output you need to
add \texttt{-pdf} as an argument. This is done by clicking on the
button marked with a plus to the right of the \dq{Arguments} box and
type \texttt{-pdf}, as shown in \figureref{fig:texworks-toolconfig3}.

\item Another argument needs to be added that specifies the
basename of the \LaTeX\ file. This is done by again clicking on the
plus button and typing \verb|$basename|, as shown in
\figureref{fig:texworks-toolconfig4} .

\item Click on \dq{OK} to close the Tool Configuration dialog.

\item If you want to set \appname{latexmk} to be your default
processing tool, you can select it from the drop-down list labelled
\dq{Default}.

\item Click \dq{OK} when you're done.
\end{enumerate}

\begin{figure}[htbp]
\figconts
 {pictures/texworks-preferences}
 {%
    \caption{TeXWorks Preferences}
 }
 {fig:texworks-pref} 
\end{figure}

\begin{figure}[htbp]
\figconts
 {pictures/texworks-toolconfig1}
 {%
    \caption{Tool Configuration Dialog}
 }
 {fig:texworks-toolconfig} 
\end{figure}

\begin{figure}[htbp]
\figconts
 {pictures/texworks-toolconfig2}
 {%
    \caption{Tool Configuration Dialog: set the name and program
location}
 }
 {fig:texworks-toolconfig2} 
\end{figure}

\begin{figure}[htbp]
\figconts
 {pictures/texworks-toolconfig3}
 {%
    \caption{Tool Configuration Dialog: adding \protect\texttt{-pdf} argument}
 }
 {fig:texworks-toolconfig3} 
\end{figure}

\begin{figure}[htbp]
\figconts
 {pictures/texworks-toolconfig4}
 {%
    \caption{Tool Configuration Dialog: adding \protect\texttt{\$basename} argument}
 }
 {fig:texworks-toolconfig4} 
\end{figure}

\begin{exercise}{Cross-Referencing}{ex:crossref}

Try modifying your code so that it has cross-references.
Again, changes made from the previous exercise are
illustrated \modification{like this}:
\begin{bcode}
\begin{alltt}
\glsni{documentclass}\oarg{12pt}\marg{scrreprt}

\glsni{usepackage}\marg{datetime}

\glsni{title}\marg{A Simple Document}
\glsni{author}\marg{Me}

\glsni{begin}\marg{document}

\glsni{maketitle}

\glsni{tableofcontents}

\glsni{begin}\marg{abstract}
A brief document to illustrate how to use \glsni{LaTeX}.
\glsni{end}\marg{abstract}

\glsni{chapter}\marg{Introduction}
\modification{\glsni{label}\marg{ch:intro}}

\glsni{section}\marg{The First Section}

This is a simple \glsni{LaTeX}\glsni{spacesym}document.  Here is the first paragraph.
\modification{The next chapter is Chapter\glsni{tildechar}\glsni{ref}\marg{ch:another}}
\modification{and is on page\glsni{tildechar}\glsni{pageref}\marg{ch:another}.}
\modification{The next section is Section\glsni{tildechar}\glsni{ref}\marg{sec:next}.}

\glsni{section}\marg{The Next Section}
\modification{\glsni{label}\marg{sec:next}}

Here is the second paragraph\glsni{footnote}\marg{with a footnote}. 
As you can see it's a rather short paragraph, but not 
as short as the previous one. This document was 
created on: \glsni{today}\glsni{spacesym}at \glsni{currenttime}.

\glsni{chapter}\marg{Another Chapter}
\modification{\glsni{label}\marg{ch:another}}

Here's another very interesting chapter.
We're going to put a picture here later.
\modification{See Chapter\glsni{tildechar}\glsni{ref}\marg{ch:intro} for an }
\modification{introduction.}

\glsni{chapter}*\marg{Acknowledgements}

I would like to acknowledge all those
very helpful people who have assisted
me in my work.

\glsni{appendix}
\glsni{chapter}\marg{Tables}

We will turn this tabular environment into a table later.

\glsni{begin}\marg{tabular}\marg{lrr}
 \glsni{ampchar} \glsni{multicolumn}\marg{2}\marg{c}\marg{\glsni{bfseries} Expenditure}\glsni{tab.dbbackslashchar}
 \glsni{ampchar} \glsni{multicolumn}\marg{1}\marg{c}\marg{Year1} \glsni{ampchar} \glsni{multicolumn}\marg{1}\marg{c}\marg{Year2}\glsni{tab.dbbackslashchar}
\glsni{bfseries} Travel \glsni{ampchar} 100,000 \glsni{ampchar} 110,000\glsni{tab.dbbackslashchar}
\glsni{bfseries} Equipment \glsni{ampchar} 50,000 \glsni{ampchar} 60,000
\glsni{end}\marg{tabular}

\glsni{end}\marg{document}
\end{alltt}
\end{bcode}
(You can \download{crossref} a copy of this file.)

\end{exercise}

%%%%%%%%%%%%%%%%%%%% Bibliography %%%%%%%%%%%%%%%%%%%%%%%%%

\setnode{biblio}
\section{Creating a Bibliography}
\label{sec:bib}

If you have a large number of citations in your document, it's best
to use an external bibliographic application\faq{Creating a BibTeX
bibliography file}{buildbib}, such as \iappname{bibtex} or
\iappname{biber}. However, that is beyond the scope of this
book (see, instead, \latexguide, \latexcomp\
or \latexthesis). Therefore this section just gives a brief
explanation of the \gls{env-thebibliography} environment, which is
usually automatically generated using \appname{bibtex} or
\appname{biber}.

\begin{definition}
\glsni{begin}\marg{thebibliography}\marg{\meta{widest tag}}
\end{definition}
This environment is very similar to the list making environments described
in \sectionref{sec:lists}, but instead of \gls{item} use 
\begin{definition}
\gls{bibitem}\oarg{\meta{tag}}\marg{\meta{key}}
\end{definition}%
where \meta{key} is a unique keyword that identifies this
item.  Your keyword can be anything you like, but as with
\glsni{label} I~recommend that you use
a short memorable keyword. I~tend to use the first author's surname
followed by the year of publication. 

The bibliography heading depends on the class file you are using.
Most of the article-style classes, such as \icls{scrreprt}, use
\glsi{refname} (which produces \dq{\refname}) in an unnumbered
section, whereas the report and book-styles, such as \icls{scrreprt}
and \icls{scrbook}, use \glsi{bibname} (which produces
\dq{\bibname}) in an unnumbered chapter. See
\tableref{tab:textlab} for the list of the common textual label
commands.

Most class files don't automatically add the bibliography to the
table of contents. The \koma\ classes provide the
\scrclsopt{bibliography} option. This can be
\scrclsopt[totoc]{bibliography} (an unnumbered unit added to the
table of contents), for example,
\begin{codeS}
\glsni{documentclass}\oarg{bibliography=totoc}\marg{scrreprt}
\end{codeS}
or \scrclsopt[totocnumbered]{bibliography} (a
numbered unit added to the table of contents), for example,
\begin{codeS}
\glsni{documentclass}\oarg{bibliography=totocnumbered}\marg{scrreprt}
\end{codeS}
If you're not using
one of the \koma\ classes, consult the 
\htmlref{documentation}{sec:texdoc} for your class to see if
there is an equivalent option. Failing that, you can use
\gls{addcontentsline} (described in \sectionref{sec:toc}).
For example (using a class that defines chapters):
\begin{code}
\glsni{addcontentsline}\marg{toc}\marg{chapter}\marg{\glsni{bibname}}\\
\glsni{begin}\marg{bibliography}\marg{1}
\end{code}

\xminisec{Example:}
(This example uses the command \gls{LaTeXe} which produces the
\LaTeXe\ logo. This indicates the current version of \LaTeX\ rather
than the old 2.09 version.)\footnote{If a friend or colleague gives
you a file containing \cmdname{documentstyle} instead of \cmdname{documentclass}
they are nearly 20~years out of date.}
The class style in use is \cls{scrbook}, so the title is
given by \glsni{bibname} (\dq{Bibliography}).
\begin{code}
\glsni{begin}\marg{thebibliography}\marg{3}\newline
\glsni{bibitem}\marg{lamport94} \glsni{quotedblleft}\glsni{LaTeX}: a document preparation 
system\glsni{quotedblright}, Leslie Lamport, 2nd edition (updated for 
\glsni{LaTeXe}), Addison-Wesley (1994).\newline
\mbox{}\newline
\glsni{bibitem}\marg{kopka95} \glsni{quotedblleft}A Guide to \glsni{LaTeX}: document 
preparation for beginners and advanced users\glsni{quotedblright}, 
Helmut Kopka and Patrick W. Daly, Addison-Wesley (1995).\newline
\mbox{}\newline
\glsni{bibitem}\marg{goossens94} \glsni{quotedblleft}The 
\glsni{LaTeX}\glsni{spacesym}Companion\glsni{quotedblright}, 
Michel Goossens, Frank Mittelbach and
Alexander Samarin, Addison-Wesley, (1994).\newline
\mbox{}\newline
\glsni{end}\marg{thebibliography}
\end{code}
\begin{result}[references.html]
\dobibexamplenumbered
\end{result}

You can cite an item in your bibliography with the command
\begin{definition}
\gls{cite}\oarg{\meta{text}}\marg{\meta{key list}}
\end{definition}%

\bookpagebreak
\xminisec{Example:}
\begin{code}
\begin{alltt}
For more information about writing bibliographies see 
Goossens \glsni{emph}\marg{et al.}\glsni{tildechar}\glsni{cite}\marg{goossens94}.
\end{alltt}
\end{code}%
Output:
\begin{resultS}[cite.html]
For more information about writing
bibliographies see 
Goossens \emph{et al.}~[3].
\end{resultS}
If you want to cite multiple works, use 
a comma-separated list:
\begin{code}
\begin{alltt}
For more information about writing bibliographies 
see\glsni{tildechar}\glsni{cite}\marg{kopka95,goossens94}.
\end{alltt}
\end{code}%
Output:
\begin{resultS}[mcite.html]
For more information about writing bibliographies 
see~[2,3].
\end{resultS}
The \gls{optional} \meta{text} to
the \glsni{cite} command can be used to add text to the
citation.  

\xminisec{Example:}
\begin{code}
\begin{alltt}
For more information about writing bibliographies see 
Goossens \glsni{emph}\marg{et al.}\glsni{tildechar}\glsni{cite}\oarg{Chapter~13}\marg{goossens94}.
\end{alltt}
\end{code}%
Output:
\begin{result}[chcite.html]
For more information about writing bibliographies see 
Goossens \emph{et al.}~[3, Chapter~13].
\end{result}

The \gls{env-thebibliography} environment has
a \gls{mandatory}:
\begin{definition}
\glsni{begin}\marg{thebibliography}\marg{\meta{widest tag}}
\end{definition}%
The argument \meta{widest tag} is the widest tag in the list of
entries.  This helps \LaTeX\ to align the references correctly.  In
the example above, the tags appeared as: [1], [2] and [3], and [3]
is the widest so that was used as the argument. These tags can be
changed from the default numbers to something else using
the optional argument to the \gls{bibitem} command. 

\xminisec{Example (Textual Tags):}

This example uses the \gls{optional} of \gls{bibitem} to use textual
rather than numerical tags.  The widest tag is [Goossens 1994]
so that is chosen to be the argument of the \gls{env-thebibliography}
environment:
\begin{code}
\glsni{begin}\marg{thebibliography}\marg{Goossens 1994}\newline
\mbox{}\newline
\glsni{bibitem}[Lamport 1994]\marg{lamport94} 
\glsni{quotedblleft}\glsni{LaTeX}\cmdname{ } : a document 
preparation system\glsni{quotedblright}, Leslie Lamport, 2nd edition 
(updated for \glsni{LaTeXe}), Addison-Wesley (1994).\newline
\mbox{}\newline
\glsni{bibitem}[Kopka 1995]\marg{kopka95} 
\glsni{quotedblleft}A Guide to \glsni{LaTeX}: document 
preparation for beginners and advanced users\glsni{quotedblright}, Helmut Kopka 
and Patrick W. Daly, Addison-Wesley (1995).\newline
\mbox{}\newline
\glsni{bibitem}[Goossens 1994]\marg{goossens94} 
\glsni{quotedblleft}The \glsni{LaTeX}\glsni{spacesym}Companion\glsni{quotedblright}, 
Michel Goossens, Frank Mittelbach and
Alexander Samarin, Addison-Wesley, (1994).\newline
\mbox{}\newline
\glsni{end}\marg{thebibliography}
\end{code}
\begin{result}[references2.html]
\dobibexamplelabelled
\end{result}

\begin{exercise}{Creating a Bibliography}{ex:biblio}

Try adding the following chapter to your document:
\begin{code}
\glsni{chapter}\marg{Recommended Reading}

For a basic introduction to \glsni{LaTeX}\glsni{spacesym}see
Lamport\glsni{tildechar}\glsni{cite}\marg{lamport94}. For
more detailed information about \glsni{LaTeX}\glsni{spacesym}and
associated applications, consult Kopka and
Daly\glsni{tildechar}\glsni{cite}\marg{kopka95} or Goossens
\glsni{emph}\marg{et
al}\glsni{tildechar}\glsni{cite}\marg{goossens94}.
\end{code}% 
and also add the bibliography shown above to the end of your
document. You can \downloadorview{biblio} the solution, but have a
go by yourself first. Remember that, as before, you will need to
\LaTeX\ the document twice to get the references up-to-date, unless
you're using \iappname{latexmk} (as described in
\sectionref{sec:crossref}) in which case it will be done
automatically.

\end{exercise}

%%%%%%%%%%%%%%%%%%%%%%%%% Page Styles %%%%%%%%%%%%%%%%%%%%%%%%%%

\setnode{pagestyle}
\section{Page Styles and Page Numbering}
\label{sec:pagestyle}

You may have noticed that the documents you have created
have all had their page numbers automatically inserted 
at the foot of most of the pages\faq{Page numbering 
\dq{\meta{n} of \meta{m}}}{nofm}.  If you have created the document
that has gradually been modified over the previous few
sections, you may have noticed that the title page has
no header or footer, the table of contents starts on page~1,
the abstract page has no page number, and the pages after
the abstract start on page~1 and continue incrementally
onwards from that point.  All the page numbers are Arabic
numerals.  This can be changed using the command:
\begin{definition}
\gls{pagenumbering}\marg{\meta{style}}
\end{definition}%
where \meta{style} can be one of:
\begin{fwlist}{\pagenumberingfmt{arabic}}
\fwitem{\ipagenumbering{arabic}} Arabic numerals (1, 2, 3, \ldots)
\fwitem{\ipagenumbering{roman}} Lower case Roman numerals (i, ii, iii, \ldots)
\fwitem{\ipagenumbering{Roman}} Upper case Roman numerals (I, II, III, \ldots)
\fwitem{\ipagenumbering{alph}} Lower case alphabetical characters (a, b, c, \ldots)
\fwitem{\ipagenumbering{Alph}} Upper case alphabetical characters (A, B, C, \ldots)
\end{fwlist}

Traditionally, the front matter (table of contents, list of figures
etc) should have lower case Roman numeral page numbering, while the
main matter should be in Arabic numerals\faq{Page numbering by
chapter}{pagebychap}.  

\xminisec{Example:}
\begin{code}
\begin{alltt}
\glsni{author}\marg{Me}
\glsni{title}\marg{A Simple Document}
\glsni{maketitle}

\glsni{pagenumbering}\marg{roman}
\glsni{tableofcontents}

\glsni{begin}\marg{abstract}
This is the abstract.
\glsni{end}\marg{abstract}

\glsni{pagenumbering}\marg{arabic}
\glsni{chapter}\marg{Introduction}
\end{alltt}
\end{code}

The \icls{scrbook} class provides:
\begin{definition}
\gls{frontmatter}
\end{definition}
which switches to lower case Roman numeral page numbering, and
\begin{definition}
\gls{mainmatter}
\end{definition}
which switches to Arabic numeral page numbering.
These two declarations also
change the way the sectioning units, such as \gls{chapter} and
\gls{section}, appear.  The former, \glsni{frontmatter},
suppresses the numbering (regardless of whether or not you've used
the \htmlref{starred version}{itm:starredcommand} of the sectioning commands). The latter,
\glsni{mainmatter}, switches the numbering back on (unless
otherwise suppressed by using the starred sectioning commands). 
In addition, \icls{scrbook} provides
\begin{definition}
\gls{backmatter}
\end{definition}
which doesn't affect the page numbering but, like \glsni{frontmatter}, suppresses the 
sectional unit numbering.

\xminisec{Note:}

The \gls{env-abstract} environment isn't defined by the \icls{scrbook} class, as 
a book summary is usually incorporated into an introductory section.

\xminisec{Example:}

\begin{code}
\begin{alltt}
\glsni{documentclass}\oarg{12pt}\marg{scrbook}

\glsni{title}\marg{A Simple Document}
\glsni{author}\marg{Me}

\glsni{begin}\marg{document}
\glsni{maketitle}

\glsni{frontmatter}
\glsni{tableofcontents}

\glsni{chapter}\marg{Summary}
A brief document to
illustrate how to use \glsni{LaTeX}.

\glsni{mainmatter}

\glsni{chapter}\marg{Introduction}
\glsni{label}\marg{ch:intro}

\glsni{end}\marg{document}
\end{alltt}
\end{code}

The headers and footers can be changed using the command
\begin{definition}
\gls{pagestyle}\marg{\meta{style}}
\end{definition}%
Individual pages can be changed using
\begin{definition}
\gls{thispagestyle}\marg{\meta{style}}
\end{definition}%
Standard styles are:
\begin{fwlist}{myheadings}
\fwitem{\ipagestyle{empty}} No header or footer. 
\fwitem{\ipagestyle{plain}} Header empty, page number in footer.
\fwitem{\ipagestyle{headings}}  Header contains page number and various information, 
 footer empty. 
\fwitem{\ipagestyle{myheadings}} Header specified by user, footer empty.
\end{fwlist}

If the \ipagestyle{myheadings} style is used, the header information
can be specified using:
\begin{definition}
\gls{markboth}\marg{\meta{left head}}\marg{\meta{right head}}
\end{definition}%
if the \clsopt{twoside} option has been passed to the
\gls{cls} (default for \icls{scrbook}), or
\begin{definition}
\gls{markright}\marg{\meta{right head}}
\end{definition}%
if the \clsopt{oneside} option has been passed to the \gls{cls}
(default for \icls{scrartcl} and \icls{scrreprt}).

The \icls{scrreprt} class file uses the \pagestylefmt{empty} style
for the \Index{title} and \Index{abstract} pages and
\pagestylefmt{plain} for the first page of each new chapter.  By
default the remaining pages are also \pagestylefmt{plain}, but
these can be changed using the \glsni{pagestyle} command.
The \icls{scrbook} class defaults to the \pagestylefmt{headings}
style instead of \pagestylefmt{plain}.

\faq{Alternative head- and footlines in LaTeX}{fancyhdr}%
The \koma\ bundle provides a way to define new page
styles, but that's beyond the scope of this introductory tutorial. 
See the \koma\ documentation for further details if you are
interested.

\casemedia
 {% A4
  This book mostly uses the \pagestylefmt{headings} page style and the
  \icls{scrbook} class with the 
  \clsopt{oneside} option, so there is no difference between odd and
  even page headers, whereas the paperback version uses the
  \clsopt{twoside} option, so the odd pages display the chapter
   number and title and the even pages display the current section
   header and title. The on-screen PDF version of this book uses a page style 
  I~defined myself that incorporates a navigation bar in the footer.
 }%
 {% screen
   This screen version of the book uses a page style I~defined myself
   that incorporates a navigation bar in the footer. 
   The \htmladdnormallink{A4 version}{\baseurl/latex/novices/novices-a4.pdf} mostly uses
   the \pagestylefmt{headings} page style. If you look at it, you
   will see that the chapter number and title appear on the top left and
   the page number appears in the top right of most pages.  The 
   \clsopt{oneside} option was used, so there is no difference between
   the formatting of odd and even numbered pages. Whereas the paperback
   version uses the \clsopt{twoside} option so the odd and even page
   headers are different.
 }
 {% book
  This book mostly\footnote{I made some modifications to the page style 
  for the footers in the summary and index.}\ uses the \pagestylefmt{headings} 
  page style and the \icls{scrbook} class with the 
  \clsopt{twoside} option, so the odd and even page headers are
  different, whereas the A4 PDF version (available from
  \url{\baseurl/}) uses the \clsopt{oneside} option, so the odd and
  even page headers are the same.
  The on-screen PDF version of this book uses a page style 
  I~defined myself that incorporates a navigation bar in the footer.
 }
 {% html
   The \htmladdnormallink{A4
   version}{\baseurl/latex/novices/novices-a4.pdf} of this book
   mostly uses the \pagestylefmt{headings} page style.  If you look at it, you
   will see that the chapter number and title appear on the top left and
   the page number appears in the top right of most pages.  The 
   \clsopt{oneside} option was used, so there is no difference between
   the formatting of odd and even numbered pages. Whereas the paperback
   version uses the \clsopt{twoside} option, so the odd pages display the chapter
   number and title and the even pages display the current section
   header and title.
   The \htmladdnormallink{on-screen PDF
   version}{\baseurl/latex/novices/novices-screen.pdf} of this book
   uses a page style I~defined myself that incorporates a navigation bar
   in the footer. 
 }

\begin{exercise}{Page Styles and Page Numbering}{ex:pagestyle}

Try modifying your code so that it uses the \icls{scrbook} class,
\gls{frontmatter} and \gls{mainmatter}. Replace the
\gls{env-abstract} environment with an unnumbered chapter, as shown
below. Again, changes made from the previous document are illustrated
\modification{like this}:\screenpagebreak
\begin{bcode}
\begin{alltt}
\glsni{documentclass}\oarg{12pt}\marg{\modification{scrbook}}

\glsni{usepackage}\marg{datetime}

\modification{\glsni{pagestyle}\marg{headings}}

\glsni{title}\marg{A Simple Document}
\glsni{author}\marg{Me}

\glsni{begin}\marg{document}

\glsni{maketitle}

\modification{\glsni{frontmatter}}
\glsni{tableofcontents}

\modification{\glsni{chapter}\marg{Summary}}
A brief document to
illustrate how to use \glsni{LaTeX}.

\modification{\glsni{mainmatter}}

\glsni{chapter}\marg{Introduction}
\glsni{label}\marg{ch:intro}

\glsni{section}\marg{The First Section}

This is a simple \glsni{LaTeX}\glsni{spacesym}document.
Here is the first paragraph.
The next chapter is Chapter\glsni{tildechar}\glsni{ref}\marg{ch:another}
and is on page\glsni{tildechar}\glsni{pageref}\marg{ch:another}.
The next section is Section\glsni{tildechar}\glsni{ref}\marg{sec:next}.

\glsni{percentchar} Rest of document unchanged but 
\glsni{percentchar} omitted for brevity.
\glsni{end}\marg{document}
\end{alltt}
\end{bcode}
(You can \downloadorview{pagestyle} the edited
document.)

\end{exercise}

%%%%%%%%%%%%%%%%%%%%%%% babel %%%%%%%%%%%%%%%%%%%%%%%%%

\setnode{babel}
\section{Multi-Lingual Support: using the 
\texorpdfstring{\sty{babel}}{babel} package}
\label{sec:babel}

You may have noticed that the \gls{tableofcontents} and
\gls{chapter} commands have produced English words like
\dq{Contents} and \dq{Chapter}\faq{How to change LaTeX's \dq{fixed
names}}{fixnam}. If you are writing in another language, this is not
appropriate. In this case, you can use the \isty{babel} package,
and specify which language you will be using, either as an option to
the \isty{babel} package, or as an option to the class file\faq{Using
a new language with Babel}{newlang}.  If you are writing in more than
one language, list all the languages that you will be using where
the last named language is the default language\faq{Parallel
setting of text}{parallel}. For example:
\begin{alltt}
\glsni{usepackage}\oarg{english,french}\marg{babel}
\end{alltt}
or
\begin{alltt}
\glsni{documentclass}\oarg{english,french}\marg{scrreprt}
\glsni{usepackage}\marg{babel}
\end{alltt}
You can then switch between the named languages either using the 
\gls{declaration}:
\begin{definition}
\gls{selectlanguage}\marg{\meta{language}}
\end{definition}%
or the \gls{env-otherlanguage} \glsterm{environment}:
\begin{definition}
\glsni{begin}\marg{otherlanguage}\marg{\meta{language}}
\end{definition}%
These will affect all translations, including the date format
and predefined names like \dq{Chapter}. This also changes the
\gls{hyphenation} patterns.\latex{ (See \sectionref{sec:hyphenation}.)}

If you only want to set a short section of text in a different
language, without affecting the date format or predefined names,
then you can either use the command:
\begin{definition}
\gls{foreignlanguage}\marg{\meta{language}}\marg{\meta{text}}
\end{definition}%
or the starred version of the \glsni{env-otherlanguage} environment:
\begin{definition}
\glsni{begin}\marg{otherlanguage*}\marg{\meta{language}}
\end{definition}%
You can test to see if a given language is currently selected
using:
\begin{definition}
\gls{iflanguage}\marg{\meta{language}}\marg{\meta{true text}}%
\marg{\meta{false text}}
\end{definition}%

\xminisec{Example:}
\begin{code}
\begin{alltt}
\glsni{documentclass}\oarg{UKenglish,USenglish,french}\marg{scrartcl}
\glsni{percentchar} french is the last named option, so that's the current language

\glsni{usepackage}\oarg{T1}\marg{fontenc}
\glsni{usepackage}\oarg{utf8}\marg{inputenc}
\glsni{usepackage}\marg{babel}

\glsni{begin}\marg{document}
Ce texte est en fran\glsni{c}\marg{c}ais. La date aujourd'hui est: \glsni{today}.

\glsni{selectlanguage}\marg{USenglish}
This text is in US English. Today's date is: \glsni{today}.

\glsni{selectlanguage}\marg{UKenglish}
This text is in UK English. Today's date is: \glsni{today}.
\glsni{end}\marg{document}
\end{alltt}
\end{code}
Result:
\begin{result}[babel.html]
Ce texte est en fran\c{c}ais. La date aujourd'hui est~: 
{\datefrench\docdate}.

This text is in US English. Today's date is: {\usdate\docdate}.

This text is in UK English. Today's date is: {\ukenglishdate\docdate}.
\end{result}

\xminisec{Note:}
\warning If you are using the \istyopt{babel}{french} option, the
colon character (\Indextt{:}) is made active (that is, it's turned
into a special character) so if you are writing in French it's best
not to use a colon in labels (so where I've used, say,
\texttt{ch:def} you may need to change the colon to something else).

%%%%%%%%%%%%%%%%%%%%%%%% graphicx %%%%%%%%%%%%%%%%%%%

\setnode{graphicx}
\chapter{The \texorpdfstring{\sty{graphicx}}{graphicx} Package}
\label{ch:graphicx}

It is possible to generate images using \LaTeX\ commands\faq{Drawing
with TeX}{drawing} (See the \isty{pgf}\slash\isty{tikz} package or
\latexgraphic) however most people find it easier to create a
picture in some other application, and include that file into their
\LaTeX\ document.

PDF\LaTeX\ can insert PDF, PNG and JPG image files into your
document. If your image file is in a different format, you may be
able to find an application to convert it. \faq{Spawning programs
from (La)TeX: \cmdname{write}18}{spawnprog}Modern
\gls{tex}-distributions can automatically convert EPS\indexEPS\
files to PDF during the \LaTeX\ run using the \gls{perl} script
\iappname{epstopdf}. If your \gls{tex}-distribution doesn't support
this, you can convert your EPS file using \appname{epstopdf}
explicitly. For example, if you have an EPS image called, say,
\texttt{sample-image.eps}, you can convert it to a PDF image called
\texttt{sample-image.pdf}, by using the following command in a
\glslink{terminal}{terminal or command prompt}:
\begin{verbatim}
epstopdf sample-image.eps
\end{verbatim}
or (full path name may be required)
\begin{verbatim}
perl epstopdf sample-image.eps
\end{verbatim}

To insert an image file into your document, you first need to specify that 
you want to use the \isty{graphicx} package. So the following must 
go in the \gls{preamble}:
\begin{codeS}
\glsni{usepackage}\marg{graphicx}
\end{codeS}
The image can then be included in your document using the command
\begin{definition}
\gls{includegraphics}\oarg{\meta{key-val options}}\marg{\meta{filename}}
\end{definition}
where \meta{filename} is the name of your image file \emph{without the file
extension}\faq{\dq{Modern} graphics file names}{grffilenames}, and
\meta{key-val options} is a comma-separated list of options that can be
used to change the way the image is displayed. Note that if you have an 
image where the file name contain spaces or multiple dots, you need to
use the \isty{grffile} package:
\begin{codeS}
\glsni{usepackage}\marg{graphicx,grffile}
\end{codeS}

\xminisec{Example:} 

Suppose you had a file called \texttt{shapes.pdf}, then to
include it in your document you would do:
\begin{codeS}
\glsni{includegraphics}\marg{shapes}
\end{codeS}%
which would produce:
\begin{resultS*}
\incGraphics[alt=Image of some shapes]{exercises/shapes}
\end{resultS*}

You can specify a full or relative pathname, but you must use a
forward slash \gls{dir.slash} as the directory divider, even if you are using
Windows. For example:
\begin{codeS}
\glsni{includegraphics}\marg{pictures\glsni{dir.slash}shapes}
\end{codeS}
means the file \texttt{pictures/shapes.pdf} on Unix-type systems,
and it means the file \texttt{pictures\backslashsym shapes.pdf} on
Windows.\footnote{Or \texttt{shapes.png} or \texttt{shapes.jpg} or
\texttt{shapes.eps}. The example assumes a PDF image file.} This is
mainly because the backslash character is a \LaTeX\ special
character indicating a command, but it also helps portability
between platforms.

You can specify the order of the file types to look for with the command
\begin{definition}
\gls{DeclareGraphicsExtensions}\marg{\meta{ext-list}}
\end{definition}%
where \meta{ext-list} is a comma-separated list of extensions.  
For example, 
you might want to search first for PDF files, then for PNG files,
then for JPG files and finally for EPS files:
\begin{codeS}
\glsni{DeclareGraphicsExtensions}\marg{.pdf,.png,.jpg,.eps}
\end{codeS}%
The default for \iPDFLaTeX\ is:
\begin{flushleft}\ttfamily
.png\comma .pdf\comma .jpg\comma .mps\comma .jpeg\comma .jbig2\comma .jb2\comma .PNG\comma .PDF\comma .JPG\comma .JPEG\comma .JBIG2\comma .JB2\comma .eps
\end{flushleft}

The \gls{optional} \meta{key-val options}
should be a comma-separated list of
\meta{key}\texttt{=}\meta{value} pairs. Common options are:

\begin{fwlist}{\ikeyvalopt{includegraphics}{trim}=\meta{l} \meta{b} \meta{r} \meta{t}}
\fwitem{\ikeyvalopt{includegraphics}{angle}=\meta{\mathorit{x}}} rotate the image
by \degrees{x} anticlockwise.

\fwitem{\ikeyvalopt{includegraphics}{width}=\meta{length}} scale the image so that the
width is \meta{length}.  (Remember to specify the
\htmlref{units}{tab:units}.)

\fwitem{\ikeyvalopt{includegraphics}{height}=\meta{length}}  scale the image so that the
height is \meta{length}.  (Remember to specify the
\htmlref{units}{tab:units}.)

\fwitem{\ikeyvalopt{includegraphics}{scale}=\meta{value}} Scale the image by \meta{value}

\fwitem{\ikeyvalopt{includegraphics}{trim}=\meta{l} \meta{b} \meta{r} \meta{t}}  Specifies
the amount to remove from each side.  For example
\begin{codeS}
 \glsni{includegraphics}\oarg{trim=1 2 3 4}\marg{shapes}
\end{codeS}
crops the image by 1bp from the left, 2bp from the bottom, 3bp from
the right and 4bp from the top. (Recall the \iunit{bp}{big point} unit
from \tableref{tab:units}.)\reportpagebreak

\fwitem{\ikeyvalopt{includegraphics}{draft}}  Don't actually print the image, just draw a
box of the same size and print the filename inside it. 

\end{fwlist}

\xminisec{Example:}

This example first rotates the image by \degrees{45} anticlockwise,
then scales it so that the width is 1~inch.

\begin{codeS}
\glsni{includegraphics}\oarg{angle=45,width=1in}\marg{shapes}
\end{codeS}

\begin{resultS}[Image of shapes rotated then scaled]
\includegraphics[angle=45,width=1in]{exercises/shapes}
\end{resultS}

Note that this isn't the same as scaling and then rotating:
\begin{codeS}
\glsni{includegraphics}\oarg{width=1in,angle=45}\marg{shapes}
\end{codeS}

\begin{resultS}[Image of shapes scaled then rotated]
\includegraphics[width=1in,angle=45]{exercises/shapes}
\end{resultS}

You can also scale an image relative to the text area using the
\gls{length} registers \gls{textwidth} and \gls{textheight}. For
example, to scale a portrait image so that its height is
three-quarters of the text area height, you can do:
\begin{codeS}
\glsni{includegraphics}\oarg{height=0.75\glsni{textheight}}\marg{shapes}
\end{codeS}
or to scale a landscape image so that its width is half the text
area width, you can do:
\begin{codeS}
\glsni{includegraphics}\oarg{height=0.5\glsni{textwidth}}\marg{shapes}
\end{codeS}

\xminisec{Note:}

The \gls{includegraphics} command is another form of box (see
\sectionref{sec:boxes}), and can be used in the middle of a line of
text, just like the \gls{env-tabular} environment. See
\sectionref{sec:figures} to find out how to put the image in a
figure with a caption.

\xminisec{Example:}

\refstepcounter{object}\label{example:inlinegraphic}%
Recall the \iunit{ex}{relative unit} unit of measure from
\tableref{tab:units}. This can be used to scale
an image relative to the font size:
\begin{code}
An image can be inserted into a line
of text like this:
\glsni{includegraphics}\oarg{height=2ex}\marg{shapes}
\end{code}
\begin{resultS}[Image of typeset output: the included image has been
scaled so that it fits into the line of text]
An image can be inserted into a line
of text like this:
\includegraphics[height=2ex]{exercises/shapes}
\end{resultS}

\setnode{graphictrans}
\section{Graphical Transformations}
\label{sec:graphtrans}

The \isty{graphicx} package also provides commands to rotate, resize, 
reflect and scale text.  They are as follows:
\begin{definition}
\gls{rotatebox}\oarg{\meta{option list}}\marg{\meta{angle}}\marg{\meta{text}}
\end{definition}
Rotates \meta{text} by \meta{angle} (degrees anti-clockwise by
default). The optional argument \meta{option list} is a
comma-separated list of any of the following options:

\begin{itemize}
  \item \ikeyvalopt{rotatebox}{units}=\meta{number}

    The number of units in one full anti-clockwise rotation. So
    \texttt{units=-360} means that \meta{angle} specifies degrees
    clockwise whereas \texttt{units=6.283185} means that
    \meta{angle} is in radians.

  \item \ikeyvalopt{rotatebox}{origin}=\meta{label}

    The point of rotation.
    The value \meta{label} may contain one from either or both of
    the two lists: \texttt{lrc} (left, right, centre) and \texttt{tbB}
    (top, bottom, baseline). Alternatively the origin may be
    specified using the following two keys:

  \item \ikeyvalopt{rotatebox}{x}=\meta{dimen}

  \item \ikeyvalopt{rotatebox}{y}=\meta{dimen}
\end{itemize}

\xminisec{Example:}
\begin{code}
\begin{alltt}
base line
\glsni{rotatebox}\marg{45}\marg{Some text}
\glsni{rotatebox}\oarg{units=-360}\marg{45}\marg{Some text}
\glsni{rotatebox}\oarg{units=-360,origin=rB}\marg{45}\marg{Some text}
\glsni{rotatebox}\oarg{x=3em,y=3em}\marg{45}\marg{Some text}
base line
\end{alltt}
\end{code}\screenpagebreak
\begin{result}[Image of typeset output: each occurrence of 'Some text' has been rotated]
base line
\rotatebox{45}{Some text}
\rotatebox[units=-360]{45}{Some text}
\rotatebox[units=-360,origin=rB]{45}{Some text}
\rotatebox[x=3em,y=3em]{45}{Some text}
base line
\end{result}

\begin{definition}
\gls{scalebox}\marg{\meta{h scale}}\oarg{\meta{v scale}}\marg{\meta{text}}
\end{definition}
Scales \meta{text} by \meta{h scale} in both directions if \meta{v
scale} omitted, otherwise scales \meta{text} by \meta{h scale}
horizontally and \meta{y scale} vertically.

\xminisec{Example:}
\begin{codeS}
\glsni{scalebox}\marg{0.8}\marg{Some text}
\end{codeS}
\begin{resultS}[Some text (scaled by a factor of 0.8)]
\scalebox{0.8}{Some text}
\end{resultS}
Compare with:
\begin{codeS}
\glsni{scalebox}\marg{0.8}\oarg{1.2}\marg{Some text}
\end{codeS}
\begin{resultS}[Some text (aspect ratio changed)]
\scalebox{0.8}[1.2]{Some text}
\end{resultS}

\begin{definition}
\gls{reflectbox}\marg{\meta{text}}
\end{definition}
Reflects \meta{text} (equivalent to \glsni{scalebox}\marg{-1}\oarg{1}\marg{\meta{text}}).

\xminisec{Example:}
\begin{codeS}
\glsni{reflectbox}\marg{Some text}
\end{codeS}
\begin{resultS}[Some text (mirror image)]
\reflectbox{Some text}
\end{resultS}

\begin{definition}
\gls{resizebox}\marg{\meta{h length}}\marg{\meta{v length}}\marg{\meta{text}}
\end{definition}
Scales \meta{text} so that it is \meta{h length} wide and \meta{v
length} high. To preserve the aspect ratio, use \gls{resizebox.exclamchar}
instead of one of the dimensions.

\xminisec{Example:}
\begin{code}
\begin{alltt}
\glsni{resizebox}\marg{12mm}\marg{1cm}\marg{Some text}
\glsni{resizebox}\marg{\glsni{resizebox.exclamchar}}\marg{1cm}\marg{Some text}
\end{alltt}
\end{code}%
\begin{resultS}[Image of typeset output: two instances of 'Some text' scaled with and without changing the aspect ratio.]
\dohtmlcolorbox{%
\resizebox{12mm}{1cm}{Some text}
\resizebox{!}{1cm}{Some text}
}
\end{resultS}

\setnode{graphicxoptions}
\section{Package Options}

The \isty{graphicx} package can have the following options passed to 
it:
\begin{fwlist}{\optfmt{hiderotate}}
\fwitem{\mdseries \istyopt{graphicx}{draft}} Don't actually display the
images, just print the filename in a box of the correct size.  This
is useful if you want to print out a draft copy of a document to
check the text rather than the images.

\fwitem{\mdseries \istyopt{graphicx}{final}} Opposite of 
\istyopt{graphicx}{draft} (default).

\fwitem{\mdseries \istyopt{graphicx}{hiderotate}} Don't show rotated
text.

\fwitem{\mdseries \istyopt{graphicx}{hidescale}} Don't show scaled text.

\end{fwlist}

Remember that relevant options passed to the class file\indexCLS\ also 
affect packages.

\xminisec{Example (Draft Mode):}

Draft mode helps to speed up compilation of a large document when
you are editing the text. In the preamble:
\begin{codeS}
\glsni{usepackage}\oarg{draft}\marg{graphicx}
\end{codeS}
or
\begin{code}
\begin{alltt}
\glsni{documentclass}\oarg{draft}\marg{scrbook}

\glsni{usepackage}\marg{graphicx}
\end{alltt}
\end{code}
Later in the document:
\begin{codeS}
\glsni{includegraphics}\oarg{width=1in,angle=45}\marg{pictures/shapes}
\end{codeS}
\begin{resultS}[Only bounding box and file name drawn]
 \incGraphics{pictures/draftimage}
\end{resultS}

\begin{exercise}{Using the \sty{graphicx} Package}{ex:graphicx}

Download the image file
\htmladdnormallink{\texttt{shapes.pdf}}{exercises/shapes.pdf}%
\latex{ from \url{\exerciseurl/}} (or create your own image),
and include it into your document. (You
can \downloadorview{graphic} an example solution.)
\end{exercise}

\bookpagebreak\screenpagebreak
For more information on the \isty{graphicx} package see
\latexgraphic\ or the \isty{graphicx} documentation.

\xminisec{Related \gls{ukfaq} topics:}
\begin{itemize}
\item \faqlink{How to import graphics into (La)TeX documents}{impgraph}
\item \faqlink{Imported graphics in PDFLaTeX}{pdftexgraphics}
\item \faqlink{Imported graphics in dvips}{dvipsgraphics}
\item \faqlink{Imported graphics in dvipdfm}{dvipdfmgraphics}
\item \faqlink{Importing graphics from \dq{somewhere else}}{graphicspath}
\item \faqlink{Portable imported graphics}{graph-pspdf}
\item \faqlink{Repeated graphics in a document}{repeatgrf}
\item \faqlink{Limit the width of imported graphics}{grmaxwidth}
\item \faqlink{Top-aligning imported graphics}{topgraph}
\item \faqlink{Labelling graphics}{labelfig}
\item \faqlink{Graphics division by zero}{divzero}
\end{itemize}

%%%%%%%%%%%%%%%%%%%%%%%%%%% FLOATS %%%%%%%%%%%%%%%%%%%%%%%%%%%%%%%%%

\setnode{floats}
\chapter{Floats}
\label{ch:floats}

Figures and tables are referred to as \dq{floats} because they are
\emph{floated} to the nearest location. This prevents ugly large
spaces appearing on the page if there isn't enough room for the
figure or table before the page break. Floats have a caption and
associated number\faq{The style of captions}{captsty}. It is
customary for captions to appear at the bottom of figures but at the
top of tables~\cite{turabian96,oxford}.

For both figures and tables, the caption is generated using the 
command:
\begin{definition}
\gls{caption}\oarg{\meta{short caption}}\marg{\meta{text}}
\end{definition}%
Note that the \cmdname{caption} command has a
moving argument\index{argument!moving|hyperpage}, so
\glspl{fragile}\faq{Footnotes in captions}{ftncapt} will need to be
protected using \gls{protect}.  The \glsterm{optional}
\meta{short caption} is used to provide an alternative shorter
caption for the list of figures or list of tables, akin to the
optional argument to the \htmlref{sectioning
commands}{sec:section}\latex{\ifscreenorother{}{\ described in
\sectionref{sec:section}}}.

\xminisec{Note:}

Although the \glsni{caption} command can have an optional
short title, in general, captions should be brief. They should not
contain lots of description or background detail~\cite{turabian96}.
That type of information should be placed in the main text not the
caption.

\xminisec{Positioning:}

\refstepcounter{object}\label{obj:positioning}%
\faq{Wide figures in two-column documents}{widefigs}Both the 
\glsni{env-figure} and \glsni{env-table} environments have
an optional argument \meta{placement specifiers}, which indicates
permissible locations for the float. This may be a combination of
\texttt{h} (\dq{here}), \texttt{t} (top), \texttt{b} (bottom) and
\texttt{p} (page of floats.) Note that this only gives a general
guideline as to where the float will end up. The final location is
governed by other factors, such as space left on the page and the
proportion of text to floats on the page. If you omit one or more of the
placement specifiers, then you are prohibiting the float from being
placed in that location. A common mistake is to do
\begin{alltt}
\wrong\glsni{begin}\marg{figure}\oarg{h}
\end{alltt}
which says \dq{I~want the figure here and it can't go anywhere
else!} If the figure \emph{can't} be placed exactly here (for
example, there may not be enough room on the page), then you have
given it no alternative location, which can result in this and all
subsequent figures being dumped at the end of the chapter or
document, or can result in a fatal error when running
\LaTeX\faq{\dq{Too many unprocessed floats}}{tmupfl}. You may be
able to manage with only one of the other options, for example,
\begin{alltt}
\glsni{begin}\marg{figure}\oarg{t}
\end{alltt}
(In fact, modern \TeX\ distributions now replace \oarg{h} with
\oarg{t} if the float can't be placed.)
However, if you have a large number of floats it is advisable to
provide as many options as possible:
\begin{alltt}
\glsni{begin}\marg{figure}\oarg{htbp}
\end{alltt}
Similarly for tables.

If you are absolutely adamant that an image must go \dq{right here},
then it's not a float, and you shouldn't be using the 
\glsni{env-figure} environment. It's just a horizontal box, like the
example \xpageref{earlier}{example:inlinegraphic}. Similarly for
tabulated material.

It's worth bearing in mind what the Oxford Style Manual~\cite{oxford} has to say:
\begin{quote}
\dq{Text must not be read into it so as to give [the figure] an
explicit and fixed introduction, for example \sq{in the following
figure}: the final placement is determined by page breaks, which
cannot be anticipated before setting, and this makes rewording the
text necessary if the illustration does not fit the make-up of the
page.}
\end{quote}
Turabian~\cite{turabian96} gives the same advice (and reiterates it for 
figures):
\begin{quote}
\dq{All text references to a table should be by a number, not by an
introductory phrase such as \sq{in the following table}.}
\end{quote}

\setnode{figures}
\section{Figures}
\label{sec:figures}

Figures are created using the \gls{env-figure} environment.
\begin{definition}
\glsni{begin}\marg{figure}\oarg{\meta{placement specifiers}}
\end{definition}%
This environment may contain one or more captions (specified, as
described \htmlref{above}{ch:floats}, with
the \glsni{caption} command) but page breaks are not allowed in the
contents of a \glsni{env-figure} environment. The optional argument
\meta{placement specifiers} is as described
\htmlref{above}{obj:positioning}.

Recall from \chapterref{ch:graphicx} that we
can include an image in our document with
the command \gls{includegraphics} defined in the \isty{graphicx}
package.  We can put our
\htmladdnormallink{\texttt{shapes.pdf}}{exercises/shapes.pdf} image
into a figure as follows:
\begin{code}
\glsni{begin}\marg{figure}\oarg{htbp}\newline
\strut~~\glsni{includegraphics}\marg{shapes}\newline
\strut~~\glsni{caption}\marg{Some Shapes}\newline
\glsni{end}\marg{figure}
\end{code}%
So far so good, but our picture needs to be centred.  This can
be done using the \gls{centering} declaration mentioned in
\sectionref{sec:declaration}:
\begin{code}
\glsni{begin}\marg{figure}\oarg{htbp}\newline
\strut~~\glsni{centering}\newline
\strut~~\glsni{includegraphics}\marg{shapes}\newline
\strut~~\glsni{caption}\marg{Some Shapes}\newline
\glsni{end}\marg{figure}
\end{code}
The \gls{caption} command generates a number, just like
\gls{section}, so we can
\htmlref{cross-reference}{sec:crossref} it with \gls{ref} and
\gls{label}.  First, let's label the figure:
\begin{code}
\glsni{begin}\marg{figure}\oarg{htbp}\newline
\strut~~\glsni{centering}\newline
\strut~~\glsni{includegraphics}\marg{shapes}\newline
\strut~~\glsni{caption}\marg{Some Shapes}\newline
\strut~~\glsni{label}\marg{fig:shapes}\newline
\glsni{end}\marg{figure}
\end{code}%
Now we can reference it:
\begin{codeS}
Figure\glsni{tildechar}\glsni{ref}\marg{fig:shapes} shows some shapes.
\end{codeS}%
(As \htmlref{before}{cmd:tilde} we use \gls{tildechar}
to make an unbreakable space.)  This produces the following output
in the text:
\begin{resultS}[Figure 7.1 shows some shapes.]
Figure~\ref*{fig:shapes} shows some shapes.
\end{resultS}%
and produces \figureref{fig:shapes}.

\begin{figure}[hbtp]
\figconts
  {exercises/shapes}
  {%
    \caption{Some Shapes}%
  }%
  {fig:shapes}
\end{figure}

\xminisec{Important Note:}

If you want to change the caption font, \textbf{\em don't} do, e.g.:
\begin{alltt}
\wrong\glsni{caption}\marg{\gls{bfseries} Some Shapes}
\end{alltt}
Recall \gls{addtokomafont} from \sectionref{sec:section}. This can
also be used to change the fonts used by the caption.\bookpagebreak
\begin{alltt}
\correct\glsni{addtokomafont}\marg{caption}\marg{\glsni{bfseries}}
\end{alltt}
Similarly for the caption label. For example:
\begin{codeS}
\glsni{addtokomafont}\marg{captionlabel}\marg{\glsni{scshape}}
\end{codeS}

\xminisec{List of Figures}

Just as we were able to generate a \htmlref{table of
contents}{sec:toc} using \gls{tableofcontents}, we can also
generate a list of figures using the command
\begin{definition}
\gls{listoffigures}
\end{definition}
This creates a file with the extension \texttt{.lof}\indexLOF\ (see
\sectionref{sec:auxiliary}).
As with \reportlinebreak\glsni{tableofcontents} you will need to \LaTeX\ your document twice to
get the list of figures up-to-date, unless
you're using \iappname{latexmk} (as described in
\sectionref{sec:crossref}) in which case it will be done
automatically.

\begin{exercise}{Creating Figures}{ex:figures}

If you did \exerciseref{ex:graphicx}, you should have a document
with an image in it. You now need to put this image into a
\gls{env-figure} environment. Remember to centre the image, and give
the figure a caption. Next, try labelling the figure and referencing
it in the text. You could also put in a list of figures after the
table of contents.
You can \downloadorview{figures} an example.

\end{exercise}

\setnode{sidebysidefigs}
\subsection{Side-By-Side Figures}
\label{sec:sidebysidefigs}

Recall at the start of \sectionref{sec:figures}, I mentioned that
the \gls{env-figure} environment may contain one or more captions.
In most cases, you'll just have a single caption per
\glsni{env-figure} environment, but sometimes you may want to have
two figures side-by-side, in which case you'll need two captions
within the same \gls{env-figure} environment in order to keep the
figures together.

To do this, we can use the 
\gls{env-minipage} environment, which was covered 
in \sectionref{sec:boxes}. Recall that 
the \gls*{env-minipage} environment creates a 
\Index{horizontal box}, which means that two mini-pages can
be placed side-by-side on the same line. All you need to do 
now, is place one image and caption in one mini-page, and the
other image and caption in the neighbouring mini-page. (Do you remember what effect is obtained
by placing a \htmlref{percent symbol}{obj:suppresseol} at the end of a
line?)\doifbook{\footnote{See page~\pageref{obj:suppresseol}.}}\bookpagebreak
\begin{code}
\glsni{begin}\marg{figure}\oarg{htbp}\newline
\strut~\glsni{begin}\marg{minipage}\marg{0.5\glsni{linewidth}}\newline
\strut~~\glsni{centering}\newline
\strut~~\glsni{includegraphics}\marg{circle}\newline
\strut~~\glsni{caption}\marg{A Circle}\newline
\strut~~\glsni{label}\marg{fig:circle}\newline
\strut~\glsni{end}\marg{minipage}\glsni{percentchar}\newline
\strut~\glsni{begin}\marg{minipage}\marg{0.5\glsni{linewidth}}\newline
\strut~~\glsni{centering}\newline
\strut~~\glsni{includegraphics}\marg{rectangle}\newline
\strut~~\glsni{caption}\marg{A Rectangle}\newline
\strut~~\glsni{label}\marg{fig:rectangle}\newline
\strut~\glsni{end}\marg{minipage}\newline
\glsni{end}\marg{figure}
\end{code}%
\latex{The above code produces 
Figures~\ref{fig:circle-sbs} and~\vref{fig:rectangle-sbs}. }Note that each
mini-page uses \gls{centering} to centre its contents, and the
label is also placed in the same mini-page, after the
\gls{caption} command. If the \gls{label} was not in the same
\htmlref{scope}{sec:group} as the \gls{caption}, the
cross-reference would be incorrect.

A common mistake when trying to create side-by-side figures is
to do:
\begin{alltt}
\glsni{begin}\marg{figure}[htbp]
 \glsni{begin}\marg{minipage}\marg{0.5\glsni{linewidth}}
  \glsni{centering}
  \glsni{includegraphics}\marg{circle}
  \glsni{caption}\marg{A Circle}
  \glsni{label}\marg{fig:circle}
 \glsni{end}\marg{minipage}
\wrong
 \glsni{begin}\marg{minipage}\marg{0.5\glsni{linewidth}}
  \glsni{centering}
  \glsni{includegraphics}\marg{rectangle}
  \glsni{caption}\marg{A Rectangle}
  \glsni{label}\marg{fig:rectangle}
 \glsni{end}\marg{minipage}
\glsni{end}\marg{figure}
\end{alltt}
This produces one figure on top of the other, instead of
side-by-side\html{, as illustrated in \figureref{fig:circle-sbs} and
\figureref{fig:rectangle-sbs}}. This is because the blank
line indicates a paragraph break, so each minipage is in a separate
paragraph, so it's not possible for them to be on the same line.

If you want a bit of spacing in your code to make it more readable,
use \gls{percentchar} to comment out the paragraph break. For
example:
\begin{alltt}
\glsni{end}\marg{minipage}%
\correct\glsni{percentchar}
\glsni{begin}\marg{minipage}\marg{0.5\glsni{linewidth}}
\end{alltt}


\begin{figure}[htbp]
  \begin{makeimage}\end{makeimage}
  \begin{minipage}{0.5\linewidth}
    \figuretop{fig:circle-sbs}
    \centering
    \incGraphics[alt=Image of a circle]{exercises/circle}
    \caption{A Circle}
    \label{fig:circle-sbs}
  \end{minipage}%
  \begin{minipage}{0.5\linewidth}
    \figuretop{fig:rectangle-sbs}
    \centering
    \incGraphics[alt=Image of a rectangle]{exercises/rectangle}
    \caption{A Rectangle}
    \label{fig:rectangle-sbs}
  \end{minipage}
\end{figure}

%%%%%%%%%%%%%%%%%%%%%% Tables %%%%%%%%%%%%%%%%%%%%%%%%%

\setnode{tables}
\section{Tables}
\label{sec:tables}

Tables are produced in much the same way as figures, except that the
\gls{env-table} environment is used instead.
\begin{definition}
\glsni{begin}\marg{table}\oarg{\meta{placement specifiers}}
\end{definition}
Where the optional argument \meta{placement specifiers} is as described
\xpageref{earlier}{obj:positioning}.

As mentioned at the \htmlref{start of this chapter}{ch:floats},
tables typically have the caption at the top of the
table~\cite{turabian96}. With the \koma\ classes, such as
\icls{scrartcl}, \icls{scrreprt} and \icls{scrbook}, use the class
option \scrclsopt[tableabove]{captions} to ensure that the vertical
spacing appears correctly between the caption and the table content
and put \gls{caption} at the start of the \envname{table}
environment. \faq{Tables longer than a single page}{longtab}Page
breaks are not permitted in the \glsni{env-table}
environment. (The \isty{longtable} package can be used for that
instead. See the \sty{longtable} \htmlref{documentation}{sec:texdoc} for 
further details.)

\xminisec{Example:}

In the \glsterm{preamble}:
\begin{codeS}
\glsni{documentclass}\oarg{captions=tableabove}\marg{scrbook}
\end{codeS}
Later in the document:
\begin{code}
\begin{alltt}
\glsni{begin}\marg{table}\oarg{htbp}
 \glsni{caption}\marg{A Sample Table} 
 \glsni{label}\marg{tab:sample}
 \glsni{centering}
 \glsni{begin}\marg{tabular}\marg{lr}
 Item  \glsni{ampchar} Cost\glsni{tab.dbbackslashchar}
 Video \glsni{ampchar} 8.99\glsni{tab.dbbackslashchar}
 CD    \glsni{ampchar} 9.99\glsni{tab.dbbackslashchar}
 DVD   \glsni{ampchar} 15.00
 \glsni{end}\marg{tabular}
\glsni{end}\marg{table}
\end{alltt}
\end{code}%
This produces \tableref{tab:sample}.

\begin{table}[htbp]
\caption{A Sample Table}
\label{tab:sample}
\centering
\begin{makeimage}
\begin{tabular}{lr}
Item & Cost\\
Video & 8.99\\
CD    & 9.99\\
DVD   & 15.00
\end{tabular}
\end{makeimage}
\end{table}

Again, the \gls{centering} declaration is used to centre the
\gls{env-tabular} environment.
As with figures, you can create a list of tables using the command
\begin{definition}
\gls{listoftables}
\end{definition}
This creates a file with the extension \texttt{.lot}\indexLOT\ (see
\sectionref{sec:auxiliary}).
As with the table of contents and list of figures, you will need to
\LaTeX\ your document twice to get the list of tables up-to-date, unless
you're using \iappname{latexmk} (as described in
\sectionref{sec:crossref}) in which case it will be done
automatically.

\begin{exercise}{Creating Tables}{ex:tables}
If you did \exerciseref{ex:tabular}, you should have 
a \gls{env-tabular} environment in your document.  Try
turning this into a table, and add \tableref{tab:sample}.  
You could also try adding a list of tables.
As before, you can \downloadorview{tables} the solution.
\end{exercise}

\setnode{sidebysidetab}
\subsection{Side-by-Side Tables}
\label{sec:sidebysidetab}

You can create side-by-side tables using an analogous method to 
the side-by-side figures approach described
\htmlref{above}{sec:sidebysidefigs}.

\xminisec{Example:}

This example is similar to the one in \sectionref{sec:sidebysidefigs}.
Again, take care to ensure that there is no paragraph break between the
two \gls{env-minipage} environments.
\begin{code}
\begin{alltt}
\glsni{begin}\marg{table}
 \glsni{begin}\marg{minipage}\marg{0.5\glsni{linewidth}}
  \glsni{caption}\marg{Prices for 2011}
  \glsni{label}\marg{tab:prices2011}
  \glsni{centering}
  \glsni{begin}\marg{tabular}\marg{lr}
   Item \glsni{ampchar} Price (\pounds)\glsni{tab.dbbackslashchar}
   Widgets \glsni{ampchar} 10.99\glsni{tab.dbbackslashchar}
   Whatsits \glsni{ampchar} 5.99
  \glsni{end}\marg{tabular}
 \glsni{end}\marg{minipage}\glsni{percentchar}
\glsni{percentchar}
 \glsni{begin}\marg{minipage}\marg{0.5\glsni{linewidth}}
  \glsni{caption}\marg{Prices for 2012}
  \glsni{label}\marg{tab:prices2012}
  \glsni{centering}
  \glsni{begin}\marg{tabular}\marg{lr}
   Item \glsni{ampchar} Price (\glsni{pounds})\glsni{tab.dbbackslashchar}
   Widgets \glsni{ampchar} 11.99\glsni{tab.dbbackslashchar}
   Whatsits \glsni{ampchar} 6.99
  \glsni{end}\marg{tabular}
 \glsni{end}\marg{minipage}\glsni{percentchar}
\glsni{end}\marg{table}
\end{alltt}
\end{code}
This produces Tables~\ref{tab:prices2011} and~\vref{tab:prices2012}.

\begin{table}[htbp]
 \begin{minipage}{0.5\linewidth}
  \caption{Prices for 2011}
  \label{tab:prices2011}
  \centering
  \begin{makeimage}
  \begin{tabular}{lr}
   Item & Price (\pounds)\\
   Widgets & 10.99\\
   Whatsits & 5.99
  \end{tabular}
  \end{makeimage}
 \end{minipage}%
%
 \begin{minipage}{0.5\linewidth}
  \caption{Prices for 2012}
  \label{tab:prices2012}
  \centering
  \begin{makeimage}
  \begin{tabular}{lr}
   Item & Price (\pounds)\\
   Widgets & 11.99\\
   Whatsits & 6.99
  \end{tabular}
  \end{makeimage}
 \end{minipage}%
\end{table}

\setnode{sidewaysfloats}
\section{Sideways Floats}
\label{sec:sidewaysfloats}

The \isty{rotating} package provides the \gls{env-sidewaysfigure}
environment:
\begin{definition}
\glsni{begin}\marg{sidewaysfigure}
\end{definition}
and the \gls{env-sidewaystable} environment:
\begin{definition}
\glsni{begin}\marg{sidewaystable}
\end{definition}
which are like \gls{env-figure} and \gls{env-table}, respectively, but rotate the
entire float (including caption) sideways. This sideways float is
always placed on a page of its own.

If you have used the \clsopt{twoside} class option (or you are
using a class like \icls{scrbook}, which defaults to that option)
then the sideways floats will be rotated clockwise or
anti-clockwise, depending on whether they fall on an even (verso) or
odd (recto) numbered page. (Requires a second \LaTeX\ run to get it
correct.) 

\xminisec{Example:}
\begin{code}
\begin{alltt}
\glsni{begin}\marg{sidewaysfigure}
 \glsni{centering}
 \glsni{includegraphics}\oarg{width=0.75\glsni{textheight}}\marg{shapes}
 \glsni{caption}\marg{A Sideways Figure}
 \glsni{label}\marg{fig:sideways}
\glsni{end}\marg{sidewaysfigure}
\end{alltt}
\end{code}
The above code produces \figureref{fig:sideways}.

\dosidewaysfigure
 {fig:sideways}
 {\includegraphics[width=0.75\textheight]{exercises/shapes}}
 {A Sideways Figure}

\begin{htmlonly}
  \begin{figure}
    \figuretop{fig:sideways}
    \centering
\makeimg{A sideways figure}{%
\colorbox{white}{%
    \rotatebox[br]{90}{%
      \begin{tabular}{c}
      \includegraphics{exercises/shapes}\\
      {\usekomafont{captionlabel}Figure~\ref*{fig:sideways}}\captionformat A Sideways Figure
      \end{tabular}
    }%
}%
}
    \caption{A Sideways Figure}
    \label{fig:sideways}
  \end{figure}
\end{htmlonly}

\setnode{subfloats}
\section{Sub-Floats}
\label{sec:subfloats}

Some floats have sub-floats within them. For example, a figure may
contain several sub-figures, each of which requires a caption.
The simplest way to do this is to use the \isty{subcaption} package
that provides the \gls{env-subfigure} and \gls{env-subtable} environments:
\begin{definition}
\glsni{begin}\marg{subfigure}\oarg{\meta{pos}}\marg{\meta{width}}
\end{definition}
\begin{definition}
\glsni{begin}\marg{subtable}\oarg{\meta{pos}}\marg{\meta{width}}
\end{definition}
Within these environments, you can use \gls{caption} to create a
sub-caption. (In addition to the main \gls{caption} for the
containing \gls{env-figure} or \gls{env-table}
environment.)

\xminisec{Note:}
The \isty{subcaption} package requires the \isty{caption} package,
but doesn't automatically load it, so you'll need to load
both:\reportpagebreak
\begin{codeS}
\glsni{usepackage}\marg{caption,subcaption}
\end{codeS}

\xminisec{Example:}

This is very similar to the side-by-side figures example from
\sectionref{sec:sidebysidefigs}.
 
\begin{code}[0.9\linewidth]
\glsni{begin}\marg{figure}\oarg{htbp}\newline
\strut~\glsni{begin}\marg{subfigure}\oarg{b}\marg{0.5\glsni{linewidth}}\newline
\strut~~\glsni{centering}\newline
\strut~~\glsni{includegraphics}\marg{rectangle}\newline
\strut~~\glsni{caption}\marg{Rectangle}\glsni{label}\marg{fig:rectangle}\newline
\strut~\glsni{end}\marg{subfigure}\glsni{percentchar}\newline
\glsni{percentchar}\newline
\strut~\glsni{begin}\marg{subfigure}\oarg{b}\marg{0.5\glsni{linewidth}}\newline
\strut~~\glsni{centering}\newline
\strut~~\glsni{includegraphics}\marg{circle}\newline
\strut~~\glsni{caption}\marg{Circle}\glsni{label}\marg{fig:circle}\newline
\strut~\glsni{end}\marg{subfigure}\glsni{percentchar}\newline
\glsni{caption}\marg{Two Shapes}\newline
\glsni{label}\marg{fig:shape}\newline
\glsni{end}\marg{figure}
\end{code}

This produces \figureref{fig:shapes2}.
Elsewhere in the document, the figure and its components can be
referenced:
\begin{code}
Figure\gls{tildechar}\gls{ref}\marg{fig:shapes2} shows some shapes.  
Figure\gls{tildechar}\gls{ref}\marg{fig:rectangle} shows a rectangle and 
Figure\gls{tildechar}\gls{ref}\marg{fig:circle} shows a circle.
\end{code}
which produces the following text:
\begin{result}[Figure 7.5 shows some shapes. Figure 7.5a shows a
rectangle and Figure 7.5b shows a circle.]
Figure~\ref*{fig:shapes2} shows some shapes.
Figure~\ref*{fig:rectangle} shows a rectangle and
Figure~\ref*{fig:circle} shows a circle.\relax
\end{result}

You can also reference just the sub-float using
\begin{definition}
\gls{subref}\marg{\meta{label}}
\end{definition}
which is analogous to \gls{ref}, but only displays the sub-float
number without the number associated with its containing float.

\xminisec{Example:}
\begin{code}
Figure\glsni{tildechar}\cmdname{ref}\marg{fig:shapes2}
shows: (\glsni{subref}\marg{fig:rectangle}) a rectangle and
(\glsni{subref}\marg{fig:circle}) a circle.
\end{code}
produces
\begin{resultS}[Figure 7.5 shows: (a) a rectangle and (b) a circle.]
Figure~\ref*{fig:shapes2} shows: (\subref*{fig:rectangle}) a rectangle and (\subref*{fig:circle}) a circle.
\end{resultS}

\begin{figure}[htbp]
\begin{makeimage}\end{makeimage}
\figuretop{fig:shapes2}
\centering
\begin{latexonly}
 \begin{subfigure}[b]{0.5\linewidth}
   \centering
   \includegraphics{exercises/rectangle}
   \caption{Rectangle}\label{fig:rectangle}
 \end{subfigure}%
 \begin{subfigure}[b]{0.5\linewidth}
   \centering
   \includegraphics{exercises/circle}
   \caption{Circle}\label{fig:circle}
 \end{subfigure}%
\end{latexonly}
\begin{htmlonly}
\begin{tabular}{cc}
\htmladdimg{exercises/rectangle.png} \hspace{0.8in}&
\htmladdimg{exercises/circle.png}\\
(a) Rectangle & (b) Circle
\end{tabular}
\end{htmlonly}
\caption{Two Shapes}
\label{fig:shapes2}
\end{figure}

\xminisec{Note:}

The subfigure labels (a, b, etc) should typically be in
italic~\cite{turabian96}. This can be achieved with the
\isty{caption} package using:
\begin{definition}
\gls{DeclareCaptionLabelFormat}\marg{\meta{format-name}}\marg{\meta{code}}
\end{definition}
where \meta{format-name} is the name for this new format and
\meta{code} is the code used to format the label
where \texttt{\glsni{hashchar}2} gets replaced by the reference number.

Once you have defined a new format, you can then use
\begin{definition}
\gls{captionsetup}\oarg{\meta{type}}\marg{\meta{options}}
\end{definition}
to switch to that new format. For subfloats, \meta{type} needs to be
set to \texttt{sub}. The second argument \meta{options} is a
\meta{key}=\meta{value} comma-separated list. The key that sets the format is
\ikeyvalopt{captionsetup}{labelformat}. (For further details about
both \glsni{DeclareCaptionLabelFormat} and \glsni{captionsetup}, see
the \isty{caption} package \htmlref{documentation}{sec:texdoc}.)

For example, to create a format called \texttt{em-noparens} that
displays the number in an emphasized font without parentheses:
\begin{codeS}
\glsni{DeclareCaptionLabelFormat}\marg{em-noparens}\marg{\glsni{emph}\marg{\glsni{hashchar}2}}
\end{codeS}
Now switch to that new format:
\begin{codeS}
\glsni{captionsetup}\oarg{sub}\marg{labelformat=em-noparens}
\end{codeS}

Note that this only changes the caption label format. It doesn't affect the font 
used by \gls{ref} or \glsni{subref}.  For \gls{ref}, you can use the 
\isty{fncylab} package, which provides the command:
\begin{definition}
\gls{labelformat}\marg{\meta{ctr}}\marg{\meta{defn}}
\end{definition}
Within \meta{defn}, use \glsni{hashchar}1 to represent the subfigure value
and use \glsi{thefigure} for the encapsulating figure number. For
example:
\begin{codeS}
\glsni{labelformat}\marg{subfigure}\marg{\glsni{thefigure}\glsni{emph}\marg{\glsni{hashchar}1}}
\end{codeS}
Now
\begin{codeS}
\glsni{ref}\marg{fig:circle}
\end{codeS}
will produce
\begin{resultS}[7.5a (the letter a is emphasized)]
\ref*{fig:shapes2}\emph{a}
\end{resultS}
Unfortunately, this doesn't work for \glsni{subref}. Instead you will have 
to do, for example, the following in the text:
\begin{codeS}
\glsni{emph}\marg{\glsni{subref}\marg{fig:circle}}
\end{codeS}

If you want to add parentheses, the above can be modified
to:
\begin{code}
\glsni{DeclareCaptionLabelFormat}\marg{em-parens}\marg{(\glsni{emph}\marg{\glsni{hashchar}2})}\newline
\glsni{captionsetup}\oarg{sub}\marg{labelformat=em-parens}\newline
\glsni{labelformat}\marg{subfigure}\marg{\glsni{thefigure}(\glsni{emph}\marg{\glsni{hashchar}1})}
\end{code}
For \glsni{subref}, you will have to do, for example, the following
in the text:
\begin{codeS}
(\glsni{emph}\marg{\glsni{subref}\marg{fig:circle}})
\end{codeS}

\begin{exercise}{Creating Sub-Figures}{ex:subfloats}

Download the image files \htmladdnormallink{\texttt{rectangle.pdf}}{exercises/rectangle.pdf}
and \htmladdnormallink{\texttt{circle.pdf}}{exercises/circle.pdf}
\latex{from \url{\exerciseurl/}} (or create your own images) and 
add \figureref{fig:shapes2} to your document.  You can
\downloadorview{subfloats} the solution.
\end{exercise}

%%%%%%%%%%%%%%%%%%% DEFINING COMMANDS %%%%%%%%%%%%%%%%%%%

\setnode{newcom}
\chapter{Defining Commands}
\label{ch:newcom}

It's possible to define your own \glspl{command} or
redefine existing ones.  Be very careful about redefining existing
commands; don't redefine a command simply because you want to use the
name, only redefine it if you are making a modification.  For
example, if you want to change the format of the current date, you
would redefine \gls{today}, but if you want to define a command
to display a specific date, you should define a new command with a
different name.

There are several reasons why you might want to define
a new command:
\begin{enumerate}

\item Reduce typing:

Suppose you have a series of commands or text that
you find yourself frequently using, then you could define
a command to do all these other commands for you.

\xminisec{Example:}  

Suppose you want a lot of large bold slanted sans-serif
portions of text within your document.  Every time you type those
portions of text, you will have to do something like:
\begin{codeS}
\glsni{textsf}\marg{\glsni{large}\glsni{bfseries}\glsni{slshape} Some text}
\end{codeS}%
It would be much easier if you could use just one command
to do all that, called, say, \cmdname{largeboldsfsl}:
\begin{codeS}
\begin{verbatim}
\largeboldsfsl{Some text}
\end{verbatim}
\end{codeS}%
or you could call it, say, \cmdname{lbsfsl} which is shorter,
but slightly less memorable:
\begin{codeS}
\begin{verbatim}
\lbsfsl{Some text}
\end{verbatim}
\end{codeS}

\item Ensure consistency:

You may find that you want to format an object a certain
way.

\xminisec{Example:}

Recall near the end of \sectionref{sec:subfloats}, I~suggested the
following to reference a subfigure (when using \gls{subref} instead
of \gls{ref}):
\begin{codeS}
(\glsni{emph}\marg{\glsni{subref}\marg{fig:circle}})
\end{codeS}
For consistency, you might want to define a command, say,
\begin{alltt}
\cmdname{formattedsubref}\marg{\meta{label}}
\end{alltt}
that was the same as 
(\glsni{emph}\marg{\glsni{subref}\marg{\meta{label}}}).

\xminisec{Another Example}

\refstepcounter{object}\label{example:keywords}%
Suppose your document has a lot of keywords
in it, and you want to format these keywords in a different font,
say sans-serif, so that they stand out.  You could just
do:
\begin{codeS}
A \glsni{textsf}\marg{command} usually begins with a backslash.
\end{codeS}%
however, it is better to define a new command called, say,
\cmdname{keyword} that will typeset its argument in a
sans-serif font.  That way it becomes a lot easier
to change the format at some later date.  For example,
you may decide to splash out and have your keywords
typed in a particular colour.  In which case, all you
need to do is simply change the definition of the
command \cmdname{keyword}, otherwise you'll have to 
go through your entire document looking for keywords
and changing each one which could be very time consuming if
you have a large document.  You might also decide at some
later date to make an index for your document.  Indexing
all the keywords then becomes much simpler, as again all you'll
need to do is modify the \cmdname{keyword} command.

\end{enumerate}

New commands are defined using the command:
\begin{definition}
\gls{newcommand}\marg{\meta{cmd}}\oarg{\meta{n-args}}\oarg{\meta{default}}\marg{\meta{text}}
\end{definition}%
The first \gls{mandatory} \meta{cmd} is
the name of your new command, which must start with a backslash.
The \gls{optional} \meta{n-args} specifies
how many arguments your new command must take.  The next
optional argument \meta{default} will be discussed later.
The final mandatory argument \meta{text} specifies what \LaTeX\ 
should do every time it encounters this command.

\xminisec{Example (No Parameters):}

Let's begin with a trivial example.  Suppose I~wanted to write a
document about a particular course, say \dq{Programming \textemdash\  Languages
and Software Construction}, and I~had to keep writing the course
title, then I~might decide to define a command that prints the course
title rather than having to laboriously type it out every time.
Let's call our new command \cmdname{coursetitle}.  We want the
following code:
\begin{codeS}
The course \glsni{emph}\marg{\cmdname{coursetitle}} is an undergraduate course.
\end{codeS}%
to produce the following output:
\begin{result}[coursetitle.html]
The course \emph{Programming \textemdash\ Languages and Software Construction}
is an undergraduate course.
\end{result}%
Clearly this command doesn't need any arguments, so we don't need to
worry about the optional argument \meta{n-args} to
\gls{newcommand}, and the only thing our new command needs to do
is print:
\begin{alltt}
Programming \glsni{emdash} Languages and Software Construction
\end{alltt}
so we would define our new command as follows:
\begin{code}
\glsni{newcommand}\marg{\cmdname{coursetitle}}\marg{Programming \glsni{emdash} Languages and Software Construction}
\end{code}%
Commands must always be defined before they are used.  The
best place to define commands is in the 
\gls{preamble}:
\begin{code}
\begin{alltt}
\glsni{documentclass}\marg{scrartcl}

\glsni{newcommand}\marg{\cmdname{coursetitle}}\marg{Programming \glsni{emdash} Languages 
and Software Construction}

\glsni{begin}\marg{document}

\glsni{section}\marg{\cmdname{coursetitle}}

The course \glsni{emph}\marg{\cmdname{coursetitle}} is an undergraduate course.

\glsni{end}\marg{document}
\end{alltt}
\end{code}

\xminisec{Example (One Mandatory Argument):}

\refstepcounter{object}\label{sec:newcomarg}%
Now let's try defining a command that takes an \gls{argument} (or
parameter).  Let's go back to our \htmlref{\cmdname{keyword}
example}{example:keywords}\latex{ \vpageref{example:keywords}}.  This
command needs to take one argument that is the keyword.  Let's
suppose we want keywords to come out in
\htmlref{sans-serif}{sec:fontstyle}, then we could do:
\begin{codeS}
\glsni{newcommand}\marg{\cmdname{keyword}}\oarg{1}\marg{\glsni{textsf}\marg{\glsni{hashchar}1}}
\end{codeS}%
In this case we have used the optional argument \meta{n-args} to
\gls{newcommand}.  We want our command \cmdname{keyword} to have
one argument, so we have \texttt{[1]}.  In \texttt{\glsni{textsf}\marg{\glsni{hashchar}1}} the
\gls{hashchar}\texttt{1} represents the first argument.  (If we had
more than one argument, \glsni{hashchar}\texttt{2} would represent the second
argument, \glsni{hashchar}\texttt{3} would represent the third argument etc.\ up to a
maximum of~9\faq{How to break the 9-argument limit}{moren9}.) So
\begin{verbatim}
\keyword{commands}
\end{verbatim}
will be equivalent to 
\begin{alltt}
\glsni{textsf}\marg{commands}
\end{alltt}
and
\begin{verbatim}
\keyword{environment}
\end{verbatim}
will be equivalent to
\begin{alltt}
\glsni{textsf}\marg{environment}
\end{alltt}
and so on.
Again, it's best to put the command definition in the preamble to
ensure the command won't be used before it's defined.

\begin{code}
\begin{alltt}
\glsni{documentclass}\marg{scrartcl}

\glsni{newcommand}\marg{\cmdname{keyword}}\oarg{1}\marg{\glsni{textsf}\marg{\glsni{hashchar}1}}

\glsni{begin}\marg{document}

A \cmdname{keyword}\marg{command} usually begins with a backslash.

\glsni{end}\marg{document}
\end{alltt}
\end{code}

Now if we want to change the way the keywords are formatted,
we can simply change the definition of \cmdname{keyword}.
Let's modify our code so that the keyword is now in a 
slanted sans-serif font:
\begin{code}
\begin{alltt}
\glsni{documentclass}\marg{scrartcl}

\glsni{newcommand}\marg{\cmdname{keyword}}\oarg{1}\marg{\glsni{textsf}\marg{\glsni{slshape} \glsni{hashchar}1}}

\glsni{begin}\marg{document}

A \cmdname{keyword}\marg{command} usually begins with a backslash.

\glsni{end}\marg{document}
\end{alltt}
\end{code}%
Let's go one stage further.  The \isty{color} 
\htmlref{package}{sec:packages} provides the \gls{declaration}:
\begin{definition}
\gls{color}\marg{\meta{col-name}}
\end{definition}
which switches the foreground colour to \meta{col-name}. It also
provides the text-block command:
\begin{definition}
\gls{textcolor}\marg{\meta{col-name}}\marg{\meta{text}}
\end{definition}
which sets \meta{text} in the colour given by \meta{col-name}.

So let's use the \isty{color} package to make our keywords
blue:\screenpagebreak
\begin{code}
\begin{alltt}
\glsni{documentclass}\marg{scrartcl}

\glsni{usepackage}\marg{color}

\glsni{newcommand}\marg{\cmdname{keyword}}\oarg{1}\marg{\glsni{textsf}\marg{\glsni{slshape}\glsni{color}\marg{blue}\glsni{hashchar}1}}

\glsni{begin}\marg{document}

A \cmdname{keyword}\marg{command} usually begins with a backslash.

\glsni{end}\marg{document}
\end{alltt}
\end{code}
Or we could index the keywords.  To do this we need the
\isty{makeidx} \htmlref{package}{sec:packages} and the commands 
\gls{makeindex}, 
\begin{inlinedef}\gls{index}\marg{\meta{text}}\end{inlinedef}
and \gls{printindex}:
\begin{code}
\begin{alltt}
\glsni{documentclass}\marg{scrartcl}

\glsni{usepackage}\marg{makeidx}

\glsni{makeindex}

\glsni{newcommand}\marg{\cmdname{keyword}}\oarg{1}\marg{\glsni{textsf}\marg{\glsni{slshape} \glsni{hashchar}1}\glsni{index}\marg{\glsni{hashchar}1}}

\glsni{begin}\marg{document}

A \cmdname{keyword}\marg{command} usually begins with a backslash.

\glsni{printindex}

\glsni{end}\marg{document}
\end{alltt}
\end{code}
For further information about how to create an index, 
see \latexguide\ or \latexcomp.
Alternatively, if you want a brief overview, try \latexthesis.

Since it is unlikely that the keyword will contain a paragraph break,
we should indicate that this is a 
\gls{short} using the \htmlref{starred form}{itm:starredcommand}:
\begin{codeS}
\glsni{newcommand}*\marg{\cmdname{keyword}}\oarg{1}\marg{\glsni{textsf}\marg{\glsni{slshape} \glsni{hashchar}1}\glsni{index}\marg{\glsni{hashchar}1}}
\end{codeS}
Now if you forget to add the closing brace, for example,
\verb|\keyword{command|, then \gls{tex}['s] error checking mechanism will 
pick up the error sooner. This will give an error message that
looks like:
\begin{verbatim}
! Paragraph ended before \keyword was complete.
<to be read again>
                   \par
l.604
\end{verbatim}
This at least gives you the line number (604 in this example) of the
end of the paragraph where the error has occurred.

If you don't used the starred form of \gls{newcommand}, then
you will get the somewhat less than helpful error:
\begin{verbatim}
! File ended while scanning use of \keyword.
\end{verbatim}
If you have a very large document, it may take a while to track
down where exactly you have missed a brace.

\begin{exercise}{Defining a New Command}{ex:newcom}
Try typing up the following code:
\begin{bcode}
\begin{alltt}
\glsni{documentclass}\marg{scrartcl}

\glsni{newcommand}*\marg{\cmdname{keyword}}\oarg{1}\marg{\glsni{textsf}\marg{\glsni{hashchar}1}}

\glsni{begin}\marg{document}

A \cmdname{keyword}\marg{command} usually begins with a backslash.

Segments of code may be \cmdname{keyword}\marg{grouped}.

Some \cmdname{keyword}\marg{commands} take one or more \cmdname{keyword}\marg{arguments}.
\glsni{end}\marg{document}
\end{alltt}
\end{bcode}

Then modify your code so that the keywords are in a slanted
sans-serif font or modify your code so that the keywords come out in
blue (using the \isty{color} package as in the example earlier).
Again you can \downloadorview{newcom} the result.

\xminisec{For the more adventurous:}

If you want to create an index as in the previous example, you will
need to use the application \iappname{makeindex}. If you used
\iappname{latexmk} back in \sectionref{sec:crossref}, you can just
carry on using that as before. If not you need to do the following
in TeXworks:
\begin{enumerate}
  \item Create the PDF as described in \sectionref{sec:texworks}.

  \item Select \menu{MakeIndex} from the drop-down list next to the
   build (typeset) button (see \figureref{fig:texworks9}).

  \item Click on the build button. If all goes well, you won't see
    anything different. If you see something like the following:
\begin{alltt}
Couldn't find input index file exercise\theexercise\space nor exercise\theexercise.idx.
\end{alltt}
    then you probably forgot to add the command \gls{makeindex} to
    the preamble. Add it in and go back to Step~1.

  \item Select \menu{pdfLaTeX} from the drop-down list and build the
    PDF file again. Move to the last page of the PDF, and you should
    see the index.
\end{enumerate}

\begin{figure}[htbp]
\figconts
  {pictures/texworks9}
  {%
    \caption{Selecting MakeIndex in TeXWorks}
  }%
  {fig:texworks9}
\end{figure}

\end{exercise}

%%%%%%%%%%%%% Defining Commands with an Optional Argument %%%%%%%%%%%%%

\setnode{newcomopt}
\section{Defining Commands with an Optional Argument}
\label{sec:newcomopt}

As mentioned earlier, the \gls{newcommand} command has a second
optional argument \meta{default}.  This allows you to define a
command with an optional argument\faq{More than one optional
argument}{twooptarg}.  For example, suppose we want a command called,
say, \cmdname{price}.  Suppose we want the following code:
\begin{codeS}
\begin{verbatim}
\price{100}
\end{verbatim}
\end{codeS}%
to produce the following output:
\begin{resultS}[price.html]
\pounds100 excl VAT @ 17.5\%\relax
\end{resultS}%
and let's suppose we want an optional argument
so that we can change the VAT.  That is, we
would want the following code:
\begin{codeS}
\begin{verbatim}
\price[20]{30}
\end{verbatim}
\end{codeS}%
to produce the following output:
\begin{resultS}[price2.html]
\pounds30 excl VAT @ 20\%
\end{resultS}%
Therefore we want to define a command
such that if the optional argument is absent 
we will have \texttt{17.5}, and if it is present
the optional argument will be substituted instead.
This command can be defined as follows:
\begin{codeS}
\gls{newcommand}\marg{\cmdname{price}}\oarg{2}\oarg{17.5}\marg{\glsni{pounds} \gls{hashchar}2 excl VAT @ \glsni{hashchar}1\glsni{percent}}
\end{codeS}%
Here, \texttt{\glsni{hashchar}1} represents the optional argument (by default
\texttt{17.5}) and \texttt{\glsni{hashchar}2} represents the mandatory argument (the
second argument if the optional argument is present, or the only
argument if the optional argument is absent.)

As before, since the argument is unlikely to contain a paragraph
break, we should indicate that it is a \gls{short} using the \htmlref{starred form}{itm:starredcommand}:
\begin{codeS}
\gls{newcommand}*\marg{\cmdname{price}}\oarg{2}\oarg{17.5}\marg{\glsni{pounds} \glsni{hashchar}2 excl VAT @ \glsni{hashchar}1\gls{percent}}
\end{codeS}%

\begin{exercise}{Defining Commands with an Optional Argument}{ex:newcomopt}
In this exercise, you will need to define a slightly modified
version of the above example.  Try defining a command called,
say, \cmdname{cost}.  It should take one optional argument and
one mandatory argument.  Without the optional argument, it
behaves in the same way as the \cmdname{price} example above,
so that, say,
\begin{codeS}
\begin{verbatim}
\cost{50}
\end{verbatim}
\end{codeS}%
will produce
\begin{resultS}[price3.html]
\pounds50 excl VAT @ 17.5\%\null
\end{resultS}%
but with the optional argument, you can change the
\verb|excl VAT @ 17.5\%| bit.  So that, say,
\begin{codeS}
\begin{verbatim}
\cost[inc VAT]{50}
\end{verbatim}
\end{codeS}%
will produce
\begin{resultS}[price4.html]
\pounds50 inc VAT\mbox{}% LaTeX2HTML is getting confused!
\end{resultS}

You can \downloadorview{newcomopt} the solution.

\xminisec{For the more adventurous:}
If you did \exerciseref{ex:newcom} and you modified \cmdname{keyword}
so that it indexed the keyword, you may have noticed that
\begin{codeS}
\verb|\keyword{command}|
\end{codeS}
and
\begin{codeS}
\verb|\keyword{commands}|
\end{codeS}
produced separate entries in the index. It would be better to have an
optional argument to override the indexing mechanism. For example,
\begin{codeS}
\verb|\keyword{command}|
\end{codeS}
should print and index the word \dq{command},
whereas
\begin{codeS}
\verb|\keyword[command]{commands}|
\end{codeS}
should print \dq{commands}
and index \dq{command}. In other words, we need an optional
argument that defaults to the mandatory argument if it is not
present.
This is how to achieve that type of effect:\footnote{Recall from \chapterref{ch:def} the percent symbol discards the space resulting from the end
of line character.}
\begin{code}
\glsni{newcommand}*\marg{\cmdname{keyword}}\oarg{2}\oarg{\cmdname{keywordentry}}\marg{\glsni{percentchar}\newline
\strut~~\icmdname{def}\cmdname{keywordentry}\marg{\glsni{hashchar}2}\glsni{percentchar}\newline
\strut~~\glsni{textsf}\marg{\glsni{hashchar}2}\glsni{percentchar}\newline
\strut~~\gls{index}\marg{\glsni{hashchar}1}\glsni{percentchar}\newline
}
\end{code}
In this example, the default value for the optional argument
is set to the command \cmdname{keywordentry}. At the start of 
\cmdname{keyword} this is defined to be the mandatory argument 
(as specified by \verb|#2|) using \gls{tex}['s] \cmdname{def}
command:\footnote{\cmdname{def} is too complicated for an
introductory \LaTeX\ guide but, if you're interested, read \emph{The
\TeX{}book}~\cite{texbook}.}
\begin{alltt}
\cmdname{def}\cmdname{keywordentry}\marg{\glsni{hashchar}2}
\end{alltt}
Then typeset the keyword (given in the mandatory
argument \texttt{\glsni{hashchar}2}) in a sans-serif font:
\begin{alltt}
\glsni{textsf}\marg{\glsni{hashchar}2}
\end{alltt}
Now index the term using the optional argument (\texttt{\glsni{hashchar}1}):
\begin{alltt}
\glsni{index}\marg{\glsni{hashchar}1}
\end{alltt}
If an optional argument is specified, \texttt{\glsni{hashchar}1} will be the given
argument, but if the optional argument is missing, \texttt{\glsni{hashchar}1} will
be \cmdname{keywordentry}, which has earlier been set to
the mandatory argument \texttt{\glsni{hashchar}2}.
\end{exercise}

%%%%%%%%%%%%%%%%%%%%%%%%% Redefining Commands %%%%%%%%%%%%%%%%%%%%%

\setnode{renewcom}
\section{Redefining Commands}
\label{sec:renewcom}

Commands can be redefined using the command:
\begin{definition}
\gls{renewcommand}\marg{\meta{cmd}}\oarg{\meta{n-args}}\oarg{\meta{default}}\marg{\meta{text}}
\end{definition}\reportpagebreak\noindent
This has exactly the same format as \gls{newcommand} but is
used for redefining existing commands.
Again there is a \htmlref{starred version}{itm:starredcommand} to indicate that the command
is a \gls{short}.

\xminisec{Caveat:}

\warning Never redefine a command whose existing function is unknown
to you or just because you want to use a particular command name,
regardless of its previous function. By way of illustration: as a
production editor, I~have to combine articles by different authors
into a single book. Each author supplies the \LaTeX\ code for their
own article. Every so often, I~get code that redefines a command for
the convenience of the author. Later on another author tries to use
the same command, on the assumption that the command behaves
according to its original definition. This tends to involve the
accent commands as they are short and that saves the author typing.
It usually goes along these lines: author~A redefines \gls{c} (the
cedilla accent command) to display a maths bold
\dq{\boldc} to indicate a
vector. Later, author~B, uses the cedilla accent, say, in the name
Fran\c{c}ois:
\begin{alltt}
Fran\gls{c}\marg{c}ois
\end{alltt}
Author~A's hack turns this into
Fran\boldc cois.

\xminisec{Example (Redefining List Labels):}

Recall the \gls{env-itemize} environment discussed in
\sectionref{sec:itemize}.  You may have up to four nested
\gls{env-itemize} environments, the labels for the outer environment
are specified by the command \gls{labelitemi}, the labels for
the second level are specified by \gls{labelitemii}, the third
by \gls{labelitemiii} and the fourth by \gls{labelitemiv}.
By default, \glsni{labelitemi} is a bullet point (\makeimg{bullet
symbol}{\labelitemi}),
\glsni{labelitemii} is an en dash (\makeimg{en dash}{\labelitemii}), 
\glsni{labelitemiii} is an
asterisk (\makeimg{asterisk}{\labelitemiii}) and 
\glsni{labelitemiv} is a centred dot (\makeimg{centered
dot}{\labelitemiv}).  These can be changed by redefining
\glsni{labelitemi} etc.

Recall from \tableref{tab:symbols} that the command
\gls{dag} produces a dagger symbol, we can use this symbol
instead of a bullet point:
\begin{code}
\glsni{renewcommand}*\marg{\glsni{labelitemi}}\marg{\gls{dag}}\newline
\glsni{begin}\marg{itemize}\newline
\mbox{}\newline
\gls{item} Animal\newline
\mbox{}\newline
\glsni{item} Mineral\newline
\mbox{}\newline
\glsni{item} Vegetable\newline
\mbox{}\newline
\glsni{end}\marg{itemize}
\end{code}%
Output:
\begin{result}[List uses dagger symbol instead of a bullet point for
the marker]
\renewcommand*{\labelitemi}{\dag}

\begin{itemize}

\item Animal

\item Mineral

\item Vegetable

\end{itemize}
\end{result}

Here's another example, it uses the \Index{PostScript}
font ZapfDingbats via the \isty{pifont} package:
\begin{code}
\glsni{renewcommand}*\marg{\glsni{labelitemi}}\marg{\gls{ding}\marg{43}}\newline
\glsni{begin}\marg{itemize}\newline
\mbox{}\newline
\glsni{item} Animal\newline
\mbox{}\newline
\glsni{item} Mineral\newline
\mbox{}\newline
\glsni{item} Vegetable\newline
\mbox{}\newline
\glsni{end}\marg{itemize}
\end{code}\bookpagebreak\noindent
Output:
\begin{result}[List uses a pointing finger instead of a bullet point
for the marker]
\renewcommand*{\labelitemi}{\ding{43}}

\begin{itemize}

\item Animal

\item Mineral

\item Vegetable

\end{itemize}
\end{result}%
In the above example, it would actually be easier to use the
\gls{env-dinglist} environment defined in the \isty{pifont} package:
\begin{code}
\begin{alltt}
\glsni{begin}\marg{dinglist}\marg{43}

\glsni{item} Animal

\glsni{item} Mineral

\glsni{item} Vegetable\screenpagebreak

\glsni{end}\marg{dinglist}
\end{alltt}
\end{code}


\xminisec{Example (Redefining the Default Font):}

Recall from \sectionref{sec:changingfonts} that the default font
family is usually the serif (Roman) family. So what happens if you
need to write your entire document in, say, Helvetica?  The default
font family name is stored in:
\begin{definition}
\gls{familydefault}
\end{definition}
This command is usually
defined to be just \gls{rmdefault}, which in turn stores the name of the
default serif font (initially \texttt{cmr}, Computer Modern Roman). If 
you want the default font to be sans-serif, all you
need do is add the following line to the
\glsterm{preamble}:
\begin{codeS}
\glsni{renewcommand}\marg{\glsni{familydefault}}\marg{\gls{sfdefault}}
\end{codeS}
\glsni{sfdefault} stores the name of the default sans-serif font (initially \texttt{cmss},
Computer Modern Sans-Serif) and
the \isty{helvet} package redefines \glsni{sfdefault} to \texttt{phv},
which is the identifier for the Helvetica font.
So the following document will be in Helvetica:
\begin{alltt}
\glsni{documentclass}\marg{scrartcl}

\glsni{usepackage}\marg{helvet}
\glsni{renewcommand}\marg{\glsni{familydefault}}\marg{\glsni{sfdefault}}

\glsni{begin}\marg{document}
This is a sample document.
\glsni{end}\marg{document}
\end{alltt}
Similarly, if you want the default font to be monospaced (typewriter) then you'd need to do:
\begin{codeS}
\glsni{renewcommand}\marg{\glsni{familydefault}}\marg{\gls{ttdefault}}
\end{codeS}

Incidentally, you may have noticed in \sectionref{sec:changingfonts}
that although I~said I'd used the \isty{anttor} and
\isty{libris} packages to set the serif and sans-serif families for
\latexhtml{this book}{the PDF version of this book}, I
didn't mention anything about the typewriter (monospaced) font. I~used
the TXTT font, but that doesn't have a corresponding package. You
just redefine \glsni{ttdefault} to \texttt{txtt}:
\begin{codeS}
\glsni{renewcommand}*\marg{\glsni{ttdefault}}\marg{txtt}
\end{codeS}

\xminisec{Example (Redefining Fixed Names):}

You may have noticed that \LaTeX\ automatically generates pieces of
text such as \dq{Chapter}, \dq{Figure}, \dq{Bibliography}.  These are
generated by the commands shown in \tableref{tab:textlab}.

\begin{table}[hbtp]
\caption[Predefined Names]{Predefined Names 
(\supdag Book and report style classes (such as \icls{scrreprt} and \icls{scrbook}), 
\supddag article-style classes (such as \icls{scrartcl}), 
remainder book, report and article-style classes)}
\label{tab:textlab}
\centering
\begin{tabular}{ll}
\toprule
\bfseries Command & \bfseries Default Text\\
\midrule
\gls{contentsname} & \contentsname\\
\gls{listfigurename} & \listfigurename\\
\gls{listtablename} & \listtablename\\
\gls{bibname}\supdag & \bibname\\
\gls{refname}\supddag & \refname\\
\gls{indexname} & \indexname\\
\gls{figurename} & \figurename\\
\gls{tablename} & \tablename\\
\gls{partname} & \partname\\
\gls{chaptername}\supdag & \chaptername\\
\gls{appendixname} & \appendixname\\
\gls{abstractname} & \abstractname
\\\bottomrule
\end{tabular}
\end{table}

You can change the defaults using \gls{renewcommand}.  For
example, suppose you want the table of contents to be labelled
\dq{Table of Contents}, instead of the default \dq{Contents}, you
would need to do:
\begin{codeS}
\glsni{renewcommand}*\marg{\glsni{contentsname}}\marg{Table of Contents}
\end{codeS}
\faq{Changing the words babel uses}{latexwords}In fact, the
\isty{babel} package (see \sectionref{sec:babel}) uses this method
to redefine the commands in \tableref{tab:textlab} whenever you
switch language using \reportlinebreak\screenlinebreak
\gls{selectlanguage} or within the contents of
the \gls{env-otherlanguage} environment. This unfortunately has the
side-effect that means if you try to redefine these commands,
\sty{babel} will automatically overwrite your definition whenever
there's a language change, which includes at the beginning of the
\gls{env-document} environment. Instead you need to use
\sty{babel}'s \glsni{addto} mechanism.
\begin{definition}
\gls{addto}\marg{\meta{command}}\marg{\meta{code}}
\end{definition}
This patches the definition of a command (specified in the first
\gls{argument}) adding \meta{code} to the end of the command
definition. Whenever \sty{babel} switches the current language, it
uses the command \cmdname{captions}\meta{language}, which performs
all the redefinitions of commands like those listed in
\tableref{tab:textlab}. For example, if you are using \sty{babel}
with the \istyopt{babel}{english} option and you want to change
\glsni{contentsname} so that it does \dq{Table of Contents} instead
of \dq{Contents}, you need to do:
\begin{code}
\begin{alltt}
\glsni{addto}\marg{\cmdname{captionsenglish}}\marg{\glsni{percentchar}
  \glsni{renewcommand}\marg{\glsni{contentsname}}\marg{Table of Contents}\glsni{percentchar}
}
\end{alltt}
\end{code}

\xminisec{Notes:}

\warning Take care if you want to patch an existing command. For
example, suppose you want to append something to the action of a
command. You might be tempted to do
\begin{alltt}
\wrong\glsni{renewcommand}*\marg{\cmdname{foo}}\marg{\cmdname{foo} Something else}
\end{alltt}
This will cause an infinite loop where \cmdname{foo} recursively
calls itself. Instead you should use one of the commands provided by
the \isty{etoolbox} package (such as \gls{appto}, which has the same
syntax as \sty{babel}'s \glsni{addto} described above). For further details,
read the \sty{etoolbox} \htmlref{documentation}{sec:texdoc}.

\begin{exercise}{Renewing Commands}{ex:renewcom}
If you did Exercises~\ref{ex:figures} and \ref{ex:tables}, go back to
that document and changed the figures and tables so that they
are labelled \dq{Fig} and \dq{Tab} instead of \dq{Figure} and \dq{Table}.
Hint: you need to redefine \glsni{tablename} and \glsni{figurename}.

You can \downloadorview{renewcom} the solution.
\end{exercise}

%%%%%%%%%%%%%%%%%%%%%%%% MATHEMATICS %%%%%%%%%%%%%%%%%%%%%%%%%%%%%%%

\setnode{mathematics}
\chapter{Mathematics}
\label{ch:maths}

As mentioned in the \htmlref{introduction}{ch:intro}, \LaTeX\ is
particularly good at typesetting mathematics.  In order to
use any of the maths commands we need to be in one of the
mathematics \glspl{environment}.  There are two
basic types of mathematics: \keyword{in-line maths} and
\keyword{displayed maths}. In-line maths is mathematics that
occurs within a line of text, for example:
\begin{resultS}[inlinemaths.html]
The variable $x$ is transformed by the function $f(x)$.
\end{resultS}%
Displayed maths is mathematics that occurs on a line
of its own.  For example:
\begin{result}[displaymath.html]
A polynomial is a function of the form
\[
f(x) = \sum_{i=0}^na_ix^i
\]
\end{result}
The maths environments switch to \LaTeX's \dq{math mode}, which uses
specialist maths fonts and spacing rather than just using an italic
font.

If you want to typeset any mathematics, I~strongly advise using the
\istystart{amsmath} package:
\begin{alltt}
\glsni{usepackage}\marg{amsmath}
\end{alltt}
This patches some existing \LaTeX\ commands and environments and
also provides many useful additions.

This chapter is just an introduction to typesetting mathematics in
\LaTeX. If you want a comprehensive guide, I~recommend you read \emph{Math Mode}
by Herbert~Vo\ss~\cite{voss}, which can be access via
\appname{texdoc} (see \sectionref{sec:texdoc}):
\begin{verbatim}
texdoc mathmode
\end{verbatim}

%%%%%%%%%%%%%%%% In-Line Mathematics %%%%%%%%%%%%%%%%%%%

\setnode{inlinemaths}
\section{In-Line Mathematics}
\label{sec:inline}

In-line mathematics is created using the \gls{env-math}
environment. (Note U.S. spelling --- \dq{math} not \dq{maths}).

\xminisec{Example:}
\begin{code}
The variable \glsni{begin}\marg{math}x\glsni{end}\marg{math} is transformed
by the function \glsni{begin}\marg{math}f(x)\glsni{end}\marg{math}.
\end{code}
It's somewhat cumbersome having to type
\glsni{begin}\marg{math} and \glsni{end}\marg{math} and it also
makes the \gls{source} a little difficult to read so there
are shorthand notations that can be used instead:
\gls{beginmath} is equivalent to \glsni{begin}\marg{math} and
\gls{endmath} is equivalent to \glsni{end}\marg{math}.  So
the example above can be rewritten:
\begin{codeS}
The variable \glsni{beginmath}x\glsni{endmath} is transformed by the function 
\glsni{beginmath}f(x)\glsni{endmath}.
\end{codeS}%
There is an even shorter notation:
The \htmlref{special character}{sec:chars} \gls{dollarchar} is equivalent
to both \glsni{begin}\marg{math} and \glsni{end}\marg{math}:
\begin{codeS}
The variable \glsni{dollarchar}x\glsni{dollarchar} is transformed by the 
function \glsni{dollarchar}f(x)\glsni{dollarchar}.
\end{codeS}%
This is considerably easier to type and
to read, but you need to make sure that all your \glsni{dollarchar}
symbols have matching pairs.  The above code will look like:
\begin{resultS}[inlinemaths2.html]
The variable $x$ is transformed
by the function $f(x)$.
\end{resultS}%
The other advantage in using \glsni{dollarchar} over \glsni{beginmath} and
\glsni{endmath} is that \glsni{dollarchar} is a \gls{robust}, 
whereas \glsni{beginmath} and \glsni{endmath} are 
\glspl{fragile} and
will need to be protected if they occur in a \glslink{fragile}{moving
argument}.\casemedia{}{}{ (See \sectionref{sec:fragile}.)}{}

Note: you should always make sure you are in maths mode to
typeset any variables (such as $x$, $y$, $z$), as this will
ensure that the correct maths fonts are used, as well as the
appropriate spacing. Similarly, don't use \glsni{dollarchar}
as a short cut for an italic font.
\begin{codeS}[\linewidth]
Notice the \glsni{dollarchar}difference\glsni{dollarchar} between 
\glsni{dollarchar}(x', y', z')\glsni{dollarchar} and
\glsni{textit}\marg{(x', y', z')}.
\end{codeS}%
\begin{resultS}[Image of result: the brackets are upright in the
maths font but slanted in the italic font. The apostrophe becomes a
prime in the maths font. The spacing is also
different.]
Notice the $difference$ between $(x', y', z')$ and
\textit{(x', y', z')}.
\end{resultS}

%%%%%%%%%%%%%%%% Displayed Mathematics %%%%%%%%%%%%%%%%%%%%%%

\setnode{displayedmaths}
\section{Displayed Mathematics}
\label{sec:displayed}

One-line unnumbered displayed mathematics can be created using:
\begin{definition}
\gls{begindispmath}\meta{maths}\gls{enddispmath}
\end{definition}
where \meta{maths} is the mathematics to be displayed.

\xminisec{Example:}
\begin{code}
\begin{alltt}
A linear function is a function of the form
\glsni{begindispmath} y = mx + c \glsni{enddispmath}
\end{alltt}
\end{code}%
Output:
\begin{result}[displaymath2.html]
A linear function is a function of the form
\[ y = mx + c \]
\end{result}%

\warning Don't use the \envname{displaymath} environment or \verb|$$|\ldots\verb|$$|~\cite{l2tabu}.
Use \glsni{begindispmath} and \glsni{enddispmath} with
the \isty{amsmath} package.

The \glsni{env-equation}\refstepcounter{object}\label{env:equation} 
environment provides something similar to
\glsni{begindispmath} \glsni{enddispmath}, except
that the equation is numbered.  Modifying the above example:
\begin{code}
\begin{alltt}
A linear function is a function of the form
\glsni{begin}\marg{equation}
y = mx + c
\glsni{end}\marg{equation}
\end{alltt}
\end{code}%
results in the following output:
\begin{result}[As before but the equation now has a number in
brackets on the right.]
A linear function is a function of the form
\begin{equation}
y = mx + c
\end{equation}
\end{result}
Normal text can be inserted into the equation using
\begin{definition}
\gls{text}\marg{\meta{text}}
\end{definition}
which is provided by the \isty{amsmath} package.

\xminisec{Example:}
\begin{codeS}
\glsni{begindispmath} x = 2 \glsni{text}\marg{ and } y = -1 \glsni{enddispmath}
\end{codeS}%
results in the following output:
\begin{result}[x equals 2 and y equals minus 1 (the word 'and' is in normal text font)]
\[ x = 2 \text{ and } y = -1 \]
\end{result}

\faq{Re-using an equation}{reuseq}%
Recall from \sectionref{sec:crossref} that we can
\htmlref{cross-reference}{sec:crossref} most things that \LaTeX\
automatically numbers using \gls{ref} and \gls{label}.
Equations can be cross-referenced in the same way:\reportpagebreak
\begin{code}
\begin{alltt}
Equation\glsni{tildechar}\glsni{ref}\marg{eqn:linear} is a linear function.
\glsni{begin}\marg{equation}
\glsni{label}\marg{eqn:linear}
f(x) = mx + c
\glsni{end}\marg{equation}\screenpagebreak
\end{alltt}
\end{code}
\begin{result}[eqref.html]
Equation~\ref*{eqn:linear} is a linear function.
\begin{equation}
\label{eqn:linear}
f(x) = mx + c
\end{equation}
\end{result}
Equation numbers are usually given in parentheses, which can be done
using:
\begin{codeS}
Equation\glsni{tildechar}(\glsni{ref}\marg{eqn:linear})
\end{codeS}
The \isty{amsmath} package provides a convenient short cut:
\begin{definition}
\gls{eqref}\marg{\meta{label}}
\end{definition}
So the above can be written as:
\begin{codeS}
Equation\glsni{tildechar}\glsni{eqref}\marg{eqn:linear}
\end{codeS}
\begin{resultS}[Equation (9.2)]
% fake it to suppress hyperlink
Equation~(\ref*{eqn:linear})
\end{resultS}

\xminisec{Note:}

Both the \gls{env-equation} environment and
\gls{begindispmath}\ldots\gls{enddispmath} are only designed for one
line of maths. Therefore you must not have any line breaks or
paragraph breaks within them.
The following will cause an error:
\begin{alltt}
\glsni{begin}\marg{equation}
\wrong
f(x) = mx + c
\wrong
\glsni{end}\marg{equation}
\end{alltt}
Either remove the blank lines or comment them out:
\begin{alltt}
\glsni{begin}\marg{equation}
\glsni{percentchar}\correct
f(x) = mx + c
\glsni{percentchar}\correct
\glsni{end}\marg{equation}
\end{alltt}

\setnode{align}
\section{Multiple Lines of Displayed Maths}
\label{sec:align}

The \isty{amsmath} package provides the \gls{env-align}
and \gls{env-align*} environments for aligned equations.
The starred version doesn't number the equations. These environments 
provide pairs of left- and right-aligned columns.
As with the \gls{env-tabular} environment, use \gls{ampchar} to
separate columns and \gls{tab.dbbackslashchar} to separate rows.
Unlike the \gls{env-tabular} environment, there is no argument as
the column specifiers are predefined. Another difference is that no
page breaks can occur in the \gls{env-tabular} environment, but it's
possible to allow a page break in \glsni{env-align} or
\glsni{env-align*} using
\begin{definition}
\gls{displaybreak}\oarg{\meta{n}}
\end{definition}
immediately before the \gls{tab.dbbackslashchar} where it is to take
effect. The optional argument is a number from~0 to~4 indicating the
desirability to break the page (from~0 the least to~4 the most).

If you want to mix numbered and unnumbered rows, you can use
\begin{definition}
\gls{notag}
\end{definition}
to suppress the numbering for a particular row in the \glsni{env-align} environment.
This command must go before \gls{tab.dbbackslashchar} at the end of the
row. The default equation numbering can be overridden for a
particular row using:
\begin{definition}
\gls{tag}\marg{\meta{tag}}
\end{definition}
where \meta{tag} is the replacement for the equation number.

\screenpagebreak\warning Don't use the \envname{eqnarray} or \envname{eqnarray*}
environments. They're obsolete~\cite{l2tabu}.


\xminisec{Example (Unnumbered):}

\begin{code}
\glsni{begin}\marg{align*}\newline
y \glsni{ampchar}= 2x + 2\glsni{tab.dbbackslashchar}\newline
\strut~~\glsni{ampchar}= 2(x+1)\newline
\glsni{end}\marg{align*}
\end{code}
\begin{result}[align.html]
\begin{align*}
y &= 2x + 2\\
  &= 2(x+1)
\end{align*}
\end{result}
Note that the equals sign is placed at the start of the second
column, \emph{after} the ampersand \gls{ampchar}. This ensures the
correct amount of spacing on either side. If the first line of the
above equation was changed to:\bookpagebreak
\begin{alltt}
\wrong y =\gls{ampchar} 2x + 2\gls{tab.dbbackslashchar}
\end{alltt}
there wouldn't be enough space on the right of the equal sign:
\begin{htmlresult}{y equals 2x + 2}
\begin{makeimage}
\begin{align*}
y =& 2x + 2
\end{align*}
\end{makeimage}
\end{htmlresult}

\xminisec{Example (One Row Numbered):}

\begin{code}
\glsni{begin}\marg{align}\newline
y \glsni{ampchar}= 2x + 2\glsni{notag}\glsni{tab.dbbackslashchar}\newline
\strut~~\glsni{ampchar}= 2(x+1)\newline
\glsni{end}\marg{align}
\end{code}
\begin{result}[Image of result: as the earlier example but the second equation has been numbered]
\begin{align}
y &= 2x + 2\notag\\
  &= 2(x+1)
\end{align}
\end{result}\screenpagebreak

\xminisec{Example (Four Columns):}

\begin{code}
\glsni{begin}\marg{align*}\newline
y \glsni{ampchar}= 2x + 2 \glsni{ampchar} z 
 \glsni{ampchar}= 6x + 3\glsni{tab.dbbackslashchar}\newline
\strut~~\glsni{ampchar}= 2(x+1) \glsni{ampchar}~~~\glsni{ampchar}= 3(2x+1)\newline
\glsni{end}\marg{align*}
\end{code}
\begin{result}[4colalign.html]
\begin{align*}
y &= 2x + 2 & z &= 6x + 3\\
  &= 2(x+1) &  &= 3(2x+1)
\end{align*}
\end{result}

As with \gls{env-equation}, you can cross-reference individual rows
of an \glsni{env-align} environment, but you must remember to
put \gls{label} \emph{before} the end of row \gls{tab.dbbackslashchar}
separator. You can reference a row in the \glsni{env-align*}
environment if you have assigned it a tag with \gls{tag}, but don't
try labelling a row in the \glsni{env-align} environment where the
numbering has been suppressed with \gls{notag}.

\xminisec{Example (Cross-Referenced):}

This example has two numbered equations in an \gls{env-align} environment,
both of which are labelled and referenced:
\begin{code}
The function \glsni{dollarchar}f(x)\glsni{dollarchar} is given in 
Equation\glsni{tildechar}\glsni{eqref}\marg{eq:fx},
and its derivative \glsni{dollarchar}f'(x)\glsni{dollarchar} is given in 
Equation\glsni{tildechar}\glsni{eqref}\marg{eq:dfx}.
\newline
\glsni{begin}\marg{align}\newline
f(x) \glsni{ampchar}= 2x + 1
\glsni{label}\marg{eq:fx}\glsni{tab.dbbackslashchar}\newline
f'(x) \glsni{ampchar}= 2 \glsni{label}\marg{eq:dfx}\newline
\glsni{end}\marg{align}
\end{code}
\begin{result}[alignref.html]
The function $f(x)$ is given in Equation~(\ref*{eq:fx}), and its
derivative $f'(x)$ is given in Equation~(\ref*{eq:dfx}).
\begin{align}
f(x) &= 2x + 1\label{eq:fx}\\
f'(x) &= 2 \label{eq:dfx}
\end{align}
\end{result}

Recall the command \gls{text}\marg{\meta{text}} from the
\htmlref{previous section}{sec:displayed}. This can be used
within cells of the \glsni{env-align} and \glsni{env-align*}
environments, but the \isty{amsmath} package also provides
\begin{definition}
\gls{intertext}\marg{\meta{text}}
\end{definition}
which can be used for a line of interjection between the rows. This
command may only go right after \gls{tab.dbbackslashchar}.

\xminisec{Example}

\begin{code}
\glsni{begin}\marg{align*}\newline
y \glsni{ampchar}= 2x +
2\glsni{tab.dbbackslashchar}\newline
\glsni{intertext}\marg{Using the distributive law:}\newline
  \glsni{ampchar}= 2(x+1)\newline
\glsni{end}\marg{align*}
\end{code}
\begin{result}[intertext.html]
\begin{align*}
y &= 2x + 2\\
\intertext{Using the distributive law:}
  &= 2(x+1)
\end{align*}
\end{result}\bookpagebreak

There are other environments for multiple-line displayed maths, but
they are beyond the scope of this book. See
the \isty{amsmath} documentation for further details.

%%%%%%%%%%%%%%%% Mathematical Commands %%%%%%%%%%%%%%%%%%%%

\setnode{mathscom}
\section{Mathematical Commands}
\label{sec:mathscom}

Most of the \glspl{command} described in this
section may only be used in one of the mathematics environments.  If
you try to use a mathematics command outside a maths environment you
will get a \dq{\texttt{Missing \$ inserted}} error message.

%------------------ Maths Fonts ------------------

\setnode{mathsfonts}
\subsection{Maths Fonts}
\label{sec:mathsfonts}

Just as we are able to \htmlref{change text fonts}{sec:fonts} using
the commands \gls{textrm}, \gls{textbf} etc, we can also
use commands to change the maths font.  Basic maths font changing
commands are shown in \tableref{tab:mathsfonts}.

\begin{table}[tbhp]
\caption{Maths Font Changing Commands}
\label{tab:mathsfonts}
\centering
\begin{tabular}{lll}
\toprule
\bfseries Command & \bfseries Example Input & \bfseries Corresponding Output\\
 & & (Computer Modern)\\
\midrule
\begin{inlinedef}\gls{mathrm}\marg{\meta{maths}}\end{inlinedef} &
\verb|$\mathrm{xyz}$| & $\cmmathrm{xyz}$\\
\begin{inlinedef}\gls{mathsf}\marg{\meta{maths}}\end{inlinedef} &
\verb|$\mathsf{xyz}$| & $\cmmathsf{xyz}$\\
\begin{inlinedef}\gls{mathtt}\marg{\meta{maths}}\end{inlinedef} &
\verb|$\mathtt{xyz}$| & $\cmmathtt{xyz}$\\
\begin{inlinedef}\gls{mathit}\marg{\meta{maths}}\end{inlinedef} &
\verb|$\mathit{xyz}$| & $\cmmathit{xyz}$\\
\begin{inlinedef}\gls{mathbf}\marg{\meta{maths}}\end{inlinedef} &
\verb|$\mathbf{xyz}$| & $\cmmathbf{xyz}$\\
\begin{inlinedef}\gls{mathcal}\marg{\meta{maths}}\end{inlinedef} &
\verb|$\mathcal{XYZ}$| & $\cmmathcal{XYZ}$
\\\bottomrule
\end{tabular}
\end{table}

The calligraphic fonts via \glsni{mathcal} are only available for upper-case 
characters\faq{Better script fonts for maths}{scriptfonts}.
\tableref{tab:amsmathfont} lists additional font commands supplied 
with the \isty{amsmath} and \isty{amsfonts} 
\htmlref{packages}{sec:packages}.

\begin{table}[hbtp]
\caption[The amsfonts and amsmath Font Commands]{The 
\isty{amsfonts}\supddag\ and \isty{amsmath}\supdag\ Font Commands}
\label{tab:amsmathfont}
\centering
\begin{tabular}{lll}
\toprule
\bfseries Command & \bfseries Example Input & \bfseries Example Output\\
\midrule
\begin{inlinedef}\supddag\gls{mathbb}\marg{\meta{maths}}\end{inlinedef} &
\verb|$\mathbb{A+B=C}$| & $\mathbb{A+B=C}$\\
\begin{inlinedef}\supddag\gls{mathfrak}\marg{\meta{maths}}\end{inlinedef} &
\verb|$\mathfrak{A+B=C}$| & $\mathfrak{A+B=C}$\\
\begin{inlinedef}\supdag\gls{boldsymbol}\marg{\meta{maths}}\end{inlinedef} &
\verb|$\boldsymbol{A+B=C}$| & $\boldsymbol{A+B=C}$\\
\begin{inlinedef}\supdag\gls{pmb}\marg{\meta{symbol}}\end{inlinedef} &
\verb|$\pmb{+-=}$| & $\pmb{+-=}$
\\\bottomrule
\end{tabular}
\end{table}

%----------------- Greek Letters -----------------------

\setnode{greek}
\subsection{Greek Letters}
\label{sec:greek}

Greek letters that differ from the corresponding Roman letters are
obtained by placing a backslash in front of the name.\footnote{So,
for example, there is no omicron since it looks the same as a Roman
o.} Lower case and upper case Greek letters are shown in
\tableref{tab:greek} and \tableref{tab:Greek}, respectively. There are
also some variants of certain symbols, such as \glsni{vartheta} as
opposed to \glsni{theta}.

\begin{table}[htbp]
\caption{Lower Case Greek Letters}
\label{tab:greek}
\centering
\begin{tabular}{llllll}
\gls{alpha} & $\alpha$ &
\gls{beta} & $\beta$ &
\gls{gamma} & $\gamma$ \\
\gls{delta} & $\delta$  &
\gls{epsilon} & $\epsilon$ &
\gls{varepsilon} & $\varepsilon$\\
\gls{zeta} & $\zeta$ &
\gls{eta} & $\eta$ &
\gls{theta} & $\theta$ \\
\gls{vartheta} & $\vartheta$ &
\gls{iota} & $\iota$ & 
\gls{kappa} & $\kappa$ \\
\gls{lambda} & $\lambda$ &
\gls{mu} & $\mu$ &
\gls{nu} & $\nu$ \\
\gls{xi} & $\xi$ &
\gls{pi} & $\pi$ &
\gls{varpi} & $\varpi$\\
\gls{rho} & $\rho$ &
\gls{varrho} & $\varrho$ &
\gls{sigma} & $\sigma$ \\
\gls{varsigma} & $\varsigma$ &
\gls{tau} & $\tau$ &
\gls{upsilon} & $\upsilon$ \\
\gls{phi} & $\phi$ &
\gls{varphi} & $\varphi$ &
\gls{chi} & $\chi$ \\
\gls{psi} & $\psi$ &
\gls{omega} & $\omega$ 
\end{tabular}
\end{table}

\begin{table}[htbp]
\caption{Upper Case Greek Letters}
\label{tab:Greek}
\centering
\begin{tabular}{llllll}
\gls{Gamma} & $\Gamma$ &
\gls{Delta} & $\Delta$ &
\gls{Theta} & $\Theta$ \\
\gls{Lambda} & $\Lambda$ &
\gls{Xi} & $\Xi$ &
\gls{Pi} & $\Pi$ \\
\gls{Sigma} & $\Sigma$ &
\gls{Upsilon} & $\Upsilon$ &
\gls{Phi} & $\Phi$ \\
\gls{Psi} & $\Psi$ &
\gls{Omega} & $\Omega$
\end{tabular}
\end{table}

\xminisec{Example:}

The following code
\begin{codeS}
\glsni{begindispmath} x' = x + \glsni{Delta} x \glsni{enddispmath}
\end{codeS}
produces:
\begin{resultS}[x prime equals x plus capital Delta x]
\[ x' = x + \Delta x \]
\end{resultS}

%------------- Subscripts and Superscripts ------------

\setnode{subsupscripts}
\subsection{Subscripts and Superscripts}
\label{sec:scripts}

Subscripts are obtained either by the command
\begin{definition}
\gls{sb}\marg{\meta{maths}}
\end{definition}%
or by the \htmlref{special character}{sec:chars}:
\begin{definition}
\gls{underscorechar}\marg{\meta{maths}}
\end{definition}%
Superscripts are obtained either by the command
\begin{definition}
\gls{sp}\marg{\meta{maths}}
\end{definition}%
or by the special character:
\begin{definition}
\gls{circumchar}\marg{\meta{maths}}
\end{definition}%

\xminisec{Examples:}
\begin{enumerate}
\item  This example uses \glsni{sb} and \glsni{sp}:
\begin{codeS}
\glsni{begindispmath}y = x\glsni{sb}\marg{1}\glsni{sp}\marg{2} +
x\glsni{sb}\marg{2}\glsni{sp}\marg{2}\glsni{enddispmath}
\end{codeS}%

\item This example uses \glsni{underscorechar} and \glsni{circumchar}
\begin{codeS}
\glsni{begindispmath}y = x\glsni{underscorechar}\marg{1}\glsni{circumchar}\marg{2}
 + x\glsni{underscorechar}\marg{2}\glsni{circumchar}\marg{2}\glsni{enddispmath}
\end{codeS}%

\item Recall from 
\latexhtml{page~\pageref{itm:argnogrp}}{\htmlref{earlier}{itm:argnogrp}} that
\glspl{mandatory} only consisting of one
character don't need to be grouped, so the above code can also be
written as:
\begin{codeS}
\glsni{begindispmath}y = x\glsni{underscorechar}1\glsni{circumchar}2 
 + x\glsni{underscorechar}2\glsni{circumchar}2\glsni{enddispmath}
\end{codeS}%
This is simpler than the first two examples. However it's a good idea to be in the
habit of always using braces in case you forgot them when they're
needed.
\end{enumerate}
All three of the above examples produce the same output:
\begin{resultS}[subsup.html]
\[y = x_1^2 + x_2^2\]
\end{resultS}%

Notice how the subscript gets tucked under the slope of the $Y$ in:
\begin{codeS}
\glsni{begindispmath}
Y\glsni{underscorechar}\marg{1}\glsni{circumchar}\marg{2}
\glsni{enddispmath}
\end{codeS}
\begin{resultS}[ysubsup1.html]
\[ Y_{1}^{2} \]
\end{resultS}
Compare with
\begin{codeS}
\glsni{begindispmath}
Y\marg{}\glsni{underscorechar}\marg{1}\glsni{circumchar}\marg{2}
\glsni{enddispmath}
\end{codeS}
\begin{resultS}[ysubsup2.html]
\[ Y{}_{1}^{2} \]

\end{resultS}

\xminisec{Example (Nested)}

Subscripts and superscripts can also be nested (note that
it is now necessary to group the argument to the superscript command):
\begin{codeS}
\glsni{begindispmath}
 f(x) = e\glsni{circumchar}\marg{x\glsni{underscorechar}1}
\glsni{enddispmath}
\end{codeS}%
which produces
\begin{resultS}[nestedsup.html]
\[
f(x) = e^{x_1}
\]
\end{resultS}

This example isn't quite right as e isn't actually a variable
and shouldn't be typeset in italic.  The correct way to do this
is:\bookpagebreak
\begin{codeS}
\glsni{begindispmath}
 f(x) =
\gls{mathrm}\marg{e}\glsni{circumchar}\marg{x\glsni{underscorechar}1}
\glsni{enddispmath}
\end{codeS}%
which results in:
\begin{resultS}[Image: as above but the letter e is upright]
\[ f(x) = \mathrm{e}^{x_1} \]
\end{resultS}

\refstepcounter{object}\label{obj:e}If you are going to use e a lot,
it will be simpler to \htmlref{define a new command}{ch:newcom} to do
this. The definition should go in the
\glsterm{preamble}:
\begin{codeS}
\gls{newcommand}\marg{\cmdname{e}}\marg{\glsni{mathrm}\marg{e}}
\end{codeS}%
Then in the document:
\begin{codeS}
\glsni{begindispmath} f(x\glsni{underscorechar}1, x\glsni{underscorechar}2) 
 = \cmdname{e}\glsni{circumchar}\marg{x\glsni{underscorechar}1\glsni{circumchar}2} 
  +
\cmdname{e}\glsni{circumchar}\marg{x\glsni{underscorechar}2\glsni{circumchar}2}
\glsni{enddispmath}
\end{codeS}%
\begin{resultS}[nestedsup2.html]
\[
f(x_1, x_2) = \e^{x_1^2} + \e^{x_2^2}
\]
\end{resultS}

Take care when nesting subscripts or superscripts. The following
\begin{alltt}\wrong
x\glsni{underscorechar}1\glsni{underscorechar}2
\end{alltt}
will give a \texttt{!\ Double subscript} error.

%-------------- Functional Names --------------

\setnode{functionnames}
\subsection{Functional Names}
\label{sec:funcnames}

Functions such as log and tan can't simply be typed in as
\texttt{log} or \texttt{tan} otherwise they will come out looking
like the variables \mathorit{l} times
\mathorit{o} times \mathorit{g} (\makeimg{l o g}{$ l o g$})
or \mathorit{t} times \mathorit{a} times \mathorit{n} 
(\makeimg{t a n}{$ t a n$}).  Instead you should use
one of the commands listed in \tableref{tab:fnnames}.
The functions denoted with \supdag\ can have limits by using
the subscript command \gls{underscorechar} or the superscript command \gls{circumchar}\faq{Sub- and
superscript positioning for operators}{limits}.
In addition, the modulo commands listed in \tableref{tab:mod} are also available.

\begin{table}[bhtp]
\caption[Function Names]{Function Names (\supdag
indicates command may have limits, \supddag defined
by \protect\isty{amsmath}).}
\label{tab:fnnames}
\centering
\begin{tabular}{llllll}
\gls{arccos} & $\arccos$ &
\gls{arcsin} & $\arcsin$ &
\gls{arctan} & $\arctan$ \\
\gls{arg} & $\arg$ &
\gls{cos} & $\cos$ &
\gls{cosh} & $\cosh$ \\
\gls{cot} & $\cot$ &
\gls{coth} & $\coth$ &
\gls{csc} & $\csc$ \\
\gls{deg} & $\deg$ &
\gls{det}\supdag & $\det$ &
\gls{dim} & $\dim$ \\
\gls{exp} & $\exp$ &
\gls{gcd}\supdag & $\gcd$ &
\gls{hom} & $\hom$ \\
\gls{inf}\supdag & $\inf$ &
\gls{injlim}\supdagddag & $\injlim$ &
\gls{ker} & $\ker$ \\
\gls{lg} &  $\lg$ &
\gls{lim}\supdag & $\lim$ &
\gls{liminf}\supdag & $\liminf$ \\
\gls{limsup}\supdag & $\limsup$ &
\gls{ln} & $\ln$ &
\gls{log} & $\log$ \\
\gls{max}\supdag &  $\max$ &
\gls{min}\supdag & $\min$ &
\gls{Pr}\supdag & $\Pr$ \\
\gls{projlim}\supdagddag & $\projlim$ &
\gls{sec} & $\sec$ &
\gls{sin} & $\sin$ \\
\gls{sinh} & $\sinh$ &
\gls{sup}\supdag & $\sup$ &
\gls{tan} & $\tan$ \\
\gls{tanh} & $\tanh$ &
\gls{varinjlim}\supdagddag & $\varinjlim$ &
\gls{varliminf}\supdagddag & $\varliminf$ \\
\gls{varlimsup}\supdagddag & $\varlimsup$ &
\gls{varprojlim}\supdagddag &  $\varprojlim$ &
\end{tabular}
\end{table}

\begin{table}[htbp]
\caption[Modulo Commands]{Modulo Commands (\supddag
defined by \protect\isty{amsmath} package)}
\label{tab:mod}
\centering
\begin{tabular}{lll}
\toprule
\bfseries Command & \bfseries Example Input & \bfseries Example Output\\
\midrule
\gls{bmod} & \verb|$m \bmod n$| & $m \bmod n$\\
\gls{pmod}\marg{\meta{maths}} & \verb|$m \pmod{n}$| &$m \pmod{n}$\\
\gls{mod}\marg{\meta{maths}}\supddag & \verb|$m \mod{n}$| &$m \mod{n}$\\
\gls{pod}\marg{\meta{maths}}\supddag & \verb|$m \pod{n}$| &$m \pod{n}$
\\\bottomrule
\end{tabular}
\end{table}

\xminisec{Example (Trigonometric Functions):}

This example uses the cos and sin functions and also the 
\htmlref{Greek letter}{sec:greek} theta.
\begin{codeS}
\glsni{begindispmath} z = r(\glsni{cos}\glsni{theta} + i\glsni{sin}\glsni{theta})
\glsni{enddispmath}
\end{codeS}%
\begin{resultS}[Image: z equals r (cos theta plus i sign theta)]
\[
z = r(\cos\theta + i\sin\theta)
\]
\end{resultS}

\xminisec{Example (Limit):}

The command \gls{infty} is the 
infinity symbol \makeimg{image of infinity symbol}{$\infty$}, and the command \gls{to} displays an 
arrow pointing to the right.  Note the use of \gls{underscorechar} since the
limit is a subscript.
\begin{codeS}
\glsni{begindispmath}
  \glsni{lim}\glsni{underscorechar}\marg{x\glsni{to}\glsni{infty}} f(x) 
\glsni{enddispmath}
\end{codeS}%
\begin{resultS}[Image: limit as x tends to infinity of f of x]
\[
\lim_{x\to\infty} f(x)
\]
\end{resultS}%
The operators with limits behave differently depending on whether
they are in \htmlref{displayed}{sec:displayed} or 
\htmlref{in-line}{sec:inline} maths.
Notice the difference when the same code appears in-line:
\begin{codeS}
In a line of text 
\glsni{dollarchar}\glsni{lim}\glsni{underscorechar}\marg{x\glsni{to}\glsni{infty}}
f(x)\glsni{dollarchar}
\end{codeS}%
which now displays as:
\begin{resultS}[Image: the limit appears to the side of lim instead
of below it.]
In a line of text $\lim_{x\to\infty} f(x)$
\end{resultS}

\xminisec{Example (With Subscript):}

This is another example of a functional name using a subscript:
\begin{codeS}
\glsni{begindispmath} \glsni{min}\glsni{underscorechar}x f(x) \glsni{enddispmath}
\end{codeS}%
\begin{resultS}[Image: the x appears below the word min.]
\[
\min_x f(x)
\]
\end{resultS}%
Again, notice the difference when it is used in-line:
\begin{codeS}
In a line of text \glsni{dollarchar}\glsni{min}\glsni{underscorechar}x 
f(x)\glsni{dollarchar}
\end{codeS}%
\begin{resultS}[Image: the x appears to the side of the word min.]
In a line of text $\min_x f(x)$
\end{resultS}
\bookpagebreak

\setnode{declaremathop}
\subsubsection{Defining New Functional Operators}
\label{sec:declaremathop}

It may be that you want a function that isn't specified in
\tableref{tab:fnnames}. In this case, the \isty{amsmath} provides the
\textbf{\em \gls{preamble} only} command
\begin{definition}
\gls{DeclareMathOperator}\marg{\meta{cmd}}\marg{\meta{operator
name}}
\end{definition}
or its \htmlref{starred variant}{itm:starredcommand}
\begin{definition}
\glsni{DeclareMathOperator}*\marg{\meta{cmd}}\marg{\meta{operator
name}}
\end{definition}
\faq{Defining a new log-like function in LaTeX}{newfunction}%
Both versions define a command called \meta{cmd}, which must start
with a backslash, that typesets \meta{operator name} as a function
name.  The starred version is for function names that can take
limits (like \gls{lim} and \gls{min} described above).

\xminisec{Example (Operator Without Limits):}

Suppose I want a function called card, which represents the
cardinality of a set $\mathcal{S}$. First I~need to define the new
operator command (which I'm going to call \cmdname{card}) \emph{in
the \gls{preamble}}:
\begin{codeS}
\glsni{DeclareMathOperator}\marg{\cmdname{card}}\marg{card}
\end{codeS}
This operator doesn't take any limits, so I~have used the unstarred
version.

Later in the document, I~can use this new operator command:
\begin{codeS}
\glsni{begindispmath} n = \cmdname{card}(\glsni{mathcal}\marg{S}) \glsni{enddispmath}
\end{codeS}%
\begin{resultS}[Image: n equals card (S)]
\[
n = \card(\mathcal{S})
\]
\end{resultS}%
In this example \gls{mathcal} is used as sets are
typically represented in a calligraphic font.

\xminisec{Example (Operator With Limits):}

Suppose I~now want a function called mode, which represents the
mode of a set of numbers. First, I~define the operator command in the
preamble:
\begin{codeS}
\glsni{DeclareMathOperator}*\marg{\cmdname{mode}}\marg{mode}
\end{codeS}
This operator needs to be able to have a subscript, so I~have used
the starred version.

Later in the document, I~can use this new operator command:
\begin{codeS}
\glsni{begindispmath} x\glsni{underscorechar}m = 
\cmdname{mode}\glsni{underscorechar}\marg{x \glsni{in} \glsni{mathcal}\marg{S}}(x)
\glsni{enddispmath}
\end{codeS}
\begin{resultS}[Image: x subscript m equals mode of x in S of x]
\[
x_m = \mode_{x \in \mathcal{S}}(x)
\]
\end{resultS}

%----------- Fractions ------------

\bookpagebreak
\setnode{fractions}
\subsection{Fractions}
\label{sec:fractions}

Fractions are created using the command
\begin{definition}
\gls{frac}\marg{\meta{numerator}}\marg{\meta{denominator}}
\end{definition}
The \isty{amsmath} package also provides the command
\begin{definition}
\gls{cfrac}\oarg{\meta{pos}}\marg{\meta{numerator}}\marg{\meta{denominator}}
\end{definition}
which is designed for continued fractions. The optional argument
\texttt{pos} can be used for left (\texttt{l}) or right (\texttt{r})
placement of any of the numerators. (The default is centred.)

\xminisec{Example:}

A simple fraction:
\begin{codeS}
\glsni{begindispmath} 
 \glsni{frac}\marg{1}\marg{1+x} 
\glsni{enddispmath}
\end{codeS}
Produces:
\begin{resultS}[Image of fraction with 1 as the numerator and 1+x as
the denominator.]
\[ \frac{1}{1+x} \]
\end{resultS}
Compare with:
\begin{codeS}
In-line:
\glsni{dollarchar} 
 \glsni{frac}\marg{1}\marg{1+x} 
\glsni{dollarchar}
\end{codeS}
which produces:
\begin{resultS}[Image: as before but fraction is smaller.]
In-line: $\frac{1}{1+x}$
\end{resultS}

\xminisec{Example (Nested):}
\begin{codeS}
\glsni{begindispmath} 
\glsni{frac}\marg{1+\glsni{frac}\marg{1}\marg{x}}\marg{1+x+x\glsni{circumchar}2}
\glsni{enddispmath}
\end{codeS}%
\begin{resultS}[Image of fraction with (1+1/x) as the numerator and
(1 + x + x squared as the denominator.]
\[
\frac{1+\frac{1}{x}}{1+x+x^2}
\]
\end{resultS}

\xminisec{Example (Continued Fraction);}
A continued fraction (example taken from \isty{amsmath}
documentation and uses \glsni{sqrt}, described in \sectionref{sec:roots},
and \glsni{dotsb}, described in
\sectionref{sec:mathssym}):
\begin{code}
\begin{alltt}
\glsni{begindispmath}
\glsni{cfrac}\marg{1}\marg{\glsni{sqrt}\marg{2}+
\glsni{cfrac}\marg{1}\marg{\glsni{sqrt}\marg{2}+
\glsni{cfrac}\marg{1}\marg{\glsni{sqrt}\marg{2}+\glsni{dotsb}
}}}
\glsni{enddispmath}
\end{alltt}
\end{code}\screenpagebreak
\begin{result}[Image of a continued fraction (the denominator has a
fraction with a fraction in its denominator etc).]
\[
\cfrac{1}{\sqrt{2}+
\cfrac{1}{\sqrt{2}+
\cfrac{1}{\sqrt{2}+\dotsb
}}}
\]
\end{result}

\xminisec{Example (A Derivative):}
\begin{codeS}
\glsni{begindispmath} f'(x) = \glsni{frac}\marg{df}\marg{dx} \glsni{enddispmath}
\end{codeS}%
\begin{resultS}[Image: f'(x) = d f by d x]
\[
f'(x) = \frac{df}{dx}
\]
\end{resultS}%
As with \dq{e}, the differential operator \dq{d} should
be in an upright font as it is not a variable:
\begin{code}
\begin{alltt}
\glsni{begindispmath}
  f'(x) = \glsni{frac}\marg{\glsni{mathrm}\marg{d}f}\marg{\glsni{mathrm}\marg{d}x}
\glsni{enddispmath}
\end{alltt}
\end{code}%
\begin{resultS}[Image: as above but the letter d is upright]
\[
f'(x) = \frac{\mathrm{d}f}{\mathrm{d}x}
\]
\end{resultS}

The above example is rather cumbersome, particularly
if you have a lot of derivatives, so it might be easier to
\htmlref{define a new command}{ch:newcom}\latex{ (see \chapterref{ch:newcom})}. In the 
\glsterm{preamble} define:
\begin{codeS}
\gls{newcommand}\marg{\cmdname{deriv}}\oarg{2}\marg{\glsni{frac}\marg{\glsni{mathrm}\marg{d}\glsni{hashchar}1}\marg{\glsni{mathrm}\marg{d}\glsni{hashchar}2}}
\end{codeS}%
Then in the document:
\begin{codeS}
\glsni{begindispmath} f'(x) = \cmdname{deriv}\marg{f}\marg{x} \glsni{enddispmath}
\end{codeS}%
\begin{resultS}[Image: same as before]
\newcommand{\deriv}[2]{\frac{\mathrm{d}#1}{\mathrm{d}#2}}

\[ f'(x) = \deriv{f}{x} \]
\end{resultS}

\xminisec{Example (Partial Derivative):}
\label{itm:pderiv}%

Partial derivatives can be obtained similarly
using the command \gls{partial} to display the partial
derivative symbol. As in the previous example, first define a new
command to format a partial derivative in the
\glsterm{preamble}:
\begin{codeS}
\gls{newcommand}\marg{\cmdname{pderiv}}\oarg{2}\marg{\glsni{frac}\marg{\glsni{partial} \glsni{hashchar}1}{\glsni{partial} \glsni{hashchar}2}}
\end{codeS}%
Then in the document:
\begin{codeS}
\glsni{begindispmath} f\glsni{underscorechar}x = \cmdname{pderiv}\marg{f}\marg{x} \glsni{enddispmath}
\end{codeS}%
\begin{resultS}[Image: f subscript x = partial d f by partial d x]
\[
f_x = \pderiv{f}{x}
\]
\end{resultS}

\xminisec{Example (Double Partial Derivative):}

\begin{code}
\begin{alltt}
\glsni{begindispmath}
  f\glsni{underscorechar}\marg{xy} = \glsni{frac}\marg{\glsni{partial}\glsni{circumchar}2 f}\marg{\glsni{partial} x \glsni{partial} y}
\glsni{enddispmath}
\end{alltt}
\end{code}%
\begin{resultS}[Image: f subscript xy = partial d 2 f by partial d x partial d y.]
\[
f_{xy} = \frac{\partial^2 f}{\partial x \partial y}
\]
\end{resultS}

\xminisec{Example (First principles):}
\begin{code}
\glsni{begindispmath}\newline
\mbox{}\qquad f'(x) = \gls{lim}\gls{underscorechar}\marg{\gls{Delta} x \gls{to}
0}\newline
\mbox{}\qquad \glsni{frac}\marg{f(x + \gls{Delta} x)-f(x)}\marg{\gls{Delta}
x}\newline
\glsni{enddispmath}
\end{code}
\begin{resultS}
\[
f'(x) = \lim_{\Delta x\to0}
  \frac{f(x+\Delta x) - f(x)}{\Delta x}
\]
\end{resultS}

%------------ Roots --------------

\setnode{roots}
\subsection{Roots}
\label{sec:roots}

Roots are obtained using the command
\begin{definition}
\gls{sqrt}\oarg{\meta{order}}\marg{\meta{maths}}
\end{definition}%
without the optional argument \meta{order}
it will produce a simple square root. Cubic roots etc
can be obtained using the optional argument.

\xminisec{Examples:}

\begin{enumerate}
\item A square root:
\begin{codeS}
\glsni{begindispmath} \glsni{sqrt}\marg{a+b} \glsni{enddispmath}
\end{codeS}%
\begin{resultS}[Image: square root of (a + b)]
\[
\sqrt{a+b}
\]
\end{resultS}

\item A cubic root:
\begin{codeS}
\glsni{begindispmath} \glsni{sqrt}\oarg{3}\marg{a+b} \glsni{enddispmath}
\end{codeS}%
\begin{resultS}[Image: cubic root of (a + b)]
\[
\sqrt[3]{a+b}
\]
\end{resultS}

\item An \ensuremath{n}th root:
\begin{codeS}
\glsni{begindispmath} \glsni{sqrt}\oarg{n}\marg{a+b} \glsni{enddispmath}
\end{codeS}%
\begin{resultS}[Image: the nth root of (a + b)]
\[
\sqrt[n]{a+b}
\]
\end{resultS}
\end{enumerate}

%-------- Mathematical Symbols -------------

\setnode{mathssym}
\subsection{Mathematical Symbols}
\label{sec:mathssym}

Relational symbols\faq{Where can I find the symbol for
\ldots}{symbols} are shown in \tableref{tab:relsym}.  If you want a
negation that is not shown, you can obtain it by preceding the
symbol with the command \gls{not}.  For example:
\glsni{not}\glsni{subset} produces the symbol 
\makeimg{image of subset symbol with slash through it}{$\not\subset$}.

\begin{table}[hbtp]
\caption{Relational Symbols}
\label{tab:relsym}
\centering
\begin{tabular}{llllll}
\gls{approx} & $\approx$ &
\gls{asymp} & $\asymp$ &
\gls{bowtie} & $\bowtie$ \\
\gls{cong} & $\cong$ &
\gls{dashv} & $\dashv$ &
\gls{doteq} & $\doteq$ \\
\gls{equiv} & $\equiv$ &
\gls{frown} & $\frown$ &
\gls{ge} or \gls{geq} & $\geq$ \\
\gls{gg} & $\gg$ &
\gls{in} & $\in$ &
\gls{le} or \gls{leq} & $\leq$ \\
\gls{ll} & $\ll$ &
\gls{mid} or \verb"|" & $\mid$ &
\gls{models} & $\models$ \\ 
\gls{neq} & $\neq$ &
\gls{ni} & $\ni$ &
\gls{notin} & $\notin$ \\
\gls{parallel} & $\parallel$ &
\gls{prec}& $\prec$ &
\gls{preceq}& $\preceq$ \\
\gls{perp} & $\perp$ &
\gls{propto} & $\propto$ &
\gls{sim} & $\sim$ \\
\gls{simeq} & $\simeq$ &
\gls{smile} & $\smile$ &
\gls{sqsubseteq} & $\sqsubseteq$ \\
\gls{sqsupseteq} & $\sqsupseteq$ &
\gls{subset} & $\subset$ &
\gls{subseteq} & $\subseteq$ \\
\gls{succ} & $\succ$ &
\gls{succeq} & $\succeq$ &
\gls{supset} & $\supset$ \\
\gls{supseteq} & $\supseteq$ &
\gls{vdash} & $\vdash$ &
\end{tabular}
\end{table}

Binary operator symbols are shown in \tableref{tab:binop}, and arrow
symbols are shown in \xtableref{tab:arrows}. There are also over and
under arrows (\xtableref{tab:overunderarrows}) that have an argument.
The over arrows put an extendible arrow over their argument, and the
under arrows put an extendible arrow under their argument. In
addition, the \isty{amsmath} package provides extensible arrows that
take a superscript and, optionally, a subscript:
\begin{definition}
\gls{xleftarrow}\oarg{\meta{subscript}}\marg{\meta{superscript}}
\end{definition}
\begin{definition}
\gls{xrightarrow}\oarg{\meta{subscript}}\marg{\meta{superscript}}
\end{definition}

\xminisec{Example:}
\begin{code}
\glsni{begindispmath}\newline
\mbox{}\qquad  A \glsni{xleftarrow}\marg{n+m-p} B  
 \glsni{xrightarrow}\oarg{X}\marg{n+p} C\newline
\glsni{enddispmath}
\end{code}
\begin{result}[Image of result: the arrow subscripts and superscript are centred over the arrow.]
\[
  A \xleftarrow{n+m-p} B \xrightarrow[X]{n+p} C
\]
\end{result}

\begin{table}[tbhp]
\caption{Binary Operator Symbols}
\label{tab:binop}
\centering
\begin{tabular}{llllll}
\gls{amalg} & $\amalg$ &
\gls{ast} & $\ast$ &
\gls{bullet} & $\bullet$ \\
\gls{bigcirc} & $\bigcirc$ &
\gls{bigtriangledown} & $\bigtriangledown$ &
\gls{bigtriangleup} & $\bigtriangleup$ \\
\gls{cap} & $\cap$ &
\gls{cdot} & $\cdot$ &
\gls{circ} & $\circ$ \\
\gls{cup} & $\cup$ &
\gls{dagger} & $\dagger$ &
\gls{ddagger} & $\ddagger$ \\
\gls{diamond} & $\diamond$ &
\gls{div} & $\div$ &
\gls{mp} & $\mp$ \\
\gls{odot} & $\odot$ &
\gls{ominus} & $\ominus$ &
\gls{oplus} & $\oplus$ \\
\gls{oslash} & $\oslash$ &
\gls{otimes} & $\otimes$ &
\gls{pm} & $\pm$ \\
\gls{setminus} & $\setminus$ &
\gls{sqcap} & $\sqcap$ &
\gls{sqcup} & $\sqcup$ \\
\gls{star} & $\star$ &
\gls{times} & $\times$ &
\gls{triangleleft} & $\triangleleft$ \\
\gls{triangleright} & $\triangleright$ &
\gls{uplus} & $\uplus$ &
\gls{vee} & $\vee$ \\
\gls{wedge} & $\wedge$ &
\gls{wr} & $\wr$
\end{tabular}
\end{table}

\begin{table}[tbhp]
\caption{Arrow Symbols}
\label{tab:arrows}
\centering
\begin{tabular}{llll}
\gls{downarrow} & $\downarrow$ &
\gls{Downarrow} & $\Downarrow$ \\
\gls{hookleftarrow} & $\hookleftarrow$ &
\gls{hookrightarrow} & $\hookrightarrow$ \\
\gls{leftarrow} or \gls{gets} & $\leftarrow$ &
\gls{Leftarrow} & $\Leftarrow$ \\
\gls{leftharpoondown} & $\leftharpoondown$ &
\gls{leftharpoonup} & $\leftharpoonup$ \\
\gls{leftrightarrow} & $\leftrightarrow$ &
\gls{Leftrightarrow} & $\Leftrightarrow$ \\
\gls{longleftarrow} & $\longleftarrow$ &
\gls{Longleftarrow} & $\Longleftarrow$ \\
\gls{longleftrightarrow} & $\longleftrightarrow$ &
\gls{Longleftrightarrow} & $\Longleftrightarrow$ \\
\gls{longmapsto} & $\longmapsto$ &
\gls{longrightarrow} & $\longrightarrow$ \\
\gls{Longrightarrow} & $\Longrightarrow$ &
\gls{mapsto} & $\mapsto$ \\
\gls{nearrow} & $\nearrow$ &
\gls{nwarrow} & $\nwarrow$ \\
\gls{rightarrow} or \gls{to} & $\to$ &
\gls{Rightarrow} & $\Rightarrow$ \\
\gls{rightharpoondown} & $\rightharpoondown$ &
\gls{rightharpoonup} & $\rightharpoonup$ \\
\gls{rightleftharpoons} & $\rightleftharpoons$ &
\gls{searrow} & $\searrow$ \\
\gls{swarrow} & $\swarrow$ &
\gls{uparrow} & $\uparrow$ \\
\gls{Uparrow} & $\Uparrow$ &
\gls{updownarrow} & $\updownarrow$ \\
\gls{Updownarrow} & $\Updownarrow$ &
\end{tabular}
\end{table}

\begin{table}[htbp]
\caption[Over and Under Arrows]{Over and Under Arrows
(\supdag defined by \protect\isty{amsmath})}
\label{tab:overunderarrows}
\centering
\begin{tabular}{@{}lll@{}}
\toprule
\bfseries Definition & \multicolumn{2}{c}{\bfseries Example}\\
\midrule
\begin{inlinedef}\gls{overleftarrow}\marg{\meta{maths}}\end{inlinedef}
  & \verb|\overleftarrow{ABC}| & $\overleftarrow{ABC}$\\
\begin{inlinedef}\gls{overrightarrow}\marg{\meta{maths}}\end{inlinedef}
  & \verb|\overrightarrow{ABC}| & $\overrightarrow{ABC}$\\
\begin{inlinedef}\gls{overleftrightarrow}\marg{\meta{maths}}\end{inlinedef}\supdag
  & \verb|\overleftrightarrow{ABC}| & $\overleftrightarrow{ABC}$\\
\begin{inlinedef}\gls{underleftarrow}\marg{\meta{maths}}\end{inlinedef}\supdag
  & \verb|\underleftarrow{ABC}| & $\underleftarrow{ABC}$\\
\begin{inlinedef}\gls{underrightarrow}\marg{\meta{maths}}\end{inlinedef}\supdag
  & \verb|\underrightarrow{ABC}| & $\underrightarrow{ABC}$\\
\begin{inlinedef}\gls{underleftrightarrow}\marg{\meta{maths}}\end{inlinedef}\supdag
  & \verb|\underleftrightarrow{ABC}| & $\underleftrightarrow{ABC}$
\\\bottomrule
\end{tabular}
\end{table}

\bookpagebreak
\begin{table}[htbp]
\caption{Symbols with Limits}
\label{tab:symlim}
\centering
\begin{tabular}{llllll}
\gls{sum} & $\sum$ &
\gls{int} & $\int$ &
\gls{oint} & $\oint$ \\[5pt]
\gls{prod} & $\prod$ &
\gls{coprod} & $\coprod$ &
\gls{bigcap} & $\bigcap$ \\[5pt]
\gls{bigcup} & $\bigcup$ &
\gls{bigsqcup} & $\bigsqcup$ &
\gls{bigvee} & $\bigvee$ \\[5pt]
\gls{bigwedge} & $\bigwedge$ &
\gls{bigodot} & $\bigodot$ &
\gls{bigotimes} & $\bigotimes$ \\[5pt]
\gls{bigoplus} & $\bigoplus$ &
\gls{biguplus} & $\biguplus$
\end{tabular}
\end{table}

Symbols that can have limits are shown in \tableref{tab:symlim}.
The size of these symbols depends on whether they are in displayed 
maths or in-line maths. 

\xminisec{Example (Displayed Summation and Product):}

The limits of summations and products are placed above and below the
symbol in displayed maths:
\begin{code}
\begin{alltt}
\glsni{begindispmath}
  f(x) = \glsni{sum}\glsni{underscorechar}\marg{i=1}\glsni{circumchar}\marg{n} x\glsni{underscorechar}i + \glsni{prod}\glsni{underscorechar}\marg{i=1}\glsni{circumchar}\marg{n} x\glsni{underscorechar}i
\glsni{enddispmath}
\end{alltt}
\end{code}
\begin{resultS}[Image: f(x) = sum from i equals 1 to n x subscript i
+ product from i equals 1 to n x subscript i.]
\[
f(x) = \sum_{i=1}^{n} x_i + \prod_{i=1}^{n} x_i
\]
\end{resultS}

\xminisec{Example (In-line Summation and Product):}

The limits of summations and products are placed to the right of the
symbol in in-line maths:
\begin{code}
\begin{alltt}
In a line of text:
\glsni{begin}\marg{math}
  f(x) = \glsni{sum}\glsni{underscorechar}\marg{i=1}\glsni{circumchar}\marg{n} x\glsni{underscorechar}i + \glsni{prod}\glsni{underscorechar}\marg{i=1}\glsni{circumchar}\marg{n} x\glsni{underscorechar}i
\glsni{end}\marg{math}
\end{alltt}
\end{code}
\begin{resultS}[Image: as before but summation and product symbols
are smaller and limits are to the side instead of above and below]
In a line of text:
\begin{math}
f(x) = \sum_{i=1}^{n} x_i + \prod_{i=1}^{n} x_i
\end{math}
\end{resultS}

\xminisec{Multiline Sub- or Superscripts}

The \isty{amsmath} package provides the command:
\begin{definition}
\gls{substack}\marg{\meta{maths}}
\end{definition}
which can be used for multiline sub- or superscripts. Within the
argument \meta{maths} use \gls{tab.dbbackslashchar} to separate
rows. For example:
\begin{code}
\glsni{begindispmath}\newline
\strut~\glsni{sum}\glsni{underscorechar}\marg{\newline
\strut~~\glsni{substack}\newline
\strut~~\marg{\newline
\strut~~~~i \glsni{in}
\glsni{mathcal}\marg{I}\glsni{tab.dbbackslashchar}\newline
\strut~~~~i \glsni{neq} 0\newline
\strut~~}\newline
\strut~}\newline
\strut~x\glsni{underscorechar}i\newline
\glsni{enddispmath}
\end{code}
\begin{result}[substack.html]
\[
 \sum_{
   \substack
   {
     i \in \mathcal{I}\\
     i \neq 0
   }
 }
 x_i
\]
\end{result}

\setnode{ellipses}
\subsection{Ellipses}
\label{sec:ellipses}

Ellipsis (omission mark)\index{ellipses (omission marks)|hyperpage} commands are 
shown in \tableref{tab:ellipsis}. The \isty{amsmath} package also
provides: \glsni{dotsc} for dots with commas, \glsni{dotsb} for dots
with binary operators/relations, \glsni{dotsm} for multiplication
dots, \glsni{dotsi} for dots with integrals and \glsni{dotso} for other
dots, which can be used as replacements for \glsni{ldots} and
\glsni{cdots}.

\begin{table}[hbtp]
\caption[Ellipses]{Ellipses (\supdag\ provided by \sty{amsmath}
package)}
\label{tab:ellipsis}
\centering
\begin{tabular}{ll@{\hspace{4\tabcolsep}}ll@{\hspace{4\tabcolsep}}llllll}
\gls{vdots} & $\vdots$ &
\gls{cdots} & $\cdots$ &
\gls{dotsb}\supdag & $\dotsb$ &
\gls{dotsi}\supdag & $\dotsi$ &
\gls{dotsm}\supdag & $\dotsm$\\
\gls{ddots} & $\ddots$ &
\gls{ldots} & $\ldots$ &
\gls{dotsc}\supdag & $\dotsc$ &
\gls{dotso}\supdag & $\dotso$ 
\end{tabular}
\end{table}

\xminisec{Example (Low Ellipsis):}

This example uses the command 
\refstepcounter{object}\label{cmd:forall}\gls{forall} to produce
the \dq{for all} symbol $\forall$, and it also uses 
\glsni{spacesym} (backslash space) to make a space before the for all symbol. The
\isty{amsmath} \dq{dots with commas} ellipsis \glsni{dotsc} is used rather
than the standard \glsni{ldots}:
\begin{code}
\begin{alltt}
\glsni{begindispmath}
a\glsni{underscorechar}ix\glsni{underscorechar}i = b\glsni{underscorechar}i\glsni{spacesym}\glsni{forall} i = 1,\glsni{dotsc}, n
\glsni{enddispmath}
\end{alltt}
\end{code}%
\begin{resultS}[forall.html]
\[
a_ix_i = b_i\ \forall i = 1,\dotsc, n
\]
\end{resultS}

\xminisec{Example (Centred ellipsis):}

This example uses the \isty{amsmath} \dq{dots with binary operators\slash relations}
\glsni{dotsb} instead of the standard \glsni{cdots}:
\begin{code}
\begin{alltt}
\glsni{begindispmath}
y = a\glsni{underscorechar}1 + a\glsni{underscorechar}2 + \glsni{dotsb} + a\glsni{underscorechar}n 
\glsni{enddispmath}
\end{alltt}
\end{code}%
\begin{resultS}[cdots.html]
\[
y = a_1 + a_2 + \dotsb + a_n 
\]
\end{resultS}

\begin{exercise}{Maths: Fractions and Symbols}{ex:mathssym}
This exercise uses a fraction, a square root, subscripts, 
superscripts and symbols.
Try to reproduce the following output:
\begin{result}[maths1ex.html]
The quadratic equation
\[
\sum_{i=0}^{2} a_i x^i = 0
\]
has solutions given by
\[
x = \frac{-a_1 \pm \sqrt{a_1^2 - 4 a_2 a_0}}{2a_2}
\]
\end{result}%
Again you can \downloadorview{mathssym} the solution.
\end{exercise}

%---------- Delimiters ------------

\setnode{delimiters}
\subsection{Delimiters}
\label{sec:delimiters}

Placing brackets around a tall object in maths mode, such as
fractions, does not look right if you use normal sized brackets.
For example:
\begin{code}
\glsni{begindispmath}\newline
\strut~~(\glsni{frac}\marg{1}\marg{1+x})\newline
\glsni{enddispmath}
\end{code}%
results in:
\begin{resultS}[Image: normal sized brackets around fraction.]
\[
(\frac{1}{1+x})
\]
\end{resultS}\bookpagebreak\noindent
Instead, you can automatically resize the delimiters using the
commands:
\begin{definition}
\gls{left}\meta{delimiter}
\end{definition}%
and
\begin{definition}
\gls{right}\meta{delimiter}
\end{definition}%
Rewriting the above example:
\begin{code}
\begin{alltt}
\glsni{begindispmath}
  \glsni{left}( \glsni{frac}\marg{1}\marg{1+x} \glsni{right})
\glsni{enddispmath}
\end{alltt}
\end{code}%
produces:
\begin{resultS}[Image: brackets now the same size as the fraction.]
\[
  \left( \frac{1}{1+x} \right)
\]
\end{resultS}%

Note that you must always have matching \cmdname{left} and
\cmdname{right} commands, although the delimiters used may be
different. \refstepcounter{object}\label{obj:invisdelim}If you want
one of the delimiters to be invisible, use a
\gls{delimiter.periodchar}\ (full stop) as the delimiter. Available
delimiters are shown in
\tableref{tab:delimiters}. (Note for a vertical bar
delimiter it's best to use \isty{amsmath}'s \glsni{lvert} command instead of
\gls{delim.barchar} and \glsni{lVert} instead of \gls{csbar}.) Sometimes using \glsni{left} and
\glsni{right} doesn't produce the optimal sized delimiters. In which
case you can use additional commands provided by the \isty{amsmath}
package shown in \tableref{tab:big}.

\begin{table}[hbtp]
\caption[Delimiters]{Delimiters (\supdag defined by
\protect\isty{amsmath})}
\label{tab:delimiters}
\centering
\begin{tabular}{llllllll}
\begin{inlinedef}\gls{openparen}\end{inlinedef} &
\makeimg{open round bracket}{$($} &
\begin{inlinedef}\gls{closeparen}\end{inlinedef} &
\makeimg{close round bracket}{$)$} &
\begin{inlinedef}\gls{delimiter.opensq}\end{inlinedef} & 
\makeimg{open square bracket}{$[$} &
\begin{inlinedef}\gls{delimiter.closesq}\end{inlinedef} & 
\makeimg{close square bracket}{$]$} \\
\begin{inlinedef}\gls{leftbrace}\end{inlinedef} & 
\makeimg{open curly brace}{$\{$} &
\begin{inlinedef}\gls{rightbrace}\end{inlinedef} & 
\makeimg{close curly brace}{$\}$} &
\begin{inlinedef}\gls{lvert}\supdag\end{inlinedef} & 
\makeimg{left vertical bar}{$\lvert$} &
\begin{inlinedef}\gls{rvert}\supdag\end{inlinedef} & 
\makeimg{right vertical bar}{$\rvert$} \\
\begin{inlinedef}\gls{lVert}\supdag\end{inlinedef} & 
\makeimg{left double vertical bar}{$\lVert$} &
\begin{inlinedef}\gls{rVert}\supdag\end{inlinedef} & 
\makeimg{right double vertical bar}{$\rVert$} &
\begin{inlinedef}\gls{langle}\end{inlinedef} & 
\makeimg{left angle bracket}{$\langle$} &
\begin{inlinedef}\gls{rangle}\end{inlinedef} & 
\makeimg{right angle bracket}{$\rangle$} \\
\begin{inlinedef}\gls{lfloor}\end{inlinedef} & 
\makeimg{left floor}{$\lfloor$} &
\begin{inlinedef}\gls{rfloor}\end{inlinedef} & 
\makeimg{right floor}{$\rfloor$} &
\begin{inlinedef}\gls{lceil}\end{inlinedef} & 
\makeimg{left ceiling}{$\lceil$} &
\begin{inlinedef}\gls{rceil}\end{inlinedef} & 
\makeimg{right ceiling}{$\rceil$}\\
\begin{inlinedef}\gls{uparrow}\end{inlinedef} & 
\makeimg{up arrow}{$\uparrow$} &
\begin{inlinedef}\gls{downarrow}\end{inlinedef} &
\makeimg{down arrow}{$\downarrow$} &
\begin{inlinedef}\gls{Uparrow}\end{inlinedef} & 
\makeimg{double lined up arrow}{$\Uparrow$} &
\begin{inlinedef}\gls{Downarrow}\end{inlinedef} & 
\makeimg{double lined down arrow}{$\Downarrow$}\\
\begin{inlinedef}\gls{updownarrow}\end{inlinedef} & 
\makeimg{up down arrow}{$\updownarrow$} &
\begin{inlinedef}\gls{Updownarrow}\end{inlinedef} & 
\makeimg{double lined up down arrow}{$\Updownarrow$} &
\begin{inlinedef}\gls{delim.slash}\end{inlinedef} & 
\makeimg{forward slash}{$/$} &
\begin{inlinedef}\gls{backslash}\end{inlinedef} & 
\makeimg{backslash}{$\backslash$} 
\end{tabular}
\end{table}

\begin{table}
\caption[Delimiter Sizing]{Additional Commands Provided by
\protect\isty{amsmath} for Delimiter Sizing}
\label{tab:big}
\centering
\renewcommand*{\arraystretch}{2.2}
\begin{tabular}{@{}ll@{\hspace{4\tabcolsep}}ll@{}}
\toprule
\multicolumn{2}{c@{\hspace{4\tabcolsep}}}{\bfseries Definitions} &
\multicolumn{2}{c}{\bfseries Example}\\
\midrule
\multicolumn{2}{c@{\hspace{4\tabcolsep}}}{\em Default Size} & \verb|$( X )$| & $(X)$\\
\begin{inlinedef}\gls{bigl}\meta{delim}\end{inlinedef} &
\begin{inlinedef}\gls{bigr}\meta{delim}\end{inlinedef} &
\verb|$\bigl( X \bigr)$| & 
\makeimg{(X)}{$\bigl(X\bigr)$} \\
\begin{inlinedef}\gls{Bigl}\meta{delim}\end{inlinedef} &
\begin{inlinedef}\gls{Bigr}\meta{delim}\end{inlinedef} &
\verb|$\Bigl( X \Bigr)$| & 
\makeimg{(X)}{$\Bigl(X\Bigr)$} \\
\begin{inlinedef}\gls{biggl}\meta{delim}\end{inlinedef} &
\begin{inlinedef}\gls{biggr}\meta{delim}\end{inlinedef} &
\verb|$\bigl( X \biggr)$| & 
\makeimg{(X)}{$\biggl(X\biggr)$} \\
\begin{inlinedef}\gls{Biggl}\meta{delim}\end{inlinedef} &
\begin{inlinedef}\gls{Biggr}\meta{delim}\end{inlinedef} &
\verb|$\Biggl( X \Biggr)$| & 
\makeimg{(X)}{$\Biggl(X\Biggr)$} \\
\bottomrule
\end{tabular}
\renewcommand*{\arraystretch}{1}% Keep LaTeX2HTML happy
\end{table}

\xminisec{Example (Vertical Bar Delimiters):}

\begin{code}
\glsni{begindispmath}\newline
\glsni{left}\glsni{lvert}\newline
\glsni{frac}\marg{1}\marg{1+x}\newline
\glsni{right}\glsni{rvert}\newline
\glsni{enddispmath}
\end{code}%
\begin{resultS}[Image: vertical bars same height as fraction.]
\[
\left\lvert\frac{1}{1+x}\right\rvert
\]
\end{resultS}

\xminisec{Example (Delimiter with Subscript):}

Delimiters can take limits:
\begin{code}
\glsni{begindispmath}\newline
\glsni{left}\glsni{lvert}\newline
\glsni{frac}\marg{1}\marg{1+x}\newline
\glsni{right}\glsni{rvert}\glsni{underscorechar}\marg{x=0}\newline
\glsni{enddispmath}
\end{code}\reportpagebreak
\begin{resultS}[Image: as before but the right vertical bar has a subscript.]
\[
\left\lvert\frac{1}{1+x}\right\rvert_{x=0}
\]
\end{resultS}

\xminisec{Example (Mismatch):}
The left and right delimiters don't have to match:
\begin{code}
\begin{alltt}
\glsni{begindispmath}
\glsni{left}\glsni{delimiter.opensq}\glsni{frac}\marg{1}\marg{1+x}\glsni{right}\glsni{rangle}
\glsni{enddispmath}
\end{alltt}
\end{code}%
\begin{resultS}[Image: large open square bracket followed by
fraction followed by large right angle bracket.]
\[
\left[\frac{1}{1+x}\right\rangle
\]
\end{resultS}

\xminisec{Example (An invisible delimiter):}

Every \gls{right} must have a matching \gls{left} (and vice versa), so 
use a \gls{delimiter.periodchar}\ (full stop) for an invisible
delimiter.
\begin{code}
\begin{alltt}
\glsni{begindispmath}
 \glsni{left}\glsni{delimiter.periodchar}
  \glsni{frac}\marg{\glsni{partial} f}\marg{\glsni{partial} x}
 \glsni{right}\glsni{rvert}\glsni{underscorechar}\marg{x=0}
\glsni{enddispmath}
\end{alltt}
\end{code}
\begin{resultS}[Image: partial fraction followed by large vertical
line with a subscript.]
\[
\left.
\frac{\partial f}{\partial x}
\right\rvert_{x=0}
\]
\end{resultS}

\refstepcounter{object}\label{introeq}
We have now covered enough to reproduce the equation shown in 
\chapterref{ch:intro}:
\begin{code}
\begin{alltt}
\glsni{newcommand}*\marg{\cmdname{pderiv}}\oarg{2}\marg{\glsni{frac}\marg{\glsni{partial} \glsni{hashchar}1}\marg{\glsni{partial} \glsni{hashchar}2}}
\glsni{newcommand}*\marg{\cmdname{e}}\marg{\glsni{mathrm}\marg{e}}

\glsni{begindispmath}
\cmdname{pderiv}\marg{\glsni{circumchar}2\glsni{mathcal}\marg{L}}\marg{\marg{z\glsni{underscorechar}i\glsni{circumchar}\glsni{rho}}\glsni{circumchar}2} =
-\cmdname{pderiv}\marg{\glsni{rho}\glsni{underscorechar}i}\marg{z\glsni{underscorechar}i\glsni{circumchar}\glsni{rho}}
\glsni{left}(
  \cmdname{pderiv}\marg{v\glsni{underscorechar}i}\marg{\glsni{rho}\glsni{underscorechar}i} \glsni{frac}\marg{\cmdname{e}\glsni{circumchar}\marg{v\glsni{underscorechar}i}}\marg{1-\cmdname{e}\glsni{circumchar}\marg{v\glsni{underscorechar}i}}
  + v\glsni{underscorechar}i \glsni{frac}\marg{\cmdname{e}\glsni{circumchar}\marg{v\glsni{underscorechar}i}\cmdname{pderiv}\marg{v\glsni{underscorechar}i}\marg{\glsni{rho}\glsni{underscorechar}i}(1-\cmdname{e}\glsni{circumchar}\marg{v\glsni{underscorechar}i})
             +\cmdname{e}\glsni{circumchar}\marg{2v\glsni{underscorechar}i}\cmdname{pderiv}\marg{v\glsni{underscorechar}i}\marg{\glsni{rho}\glsni{underscorechar}i}}\marg{(1-\cmdname{e}\glsni{circumchar}\marg{v\glsni{underscorechar}i})\glsni{circumchar}2}
\glsni{right})
\glsni{enddispmath}
\end{alltt}
\end{code}%
\begin{result}[Image showing a second order partial derivative.]
\[
\pderiv{^2\mathcal{L}}{{z_i^\rho}^2} =
-\pderiv{\rho_i}{z_i^\rho}
\left(
  \pderiv{v_i}{\rho_i} \frac{\e^{v_i}}{1-\e^{v_i}}
  + v_i \frac{\e^{v_i}\pderiv{v_i}{\rho_i}(1-\e^{v_i})+\e^{2v_i}\pderiv{v_i}{\rho_i}}{(1-\e^{v_i})^2}
\right)
\]
\end{result}

\xminisec{Note:}
The above code looks a bit complicated, and there are
so many braces that it can be easy to lose track, so here are some
ways of making it a little easier to type:
\begin{enumerate}
\item Whenever you start a new environment type in the
\gls{begin} and \gls{end} bits first, and then insert
whatever goes inside the environment.  This ensures that you always
have a matching \gls{begin} and \gls{end}. The same goes for 
\glsni{begindispmath} and \glsni{enddispmath}.

\item Whenever you type any braces, always type the opening and
closing braces first, and then insert whatever goes in between.  This
will ensure that your braces always match up.
\end{enumerate}

So keeping these notes in mind, let's try typing in the code in a
methodical manner:\screenpagebreak
\begin{enumerate}
\item Start and end the displayed maths mode:
\begin{code}
\glsni{begindispmath}\newline
\glsni{enddispmath}
\end{code}

\item We now need a partial derivative.
(The command \cmdname{pderiv} is defined as described
\htmlref{earlier}{itm:pderiv}\latex{\ifscreenorother{}{\ on 
page~\pageref{itm:pderiv}}}. Make sure you remember to define it,
preferably in the \gls{preamble}.)
\begin{code}
{\color{midgray}
\glsnl{begindispmath}}\newline
\cmdname{pderiv}\marg{}\marg{}\newline
{\color{midgray}
\glsnl{enddispmath}}
\end{code}

\item Let's do the first argument. This partial derivative is
actually a double derivative, which means we need a squared bit on
the top along with a calligraphic L:
\begin{code}
\begin{alltt}
{\color{midgray}\glsnl{begindispmath}
\cmdname{pderiv}\marg*{\textcolor{black}{\glsni{circumchar}2 \glsni{mathcal}\marg{L}}}\marg*{}
\glsnl{enddispmath}}
\end{alltt}
\end{code}

\item The second argument is the $z_i^\rho$ squared bit.  This is a 
nested superscript
\verb|{z_i^\rho}^2|:
\begin{code}
\begin{alltt}
\color{midgray}\glsnl{begindispmath}
\cmdname{pderiv}\marg*{\textasciicircum2 \cmdname{mathcal}\marg{L}}\marg*{\textcolor{black}{\marg{z\glsni{underscorechar}i\glsni{circumchar}\glsni{rho}}\glsni{circumchar}2}}
\glsnl{enddispmath}
\end{alltt}
\end{code}

\item We can do the next partial derivative in the same way.  This 
one is slightly easier to do:
\begin{code}
{\color{midgray}\glsnl{begindispmath}\newline
\cmdname{pderiv}\marg*{\glsnl{circumchar}2 \glsnl{mathcal}{L}}\marg*{\marg*{z\glsnl{underscorechar}i\glsnl{circumchar}\glsnl{rho}}\glsnl{circumchar}2}} =\newline
-\cmdname{pderiv}\marg{\glsni{rho}\glsni{underscorechar}i}\marg{z\glsni{underscorechar}i\glsni{circumchar}\glsni{rho}}\newline
{\color{midgray}\glsnl{enddispmath}}
\end{code}

\item Delimiters also need to occur in pairs, like curly braces and 
\glsni{begin} and \glsni{end}, so let's do them next:
\begin{code}
{\color{midgray}\verb|\[|\\
\verb|\pderiv{^2 \mathcal{L}}{{z_i^\rho}^2} =|\\
\verb|-\pderiv{\rho_i}{z_i^\rho}|}\\
\strut~~\glsni{left}(\\
\strut~~\glsni{right})\\
{\color{midgray}\verb|\]|}
\end{code}\screenpagebreak

\item Now we need to do the bits inside the brackets.  First of all 
we have yet another partial derivative:
\begin{code}
{\color{midgray}\verb|\[|\\
\verb|\pderiv{^2 \mathcal{L}}{{z_i^\rho}^2} =|\\
\verb|-\pderiv{\rho_i}{z_i^\rho}|\\
\strut~~\glsnl{left}\glsnl{openparen}}\\
\strut~~~~\cmdname{pderiv}\marg{v\glsni{underscorechar}i}\marg{\glsni{rho}\glsni{underscorechar}i}\\
{\color{midgray}~~\glsnl{right}\glsnl{closeparen}\\
\verb|\]|}
\end{code}

\item Now we have a fraction following the partial derivative from
the previous step.
(Make sure you use braces for the exponential bit: \verb|\e^{v_i}|
($\e^{v_i}$) is not the same as \verb|\e^v_i| ($\e^v_i$). The
command \cmdname{e} is defined as described
\htmlref{earlier}{sec:scripts}\latex{ in \sectionref{sec:scripts}}.
Make sure you define it, preferably in the \gls{preamble}.)
\begin{code}
{\color{midgray}\verb|\[|\\
\verb|\pderiv{^2 \mathcal{L}}{{z_i^\rho}^2} =|\\
\verb|-\pderiv{\rho_i}{z_i^\rho}|\\
\strut~~\glsnl{left}\glsnl{openparen}\\
\strut~~~~\cmdname{pderiv}\marg*{v\glsnl{underscorechar}i}\marg*{\glsnl{rho}\glsnl{underscorechar}i}} \glsni{frac}\marg{\cmdname{e}\glsni{circumchar}\marg{v\glsni{underscorechar}i}}\marg{1-\cmdname{e}\glsni{circumchar}\marg{v\glsni{underscorechar}i}}\\
{\color{midgray}~~\glsnl{right}\glsnl{closeparen}\\
\verb|\]|}
\end{code}

\item This is followed by $v_i$ times another fraction:
\begin{code}\par\noindent
{\color{midgray}\verb|\[|\\
\verb|\pderiv{^2 \mathcal{L}}{{z_i^\rho}^2} =|\\
\verb|-\pderiv{\rho_i}{z_i^\rho}|\\
\strut~~\glsnl{left}\glsnl{openparen}\\
\strut~~~~\cmdname{pderiv}\marg*{v\glsnl{underscorechar}i}\marg*{\glsnl{rho}\glsnl{underscorechar}i} \glsnl{frac}\marg*{\cmdname{e}\glsnl{circumchar}\marg*{v\glsnl{underscorechar}i}}\marg*{1-\cmdname{e}\glsnl{circumchar}\marg*{v\glsnl{underscorechar}i}}}\\
\strut~~~~+ v\glsni{underscorechar}i \glsni{frac}\marg{}\marg{}\\
{\color{midgray}~~\glsnl{right}\glsnl{closeparen}\\
\color{midgray}\verb|\]|}
\end{code}

\item The bottom part of the
fraction (the denominator) is easier than the top, so let's do that first:
\begin{code}
\color{midgray}\glsnl{begindispmath}\newline
\cmdname{pderiv}\marg*{\textasciicircum2 \cmdname{mathcal}\marg*{L}}\marg*{\marg*{z\_i\textasciicircum\cmdname{rho}}\textasciicircum2} =\newline
-\cmdname{pderiv}\marg*{\cmdname{rho}\_i}\marg*{z\_i\textasciicircum\cmdname{rho}}\newline
\strut~~\cmdname{left}(\newline
\strut~~~~\cmdname{pderiv}\marg*{v\_i}\marg*{\cmdname{rho}\_i} \cmdname{frac}\marg*{\cmdname{e}\textasciicircum\marg*{v\_i}}\marg*{1-\cmdname{e}\textasciicircum\marg*{v\_i}}\newline
\strut~~~~+ v\_i \cmdname{frac}\marg*{}\marg*{\textcolor{black}{(1-\cmdname{e}\glsni{circumchar}\marg{v\glsni{underscorechar}i})\glsni{circumchar}2}}\newline
\strut~~\cmdname{right})\newline
\glsnl{enddispmath}
\color{black}
\end{code}

% Something weird is causing a shift in the 
% frame here.
% Compensating by ending and restarting enumerate
% environment
\breakenumi

\item Now for the top part of the fraction (the numerator).
To refresh your memory, it should look like:
\[
\e^{v_i}\pderiv{v_i}{\rho_i}(1-\e^{v_i})
+ \e^{2v_i}\pderiv{v_i}{\rho_i}
\]
That's a bit
complicated, so let's break it down:

\begin{enumerate}

\item The first term is:
\begin{alltt}
\cmdname{e}\glsni{circumchar}\marg{v\glsni{underscorechar}i}
\end{alltt}

\item The next term is another partial derivative:
\begin{alltt}
\cmdname{pderiv}\marg{v\glsni{underscorechar}i}\marg{\glsni{rho}\glsni{underscorechar}i}
\end{alltt}

\item Then we have:
\begin{alltt}
(1-\cmdname{e}\glsni{circumchar}\marg{v\glsni{underscorechar}i})
\end{alltt}

\item Next we have to add on:
\begin{alltt}
+\cmdname{e}\glsni{circumchar}\marg{2v\glsni{underscorechar}i}
\end{alltt}

\item And finally we have:
\begin{alltt}
\cmdname{pderiv}\marg{v\glsni{underscorechar}i}\marg{\glsni{rho}\glsni{underscorechar}i}
\end{alltt}
\end{enumerate}

So the numerator is:
\begin{code}
\begin{verbatim}
\e^{v_i}\pderiv{v_i}{\rho_i}(1-\e^{v_i})
  + \e^{2v_i}\pderiv{v_i}{\rho_i}
\end{verbatim}
\end{code}

Inserting this into our code:\reportpagebreak
\begin{code}
\begin{alltt}
\color{midgray}\glsnl{begindispmath}
\cmdname{pderiv}\marg*{\textasciicircum2\cmdname{mathcal}\marg*{L}}\marg*{\marg*{z\_i\textasciicircum\cmdname{rho}}\textasciicircum2} =
-\cmdname{pderiv}\marg*{\cmdname{rho}\_i}\marg*{z\_i\textasciicircum\cmdname{rho}}
\cmdname{left}(
  \cmdname{pderiv}\marg*{v\_i}\marg*{\cmdname{rho}\_i} \cmdname{frac}\marg*{\cmdname{e}\textasciicircum\marg*{v\_i}}\marg*{1-\cmdname{e}\textasciicircum\marg*{v\_i}}
  \color{midgray}+ v\_i \cmdname{frac}\marg*{\textcolor{black}{\cmdname{e}\textasciicircum\marg{v\_i}\cmdname{pderiv}\marg{v\_i}\marg{\cmdname{rho}\_i}(1-\cmdname{e}\textasciicircum\marg{v\_i})
             +\cmdname{e}\textasciicircum\marg{2v\_i}\cmdname{pderiv}\marg{v\_i}\marg{\cmdname{rho}\_i}}}\marg*{(1-\cmdname{e}\textasciicircum\marg*{v\_i})\textasciicircum2}
\cmdname{right})
\glsnl{enddispmath}
\end{alltt}
\color{black}
\end{code}

\end{enumerate}


%----------- Arrays -------------

\setnode{arrays}
\subsection{Arrays}
\label{sec:arrays}

Mathematical structures such as matrices and vectors require elements
to be arranged in rows and columns.  Just as we can align material in
rows and columns in text mode using the
\gls{env-tabular} environment (\sectionref{sec:tabular}), we can do the
same in maths mode using the \gls{env-array} environment.  The
\glsni{env-array} environment has the same format as the
\gls{env-tabular} environment, however it must be in maths mode.
The column half-gaps are given by the \gls{length} register
\gls{arraycolsep} (analogous to \gls{tabcolsep}). 

\xminisec{Example:}

\begin{code}
\glsni{begindispmath}\newline
\glsni{begin}\marg{array}\marg{rrr}\newline
0 \gls{ampchar} 1 \gls{ampchar} 19\gls{tab.dbbackslashchar}\newline
-6 \gls{ampchar} 10 \gls{ampchar} 200\newline
\glsni{end}\marg{array}\newline
\glsni{enddispmath}
\end{code}

\begin{result}[Image shows 3 columns and 2 rows.]
\[
\begin{array}{rrr}
0 & 1 & 19\\
-6 & 10 & 200
\end{array}
\]
\end{result}

\xminisec{Example (Adding Delimiters):}

\begin{code}
\begin{alltt}
\glsni{begindispmath}
\glsni{left}\glsni{openparen}
 \glsni{begin}\marg{array}\marg{rrr}
   0 \glsni{ampchar} 1 \glsni{ampchar} 19\glsni{tab.dbbackslashchar}
  -6 \glsni{ampchar} 10 \glsni{ampchar} 200
 \glsni{end}\marg{array}
\glsni{right}\glsni{closeparen}
\glsni{enddispmath}
\end{alltt}
\end{code}%
\begin{resultS}[Image: as before but has large round brackets on either
side.]
\[
\left(
\begin{array}{rrr}
0 & 1 & 19\\
-6 & 10 & 200
\end{array}
\right)
\]
\end{resultS}

\xminisec{Adding a Vertical Rule:}

A vertical rule can be added using \gls{array.barchar} in the column
specifier. For example:
\begin{code}
\begin{alltt}
\glsni{begindispmath}
\glsni{left}\glsni{openparen}
 \glsni{begin}\marg{array}\marg{rr\glsni{array.barchar}r}
   0 \glsni{ampchar} 1 \glsni{ampchar} 19\glsni{tab.dbbackslashchar}
  -6 \glsni{ampchar} 10 \glsni{ampchar} 200
 \glsni{end}\marg{array}
\glsni{right}\glsni{closeparen}
\glsni{enddispmath}
\end{alltt}
\end{code}%
\begin{resultS}[Image: 3 columns and 2 rows with a vertical line
between the second and third column.]
\[
\left(
\begin{array}{rr|r}
0 & 1 & 19\\
-6 & 10 & 200
\end{array}
\right)
\]
\end{resultS}

\xminisec{Example (Cases):}
This example uses an 
\htmlref{invisible delimiter}{obj:invisdelim}:
\begin{code}
\begin{alltt}
\glsni{begindispmath}
f(x) = 
 \glsni{left}\glsni{leftbrace}
  \glsni{begin}\marg{array}\marg{rl}
   -1 \glsni{ampchar} x \glsni{ltchar} 0\glsni{tab.dbbackslashchar}
   0  \glsni{ampchar} x = 0\glsni{tab.dbbackslashchar}
   +1 \glsni{ampchar} x \glsni{gtchar} 0
  \glsni{end}\marg{array}
 \glsni{right}\glsni{delimiter.periodchar}
\glsni{enddispmath}
\end{alltt}
\end{code}%
\begin{resultS}[Image: array has large open brace on the left and
nothing on the right of it.]
\[
f(x) = 
\left\{
\begin{array}{rl}
-1 & x < 0\\
0 & x = 0\\
+1 & x > 0
\end{array}
\right.
\]
\end{resultS}
This can be rewritten more compactly using the \isty{amsmath}
\gls{env-cases} environment:
\begin{code}
\gls{begindispmath}\newline
f(x) = \newline
\strut~\glsni{begin}\marg{cases}\newline
\strut~~-1 \glsni{ampchar} x < 0\glsni{tab.dbbackslashchar}\newline
\strut~~0  \glsni{ampchar} x = 0\glsni{tab.dbbackslashchar}\newline
\strut~~+1 \glsni{ampchar} x > 0\newline
\strut~\glsni{end}\marg{cases}\newline
\gls{enddispmath}
\end{code}
\begin{resultS}[Image: as previous example but first column is left
aligned.]
\[
f(x) = 
\begin{cases}
-1 & x < 0\\
0 & x = 0\\
+1 & x > 0
\end{cases}
\]
\end{resultS}

The \isty{amsmath} package provides some convenient environments to
typeset matrices: \gls{env-pmatrix}, \gls{env-bmatrix},
\gls{env-Bmatrix}, \gls{env-vmatrix} and \gls{env-Vmatrix}. These
are similar to the \gls{env-array} environment except there is no
argument, and they add (respectively) $\bigl(~\bigr)$,
$\bigl[~\bigr]$, $\bigl\{~\bigr\}$, $\bigl\lvert~\bigr\rvert$ and
$\bigl\lVert~\bigr\rVert$ delimiters. There is also the
\gls{env-matrix} environment that doesn't have any delimiters.

\xminisec{Example:}
\begin{code}
\begin{alltt}
\glsni{begin}\marg{equation}
 \glsni{begin}\marg{pmatrix}
  a \glsni{ampchar} b\glsni{tab.dbbackslashchar}
  c \glsni{ampchar} d
 \glsni{end}\marg{pmatrix}
\glsni{end}\marg{equation}
\end{alltt}
\end{code}
\begin{result}[Image: 2 by 2 matrix with round delimiters and an
equation number.]
\begin{equation}
\begin{pmatrix}
a & b\\
c & d
\end{pmatrix}
\end{equation}
\end{result}

The \isty{amsmath} package also provides the environment
\gls{env-smallmatrix} designed for in-line use. You need to add any
delimiters explicitly.

\xminisec{Example:}
\begin{code}
\begin{alltt}
Here is a small matrix
\glsni{begin}\marg{math}
 \glsni{left}\glsni{openparen}
  \glsni{begin}\marg{smallmatrix}
  a \glsni{ampchar} b\glsni{tab.dbbackslashchar}
  c \glsni{ampchar} d
  \glsni{end}\marg{smallmatrix}
 \glsni{right}\glsni{closeparen}
\glsni{end}\marg{math}
in a line of text.
\end{alltt}
\end{code}
\begin{resultS}[Image: matrix is small enough to fit into the line
of text.]
Here is a small matrix
\begin{math}
\left(
\begin{smallmatrix}
a & b\\
c & d
\end{smallmatrix}
\right)
\end{math}
in a line of text.
\end{resultS}


%----------- Vectors ------------

\setnode{vectors}
\subsection{Vectors}
\label{sec:vec}

A variable representing a vector can be typeset using the command:
\begin{definition}
\gls{vec}\marg{\meta{variable}}
\end{definition}\screenpagebreak

\xminisec{Example:}
\begin{codeS}
\glsni{begindispmath}
\glsni{vec}\marg{x}
\glsni{enddispmath}
\end{codeS}
\begin{resultS}[Image: x with a small right arrow above it.]
\[ \vec{x} \]
\end{resultS}

Vectors are often typeset in bold.  This can be 
done by \htmlref{redefining}{sec:renewcom} the \glsni{vec}
command\doifbook{ (see \sectionref{sec:renewcom})}.
You could use \gls{mathbf}, for example:
\begin{code}
\begin{alltt}
\gls{renewcommand}\marg{\glsni{vec}}\oarg{1}\marg{\glsni{mathbf}\marg{\glsni{hashchar}1}}

\glsni{begindispmath}
  \glsni{vec}\marg{x}\glsni{cdot}\glsni{vec}\marg{\glsni{xi}} = z
\glsni{enddispmath}
\end{alltt}
\end{code}\bookpagebreak\noindent
\begin{resultS}[Image: upright bold x centred dot xi equals z]
\[
\mathbf{x}\cdot\mathbf{\xi} = z
\]
\end{resultS}
\refstepcounter{object}\label{cmd:vec}However, as you may have
noticed, the 
Greek letter $\xi$ has not come out in bold.  Here's an alternative 
(using \gls{boldsymbol} defined in the \isty{amsfonts} \htmlref{package}{sec:packages}):
\begin{code}
\gls{renewcommand}\marg{\glsni{vec}}\oarg{1}\marg{\glsni{boldsymbol}\marg{\glsni{hashchar}1}}\newline
\glsni{begindispmath}\newline
\strut~~\glsni{vec}\marg{x}\glsni{cdot}\glsni{vec}\marg{\glsni{xi}} = z\newline
\glsni{enddispmath}
\end{code}%
\begin{resultS}[As before but x and xi in italic bold]
\[
\boldsymbol{x}\cdot\boldsymbol{\xi} = z
\]
\end{resultS}

Located (or position) vectors, on the other hand, are usually typeset with a right
arrow, but the default definition of \glsni{vec} produces an
arrow that is too small:
\begin{codeS}
\glsni{begindispmath}
\glsni{vec}\marg{OP}
\glsni{enddispmath}
\end{codeS}
\begin{resultS}[Image: the letters O P with a small right arrow
above them.]
\[
\vec{OP}
\]
\end{resultS}
Instead, use \gls{overrightarrow} (\tableref{tab:overunderarrows}):
\begin{codeS}
\glsni{begindispmath}
\glsni{overrightarrow}\marg{OP}
\glsni{enddispmath}
\end{codeS}
\begin{resultS}[Image: as above but with a longer arrow that spans
both letters.]
\[
\overrightarrow{OP}
\]
\end{resultS}
You might prefer to define separate commands for a located vector and a vector
variable.

\xminisec{Example:}
In the \glsterm{preamble}, define \cmdname{lvec} for a
located vector and \cmdname{bvec} for a vector variable:
\begin{code}
\begin{alltt}
\gls{newcommand}*\marg{\cmdname{lvec}}\oarg{1}\marg{\glsni{overrightarrow}\marg{\glsni{hashchar}1}}
\glsni{newcommand}*\marg{\cmdname{bvec}}\oarg{1}\marg{\glsni{boldsymbol}\marg{\glsni{hashchar}1}}
\end{alltt}
\end{code}
Later in the document:
\begin{code}
Let \glsni{dollarchar}\cmdname{bvec}\marg{u}=(x, y)\glsni{dollarchar}
represent \glsni{dollarchar}\cmdname{lvec}\marg{OP}\glsni{dollarchar}, then\newline
\glsni{begindispmath}\newline
\strut~~\glsni{lVert} \cmdname{bvec}\marg{u} \glsni{rVert}
    = \glsni{sqrt}\marg{x\glsni{circumchar}2 + y\glsni{circumchar}2}\newline
\glsni{enddispmath}
\end{code}
\begin{result}[Image: the vector u appeas in italic bold and the
letters O P are as above.]
\newcommand*{\lvec}[1]{\overrightarrow{#1}}%
\newcommand*{\bvec}[1]{\boldsymbol{#1}}%
Let $\bvec{u}=(x, y)$ represent $\lvec{OP}$, then
\[
  \lVert \bvec{u} \rVert = \sqrt{x^2 + y^2}
\]
\end{result}
\bookpagebreak

\begin{exercise}{Maths: Vectors and Arrays}{ex:vectors}
Try to produce the following:
\renewcommand{\vec}[1]{\boldsymbol{#1}}
\begin{result}[maths2ex.html]
\[
\boldsymbol{A}\vec{x} =
\begin{pmatrix}
0 & 1\\
2 & 3
\end{pmatrix}
\begin{pmatrix}
1\\
2
\end{pmatrix}
=
\begin{pmatrix}
2\\
8
\end{pmatrix}
= \vec{y}
\]
\end{result}%
As before, you can \downloadorview{vectors} the solution.
\end{exercise}

%--------------- Mathematical Spacing -------------------

\setnode{mathssp}
\subsection{Mathematical Spacing}
\label{sec:mathsspacing}

\LaTeX\ deals with mathematical spacing fairly well, but sometimes
you may find you want to adjust the spacing yourself.  Available
spacing commands are listed in \tableref{tab:spacing}.

\begin{table}[hbtp]
\caption[Mathematical Spacing Commands]{Mathematical Spacing Commands
(\supdag provided by \protect\istyend{amsmath})}
\label{tab:spacing}
\centering
\begin{tabular}{lll}
\toprule
\bfseries Command & \bfseries Example Input & \bfseries Example Output\\
\midrule
 & \verb|$AB$| & $AB$ \\
\gls{thinspace} or \gls{comma} & \verb|$A\,B$| & $A\,B$\\
\gls{medspace}\supdag\ or \gls{colon} & \verb|$A\:B$| & $A\:B$\\
\gls{thickspace}\supdag\ or \gls{semicolon} & \verb|$A\;B$| & $A\;B$\\
\gls{quad} & \verb|$A\quad B$| & $A\quad B$\\
\gls{qquad} & \verb|$A\qquad B$| & $A\qquad B$\\
\gls{negthinspace} or \gls{exclam} & \verb|$A\!B$| & $A\!B$ \\
\gls{negmedspace}\supdag & \verb|$A\negmedspace B$| & $A\negmedspace B$ \\
\gls{negthickspace}\supdag & \verb|$A\negthickspace B$| & $A\negthickspace B$ \\
\bottomrule
\end{tabular}
\end{table}

\begin{exercise}{More Mathematics}{ex:moremaths}
This exercise uses the spacing command \glsni{qquad}. In addition,
it has a
\htmlref{function name}{sec:declaremathop}, diag, and it uses the
\gls{forall} and \htmlref{ellipses}{sec:ellipses} symbols. It also
\htmlref{redefines the \cmdname{vec}
command}{cmd:vec}\indexCom{vec}, as was done in the previous
section, uses the
\gls{env-bmatrix} environment (see \sectionref{sec:arrays}), and has \htmlref{subscripts and superscripts}{sec:scripts}.

Try to reproduce the following output:
\begin{result}[maths3ex.html]
\renewcommand{\vec}[1]{\boldsymbol{#1}}
The set of linear equations:
\[
a_ix_i = b_i  \qquad \forall i=1, \ldots, n
\]
can be written as a matrix equation:
\[
\diag(\vec{a})\vec{x} = \vec{b}
\]
where $\vec{x} = (x_1, \ldots, x_n)^T$, $\vec{b} = (b_1, \ldots, b_n)^T$
and
\[
\diag(\vec{a}) = 
\begin{bmatrix}
a_1 & 0 & \cdots & 0\\
0 & a_2 & \ddots & \vdots\\
\vdots & \ddots & \ddots & 0\\
0 & \cdots & 0 & a_n
\end{bmatrix}
\]
\end{result}

Again, you can \downloadorview{moremaths} the solution.
\end{exercise}

%%%%%%%%%%%%%%%%%% DEFINING ENVIRONMENTS %%%%%%%%%%%%%%%%

\setnode{newenv}
\chapter{Defining Environments}
\label{ch:newenv}

Just as you can \htmlref{define new commands}{ch:newcom}, you can 
also define new \glspl{environment}.  The command
\begin{definition}
\gls{newenvironment}\marg{\meta{env-name}}\oarg{\meta{n-args}}\oarg{\meta{default}}\marg{\meta{begin-code}}\marg{\meta{end-code}}
\end{definition}%
is used to define a new environment. As with new commands, you can
use the optional argument \meta{n-args} to define an environment with
arguments, and \meta{default} to define an environment with an
optional argument.

The first argument \meta{env-name} is the name of your new
environment.  Remember that the environment name must not have a
backslash.  The mandatory arguments \meta{begin-code} and
\meta{end-code} indicate what \LaTeX\ \reportlinebreak should do at the beginning and
end of the environment. Note that although \meta{begin-code} can
reference the arguments using \gls{hashchar}\texttt{1} etc, the
\meta{end-code} part can't.

\xminisec{Example (An Exercise Environment):}

Let's first consider an example of an
environment without any arguments.  Let's make an environment called,
say, \envname{exercise} that prints \textbf{Exercise} in bold and
typesets the contents of the environment in italic, with a gap
between the title and the contents.  In other words,
we want the following code:
\begin{code}
\glsni{begin}\marg{exercise}\newline
This is a sample.\newline
\glsni{end}\marg{exercise}
\end{code}
to produce the following output:
\begin{result}[newenv1.html]
\par\noindent
\textbf{Exercise}
\begin{itshape}
\par
\vspace{\baselineskip}%
\noindent
This is a sample.
\end{itshape}
\end{result}
(In the \htmlref{next chapter}{ch:counters} we will add numbering.)

Let's first consider what we want this environment to do: we can get
the word \dq{Exercise} in bold using \gls{textbf},
and the italic font can be
obtained by using the \gls{env-itshape} environment (recall
\sectionref{sec:fonts}).  
So, at the start of our
new environment we need
\begin{codeS}
\glsni{textbf}\marg{Exercise}\glsni{begin}\marg{itshape}
\end{codeS}
and at the end of our new
environment we need to end the \glsni{env-itshape} environment:
\begin{codeS}
\glsni{end}\marg{itshape}
\end{codeS}
Putting the above together into the new environment definition:
\begin{code}\obeyspaces
\glsni{newenvironment}\marg{exercise}\glsni{percentchar} environment name\newline
\marg{\glsni{percentchar} begin code\newline
\strut~~\glsni{textbf}\marg{Exercise}\glsni{begin}\marg{itshape}\glsni{percentchar}\newline
}\glsni{percentchar}\newline
\marg{\glsni{end}\marg{itshape}}\glsni{percentchar} end code
\end{code}
Let's try it out:
\begin{code}
\begin{verbatim}
\begin{exercise}
This is a sample.
\end{exercise}
\end{verbatim}
\end{code}%
\begin{resultS}[newenv2.html]
\textbf{Exercise}
\begin{itshape}
This is a sample.
\end{itshape}
\end{resultS}
Not quite right.  Let's put a paragraph break after
\textbf{Exercise}, and put one before it as well. The command
\gls{par} can be used to make a paragraph break and the extra bit of
vertical spacing can be produced using \glsni{vspace}. The
\gls{length} \gls{baselineskip} is the interline spacing.
Modifications are shown in bold \modification{like this}.
\begin{code}\obeyspaces
\glsnl{newenvironment}\marg*{exercise}\glsnl{percentchar} environment name\newline
\marg*{\glsnl{percentchar} begin code\newline
\strut~~\modification{\glsni{par}\glsni{vspace}\marg{\glsni{baselineskip}}\glsni{percentchar}}\newline
\strut~~\glsnl{textbf}\marg*{Exercise}\glsnl{begin}\marg*{itshape}\glsnl{percentchar}\newline
\strut~~\modification{\glsni{par}\glsni{vspace}\marg{\glsni{baselineskip}}\glsni{percentchar}}\newline
}\glsnl{percentchar}\newline
\marg*{\glsnl{end}\marg*{itshape}}\glsnl{percentchar} end code
\end{code}%
Let's have a look at the output now:
\begin{result}
\dohtmlcolorbox{%
\par\vspace{\baselineskip}%
\textbf{Exercise}\begin{itshape}%
\par\vspace{\baselineskip}
This is a sample.
\end{itshape}
}
\end{result}
\faq{There's a space added after my environment}{spaftend}%
The indent at the start of each line is caused by the normal
paragraph indentation. This can be suppressed using \gls{noindent}.
It's also a good idea to suppress
any spaces immediately following \verb|\begin{exercise}| and \verb|\end{exercise}|, which can be done
using \gls{ignorespaces} and \reportlinebreak\screenlinebreak\gls{ignorespacesafterend}.
Modifications are again shown in bold \modification{like this}.
\begin{code}\obeyspaces
\glsnl{newenvironment}\marg*{exercise}\glsnl{percentchar} environment name\newline
\marg*{\glsnl{percentchar} begin code\newline
\strut~~\glsnl{par}\glsnl{vspace}\marg*{\glsnl{baselineskip}}\modification{\glsni{noindent}}\newline
\strut~~\glsnl{textbf}\marg*{Exercise}\glsnl{begin}\marg*{itshape}\glsnl{percentchar}\newline
\strut~~\glsnl{par}\glsnl{vspace}\marg*{\glsnl{baselineskip}}\modification{\glsni{noindent}\glsni{ignorespaces}}\newline
}\glsnl{percentchar}\newline
\marg*{\glsnl{percentchar} end code\newline
\strut~~\glsnl{end}\marg{itshape}\modification{\glsni{ignorespacesafterend}}\newline 
}
\end{code}%
The \envname{exercise} environment now appears as:
\begin{result}
\par\vspace{\baselineskip}\noindent
\textbf{Exercise}\begin{itshape}%
\par\vspace{\baselineskip}\noindent
This is a sample.
\end{itshape}
\end{result}

Now let's modify our code so that the environment takes an
argument.  The argument should indicate the exercise topic.
For example, the following code:
\begin{code}
\begin{verbatim}
\begin{exercise}{An Example}
This is a sample.
\end{exercise}
\end{verbatim}
\end{code}%
should produce the following result:
\begin{result}[newenv3.html]
\par\vspace{\baselineskip}\noindent\textbf{Exercise (An Example)}
\begin{itshape}\par\vspace{\baselineskip}\noindent
This is a sample.
\end{itshape}
\end{result}%
As with \gls{newcommand}, \gls{hashchar}\texttt{1} is
used to indicate the first argument.  We can now modify the code as
follows:
\begin{code}\obeyspaces
\glsnl{newenvironment}\marg*{exercise}\modification{\oarg{1}}\glsnl{percentchar} environment name\newline
\marg*{\glsnl{percentchar} begin code\newline
\strut~~\glsnl{par}\glsnl{vspace}\marg*{\glsnl{baselineskip}}\glsnl{noindent}\newline
\strut~~\glsnl{textbf}\marg*{Exercise\modification{ (\glsni{hashchar}1)}}\glsnl{begin}\marg*{itshape}\glsnl{percentchar}\newline
\strut~~\glsnl{par}\glsnl{vspace}\marg*{\glsnl{baselineskip}}\glsnl{noindent}\glsnl{ignorespaces}\newline
}\glsnl{percentchar}\newline
\marg*{\glsnl{percentchar} end code\newline
\strut~~\glsnl{end}\marg*{itshape}\glsnl{ignorespacesafterend}\newline
}
\end{code}

\setnode{renewenv}
\section{Redefining Environments}
\label{sec:renewenv}

It is also possible to redefine an environment using:
\begin{definition}[\linewidth]
\gls{renewenvironment}\marg{\meta{env-name}}\oarg{\meta{n-args}}\oarg{\meta{default}}\marg{\meta{begin-code}}\marg{\meta{end-code}}
\end{definition}%
As with \gls{renewcommand}, only redefine an existing environment if
you want a modified version of that environment rather than because
you like the environment name.

\begin{exercise}{Defining a New Environment}{ex:newenv}
If you did any of the exercises from \exerciseref{ex:section}
to \exerciseref{ex:tables}, go back to the document you
created and define the \envname{exercise} environment as in
the example above.  Then try creating some exercises using this
environment. You could, maybe, put an exercise in the first
chapter, and then another one in the second chapter.
Again you can \downloadorview{newenv} an example.
\end{exercise}

%%%%%%%%%%%%%%%%% Counters %%%%%%%%%%%%%%%%%%%%%%%%%

\setnode{counters}
\chapter{Counters}
\label{ch:counters}

As we have seen, \LaTeX\ automatically generates numbers for
chapters, sections, equations etc. These numbers are stored in
\keyword{counters}. The names of these counters are usually the same
as the name of the object with which it is associated but without
any backslash. For example, the \gls{chapter} command has an associated
counter called \icounter{chapter}, the \gls{footnote} command has an
associated counter called \icounter{footnote}, the
\gls{env-equation} environment has an associated counter called
\icounter{equation}, the \gls{env-figure} environment has an associated counter
called \icounter{figure} and the \gls{env-table} environment has an associated
counter called \icounter{table}. There is also a counter called
\icounter{page} that keeps track of the current page number.

The value of a counter can be displayed using the command
\begin{definition}
\icmdname{the}\meta{counter}
\end{definition}%
where \meta{counter} is the name of the associated counter.  Note
that \meta{counter} does not go in curly braces and adjoins
\cmdname{the} (for example, \gls{thepage}\faq{Page number is wrong at
start of page}{wrongpn}, \reportlinebreak\gls{thesection} or
\gls{thechapter}).   In fact, we have already encountered
\reportlinebreak\screenlinebreak\gls{thefigure} in \sectionref{sec:subfloats}.

\xminisec{Example:}
\begin{code}
\begin{alltt}
This page is Page\glsni{tildechar}\glsni{thepage}.
The current chapter is Chapter\glsni{tildechar}\glsni{thechapter}.
\end{alltt}
\end{code}\label{example:thepage}%
\begin{resultS}[This page is Page 226. The current chapter is Chapter 11.]
% LaTeX2HTML doesn't store the page number in the image file
% so use \pageref to fake \thepage instead
This page is Page~\pageref*{example:thepage}.
The current chapter is Chapter~\thechapter.
\end{resultS}%
New counters can be created using the command:
\begin{definition}
\gls{newcounter}\marg{\meta{counter-name}}\oarg{\meta{outer-counter}}
\end{definition}%
The \gls{mandatory} \meta{counter-name}
is the name of your new counter (no backslash in the name).  For
example, let's define a counter called \texttt{exercise} to keep
track of each exercise. (Recall the exercise example from
\chapterref{ch:newenv}.)
\begin{codeS}
\glsni{newcounter}\marg{exercise}
\end{codeS}%
We can now display the value of the counter using
the command \cmdname{theexercise}.  At the moment the counter
has the value zero, the value can be changed using one
of the following commands:

\begin{fwlist}{\cmdname{addtocounter}\marg{\meta{counter}}\marg{\meta{num}}}
\fwitem{\gls{stepcounter}\marg{\meta{counter}}} Increments 
\meta{counter} by 1

\fwitem{\gls{refstepcounter}\marg{\meta{counter}}} As above, but 
allows you to cross-\reportlinebreak reference the counter using \gls{label} and 
\gls{ref}

\fwitem{\gls{setcounter}\marg{\meta{counter}}\marg{\meta{num}}} 
Sets the counter to \meta{num}

\fwitem{\gls{addtocounter}\marg{\meta{counter}}\marg{\meta{num}}} 
Adds \meta{num} to \meta{counter}
\end{fwlist}

A couple of the commands above take a number \meta{num} as one of 
the arguments.  If you want to use another counter for this argument, 
you need to use
\begin{definition}
\gls{value}\marg{\meta{counter}}
\end{definition}%
For example, if you want to set our new \texttt{exercise} counter to 
the same value as the \texttt{page} counter, you would do
\begin{codeS}
\glsni{setcounter}\marg{exercise}\marg{\glsni{value}\marg{page}}
\end{codeS}%
Let's go back to the \envname{exercise} environment you
created in \exerciseref{ex:newenv}.  The exercises really ought to have
an associated number, and this number should be incremented each time 
we use the exercise environment.  So let's modify our code to do this.
Modifications are illustrated in bold \modification{like this}:
\begin{code}\obeyspaces
\modification{\cmdname{newcounter}\marg{exercise}}\newline
\strut\newline
\glsni{newenvironment}\marg{exercise}\oarg{1}\glsni{percentchar} environment name\newline
\marg{\glsni{percentchar} begin code\newline
\strut~~\glsni{par}\glsni{vspace}\marg{\glsni{baselineskip}}\glsni{noindent}\newline
\strut~~\modification{\glsni{refstepcounter}\marg{exercise}}\glsni{percentchar}\newline
\strut~~\glsni{textbf}\marg{Exercise \modification{\cmdname{theexercise}\glsni{spacesym}}(\glsni{hashchar}1)}\glsni{percentchar}\newline
\strut~~\glsni{begin}\marg{itshape}\glsni{percentchar}\newline
\strut~~\glsni{par}\glsni{vspace}\marg{\glsni{baselineskip}}\glsni{percentchar}\newline
\strut~~\glsni{noindent}\glsni{ignorespaces}\newline
}\glsni{percentchar}\newline
\marg{\glsni{percentchar} end code\newline
\strut~~\glsni{end}\marg{itshape}\glsni{percentchar}\newline
\strut~~\modification{\glsni{par}\glsni{vspace}\marg{\glsni{baselineskip}}\glsni{percentchar}}\newline
\strut~~\modification{\glsni{noindent}}\glsni{ignorespacesafterend}\newline
}
\end{code}%
Note that the counter needs to be incremented before it is used.  
I've also added an extra \gls{vspace} at the end of the
environment and a paragraph break.
Since we've used \glsni{refstepcounter} instead of 
\glsni{stepcounter} we can cross-reference our
\envname{exercise} environment:\reportpagebreak
\begin{code}
\begin{alltt}
Exercise\glsni{tildechar}\glsni{ref}\marg{ex:simple} is a simple exercise.

\glsni{begin}\marg{exercise}\marg{Simple Exercise}
\glsni{label}\marg{ex:simple}\glsni{percentchar}
This is a simple exercise.
\glsni{end}\marg{exercise}
\end{alltt}
\end{code}%
This produces 
the following output:
\begin{result}[newcounter.html]
Exercise~1 is a simple exercise.
\par\vspace{\baselineskip}\noindent\textbf{Exercise 1 (Simple Exercise)}%
\begin{itshape}\par\vspace{\baselineskip}\noindent
This is a simple exercise.
\end{itshape}
\end{result}

The counter representation can be changed by redefining
\cmdname{theexercise}\faq{Redefining counters'
\cmdname{the}-commands}{the-commands} using the \gls{renewcommand}
command described in \sectionref{sec:renewcom}.
The following commands can be used to display the counter:

\begin{fwlist}{\cmdname{fnsymbol}\marg{\meta{counter}}}
\fwitem{\gls{arabic}\marg{\meta{counter}}} Arabic numeral (1, 2, 3, \ldots)

\fwitem{\gls{Roman}\marg{\meta{counter}}} Upper case Roman numeral 
(I, II, III, \ldots)

\fwitem{\gls{roman}\marg{\meta{counter}}} Lower case Roman numeral 
(i, ii, iii, \ldots)

\fwitem{\gls{alph}\marg{\meta{counter}}} Lower case letter 
(a, b, c, \ldots, z)

\fwitem{\gls{Alph}\marg{\meta{counter}}} Upper case letter 
(A, B, C, \ldots, Z)

\fwitem{\gls{fnsymbol}\marg{\meta{counter}}}
A footnote symbol 
(\makeimg{asterisk dagger double dagger section mark
paragraph mark double bar double asterisk two single daggers
two double daggers}{\footnotesymbols})
\end{fwlist}

\xminisec{Example:}

To make the chapter numbers appear as upper case Roman 
numerals you would do:
\begin{codeS}
\gls{renewcommand}\marg{\gls{thechapter}}\marg{\glsni{Roman}\marg{chapter}}
\end{codeS}

You may have noticed that \gls{newcounter} has an optional
argument \meta{outer-counter}.  This is for use if you require the
new counter to be reset every time \meta{outer-counter} is
incremented\faq{Master and slave counters}{addtoreset}.  For example,
the section numbers in the \icls{scrbook} class are dependent on the
chapter numbers.  Each time a new chapter is started, the section
numbers are reset.  Suppose we want our \texttt{exercise} counter to
be dependent on the \icounter{chapter} counter, we would do
\begin{codeS}
\gls{newcounter}\marg{exercise}\oarg{chapter}
\end{codeS}

Note that if you make a counter dependent on another counter like
this, the default action of \cmdname{the}\meta{counter} remains the
same, so \cmdname{theexercise} won't print the chapter number.
To make the chapter number appear as well, we need to redefine
\cmdname{theexercise} (recall \sectionref{sec:renewcom}):
\begin{codeS}
\gls{renewcommand}\marg{\cmdname{theexercise}}\marg{\glsni{thechapter}.\glsni{arabic}\marg{exercise}}
\end{codeS}%
Notice the use of \gls{thechapter} instead of,
say, \verb|\arabic{chapter}|.  This way
we don't need to keep track of the \icounter{chapter} counter format.

\xminisec{Example (Footnote Markers):}

The \icounter{footnote} counter is reset at the start of each
chapter but by default the chapter number isn't displayed in
\gls{thefootnote}. In \latexhtml{this book}{the PDF version of
this document} \glsni{thefootnote} was redefined so that it
displays the chapter number:
\begin{codeS}
\gls{renewcommand}\marg{\glsni{thefootnote}}\marg{\glsni{thechapter}.\glsni{arabic}\marg{footnote}}
\end{codeS}

\begin{exercise}{Using Counters}{ex:counters}
Modify the document from \exerciseref{ex:newenv} so that
the \envname{exercise} environment has a counter.
Make the counter dependent on the chapter.
You can \downloadorview{counters} an example.
\end{exercise}

%%%%%%%%%%%%%%%%%%%%%% APPENDICES %%%%%%%%%%%%%%%%%%

\appendix

%%%%%%%%%%%%%%%%%%% Installing Packages %%%%%%%%%%%%%%

\setnode{installsty}
\chapter{Downloading and Installing Packages}
\label{ch:installsty}

New \LaTeX\ packages are being created all the time\faq{Installing things
on a (La)TeX system}{instpackages}, so you may find that there are some
packages that you don't have on your installation. In this case, if
you don't have the package you want, you can download
it\faq{Installation using MiKTeX package manager}{inst-miktex*} from
\gls{ctan}. Before discussing installing new packages, it is
a good idea for you to understand the \gls{tds}\faq{What is the TDS?}{tds}.

All the files that make up the \TeX\ distribution are stored in a
standard hierarchical structure. The root directory of the main
distribution is usually called \texttt{texmf} or \texttt{texmf-dist}. Its location
depends on your system. For example, if you are using
\itexdistro{TeX~Live}~2012 on UNIX/Linux, it will probably be located in
\texttt{/usr\slash local\slash texlive\slash 2012\slash texmf-dist} or if you are using \itexdistro{MiKTeX} it
may be located in \verb|c:\texmf| or \texttt{c:\backslashsym Program
Files\backslashsym texmf}. Whichever system you are using, I~shall
refer to this directory as \meta{TEXMF}. So, if you are using
\texdistro{TeX~Live}~2012, \meta{TEXMF}\texttt{\slash doc} refers to the directory
\texttt{/usr\slash local\slash texlive\slash 2012\slash texmf-dist\slash doc}, or if you are using
\texdistro{MiKTeX}, \meta{TEXMF}\verb|\doc| refers to the folder
\verb|c:\texmf\doc| or \texttt{c:\backslashsym Program
Files\backslashsym texmf\backslashsym doc}.
In general, you should not make any modifications to the
\meta{TEXMF} directory tree as it will get overridden whenever you
update your \TeX\ distribution.

You should also have a local texmf tree. Again, the location of the
local texmf tree depends on your system. If you are using
\texdistro{TeX~Live}, it may be 
\texttt{/usr\slash local\slash texlive\slash texmf-local}. 
If you are using \texdistro{MiKTeX}, it
may be \texttt{c:\backslashsym localtexmf} or
\texttt{c:\backslashsym Program Files\backslashsym localtexmf}.
Whichever system you are using, I~shall refer to this directory as
\meta{TEXMF-LOCAL}. There is also the \meta{TEXMF-HOME} directory.
On UNIX-like systems this is usually \verb|~/texmf|. On Windows it's
usually in your user folder. This is the one where you typically
install any new classes or packages.

These directories must all have the same structure. The principle
sub-directories relating to \LaTeX\ are illustrated in
\figureref{fig:tds}. It may be that your \meta{TEXMF-HOME} directory
doesn't exist or doesn't contain some of these sub-directories, if
so, you will need to create them.

You can use the \iappname{kpsewhich} application to find out the 
locations of \meta{TEXMF-LOCAL} and
\meta{TEXMF-HOME}. Since \appname{kpsewhich} is a command-line
application, you will need a \htmlref{command prompt or terminal}{sec:terminal} 
open\latex{ (see \sectionref{sec:terminal})}.
At the command prompt, type
\begin{verbatim}
kpsewhich -var-value=TEXMFHOME
\end{verbatim}
to display the location of \meta{TEXMF-HOME} or
\begin{verbatim}
kpsewhich -var-value=TEXMFLOCAL
\end{verbatim}
to display the location of \meta{TEXMF-LOCAL}.
(Remember to press the enter \enter\ key at the end of the line.)

\refstepcounter{object}\label{obj:docdir}%
The documentation for \LaTeX\ classes and packages can be found
in the \texttt{doc\slash latex} sub-directories: \meta{TEXMF}\texttt{\slash doc\slash latex},
\meta{TEXMF-LOCAL}\texttt{\slash doc\slash latex} and \meta{TEXMF-HOME}\texttt{\slash doc\slash latex}.

\begin{figure}[htbp]
\figconts
  {pictures/tds}
  {%
    \caption{The \protect\TeX\protect\space Directory Structure (TDS) Showing 
    the Main \protect\LaTeX-Related Sub-Directories.}%
  }%
  {fig:tds}
\end{figure}

Some packages are supplied in this format.\footnote{Complete list at \url{http://mirror.ctan.org/install/macros/latex/contrib/}.} For example, the package
\sty{sample-package} may be distributed in a compressed file \texttt{sample-package.tds.zip},
which contains the files
\begin{verbatim}
doc/latex/sample-package/sample-package.pdf
tex/latex/sample-package/sample-package.sty
tex/latex/sample-package/sample-foo.sty
tex/latex/sample-package/sample-bar.sty
\end{verbatim}
In this case all you need to do is decompress the contents of 
the archive into the
\meta{TEXMF-LOCAL} or \meta{TEXMF-HOME} directory.

\refstepcounter{object}\label{obj:runtexhash}%
On older \TeX-distributions, you would then need to \htmlref{refresh the
\TeX\ database}{sec:refresh}\latex{ (described in
\sectionref{sec:refresh}\ifscreenorother{}{\ on
page~\pageref{sec:refresh}})}. With new distributions, you don't
need to do this if you are installing a new package into your
\meta{TEXMF-HOME} directory.

\xminisec{Example (Unix-Like):}
To install \texttt{sample-package.tds.zip} (assuming you're in the
same directory as that file):
\begin{verbatim}
unzip -d ~/texmf sample-package.tds.zip
\end{verbatim}

\setnode{dtxins}
\section{DTX and INS Files}
\label{sec:dtx}
\faq{Documented LaTeX sources (.dtx files)}{dtx}Not all packages are
provided in the \gls{tds} format. Instead (or additionally) many are supplied
with the code and documentation all bundled together in one file.
This file usually has the extension \texttt{.dtx}, and it usually
comes with an installation script that has the extension
\texttt{.ins}. Once you have downloaded the \texttt{.dtx} and
\texttt{.ins} files, you will then have to extract the code before
you can use it. Let's go back to the previous example. The package
\sty{sample-package} is now distributed in a DTX file, so the
\texttt{sample-package.zip} archive contains the files
\begin{verbatim}
sample-package.dtx sample-package.ins
\end{verbatim}
(with hopefully a \texttt{README} or \texttt{INSTALL} file). Note
that this archive, unlike the TDS one, doesn't contain any
\texttt{.sty} files. The documentation source and the package code
(\texttt{sample-package.sty}, \texttt{sample-foo.sty} and
\texttt{sample-bar.sty}) are all contained in the file
\texttt{sample-package.dtx}. This is how to extract them:
\begin{enumerate}
\item Extract the contents of \texttt{sample-package.zip} to a
temporary directory.

\item\label{itm:ins} Run \LaTeX\ on the file \texttt{sample-package.ins}. If
you are using a \latexhtml{terminal}{\htmlref{terminal}{sec:terminal}}, you can type the
following at the command prompt:
\begin{verbatim}
latex sample-package.ins
\end{verbatim}
If you are using a front-end, such as TeXWorks, open the 
\texttt{.ins} file
(for example \texttt{sample-package.ins}), and click on the build/typeset button.

This will create the files containing the package code. In this
example it will create the main package file \texttt{sample-package.sty}
and supplementary packages \texttt{sample-foo.sty} and \texttt{sample-bar.sty}.

\item Make a sub-directory of \meta{TEXMF-LOCAL}\texttt{/tex/latex}%
\footnote{or \meta{TEXMF-LOCAL}\texttt{\backslashsym tex\backslashsym latex} on 
Windows} in which to place these files. In this example, the package is
called \dq{sample-package}, so make a sub-directory called \texttt{sample-package}.

\item Move the files created in \objectref{Step}{itm:ins} into
the new sub-directory you created in the previous step.

\item\label{itm:dtx} Run \iPDFLaTeX\ on the file \texttt{sample-package.dtx}.
(The same as in \objectref{Step}{itm:ins}, but use the file
\texttt{sample-package.dtx} instead of \texttt{sample-package.ins}.) This will 
create a file called \texttt{sample-package.pdf}. You may need to repeat
this step to ensure that the cross references are up-to-date.
Check the \texttt{README} file or \texttt{INSTALL} file to see
if there is anything else you need to do. (If you have downloaded
the package from CTAN, it's possible that the documentation has 
already been supplied, as package authors are encouraged to supply
a PDF version of the documentation for on-line viewing. If so,
you can omit this step.)

\item  Make a sub-directory of \meta{TEXMF-LOCAL}\texttt{/doc/latex}%
\footnote{or \meta{TEXMF-LOCAL}\texttt{\backslashsym doc\backslashsym latex} on 
Windows} in which to place the documentation. In this example, the package is
called \dq{sample-package}, so make a sub-directory called
\texttt{sample-package}.

\item Move the files created in \objectref{Step}{itm:dtx} into
the new sub-directory you created in the previous step.

\end{enumerate}

As mentioned \htmlref{above}{obj:runtexhash}, on older
\TeX-distributions, you would then need to \htmlref{refresh the
\TeX\ database}{sec:refresh}, but this isn't required for
\meta{TEXMF-HOME} installs on new distributions.

\setnode{texhash}
\section{Refreshing the \texorpdfstring{\protect\TeX}{TeX}\protect\space Database}
\label{sec:refresh}

On older \gls{tex} distributions you had to refresh the \TeX\
database whenever you installed new classes or packages. With newer
installations you don't need to do this if you install them in your
\meta{TEXMF-HOME} directory, except under certain circumstances (for
example, you're using using a networked drive). If it turns out that
\TeX\ can't find a new class or package you have installed in
\meta{TEXMF-HOME} you will need to update the database using the
\iappname{texhash} (or \iappname{mktexlsr}) application. This is a
command-line application, so you need a \htmlref{terminal or command
prompt}{sec:terminal}\latex{ (see \sectionref{sec:terminal})}.

For example, on UNIX/Linux, to update \meta{TEXMF-HOME} (the
directory \verb|~/texmf|) you need to type the following at the
command prompt:
\begin{verbatim}
texhash ~/texmf
\end{verbatim}

If you are using a modern \TeX\ distribution, such as
\itexdistro{MiKTeX}, \itexdistro{TeX~Live} or \itexdistro{MacTeX} there 
should be a package manager that has a package installation
and refresh facility. For example, \texdistro{TeX~Live} comes with
the TeX~Live Manager (\iappname{tlmgr} or \iappname{mactlmgr}) and
recent versions of \texdistro{MiKTeX} have an application called
\appname{MiKTeX Update Wizard} which can automatically download and
install known packages.

If you experience any problems, contact your system administrator for help
or try one of the resources listed in \appendixref{ch:help}.

\xminisec{Related \gls{ukfaq} topics:}
\begin{itemize}
\item \faqlink{Installing things on a (La)TeX system}{instpackages}
\item \faqlink{Installing files \dq{where (La)TeX can find them}}{wherefiles}
\item \faqlink{Installation using MiKTeX package manager}{inst-miktex*}
\item \faqlink{\dq{Temporary} installation of (La)TeX files}{tempinst}
\item \faqlink{\dq{Private} installations of files}{privinst}
\end{itemize}


%%%%%%%%%%%%%%%%%%%%%% ERRORS %%%%%%%%%%%%%%%%%%%%%%%%%%%

\setnode{commonerrors}
\chapter{Common Errors}
\label{ch:errors}

\begin{itemize}
\item
If you're running \LaTeX\ from a \gls{terminal} and the only message that gets displayed is:
\begin{verbatim}
latex: Command not found.
\end{verbatim}
or
\begin{verbatim}
Bad command or file name
\end{verbatim}
then you have either mistyped the command name, or you
don't have \LaTeX\ installed on your computer, or your path hasn't
been set up correctly.  First check that you have typed the command
correctly, then check to see if you have \TeX\ installed.  Failing
that, contact your system administrator for help or try one of the
resources listed in \appendixref{ch:help}.

\item
If you're running \LaTeX\ from a \gls{terminal} and you get the message (or something similar):
\begin{verbatim}
This is TeX, Version 3.14159 (Web2C 7.3.1)
! I can't find file `sample'.
<*> sample
          
Please type another input file name: 
\end{verbatim}
then you have either misspelt the filename or you are in the wrong
directory.  If you have misspelt the filename, simply type in the
correct name at the prompt.  If you are in the wrong directory or you
want to quit, type \texttt{X} followed by the return character \enter.
To check you are in
the right directory, on a Unix-like system you can type:
\begin{verbatim}
ls
\end{verbatim}
This will list the contents of
the directory.  If you are certain that you have spelt the filename
correctly and that you are in the right directory, there may be
something wrong with your path, in which case contact your system
administrator.

\item
Error messages will usually look something like:
\begin{verbatim}
! Undefined control sequence.
l.1 \docmentclass
                 [12pt]{scrartcl}
? 
\end{verbatim}
The first line is the error message.  In this example I~have misspelt
the command \cmdname{documentclass}.  The next line begins with
\texttt{l.}\ followed by a number.  This is the line number in the
source code where the error occurred.  In this case the error
occurred on line~1.  Following the line number is the input line
\LaTeX\ has processed so far, and staggered on the next line is the
remainder of the input line.

Here's another example.  Suppose line~8 of my 
\gls{source}
looks like:
\begin{verbatim}
The date today is: \toady, which is nice to know.
\end{verbatim}
The error in this case is the misspelling of the command 
\cmdname{today}.  The error message will appear as follows:

\begin{latexonly}
\noindent
\begin{tabular}{l@{}l}
\multicolumn{2}{l}{%
\texttt{!\ Undefined control sequence.}\rlap{\color{blue}%
\tikz[baseline=-1ex]{\draw[latex-] (0em,0ex) -- (4em,2ex);
\pgftext[left,center,at={\pgfpoint{4em}{2ex}}]{Error Message}
}}}\\
\texttt{l.}\smash{\llap{\color{blue}%
\tikz[baseline]{\draw[latex-] (0em,0ex) -- (-3em,-4ex);
\pgftext[right,top,at={\pgfpoint{-3em}{-3ex}}]{\shortstack{Line\\Number}}}}}%
\texttt{8 The date today is }%
\smash{\llap{\color{blue}%
\tikz[baseline]{\draw[latex-] (0em,0ex) -- (-3em,-3ex);
\pgftext[right,top,at={\pgfpoint{-3em}{-3ex}}]{Error}}}}%
\cmdname{toady}%
\smash{\rlap{\color{blue}%
\tikz[baseline]{\draw[latex-] (0em,0ex) -- (0em,-6ex);
\pgftext[left,top,at={\pgfpoint{0em}{-6.5ex}}]{This is how far \LaTeX\ has got}}}}\\
 &\texttt{, which is nice to know.}\smash{\rlap{\color{blue}%
\tikz[baseline]{\draw[latex-] (0em,0.5ex) -- (3em,0.5ex);
\pgftext[left,center,at={\pgfpoint{3em}{0.5ex}}]{\shortstack{Rest of\\the line}}}}}\\
\texttt{?}\color{blue}%
\tikz[baseline]{\draw[latex-] (0em,0ex) -- (2em,-4ex);
\pgftext[left,top,at={\pgfpoint{2em}{-4ex}}]{\LaTeX\ prompt}}
\end{tabular}

\vspace{\baselineskip}
\end{latexonly}
\html{\htmladdimg[alt="annotated error message"]{pictures/errormessage.png}}

At the \LaTeX\ prompt, you can either type \texttt{h} for a help 
message, or type \texttt{x} to exit \LaTeX\ and go back to your 
source code and fix the problem.
\end{itemize}

There follows below a list of common error messages. If your
problem isn't listed there, try the \gls{ukfaq}.

%%%%%%%%%%%%%%%%%%%% * (No message) %%%%%%%%%%%%%%%%%%%%%%%%%%%

\setnode{asteriskprompt}
\section{* (No message, just an asterisk prompt!)}

You've gone into \gls{tex}!  This is probably because 
you've forgotten the \texttt{\gls{end}\marg{document}}.  The asterisk is the 
\TeX\ prompt.  At this point the best thing to do is to abort the
\TeX\ run.

%%%%%%%%%%%%% Argument of \cline has an extra } %%%%%%%%%%%%%%%%

\setnode{clineextrabrace}
\section{Argument of \textbackslash cline has an extra \}}

If this error occurred on the first line in your \gls{env-tabular}
\glsterm{environment}, you may have forgotten the
\glsterm{argument} to the \gls*{env-tabular} environment.

%%%%%%%%%%%%% Argument of \multicolumn has an extra } %%%%%%%%%%%%%%%%

\setnode{multicolumnextrabrace}
\section{Argument of \textbackslash multicolumn has an extra \}}

If this error occurred on the first line in your \gls{env-tabular}
\glsterm{environment}, you may have forgotten the
\glsterm{argument} to the \gls*{env-tabular} environment.

%%%%%%%%%%%%%%%% \begin{...} ended by \end{...} %%%%%%%%%%%%%%%%%%%%

\setnode{mismatchedend}
\section{\textbackslash begin\{\texorpdfstring{\ldots}{...}\} ended by 
\textbackslash end\{\texorpdfstring{\ldots}{...}\}}

The beginning of your environment doesn't have a matching end.
\begin{itemize}
\item Check to make sure you have spelt the name of the environment 
correctly.

You will get this error message if you do, say,
\begin{alltt}\wrong
\glsni{end}\marg{docment}
\end{alltt}
instead of
\begin{alltt}\correct
\glsni{end}\marg{document}
\end{alltt}

\item Check that for every \gls{begin} you have a corresponding 
\gls{end} with the same name.
\end{itemize}

%%%%%%%%%%%%%%%%%%% Bad math environment delimiter %%%%%%%%%%%%%%%%

\setnode{badmathdelim}
\section{Bad math environment delimiter}

Only a certain type of character may be used as a 
\htmlref{delimiter}{sec:delimiters} (for example,
\gls{openparen} \gls{closeparen} \gls{delimiter.opensq}
\gls{delimiter.closesq}), check which one you have specified.  This error may 
also occur if you have forgotten a \gls{right} or not used it in the
same \glslink{group}{scope}.
(Remember to use a \gls{delimiter.periodchar} if you want an invisible delimiter) or 
you may have forgotten to end your array environment with 
\verb|\end{array}|.

%%%%%%%%%%%%%%%%%%%% Can only be used in preamble %%%%%%%%%%%%%%%%%%

\setnode{onlypreamble}
\section{Can only be used in preamble.}

Some commands, such as \gls{usepackage} may only appear in the 
\gls{preamble}.  Check to see where you have put 
it.  For example, this error will be caused by doing:
\begin{alltt}
\glsni{documentclass}\marg{scrartcl}

\glsni{begin}\marg{document}\wrong

\glsni{usepackage}\marg{graphicx}
\end{alltt}
instead of
\begin{alltt}
\glsni{documentclass}\marg{scrartcl}

\glsni{usepackage}\marg{graphicx}\correct

\glsni{begin}\marg{document}
\end{alltt}

%%%%%%%%%%%%%%%%%% Command ... already defined %%%%%%%%%%%%%%%%%%%%%

\setnode{alreadydef}
\section{Command \texorpdfstring{\protect\ldots}{...}\protect\space already defined}

You have tried to define a \gls{command} which 
already exists.  Try giving it a different name.  Remember never to
redefine a command if you don't know what the command originally
does.

Alternatively, you have tried to define an
\gls{environment} which already exists.  Give the new
environment a different name.  Again, never redefine an environment
where you don't know what the original environment does.

%%%%%%%%%%%%% Display math should end with $$ %%%%%%%%%%%%%%%%%%%%%

\setnode{missingenddispmath}
\section{Display math should end with \$\$}

You may have a dollar sign (\gls{dollarchar}) in a displayed maths
environment, such as the \gls{env-equation} environment.  
Remember that \verb|$| is short hand for \verb|\begin{math}| or 
\verb|\end{math}|, so you can't end one of the other environments 
with a \verb|$|. (This error message is in fact a bit confusing,
as it seems to be suggesting that you end a displayed maths
environment with \verb|$$| which you also shouldn't do\faq{Why use
\glsnl{begindispmath}\ldots\glsnl{enddispmath} in place of \$\$\ldots\$\$}{dolldoll}.)

%%%%%%%%%%%%%% Environment ... undefined %%%%%%%%%%%%%%%%%%%%%%%%%

\setnode{envundefined}
\section{Environment \texorpdfstring{\protect\ldots}{...}\protect\space undefined}

\LaTeX\ doesn't recognise the environment you have specified.

\begin{itemize}
\item Check you have spelt the environment name correctly.

You will get this error if you do, say,
\begin{alltt}\wrong
\glsni{begin}\marg{docment}
\end{alltt}
instead of
\begin{alltt}\correct
\glsni{begin}\marg{document}
\end{alltt}

\item If it's your own environment, check you have defined the
environment before using it.

\item If the environment is defined in a package, check you 
have included the package using the 
\htmlref{\cmdname{usepackage}}{sec:packages} command.

\end{itemize}

%%%%%%%%%% Extra alignment tab has been changed to \cr %%%%%%%%%%%%

\setnode{tab2cr}
\section{Extra alignment tab has been changed to \textbackslash cr}

You have too many ampersands (\gls{ampchar}) in one row.  The most 
probable cause is that you have forgotten the end of row command
\gls{tab.dbbackslashchar} on the previous row.  Remember also that if you have a
\gls{multicolumn} command to span more than one column, you
should have fewer \glsni{ampchar}s in that row. This error can also occur
from a confusion over the two meanings of \gls{dbbackslashchar}: a
line break within a paragraph cell and a row break.  In which case,
you need to use \gls{tabularnewline} instead of
\gls{tab.dbbackslashchar}.

%%%%%%%%%%%%%%%%%%%%%%%%% Extra \right %%%%%%%%%%%%%%%%%%%%%%%%%%

\setnode{extraright}
\section{Extra \textbackslash right}

There are a number of possible causes.  The most probable is that you
have a \gls{right} that doesn't have a matching \gls{left}.
(Remember left comes before right.)  Another possible cause is that
you have missed out \verb|\end{array}|.  (Remember that
\glspl{environment} provide implicit
\glsing{group}, and \glsni{left} and its matching
\glsni{right} must appear within the same group level.)

%%%%%%%%%%%%%% File ended while scanning use of ... %%%%%%%%%%%%%

\setnode{fileendedwhilescanning}
\section{File ended while scanning use of \texorpdfstring{\ldots}{...}}

The most usual cause of this error is a missing closing brace.

You will get this error if you do, say,
\begin{alltt}\wrong
\glsni{end}\glsni{leftbracechar}document
\end{alltt}
instead of
\begin{alltt}\correct
\glsni{end}\marg{document}
\end{alltt}

%%%%%%%%%%%%%%%%%%%%%% File not found %%%%%%%%%%%%%%%%%%%%%%%%%%

\setnode{filenotfound}
\section{File not found}

\LaTeX\ can't find the file you have specified.  You will be given
the opportunity to type in the correct filename at the prompt.  If
you want to quit, simply type \texttt{X} followed by the return
character \enter.

\begin{itemize}

\item Make sure that you have spelt the filename correctly.

This error will be caused by, say,
\begin{alltt}\wrong
\gls{documentclass}\marg{scrarticle}
\end{alltt}
instead of
\begin{alltt}\correct
\gls{documentclass}\marg{scrartcl}
\end{alltt}

If this is the case, simply type in the correct name at the prompt 
(followed by the return character \enter) then go back and fix the
spelling in the \gls{source}.

\item Make sure that the file is in the same directory as your
document or in the \LaTeX\ path.  If the file is in another directory
(not in the \LaTeX\ path), you will need to specify the pathname, but
remember that when using \LaTeX\ under Windows, you need to use a
forward slash (\glsi{dir.slash}) as the directory divider, as a backslash would be
interpreted as a command.  For example, if you have a file called
\texttt{shapes.pdf} in the subdirectory \texttt{pictures} then you
would get a \dq{file not found} error message if you did
\begin{alltt}\wrong
\gls{includegraphics}\marg{shapes}
\end{alltt}
instead of
\begin{alltt}\correct
\gls{includegraphics}\marg{pictures/shapes}
\end{alltt}

\item If the file is a \htmlref{package}{sec:packages} or
\htmlref{class file}{sec:cls}, it's possible that you don't have that
file installed on your computer.  If this is the case you will need
to download and install it as described in
\appendixref{ch:installsty}.  Remember that you need to
refresh the database after installing a new package or class file.

\end{itemize}

%%%%%%%%%%%%%%%%%%%% Illegal character in array arg %%%%%%%%%%%%%%%%%%

\setnode{illegalarraychar}
\section{Illegal character in array arg}

You have used a character in the \gls{argument} of a
\gls{env-tabular} or \gls{env-array} environment that is not
allowed. The standard available characters are: \texttt{r} (right
justified), \texttt{l} (left justified), \texttt{c} (centred) and
\texttt{p}, as well as \gls{atchar}\marg{\meta{inter-col text}}.
Remember that if you want to use the \gls{gtcol}\marg{\meta{decl}} or 
\gls{ltcol}\marg{\meta{decl}} specifiers, you must include the
\isty{array} package.

This error will also occur if you have forgotten the argument to the
\gls*{env-tabular} or \gls*{env-array} environment.

%%%%%%%%%%%%%%%%%%% Illegal parameter number in definition %%%%%%%%%%

\setnode{illegalparamnum}
\section{Illegal parameter number in definition}

You have referred to a \htmlref{parameter (argument)
number}{sec:newcomarg} that is greater than the number of parameters
you have specified. For example, suppose you defined the command to
have only one parameter, then you can't use \gls{hashchar}\texttt{2}
which refers to the second, non-existent, parameter. Remember that
you need to specify how many parameters you want in the
\gls{optional} to \glsni{newcommand}, otherwise it will be assumed
that the command has no arguments.

%%%%%%%%%%%%%%%%%%%%% Illegal unit of measure %%%%%%%%%%%%%%%%%%%%%%%

\setnode{illegalunit}
\section{Illegal unit of measure (pt inserted)}

You have either not specified a unit when giving a \gls{length} (even zero
lengths must have a unit) or you have specified an invalid unit or
you have misspelt the unit.  Available units are listed in
\tableref{tab:units}.

%%%%%%%%%%%%%%%%% Lonely \item %%%%%%%%%%%%%%%%%%%%%%%%%

\setnode{lonelyitem}
\section{Lonely \textbackslash item}

The command \gls{item} may only appear in one of the list making
environments (such as \gls{env-itemize}).  Make sure you haven't
forgotten your environment.

%%%%%%%%%%%%%%% Misplaced alignment tab character & %%%%%%%%%%%%%%%%%

\setnode{misplacedtab}
\section{Misplaced alignment tab character \&}

You have used the special character \gls{ampchar} where you shouldn't
have.  Recall from \sectionref{sec:chars} that if you want an \& sign 
to appear you need to do \gls{amp} not just \glsni{ampchar}.

You would have got this error message if you had done, say,
\begin{alltt}\wrong
\glsni{ampchar} our equipment
\end{alltt}
instead of
\begin{alltt}\correct
\glsni{amp} our equipment
\end{alltt}

%%%%%%%%%%%%%%%%%%%%% Missing } inserted %%%%%%%%%%%%%%%%%%%%%%%%%%%

\setnode{missingendbrace}
\section{Missing \} inserted}

You have missed a closing curly brace, or you may have missed out an
argument.

\xminisec{Example:} 
If the following line occurs in a tabular environment:
\begin{alltt}\wrong
 \gls{ampchar} \gls{multicolumn}\marg{2}\marg{c}\gls{tab.dbbackslashchar}
\end{alltt}
this will produce the error. (The third argument to 
\gls{multicolumn} has been omitted.)

%%%%%%%%%%%%%%%%%%%%% Missing $ inserted %%%%%%%%%%%%%%%%%%%%%%%%%%%

\setnode{missingdollar}
\section{Missing \$ inserted}

This message can be caused by a number of errors:
\begin{itemize}
\item You might have missed the beginning of one of the mathematics
\glspl{environment} (that is, you've used a
\gls{command} that must only appear in maths mode).

\item You may have typed \gls{dollarchar} instead of \gls{dollar} (you
actually want a dollar symbol to appear).  Recall from
\sectionref{sec:chars} that if you
want a \$ sign to appear you need to do \glsni{dollar} not just \glsni{dollarchar}.

You would have got this error message if you had done, say,
\begin{alltt}\wrong
expenditure came to \glsni{dollarchar}2000.00
\end{alltt}
instead of
\begin{alltt}\correct
expenditure came to \glsni{dollar}2000.00
\end{alltt}

\item You may have missed the end of a mathematics environment, or
you may have a paragraph break within an in-line or
displayed maths environment, where it is
not permitted.  Make sure you don't have any blank lines within the
environment. If you want a blank line in your code to make it easier
to edit, try having a percent sign at the start of an empty line to
ensure that the line is ignored by \LaTeX. For example:
\begin{alltt}
\glsni{begin}\marg{equation}
\glsni{percentchar}
E = mc\glsni{circumchar}2
\glsni{percentchar}
\glsni{end}\marg{equation}
\end{alltt}
\end{itemize}

%%%%%%%%%%%%%%% Missing \begin{document} %%%%%%%%%%%%%%%%%

\setnode{missingbegindoc}
\section{Missing \textbackslash begin\{document\}}

You have put some text outside of the document 
\glsterm{environment}.  Check the following:
\begin{itemize}
\item You have remembered \verb|\begin{document}|

This error would be caused by, say,
\begin{alltt}
\glsni{documentclass}\marg{scrartcl}
\wrong
This is a simple document
\end{alltt}
instead of
\begin{alltt}
\glsni{documentclass}\marg{scrartcl}

\glsni{begin}\marg{document}\correct
This is a simple document.
\end{alltt}

\item You haven't placed any text before \verb|\begin{document}|. 
For example:
\begin{alltt}
\glsni{documentclass}\marg{scrartcl}

This is a simple document\wrong 
\glsni{begin}\marg{document}
\end{alltt}
instead of
\begin{alltt}
\glsni{documentclass}\marg{scrartcl}

\glsni{begin}\marg{document}\correct
This is a simple document
\end{alltt}

\item You haven't missed out the backslash at the start of either 
\gls{documentclass} or \glsni{begin}\marg{document}

This error would be caused by, say,
\begin{alltt}\wrong
documentclass\marg{scrartcl}
\end{alltt}
instead of
\begin{alltt}\correct
\glsni{documentclass}\marg{scrartcl}
\end{alltt}
\end{itemize}

%%%%%%%%%%%%%%%%%%% Missing delimiter %%%%%%%%%%%%%%%%%%%%%%%

\setnode{missingdelim}
\section{Missing delimiter}

You have forgotten to specify the type of delimiter you want (for
example, 
\gls{openparen} \gls{closeparen} \gls{delimiter.opensq} \gls{delimiter.closesq}
\gls{leftbrace} \gls{rightbrace}). (Remember to use a \gls{delimiter.periodchar} if you want 
an invisible delimiter, and remember that if you want a 
curly brace, you must have a backslash followed by the curly brace.)

\xminisec{Example:}

This error will occur if you do, say,
\begin{alltt}\wrong
f(x) = \glsni{left}\glsni{leftbracechar}
\glsni{begin}\marg{array}\marg{ll}
0 \glsni{ampchar} x \cmdname{leq} 0\gls{tab.dbbackslashchar}
1 \glsni{ampchar} x > 1
\glsni{end}\marg{array}
\glsni{right}.
\end{alltt}
instead of
\begin{alltt}\correct
f(x) = \glsni{left}\glsni{leftbrace}
\glsni{begin}\marg{array}\marg{ll}
0 \glsni{ampchar} x \cmdname{leq} 0\gls{tab.dbbackslashchar}
1 \glsni{ampchar} x > 1
\glsni{end}\marg{array}
\glsni{right}.
\end{alltt}

%%%%%%%%%%%%%%%% Missing \endcsname inserted %%%%%%%%%%%%%%%%

\setnode{missingendcsname}
\section{Missing \textbackslash endcsname inserted}

This is a \gls{tex} error rather than a \LaTeX\ error
which makes it harder to determine the cause, however it can be
caused by placing a backslash in front of the name of an
\gls{environment}.  (Remember that
\gls{environment} names do not contain a backslash.)

This error will be caused by, say,
\begin{alltt}\wrong
\glsni{begin}\marg{\glsni{sffamily}}
\end{alltt}
instead of
\begin{alltt}\correct
\glsni{begin}\marg{sffamily}
\end{alltt}

%%%%%%%%%%%%%% Missing \endgroup inserted %%%%%%%%%%%%%%%%%%

\setnode{missingendgroup}
\section{Missing \textbackslash endgroup inserted}

A number of things could have caused this.  You may have missed out
the end of an \gls{environment}, or you may have an
environment inside of another environment it's not allowed to be in.
For example, this error can be caused by placing an
\envname{eqnarray} environment inside a \envname{displaymath}
environment, which is not allowed. (But, of course, you haven't
used either of those obsolete environments~\cite{l2tabu}, have you?) 

%%%%%%%%%%%%%%% Missing number %%%%%%%%%%%%%%%%%%

\setnode{missingnum}
\section{Missing number, treated as zero}

\LaTeX\ is expecting a number. If your command takes more than one
\gls{argument}, check to make sure the arguments are in
the correct order.  For example, if you are using a
\gls{env-minipage} environment, you might have
omitted the \gls{mandatory} which
specifies the width of the minipage, or you may have the
\glspl{optional} the wrong way round.  The placement specifier
should come first, followed by the height.

If you are using \gls{addtocounter} or \gls{setcounter} 
remember that the second argument must be a number, so if you want 
the value of a counter as the argument you must use \gls{value}.
This error can be caused by, say,
\begin{alltt}\wrong
\glsni{setcounter}\marg{exercise}\marg{chapter}
\end{alltt}
instead of
\begin{alltt}\correct
\glsni{setcounter}\marg{exercise}\marg{\glsni{value}\marg{chapter}}
\end{alltt}

%%%%%%%%%%%%%% Paragraph ended %%%%%%%%%%

\setnode{parended}
\section{Paragraph ended before \textbackslash begin was complete}

You've probably missed a closing brace at the end of the argument to
\gls{begin}.  This error will be caused by, say,
\begin{alltt}\wrong
\glsni{begin}\glsni{leftbracechar}document
\end{alltt}
instead of
\begin{alltt}\correct
\glsni{begin}\marg{document}
\end{alltt}

%%%%%%%%%%%%%%%%%%% Runaway argument %%%%%%%%%%%%%%%%%%%%%%%%%%%%

\setnode{runawayarg}
\section{Runaway argument}

There are a number of possible causes of this error:

\begin{itemize}

\item Paragraph breaks are not permitted in the
\glspl{argument} of \glspl{short}.
If there is a corresponding \gls{environment} then you
should use that instead.  For example, this error message will be
generated by doing, say,
\begin{alltt}
\glsni{textbf}\glsni{leftbracechar}This is a simple document.
Here is the first paragraph.
\wrong
Here is the second paragraph.\glsni{rightbracechar}
\end{alltt}
instead of\bookpagebreak
\begin{alltt}
\glsni{begin}\marg{bfseries}
This is a simple document.
Here is the first paragraph.
\correct
Here is the second paragraph.
\glsni{end}\marg{bfseries}
\end{alltt}

\item The closing brace of a 
\gls{mandatory} is missing:
This error will be caused by, say,
\begin{alltt}\wrong
\glsni{title}\glsni{leftbracechar}A Simple Document
\end{alltt}
instead of
\begin{alltt}\correct
\glsni{title}\marg{A Simple Document}
\end{alltt}

\item This error can also be caused by omitting the 
\gls{mandatory} of an \gls{environment}.  For example, 
this error will occur
if you do, say,
\begin{alltt}\wrong
\glsni{begin}\marg{thebibliography}
\glsni{bibitem}\marg{kopka95} A Guide to \glsni{LaTeXe}
\end{alltt}
instead of
\begin{alltt}\correct
\glsni{begin}\marg{thebibliography}\marg{1}
\glsni{bibitem}\marg{kopka95} A Guide to \glsni{LaTeXe}
\end{alltt}

\end{itemize}

%%%%%%%%%%%%%%% Something's wrong--perhaps a missing \item %%%%%%%%%%%%%%%%

\setnode{missingitem}
\section{Something's wrong--perhaps a missing \textbackslash item}

You may have missed an \gls{item} command.  The first object in
a list environment must either be an \gls{item} command, or another
list environment.

This error will be caused by, say,
\begin{alltt}\wrong
\glsni{begin}\marg{itemize}
Animal
\glsni{item} Vegetable
\glsni{item} Mineral
\glsni{end}\marg{itemize}
\end{alltt}
instead of
\begin{alltt}\correct
\glsni{begin}\marg{itemize}
\glsni{item} Animal
\glsni{item} Vegetable
\glsni{item} Mineral
\glsni{end}\marg{itemize}
\end{alltt}

\bookpagebreak
This error can also be caused by a missing \gls{bibitem} in the
\htmlref{bibliography}{sec:bib}.  For example, the error will occur
if you do, say,
\begin{alltt}\wrong
\glsni{begin}\marg{thebibliography}\marg{1}
A Guide to \glsni{LaTeXe}
\end{alltt}
instead of
\begin{alltt}\correct
\glsni{begin}\marg{thebibliography}\marg{1}
\glsni{bibitem}\marg{kopka95} A Guide to \glsni{LaTeXe}
\end{alltt}

See also \gls{ukfaq} entry: \faqlink{Perhaps a missing \cmdname{item}?}{errmissitem}.

%%%%%%%%%%%%%%%%% There's no line here to end %%%%%%%%%%%%%%%%%%%%%

\setnode{nolinetoend}
\section{There's no line here to end}

You have placed a line breaking command (such as \gls{newline.dbbackslashchar},
\gls{newline} or \gls{linebreak}) where it doesn't make
sense to have one.

%%%%%%%%%%%%%%%%% Undefined control sequence %%%%%%%%%%%%%%%%%%%%%%

\setnode{csundefined}
\section{Undefined control sequence}

\LaTeX\ doesn't understand the \gls{command} you
have used.  

\begin{itemize}
\item Check to see if you have misspelt the command name (remember
that all \LaTeX\ command names are case-sensitive.)

You will get this error if you do, say,
\begin{alltt}\wrong
This is a simple \cmdname{Latex}\glsni{spacesym}document
\end{alltt}
instead of
\begin{alltt}\correct
This is a simple \gls{LaTeX}\glsni{spacesym}document
\end{alltt}

\item Check that you have remembered the space when typing
\gls{spacesym} (backslash space).
For example, this error will occur if you do, say,
\begin{alltt}\wrong
This is a \cmdname{LaTeX}\cmdname{sample} document.
\end{alltt}
instead of
\begin{alltt}\correct
This is a \glsni{LaTeX}\glsni{spacesym}sample document
\end{alltt}

\item If you are using a command that is defined in a
\htmlref{package}{sec:packages} make sure you have included the
package using \gls{usepackage}.

\item Check that your command name hasn't run into the next piece
of text.  For example,
you can do
\begin{alltt}
man\glsni{oe}\marg{}uvre
\end{alltt}
or
\begin{alltt}
man\glsni{oe} uvre
\end{alltt}
or (not recommended)
\begin{alltt}
man\marg{\glsni{oe}}uvre
\end{alltt}
but not
\begin{alltt}\wrong
man\cmdname{oeuvre}
\end{alltt}

\item Check if you have used a backslash instead of a forward slash as a
directory divider.  (Remember that when using \LaTeX\ under Windows,
you need to use a forward slash (\glsi{dir.slash}) as the directory divider, as a
backslash would be interpreted as a command.)

For example, suppose you have a file called \texttt{shapes.pdf} in a 
subdirectory called \texttt{pictures}, then you would get an error 
if you did
\begin{alltt}\wrong
\glsni{includegraphics}\marg{pictures\cmdname{shapes}}
\end{alltt}
instead of
\begin{alltt}\correct
\glsni{includegraphics}\marg{pictures/shapes}
\end{alltt}

\end{itemize}

%%%%%%%%%% You can't use `macro parameter character #' in horizontal mode %%%%%%%%%%%%%%%

\setnode{hashinhmode}
\section{You can't use `macro parameter character \#' in horizontal mode}

You have used the special character \gls{hashchar} where you shouldn't
have.  Recall from \sectionref{sec:chars} that if you want a \# sign to appear you need
to do \gls{hash} not just \glsni{hashchar}.

This error message will be caused by doing, say,
\begin{alltt}\wrong
Item \glsni{hashchar}1
\end{alltt}
instead of
\begin{alltt}\correct
Item \glsni{hash}1
\end{alltt}

\setnode{help}
\chapter{Need More Help?}
\label{ch:help}

First, try to find your query in the \gls{ukfaq}. \Gls{tug}
also has a list of useful resources at
\url{http://tug.org/interest.html}. If you're still stuck, you can
post your question on a (La)TeX forum, newsgroup or mailing list,
such as those listed below. If you do post a question, remember
you're asking people who only have an altruistic interest in
helping. No one is paying them to help you. Most of the class files
and packages were written for free by people who had a need to solve
a particular problem and decided to make their work publicly
available. So no matter how frustrated you're feeling, stick to
being polite. If you can't work out how to use a particular class or
package, don't start by heaping offensive, unconstructive criticism on it as there's a chance
the author will read the message. There's no sense in alienating the
person most qualified to answer your question. In your message,
stick to the following guidelines:
\begin{enumerate}
\item Cut to the chase. In other words, be concise about the nature
of the problem. Don't write lots of long-winded paragraphs.

\item Provide a minimal
example\footnote{see \url{http://www.dickimaw-books.com/latex/minexample/}} that 
illustrates the problem.

\end{enumerate}

\xminisec{Example:}
\begin{verbatim}
I'm trying to use the \foo command in the "bar" package, 
but I'm getting the following error message:
! Undefined control sequence.
l.4 \foo

Here's a minimal example:
\documentclass{scrartcl}
\usepackage{bar}
\begin{document}
\foo{Blah}
\end{document}

I'm using bar version 1.0 (2012/06/30).
\end{verbatim}

\xminisec{Another example:}
\begin{verbatim}
I'm using the \foo command in the "bar" package. According
to the documentation, this command should display its argument
in a bold font, but it's coming out in italic instead. 
Anyone know why?

Here's a minimal example:
\documentclass{scrartcl}
\usepackage{bar}
\begin{document}
\foo{Blah}
\end{document}

I'm using bar version 1.1 (2012/07/30).
\end{verbatim}

There's no guarantee that you will get an answer, but if you follow
the above guidelines, you will increase your chances.

\xminisec{Resources}
\begin{itemize}

\item The \LaTeX\ Community (\url{http://www.latex-community.org/}).

\item \TeX/\LaTeX\ on StackExchange
(\url{http://tex.stackexchange.com/}).

\item \htmladdnormallink{\texttt{comp.text.tex}
newsgroup}{http://groups.google.com/group/comp.text.tex} (use a
newsreader rather than the Google interface if you want to avoid the
spam).

\item \htmladdnormallink{\texttt{texhax} archives}{http://tug.org/pipermail/texhax/}\doifbook{~\cite{texhax}}.

\end{itemize}
I~strongly recommend that you also have a look at the On-Line
Catalogue~\cite{texcat}.
It's also a good idea to look at the documentation that was
installed with your \TeX/\LaTeX\ distribution (see
\sectionref{sec:texdoc}). If you are using \itexdistro{MiKTeX} you can access the
on-line help via the Start Menu: \startmenu{MiKTeX \menuto\ Help}

(Please don't send your problems to me, unless you want to hire a
consultant. I~read both the \LaTeX\ Community Forum and
\texttt{comp.text.tex} and answer relevant questions if I~have time,
but it clogs up my inbox if people keep sending attachments that are in
the order of several megabytes in size.) Besides, you'll reach a
wide group of experts if you post to a newsgroup, forum or mailing
list, rather than a single busy individual.

%%%%%%%%%%%%%%%%%%%%%%%%%% Supplemental Material %%%%%%%%%%%%%%%%%%%

\ifpdf\else\setnode{supplemental}
\chapter{Supplementary Material}
\label{ch:supp}

This appendix chapter contains three sections from earlier versions of this
book. The information here is a little dated, but is provided in
case it's still of any use.

%%%%%%%%%%%%%%%%%%% Terminal + Text Editor %%%%%%%%%%%%%%%%%%%

\setnode{editorandterminal}
\section{Text Editor and Terminal Approach}
\label{sec:editorandterminal}

Creating a \LaTeX\ document using a text editor and a terminal
is an approach often favoured by UNIX-type users. If you have never
used a terminal (i.e.\ you have only ever used point-and-click menu
driven applications) then you will be better off using a
front-end, in which case I suggest you
\begin{latexonly}%
turn to Sections~\ref{sec:TeXnicCenter} and~\ref{sec:winedt} which 
describe \appname{TeXnicCenter} and \appname{WinEdt}, respectively.
\end{latexonly}
\begin{htmlonly}
look at the \htmlref{\appname{TeXnicCenter}}{sec:TeXnicCenter} and 
\htmlref{\appname{WinEdt}}{sec:winedt} sections instead.
\end{htmlonly}
\par
To begin, you will first need a text editor\faq{TeX-friendly editors
and shells}{editors}. There are a number available that are suited 
to using with \LaTeX, some people advocate
\iappnamelink{Emacs}{http://www.gnu.org/software/emacs/emacs.html},
others advocate \iappnamelink{vim}{http://vim.sourceforge.net/}, and
there are various others such as
\iappnamelink{NEdit}{http://www.nedit.org/}.  I prefer to use
\iappname{vim}---I'm not overly keen on using the mouse, and I prefer
being able to issue all commands via the keyboard (although 
there is a GUI version of \iappname{vim}). As with some other
editors, it comes with syntax highlighting, regular expression search
and replace, auto-insertion, and a brace matching mechanism which I
find useful. If you are using version~7 of \iappname{vim}, there is an
integrated spell checker, otherwise there is a spell checker plug-in
called
\htmladdnormallink{vimspell}{http://www.vim.org/scripts/script.php?script_id=465},
so you can check your spelling as you type.  If there is already a
text editor that you are comfortable with, then stick with that,
otherwise try out available editors, and decide which one you prefer.

When using the terminal and text editor approach, most people usually
have at least two terminals open: one to run the editor, the other to
run \LaTeX. This means that you don't have to keep quitting the
editor every time you want to \htmlref{\LaTeX\ your
document}{itm:step2}. Some editors allow you to run commands, but
personally I don't like to use this approach. If your editor has a
GUI interface, then you'll probably only need one terminal open.

Let's get started: start up your text editor. This is usually done by
entering the name of the editor at the command prompt in your
terminal. With most editors you can also specify the filename as
well. If the file doesn't exist, a new one will be created when you
save your document. \figureref{fig:startvim} shows my terminal.  The
command prompt looks like \verb|[nlct@nlctltpc examples]$|.  It will
be different for your system. I have typed \texttt{vim sample1.tex}
at the command prompt. This will start \iappname{vim} with a new file
called \texttt{sample1.tex}.  (In this section, I will be using
\iappname{vim} as the text editor, if you are not using
\iappname{vim}, then substitute the editor of your choice.)

\begin{figure}[htbp]
\figconts[alt=Image of a terminal]
 {pictures/terminal1}
 {\caption{Starting vim from a terminal}}
 {fig:startvim}
\end{figure}

Once I have pressed the return key, my terminal looks like
\figureref{fig:terminal2}. Normally \iappname{vim} starts in visual
command mode, which means that when you start typing text, it will be
interpreted as part of a command. In order to type text into your
file, you will need to enter input mode. There are a number of ways
of doing this, but pressing \texttt{i} will do for now\footnote{For a
complete set of available commands, see the \appname{vim} manual at
\url{http://vimdoc.sourceforge.net/}}. \figureref{fig:terminal3}
shows how my terminal looks when I am in input mode. I can now go
ahead and type in my text (\figureref{fig:terminal4}). To go back to
the visual command mode, press the escape key (Esc). Now that you are
back in the visual command mode, you can save your document, either
using the command \texttt{:w} if you have already given your file a
name, or \texttt{:w }\meta{filename} (e.g.\ \texttt{:w sample1.tex})
if you started \iappname{vim} without specifying a file, see
\figureref{fig:terminal5}. When you want to quit \iappname{vim} you
can do \texttt{:wq} to save and quit or \texttt{:q!} to quit without
saving, but I suggest you don't do this just yet if you have another
terminal available.

\begin{figure}[htbp]
 \figconts[alt=Image of a terminal with vim running]
 {pictures/terminal2}
 {\caption{Starting a new file in vim}}
 {fig:terminal2}
\end{figure}

\begin{figure}[htbp]
  \figconts[alt=Image of a terminal with vim running in input mode]
  {pictures/terminal3}
  {\caption{Input mode in vim}}
  {fig:terminal3}
\end{figure}

\begin{figure}[htbp]
  \figconts[alt=Image of a terminal with LaTeX code typed in vim]
  {pictures/terminal4}
  {\caption{Creating a sample document in vim}}
  {fig:terminal4}
\end{figure}

\begin{figure}[htbp]
  \figconts[alt=Image of new file being saved in vim]
  {pictures/terminal5}
  {\caption{Saving your document in vim (the file name should be 
    omitted if the file already has a name)}}
  {fig:terminal5}
\end{figure}

\objectref{Step}{itm:step1} is now complete, and you are now ready to
move on to \objectref{Step}{itm:step2}: using \LaTeX. Go to your
other terminal (or quit your editor if you only have access to one
terminal) and make sure that you are in the same directory as the
file you just created.  Typing \texttt{ls} at the command prompt will
list the contents of your current directory. If you do this, you
should see the file that you have just created (see
\figureref{fig:terminal6}).  At the command prompt type\faq{Makefiles
for LaTeX documents}{make}:
\refstepcounter{object}\label{obj:termlatex}
\begin{verbatim}
latex sample1.tex
\end{verbatim}
as shown in \figureref{fig:terminal7}. You can omit the \texttt{.tex}
extension if you like, \LaTeX\ will automatically add this if
it has been omitted. If you prefer to use \iPDFLaTeX, type
\begin{verbatim}
pdflatex sample1.tex
\end{verbatim}
instead (again the \texttt{.tex} extension may be omitted.)
You should now see something like \figureref{fig:terminal8}.

\begin{figure}[htbp]
  \figconts[alt=Image of directory listing in a terminal]
  {pictures/terminal6}
  {\caption{Listing the contents of the current directory}}
  {fig:terminal6}
\end{figure}

\begin{figure}[htbp]
  \figconts[alt=Image of a terminal where user has typed the command
   to run LaTeX]
  {pictures/terminal7}
  {\caption{Running \LaTeX}}
  {fig:terminal7}
\end{figure}

\begin{figure}[htbp]
  \figconts[alt=Image of a terminal showing normal LaTeX messages]
  {pictures/terminal8}
  {\caption{Running \LaTeX}}
  {fig:terminal8}
\end{figure}

\refstepcounter{object}\label{obj:transcript}
Numbers appearing in square brackets, e.g.\ \texttt{[1]}, indicate
which page \LaTeX\ is currently processing.  In this case, there is
only one page.  The last line to appear on screen
indicates that information about this \LaTeX\ run has been written
to the log file \texttt{sample1.log}, which you can look at using 
your text editor.

The most important thing to note is the penultimate line\footnote{if
you are using \PDFLaTeX, it will have \texttt{sample1.pdf} instead of
\texttt{sample1.dvi}}:
\begin{verbatim}
Output written on sample1.dvi (1 page, 248 bytes).
\end{verbatim}
This means that the document has been successfully created, and
is one page long.

If you have made a mistake in your \htmlref{source code}{sec:source},
for example suppose you have missed the starting backslash in
\cmdname{documentclass}, then the output will look something like:
\begin{verbatim}
! LaTeX Error: Missing \begin{document}.

See the LaTeX manual or LaTeX Companion for explanation.
Type  H <return>  for immediate help.
 ...

l.1 d
     ocumentclass[a4paper]{article}
?
\end{verbatim}
There are several things you can do at this point, but the easiest
thing to do is to exit \LaTeX\ by typing \texttt{X} followed by the
return key. Go back to your editor, fix the mistake, save the
document, and then try again. (If you do get an error message, check
\chapterref{ch:errors}.) Note that it is important to always save your
document before running \LaTeX.

\refstepcounter{object}\label{obj:xdvi}
You can view the typeset output by loading the file\indexDVI\  
\texttt{sample1.dvi} into a DVI viewer, such as \iappname{xdvi}
or \iappname{kdvi}. To do this type
\begin{verbatim}
xdvi sample1.dvi
\end{verbatim}
or
\begin{verbatim}
kdvi sample1.dvi
\end{verbatim}
at the command prompt (see
\figureref{fig:terminal9}). You will then see the final output,
as shown in \figureref{fig:terminal10}.

\begin{figure}[htbp]
 \figconts[alt=Image of a terminal where user has typed command to
  load DVI file into kdvi]
  {pictures/terminal9}
  {\caption{Load a DVI file into a DVI viewer}}
  {fig:terminal9}
\end{figure}

\begin{figure}[htbp]
  \figconts[alt=Image of DVI file being viewed using kdvi]
  {pictures/terminal10}
  {\caption{Viewing a DVI file in kdvi}}
  {fig:terminal10}
\end{figure}

If you have used \iPDFLaTeX\ instead of \LaTeX, you should have a
file called \texttt{sample1.pdf} instead of \texttt{sample1.dvi}.
You can view this using a PDF viewer, such as \iappname{acroread}
or \iappname{kpdf}.

Some viewers, such as \iappname{kdvi} and \iappname{kpdf} will
automatically reload the file whenever it is modified, in which case
you may like to keep the viewer open, and as you keep editing and
\LaTeX{}ing your document, the viewer will automatically reload the
new versions. Some viewers, such as \iappname{xpdf} don't
automatically reload, but have a reload facility, which you can use
whenever you \htmlref{\LaTeX\ your document}{itm:step2}.

\refstepcounter{object}\label{obj:dvips}
If you like, you can convert your DVI file\indexDVI\ to PostScript 
using \iappname{dvips}. To do this, type the following at the command
prompt in your terminal:
\begin{verbatim}
dvips -o sample1.ps sample1.dvi
\end{verbatim}
(The \texttt{.dvi} extension may be omitted.)
You can then view the PostScript file using \iappname{ghostscript}
or one of its associated applications, such as \iappname{ghostview}
or \iappname{GSview}. I have \iappname{kghostview} installed on
my laptop, so to view the PostScript file, \texttt{sample1.ps},
I would need to type:
\begin{verbatim}
kghostview sample1.ps
\end{verbatim}
(See \figureref{fig:terminal11}.)

\begin{figure}[htbp]
  \figconts[alt=Image of terminal where user has typed command to
   load PostScript file into kghostview]
   {pictures/terminal11}
  {\caption{Loading a PostScript file}}
  {fig:terminal11}
\end{figure}

%%%%%%%%%%%%%%%%% TeXnicCenter %%%%%%%%%%%%%%%%%%%%%%%%%%%%

\setnode{TeXnicCenter}
\section{TeXnicCenter}
\label{sec:TeXnicCenter}

\iappname{TeXnicCenter} is an application that enables you to edit
\LaTeX\ source code, and simply click on a button to pass the source
code to \LaTeX, and then click on another button to view the
resulting typeset document.  Many people prefer this approach to the
\latexhtml{text editor and terminal approach described in the
previous section}{\htmlref{text editor and terminal}{sec:editorandterminal}
approach}. This section gives a brief overview of 
\appname{TeXnicCenter}, however it has been several years since
I last used it\footnote{I used to use TeXnicCenter and
MiKTeX when I was teaching \LaTeX, but that was my limit of using
\TeX\ under Windows}, so this information may be dated.

\appname{TeXnicCenter} is free and can be downloaded from \gls{ctan} in the
\htmladdnormallink{\texttt{systems/win32/TeXnicCenter/}}{http://www.tex.ac.uk/tex-archive/systems/win32/TeXnicCenter/}
directory or from \url{http://www.toolscenter.org/}.  Note that you
must have a \htmlref{\TeX/\LaTeX\ distribution}{obj:MiKTeX} installed
before you install \appname{TeXnicCenter}. If you installed
\itexdistrolink{proTeXt}{http://www.tug.org/protext}, you should
already have \appname{TeXnicCenter} installed. If you have any
problems with installing or running \appname{TeXnicCenter}, go to
their \latexhtml{help page at
\url{http://www.texniccenter.org/help.html}}{\htmladdnormallink{help
page}{http://www.texniccenter.org/help.html}}.

Once the installation is complete, you can then run
\appname{TeXnicCenter} from the Start Menu: \startmenu{TeXnicCenter
\menuto\ TeXnicCenter} Firstly you should see the tip of the day
window (\figureref{fig:TXCtip}.)

\begin{figure}[hbtp]
  \figconts[alt=Image of TeXnicCenter's tip of the day window]
  {pictures/texniccenter/tip}
  {\caption{TeXnicCenter Tip of the Day Window}}
  {fig:TXCtip}
\end{figure}

You can close this window, and then, if this is the first time you
are using \appname{TeXnicCenter} you will have to use the
configuration wizard to set up \appname{TeXnicCenter} correctly.  I
would recommend that you choose the default settings. (Select
\menu{\underline{N}ext}, \menu{\underline{N}ext} and then
\menu{Finish}.)

\begin{figure}[hbtp]
  \figconts[alt=Image of TeXnicCenter's configuration wizard: welcome
   screen]
  {pictures/texniccenter/setup01}
  {\caption{TeXnicCenter Configuration Wizard}}
  {fig:TXCsetup1}
\end{figure}

\begin{figure}[hbtp]
  \figconts[alt=Image of TeXnicCenter's configuration wizard: use
MiKTeX option selected]
  {pictures/texniccenter/setup02}
  {\caption{TeXnicCenter Configuration Wizard}}
  {fig:TXCsetup2}
\end{figure}

\begin{figure}[hbtp]
\figconts[alt=Image of TeXnicCenter's configuration wizard: finish
configuration]
  {pictures/texniccenter/setup03}
  {\caption{TeXnicCenter Configuration Wizard}}
  {fig:TXCsetup3}
\end{figure}

Now you are ready to use \appname{TeXnicCenter}.  It should look like
\figureref{fig:TXCuse1}.

\begin{figure}[hbtp]
\figconts[alt=Image of TeXnicCenter: no files loaded]
  {pictures/texniccenter/use01}
  {\caption{TeXnicCenter}}
  {fig:TXCuse1}
\end{figure}

\mbox{}\refstepcounter{object}\label{TXCnewproject}%
To start a new project select \menu{File \menuto\ New Project}.
This will open the window shown in \figureref{fig:TXCuse2}.

\begin{figure}[hbtp]
\figconts[alt=Image of TeXnicCenter's New Project dialog box]
  {pictures/texniccenter/use02}
  {\caption{New Project Dialog Box}}
  {fig:TXCuse2}
\end{figure}

Enter a name for your project, and specify the directory where you
want to save your work.  For example, I shall call my project
\dq{example} and I want to save it in 
\verb|c:\My Documents\Nicky\example| (see \figureref{fig:TXCuse3}.)

\begin{figure}[hbtp]
\figconts[alt=Image of TeXnicCenter's New Project dialog box:
project name has been entered]
  {pictures/texniccenter/use03}
  {\caption{New Project Dialog Box}}
  {fig:TXCuse3}
\end{figure}

Select the \dq{Empty Project} icon, and click on \dq{Okay}.  You
should now see something like \figureref{fig:TXCuse4}.

\begin{figure}[hbtp]
\figconts[alt=Image of TeXnicCenter: new project on view]
  {pictures/texniccenter/use04}
  {\caption{TeXnicCenter --- New Project Started}}
  {fig:TXCuse4}
\end{figure}

You can now start typing the source code (we'll cover this
\htmlref{later}{ch:simpledoc}).  See \figureref{fig:TXCuse5}.

\begin{figure}[hbtp]
\figconts[alt=Image of TeXnicCenter: user has typed source code in
the edit panel]
  {pictures/texniccenter/use05}
  {\caption{TeXnicCenter --- Typing in Source Code}}
  {fig:TXCuse5}
\end{figure}

Save it by either clicking on the save icon 
\incGraphics[alt=Image of save icon]{pictures/texniccenter/saveicon} or 
select \menu{File \menuto\ Save}

Now select what type of output you want (DVI, PDF or
\Index{PostScript}\refstepcounter{object}\label{obj:txcdvips}) see
\figureref{fig:TXCuse6}.  If this box is blank, then it's possible
that you didn't complete all the steps in the configuration wizard
described above.

\begin{figure}[hbtp]
\figconts[alt=Image of TeXnicCenter: drop-down menu showing]
  {pictures/texniccenter/use06}
  {\caption{TeXnicCenter --- Selecting Output Type}}
  {fig:TXCuse6}
\end{figure}

\refstepcounter{object}\label{TXC:latex}
Now click on the build output icon 
\incGraphics[alt=Image of build icon]{pictures/texniccenter/build} or 
select \menu{Build \menuto\ Build Output}.  The transcript will be
written in the window at the bottom (see \xfigureref{fig:TXCuse7})
This transcript should be the same as described \latexhtml{on
page~\pageref{obj:transcript}
onwards}{\htmlref{earlier}{obj:transcript}}.  If you have selected
\verb|LaTeX => PDF|, then \appname{TeXnicCenter} will use \iPDFLaTeX\
instead of \LaTeX.  If you have selected \verb|LaTeX => PS|, then
\appname{TeXnicCenter} will use \LaTeX\ followed by \iappname{dvips}
(as in \xfigureref{fig:TXCuse7}).  The \appname{dvips} messages will
follow on from the \LaTeX\ messages.  (If you selected the BibTeX or
MakeIndex features when you initialised the project,
\figureref{fig:TXCuse3}, then \iappname{TeXnicCenter} will also use
the \iBiBTeX\ and \appname{MakeIndex} applications.)

\begin{figure}[hbtp]
\figconts[alt=Image of TeXnicCenter: messages have been printed on
the build panel]
  {pictures/texniccenter/use07}
  {\caption[TeXnicCenter (using LaTeX and dvips)]{TeXnicCenter (using \LaTeX\ and \appname{dvips})}}
  {fig:TXCuse7}
\end{figure}

\refstepcounter{object}\label{TXC:yap}
To view the document, click the View Output button 
\incGraphics[alt=Image of view icon]{pictures/texniccenter/view}.
(Note that if you have selected \verb|LaTeX => PDF| or
\verb|LaTeX => PS| you will need \iappname{Adobe Reader}
or \iappname{GSview}, respectively, to view the output file.)

\refstepcounter{object}\label{TXCerrors}
If there are any errors, you can select \menu{Build \menuto\ Next
Error} and it will show you where the error has occured (See
\figureref{fig:TXCuse8}).  If you do have any errors, check
\chapterref{ch:errors}.

\begin{figure}[hbtp]
\figconts[alt=Image of TeXnicCenter: error message displayed in
build panel. A red arrow points to line containing the error in the
edit panel.]
  {pictures/texniccenter/use08}
  {\caption{TeXnicCenter --- Showing Error}}
  {fig:TXCuse8}
\end{figure}


%%%%%%%%%%%%%%%%%%%%%%%%%%%%% WinEdt %%%%%%%%%%%%%%%%%%%%%%%%%%%%%

\setnode{winedt}
\section{WinEdt}
\label{sec:winedt}

\iappname{WinEdt} (not to be confused with \appname{WinEdit} which is
a completely different application) is an application that enables
you to edit \LaTeX\ source code, and simply click on a button to pass
the source code to \LaTeX, and then click on another button to view
the resulting typeset document.  This section gives a brief overview
of \iappname{WinEdt}, however it has been several years since I
last used it, so this information may be dated.

\iappname{WinEdt} is shareware: it can be downloaded from \gls{ctan} in the
\htmladdnormallink{\texttt{systems/win32/winedt}}{http://www.tex.ac.uk/tex-archive/systems/win32/winedt/}
directory or from \url{http://www.winedt.com/} and evaluated for a
trial period of 31 days, after which, if you want to continue to use
it, you must pay the registration fee.  Details of prices and types
of licence available can be found at \url{http://www.winedt.com/}.
Note that members of the \gls{uktug} can benefit from their 
\iappname{WinEdt} licence
scheme (see \url{http://uk.tug.org/membership} for further details).

Again, you must have a \htmlref{\TeX/\LaTeX\
distribution}{obj:MiKTeX} installed before you start.
\appname{WinEdt} is fairly easy to install.  First unpack all the
files, and then run the \appname{setup.exe} application.  I recommend
that you use the default settings.  If you have any problems
installing or using \appname{WinEdt}, go to
\url{http://www.winedt.com/support.html}.

To run \iappname{WinEdt}, select \appname{WinEdt} from the start menu:
\startmenu{WinEdt \menuto\ WinEdt}
It should look like \figureref{fig:winedt1}.

\begin{figure}[hbtp]
\figconts[alt=Image of WinEdt: no files loaded]
  {pictures/winedt/01}%
  {\caption{WinEdt}}
  {fig:winedt1}
\end{figure}

Click on the \dq{New Document} button or select \menu{File \menuto\
New}.  You can now start typing your \htmlref{source
code}{sec:source} into the \appname{WinEdt} window, as shown in
\figureref{fig:winedt3}.

\begin{figure}[htbp]
\figconts[alt=Image of WinEdt: user has typed source code]
  {pictures/winedt/03}
  {\caption{WinEdt}}
  {fig:winedt3}
\end{figure}

You can now save your document using the \menu{File \menuto\ Save as}
menu.  Select the file type to be \texttt{TeX}, and type in the name
of your file, e.g.\ \texttt{sample1.tex}.  See
\figureref{fig:winedt4}.

\begin{figure}[htbp]
  \figconts[alt=Image save file dialog]
  {pictures/winedt/04}
  {\caption{WinEdt --- Saving the File}}
  {fig:winedt4}
\end{figure}

\refstepcounter{object}\label{winedt:latex}
To \LaTeX\ your document, simply click on the \LaTeX\ button
\incGraphics[alt=Image of LaTeX icon]{pictures/winedt/latex}.
The output will appear in an MSDOS Prompt window (see
\figureref{fig:winedt5}).

\begin{figure}[htbp]
\figconts[alt=Image of MSDOS prompt showing standard LaTeX messages]
  {pictures/winedt/05}%
  {\caption[WinEdt --- LaTeX Output]{WinEdt --- \LaTeX\ Output}}
  {fig:winedt5}
\end{figure}

\refstepcounter{object}\label{winedt:yap}
To view your typeset document, click on the \dq{view DVI}\indexDVI\ 
button \incGraphics[alt=Image of previewer icon]{pictures/winedt/yap}.

\refstepcounter{object}\label{obj:winedtdvips}
You can convert your DVI file\indexDVI\ to \Index{PostScript} by clicking 
the \incGraphics[alt=Image of dvips icon]{pictures/winedt/dvips} button.  
If you have \iappname{GSview} installed, you can then view
the PostScript file by clicking on the 
\incGraphics[alt=Image of GSView icon]{pictures/winedt/gsview2} button.  

Depending on which version of \iappname{WinEdt} you have installed,
there may also be a \iPDFLaTeX\ button which you can click on to
create a Portable Document Format (\texttt{.pdf}) document.  If not,
you can click on the 
\incGraphics[alt=Image of MSDOS icon]{pictures/winedt/msdos} button to 
open up an MS-DOS Prompt window, and type \texttt{pdflatex} followed by 
the filename. For example:
\begin{verbatim}
pdflatex sample1.tex
\end{verbatim}
Note that if the filename contains a
space, you will need to use double quotes:
\begin{verbatim}
pdflatex "my file.tex"
\end{verbatim}

\fi

\backmatter

%%%%%%%%%%%%%%%%%%%%%%%%%%% Bibliography %%%%%%%%%%%%%%%%%%%%%%%%%%%%%%%%

\clearpage\phantomsection
\setnode{bibliography}
\bibliographystyle{plain}
\bibliography{novices}

%%%%%%%%%%%%%%%%%%%%%%%%%%%% GLOSSARIES %%%%%%%%%%%%%%%%%%%%%%%%%%%%

\setnode{acronyms}
\printglossary[type=acronym,nonumberlist,style=acronyms]

\setnode{summary}
\setdoublecolumnglo
\printglossary[style=summary]

\setnode{bookindex}
\printindex

%%%%%%%%%%%%%%%%%%%%%%%%% LICENCE %%%%%%%%%%%%%%%%%%%%%%%%%%%%%%%%%%

\setnode{licence}
\input{fdl}

%%%%%%%%%%%%%%%%%%%%%%%%%% HISTORY %%%%%%%%%%%%%%%%%%%%%%%%%%%%

\setnode{history}
\chapter{History}\markright{History}

\historyitem{25th Sept 2012 (Version 1.4)}

\begin{itemize}
\item Paperback edition~1 produced.
\item Added TeXWorks section.
\item Moved \dq{Text editor and Terminal approach}, \dq{TeXnicCenter}
and \dq{WinEdt} sections to new supplementary material.
\item Added hardcopy-related code.
\item Change to \koma\ classes (both for examples and for pdf
versions of this document).
\item Changed from using \sty{subfloat} to \sty{subcaption} package.
\item Added section on inter-sentence spacing.
\item Moved \dq{Errors} and \dq{Where to get Help?} to appendices.
\item Moved section \dq{Downloading and Installing Packages} to
new appendix chapter.
\item Moved introduction to packages to \dq{Creating a Simple
Document}.
\item Moved datetime section to \dq{Creating a Simple Document}
chapter. (Removed reference to ukdate package.)
\item Moved babel section to \dq{Structuring Your Document} chapter.
\item Moved graphicx section to its own chapter.
\item Added section on align.
\item Added \cmdname{cfrac}, \cmdname{substack} and \sty{amsmath} ellipses to maths chapter.
\item Added extensible arrows and \cmdname{bigl} etc to maths
chapter.
\item Added booktabs.
\item Moved lengths chapter to section in definitions.
\item Added summary chapter with commands hyperlinked to their
definitions in the summary.
\item Changed definitions chapter to use a glossary structure.
\item Moved bibliography into bib file.
\item Added \sty{varioref}.
\item Removed dependency on \sty{html} package (for pdf versions)
to avoid conflict between \sty{html} and \sty{varioref}
(\sty{html} package functions not defined by \sty{hyperref}
now emulated; \sty{comment} package loaded to provide
\envname{htmlonly} environment).
\item Removed image of equation written in Word (Microsoft have
improved their equation rendering) and added link to
Murray Sargent~III blog~\cite{sargentiii}.
\item Added section on what a terminal/command prompt is.
\item Added section on auxiliary files.
\item Added section on Perl.
\item Added information about latexmk
\item Mentioned grffile package.
\item Mentioned on-the-fly EPS conversion.
\item Mentioned \sty{etoolbox}'s \cmdname{appto} and \sty{babel}'s
\cmdname{addto}.
\item Changed to UTF-8 and mostly changed to using code points
instead of named entities in HTML files.
\item Moved the document's home page from
\url{http://theoval.cmp.uea.ac.uk/~nlct/latex/novices/} to
\url{http://www.dickimaw-books.com/latex/novices/}.
\end{itemize}

\historyitem{15th Jan 2008 (Version 1.3)}

The main reason behind this change was to increase accessibility and
conform to W3C guidelines. If you are experiencing problems relating
to accessibility, please let me know (clearly stating the problem).
\begin{itemize}
\item Corrected error in the university's post code on the
title page
\item Added alternative text tags to more of the images, and made
some of the images hyperlinks to a more detailed description of the
image.
\item Added information on how to break ligatures.
\item Moved information on TeX to the introduction, and removed
section on TeX that was in the "Some Definitions" chapter.
\item Document nodes now have permanent names instead of the generic
node\meta{n}.html which \LaTeX2HTML generates by default.
\item Went back to using straight double quotes in the HTML document
as the fancy typographic double quotes are nonstandard.

\end{itemize}

\historyitem{8th May 2007 (Version 1.2)}

\begin{itemize}
\item Links to \gls{ukfaq} added.
\item Overview made into a separate section, and tidied up a bit.
\item Added some extra definitions: moving arguments and fragile
commands, robust commands, short and long commands.
\item Changed \dq{Text editor and Terminal approach} to deal with
Unix-type systems rather than MS-DOS.
\item Moved section on \envname{tabular} environment.
\item Added section on boxes and mini-pages.
\item Segmented section on font changing commands.
\item Segmented section describing \sty{graphicx}.
\item Added section on the \sty{babel} package.
\item Updated and segmented section on downloading and installing
new packages.
\item Added section on side-by-side figures.
\item Updated section on sub-figures to use the new \sty{subfloat}
package instead of the obsolete \sty{subfigure} package.
\item Added \dq{Need More Help?} chapter.
\end{itemize}

\backcoverheading
(See \url{http://www.gnu.org/licenses/fdl-howto-opt.html#SEC2}.)

\backcovertext
\end{document}

}


\glsaddall[types=keywords]

%%%%%%%%%%%%%%%%%%%%%% SOURCE TO PDF %%%%%%%%%%%%%%%%%%%%%%%%%%

\setnode{fromsource2output}
\chapter{From Source Code to Typeset Output}
\label{ch:tex2pdf}

Every time you want to create or edit a \LaTeX\ document, there are three
basic steps you will always need to follow:

\begin{enumerate}
\item\label{itm:step1} Write or edit the \gls{source}.

\item\label{itm:step2} Pass the source code to the \iappname{latex} or \iappname{pdflatex} application (\dq{\LaTeX\ the document}).

\begin{itemize}
\item If there are any error messages, return to \objectref{Step}{itm:step1}.
\item If there are no error messages, a 
\glslink{output}{PDF file}\indexPDF\ is created.
\end{itemize}

\item\label{itm:step3} View the \glslink{output}{PDF file} to check
the result.  If you need to modify your document, go back to
\objectref{Step}{itm:step1}. 
\end{enumerate}

You will therefore need:

\begin{enumerate}
\item A text editor  (to perform \objectref{Step}{itm:step1}). For
example \iappnamelink[,]{Vim}{http://www.vim.org/} 
\iappnamelink{Emacs}{http://www.gnu.org/software/emacs/} or
\iappnamelink[.]{Gedit}{http://projects.gnome.org/gedit/}

\refstepcounter{object}\label{obj:MiKTeX}%
\item The \gls{tex} software\faq{(La)TeX for different
machines}{TeXsystems} (to perform \objectref{Step}{itm:step2}). If
you don't already have \TeX\ on your machine, you will need to
install it. The most convenient way to do this is
to install from the \TeX\ Collection DVD ROM, which is distributed
to all \gls{tug} members, but you can also download and install free
\TeX\ distributions, such as \itexdistro{TeX~Live},
\itexdistro{MiKTeX} or \itexdistro{MacTeX}, from the Internet (see
\xpageref{below}{obj:install}).  There is also \itexdistro{proTeXt},
an enhancement of \texdistro{MiKTeX} that aims to be an
easy-to-install \TeX\ Distribution.  For more information including
up-to-date links, go to \url{http://www.ctan.org/starter.html}.

\item A PDF viewer (to perform
\objectref{Step}{itm:step3}). For example 
\iappnamelink[,]{Adobe Reader}{http://www.adobe.com/} 
\iappnamelink[,]{Sumatra}{http://blog.kowalczyk.info/software/sumatrapdf/free-pdf-reader.html}
\iappnamelink{Evince}{http://projects.gnome.org/evince/} or
\iappnamelink[.]{Okular}{http://okular.kde.org/}

\end{enumerate}

This can be rather complicated for a beginner, especially for those
with no experience writing computer code. Fortunately, there are
some all-in-one applications (often called a \keyword{front-end})
that provide a text editor (for \objectref{Step}{itm:step1}),
buttons or menu items to run the \appname{latex} or
\appname{pdflatex} \glslink{terminal}{command-line application} (for
\objectref{Step}{itm:step2}) and, in some cases, a viewer to perform
\objectref{Step}{itm:step3}.

\sectionref[The next section]{sec:texworks} describes one such front-end called
\iappnamelink[.]{TeXWorks}{http://www.tug.org/texworks/} I~have chosen
to describe TeXWorks because it is a free, cross-platform
application. Binaries are available for Microsoft Windows, Mac~OS~X
and GNU/Linux. The screen shots of TeXWorks in this
book were taken from the Linux version running
under Fedora. If you run TeXWorks on other operating systems, it may
have a slightly different look, but it has the same functionality.

New versions of \texdistro{TeX~Live} and \texdistro{MiKTeX} include TeXWorks for MS~Windows,
and new versions of \texdistro{MacTeX} include TeXWorks for Mac~OS~X users.
GNU/Linux users can use their Add\slash Remove Software utility to 
install TeXWorks. Alternatively, you can download TeXWorks by
following the links provided at \url{http://www.tug.org/texworks/}.

\refstepcounter{object}\label{obj:install}%
If you're confused by all the options, let's keep
things as simple as possible:

\begin{itemize}

\item MS~Windows:

 You have a choice between \texdistro{MiKTeX} (or \texdistro{proTeXt}) and 
 \texdistro{TeX~Live}. \texdistro{MiKTeX} provides a
 smaller and quicker installation, but the downside is that you may
 not have the classes or packages you want to use. \texdistro{MiKTeX} can
 install these whenever you try to \LaTeX\ a document that uses them, 
 but you need an Internet connection while it does this. 
 \texdistro{TeX~Live} installs everything, so it takes longer and needs 
 more space, but you should have the majority of packages and classes that you need.

 \begin{description}

   \item[TeX~Live:]\mbox{}

   \begin{enumerate}
    \item Fetch and unpack
     \url{http://mirror.ctan.org/systems/texlive/tlnet/install-tl.zip}

    \item Run \iappname{install-tl} and follow the instructions.
     This can take an hour or more.
   \end{enumerate}

   \item[proTeXt:]\mbox{}

   \begin{enumerate}
    \item Go to \url{http://tug.org/protext/}

    \item Click on the \dq{download the self-extract protext.exe} link
     to download and run the executable.
   \end{enumerate}

   \item[MiKTeX:]\mbox{}

   \begin{enumerate}
    \item Go to \url{http://www.miktex.org/}

    \item In the left-hand panel, there is a link to the download page
     for the latest version. At time of writing, it is \texdistro{MiKTeX}~2.9.
     Click on that link.

    \item Scroll down to the section \dq{Installing a basic
      \texdistro{MiKTeX} system}.

    \item If you're happy with the selected mirror location, click on the 
      \dq{Download} button.

    \item Run the executable.
   \end{enumerate}
 \end{description}

\item Mac~OS~X:

 \begin{enumerate}
  \item Go to \url{http://tug.org/mactex/}

  \item Follow the instructions on that page.
 \end{enumerate}

\item GNU/Linux:

 \begin{enumerate}
  \item Fetch and unpack
    \url{http://mirror.ctan.org/systems/texlive/tlnet/install-tl-unx.tar.gz}

  \item Follow the instructions at
    \url{http://tug.org/texlive/quickinstall.html}

  \item Once \texdistro{TeX~Live} has finished installing, run your
   system's \dq{Add\slash Remove Software} tool.

  \item Find \dq{texworks}, select the newest version and install.
 \end{enumerate}

\end{itemize}

If you run into problems, there are mailing lists at
\url{http://tug.org/mailman/listinfo/tex-live} and
\url{http://docs.miktex.org/manual/lists.html} for
\itexdistro{TeX~Live} and \itexdistro{MiKTeX}, and
\itexdistro{MacTeX} help at \url{http://www.tug.org/mactex/help/}.
There is also a list of places where you can ask for help in
\appendixref{ch:help}.

%%%%%%%%%%%%%%%%%%% TeXWorks %%%%%%%%%%%%%%%%%%%%%

\setnode{texworks}
\section{TeXWorks}
\label{sec:texworks}

Hopefully you've managed to successfully install \gls{tex} and TeXWorks as
described \latexhtml{above}{in the \htmlref{previous
section}{ch:tex2pdf}}, so let's test it out.

First run TeXWorks. On Windows, you can access it via the
\menu{Start} menu. On GNU/Linux, it's probably located in
\menu{Applications}\menuto\menu{Office}, or you can type
\iappname{texworks} in a \gls{terminal}. You should now see the TeXWorks
window. The button marked
with a grey triangle in a green circle is the build or typeset function. It
runs the application in the drop-down list next to it. This is set to
pdfLaTeX, which is what we want for now.

It's a good idea to switch on the syntax highlighting, if it isn't
already on. This is done via the \menu{Format}\menuto\menu{Syntax
Coloring} sub-menu. Make sure the \menu{LaTeX} item is selected.

Next, type in the following sample source code, as shown in \figureref{fig:texworks2}
(the commands used here will be described in more detail in
\chapterref{ch:simpledoc}):
\begin{code}
\begin{verbatim}
\documentclass{scrartcl}

\begin{document}

This is an example document.

\end{document}
\end{verbatim}
\end{code}
\textbf{Pay close attention to the backslashes at the start of each
command name.} If you find the font is a bit too small for you,
you can make it larger using the \menu{Format}\menuto\menu{Font} menu item.
\emph{This doesn't affect the font size in your PDF file, just the
font size of your code.} This displays the \dq{Select Font} dialog
box. Set the font size as appropriate.

Then save the document, using the
\menu{File}\menuto\menu{Save As} menu item. I~called my document
\texttt{example1.tex} (remember the \texttt{.tex} extension and
stick to file names that only consist of alphabetical characters,
digits and hyphens\dash don't uses spaces or underscores).

\begin{figure}[htbp]
\figconts
  {pictures/texworks2}
  {%
    \caption[Source Code for an Example Document]{Source Code for an 
     Example Document. (Syntax highlighting switched on.)}%
  }%
  {fig:texworks2}
\end{figure}

Now that you have saved the file, you can run pdfLaTeX. Make sure
the drop-down list next to the build button has \dq{pdfLaTeX}
selected and click on the build button. If all goes well, a new
window should open displaying the typeset document
(\figureref{fig:texworks3}).

\begin{figure}[htbp]
\figconts
  {pictures/texworks3}
  {%
    \caption{The Typeset Document}
  }%
  {fig:texworks3}
\end{figure}

\begin{figure}[htbp]
\figconts
  {pictures/texworks4}
  {%
    \caption{The Source Code Has a Misspelt Command}
  }%
  {fig:texworks4}
\end{figure}

Now let's see what happens if there is an error in the source code.
In \figureref{fig:texworks4} I~have misspelt the
\glsni{documentclass} command. This time, when I~click on the
build button, I~get the error message:
\begin{verbatim}
! Undefined control sequence.
l.1 \documentclas
                 {scrartcl}
?
\end{verbatim}
(Shown in \figureref{fig:texworks5}.)

\begin{figure}[htbp]
\figconts
  {pictures/texworks5-annote}
  {%
     \caption{An Error Message is Displayed}%
  }%
  {fig:texworks5}
\end{figure}


Here \dq{Undefined control sequence} means an unrecognised command,
and below that message, \dq{\texttt{l.1}} means the error was encountered on
line~1. An input line at the bottom of the window has appeared with a
cursor. \LaTeX\ is in interactive mode
and is awaiting a response. There are several responses, but I'm
only going to mention two of them:
\begin{enumerate}
\item Type \texttt{h} and press the Return/Enter key \enter. This displays
a short help message and awaits a new response (see
\figureref{fig:texworks6}).

\item Type \texttt{x} and press Return/Enter. This aborts the
\LaTeX\ run.
\end{enumerate}

Notice that the green circle button with the grey triangle has turned
into a red stop button. This button can be used to abort the
process instead of typing \texttt{x}.

\begin{figure}[htbp]
\figconts
  {pictures/texworks6}
  {%
    \caption{A Short Help Message}%
  }%
  {fig:texworks6}
\end{figure}

Now, there is a second tab at the bottom of the TeXWork's window
(\figureref{fig:texworks7}). This lists the error message and
provides a link to the line where the error occurred. Clicking on
this link highlights line~1. Next I~need to fix the error by
correcting the spelling of the command. Once it's fixed, I~can click
on the build button.

\begin{figure}[htbp]
\figconts
  {pictures/texworks7}
  {%
    \caption{Error Tab}%
  }%
  {fig:texworks7}
\end{figure}

Here's another error you might encounter: I'm now going to misspell
the class name. It should be \texttt{scrartcl}, but in
\figureref{fig:texworks8} it has been misspelt. This time, when I
click on the build button, I~get the error:
\begin{verbatim}
! LaTeX Error: File `scrartc.cls' not found.
\end{verbatim}
I~have two choices: type in the correct name on the line below
\dq{Enter file name:} or I~can abort the process using the red abort
button. In either case I~need to go back and fix the error in my
code.

\begin{figure}[htbp]
\figconts
  {pictures/texworks8}
  {%
   \caption{Misspelt Class File}%
  }%
  {fig:texworks8}
\end{figure}


%%%%%%%%%%%%% CREATING A SIMPLE DOCUMENT %%%%%%%%%%%%%%%%%%

\setnode{simpledoc}
\chapter{Creating a Simple Document}
\label{ch:simpledoc}

Having installed and tested the software, let's now look at how to
actually write the \gls{source}.  The very first line of any document that you create must have
the \gls{command}:
\begin{definition}
\gls{documentclass}\oarg{\meta{option-list}}\marg{\meta{class-name}}
\end{definition}%
This tells \LaTeX\ what type of document you want to create
(for example an article, a technical report or correspondence).  The
\glsni{documentclass} command takes one \gls{mandatory}, \meta{class-name},
that specifies the \gls{cls}.

There are many class files available, and some publishers,
institutions or journals provide their own custom classes (for
example, the \icls{jmlr} class for the Journal of Machine Learning
Research). Popular classes include \icls{memoir} (for books and
reports) and those supplied in the \koma\ bundle (for books,
reports, articles and correspondence). There's also \icls{beamer}
(for presentations) as well as classes for typesetting exams,
flashcards, concert programmes etc. For simplicity, this
book will concentrate on three of the
\koma\ classes \icls{scrartcl} (for articles), \icls{scrreprt}
(for technical reports, theses etc) and \icls{scrbook} (for books).

We'll start with a very simple document, so let's use the
\icls{scrartcl} class file.  In this case the very first line of the
\gls{source} should be:
\begin{codeS}
\glsni{documentclass}\marg{scrartcl}
\end{codeS}
The \glsni{documentclass} command also takes an 
\gls{optional}, \meta{option-list},
which should be a comma separated list of options to be passed to the
class file.  This allows you to override the class file defaults.
For example, the \icls{scrartcl} class file by default uses A4 
paper, but if you are in the USA you will probably want to use
letter paper.  This can be achieved
using the option \clsopt{letterpaper}.  So you would need to edit the
above line to:
\begin{codeS}
\glsni{documentclass}\oarg{letterpaper}\marg{scrartcl}
\end{codeS}
\refstepcounter{object}\label{obj:sizeopt}%
Let's change another option.  The normal font size is \clsopt{11pt}
by default, but we have the option to change it, so let's use
\clsopt{12pt}:
\begin{codeS}
\glsni{documentclass}\oarg{letterpaper,12pt}\marg{scrartcl}
\end{codeS}
You can also change your document so that it is in a two-column
format using the \clsopt{twocolumn} option:
\begin{codeS}
\glsni{documentclass}\oarg{letterpaper,12pt,twocolumn}\marg{scrartcl}
\end{codeS}

After deciding what type of document you want, you now need to specify
the contents of the document.  This is done inside the \gls{env-document}
\gls{environment}.  The document is started with the
command:
\begin{codeS}
\glsni{begin}\marg{document}
\end{codeS}
and ended with
\begin{codeS}
\glsni{end}\marg{document}
\end{codeS}
(\LaTeX\ stops reading the file when it reaches the above line, so
anything occurring after it is ignored.)

My \gls{source} now looks like:
\begin{code}
\glsni{documentclass}[12pt]\marg{scrartcl}\newline
\mbox{}\newline
\glsni{begin}\marg{document}\newline
\mbox{}\newline
\glsni{end}\marg{document}
\end{code}
Every document you create must have this form.  You can't simply
start typing the document text.  You must first specify
your class file\indexCLS, and then place the contents of the document
inside the \glsni{env-document} environment.\screenpagebreak

So far so good, but at the moment we have an empty document, so we
won't get any output.  Let's now put some text into our document:
\begin{code}
\begin{alltt}
\glsni{documentclass}\oarg{12pt}\marg{scrartcl}

\glsni{begin}\marg{document}

This is a simple document.
Here is the first paragraph.

Here is the second paragraph. As you
can       see it's 
a rather 
short paragraph, but not as short as the previous one.

\glsni{end}\marg{document}
\end{alltt}
\end{code}

\xminisec{Top Five Mistakes Made by New Users}
I first started teaching \LaTeX\ in 1998, and these are the most
common errors I've seen when people start learning \LaTeX:
\begin{enumerate}
\item Missing out the backslash \glsni{backslashchar} at the start
of one or more of the \glspl{command}.

\item Using a forward slash \glsni{slashchar} instead of a backslash
\glsni{backslashchar}.

\item Forgetting \texttt{\glsni{end}\marg{document}}.

\item Misspelling \dq{document} (in
\texttt{\glsni{begin}\marg{document}} and
\texttt{\glsni{end}\marg{document}}).

\item Missing a closing brace \glsni{rightbracechar}.

\end{enumerate}
If you encounter any problems when you start out, go through that check list first.
Then check \appendixref{ch:errors}.

Whenever you start a new document, always type in the
\glsni{documentclass}, \texttt{\glsni{begin}\marg{document}} and
\texttt{\glsni{end}\marg{document}} commands first
(\figureref{fig:newdoc-1}). Then move your cursor between the
\glsni{begin} and \glsni{end} lines and type the document text
(\figureref{fig:newdoc-2}).

\begin{figure}[htbp]
\figconts
  {pictures/newdoc-1}
  {%
    \caption[Starting a New Document (1)]{Starting a New Document: 
    always type these three lines first.}%
  }%
  {fig:newdoc-1}

\ifscreenorother
{%
  \end{figure}
  \begin{figure}[htbp]
}%
{%
  \strut
}

\figconts
  {pictures/newdoc-2}
  {%
    \caption[Starting a New Document (2)]{Starting a New Document: 
     move the cursor inside the \glsni{env-document} environment and 
     start typing the document text.}%
  }%
  {fig:newdoc-2}
\end{figure}

\begin{exercise}{Simple Document}{ex:simpledoc}

Try typing the code in the above example into TeXWorks or the editor
of your choice (see \chapterref[the previous section]{ch:tex2pdf} if
you can't remember what to do.) You can also \download{simpledoc} a
copy of this file, but I~recommend that you try typing it in to give
yourself some practice.

Things to note while you are typing: firstly, when you press the
return character at the end of a line this end-of-line character is
converted into a space in the \gls{output}.  So the
fact that I~have some very ragged lines in my \gls{source} has
no effect on the final result. (Note that some front-ends will 
reformat your lines as you type.) Whereas a completely blank line
will be converted into a paragraph break (\gls{par} has the same
effect).

Secondly, multiple \Index{spaces} are converted into a single space,
so the large gap between the words \dq{\texttt{can}} and \dq{\texttt{see}} is
no different from having a single space.

Once you have typed up your \gls{source}, save your file 
(called, for example, \texttt{exercise1.tex}), and run PDF\LaTeX\ as described in 
\sectionref{sec:texworks}. If all goes well, TeXWorks should display
the resulting PDF file in a new window, usually alongside the window
containing the source code.

\xminisec{Notes:}
\begin{enumerate}

\item Each paragraph automatically starts with an indentation in the
PDF.

\item There is no blank line between the paragraphs in the PDF
document. (See what happens if you add the \koma\ class option
\scrclsopt[full]{parskip}:
\begin{codeS}
\glsni{documentclass}\oarg{12pt,parskip=full}\marg{scrartcl}
\end{codeS}
and rebuild the PDF.)

\item Move the mouse over one of paragraphs in the PDF viewer and
pop-up the context menu (usually a right mouse click). Select
\menu{Jump to Source}. The window containing the source code should
now gain the focus and the line of code matching the typeset line you clicked on
in the PDF should now be highlighted.

\end{enumerate}

\end{exercise}

%%%%%%%%%%%%%%%%%% Simple Commands %%%%%%%%%%%%%%%%%%

\setnode{simplecommands}
\section{Using Simple Commands}
\label{sec:simple}

\faq{Typesetting all those TeX-related logos}{logos}Now let's try
adding a few simple \glspl{command} to our
document. The command \gls{LaTeX} produces the
\makeimg{ALT=LaTeX}{\LaTeX} logo and the command \gls{today} prints
the current date. \LaTeX\ always ignores any spaces that follow a
command name that consists of letters, as it uses the space to
indicate the end of the command name. This means that if we want a
space to occur immediately after the command, we need to explicitly
say so using the command \gls{spacesym} (recall from
\ifscreenorother{\htmlref{earlier}{obj:visiblespace}}{page~\pageref{obj:visiblespace}}
\gls{visiblespace} indicates a space character). So, for example:
\begin{codeS}
\glsni{LaTeX}\glsni{spacesym}logo
\end{codeS}%
produces the output:
\begin{resultS}[LaTeX logo]
\LaTeX\ logo
\end{resultS}%
Some people when starting out can get a bit confused by this and
read it as the entity \dq{\cmdname{LaTeX}\cmdname{}} whereas it is
in fact two commands: \dq{\cmdname{LaTeX}} (print the \LaTeX\ logo)
followed \dq{\cmdname{\ }} (print a space.)

Let's also try using a command that takes an 
\gls{argument}.  The command 
\begin{definition}
\gls{footnote}\refstepcounter{object}\label{cmd:footnote}\marg{\meta{text}}
\end{definition}%
takes one argument that specifies the text that should
appear in the footnote.  This command must be placed where you want
the footnote marker to appear.

\begin{exercise}{Using Simple Commands}{ex:simple}

Edit the document you created in \exerciseref{ex:simpledoc},
so that it looks like the following: (You can \download{simple} it if
you like, but again it is better if you try typing it in yourself.)
\begin{bcode}
\glsni{documentclass}\oarg{12pt}\marg{scrartcl}\newline
\mbox{}\newline
\glsni{begin}\marg{document}\newline
\mbox{}\newline
This is a simple \glsni{LaTeX}\glsni{spacesym}document.
Here is the first paragraph.\newline
\mbox{}\newline
Here is the second paragraph\glsni{footnote}\marg{with a footnote}. 
As you can see it's a rather short paragraph, but not 
as short as the previous one. This document was 
created on: \glsni{today}.\newline
\glsni{end}\marg{document}
\end{bcode}

Now \htmlref{\LaTeX\ your document}{ch:tex2pdf} and view the result.
(Remember to check \appendixref{ch:errors}
if you have a problem.) You should see the \LaTeX\ logo, the footnote
marker and the current date.  If you scroll down to the bottom of the
page, you should see the footnote.
\end{exercise}

%%%%%%%%%%%%%%%%%%%%%%%%%% PACKAGES %%%%%%%%%%%%%%%%%%%%%%%%

\setnode{packages}
\section{Packages}
\label{sec:packages}

Packages\indexSTY\ are files with the extension \texttt{.sty} that
either define new \glspl{command} or redefine
existing commands\faq{What are LaTeX classes and packages?}{clsvpkg}.
They're like a type of add-on or plug-in. Most of the commonly used
packages should have been installed when you installed your \gls{tex}
distribution (see \chapterref{ch:tex2pdf}).
\appendixref{ch:installsty} covers how to install new packages.
Most packages come with documentation that can be accessed using the
\iappname{texdoc} application described in \sectionref{sec:texdoc}.

Packages are loaded in the \gls{preamble} (after
\gls{documentclass} and before \glsni{begin}\marg{document}) using
\begin{definition}
\gls{usepackage}\oarg{\meta{option list}}\marg{\meta{package}}
\end{definition}
where \meta{package} is the name of the package and \meta{option
list} is a list of comma-separated options. For example, to load the
package \isty{graphicx} with the \istyopt{graphicx}{draft} option:
\begin{codeS}
\glsni{usepackage}\oarg{draft}\marg{graphicx}
\end{codeS}
Any applicable class options are also passed to packages, so in
\begin{code}
\gls{documentclass}\oarg{draft}\marg{scrartcl}\newline
\glsni{usepackage}\marg{graphicx}
\end{code}
the \clsopt{draft} option is set for both the \icls{scrartcl} class
and the \isty{graphicx} package.

You can specify more than one package in the argument of
\glsni{usepackage}, where each package name is separated by a
comma. For example:
\begin{codeS}
\glsni{usepackage}\marg{amsmath,amsfonts}
\end{codeS}
The \isty{graphicx} package is covered in \chapterref{ch:graphicx}
and the \isty{amsmath} package is covered in \chapterref{ch:maths},
so let's start out with a~relatively simple example.

\setnode{datetime}
\subsection{Changing the Format of
\latexhtml{\texorpdfstring{\protect\cmdname{today}}{\protect\textbackslash
today}}{\protect\cmdname{today}}}
\label{sec:today}

In the \htmlref{previous exercise}{ex:simple}, we used the
\gls{today} command to produce the current date. By default, this
command displays the date in US format. To illustrate how to use
packages, this section will look at how to use the \isty{datetime}
package to change the way that \gls{today} displays the date.

The \isty{datetime} package has various options that can be used to
change the format of \gls{today}.  For example, by default the
\isty{datetime} package redefines \gls{today} to display the
date in the form: {\showdowtrue\longdate
\renewcommand{\fmtord}[1]{\textsuperscript{#1}}\docdate}. The option
\istyopt{datetime}{short} will produce an abbreviated form, (for
example {\shortdate\docdate}) and the option \istyopt{datetime}{nodayofweek}
won't display the day of the week (for example
{\showdowfalse\docdate}).
For those who don't like the raised ordinal, there is the
\istyopt{datetime}{level} option.
These can be passed as a comma separated list in the
\gls{optional} to the
\gls{usepackage} command.  It is also possible to use a
declaration instead. For example, to redefine \gls{today} to
display the date in the form {\ddmmyyyydate\docdate}, you can either do
\begin{codeS}
\glsni{usepackage}\oarg{ddmmyyyy}\marg{datetime}
\end{codeS}%
or
\begin{code}
\begin{alltt}
\glsni{usepackage}\marg{datetime}
\gls{ddmmyyyydate}
\end{alltt}
\end{code}

The \isty{datetime} package also defines the command
\begin{definition}
\gls{currenttime}
\end{definition}
which displays the current time, where again
the format can be changed by the package options. So the option
\istyopt{datetime}{12hr} will cause \glsni{currenttime} to display
the date in 12 hour format (for example, \ampmtime) and the option
\istyopt{datetime}{24hr} will cause \glsni{currenttime} to display
the date in 24 hour format (for example, \xxivtime).

\begin{exercise}{Using the datetime Package}{ex:datetime}

Edit your document from \exerciseref{ex:simple} so that it uses the \isty{datetime}
package. Experiment with the different package options, for example
\begin{codeS}
\glsni{usepackage}\oarg{short,nodayofweek,level,12hr}\marg{datetime}
\end{codeS}
and add the current time
\begin{codeS}
This document was created on: \glsni{today}\glsni{spacesym}at \glsni{currenttime}.
\end{codeS}
You can \downloadorview{datetime} an example.
For a full list \isty{datetime} of package options, see the
\sty{datetime} \htmlref{documentation}{sec:texdoc}.
\latex{(Refer to \sectionref{sec:texdoc} on how to find package
documentation.)}

\end{exercise}

%%%%%%%%%%%%%%%%%%%%%% Special Characters %%%%%%%%%%%%%%%%%%%%

\setnode{symbols}
\section{Special Characters and Symbols}
\label{sec:chars}

You can use most of the standard characters that you find on your keyboard, but the 
10 symbols shown in \tableref{tab:specialchars} have a special
meaning.

\begin{table}[htbp]
\caption{Special Characters}
\label{tab:specialchars}
\centering
\begin{makeimage}\end{makeimage}
\glsi{leftbracechar} \glsi{rightbracechar} 
\glsi{percentchar}
\glsi{ampchar} 
\glsi{dollarchar} 
\glsi{hashchar} 
\glsi{underscorechar}
\glsi{circumchar} 
\glsi{tildechar} 
\glsi{backslashchar}
\end{table}

We have already used the curly braces \gls{leftbracechar} and
\gls{rightbracechar}. The percent symbol \gls{percentchar} is a
comment character. Everything from the percent symbol up to the end
of line is ignored by \LaTeX. This means you can have comments in
your \gls{source} to remind you what a particular part of your code
is doing. We have also used the backslash symbol
\glsni{backslashchar} which indicates that we are using a \LaTeX\
command, as in \glstext{LaTeX} or \glstext{today}. The meaning of
the other special characters will be covered later.

So what do you do if you want one of these symbols to actually appear
in your document? \faq{Where can I find the symbol for
\ldots?}{symbols} \tableref{tab:symbols} lists commands that produce
these and other symbols\faq{How to get copyright, trademark,
etc}{tradesyms}. Note that some of the commands have short cuts, such
as \glsni{emdash} instead of \glsni{textemdash} and
\glsni{questiondown} instead of \glsni{textquestiondown}.

The symbol \texttt{`} is the backtick (or grave) symbol, as
opposed to the apostrophe symbol~\texttt{'}.  \latexhtml{The backtick
symbol usually looks like~\texttt{\`{}} on a keyboard, and on}{On}
most UK keyboards it is situated to the left of the \texttt{1}~key.
The opening double quote is created using two adjacent backtick
symbols and the closing double quote with two adjacent
apostrophe symbols. This gives 66 and 99 style quotes, which you
wouldn't get using the double quote symbol on your keyboard.

Note that the symbols \gls{barchar} \gls{ltchar} and \gls{gtchar} have to
be created using \glsni{textbar}, \glsni{textless} and
\glsni{textgreater} when in normal text mode. If you try to enter
them using the corresponding keyboard characters you may get
\makeimg{em dash}{\textemdash} \makeimg{upside down exclamation
mark}{\textexclamdown} and \makeimg{upside down question
mark}{\textquestiondown}. (They do however work if you are in
\htmlref{maths mode}{ch:maths}.\footnote{There are also some text
fonts that will display them correctly, but don't rely on it.}) The
slash character \gls{text.slash} may be used directly, as in
\dq{\texttt{and\glsni{text.slash}or}}, but no line break will be permitted at the slash,
whereas \glsni{slash} (as in \dq{\texttt{and\glsni{slash}\glsni{visiblespace}or}}) will allow a
line break at that point.

\begin{table}[htbp]
\caption{Symbols}
\label{tab:symbols}
\centering
\begin{latexonly}
\begin{tabular}{@{}ll@{\hspace{4\tabcolsep}}ll@{\hspace{4\tabcolsep}}ll@{}}
\begin{inlinedef}\gls{textbackslash}\end{inlinedef} & \textbackslash&
\begin{inlinedef}\gls{slash}\end{inlinedef} & \slash &
\begin{inlinedef}\gls{textgreater}\end{inlinedef} & \textgreater \\
\begin{inlinedef}\gls{textasciicircum}\end{inlinedef} & \textasciicircum & 
\begin{inlinedef}\gls{dollar}\end{inlinedef} & \$ &
\begin{inlinedef}\gls{textbar}\end{inlinedef} & \textbar \\
\begin{inlinedef}\gls{textasciitilde}\end{inlinedef} & \textasciitilde&
\begin{inlinedef}\gls{leftbrace}\end{inlinedef} & \{ &
\begin{inlinedef}\gls{textless}\end{inlinedef} & \textless \\
\begin{inlinedef}\gls{pounds}\end{inlinedef} & \pounds &
\begin{inlinedef}\gls{rightbrace}\end{inlinedef} & \} &
\begin{inlinedef}\gls{dag}\end{inlinedef} & \dag \\
\begin{inlinedef}\gls{textregistered}\end{inlinedef}& 
\textregistered&
\begin{inlinedef}\gls{hash}\end{inlinedef} & \# &
\begin{inlinedef}\gls{ddag}\end{inlinedef} & \ddag \\
\begin{inlinedef}\gls{texttrademark}\end{inlinedef} & \texttrademark &
\begin{inlinedef}\gls{percent}\end{inlinedef} & \% &
\begin{inlinedef}\gls{quoteright}\end{inlinedef}
or \begin{inlinedef}\gls{textquoteright}\end{inlinedef} & \textquoteright\\
\begin{inlinedef}\gls{copyright}\end{inlinedef} & \copyright &
\begin{inlinedef}\gls{amp}\end{inlinedef} & \& &
\begin{inlinedef}\gls{quoteleft}\end{inlinedef} 
or \begin{inlinedef}\gls{textquoteleft}\end{inlinedef} & \textquoteleft \\
\begin{inlinedef}\gls{textbullet}\end{inlinedef} & \textbullet & 
\begin{inlinedef}\gls{i}\end{inlinedef} & \i & 
\begin{inlinedef}\gls{quotedblright}\end{inlinedef}
or \begin{inlinedef}\gls{textquotedblright}\end{inlinedef} & \textquotedblright\\
\begin{inlinedef}\gls{questiondown}\end{inlinedef}
or \begin{inlinedef}\gls{textquestiondown}\end{inlinedef} & ?` &
\begin{inlinedef}\gls{j}\end{inlinedef} & \dotlessj & 
\begin{inlinedef}\gls{quotedblleft}\end{inlinedef} 
or \begin{inlinedef}\gls{textquotedblleft}\end{inlinedef} & \textquotedblleft\\
\begin{inlinedef}\gls{exclamdown}\end{inlinedef}
or \begin{inlinedef}\gls{textexclamdown}\end{inlinedef} & !` &
\begin{inlinedef}\gls{dash}\end{inlinedef} & - &
\begin{inlinedef}\gls{endash}\end{inlinedef}
or \begin{inlinedef}\gls{textendash}\end{inlinedef} & \textendash \\
\begin{inlinedef}\gls{emdash}\end{inlinedef}
or \begin{inlinedef}\gls{textemdash}\end{inlinedef} & \textemdash &
\begin{inlinedef}\gls{S}\end{inlinedef} & \S &
\begin{inlinedef}\gls{textperiodcentered}\end{inlinedef} &
\textperiodcentered\\
\begin{inlinedef}\gls{ldots}\end{inlinedef} & \ldots  &
\begin{inlinedef}\gls{P}\end{inlinedef} & \P &
\begin{inlinedef}\gls{underscore}\end{inlinedef} 
or \begin{inlinedef}\gls{textunderscore}\end{inlinedef}& \_ 
\end{tabular}
\end{latexonly}
\begin{htmlonly}
\begin{tabular}{ll|ll}
\begin{inlinedef}\gls{textbackslash}\end{inlinedef} &
\textbackslash&
\begin{inlinedef}\gls{underscore}\end{inlinedef}\space 
or \begin{inlinedef}\gls{textunderscore}\end{inlinedef}& \_ \\
\begin{inlinedef}\gls{dash}\end{inlinedef} & - &
\begin{inlinedef}\gls{endash}\end{inlinedef}\space
or \begin{inlinedef}\gls{textendash}\end{inlinedef} & \makeimg{en dash character}{\textendash} \\
\begin{inlinedef}\gls{emdash}\end{inlinedef}\space  
or \begin{inlinedef}\gls{textemdash}\end{inlinedef} & \makeimg{em dash character}{\textemdash} &
\begin{inlinedef}\gls{P}\end{inlinedef} & \P \\
\begin{inlinedef}\gls{textasciicircum}\end{inlinedef} & \textasciicircum & 
\begin{inlinedef}\gls{dollar}\end{inlinedef} & \$ \\
\begin{inlinedef}\gls{S}\end{inlinedef} & \S &
\begin{inlinedef}\gls{textasciitilde}\end{inlinedef} & \textasciitilde\\
\begin{inlinedef}\gls{leftbrace}\end{inlinedef} & \{ &
\begin{inlinedef}\gls{ldots}\end{inlinedef} & \ldots \\
\begin{inlinedef}\gls{pounds}\end{inlinedef} & \pounds &
\begin{inlinedef}\gls{rightbrace}\end{inlinedef} & \} \\
\begin{inlinedef}\gls{dag}\end{inlinedef} & 
\makeimg{dagger symbol}{\dag} &
\begin{inlinedef}\gls{questiondown}\end{inlinedef}\space  
or \begin{inlinedef}\gls{textquestiondown}\end{inlinedef} & 
\makeimg{upside-down question mark}{?`} \\
\begin{inlinedef}\gls{textregistered}\end{inlinedef}& 
\makeimg{registered symbol}{\textregistered}&
\begin{inlinedef}\gls{hash}\end{inlinedef} & \# \\
\begin{inlinedef}\gls{ddag}\end{inlinedef} & 
\makeimg{double dagger symbol}{\ddag} &
\begin{inlinedef}\gls{exclamdown}\end{inlinedef}\space  
or \begin{inlinedef}\gls{textexclamdown}\end{inlinedef} & 
\makeimg{upside-down exclamation mark}{!`} \\
\begin{inlinedef}\gls{texttrademark}\end{inlinedef} & \texttrademark &
\begin{inlinedef}\gls{percent}\end{inlinedef} & \% \\
\begin{inlinedef}\gls{quoteright}\end{inlinedef}\space  
or \begin{inlinedef}\gls{textquoteright}\end{inlinedef} &
\makeimg{closing single quote}{\textquoteright}&
\begin{inlinedef}\gls{quotedblright}\end{inlinedef}\space  
or \begin{inlinedef}\gls{textquotedblright}\end{inlinedef} &
\makeimg{closing double quote}{\textquotedblright}\\
\begin{inlinedef}\gls{copyright}\end{inlinedef} & \copyright &
\begin{inlinedef}\gls{amp}\end{inlinedef} & \& \\
\begin{inlinedef}\gls{quoteleft}\end{inlinedef}\space  
or \begin{inlinedef}\gls{textquoteleft}\end{inlinedef} &
\makeimg{opening single quote}{\textquoteleft} &
\begin{inlinedef}\gls{quotedblleft}\end{inlinedef}\space  
or \begin{inlinedef}\gls{textquotedblleft}\end{inlinedef} &
\makeimg{opening double quote}{\textquotedblleft}\\
\begin{inlinedef}\gls{textbullet}\end{inlinedef} & \makeimg{bullet point}{\textbullet} & 
\begin{inlinedef}\gls{i}\end{inlinedef} & \makeimg{dotless i}{\i} \\ 
\begin{inlinedef}\gls{j}\end{inlinedef} & \makeimg{dotless j}{\dotlessj} & 
\begin{inlinedef}\gls{textbar}\end{inlinedef} & \textbar\\
\begin{inlinedef}\gls{textperiodcentered}\end{inlinedef} &
\middot &
\begin{inlinedef}\gls{textless}\end{inlinedef} & \textless\\
\begin{inlinedef}\gls{textgreater}\end{inlinedef} & \textgreater&
\begin{inlinedef}\gls{slash}\end{inlinedef} & / \\
\end{tabular}
\end{htmlonly}

\end{table}

Ligatures and special symbols are shown in \tableref{tab:ligatures}.
(Note that, as mentioned in the \htmlref{introduction}{ch:intro}, the
f-ligatures are automatically converted.)  When using a command in
the middle of a word, take care that the command doesn't run into the
rest of the word.  For example, the British spelling of the word
man\oe{}uvre has an
\latexhtml{\oe}{oe}-ligature in the middle of it.
You will get an error if you try:
\begin{alltt}\wrong
man\cmdname{oe}uvre
\end{alltt}
as \LaTeX\ will interpret it as the command
\verb|\oeuvre| which doesn't exist.

There are several ways to code this in \LaTeX:
\begin{enumerate}

\item Place a space after the command:
\begin{codeS}
man\glsni{oe}\glsni{visiblespace}uvre
\end{codeS}

\item Place an empty brace after the command:
\begin{codeS}
man\glsni{oe}\marg{}uvre
\end{codeS}

\item Group the command:
\begin{codeS}
man\marg{\glsni{oe}}uvre
\end{codeS}
(This can adversely affect the kerning so is best avoided.)

\end{enumerate}

\begin{table}[htbp]
\caption[Ligatures and Special Symbols]{Ligatures and Special Symbols (Computer Modern Font)}
\label{tab:ligatures}
\centering
\begin{tabular}{ll@{\hspace{3\tabcolsep}}ll@{\hspace{3\tabcolsep}}ll@{\hspace{3\tabcolsep}}ll}
\gls{AE} & 
  \makeimg{upper case AE ligature}{\fontfamily{cmr}\selectfont\AE} &
\gls{ae} & 
  \makeimg{lower case ae ligature}{\fontfamily{cmr}\selectfont\ae} &
\gls{OE} & 
  \makeimg{upper case OE ligature}{\fontfamily{cmr}\selectfont\OE} &
\gls{oe} & 
  \makeimg{lower case oe ligature}{\fontfamily{cmr}\selectfont\oe}\\
\Indextt{fi} & 
  \makeimg{fi ligature}{\fontfamily{cmr}\selectfont fi} &
\Indextt{ffi} & 
  \makeimg{ffi ligature}{\fontfamily{cmr}\selectfont ffi} &
\Indextt{fl} & 
  \makeimg{fl ligature}{\fontfamily{cmr}\selectfont fl} &
\Indextt{ffl} & 
  \makeimg{ffl ligature}{\fontfamily{cmr}\selectfont ffl}\\
\gls{AA} & \makeimg{upper case A ring}{\fontfamily{cmr}\selectfont\AA} &
\gls{aa} & \makeimg{lower case a ring}{\fontfamily{cmr}\selectfont\aa} &
\gls{L} & \makeimg{upper case L with a stroke}{\fontfamily{cmr}\selectfont\L} &
\gls{l} & \makeimg{lower case l with a stroke}{\fontfamily{cmr}\selectfont\l}\\
\gls{O} & \makeimg{upper case slashed O}{\fontfamily{cmr}\selectfont\O}&
\gls{o} & \makeimg{lower case slashed o}{\fontfamily{cmr}\selectfont\o} &
\gls{SS} & \makeimg{upper case Eszett}{\fontfamily{cmr}\selectfont\SS} &
\gls{ss} & \makeimg{lower case eszett}{\fontfamily{cmr}\selectfont\ss}
\end{tabular}
\end{table}

\refstepcounter{object}\label{obj:accents}%
English speakers are by and large very lackadaisical when it comes
to accents, but accents affect pronunciation, and so are just as
important as the correct spelling. There is a big difference
between putting your knife into someone's p\^at\'e (meat paste), and
putting your knife into someone's pate (head)!

Accented letters are created by specifying which accent you want, and
the letter on which to put the accent.  The accent commands are
listed in \xtableref{tab:accents}, and each command takes one
\gls{mandatory}.  The command indicates
what accent to use, and the argument indicates the letter on which to put
the accent.

You may have noticed in \tableref{tab:symbols} the commands \gls{i}
and \gls{j} which produce a dotless i and j (\i\ and \dotlessj).
With old versions of \LaTeX\ (or \TeX) an accent over a normal
\dq{i} or \dq{j} left the original dot in, which is incorrect, so a
dotless \dq{\i} or \dq{\dotlessj} were required. With modern
distributions, an accented \dq{i} or \dq{j} is correctly rendered.

\bookpagebreak
\xminisec{Example:}
\begin{code}
It's na\glsni{doublequote}ive to think that eating mouldy 
p\glsni{circum}at\glsni{acute}e won't result in food poisoning.
\end{code}%
Result:
\begin{resultS}[naive.html]
It's na\"{\i}ve to think that eating mouldy p\^at\'e won't result in food poisoning.
\end{resultS}%

\begin{table}[htbp]
\caption{Accent Commands}
\label{tab:accents}
\centering
\begin{tabular}{llc@{\hspace{3\tabcolsep}}llc}
\toprule
  & \multicolumn{2}{c@{\hspace{3\tabcolsep}}}{Example} &
  & \multicolumn{2}{c}{Example} \\
 Definition & Input & Output & Definition & Input & Output\\\midrule
\begin{inlinedef}\gls{acute}\marg{\meta{object}}\end{inlinedef} &
\verb|\'{c}| & \makeimg{c with an acute accent}{\'{c}} &
\begin{inlinedef}\gls{macron}\marg{\meta{object}}\end{inlinedef} &
\verb|\={c}| & \makeimg{c with a macron accent}{\={c}} \\
\begin{inlinedef}\gls{grave}\marg{\meta{object}}\end{inlinedef} &
\verb|\`{c}| & \makeimg{c with a grave accent}{\`{c}} &
\begin{inlinedef}\gls{period}\marg{\meta{object}}\end{inlinedef} &
\verb|\.{c}| & \makeimg{c with a dot over it}{\.{c}} \\
\begin{inlinedef}\gls{circum}\marg{\meta{object}}\end{inlinedef} &
\verb|\^{c}| & \makeimg{c with a circumflex}{\^{c}} &
\begin{inlinedef}\gls{tilde}\marg{\meta{object}}\end{inlinedef} &
\verb|\~{c}| & \makeimg{c with a tilde above}{\~{c}} \\
\begin{inlinedef}\gls{doublequote}\marg{\meta{object}}\end{inlinedef} &
\verb|\"{c}| & \makeimg{c with an umlaut}{\"{c}} &
\begin{inlinedef}\gls{v}\marg{\meta{object}}\end{inlinedef} &
\verb|\v{c}| & \makeimg{c with a caron accent}{\v{c}} \\
\begin{inlinedef}\gls{u}\marg{\meta{object}}\end{inlinedef} &
\verb|\u{c}| & \makeimg{c with a breve accent over it}{\u{c}} &
\begin{inlinedef}\gls{H}\marg{\meta{object}}\end{inlinedef} &
\verb|\H{c}| & \makeimg{c with double acute accent}{\H{c}} \\
\begin{inlinedef}\gls{t}\marg{\meta{object}}\end{inlinedef} &
\verb|\t{xy}| & \makeimg{x and y with a tie over them}{\t{xy}} &
\begin{inlinedef}\gls{c}\marg{\meta{object}}\end{inlinedef} &
\verb|\c{c}| & \makeimg{c with a cedilla}{\c{c}}\\
\begin{inlinedef}\gls{d}\marg{\meta{object}}\end{inlinedef} &
\verb|\d{c}| & \makeimg{c with a dot under it}{\d{c}} &
\begin{inlinedef}\gls{b}\marg{\meta{object}}\end{inlinedef} &
\verb|\b{c}| & \makeimg{c with a bar under it}{\b{c}}\\
\begin{inlinedef}\gls{r}\marg{\meta{object}}\end{inlinedef} &
\verb|\r{c}| & \makeimg{c with a ring above}{\r{c}} & & &\\
\bottomrule
\end{tabular}
\end{table}

This book only covers a very small subset of available symbol
commands. If the command you want isn't here, try Scott~Pakin's
comprehensive symbol list~\cite{pakin09}. Another useful resource is
\iappnamelink[.]{detexify}{http://detexify.kirelabs.org/classify.html}

\setnode{inputenc}
\subsection{The \sty{inputenc} Package}
\label{sec:inputenc}

Instead of using the accent and ligature commands described
\htmlref{above}{sec:chars}, you can use the \isty{inputenc} package
and enter the character directly, but you must ensure you match the
encoding with that used by your text editor. For example, this book
uses UTF8 encoding so I~have loaded the \sty{inputenc} package in
the \glsterm{preamble} with the \istyopt{inputenc}{utf8} option:
\begin{codeS}
\glsni{usepackage}\oarg{utf8}\marg{inputenc}
\end{codeS}
Note that it's a good idea to also use the \isty{fontenc} package as
well. For example, if you want to use Type~1 fonts:
\begin{code}
\begin{alltt}
\glsni{usepackage}\oarg{T1}\marg{fontenc}
\glsni{usepackage}\oarg{utf8}\marg{inputenc}
\end{alltt}
\end{code}

Returning to an earlier example, 
I~can directly enter the Unicode character (U+0153) for the
lower case \oe\ ligature:
\bookpagebreak
% actually I can't do it because latex2html won't like it, but lets
% pretend I have
\begin{codeS}
man\oe uvre
\end{codeS}
Note that if you are collaborating on a document and you want to
use this approach, you must ensure that all your co-authors use the
same input encoding. For example, suppose you decide to use
ISO~Latin~1 encoding (\istyopt{inputenc}{latin1} option):
\begin{codeS}
\glsni{usepackage}\oarg{latin1}\marg{inputenc}
\end{codeS}
but your co-author is using a UTF-8 editor and types:
\begin{codeS}
na\"ive
\end{codeS}
where \texttt{\"i} is the Unicode character U+00EF\@. UTF-8 characters
use one to four 8-bit bytes whereas ISO Latin~1 uses an 8-bit
single-byte character set. So the U+00EF binary sequence is
interpreted by ISO Latin~1 encoding as two characters: \~A\ (0xC3)
and \={ } (0xAF)\@.  Therefore the resulting PDF file will end up
containing the rather odd looking:
\begin{flushleft}\wrong
na\~A\={ }ve
\end{flushleft}\html{\par}%
(If you are using TeXWorks, you can set your preferred encoding
using \menu{Edit}\menuto\menu{Preferences} and select the \dq{Editor} tab 
where there is an \dq{Encoding} setting. Make sure this setting
matches the \isty{inputenc} option you use in your document.)

\begin{exercise}{Using Special Characters}{ex:spchar}
Start a new file in TeXworks, and
see if you can write the source code to create the output below.
\doifnotbook{(Ignore any \gls{hyphenation} that may appear below, \LaTeX\ does
that automatically where necessary, see \sectionref{sec:hyphenation}. Likewise, 
ignore where the line breaks occur, except for the paragraph break.) }Choose 
whether you want to use the \isty{inputenc} package or if you
want to use commands such as \gls{c}, but in either case you need to
be careful of the \htmlref{special characters}{tab:specialchars}\doifbook{\ 
listed in \tableref{tab:specialchars}}.

\begin{result}[spcharex.html]\setlength{\parindent}{1.5em}%
Item \#1: Our travel expenditure came to \$2000.00 
\& our equipment expenditure came
to \pounds 100.00 plus VAT @ 17.5\%.

Chlo\"e collected Zo\"e from the cr\`eche. They stopped to
admire the fa\c{c}ade of a new caf\'e and then went to a matin\'ee.
\end{result}

You can \downloadorview{spchar} the source code if you can't work out
how to do it, and remember to check \appendixref{ch:errors} if you have a problem.

\end{exercise}

%%%%%%%%%%%%%%%%%%%%%%%%% Lists %%%%%%%%%%%%%%%%%%%%%%%%%%

\setnode{lists}
\section{Lists}
\label{sec:lists}

Now you've had a go at using some \glspl{command},
let's use some \glspl{environment}\latex{ (recall
\sectionref{sec:environment})}. A good example of environments are the list
making environments. There are three basic list making environments:
\glsni{env-itemize} (for unordered lists),
\glsni{env-enumerate} (for ordered lists) and
\glsni{env-description} (for lists where you want to specify
your own label.)

In each of these environments, each item in the list must be started
with the command: 
\begin{definition}
\gls{item}\oarg{\meta{marker}}
\end{definition}
The optional argument \meta{marker} can be used to override 
the default marker for that particular item. (For example, to replace
the bullet point for an individual item in an \htmlref{unordered
list}{sec:itemize} to make that item stand out from all the other
items.) We will be looking at how to change the default marker
in \sectionref{sec:renewcom}.

\xminisec{Related \gls{ukfaq} topics:}
\begin{itemize}
\item \faqlink{Perhaps a missing \cmdname{item}?}{errmissitem}
\item \faqlink{Fancy enumeration lists}{enumerate}
\item \faqlink{How to adjust list spacing}{complist}
\item \faqlink{Interrupting enumerated lists}{interruptlist}
\item \faqlink{\dq{Too deeply nested}}{toodeep}
\end{itemize}

%---------------------- Itemize --------------------------

\setnode{itemize}
\subsection{Unordered Lists}
\label{sec:itemize}

Unordered lists are created using the \gls{env-itemize} environment.

\xminisec{Example:}
\begin{code}
\glsni{begin}\marg{itemize}\newline
\glsni{item} Animal\newline
\glsni{item} Vegetable\newline
\glsni{item} Mineral\newline
\glsni{end}\marg{itemize}
\end{code}
\begin{result}[itemize.html]
\begin{itemize}
\item Animal
\item Vegetable
\item Mineral
\end{itemize}
\end{result}

\xminisec{Another Example:}

Changing the default markers is covered in \sectionref{sec:renewcom}, 
but it's also possible to override the default marker for a
particular item, as in this example (recall the double-dagger
symbol command \gls{ddag} from \tableref{tab:symbols}):
\begin{code}
\glsni{begin}\marg{itemize}\newline
\glsni{item} Animal\newline
\glsni{item}\oarg{\glsni{ddag}} Vegetable\newline
\glsni{item} Mineral\newline
\glsni{end}\marg{itemize}
\end{code}%
\begin{result}[itemizedag.html]
\begin{itemize}
\item Animal
\item[\ddag] Vegetable
\item Mineral
\end{itemize}
\end{result}

Be careful about using square brackets
\gls{text.opensq}\gls{text.closesq} inside an optional
argument. Grouping is required, as in:
\begin{code}
\glsni{begin}\marg{itemize}\newline
\glsni{item} Animal\newline
\glsni{item}\oarg{\marg{\gls{text.opensq}X\gls{text.closesq}}} Vegetable\newline
\glsni{item} Mineral\newline
\glsni{end}\marg{itemize}
\end{code}%
\begin{result}[itemizesq.html]
\begin{itemize}
\item Animal
\item[{[X]}] Vegetable
\item Mineral
\end{itemize}
\end{result}
Similarly if the item starts with an open square bracket
\glsni{opensq}, as in:
\begin{code}
\glsni{begin}\marg{itemize}\newline
\glsni{item} Animal\newline
\glsni{item} \marg{\gls{text.opensq}sic\gls{text.closesq}} Wegetable\newline
\glsni{item} Mineral\newline
\glsni{end}\marg{itemize}
\end{code}
\begin{result}[itemizetextsqbr.html]
\begin{itemize}
\item Animal
\item {[sic]} Wegetable
\item Mineral
\end{itemize}
\end{result}

\setnode{nesteditemize}
\subsubsection{Nested Lists}
\label{sec:nesteditemise}

It is also possible to nest \glsni{env-itemize} environments.
The following example has three levels, each using its own default marker.

\begin{code}
\glsni{begin}\marg{itemize}\newline
\glsni{item} Animal\newline
\glsni{begin}\marg{itemize}\newline
\glsni{item} Mammals\newline
\glsni{item} Birds\newline
\glsni{item} Reptiles. For example:\newline
\glsni{begin}\marg{itemize}\newline
\glsni{item} dinosaurs\newline
\glsni{item} crocodiles\newline
\glsni{end}\marg{itemize}\newline
\glsni{end}\marg{itemize}\newline
\glsni{item} Vegetable\newline
\glsni{begin}\marg{itemize}\newline
\glsni{item} Cultivated\newline
\glsni{item} Wild\newline
\glsni{end}\marg{itemize}\newline
\glsni{item} Mineral\newline
\glsni{end}\marg{itemize}
\end{code}
\begin{result}[nesteditemize.html]
\begin{itemize}
\item Animal
\begin{itemize}
\item Mammals
\item Birds
\item Reptiles. For example:
\begin{itemize}
\item dinosaurs
\item crocodiles
\end{itemize}
\end{itemize}
\item Vegetable
\begin{itemize}
\item Cultivated
\item Wild
\end{itemize}
\item Mineral
\end{itemize}
\end{result}

You might have noticed the code in the above example is a little
difficult to read. Each new list item starts a new paragraph, so it
doesn't matter if we have blank lines before each item. Also, recall
from \chapterref{ch:def} that spaces at the start of each line of
code are ignored, so it's possible to make the code more readable
without affecting the final result:

\begin{code}
\begin{alltt}
\glsni{begin}\marg{itemize}

   \glsni{item} Animal

   \glsni{begin}\marg{itemize}

      \glsni{item} Mammals

      \glsni{item} Birds

      \glsni{item} Reptiles. For example:

      \glsni{begin}\marg{itemize}

         \glsni{item} dinosaurs

         \glsni{item} crocodiles

       \glsni{end}\marg{itemize}

   \glsni{end}\marg{itemize}

   \glsni{item} Vegetable

   \glsni{begin}\marg{itemize}

      \glsni{item} Cultivated

      \glsni{item} Wild

   \glsni{end}\marg{itemize}

   \glsni{item} Mineral

\glsni{end}\marg{itemize}
\end{alltt}
\end{code}

It's now a little easier to see which \verb|\begin{itemize}| matches
up with the corresponding \verb|\end{itemize}|.

\xminisec{Example (Four Levels)}

This example has four levels, which is the maximum allowed by most
classes.

\begin{code}
\begin{alltt}
\glsni{begin}\marg{itemize}
 \glsni{item} Animal

 \glsni{begin}\marg{itemize}
   \glsni{item} Mammal

   \glsni{begin}\marg{itemize}
    \glsni{item} Placental

    \glsni{item} Monotreme

    \glsni{begin}\marg{itemize}
      \glsni{item} Platypus
    \glsni{end}\marg{itemize}

    \glsni{item} Marsupial

    \glsni{begin}\marg{itemize}
     \glsni{item} Kangaroo

     \glsni{item} Koala
    \glsni{end}\marg{itemize}
   \glsni{end}\marg{itemize}

   \glsni{item} Reptile
 \glsni{end}\marg{itemize}

 \glsni{item} Vegetable

 \glsni{item} Mineral
\glsni{end}\marg{itemize}
\end{alltt}
\end{code}
\begin{result}[fourlevelitemize.html]
\begin{itemize}
 \item Animal
 \begin{itemize}
   \item Mammal

   \begin{itemize}
    \item Placental

    \item Monotreme
    \begin{itemize}
      \item Platypus
    \end{itemize}


    \item Marsupial

    \begin{itemize}
     \item Kangaroo

     \item Koala
    \end{itemize}
   \end{itemize}

   \item Reptile
 \end{itemize}

 \item Vegetable

 \item Mineral
\end{itemize}
\end{result}

If you try adding a further level, \LaTeX\ will give a \dq{Too
deeply nested} error.

%---------------- Enumerate -------------------------

\setnode{enumerate}
\subsection{Ordered Lists}
\label{sec:enumerate}

Ordered lists are created using the \gls{env-enumerate} 
environment.
It has exactly the same format as the \gls{env-itemize} environment
described \htmlref{above}{sec:itemize}.

We can use the same example as before, only this time use
\glsni{env-enumerate} instead of \gls{env-itemize}.

\begin{code}
\begin{alltt}
\glsni{begin}\marg{enumerate}
  \glsni{item} Animal
  \glsni{item} Vegetable
  \glsni{item} Mineral
\glsni{end}\marg{enumerate}
\end{alltt}
\end{code}%
The above input will produce the following output:
\begin{result}[enumerate.html]
\begin{enumerate}
\item Animal
\item Vegetable
\item Mineral
\end{enumerate}
\end{result}

As before, the marker for a particular item can be overridden:
\begin{code}
\glsni{begin}\marg{enumerate}\newline
\strut~~\glsni{item} Animal\newline
\strut~~\glsni{item}\oarg{\marg{\gls{text.opensq}X\gls{text.closesq}}} Vegetable\newline
\strut~~\glsni{item} Mineral\newline
\glsni{end}\marg{enumerate}
\end{code}%
\begin{result}[enumeratesqbr.html]
\begin{enumerate}
\item Animal
\item[{[X]}] Vegetable
\item Mineral
\end{enumerate}
\end{result}

\xminisec{Example (Nested):}

As with the \gls{env-itemize} environment, most classes allow a
maximum of four nested \gls{env-enumerate} environments.

\begin{code}
\begin{alltt}
\glsni{begin}\marg{enumerate}
 \glsni{item} Animal

 \glsni{begin}\marg{enumerate}
   \glsni{item} Mammal

   \glsni{begin}\marg{enumerate}
    \glsni{item} Placental

    \glsni{item} Monotreme

    \glsni{begin}\marg{enumerate}
      \glsni{item} Platypus
    \glsni{end}\marg{enumerate}

    \glsni{item} Marsupial

    \glsni{begin}\marg{enumerate}
     \glsni{item} Kangaroo

     \glsni{item} Koala
    \glsni{end}\marg{enumerate}
   \glsni{end}\marg{enumerate}

   \glsni{item} Reptile
 \glsni{end}\marg{enumerate}

 \glsni{item} Vegetable

 \glsni{item} Mineral
\glsni{end}\marg{enumerate}
\end{alltt}
\end{code}
\begin{result}[nestedenumerate.html]
\begin{enumerate}
 \item Animal
 \begin{enumerate}
   \item Mammal

   \begin{enumerate}
    \item Placental

    \item Monotreme
    \begin{enumerate}
      \item Platypus
    \end{enumerate}


    \item Marsupial

    \begin{enumerate}
     \item Kangaroo

     \item Koala
    \end{enumerate}
   \end{enumerate}

   \item Reptile
 \end{enumerate}

 \item Vegetable

 \item Mineral
\end{enumerate}
\end{result}

%--------------------- Description -------------------------

\setnode{description}
\subsection{Description Environment}
\label{sec:description}

The \gls{env-description} environment has the same format as
the \gls{env-itemize} environment described in
\sectionref{sec:itemize}, only this time you need to
specify a marker as an \gls{optional} to
the \gls{item} command, since there is no default marker
for this environment. The marker may be a textual label, and most
classes will typeset it in bold. The \koma\ classes, such as
\icls{scrartcl}, default to a bold sans-serif font, as illustrated
in this next example:

\begin{code}
\begin{alltt}
\glsni{begin}\marg{description}
  \glsni{item}\oarg{Animal} Living being

  \glsni{item}\oarg{Vegetable} Plant
 
  \glsni{item}\oarg{Mineral} Natural inorganic substance
\glsni{end}\marg{description}
\end{alltt}
\end{code}
\begin{result}[description.html]
\dodescriptionexample
\end{result}

The \koma\ classes provide a way of changing the font style in
the \glsni{env-description} label markers. (The font changing commands
\glsni{normalfont} and \glsni{scshape} will be covered in 
\sectionref{sec:fonts}, and the \koma\ command
\glsni{addtokomafont} in \sectionref{sec:section}.)

\begin{code}
\begin{alltt}
\glsni{addtokomafont}\marg{descriptionlabel}\marg{\glsni{normalfont}\glsni{scshape}}
\glsni{begin}\marg{description}
  \glsni{item}\oarg{Animal} Living being

  \glsni{item}\oarg{Vegetable} Plant

  \glsni{item}\oarg{Mineral} Natural inorganic substance
\glsni{end}\marg{description}
\end{alltt}
\end{code}
\begin{result}[descriptionkoma.html]
% Make sure LaTeX2HTML doesn't think this should be the same
% image as the previous example
\dodescriptionkomaexample
\end{result}

It is possible to nest all the listing environments, as long as you
don't exceed four \gls{env-itemize} and four \gls{env-enumerate}
environments. The \gls{env-description} environment has no
restriction on the number of times it can be nested. However, just
because you can do something, doesn't mean you should. In general
it's best to avoid an excessively complicated block of text in your
document.

\xminisec{Example (Assorted Nesting):}

This example uses each of the listing environments described above.

\bookpagebreak
\begin{code}
\begin{alltt}
\glsni{begin}\marg{description}

   \glsni{item}\oarg{Animal} Living being

   \glsni{begin}\marg{itemize}

      \glsni{item} Mammals

      \glsni{item} Birds

      \glsni{item} Reptiles. For example:

      \glsni{begin}\marg{enumerate}

         \glsni{item} dinosaurs

         \glsni{item} crocodiles

       \glsni{end}\marg{enumerate}

   \glsni{end}\marg{itemize}

   \glsni{item}\oarg{Vegetable} Plant

   \glsni{begin}\marg{itemize}

      \glsni{item} Cultivated. For example:

      \glsni{begin}\marg{enumerate}

         \glsni{item} Carrots

         \glsni{item} Broccoli

         \glsni{item} Potatoes

      \glsni{end}\marg{enumerate}

      \glsni{item} Wild

   \glsni{end}\marg{itemize}

   \glsni{item}\oarg{Mineral} Natural inorganic substance

\glsni{end}\marg{description}
\end{alltt}
\end{code}

\begin{result}[nestedlists.html]
\begin{description}

   \item[Animal] Living being

   \begin{itemize}

      \item Mammals

      \item Birds

      \item Reptiles. For example:
      \begin{enumerate}

         \item dinosaurs

         \item crocodiles

       \end{enumerate}

   \end{itemize}

   \item[Vegetable] Plant

   \begin{itemize}

      \item Cultivated. For example:

      \begin{enumerate}

         \item Carrots

         \item Broccoli

         \item Potatoes

      \end{enumerate}


      \item Wild

   \end{itemize}

   \item[Mineral] Natural inorganic substance

\end{description}
\end{result}

\begin{exercise}{Lists}{ex:lists}

Try writing the \gls{source} that will create
the output shown below.

\begin{result}[listex.html]
\begin{description}
\item[Village] A small collection of dwelling places.
Examples:
\begin{enumerate}
\item Marlingford
\item Saxlingham Nethergate
\end{enumerate}

\item[Town] A large collection of dwelling places.
Examples:
\begin{enumerate}
\item Great Yarmouth
\item Beccles
\end{enumerate}

\item[City] A large town, usually containing a cathedral.
Examples:
\begin{enumerate}
\item Norwich
\item Birmingham
\item London
\end{enumerate}
\end{description}
\end{result}\bookpagebreak
You can \downloadorview{lists} the answer if you can't work out how to do it.

\end{exercise}

%%%%%%%%%%%%%%%%%% Font Changing Commands %%%%%%%%%%%%%%%%%%%%%%

\setnode{fontcommands}
\section{Fonts}
\label{sec:fonts}

\LaTeX\ uses Donald~Knuth's Computer Modern fonts by
default\faq{Using PostScript fonts with TeX}{usepsfont}. This
supplies three font families: serif, sans-serif and a typewriter (or
monospaced) font (as well as the maths
fonts which are discussed in
\sectionref{sec:mathsfonts}). With each font family, you can
change the shape and weight, as well as the size.

\reportpagebreak
\setnode{fontstyle}
\subsection{Changing the Font Style}
\label{sec:fontstyle}

There are two basic ways of changing fonts: you can either change
the font for a small selection of text, for example, if you want to
\emph{emphasize} a word, or you may wish to change the font \dq{from
this point onwards}. The \glspl{command} shown in
\tableref{tab:fontsI} are of the first type (text-block
commands\index{command!text-block|hyperpage}), whereas those shown
in \tableref{tab:fontsII} are of the second type\dash a
\gls{declaration} (or modal command\index{command!modal|hyperpage}).

\xminisec{Note:}

\warning Don't be tempted to use \cmdname{bf}, \cmdname{md}, \cmdname{it}, 
\cmdname{sl}, \cmdname{sc}, \cmdname{sf}, \cmdname{tt} or \cmdname{rm}. 
These commands are \textbf{obsolete}~\cite{l2tabu}.\faq{What's wrong with \cmdname{bf}, \cmdname{it}
etc.?}{2letterfontcmd}

If you use an italic or slanted font declaration, such as
\glsni{itshape}, you will need to add an \keyword{italic correction}
\gls{itcorr} at the end of the block of text, when the last letter
of the sloping text leans too far over. This isn't necessary for
text-block commands, such as \glsni{textit}, just for the modal
commands. The effect is more noticeable when part of a word is
stressed, particularly with certain fonts.

\xminisec{Example:}

In the code below, the first instance of \dq{repeated} doesn't have an italic
correction but the second does: 
\begin{codeS}
\marg{\glsni{itshape} repeated}ly \marg{\glsni{itshape}
repeated\glsni{itcorr}}ly
\end{codeS}%
Using Computer Modern:
\begin{resultS}[repeatedly repeatedly]\fontfamily{cmr}\selectfont
{\itshape repeated}ly {\itshape repeated\/}ly
\end{resultS}
Using Helvetica:
\begin{resultS}[repeatedly repeatedly]\fontfamily{phv}\selectfont
{\itshape repeated}ly {\itshape repeated\/}ly
\end{resultS}
Using Antykwa Toru\'nska typeface:
\begin{resultS}[repeatedly repeatedly]
{\itshape repeated}ly {\itshape repeated\/}ly
\end{resultS}

\begin{table}[tp]
\caption{Font Changing Text-Block Commands}
\label{tab:fontsI}
\centering
\begin{tabular}{@{}lll@{}}
\toprule
\bfseries Command & \bfseries Example Input & \bfseries Corresponding output\\
& & (Computer Modern)\\
\midrule
\begin{inlinedef}\gls{textrm}\marg{\meta{text}}\end{inlinedef} &
\verb|\textrm{roman} text| & \makeimg{Roman text}{\fontfamily{cmr}\selectfont{\fontfamily{cmr}\selectfont roman} text}\\
\begin{inlinedef}\gls{textsf}\marg{\meta{text}}\end{inlinedef} &
\verb|\textsf{sans serif} text| & \makeimg{sans serif text}{\fontfamily{cmr}\selectfont{\fontfamily{cmss}\selectfont sans serif} text}\\
\begin{inlinedef}\gls{texttt}\marg{\meta{text}}\end{inlinedef} &
\verb|\texttt{typewriter} text| & \makeimg{typewriter text}{\fontfamily{cmr}\selectfont{\fontfamily{cmtt}\selectfont typewriter} text}\\[5pt]
\begin{inlinedef}\gls{textmd}\marg{\meta{text}}\end{inlinedef} &
\verb|\textmd{medium} text| & \makeimg{medium text}{\fontfamily{cmr}\selectfont\textmd{medium} text}\\
\begin{inlinedef}\gls{textbf}\marg{\meta{text}}\end{inlinedef} &
\verb|\textbf{bold} text| & \makeimg{bold text}{\fontfamily{cmr}\selectfont\textbf{bold} text}\\[5pt]
\begin{inlinedef}\gls{textup}\marg{\meta{text}}\end{inlinedef} &
\verb|\textup{upright} text| & \makeimg{upright text}{\fontfamily{cmr}\selectfont\textup{upright} text}\\
\begin{inlinedef}\gls{textit}\marg{\meta{text}}\end{inlinedef} &
\verb|\textit{italic} text| & \makeimg{italic text}{\fontfamily{cmr}\selectfont\textit{italic} text}\\
\begin{inlinedef}\gls{textsl}\marg{\meta{text}}\end{inlinedef} &
\verb|\textsl{slanted} text| & \makeimg{slanted text}{\fontfamily{cmr}\selectfont\textsl{slanted} text}\\[5pt]
\begin{inlinedef}\gls{textsc}\marg{\meta{text}}\end{inlinedef} &
\verb|\textsc{Small Caps} text| &
\makeimg{Small Caps text}{\fontfamily{cmr}\selectfont\textsc{Small Caps} text}\\
\begin{inlinedef}\gls{emph}\marg{\meta{text}}\end{inlinedef} &
\verb|\emph{emphasized} text| &
\makeimg{emphasized text}{\fontfamily{cmr}\selectfont\emph{emphasized} text}\\[5pt]
\begin{inlinedef}\gls{textnormal}\marg{\meta{text}}\end{inlinedef} &
\verb|\textnormal{default} text| &
\makeimg{default text}{\fontfamily{cmr}\selectfont default text}
\\\bottomrule
\end{tabular}
\end{table}

Note that if you want to typeset an URL, rather than using
\glsni{texttt} it is better to use\bookpagebreak
\begin{definition}
\gls{url}\marg{\meta{address}}
\end{definition}
which is defined in the \isty{url} package. For example:
\begin{codeS}%
\begin{alltt}
\glsni{url}\marg{http://theoval.cmp.uea.ac.uk/\glsnl{tildechar}nlct/}
\end{alltt}
\end{codeS}
produces:
\begin{resultS}[http://theoval.cmp.uea.ac.uk/\tildesym nlct/]
% simulate to prevent link, which might confuse the reader into
% thinking links automatically occur.
\texttt{http://theoval.cmp.uea.ac.uk/\tildesym nlct/}
\end{resultS}
(Note there is no need to do anything with the \gls{tildechar}
(tilde) \htmlref{special character}{sec:chars} if you use it in the argument of \glsni{url}.)

\begin{table}[hbtp]
\caption{Font Changing Declarations}
\label{tab:fontsII}
\centering
\begin{tabular}{@{}lll@{}}
\toprule
\bfseries Declaration & \bfseries Example Input & \bfseries Corresponding output\\
& & (Computer Modern)\\
\midrule
\begin{inlinedef}\gls{rmfamily}\end{inlinedef} &
\verb|\rmfamily roman text| &
\makeimg{Roman text}{\fontfamily{cmr}\selectfont roman text}\\
\begin{inlinedef}\gls{sffamily}\end{inlinedef} &
\verb|\sffamily sans serif text| &
\makeimg{sans serif text}{\fontfamily{cmss}\selectfont sans serif text}\\
\begin{inlinedef}\gls{ttfamily}\end{inlinedef} &
\verb|\ttfamily typewriter text| &
\makeimg{typewriter text}{\fontfamily{cmtt}\selectfont typewriter text}\\[5pt]
\begin{inlinedef}\gls{mdseries}\end{inlinedef} &
\verb|\mdseries medium text| &
\makeimg{medium text}{\fontfamily{cmr}\selectfont\mdseries medium text}\\
\begin{inlinedef}\gls{bfseries}\end{inlinedef} &
\verb|\bfseries bold text| &
\makeimg{bold text}{\fontfamily{cmr}\selectfont\bfseries bold text}\\[5pt]
\begin{inlinedef}\gls{upshape}\end{inlinedef} &
\verb|\upshape upright text| &
\makeimg{upright text}{\fontfamily{cmr}\selectfont\upshape upright text}\\
\begin{inlinedef}\gls{itshape}\end{inlinedef} &
\verb|\itshape italic text| &
\makeimg{italic text}{\fontfamily{cmr}\selectfont\itshape italic text}\\
\begin{inlinedef}\gls{slshape}\end{inlinedef} &
\verb|\slshape slanted text| &
\makeimg{slanted text}{\fontfamily{cmr}\selectfont\slshape slanted text}\\
\begin{inlinedef}\gls{scshape}\end{inlinedef} &
\verb|\scshape Small Caps text| &
\makeimg{Small Caps text}{\fontfamily{cmr}\selectfont\scshape Small Caps text}\\[5pt]
\begin{inlinedef}\gls{em}\end{inlinedef} &
\verb|\em emphasized text| & 
\makeimg{emphasized text}{\fontfamily{cmr}\selectfont\em emphasized text}\\[5pt]
\begin{inlinedef}\gls{normalfont}\end{inlinedef} &
\verb|\normalfont default text| &
\makeimg{default text}{\fontfamily{cmr}\selectfont default text}
\\\bottomrule
\end{tabular}

\end{table}

\refstepcounter{object}\label{obj:fontenv}\Glspl{environment} can be used instead.
Each \gls{environment} has the same name as its corresponding
declaration, but \emph{without} the preceding backslash. For example:
\begin{codeS}
\begin{alltt}
\glsni{begin}\marg{sffamily}Some sans-serif text.\glsni{end}\marg{sffamily}
\end{alltt}
\end{codeS}%
yields:
\begin{resultS}[Some sans-serif text.]
\begin{sffamily}Some sans-serif text.\end{sffamily}
\end{resultS}

You can combine a font family with a given shape and weight using
a variety of methods.

\xminisec{Examples:}

\begin{enumerate}
\item Localised declarations:

\begin{codeS}
\marg{\glsni{sffamily}\glsni{slshape} Some slanted sans-serif text.}
\end{codeS}

\item Declarations that later get explicitly reset:

\begin{codeS}
\glsni{sffamily}\glsni{slshape} Some slanted sans-serif
text.\glsni{normalfont}
\end{codeS}

\bookpagebreak
\item Mixing text-block and modal commands:

\begin{codeS}
\glsni{textsf}\marg{\glsni{slshape} Some slanted sans-serif text.}
\end{codeS}

\item Nested commands:

\begin{codeS}
\glsni{textsf}\marg{\glsni{textsl}\marg{Some slanted sans-serif text.}}
\end{codeS}

\item Mixing environments and declarations:

\begin{code}
\glsni{begin}\marg{sffamily}\glsni{slshape} Some slanted sans-serif 
text.\glsni{end}\marg{sffamily}
\end{code}
\end{enumerate}

All of the above produce the same output:
\begin{resultS}[Some slanted sans-serif text.]
\textsf{\slshape Some slanted sans-serif text.}
\end{resultS}

\faq{Warning: \dq{Font shape \ldots\ not
available}}{fontunavail}Note that some combinations are not
available, in which case \LaTeX\ will give a warning message, and
will substitute the font for what it considers to be the closest
available match.

\xminisec{Example:}
\begin{codeS}
\glsni{textsc}\marg{\glsni{bfseries} Text}
\end{codeS}
With the Antykwa Toru\'nska typeface, this appears as:
\begin{resultS}[Text (in bold small caps)]
\textsc{\bfseries Text}
\end{resultS}\bookpagebreak\noindent
whereas with Computer Modern, the result is:
\begin{resultS}[Text (in bold)]\fontfamily{cmr}\selectfont
\textsc{\bfseries Text}
\end{resultS}
This is because Computer Modern doesn't have a bold small-caps
font, so it just uses bold. \LaTeX\ gives the following warning:
\begin{verbatim}
LaTeX Font Warning: Font shape `T1/cmr/b/sc' undefined
(Font)              using `T1/cmr/b/n' instead on input line 2792.
\end{verbatim}
Most sans-serif fonts don't provide a small-caps variant, so 
\begin{codeS}
\glsni{textsf}\marg{\glsni{scshape} Text}
\end{codeS}
will either appear in regular sans-serif or small-caps serif,
depending on the font in use. Using Libris sans-serif the result is:
\begin{resultS}[Text (in sans-serif)]
\textsf{\scshape Text}
\end{resultS}
whereas using Computer Modern Sans, the result is:
\begin{resultS}[Text (in small caps)]
\fontfamily{cmss}
\scshape Text
\end{resultS}

\setnode{emph}
\subsubsection{Emphasizing Words or Phrases}
\label{sec:emph}

The \gls{command} \gls{emph},
the \gls{declaration} \gls{em} and the
\gls{environment} \gls{env-em} behave slightly
differently to the corresponding \gls{textit} command,
\gls{itshape} declaration and \gls{env-itshape} environment.
The latter group simply use an italic font, whereas the former will toggle
between sloping and upright.  So if the surrounding font is upright
then \gls{emph}, \gls{em} and \glsni{env-em} will use the
sloping font, but if the surrounding font is italic or slanted,
\gls{emph}, \gls{em} and \glsni{env-em} will use an upright
font.  This is particularly useful in abstracts where the
\Index{abstract} font varies between \glspl{cls}.  It is recommended that if your intention
is to emphasize something, you should use \gls{emph} etc.\ rather
than \gls{textit} etc.\screenpagebreak

\xminisec{Examples:}
\begin{enumerate}
\item Emphasized text in upright surrounding:
\begin{codeS}
Some \glsni{emph}\marg{emphasized} text.
\end{codeS}%
yields
\begin{resultS}[Some emphasized text.]
Some \emph{emphasized} text.
\end{resultS}

\item Emphasized text in italic surrounding:
\begin{codeS}
\marg{\glsni{itshape} Some \glsni{emph}\marg{emphasized} text.}
\end{codeS}%
yields
\begin{resultS}[Some emphasized text]
{\itshape Some \emph{emphasized} text.}
\end{resultS}

\bookpagebreak
\item Emphasized text in upright sans-serif surrounding:
\begin{codeS}
\marg{\glsni{sffamily} Some \glsni{emph}\marg{emphasized} text.}
\end{codeS}%
yields
\begin{resultS}[Some emphasized text]
{\sffamily Some \emph{emphasized} text.}
\end{resultS}
\end{enumerate}

\setnode{fontsize}
\subsection{Changing the Font Size}
\label{sec:fontsize}

When you start writing a document, you need to decide what the base
font size should be. The \koma\ classes default to 11pt, but this can
be changed using the class options \scrclsopt{8pt}, \scrclsopt{9pt},
\clsopt{10pt}, \clsopt{12pt}, \scrclsopt{14pt}, \scrclsopt{17pt} or
\scrclsopt{20pt}.

You can then change the font size \emph{relative to} the base size,
using one of the declarations shown in \tableref{tab:fontsize}. That
way, if you later decide to change the normal font size from, say,
11pt to 12pt, all you need do is change the class option (see
\latexhtml{page~\pageref{obj:sizeopt}}{\htmlref{earlier}{obj:sizeopt}})
and re-run \LaTeX. Note that there are no equivalent text-block
commands.

\begin{table}[tbhp]
\caption{Font Size Changing Declarations}
\label{tab:fontsize}
\centering
\begin{tabular}{@{}lll@{}}
\toprule
\bfseries Declaration & \bfseries Example Input & 
\bfseries Corresponding Output\\
\midrule
\begin{inlinedef}\gls{tiny}\end{inlinedef} &
\verb|\tiny tiny text| & \makeimg{tiny text}{\tiny tiny text}\\
\begin{inlinedef}\gls{scriptsize}\end{inlinedef} &
\verb|\scriptsize script size| &
\makeimg{script size}{\scriptsize script size}\\
\begin{inlinedef}\gls{footnotesize}\end{inlinedef} &
\verb|\footnotesize footnote size| &
\makeimg{footnote size}{\footnotesize footnote size}\\
\begin{inlinedef}\gls{small}\end{inlinedef} &
\verb|\small small text| &
\makeimg{small text}{\small small text}\\
\begin{inlinedef}\gls{normalsize}\end{inlinedef} &
\verb|\normalsize normal size| &
\makeimg{normal size}{\normalsize normal size}\\
\begin{inlinedef}\gls{large}\end{inlinedef} &
\verb|\large large text| &
\makeimg{large text}{\large large text}\\
\begin{inlinedef}\gls{Large}\end{inlinedef} &
\verb|\Large even larger| &
\makeimg{even larger}{\Large even larger}\\
\begin{inlinedef}\gls{LARGE}\end{inlinedef} &
\verb|\LARGE larger still| &
\makeimg{larger still}{\LARGE larger still}\\
\begin{inlinedef}\gls{huge}\end{inlinedef} &
\verb|\huge huge| & \makeimg{huge}{\huge huge}\\
\begin{inlinedef}\gls{Huge}\end{inlinedef} &
\verb|\Huge extra huge| &
\makeimg{extra huge}{\Huge extra huge}
\\\bottomrule
\end{tabular}

\end{table}

Again, \glspl{environment} can be used instead, where
each \gls{environment} has the same name as its corresponding
declaration, but \emph{without} the preceding backslash.
Font environments may be nested, for example:%
\indexEnv{itshape}\indexEnv{Large}
\begin{code}
\glsni{begin}\marg{itshape} Some italic text.
\glsni{begin}\marg{Large}This text is large.\glsni{end}\marg{Large}
\glsni{end}\marg{itshape} Back to normal.
\end{code}%
Output:
\begin{resultS}[Some italic text. This text is large. Back to
normal.]
\begin{itshape} Some italic text.
\begin{Large}This text is large.\end{Large}
\end{itshape}
Back to normal.
\end{resultS}

\setnode{documentfonts}
\subsection{Changing Document Fonts}
\label{sec:changingfonts}

\faq{Choice of scalable outline fonts}{psfchoice}What if you don't
want to use the default Computer Modern fonts? Some publishers and
institutions insist on a combination of Times Roman (serif),
Helvetica (sans-serif) and Courier (typewriter). To do this, you can
load the following packages:
\begin{description}
\item[\isty{mathptmx}] (Times) Only affects \gls{rmfamily} and
\gls{textrm}.

\item[\isty{helvet}] (Helvetica) Only affects \gls{sffamily} and
\gls{textsf}.

\item[\isty{courier}] (Courier) Only affects \gls{ttfamily}
and \gls{texttt}.
\end{description}

\xminisec{Notes:}

\begin{enumerate}

\item \warning Don't be tempted to use the \sty{times} package. It's
obsolete~\cite{l2tabu}. Use \isty{mathptmx} instead.

\item Although Times and Helvetica are commonly used together, they
don't match, as illustrated below (temporarily switching from this
book's fonts to Times and Helvetica):
\begin{codeS}
\glsni{rmfamily} xx \glsni{sffamily} xx
\end{codeS}
Results in:
\begin{resultS}[xx xx]
\fontfamily{ptm}\selectfont xx \fontfamily{phv}\selectfont xx
\end{resultS}
The first two x's are in Times Roman and the second two are in
Helvetica, which are somewhat larger. To compensate for this you
need to scale the Helvetica font using the \istyopt{helvet}{scaled}
option:
\begin{alltt}
\glsni{usepackage}\oarg{scaled=0.9}\marg{helvet}
\end{alltt}

\item Loading \isty{helvet} or \isty{courier} doesn't change the
default font family. Consider the following:
\begin{alltt}
\glsni{documentclass}\marg{scrartcl}

\glsni{usepackage}\marg{helvet}

\glsni{begin}\marg{document}
This is a sample document.
\glsni{end}\marg{document}
\end{alltt}
Here, the text \dq{This is a sample document} will be typeset in
Computer Modern Roman. This is because \gls{rmfamily} is the default
font and \isty{helvet} only affects \gls{sffamily}, which hasn't
been used. (See \sectionref{sec:renewcom} to find out how to change
the default font family.)
\end{enumerate}
\latexhtml{This book}{The PDF version of this document} has used the following packages:
\begin{alltt}
\glsni{usepackage}\oarg{T1}\marg{fontenc}
\glsni{usepackage}\oarg{math}\marg{anttor}
\glsni{usepackage}\marg{libris}
\end{alltt}
The \isty{fontenc} package is used to switch to Type~1 font
encoding\faq{Why bother with \sty{inputenc} and \sty{fontenc}?}{why-inp-font},
the \isty{anttor} package is used to set the serif family to Antykwa
Toru\'nska typeface, and the \isty{libris} package is used to set
the sans-serif family to the Libris ADF typeface.

\begin{exercise}{Fonts}{ex:fonts}

Go back to the document you created in \exerciseref{ex:simpledoc} and
change the first paragraph to a large bold font and the second
paragraph to normal size italic.  Emphasize the words \dq{simple} and
\dq{short}.  (Again, you can \downloadorview{fonts} the solution.)

If you like, you can try experimenting with loading
different font packages, such as \isty{mathptmx}, to change the
default typeface. The \LaTeX\ Font Catalogue~\cite{fontcat} provides
a useful list of fonts, although you may not have all of them
installed.
\end{exercise}

%%%%%%%%%%%%%% Aligning Material in Rows and Columns %%%%%%%%%%%%%%%%

\setnode{tabular}
\section{Aligning Material in Rows and Columns}
\label{sec:tabular}

Text can be aligned in rows and columns using the \gls{env-tabular}
environment.
\begin{definition}
\glsni{begin}\marg{tabular}\oarg{\meta{pos}}\marg{\meta{column specifiers}}
\end{definition}%
This \glsterm{environment} has a
\glsterm{mandatory} \meta{column specifiers}
that specifies how to align each column. Within \meta{column
specifiers}, there must be a specifier for each column. The three
basic are: \texttt{r} (right aligned), \texttt{l} (left aligned) and
\texttt{c} (centred).
(Make sure you don't confuse \texttt{l} (the letter \dq{ell}) with \texttt{1} 
(the digit one).)
The optional argument \meta{pos} is covered 
\latexhtml{in \sectionref{sec:boxes}}{\htmlref{later}{sec:boxes}}.

\xminisec{Example:}
Three columns (\underline{l}eft, \underline{c}entred, \underline{c}entred):
\begin{codeS}
\glsni{begin}\marg{tabular}\marg{lcc}
\end{codeS} 

\xminisec{Another Example:}
Four columns (\underline{l}eft, \underline{c}entred, \underline{c}entred, \underline{r}ight):
\begin{codeS}
\glsni{begin}\marg{tabular}\marg{lccr}
\end{codeS} 

The \texttt{r}, \texttt{l} and \texttt{c} specifiers don't allow
line breaks or paragraphs within a cell. It's not a good idea to
have too much text in a cell, but if it's required you can use
\begin{definition}
p\marg{\meta{width}}
\end{definition}
which indicates a paragraph cell of the given width.

\xminisec{Example:}
Three columns (\underline{p}aragraph of width 1in,
\underline{c}entred, \underline{r}ight):
\begin{codeS}
\glsni{begin}\marg{tabular}\marg{p\marg{1in}cr}
\end{codeS}
The paragraph cell will be formatted fully justified, which is often
inappropriate for a narrow block of text. The \isty{array} package
provides
\begin{definition}
\gls{gtcol}\marg{\meta{declaration}}
\end{definition}
which can be used directly in front of the \texttt{l}, \texttt{c},
\texttt{r} or \texttt{p} column specifiers. This inserts
\meta{declaration} in front of the entries for that column, so it
can be used to insert, say, \gls{raggedright}.

\xminisec{Example:}
Three columns, the first \underline{l}eft justified where each entry
is in bold, the second a \underline{p}aragraph column of width 1in set to
ragged right and the third \underline{c}entered:
\begin{codeS}
\glsni{begin}\marg{tabular}\marg{\glsni{gtcol}\marg{\glsni{bfseries}}l\glsni{gtcol}\marg{\glsni{raggedright}}p\marg{1in}c}
\end{codeS}

The \isty{array} package also provides
\begin{definition}
\gls{ltcol}\marg{\meta{declaration}}
\end{definition}
which can be used directly after the \texttt{l}, \texttt{c},
\texttt{r} or \texttt{p} column specifiers. This inserts
\meta{declaration} after the entries for that column.

\xminisec{Inter-Column Gap:}

The gap between columns is given by twice the value of the
\gls{length} register:
\begin{definition}
\gls{tabcolsep}
\end{definition}
A gap of \glsni{tabcolsep} is
also inserted before the first column and after the last column.
This length can be changed using one of the commands described in
\sectionref{sec:length}. For example:
\begin{codeS}
\gls{setlength}\marg{\glsni{tabcolsep}}\marg{4pt}
\end{codeS}
This will put an 8pt gap between columns and a 4pt gap before the
first column and after the last column.

The column specifiers can also include:
\begin{definition}
\gls{atchar}\marg{\meta{inter-column text}}
\end{definition}
This inserts \meta{inter-column text} at that place on each row of
the table, replacing the default inter-column gap.

\xminisec{Example:}

Suppose we want a centred first column, a right justified second
column and a left justified third column with a dot between the
second and third columns:
\begin{codeS}
\begin{alltt}
\glsni{begin}\marg{tabular}\marg{cr\glsni{atchar}\marg{.}l}
\end{alltt}
\end{codeS}
Alternatively, you may want a larger gap between groups of columns,
for example, two groups of three left justified columns:
\begin{codeS}
\begin{alltt}
\glsni{begin}\marg{tabular}\marg{lll\glsni{atchar}\marg{\glsni{hspace}\marg{4\glsni{tabcolsep}}}lll}
\end{alltt}
\end{codeS}
This uses the command:
\begin{definition}
\gls{hspace}\marg{\meta{length}}
\end{definition}
which inserts a horizontal space of a given \gls{length}. In this
case, four times the value of \glsni{tabcolsep}. This makes the gap
between the third and fourth columns twice as wide as the gap
between the other columns.

\setnode{rowscols}
\subsection{Column and Row Separation}
\label{sec:rowscols}

Remember the \latexhtml{special characters mentioned in
\sectionref{sec:chars}}{\htmlref{special characters}{sec:chars}}?
The ampersand character \gls{ampchar} is used to separate column
entries. Rows are separated using:
\begin{definition}
\gls{tab.dbbackslashchar}\oarg{\meta{vertical space}}
\end{definition}
where \meta{vertical space} is extra vertical spacing between rows,
if required. There is also a longer equivalent:
\begin{definition}
\gls{tabularnewline}
\end{definition}
\faq{Alignment tab changed to \cmdname{cr}}{altabcr}%
If you have used something like
\texttt{\gls{gtcol}\marg{\gls{raggedright}}p\marg{\meta{length}}}
as the specifier for your last column, you must use
\glsni{tabularnewline} instead of \glsni{tab.dbbackslashchar}
to indicate the row break otherwise you will get the following
error:
\begin{verbatim}
! Extra alignment tab has been changed to \cr.
<recently read> \endtemplate
\end{verbatim}
\screenpagebreak

\xminisec{Example:} 

Let's have two columns, the first left justified and the second right justified:
\begin{code}
\glsni{begin}\marg{tabular}\marg{lr}\newline
Video \glsni{ampchar}
8.99\glsni{tab.dbbackslashchar}\newline
CD \glsni{ampchar} 9.11\glsni{tab.dbbackslashchar}\newline
DVD \glsni{ampchar}
15.00\glsni{tab.dbbackslashchar}\newline
Total \glsni{ampchar} 33.10\newline
\glsni{end}\marg{tabular}
\end{code}

\begin{result}[tabular1.html]
\begin{tabular}{lr}
Video & 8.99\\
CD & 9.11\\
DVD & 15.00\\
Total & 33.10
\end{tabular}
\end{result}%
Recall from \chapterref{ch:def} that \LaTeX\ ignores spaces at the start of a line and
treats multiple spaces as a single space, so I~could just have
easily done:
\begin{code}
\glsni{begin}\marg{tabular}\marg{lr}\newline
\strut~~Video \glsni{ampchar} 8.99\glsni{tab.dbbackslashchar}\newline
\strut~~CD~~~~\glsni{ampchar} 9.11\glsni{tab.dbbackslashchar}\newline
\strut~~DVD~~~\glsni{ampchar} 15.00\glsni{tab.dbbackslashchar}\newline
\strut~~Total \glsni{ampchar} 33.10\newline
\glsni{end}\marg{tabular}
\end{code}%
and still have got the same result, but now the code is easier to
read.

Entries form implicit \glsing{group}, so
\glspl{declaration} made
within a \gls{env-tabular} environment only have an effect up to the
next \gls{ampchar} or \gls{tab.dbbackslashchar}.  

\xminisec{Example:}
\begin{code}
\glsni{begin}\marg{tabular}\marg{lr}\newline
\strut~~Video~\glsni{ampchar} 8.99\glsni{tab.dbbackslashchar}\newline
\strut~~CD~~~~\glsni{ampchar} 9.11\glsni{tab.dbbackslashchar}\newline
\strut~~DVD~~~\glsni{ampchar}
15.00\glsni{tab.dbbackslashchar}\newline
\strut~~\glsni{bfseries} Total \glsni{ampchar} 33.10\newline
\glsni{end}\marg{tabular}
\end{code}%
Output:\screenpagebreak
\begin{result}[Image: as the previous example except that the word
Total has appeared in bold]
\begin{tabular}{lr}
Video & 8.99\\
CD & 9.11\\
DVD & 15.00\\
\bfseries Total & 33.10
\end{tabular}
\end{result}%
Let's add an extra column and a header row:
\begin{code}\obeyspaces
\glsni{begin}\marg{tabular}\marg{lrr}\newline
\strut~~Item~~\glsni{ampchar} ex VAT \glsni{ampchar} inc
VAT\glsni{tab.dbbackslashchar}\newline
\strut~~Video~\glsni{ampchar}~~~8.99 \glsni{ampchar}~~10.56 \glsni{tab.dbbackslashchar}\newline
\strut~~CD~~~~\glsni{ampchar}~~~9.11 \glsni{ampchar}~~10.70 \glsni{tab.dbbackslashchar}\newline
\strut~~DVD~~~\glsni{ampchar}~~15.00 \glsni{ampchar}~~17.63 \glsni{tab.dbbackslashchar}\newline
\strut~~\glsni{bfseries} Total \glsni{ampchar}~~33.10 \glsni{ampchar}~~39.89\newline
\glsni{end}\marg{tabular}
\end{code}%
Output:
\begin{result}[tabular2.html]
\begin{tabular}{lrr}
Item  & ex VAT & inc VAT\\ 
Video & 8.99 & 10.56\\
CD & 9.11 & 10.70\\
DVD & 15.00 & 17.63\\
\bfseries Total & 33.10 & 39.89
\end{tabular}
\end{result}%

\xminisec{Example (Aligning on a Decimal Point):}

\label{xmp:decpt}If you want to align on the decimal point, it's best to use the
\isty{siunitx} package. That's beyond the scope of this
book, but for simple data this can be achieved
using the \gls{atchar} inter-column specifier. For example:
\begin{code}\obeyspaces
\glsni{begin}\marg{tabular}\marg{lr\glsni{atchar}\marg{.}l}\newline
\strut~~Video \glsni{ampchar} 8 \glsni{ampchar} 99\glsni{tab.dbbackslashchar}\newline
\strut~~CD~~~~\glsni{ampchar} 9 \glsni{ampchar} 11\glsni{tab.dbbackslashchar}\newline
\strut~~DVD~~~\glsni{ampchar} 15 \glsni{ampchar} 00\glsni{tab.dbbackslashchar}\newline
\strut~~\glsni{bfseries} Total \glsni{ampchar} 33 \glsni{ampchar} 10\newline
\glsni{end}\marg{tabular}
\end{code}%
Output:\reportpagebreak
\begin{result}[Image: as the second example but the decimal points
are aligned]
\begin{tabular}{lr@{.}l}
Video & 8 & 99\\
CD & 9 & 11\\
DVD & 15 & 00\\
\bfseries Total & 33 & 10
\end{tabular}
\end{result}%

\setnode{multicolumn}
\subsection{Spanning Columns}
\label{sec:multicolumn}

You may have noticed I~omitted the column headers in the \htmlref{above
example}{xmp:decpt}. \faq{Merging cells in
a column of a table}{multirow}
The problem with rewriting the table using 
\texttt{r\glsni{atchar}\marg{.}l} to align the decimal point is that the header
now needs to span the last two columns.
This can be done using the command:
\begin{definition}
\gls{multicolumn}\marg{\meta{cols spanned}}\marg{\meta{col specifier}}\marg{\meta{text}}
\end{definition}%
The first \gls{mandatory} \meta{cols spanned} is the number of
columns you want to span, the second argument \meta{col specifier}
indicates how to align this column-spanning entry, the third
argument \meta{text} indicates what should go in this entry. Note
that \meta{col specifier} should only have a single column
specifier, such as \texttt{\marg{c}} or \texttt{\marg{r}}.  We can
use \glsni{multicolumn} to modify an earlier example as
follows:\screenpagebreak
\begin{code}
\glsni{begin}\marg{tabular}\marg{lrr}\newline
\strut~~~~~~~~\glsni{ampchar} \glsni{multicolumn}\marg{2}\marg{c}\marg{Price (\glsni{pounds})}\glsni{tab.dbbackslashchar}\newline
\strut~~Item~~\glsni{ampchar} ex VAT \glsni{ampchar} inc
VAT\glsni{tab.dbbackslashchar}\newline
\strut~~Video~\glsni{ampchar}~~~8.99 \glsni{ampchar}~~10.56 \glsni{tab.dbbackslashchar}\newline
\strut~~CD~~~~\glsni{ampchar}~~~9.11 \glsni{ampchar}~~10.70 \glsni{tab.dbbackslashchar}\newline
\strut~~DVD~~~\glsni{ampchar}~~15.00 \glsni{ampchar}~~17.63 \glsni{tab.dbbackslashchar}\newline
\strut~~\glsni{bfseries} Total \glsni{ampchar}~~33.10 \glsni{ampchar}~~39.89\newline
\glsni{end}\marg{tabular}
\end{code}%
Output:
\begin{result}[tabular2multicol.html]
\begin{tabular}{lrr}
 & \multicolumn{2}{c}{Price (\pounds)}\\
Item  & ex VAT & inc VAT\\ 
Video & 8.99 & 10.56\\
CD & 9.99 & 11.74\\
DVD & 15.00 & 17.63\\
\bfseries Total & 33.98 & 39.93
\end{tabular}
\end{result}%
Here we are spanning two columns, so the first argument to \glsni{multicolumn} is \verb|{2}|,
we want the entry centred, so the second argument is \verb|{c}| and the text to go in this entry 
is simply \texttt{\marg{Price (\gls{pounds})}}.

\faq{How to alter the alignment
of tabular cells}{tabcellalign}The \glsni{multicolumn} command can also be used to override
the alignment of individual entries.  Consider
the following example:
\begin{code}
\glsni{begin}\marg{tabular}\marg{lrr}\newline
\strut~~~~~~~~~~\glsni{ampchar} Year1   \glsni{ampchar} Year2
\glsni{tab.dbbackslashchar}\newline
Travel~~~~\glsni{ampchar} 100,000 \glsni{ampchar}
110,000\glsni{tab.dbbackslashchar}\newline
Equipment \glsni{ampchar} 50,000  \glsni{ampchar} 60,000\newline
\glsni{end}\marg{tabular}
\end{code}
Output:
\begin{result}[tabular3.html]
\begin{tabular}{lrr}
 & Year1 & Year2 \\
Travel & 100,000 & 110,000\\
Equipment & 50,000 & 60,000
\end{tabular}
\end{result}%
In this example, the headers \dq{Year1} and \dq{Year2} would look
better centred, but the rest of the entries in the second
and third columns look best right aligned.  We can
use \glsni{multicolumn} to span just one column, and
use the second argument of \glsni{multicolumn} to override
the column specification:
\begin{code}
\glsni{begin}\marg{tabular}\marg{lrr}\newline
\strut~~~~~~~~~~\glsni{ampchar} \glsni{multicolumn}\marg{1}\marg{c}\marg{Year1}\newline
\strut~~~~~~~~~~\glsni{ampchar} \glsni{multicolumn}\marg{1}\marg{c}\marg{Year2}\glsni{tab.dbbackslashchar}\newline
Travel~~~~\glsni{ampchar} 100,000 \glsni{ampchar} 110,000\glsni{tab.dbbackslashchar}\newline
Equipment \glsni{ampchar} 50,000  \glsni{ampchar} 60,000\newline
\glsni{end}\marg{tabular}
\end{code}
Output:
\begin{result}[Image: as the previous example except that the words
'Year1' and 'Year2' have been centred in their columns]
\begin{tabular}{lrr}
 & \multicolumn{1}{c}{Year1} & \multicolumn{1}{c}{Year2} \\
Travel & 100,000 & 110,000\\
Equipment & 50,000 & 60,000
\end{tabular}
\end{result}

\setnode{tabrules}
\subsection{Rules}
\label{sec:rules}

In general, vertical rules are considered superfluous~\cite{oxford}.
Although Turabian~\cite{turabian96} allows for the possibility of
vertical rules for tabulated material containing more than two
columns but still advises against having too many and deprecates
the use of them at either end.

Horizontal rules may be used at the top and bottom of the
tabulated material, but other horizontal rules should be kept to a
minimum. In general, the top and bottom rule should be thicker than
the mid rules.

The \isty{booktabs} package provides:
\begin{definition}
\gls{toprule}\oarg{\meta{wd}}
\end{definition}
for the top horizontal rule,
\begin{definition}
\gls{bottomrule}\oarg{\meta{wd}}
\end{definition}
for the bottom horizontal rule, and
\begin{definition}
\gls{midrule}\oarg{\meta{wd}}
\end{definition}
for horizontal rules in between rows, such as after the header row.

These commands should all go at the start of the appropriate row.
This means that if you want a bottom rule, you need to add
\gls{tab.dbbackslashchar} followed by \glsni{bottomrule} at the end of
the tabulated material.

\xminisec{Example:}
\begin{code}
\glsni{begin}\marg{tabular}\marg{lrr}\newline
\glsni{toprule}\newline
\strut~~~~~~~~~~\glsni{ampchar}
\glsni{multicolumn}\marg{1}\marg{c}\marg{Year1}\newline
\strut~~~~~~~~~~\glsni{ampchar}
\glsni{multicolumn}\marg{1}\marg{c}\marg{Year2}%
\glsni{tab.dbbackslashchar}\newline
\glsni{midrule}\newline
Travel~~~~\glsni{ampchar} 100,000 \glsni{ampchar}
110,000\glsni{tab.dbbackslashchar}\newline
Equipment \glsni{ampchar} 50,000  \glsni{ampchar} 60,000%
\glsni{tab.dbbackslashchar}\newline
\glsni{bottomrule}\newline
\glsni{end}\marg{tabular}
\end{code}
results in:
\begin{result}[Image: as previous example but with horizontal rules
at the start and end and between the first and second rows.]
\begin{tabular}{lrr}
\toprule
          & \multicolumn{1}{c}{Year1} 
          & \multicolumn{1}{c}{Year2} \\
\midrule
Travel    & 100,000 & 110,000\\
Equipment & 50,000  & 60,000
\\\bottomrule
\end{tabular}
\end{result}

The thickness of the top and bottom rule is given by
the \gls{length} register \booklinebreak\gls{heavyrulewidth}, and the thickness of
the mid rule is given by the \gls{length} register
\gls{lightrulewidth}. These rule thicknesses can be overridden using the
optional argument \meta{wd} for \glsni{toprule}, \glsni{midrule}
and \glsni{bottomrule}.

\begin{exercise}{Aligning Material}{ex:tabular}
Go back to the document you created in \exerciseref{ex:simple} (and later modified in \exerciseref{ex:datetime}),
and add the following: 
\begin{result}[tabex.html]
\begin{tabular}{lrr}
 & \multicolumn{2}{c}{\bfseries Expenditure (\pounds)}\\
 & \multicolumn{1}{c}{Year1} & \multicolumn{1}{c}{Year2} \\
\bfseries Travel & 100,000 & 110,000\\
\bfseries Equipment & 50,000 & 60,000
\end{tabular}
\end{result}

Note that the \gls{env-tabular} environment
doesn't create a caption, all it does is arrange its contents in
rows and columns. You can find out how to turn your
\gls{env-tabular} environment into a table 
in \sectionref{sec:tables}.

You can \downloadorview{tabular} the solution to this exercise. (Remember to check
\appendixref{ch:errors} if you encounter an error message.)
\end{exercise}

For more information about using the \gls{env-tabular} environment see \latexbook,
\latexguide\ or \latexcomp.  \emph{The \LaTeX\ Companion} also describes how
to span rows using the \isty{multirow}
\htmlref{package}{sec:packages}.  For information on how to create
coloured tables using the \isty{colortbl}
\htmlref{package}{sec:packages}, see \latexgraphic.

\xminisec{Related \gls{ukfaq} topics:}
\begin{itemize}
\item \faqlink{How to change a whole row of a table}{wholerow}
\item \faqlink{Merging cells in a column of a table}{multirow}
\item \faqlink{Fixed width tables}{fixwidtab}
\item \faqlink{Variable-width columns in tables}{varwidcol}
\item \faqlink{Spacing lines in tables}{struttab}
\end{itemize}

%%%%%%%%%%%%%%%% Boxes and Mini-Pages %%%%%%%%%%%%%%%%%%%

\setnode{boxes}
\section{Boxes and Mini-Pages}
\label{sec:boxes}

\gls{tex} views everything on a page as a form of box. Each box has an
associated width, height and depth, and the boxes are placed together
on the page with \keyword{glue}. This is reminiscent of the days
of manual typesetting, where each letter or symbol was on a wooden
block, and the wooden blocks were glued in place. The simplest form
of box is a single letter. Some letters, such as \dq{a} only have a
height and width, whereas other letters, such as \dq{y} have a
height, width and depth (see \figureref{fig:letterbox}).

\begin{figure}[htbp]
\figconts[alt=Image of the letters y and a illustrating their
height depth and width. The letter a has zero depth as it does not
extend below the baseline.]
  {pictures/letterbox}%
  {%
    \caption{\texorpdfstring{\protect\TeX}{TeX}\protect\space Views Each Letter as a Box}%
  }%
  {fig:letterbox}%
\end{figure}

For example, the phrase \dq{cabbages and peas} is made up of 15 boxes:
\begin{center}
\makeimg{c a b b a g e s a n d p e a s}{\setlength{\fboxsep}{0.5pt}%
\setlength{\fboxrule}{0.1pt}\LARGE
\fbox{c}\fbox{a}\fbox{b}\fbox{b}\fbox{a}%
\fbox{g}\fbox{e}\fbox{s}~\fbox{a}\fbox{n}\fbox{d}~\fbox{p}\fbox{e}%
\fbox{a}\fbox{s}}
\end{center}
whereas the word \dq{cauliflower} consists of 10 boxes:\footnote{The 
fl-ligature is a single character, and so is one box not
two.}
\begin{center}
\makeimg{c a u l i fl-ligature o u r}{\setlength{\fboxsep}{0.5pt}%
\setlength{\fboxrule}{0.1pt}\LARGE
\fbox{c}\fbox{a}\fbox{u}\fbox{l}\fbox{i}%
\fbox{fl}\fbox{o}\fbox{w}\fbox{e}\fbox{r}}
\end{center}

More complicated boxes are made up of smaller boxes. We have already
encountered one of these more complicated boxes: the 
\gls{env-tabular} environment, discussed in the 
\htmlref{previous section}{sec:tabular}. This type of box is called
a \keyword{horizontal box}, which means that it can go in a line
of text. For example:
\begin{code}\obeyspaces
Here is some text.\newline
\glsni{begin}\marg{tabular}\marg{cc}\newline
A \glsni{ampchar} B\glsni{tab.dbbackslashchar}\newline
C \glsni{ampchar} D\newline
\glsni{end}\marg{tabular}\newline
The rest of the line.
\end{code}
produces:
\begin{result}[Image illustrating the typeset output: the tabulated
material is centred vertically within the line of text.]
Here is some text.
\begin{tabular}{cc}
A & B\\
C & D
\end{tabular}
The rest of the line.
\end{result}
Recall from \htmlref{the previous section}{sec:tabular} that the 
\gls{env-tabular} environment had an optional argument
\meta{pos}. This governs the vertical alignment when the
\gls{env-tabular} environment occurs within a line of text.
This can be one of \texttt{c} (centred\dash the default, as illustrated
above), \texttt{t} (top) and \texttt{b} (bottom). For example,
\begin{code}\obeyspaces
Here is some text.\newline
\glsni{begin}\marg{tabular}\oarg{b}\marg{cc}\newline
A \glsni{ampchar} B\glsni{tab.dbbackslashchar}\newline
C \glsni{ampchar} D\newline
\glsni{end}\marg{tabular}\newline
The rest of the line.
\end{code}%
produces:
\begin{result}[Image illustrating the typeset output: the last row
of the tabulated material rests on the baseline of the text outside it.]
Here is some text.
\begin{tabular}[b]{cc}
A & B\\
C & D
\end{tabular}
The rest of the line.
\end{result}

Since a box can't be broken across a line of text, you can use
the box making command:
\begin{definition}
\gls{mbox}\marg{\meta{text}}
\end{definition}
to prevent \meta{text} from spanning a line break.

\xminisec{Example:}
Compare:
\begin{code}
\gls{raggedright}
Some text at the beginning of a paragraph. Some
text in the middle of the paragraph. Some more text.
\glsni{par}
\end{code}
\begin{result}[Image of typeset output: paragraph is left-justified]
\raggedright
Some text at the beginning of a paragraph. Some text in the
middle of the paragraph. Some more text.
\par
\end{result}
with:
\begin{code}
\gls{raggedright}
Some text at the beginning of a paragraph. \glsni{mbox}\marg{Some
text in the middle of the paragraph.} Some more text.
\glsni{par}
\end{code}
\begin{result}[Image of typeset output: line break occurs at the
start of the second sentence leaving a large white space at the end
of the first line.]
\raggedright
Some text at the beginning of a paragraph. \mbox{Some text in the
middle of the paragraph.} Some more text.
\par
\end{result}
(If \gls{raggedright} had not been used, the text in the \glsni{mbox}
would've spilt out over the edge of the page.)

Another type of box which can again be placed in a line of text,
is the \gls{env-minipage} environment.
\begin{definition}
\glsni{begin}\marg{minipage}\oarg{\meta{pos}}\oarg{\meta{height}}%
\marg{\meta{width}}
\end{definition}
As the name suggests, this environment creates a \dq{mini-page}
of the given width.

\xminisec{Example:}
\begin{code}
Some text.\newline
\glsni{begin}\marg{minipage}\marg{2in}\newline
This is a mini-page. The text inside
it is formatted as usual.\newline
\mbox{}\newline
Paragraph breaks can also be used, but 
there is no indentation by default\gls{footnote}\marg{and 
this is how a footnote appears}.\newline
\glsni{end}\marg{minipage}\newline
The rest of the line.
\end{code}%
which produces:\reportpagebreak
\begin{result}[Image: the contents of the minipage including
footnote have been typeset in a rectangular block within the outer
line of text.]
Some text.
\begin{minipage}{2in}
This is a mini-page. The text inside
it is formatted as usual.

Paragraph breaks can also be used, but 
there is no indentation by default\mpexfootnote{and 
this is how a footnote appears}.
\end{minipage}
The rest of the line.
\end{result}

You can optionally specify a height, and how
the mini-page is aligned with the rest of the text. As with
the \gls{env-tabular} environment, the alignment
option \meta{pos} can be one of \texttt{t} (top), \texttt{c} (centred)
or \texttt{b} (bottom). The default is \texttt{c}, which is why
the above example has the mini-page centred vertically. This
can be changed, for example:
\begin{code}
Some text.\newline
\glsni{begin}\marg{minipage}\oarg{t}\marg{2in}\newline
This is a mini-page. The text inside
it is formatted as usual.\newline
\mbox{}\newline
Paragraph breaks can also be used, but 
there is no indentation by default\gls{footnote}\marg{and 
this is how a footnote appears}.\newline
\glsni{end}\marg{minipage}\newline
The rest of the line.
\end{code}
which produces
\begin{result}[Image: as previous example but the first line of the
minipage is aligned with the outer text.]
\begin{makeimage}
Some text.
\begin{minipage}[t]{2in}
This is a mini-page. The text inside
it is formatted as usual.

Paragraph breaks can also be used, but 
there is no indentation by default\mpexfootnote{and this
is how a footnote appears}.
\end{minipage}
The rest of the line.
\end{makeimage}
\end{result}

Note that the width can be specified relative to the current
line width, using the \gls{length} register \gls{linewidth}. For example,
\begin{alltt}
\glsni{begin}\marg{minipage}\marg{0.5\glsni{linewidth}}
\end{alltt}
will start a mini-page that is half the width of the current
line.

There is also a corresponding command
\begin{definition}
\gls{parbox}\oarg{\meta{pos}}\oarg{\meta{height}}%
\marg{\meta{width}}\marg{\meta{text}}
\end{definition}%
which behaves in a similar way. So the above example can be 
rewritten using a \glsni{parbox}:
\begin{code}
Some text.\newline
\glsni{parbox}[t]\marg{2in}\marg{This is a parbox. 
The text inside
it is formatted as usual.\newline
\mbox{}\newline
Paragraph breaks can also be used, but 
there is no indentation by default.}\newline
The rest of the line.
\end{code}
which produces
\begin{result}[Image: as previous example but without a footnote]
Some text.
\parbox[t]{2in}{This is a parbox. 
The text inside
it is formatted as usual.

Paragraph breaks can also be used, but 
there is no indentation by default.}
The rest of the line.
\end{result}
You may have noticed that the \gls{footnote} command has
not been used in the above example.  The \glsni{parbox} command 
is more restricted than the
\gls{env-minipage} environment, so you can't use the 
\gls{footnote} command in it. There are also certain 
environments, such as the 
\htmlref{list-making
environments}{sec:lists}\latex{\ifscreenorother{}{\ 
described in \sectionref{sec:lists}}}, that can be
used in a \gls{env-minipage} but not in a \glsni{parbox}.

\setnode{fbox}
\subsection{Framed Boxes}
\label{sec:fbox}

Recall the \gls{framebox} command described in \sectionref{sec:optional}:
\begin{definition}
\glsni{framebox}\oarg{\meta{width}}\oarg{\meta{align}}\marg{\meta{text}}
\end{definition}
This treats \meta{text} as a box of width \meta{width} and puts a frame around it.
The second optional argument may be one of: \texttt{c} (centred
contents), \texttt{l} (left-aligned contents), \texttt{r}
(right-aligned contents).

\xminisec{Example:}

\begin{codeS}
Some \glsni{framebox}\oarg{2in}\marg{framed} text.
\end{codeS}
\begin{resultS}[Some framed text (the word 'framed' is inside a 2
inch wide box)]
Some \framebox[2in]{framed} text.
\end{resultS}

There is a shorter related command with no optional arguments:
\begin{definition}
\gls{fbox}\marg{\meta{text}}
\end{definition}

The \isty{fancybox} package provides some additional framing
commands:
\begin{definition}
\gls{shadowbox}\marg{\meta{text}}
\end{definition}
Puts a shadow-style frame around its contents:\bookpagebreak
\begin{codeS}
Some \glsni{shadowbox}\marg{framed} text.
\end{codeS}
\begin{resultS}[Some framed text (the word 'framed' is in a 
box with a shadow)]
Some \shadowbox{framed} text.
\end{resultS}

\begin{definition}
\gls{doublebox}\marg{\meta{text}}
\end{definition}
Puts a double-lined frame around its contents:
\begin{codeS}
Some \glsni{doublebox}\marg{framed} text.
\end{codeS}
\begin{resultS}[Some framed text (the word 'framed' is in a 
box with a double lined frame)]
Some \doublebox{framed} text.
\end{resultS}

\begin{definition}
\gls{ovalbox}\marg{\meta{text}}
\end{definition}
Puts a thin-lined oval frame around its contents:
\begin{codeS}
Some \glsni{ovalbox}\marg{framed} text.
\end{codeS}
\begin{resultS}[Some framed text (the word 'framed' is in a 
a box with rounded corners)]
Some \ovalbox{framed} text.
\end{resultS}

\begin{definition}
\gls{Ovalbox}\marg{\meta{text}}
\end{definition}
Puts a thick-lined oval frame around its contents:

\begin{codeS}
Some \glsni{Ovalbox}\marg{framed} text.
\end{codeS}
\begin{resultS}[Some framed text (the word 'framed' is in a 
a thick bordered box with rounded corners)]
Some \Ovalbox{framed} text.
\end{resultS}

If you want a different frame effect, you will need to use a
graphical package, such as \isty{pgf}/\isty{tikz}.

\xminisec{Example:}

This example uses commands beyond the scope of this book,
but gives you an idea of what's possible.

\begin{code}
\begin{verbatim}
\documentclass{scrartcl}

\usepackage{tikz}
\usetikzlibrary{shapes}
\usetikzlibrary{decorations.pathmorphing}

\begin{document}

Some
\begin{tikzpicture}[baseline=(n.base),decoration=bumps]
\node[draw,ellipse,decorate] (n) {framed};
\end{tikzpicture}
text.

\end{document}
\end{verbatim}
\end{code}
\begin{resultS}[Some framed text (the word 'framed' is in a 
decorated with a wiggly border)]
Some
\begin{tikzpicture}[baseline=(n.base),decoration=bumps]
\node[draw,ellipse,decorate] (n) {framed};
\end{tikzpicture}
text.
\end{resultS}

For further details, see the \isty{pgf} \htmlref{documentation}{sec:texdoc}. 

\xminisec{Related \gls{ukfaq} topics:}
\begin{itemize}
\item \faqlink{Automatic sizing of minipage}{varwidth}
\item \faqlink{Float(s) lost}{fllost}
\item \faqlink{Perhaps a missing \cmdname{item}?}{errmissitem}
\end{itemize}

%%%%%%%%%%%%%%%%% CHAPTERS, SECTIONS ETC %%%%%%%%%%%%%%%%%%%%%%%

\setnode{structure}
\chapter{Structuring Your Document}
\label{ch:sections}

Let's go back to the document we modified in \exerciseref{ex:tabular}.
In this \latexhtml{chapter}{section} we shall edit that document
step by step until we have a fully-fledged document with title,
abstract, table of contents, sections etc.

%%%%%%%%%%%%%%%%%% Maketitle %%%%%%%%%%%%%%%%%%%%%%%%%%%%%%%%%

\setnode{title}
\section{Author and Title Information}
\label{sec:title}

The term \keyword{title page} is used to indicate the author, title and
date information that can appear either on the front cover by itself
or along the top of the first page of text.  In order to do this, you
must first specify the information.  Once this information has been
specified it can then be displayed.

The author, title and date are entered using the 
\glspl{command}:
\begin{definition}
\gls{author}\marg{\meta{author names}}
\\%\end{definition}
%\begin{definition}
\gls{title}\marg{\meta{title text}}
\\%\end{definition}
%\begin{definition}
\gls{date}\marg{\meta{document date}}
\end{definition}%
The \koma\ classes also define:
\begin{definition}
\gls{titlehead}\marg{\meta{Title heading}}\\
\gls{subject}\marg{\meta{Subject}}\\
\gls{subtitle}\marg{\meta{Subtitle}}\\
\gls{publishers}\marg{\meta{Publisher}}
\end{definition}

\faq{The style of document titles}{titlsty}All these title-related
commands only \emph{store} information, they don't actually display
anything. These commands can be put in the \gls{preamble}. With most
classes, you will typically need to use at least \glsni{author} and
\glsni{title}.

Once you have used these commands, you can then
display the information using the command:\refstepcounter{object}\label{obj:maketitle}
\begin{definition}
\gls{maketitle}
\end{definition}
This command should be placed where you want the title page to
appear, which is normally at the start of the \gls{env-document}
environment.

Note that if you don't use the \glsni{date} command, the current 
date will be inserted.  If you want no date to appear, you need to 
specify an empty argument:
\begin{codeS}
\glsni{date}\marg{}
\end{codeS}
\bookpagebreak
Multiple authors should be separated by the command \gls{and}.
For example:
\begin{code}
\begin{alltt}
\glsni{author}\marg{A. Jones\glsni{tab.dbbackslashchar}University of Somewhere \glsni{and}
B. Smith\glsni{tab.dbbackslashchar}University of Somewhere Else}
\end{alltt}
\end{code}%
Within these titling fields, you can also use the command:
\begin{definition}
\gls{thanks}\marg{text}
\end{definition}%
which produces a special type of footnote. For example:
\begin{codeS}
\glsni{title}\marg{A Great Project\glsni{thanks}\marg{funded by XYZ}}
\end{codeS}%
Note that the footnote marker produced using \glsni{thanks}
is considered to have zero width, so if it occurs in the middle
of a line, rather than the end, you will need to insert some
extra space using \gls{spacesym} (backslash space).
The argument of \glsni{thanks} is a \glslink{fragile}{moving 
argument}\index{argument!moving|hyperpage}.

\begin{exercise}{Creating Title Pages}{ex:title}

Try editing the document you modified in \exerciseref{ex:tabular} to
include title information.  Modifications are illustrated in bold
\modification{like this}:

\begin{bcode}
\begin{alltt}
\glsni{documentclass}\oarg{12pt}\marg{scrartcl}

\glsni{usepackage}\marg{datetime}

\modification{\glsni{title}\marg{A Simple Document}}
\modification{\glsni{author}\marg{Me}}

\glsni{begin}\marg{document}

\modification{\glsni{maketitle}}

This is a simple \glsni{LaTeX}\glsni{spacesym}document.
Here is the first paragraph.

Here is the second paragraph\glsni{footnote}\marg{with a footnote}.
As you can see it's a rather short paragraph, but not
as short as the previous one. This document was
created on: \glsni{today}\glsni{spacesym}at \glsni{currenttime}.

\glsni{begin}\marg{tabular}\marg{lrr}
 \glsni{ampchar} \glsni{multicolumn}\marg{2}\marg{c}\marg{\glsni{bfseries} Expenditure}\glsni{tab.dbbackslashchar}
 \glsni{ampchar} \glsni{multicolumn}\marg{1}\marg{c}\marg{Year1} \glsni{ampchar} \glsni{multicolumn}\marg{1}\marg{c}\marg{Year2}\glsni{tab.dbbackslashchar}
\glsni{bfseries} Travel \glsni{ampchar} 100,000 \glsni{ampchar} 110,000\glsni{tab.dbbackslashchar}
\glsni{bfseries} Equipment \glsni{ampchar} 50,000 \glsni{ampchar} 60,000
\glsni{end}\marg{tabular}

\glsni{end}\marg{document}
\end{alltt}
\end{bcode}

You can \download{title} this document.
\end{exercise}

%%%%%%%%%%%%%%%%%%%%%%%%% Abstract %%%%%%%%%%%%%%%%%%%%%%%%

\setnode{abstr}
\section{Abstract}
\label{sec:abstract}

The \gls{env-abstract} \glsterm{environment} is used to
create an \Index{abstract} for the document\faq{1-column abstract in
2-column document}{onecolabs}.  The way in which the
abstract is formatted depends on the class file\indexCLS.  The
\icls{scrreprt} class file will put the abstract on a page by itself,
some class files will indent the abstract and some will typeset the
abstract in italic.  Note also that some class files (such as
\icls{scrbook}) don't have an \glsni{env-abstract}
environment.  Abstracts traditionally go at the start of the document
after the title, so the \glsni{env-abstract} environment should go
after the \gls{maketitle} command.

\begin{exercise}{Creating an Abstract}{ex:abstr}

Try editing your document so that it has an abstract:
Modifications are illustrated \modification{like this}:

\begin{bcode}
\begin{alltt}
\glsni{documentclass}\oarg{12pt}\marg{scrartcl}

\glsni{usepackage}\marg{datetime}

\glsni{title}\marg{A Simple Document}
\glsni{author}\marg{Me}

\glsni{begin}\marg{document}

\glsni{maketitle}

\modification{\glsni{begin}\marg{abstract}}
\modification{A brief document to illustrate how to use \glsni{LaTeX}.}
\modification{\glsni{end}\marg{abstract}}

This is a simple \glsni{LaTeX}\glsni{spacesym}document.
Here is the first paragraph.

Here is the second paragraph\glsni{footnote}\marg{with a footnote}.
As you can see it's a rather short paragraph, but not
as short as the previous one. This document was
created on: \glsni{today}\glsni{spacesym}at \glsni{currenttime}.

\glsni{begin}\marg{tabular}\marg{lrr}
 \glsni{ampchar} \glsni{multicolumn}\marg{2}\marg{c}\marg{\glsni{bfseries} Expenditure}\glsni{tab.dbbackslashchar}
 \glsni{ampchar} \glsni{multicolumn}\marg{1}\marg{c}\marg{Year1} \glsni{ampchar} \glsni{multicolumn}\marg{1}\marg{c}\marg{Year2}\glsni{tab.dbbackslashchar}
\glsni{bfseries} Travel \glsni{ampchar} 100,000 \glsni{ampchar} 110,000\glsni{tab.dbbackslashchar}
\glsni{bfseries} Equipment \glsni{ampchar} 50,000 \glsni{ampchar} 60,000
\glsni{end}\marg{tabular}


\glsni{end}\marg{document}
\end{alltt}
\end{bcode}

You can \download{abstr} this document.
\end{exercise}

%%%%%%%%%%%%%%%%% Sections, Subsections etc %%%%%%%%%%%%%%%%%%%%%%%%%%%%

\setnode{sectionunits}
\section{Chapters, Sections, Subsections \ldots}
\label{sec:section}

Chapters, sections, subsections etc can be inserted using the
commands:
\begin{definition}
\gls{part}\oarg{\meta{short title}}\marg{\meta{title}}\\
\gls{chapter}\oarg{\meta{short title}}\marg{\meta{title}}\\
\gls{section}\oarg{\meta{short title}}\marg{\meta{title}}\\
\gls{subsection}\oarg{\meta{short title}}\marg{\meta{title}}\\
\gls{subsubsection}\oarg{\meta{short title}}\marg{\meta{title}}\\
\gls{paragraph}\oarg{\meta{short title}}\marg{\meta{title}}\\
\gls{subparagraph}\oarg{\meta{short title}}\marg{\meta{title}}
\end{definition}%
All these commands have a \latexhtml{moving argument (see
\sectionref{sec:fragile})}{\glslink{fragile}{moving
argument}\index{argument!moving|hyperpage}}, so fragile commands will need to
be protected using \gls{protect}. \faq{How to create a
\cmdname{subsubsub\-section}}{subsubsub}The final two commands in the
above list, \glsni{paragraph} and \glsni{subparagraph}, represent
subsubsubsections and subsubsubsubsections, although most
\glspl{cls} typeset their arguments as unnumbered running titles.

Note that the availability of these commands depends on the
\gls{cls} you are using. For example, the \icls{scrartcl} class
file that we have been using is designed for articles, so the
\glsni{chapter} command is not defined in that class, whereas it is
defined in the \icls{scrreprt} and \icls{scrbook} class files.

Each of the commands above has a \gls{mandatory} \meta{title} and an
\gls{optional} \meta{short title}. The mandatory argument
\meta{title} is simply the title of the chapter\slash section\slash subsection
etc. For
example:
\begin{codeS}
\glsni{section}\marg{Introduction}
\end{codeS}%
If you are using the \icls{scrartcl} class file, the output will look 
like:
\begin{resultS}[Image: The text '1 Introduction' is typeset in a
large bold font]
\textbf{\sffamily\Large 1\quad Introduction}
\end{resultS}%
Note that you don't specify the section number as \LaTeX\ does this
automatically.  This means that you can insert a new section or
chapter or swap sections around or even change a section to a
subsection etc, without having to worry about updating all the
section numbers\faq{The style of section headings}{secthead}.

\renewcommand{\fmtord}[1]{\textsuperscript{#1}}%
If you are using a class file\indexCLS\ that contains chapters as
well as sections, the section number will depend on the chapter.  So,
for example, the current section is the \ordinal{section}\ section
of chapter~\arabic{chapter}, so the section number is \thesection.
\faq{Why are my sections numbered 0.1 ...?}{zerochap}(Note that if
you are using a class file where the section number depends on the
chapter number, you must have a \gls{chapter} command before your
first \gls{section} command, otherwise your section numbers will
come out as 0.1, 0.2 etc.)

Unnumbered chapters/sections etc are produced by placing an 
asterisk \texttt{*} after the command name.  For example:
\begin{codeS}
\glsni{chapter}*\marg{Acknowledgements}
\end{codeS}

You can switch to appendices using the 
command
\begin{definition}
\gls{appendix}
\end{definition}%
then continue using \gls{chapter}, \gls{section} 
etc\faq{Appendixes}{appendix}.  For example (using the \icls{scrreprt} class file):
\begin{code}
\begin{alltt}
\glsni{appendix}
\glsni{chapter}\marg{Derivations}
Some derivations.

\glsni{chapter}\marg{Tables}
Some tables.
\end{alltt}
\end{code}

\xminisec{Note:}
The \koma\ classes have another type of sectioning command:
\begin{definition}
\gls{minisec}\marg{\meta{heading}}
\end{definition}\screenpagebreak
This provides an unnumbered heading not associated with any of the
structuring levels. For example, the above was produce using:
\begin{code}
\begin{alltt}
\glsni{minisec}\marg{Note:}
The \koma\ classes have another type of sectioning command:
\end{alltt}
\end{code}
The next note below was produced using:
\begin{code}
\begin{alltt}
\glsni{minisec}\marg{Important Note:}
If you want to change the font style used by headings, 
\glsni{emph}\marg{\glsni{bfseries} do not} use font declarations in the sectioning
command arguments.
\end{alltt}
\end{code}

\xminisec{Important Note:}
If you want to change the font style used by headings, 
\emph{\bfseries do not} use font declarations in the sectioning command
arguments. Don't do, for example:
\begin{alltt}\wrong
\glsni{chapter}\marg{\glsni{itshape} Introduction}
\end{alltt}
The \koma\ classes provide the command:
\begin{definition}
\gls{addtokomafont}\marg{\meta{element}}\marg{\meta{commands}}
\end{definition}
where \meta{element} is the name of a structuring element (no
backslash) and \meta{commands} is the list of font changing
declarations (see \tableref{tab:fontsII}) to apply to that element style.
For example, \latexhtml{this book}{the PDF version of this document}
uses the commands:
\begin{code}
\begin{alltt}
\glsni{addtokomafont}\marg{section}\marg{\glsni{rmfamily}\glsni{bfseries}}
\glsni{addtokomafont}\marg{minisec}\marg{\glsni{rmfamily}\glsni{bfseries}\glsni{scshape}}
\end{alltt}
\end{code}

\begin{exercise}{Creating Chapters, Sections etc}{ex:section}

Let's try editing our document so that it now has chapters, sections
and an appendix.  Since the \icls{scrartcl} class file doesn't have
chapters, let's change to the \icls{scrreprt} class.  Changes from our
previous document are shown \modification{like this}.

\begin{bcode}%
\begin{alltt}
\glsni{documentclass}\oarg{12pt}\marg{\modification{scrreprt}}

\glsni{usepackage}\marg{datetime}

\glsni{title}\marg{A Simple Document}
\glsni{author}\marg{Me}

\glsni{begin}\marg{document}
\glsni{maketitle}

\glsni{begin}\marg{abstract}
A brief document to illustrate how to use \glsni{LaTeX}.
\glsni{end}\marg{abstract}

\modification{\glsni{chapter}\marg{Introduction}}

\modification{\glsni{section}\marg{The First Section}}
This is a simple \glsni{LaTeX}\glsni{spacesym}document.
Here is the first paragraph.

\modification{\glsni{section}\marg{The Next Section}}
Here is the second paragraph\glsni{footnote}\marg{with a footnote}. 
As you can see it's a rather short paragraph, but not 
as short as the previous one. This document was 
created on: \glsni{today}\glsni{spacesym}at \glsni{currenttime}.

\modification{\glsni{chapter}\marg{Another Chapter}}

\modification{Here\textquoteright{}s another very interesting chapter.}
\modification{We\textquoteright{}re going to put a picture here later.}

\modification{\glsni{chapter}*\marg{Acknowledgements}}

\modification{I would like to acknowledge all those}
\modification{very helpful people who have assisted me in my work.}

\modification{\glsni{appendix}}

\modification{\glsni{chapter}\marg{Tables}}
\modification{We will turn this tabular environment into a table later.}

\glsni{begin}\marg{tabular}\marg{lrr}
 \glsni{ampchar} \glsni{multicolumn}\marg{2}\marg{c}\marg{\glsni{bfseries} Expenditure}\glsni{tab.dbbackslashchar}
 \glsni{ampchar} \glsni{multicolumn}\marg{1}\marg{c}\marg{Year1} \glsni{ampchar} \glsni{multicolumn}\marg{1}\marg{c}\marg{Year2}\glsni{tab.dbbackslashchar}
\glsni{bfseries} Travel \glsni{ampchar} 100,000 \glsni{ampchar} 110,000\glsni{tab.dbbackslashchar}
\glsni{bfseries} Equipment \glsni{ampchar} 50,000 \glsni{ampchar} 60,000
\glsni{end}\marg{tabular}

\glsni{end}\marg{document}
\end{alltt}
\end{bcode}

(You can \download{section} a copy of this file if you like, but
I~recommend that you try editing the file yourself to give you
practice.)
\end{exercise}

%%%%%%%%%%%%%%%%%%% Table of Contents %%%%%%%%%%%%%%%%%%%%%%%%

\setnode{toc}
\section{Creating a Table of Contents}
\label{sec:toc}

Once you have all your \htmlref{sectioning commands}{sec:section},
such as \gls{chapter} and \gls{section},
you can create a table of contents with the command
\begin{definition}
\gls{tableofcontents}
\end{definition}%
This command should go where you want your table of
contents to appear (usually after \gls{maketitle}).\faq{The format of the
Table of Contents, etc}{tocloft} The \koma\ classes provide two
options that govern the format of the table of contents:
\scrclsopt[graduated]{toc} and \scrclsopt[flat]{toc}. The first is the
default and indents the different sectioning levels. The second
doesn't use any indentation.

\xminisec{Example:}

\begin{codeS}
\glsni{documentclass}\oarg{12pt,toc=flat}\marg{scrreprt}
\end{codeS}

You may recall from \htmlref{the previous section}{sec:section} that
the sectioning commands all had an optional argument \meta{short
title}.  If your chapter or section title is particularly long\faq{My
section title is too wide for the page header}{runheadtoobig}, you
can use \meta{short title} to specify a shorter title that should go
in the table of contents.\footnote{\label{ftn:header}and in the page header, depending
on the page style.}  The longer title (given by the other
argument \meta{title}) will still appear in the section heading in
the main part of the document.

\LaTeX\ processes all source code sequentially, so when it first
encounters the \glsni{tableofcontents} command, it doesn't yet
know anything about the chapters, sections etc.  So the first time
the \htmlref{document is \LaTeX{}ed}{ch:tex2pdf} the necessary
information is written to the table of contents 
(\texttt{.toc})\indexTOC\ file (see \sectionref{sec:auxiliary}).  The subsequent pass reads the
information in from the \texttt{.toc}\indexTOC\ file, and generates
the table of contents.  You will therefore need to \LaTeX\ your
document twice to make sure that the table of contents is
up-to-date\faq{Numbers too large in table of contents,
etc}{tocloftwrong}.

\xminisec{Adding Extra Information}

The \htmlref{starred versions}{itm:starredcommand} of the sectional 
commands (such as \glsnl{chapter}\texttt{*})
don't get added to the table of contents. It may be that you want to
add it, in which case you need to use
\begin{definition}
\gls{addcontentsline}\marg{\meta{toc}}\marg{\meta{section
unit}}\marg{\meta{text}}
\end{definition}
after the heading. The first argument \meta{toc} is the file
extension without the dot. As mentioned above, the table of contents
file has the extension \texttt{.toc}, so the first argument should
be \texttt{toc} (\htmlref{later}{ch:floats}\latex{ in \chapterref{ch:floats}}, 
we'll be adding a list of figures and a list of tables, and those
have file extensions \texttt{.lof}\indexLOF\ and \texttt{.lot}\indexLOT, respectively).
The second argument \meta{section unit} is the name of the section
unit. This is just the name of the relevant sectioning command
\emph{without} the backslash. The final argument \meta{text} is the
entry text. For example (using \icls{scrreprt} class):
\begin{code}
\gls{chapter}*\marg{Acknowledgments}\newline
\glsni{addcontentsline}\marg{toc}\marg{chapter}\marg{Acknowledgements}
\end{code}

\begin{exercise}{Creating a Table of Contents}{ex:toc}

Try modifying your document so that it has a table of contents.
Modifications from the previous exercise are illustrated 
\modification{like this}:
\begin{bcode}
\begin{alltt}
\glsni{documentclass}\oarg{12pt}\marg{scrreprt}

\glsni{usepackage}\marg{datetime}

\glsni{title}\marg{A Simple Document}
\glsni{author}\marg{Me}

\glsni{begin}\marg{document}

\glsni{maketitle}

\modification{\glsni{tableofcontents}}

\glsni{begin}\marg{abstract}
A brief document to illustrate how to use \glsni{LaTeX}.
\glsni{end}\marg{abstract}

\glsni{chapter}\marg{Introduction}

\glsni{section}\marg{The First Section}

This is a simple \glsni{LaTeX}\glsni{spacesym}document.  Here is the first paragraph.

\glsni{section}\marg{The Next Section}

Here is the second paragraph\glsni{footnote}\marg{with a footnote}. 
As you can see it's a rather short paragraph, but not 
as short as the previous one. This document was 
created on: \glsni{today}\glsni{spacesym}at \glsni{currenttime}.

\glsni{chapter}\marg{Another Chapter}

Here's another very interesting chapter.
We're going to put a picture here later.

\glsni{chapter}*\marg{Acknowledgements}

I would like to acknowledge all those
very helpful people who have assisted
me in my work.

\glsni{appendix}
\glsni{chapter}\marg{Tables}

We will turn this tabular environment into a table later.

\glsni{begin}\marg{tabular}\marg{lrr}
 \glsni{ampchar} \glsni{multicolumn}\marg{2}\marg{c}\marg{\glsni{bfseries} Expenditure}\glsni{tab.dbbackslashchar}
 \glsni{ampchar} \glsni{multicolumn}\marg{1}\marg{c}\marg{Year1} \glsni{ampchar} \glsni{multicolumn}\marg{1}\marg{c}\marg{Year2}\glsni{tab.dbbackslashchar}
\glsni{bfseries} Travel \glsni{ampchar} 100,000 \glsni{ampchar} 110,000\glsni{tab.dbbackslashchar}
\glsni{bfseries} Equipment \glsni{ampchar} 50,000 \glsni{ampchar} 60,000
\glsni{end}\marg{tabular}


\glsni{end}\marg{document}
\end{alltt}
\end{bcode}

If your table of contents doesn't come out right, try \LaTeX{}ing it
again.  (Again, you can \download{toc} this file.)

You might want to try experimenting with the
\scrclsopt[flat]{toc} class options to
see what difference it makes:
\begin{codeS}
\glsni{documentclass}\oarg{12pt,toc=flat}\marg{scrreprt}
\end{codeS}

\end{exercise}

%%%%%%%%%%%%%%%%% Cross-Referencing %%%%%%%%%%%%%%%%%%%%%%%%%%%%

\setnode{crossref}
\section{Cross-Referencing}
\label{sec:crossref}

We have already seen that \LaTeX\ takes care of all the numbering for
the chapters etc, but what happens if you want to refer to a chapter
or section?  There's no point leaving \LaTeX\ to automatically
generate the section numbers if you have to keep track of them all,
and change all your cross-references every time you add a new
section.  Fortunately \LaTeX\ provides a way to generate the correct
number. All you have to do is label the part of the document you want
to reference, and then refer to this label when you want to
cross-reference it\faq{Referring to labels in other
documents}{extref}.  \LaTeX\ will then determine the correct number
that needs to be inserted at that point.

The first part, labelling the place you want to reference, is done 
using the command:
\begin{definition}
\gls{label}\marg{\meta{string}}
\end{definition}\refstepcounter{object}\label{obj:label}%
The \gls{argument} \meta{string} should be a
unique textual label.  This label can be anything you like as long as
it is unique, but it's a good idea to make it something obvious
so that, firstly, you can remember the label when you want to 
use it, and secondly, when you read through your code at some later
date, it's immediately apparent to you to which part of the
document you are referring.  People tend to have their own conventions
for labelling.  I usually start the label with two or three letters
that signify what type of thing I'm labelling.  For example, if I'm
labelling a chapter I'll start with \texttt{ch}, if I'm labelling a
section I'll start with \texttt{sec}. 

\xminisec{Examples:}

\begin{enumerate}
\item Labelling a chapter:

\begin{code}
\begin{alltt}
\glsni{chapter}\marg{Introduction}
\glsni{label}\marg{ch:intro}
\end{alltt}
\end{code}

\item Labelling a section:

\begin{code}
\begin{alltt}
\glsni{section}\marg{Technical Details}
\glsni{label}\marg{sec:details}
\end{alltt}
\end{code}
\end{enumerate}

Note that the \glsni{label} command doesn't produce any text, it
simply assigns a label.  You can now refer to that object
using the command:
\begin{definition}
\gls{ref}\marg{\meta{string}}
\end{definition}
which will produce the relevant number.

\xminisec{Example:}
\begin{codeS}
See Section \glsni{ref}\marg{sec:results} for an analysis
of the results.
\end{codeS}
\refstepcounter{object}\label{cmd:tilde}%
It is a typographical convention that you should never
start a new line with a number.  For example,
if you have the text \dq{Chapter~1} the \dq{1} must
be on the same line as the \dq{Chapter}.  We can 
do this by using an \keyword{unbreakable space}, which will
put a space but won't allow \LaTeX\ to break the line at
that point.  This is done using the tilde (\gls{tildechar}) 
\htmlref{special character}{sec:chars},
so the example above should actually be:\bookpagebreak
\begin{codeS}
See Section\glsni{tildechar}\glsni{ref}\marg{sec:results} for an analysis
of the results.
\end{codeS}
There is a similar command to reference the page number:
\begin{definition}
\gls{pageref}\marg{\meta{string}}
\end{definition}

\xminisec{Example:}
\begin{code}
See Chapter\glsni{tildechar}\glsni{ref}\marg{ch:def} on
page\glsni{tildechar}\glsni{pageref}\marg{ch:def} for a list of definitions.
\end{code}
The label \texttt{ch:def} obviously needs to
be defined somewhere:
\begin{code}
\begin{alltt}
\glsni{chapter}\marg{Definitions}
\glsni{label}\marg{ch:def}
\end{alltt}
\end{code}
In fact, I~have done this in my source code for \chapterref[the
definitions section]{ch:def} of this document, so the above example
would look like:
\begin{resultS}[See Chapter 2 on page 8 for a list of definitions]
See Chapter~\ref*{ch:def} on
page~\pageref*{ch:def} for a list of definitions.
\end{resultS}

\faq{Referring to things by their name}{nameref}%
It's not just chapters and sections that you can reference, most of
the numbers that \LaTeX\ automatically generates can be
cross-referenced. 

\xminisec{Example:} 
The source code for footnote~\ref{ftn:header} \latexhtml{on
page~\pageref{ftn:header}}{in \sectionref{sec:toc}} is:
\begin{code}
\glsni{footnote}\marg{\glsni{label}\marg{ftn:header}and in the page header, depending
on the page style}
\end{code}
\latexhtml{and it was referenced above using:}{and can be referenced
using:}
\begin{code}
The source code for footnote\glsni{tildechar}\glsni{ref}\marg{ftn:header} on
page\glsni{tildechar}\glsni{pageref}\marg{ftn:header} is:
\end{code}

The \isty{varioref} package provides a more convenient way of doing this using the
command:
\begin{definition}
\gls{vref}\marg{\meta{label}}
\end{definition}
This is like \gls{ref} but also adds information about the location, 
such as \dq{on page~\meta{n}} or \dq{on the following page}, if the
corresponding \gls{label} occurs on a different page, so the above
example can be changed to:
\begin{codeS}
The source code for footnote\glsni{tildechar}\glsni{vref}\marg{ftn:header} is:
\end{codeS}
\begin{latexonly}
\ifscreenorother
{%
which produces
\begin{resultS}
The source code for footnote~\vref{ftn:header} is:
\end{resultS}
}%
{%
\bookpagebreak\noindent which still produces
\begin{resultS}
The source code for footnote~\vref{ftn:header} is:
\end{resultS}
Compare with a reference to one of the labels in the next example:
\begin{codeS}
See step\glsni{tildechar}\cmdname{vref}\marg{itm:edit}.
\end{codeS}
which produces:
\begin{resultS}
See step~\vref{itm:edit}.
\end{resultS}
}%
\end{latexonly}
\begin{htmlonly}
which produces
\begin{resultS}[The source code for footnote 5.1 on page 110 is:]
The source code for footnote~\vref{ftn:header} is:
\end{resultS}
\end{htmlonly}

\xminisec{Caveat:}
You can run into trouble if the \glsni{vref} command occurs on a
page break. When it tries to insert the location information, such
as \dq{on the next page}, the information is no longer correct. This
can cause an \dq{Infinite loop} error. When this happens, either
edit your paragraph so the reference no longer falls on the page
break or use \gls{ref} instead of \glsni{vref} for that instance.

\xminisec{Another Example:}

The \gls{env-enumerate} environment \latex{\ifscreenorother{}{described 
in \sectionref{sec:enumerate} }}%
automatically numbers the items within an ordered list, so it's
possible to label list items. Recall the numbered list of
instructions at the start of
\chapterref{ch:tex2pdf}. Here's the code:

\begin{code}[0.9\linewidth]
\begin{alltt}
\glsni{begin}\marg{enumerate}

   \glsni{item}\glsni{label}\marg{itm:edit} Write or edit the source code.

   \glsni{item} Pass the source code to the \glsni{texttt}\marg{latex} or \glsni{texttt}\marg{pdflatex} 
        application (``\glsni{LaTeX}\glsni{spacesym}the document'').

      \glsni{begin}\marg{itemize}

         \glsni{item} If there are any error messages, 
               return to Step\glsni{tildechar}\glsni{ref}\marg{itm:edit}.

         \glsni{item} If there are no error messages, a PDF file
               is created, go to Step\glsni{tildechar}\glsni{ref}\marg{itm:view}.
      \glsni{end}\marg{itemize}

   \glsni{item}\glsni{label}\marg{itm:view} View the PDF file to check the result.

\glsni{end}\marg{enumerate}
\end{alltt}
\end{code}
Output:
\begin{result}[enumerateref.html]
\begin{enumerate}

   \item\label{itm:edit} Write or edit the source code.

   \item Pass the source code to the \texttt{latex} or \texttt{pdflatex} application 
        (\dq{\LaTeX\ the document}).

      \begin{itemize}

         \item If there are any error messages, 
               return to Step~\ref*{itm:edit}.

         \item If there are no error messages, a PDF file
               is created, go to Step~\ref*{itm:view}.
      \end{itemize}

   \item\label{itm:view} View the PDF file to check the result.

\end{enumerate}
\end{result}

The \gls{ref} and \gls{pageref} commands may come before or
after the corresponding \gls{label} command.  As with the table
of contents, \LaTeX\ first writes out all the cross-referencing
information to another file (the \iauxfile, see \sectionref{sec:auxiliary}) and then reads it in the
next time, so you will need to \LaTeX\ your document twice to get
everything up-to-date.

If the references aren't up-to-date, you will see the following
message at the end of the \LaTeX\ run:
\begin{verbatim}
LaTeX Warning: Label(s) may have changed. 
Rerun to get cross-references right.
\end{verbatim}
The following warning
\begin{verbatim}
LaTeX Warning: There were undefined references.
\end{verbatim}
means that \LaTeX\ found a reference to a label that
does not appear in the auxiliary file\indexAUX.  This could
mean that it's a new label, and the warning
will go away the next time you \LaTeX\ your document,
or it could mean that either you've forgotten to
define your label with the \gls{label} command,
or you've simply misspelt the label. The undefined references
will show up as two question marks \doublequestionmark\ in the \gls{output}.

\faq{\dq{Rerun} messages won't go away}{rerun}Very occasionally, if you have cross-references
and a table of contents, you might have to \LaTeX\ 
your document three times to get everything up to
date.  Just check to see if the \texttt{Label(s) may have changed}
warning appears.

If you find it inconvenient having to remember to click the typeset
button twice, you can use \iappname{latexmk}. This will run \LaTeX\
the required number of times to ensure the document is up-to-date.
To do this in TeXWorks, change the drop-down menu to \dq{LaTeXmk}, as
illustrated in \figureref{fig:latexmk}. Note that \appname{latexmk}
is a \gls{perl} script, so you need to make sure you have
\iappname{perl} installed (see \sectionref{sec:perl}).

\begin{figure}[htbp]
\figconts
 {pictures/texworks-latexmk}
 {%
   \caption{Selecting LaTeXmk in TeXWorks}
 }%
 {fig:latexmk}
\end{figure}

If \appname{latexmk} isn't listed in the drop-down menu, you can add
it via \menu{Edit}\menuto\menu{Preferences}. This opens the dialog
box shown in \figureref{fig:texworks-pref}. You can add a new tool
as follows:

\begin{enumerate}
\item To the right of the box
labelled \dq{Processing Tools} there is a button marked with a plus
(+) sign. Click on it to open the tool configuration dialog, shown
in \figureref{fig:texworks-toolconfig}.

\item Fill in the name \dq{LaTeXmk} in the box labelled \dq{Name} and either type
in the location of \appname{latexmk} in the box labelled \dq{Program} 
or use the \dq{Browse} button to locate it on your filing system.
(See \figureref{fig:texworks-toolconfig2}.) This will vary depending
on your operating system and \gls{tex}-distribution, but it will
probably be in a subdirectory (folder) called \texttt{bin} somewhere
in the \gls*{tex}-distribution tree.

\item There are lots of options that can be passed to
\appname{latexmk}, but if you want to produce PDF output you need to
add \texttt{-pdf} as an argument. This is done by clicking on the
button marked with a plus to the right of the \dq{Arguments} box and
type \texttt{-pdf}, as shown in \figureref{fig:texworks-toolconfig3}.

\item Another argument needs to be added that specifies the
basename of the \LaTeX\ file. This is done by again clicking on the
plus button and typing \verb|$basename|, as shown in
\figureref{fig:texworks-toolconfig4} .

\item Click on \dq{OK} to close the Tool Configuration dialog.

\item If you want to set \appname{latexmk} to be your default
processing tool, you can select it from the drop-down list labelled
\dq{Default}.

\item Click \dq{OK} when you're done.
\end{enumerate}

\begin{figure}[htbp]
\figconts
 {pictures/texworks-preferences}
 {%
    \caption{TeXWorks Preferences}
 }
 {fig:texworks-pref} 
\end{figure}

\begin{figure}[htbp]
\figconts
 {pictures/texworks-toolconfig1}
 {%
    \caption{Tool Configuration Dialog}
 }
 {fig:texworks-toolconfig} 
\end{figure}

\begin{figure}[htbp]
\figconts
 {pictures/texworks-toolconfig2}
 {%
    \caption{Tool Configuration Dialog: set the name and program
location}
 }
 {fig:texworks-toolconfig2} 
\end{figure}

\begin{figure}[htbp]
\figconts
 {pictures/texworks-toolconfig3}
 {%
    \caption{Tool Configuration Dialog: adding \protect\texttt{-pdf} argument}
 }
 {fig:texworks-toolconfig3} 
\end{figure}

\begin{figure}[htbp]
\figconts
 {pictures/texworks-toolconfig4}
 {%
    \caption{Tool Configuration Dialog: adding \protect\texttt{\$basename} argument}
 }
 {fig:texworks-toolconfig4} 
\end{figure}

\begin{exercise}{Cross-Referencing}{ex:crossref}

Try modifying your code so that it has cross-references.
Again, changes made from the previous exercise are
illustrated \modification{like this}:
\begin{bcode}
\begin{alltt}
\glsni{documentclass}\oarg{12pt}\marg{scrreprt}

\glsni{usepackage}\marg{datetime}

\glsni{title}\marg{A Simple Document}
\glsni{author}\marg{Me}

\glsni{begin}\marg{document}

\glsni{maketitle}

\glsni{tableofcontents}

\glsni{begin}\marg{abstract}
A brief document to illustrate how to use \glsni{LaTeX}.
\glsni{end}\marg{abstract}

\glsni{chapter}\marg{Introduction}
\modification{\glsni{label}\marg{ch:intro}}

\glsni{section}\marg{The First Section}

This is a simple \glsni{LaTeX}\glsni{spacesym}document.  Here is the first paragraph.
\modification{The next chapter is Chapter\glsni{tildechar}\glsni{ref}\marg{ch:another}}
\modification{and is on page\glsni{tildechar}\glsni{pageref}\marg{ch:another}.}
\modification{The next section is Section\glsni{tildechar}\glsni{ref}\marg{sec:next}.}

\glsni{section}\marg{The Next Section}
\modification{\glsni{label}\marg{sec:next}}

Here is the second paragraph\glsni{footnote}\marg{with a footnote}. 
As you can see it's a rather short paragraph, but not 
as short as the previous one. This document was 
created on: \glsni{today}\glsni{spacesym}at \glsni{currenttime}.

\glsni{chapter}\marg{Another Chapter}
\modification{\glsni{label}\marg{ch:another}}

Here's another very interesting chapter.
We're going to put a picture here later.
\modification{See Chapter\glsni{tildechar}\glsni{ref}\marg{ch:intro} for an }
\modification{introduction.}

\glsni{chapter}*\marg{Acknowledgements}

I would like to acknowledge all those
very helpful people who have assisted
me in my work.

\glsni{appendix}
\glsni{chapter}\marg{Tables}

We will turn this tabular environment into a table later.

\glsni{begin}\marg{tabular}\marg{lrr}
 \glsni{ampchar} \glsni{multicolumn}\marg{2}\marg{c}\marg{\glsni{bfseries} Expenditure}\glsni{tab.dbbackslashchar}
 \glsni{ampchar} \glsni{multicolumn}\marg{1}\marg{c}\marg{Year1} \glsni{ampchar} \glsni{multicolumn}\marg{1}\marg{c}\marg{Year2}\glsni{tab.dbbackslashchar}
\glsni{bfseries} Travel \glsni{ampchar} 100,000 \glsni{ampchar} 110,000\glsni{tab.dbbackslashchar}
\glsni{bfseries} Equipment \glsni{ampchar} 50,000 \glsni{ampchar} 60,000
\glsni{end}\marg{tabular}

\glsni{end}\marg{document}
\end{alltt}
\end{bcode}
(You can \download{crossref} a copy of this file.)

\end{exercise}

%%%%%%%%%%%%%%%%%%%% Bibliography %%%%%%%%%%%%%%%%%%%%%%%%%

\setnode{biblio}
\section{Creating a Bibliography}
\label{sec:bib}

If you have a large number of citations in your document, it's best
to use an external bibliographic application\faq{Creating a BibTeX
bibliography file}{buildbib}, such as \iappname{bibtex} or
\iappname{biber}. However, that is beyond the scope of this
book (see, instead, \latexguide, \latexcomp\
or \latexthesis). Therefore this section just gives a brief
explanation of the \gls{env-thebibliography} environment, which is
usually automatically generated using \appname{bibtex} or
\appname{biber}.

\begin{definition}
\glsni{begin}\marg{thebibliography}\marg{\meta{widest tag}}
\end{definition}
This environment is very similar to the list making environments described
in \sectionref{sec:lists}, but instead of \gls{item} use 
\begin{definition}
\gls{bibitem}\oarg{\meta{tag}}\marg{\meta{key}}
\end{definition}%
where \meta{key} is a unique keyword that identifies this
item.  Your keyword can be anything you like, but as with
\glsni{label} I~recommend that you use
a short memorable keyword. I~tend to use the first author's surname
followed by the year of publication. 

The bibliography heading depends on the class file you are using.
Most of the article-style classes, such as \icls{scrreprt}, use
\glsi{refname} (which produces \dq{\refname}) in an unnumbered
section, whereas the report and book-styles, such as \icls{scrreprt}
and \icls{scrbook}, use \glsi{bibname} (which produces
\dq{\bibname}) in an unnumbered chapter. See
\tableref{tab:textlab} for the list of the common textual label
commands.

Most class files don't automatically add the bibliography to the
table of contents. The \koma\ classes provide the
\scrclsopt{bibliography} option. This can be
\scrclsopt[totoc]{bibliography} (an unnumbered unit added to the
table of contents), for example,
\begin{codeS}
\glsni{documentclass}\oarg{bibliography=totoc}\marg{scrreprt}
\end{codeS}
or \scrclsopt[totocnumbered]{bibliography} (a
numbered unit added to the table of contents), for example,
\begin{codeS}
\glsni{documentclass}\oarg{bibliography=totocnumbered}\marg{scrreprt}
\end{codeS}
If you're not using
one of the \koma\ classes, consult the 
\htmlref{documentation}{sec:texdoc} for your class to see if
there is an equivalent option. Failing that, you can use
\gls{addcontentsline} (described in \sectionref{sec:toc}).
For example (using a class that defines chapters):
\begin{code}
\glsni{addcontentsline}\marg{toc}\marg{chapter}\marg{\glsni{bibname}}\\
\glsni{begin}\marg{bibliography}\marg{1}
\end{code}

\xminisec{Example:}
(This example uses the command \gls{LaTeXe} which produces the
\LaTeXe\ logo. This indicates the current version of \LaTeX\ rather
than the old 2.09 version.)\footnote{If a friend or colleague gives
you a file containing \cmdname{documentstyle} instead of \cmdname{documentclass}
they are nearly 20~years out of date.}
The class style in use is \cls{scrbook}, so the title is
given by \glsni{bibname} (\dq{Bibliography}).
\begin{code}
\glsni{begin}\marg{thebibliography}\marg{3}\newline
\glsni{bibitem}\marg{lamport94} \glsni{quotedblleft}\glsni{LaTeX}: a document preparation 
system\glsni{quotedblright}, Leslie Lamport, 2nd edition (updated for 
\glsni{LaTeXe}), Addison-Wesley (1994).\newline
\mbox{}\newline
\glsni{bibitem}\marg{kopka95} \glsni{quotedblleft}A Guide to \glsni{LaTeX}: document 
preparation for beginners and advanced users\glsni{quotedblright}, 
Helmut Kopka and Patrick W. Daly, Addison-Wesley (1995).\newline
\mbox{}\newline
\glsni{bibitem}\marg{goossens94} \glsni{quotedblleft}The 
\glsni{LaTeX}\glsni{spacesym}Companion\glsni{quotedblright}, 
Michel Goossens, Frank Mittelbach and
Alexander Samarin, Addison-Wesley, (1994).\newline
\mbox{}\newline
\glsni{end}\marg{thebibliography}
\end{code}
\begin{result}[references.html]
\dobibexamplenumbered
\end{result}

You can cite an item in your bibliography with the command
\begin{definition}
\gls{cite}\oarg{\meta{text}}\marg{\meta{key list}}
\end{definition}%

\bookpagebreak
\xminisec{Example:}
\begin{code}
\begin{alltt}
For more information about writing bibliographies see 
Goossens \glsni{emph}\marg{et al.}\glsni{tildechar}\glsni{cite}\marg{goossens94}.
\end{alltt}
\end{code}%
Output:
\begin{resultS}[cite.html]
For more information about writing
bibliographies see 
Goossens \emph{et al.}~[3].
\end{resultS}
If you want to cite multiple works, use 
a comma-separated list:
\begin{code}
\begin{alltt}
For more information about writing bibliographies 
see\glsni{tildechar}\glsni{cite}\marg{kopka95,goossens94}.
\end{alltt}
\end{code}%
Output:
\begin{resultS}[mcite.html]
For more information about writing bibliographies 
see~[2,3].
\end{resultS}
The \gls{optional} \meta{text} to
the \glsni{cite} command can be used to add text to the
citation.  

\xminisec{Example:}
\begin{code}
\begin{alltt}
For more information about writing bibliographies see 
Goossens \glsni{emph}\marg{et al.}\glsni{tildechar}\glsni{cite}\oarg{Chapter~13}\marg{goossens94}.
\end{alltt}
\end{code}%
Output:
\begin{result}[chcite.html]
For more information about writing bibliographies see 
Goossens \emph{et al.}~[3, Chapter~13].
\end{result}

The \gls{env-thebibliography} environment has
a \gls{mandatory}:
\begin{definition}
\glsni{begin}\marg{thebibliography}\marg{\meta{widest tag}}
\end{definition}%
The argument \meta{widest tag} is the widest tag in the list of
entries.  This helps \LaTeX\ to align the references correctly.  In
the example above, the tags appeared as: [1], [2] and [3], and [3]
is the widest so that was used as the argument. These tags can be
changed from the default numbers to something else using
the optional argument to the \gls{bibitem} command. 

\xminisec{Example (Textual Tags):}

This example uses the \gls{optional} of \gls{bibitem} to use textual
rather than numerical tags.  The widest tag is [Goossens 1994]
so that is chosen to be the argument of the \gls{env-thebibliography}
environment:
\begin{code}
\glsni{begin}\marg{thebibliography}\marg{Goossens 1994}\newline
\mbox{}\newline
\glsni{bibitem}[Lamport 1994]\marg{lamport94} 
\glsni{quotedblleft}\glsni{LaTeX}\cmdname{ } : a document 
preparation system\glsni{quotedblright}, Leslie Lamport, 2nd edition 
(updated for \glsni{LaTeXe}), Addison-Wesley (1994).\newline
\mbox{}\newline
\glsni{bibitem}[Kopka 1995]\marg{kopka95} 
\glsni{quotedblleft}A Guide to \glsni{LaTeX}: document 
preparation for beginners and advanced users\glsni{quotedblright}, Helmut Kopka 
and Patrick W. Daly, Addison-Wesley (1995).\newline
\mbox{}\newline
\glsni{bibitem}[Goossens 1994]\marg{goossens94} 
\glsni{quotedblleft}The \glsni{LaTeX}\glsni{spacesym}Companion\glsni{quotedblright}, 
Michel Goossens, Frank Mittelbach and
Alexander Samarin, Addison-Wesley, (1994).\newline
\mbox{}\newline
\glsni{end}\marg{thebibliography}
\end{code}
\begin{result}[references2.html]
\dobibexamplelabelled
\end{result}

\begin{exercise}{Creating a Bibliography}{ex:biblio}

Try adding the following chapter to your document:
\begin{code}
\glsni{chapter}\marg{Recommended Reading}

For a basic introduction to \glsni{LaTeX}\glsni{spacesym}see
Lamport\glsni{tildechar}\glsni{cite}\marg{lamport94}. For
more detailed information about \glsni{LaTeX}\glsni{spacesym}and
associated applications, consult Kopka and
Daly\glsni{tildechar}\glsni{cite}\marg{kopka95} or Goossens
\glsni{emph}\marg{et
al}\glsni{tildechar}\glsni{cite}\marg{goossens94}.
\end{code}% 
and also add the bibliography shown above to the end of your
document. You can \downloadorview{biblio} the solution, but have a
go by yourself first. Remember that, as before, you will need to
\LaTeX\ the document twice to get the references up-to-date, unless
you're using \iappname{latexmk} (as described in
\sectionref{sec:crossref}) in which case it will be done
automatically.

\end{exercise}

%%%%%%%%%%%%%%%%%%%%%%%%% Page Styles %%%%%%%%%%%%%%%%%%%%%%%%%%

\setnode{pagestyle}
\section{Page Styles and Page Numbering}
\label{sec:pagestyle}

You may have noticed that the documents you have created
have all had their page numbers automatically inserted 
at the foot of most of the pages\faq{Page numbering 
\dq{\meta{n} of \meta{m}}}{nofm}.  If you have created the document
that has gradually been modified over the previous few
sections, you may have noticed that the title page has
no header or footer, the table of contents starts on page~1,
the abstract page has no page number, and the pages after
the abstract start on page~1 and continue incrementally
onwards from that point.  All the page numbers are Arabic
numerals.  This can be changed using the command:
\begin{definition}
\gls{pagenumbering}\marg{\meta{style}}
\end{definition}%
where \meta{style} can be one of:
\begin{fwlist}{\pagenumberingfmt{arabic}}
\fwitem{\ipagenumbering{arabic}} Arabic numerals (1, 2, 3, \ldots)
\fwitem{\ipagenumbering{roman}} Lower case Roman numerals (i, ii, iii, \ldots)
\fwitem{\ipagenumbering{Roman}} Upper case Roman numerals (I, II, III, \ldots)
\fwitem{\ipagenumbering{alph}} Lower case alphabetical characters (a, b, c, \ldots)
\fwitem{\ipagenumbering{Alph}} Upper case alphabetical characters (A, B, C, \ldots)
\end{fwlist}

Traditionally, the front matter (table of contents, list of figures
etc) should have lower case Roman numeral page numbering, while the
main matter should be in Arabic numerals\faq{Page numbering by
chapter}{pagebychap}.  

\xminisec{Example:}
\begin{code}
\begin{alltt}
\glsni{author}\marg{Me}
\glsni{title}\marg{A Simple Document}
\glsni{maketitle}

\glsni{pagenumbering}\marg{roman}
\glsni{tableofcontents}

\glsni{begin}\marg{abstract}
This is the abstract.
\glsni{end}\marg{abstract}

\glsni{pagenumbering}\marg{arabic}
\glsni{chapter}\marg{Introduction}
\end{alltt}
\end{code}

The \icls{scrbook} class provides:
\begin{definition}
\gls{frontmatter}
\end{definition}
which switches to lower case Roman numeral page numbering, and
\begin{definition}
\gls{mainmatter}
\end{definition}
which switches to Arabic numeral page numbering.
These two declarations also
change the way the sectioning units, such as \gls{chapter} and
\gls{section}, appear.  The former, \glsni{frontmatter},
suppresses the numbering (regardless of whether or not you've used
the \htmlref{starred version}{itm:starredcommand} of the sectioning commands). The latter,
\glsni{mainmatter}, switches the numbering back on (unless
otherwise suppressed by using the starred sectioning commands). 
In addition, \icls{scrbook} provides
\begin{definition}
\gls{backmatter}
\end{definition}
which doesn't affect the page numbering but, like \glsni{frontmatter}, suppresses the 
sectional unit numbering.

\xminisec{Note:}

The \gls{env-abstract} environment isn't defined by the \icls{scrbook} class, as 
a book summary is usually incorporated into an introductory section.

\xminisec{Example:}

\begin{code}
\begin{alltt}
\glsni{documentclass}\oarg{12pt}\marg{scrbook}

\glsni{title}\marg{A Simple Document}
\glsni{author}\marg{Me}

\glsni{begin}\marg{document}
\glsni{maketitle}

\glsni{frontmatter}
\glsni{tableofcontents}

\glsni{chapter}\marg{Summary}
A brief document to
illustrate how to use \glsni{LaTeX}.

\glsni{mainmatter}

\glsni{chapter}\marg{Introduction}
\glsni{label}\marg{ch:intro}

\glsni{end}\marg{document}
\end{alltt}
\end{code}

The headers and footers can be changed using the command
\begin{definition}
\gls{pagestyle}\marg{\meta{style}}
\end{definition}%
Individual pages can be changed using
\begin{definition}
\gls{thispagestyle}\marg{\meta{style}}
\end{definition}%
Standard styles are:
\begin{fwlist}{myheadings}
\fwitem{\ipagestyle{empty}} No header or footer. 
\fwitem{\ipagestyle{plain}} Header empty, page number in footer.
\fwitem{\ipagestyle{headings}}  Header contains page number and various information, 
 footer empty. 
\fwitem{\ipagestyle{myheadings}} Header specified by user, footer empty.
\end{fwlist}

If the \ipagestyle{myheadings} style is used, the header information
can be specified using:
\begin{definition}
\gls{markboth}\marg{\meta{left head}}\marg{\meta{right head}}
\end{definition}%
if the \clsopt{twoside} option has been passed to the
\gls{cls} (default for \icls{scrbook}), or
\begin{definition}
\gls{markright}\marg{\meta{right head}}
\end{definition}%
if the \clsopt{oneside} option has been passed to the \gls{cls}
(default for \icls{scrartcl} and \icls{scrreprt}).

The \icls{scrreprt} class file uses the \pagestylefmt{empty} style
for the \Index{title} and \Index{abstract} pages and
\pagestylefmt{plain} for the first page of each new chapter.  By
default the remaining pages are also \pagestylefmt{plain}, but
these can be changed using the \glsni{pagestyle} command.
The \icls{scrbook} class defaults to the \pagestylefmt{headings}
style instead of \pagestylefmt{plain}.

\faq{Alternative head- and footlines in LaTeX}{fancyhdr}%
The \koma\ bundle provides a way to define new page
styles, but that's beyond the scope of this introductory tutorial. 
See the \koma\ documentation for further details if you are
interested.

\casemedia
 {% A4
  This book mostly uses the \pagestylefmt{headings} page style and the
  \icls{scrbook} class with the 
  \clsopt{oneside} option, so there is no difference between odd and
  even page headers, whereas the paperback version uses the
  \clsopt{twoside} option, so the odd pages display the chapter
   number and title and the even pages display the current section
   header and title. The on-screen PDF version of this book uses a page style 
  I~defined myself that incorporates a navigation bar in the footer.
 }%
 {% screen
   This screen version of the book uses a page style I~defined myself
   that incorporates a navigation bar in the footer. 
   The \htmladdnormallink{A4 version}{\baseurl/latex/novices/novices-a4.pdf} mostly uses
   the \pagestylefmt{headings} page style. If you look at it, you
   will see that the chapter number and title appear on the top left and
   the page number appears in the top right of most pages.  The 
   \clsopt{oneside} option was used, so there is no difference between
   the formatting of odd and even numbered pages. Whereas the paperback
   version uses the \clsopt{twoside} option so the odd and even page
   headers are different.
 }
 {% book
  This book mostly\footnote{I made some modifications to the page style 
  for the footers in the summary and index.}\ uses the \pagestylefmt{headings} 
  page style and the \icls{scrbook} class with the 
  \clsopt{twoside} option, so the odd and even page headers are
  different, whereas the A4 PDF version (available from
  \url{\baseurl/}) uses the \clsopt{oneside} option, so the odd and
  even page headers are the same.
  The on-screen PDF version of this book uses a page style 
  I~defined myself that incorporates a navigation bar in the footer.
 }
 {% html
   The \htmladdnormallink{A4
   version}{\baseurl/latex/novices/novices-a4.pdf} of this book
   mostly uses the \pagestylefmt{headings} page style.  If you look at it, you
   will see that the chapter number and title appear on the top left and
   the page number appears in the top right of most pages.  The 
   \clsopt{oneside} option was used, so there is no difference between
   the formatting of odd and even numbered pages. Whereas the paperback
   version uses the \clsopt{twoside} option, so the odd pages display the chapter
   number and title and the even pages display the current section
   header and title.
   The \htmladdnormallink{on-screen PDF
   version}{\baseurl/latex/novices/novices-screen.pdf} of this book
   uses a page style I~defined myself that incorporates a navigation bar
   in the footer. 
 }

\begin{exercise}{Page Styles and Page Numbering}{ex:pagestyle}

Try modifying your code so that it uses the \icls{scrbook} class,
\gls{frontmatter} and \gls{mainmatter}. Replace the
\gls{env-abstract} environment with an unnumbered chapter, as shown
below. Again, changes made from the previous document are illustrated
\modification{like this}:\screenpagebreak
\begin{bcode}
\begin{alltt}
\glsni{documentclass}\oarg{12pt}\marg{\modification{scrbook}}

\glsni{usepackage}\marg{datetime}

\modification{\glsni{pagestyle}\marg{headings}}

\glsni{title}\marg{A Simple Document}
\glsni{author}\marg{Me}

\glsni{begin}\marg{document}

\glsni{maketitle}

\modification{\glsni{frontmatter}}
\glsni{tableofcontents}

\modification{\glsni{chapter}\marg{Summary}}
A brief document to
illustrate how to use \glsni{LaTeX}.

\modification{\glsni{mainmatter}}

\glsni{chapter}\marg{Introduction}
\glsni{label}\marg{ch:intro}

\glsni{section}\marg{The First Section}

This is a simple \glsni{LaTeX}\glsni{spacesym}document.
Here is the first paragraph.
The next chapter is Chapter\glsni{tildechar}\glsni{ref}\marg{ch:another}
and is on page\glsni{tildechar}\glsni{pageref}\marg{ch:another}.
The next section is Section\glsni{tildechar}\glsni{ref}\marg{sec:next}.

\glsni{percentchar} Rest of document unchanged but 
\glsni{percentchar} omitted for brevity.
\glsni{end}\marg{document}
\end{alltt}
\end{bcode}
(You can \downloadorview{pagestyle} the edited
document.)

\end{exercise}

%%%%%%%%%%%%%%%%%%%%%%% babel %%%%%%%%%%%%%%%%%%%%%%%%%

\setnode{babel}
\section{Multi-Lingual Support: using the 
\texorpdfstring{\sty{babel}}{babel} package}
\label{sec:babel}

You may have noticed that the \gls{tableofcontents} and
\gls{chapter} commands have produced English words like
\dq{Contents} and \dq{Chapter}\faq{How to change LaTeX's \dq{fixed
names}}{fixnam}. If you are writing in another language, this is not
appropriate. In this case, you can use the \isty{babel} package,
and specify which language you will be using, either as an option to
the \isty{babel} package, or as an option to the class file\faq{Using
a new language with Babel}{newlang}.  If you are writing in more than
one language, list all the languages that you will be using where
the last named language is the default language\faq{Parallel
setting of text}{parallel}. For example:
\begin{alltt}
\glsni{usepackage}\oarg{english,french}\marg{babel}
\end{alltt}
or
\begin{alltt}
\glsni{documentclass}\oarg{english,french}\marg{scrreprt}
\glsni{usepackage}\marg{babel}
\end{alltt}
You can then switch between the named languages either using the 
\gls{declaration}:
\begin{definition}
\gls{selectlanguage}\marg{\meta{language}}
\end{definition}%
or the \gls{env-otherlanguage} \glsterm{environment}:
\begin{definition}
\glsni{begin}\marg{otherlanguage}\marg{\meta{language}}
\end{definition}%
These will affect all translations, including the date format
and predefined names like \dq{Chapter}. This also changes the
\gls{hyphenation} patterns.\latex{ (See \sectionref{sec:hyphenation}.)}

If you only want to set a short section of text in a different
language, without affecting the date format or predefined names,
then you can either use the command:
\begin{definition}
\gls{foreignlanguage}\marg{\meta{language}}\marg{\meta{text}}
\end{definition}%
or the starred version of the \glsni{env-otherlanguage} environment:
\begin{definition}
\glsni{begin}\marg{otherlanguage*}\marg{\meta{language}}
\end{definition}%
You can test to see if a given language is currently selected
using:
\begin{definition}
\gls{iflanguage}\marg{\meta{language}}\marg{\meta{true text}}%
\marg{\meta{false text}}
\end{definition}%

\xminisec{Example:}
\begin{code}
\begin{alltt}
\glsni{documentclass}\oarg{UKenglish,USenglish,french}\marg{scrartcl}
\glsni{percentchar} french is the last named option, so that's the current language

\glsni{usepackage}\oarg{T1}\marg{fontenc}
\glsni{usepackage}\oarg{utf8}\marg{inputenc}
\glsni{usepackage}\marg{babel}

\glsni{begin}\marg{document}
Ce texte est en fran\glsni{c}\marg{c}ais. La date aujourd'hui est: \glsni{today}.

\glsni{selectlanguage}\marg{USenglish}
This text is in US English. Today's date is: \glsni{today}.

\glsni{selectlanguage}\marg{UKenglish}
This text is in UK English. Today's date is: \glsni{today}.
\glsni{end}\marg{document}
\end{alltt}
\end{code}
Result:
\begin{result}[babel.html]
Ce texte est en fran\c{c}ais. La date aujourd'hui est~: 
{\datefrench\docdate}.

This text is in US English. Today's date is: {\usdate\docdate}.

This text is in UK English. Today's date is: {\ukenglishdate\docdate}.
\end{result}

\xminisec{Note:}
\warning If you are using the \istyopt{babel}{french} option, the
colon character (\Indextt{:}) is made active (that is, it's turned
into a special character) so if you are writing in French it's best
not to use a colon in labels (so where I've used, say,
\texttt{ch:def} you may need to change the colon to something else).

%%%%%%%%%%%%%%%%%%%%%%%% graphicx %%%%%%%%%%%%%%%%%%%

\setnode{graphicx}
\chapter{The \texorpdfstring{\sty{graphicx}}{graphicx} Package}
\label{ch:graphicx}

It is possible to generate images using \LaTeX\ commands\faq{Drawing
with TeX}{drawing} (See the \isty{pgf}\slash\isty{tikz} package or
\latexgraphic) however most people find it easier to create a
picture in some other application, and include that file into their
\LaTeX\ document.

PDF\LaTeX\ can insert PDF, PNG and JPG image files into your
document. If your image file is in a different format, you may be
able to find an application to convert it. \faq{Spawning programs
from (La)TeX: \cmdname{write}18}{spawnprog}Modern
\gls{tex}-distributions can automatically convert EPS\indexEPS\
files to PDF during the \LaTeX\ run using the \gls{perl} script
\iappname{epstopdf}. If your \gls{tex}-distribution doesn't support
this, you can convert your EPS file using \appname{epstopdf}
explicitly. For example, if you have an EPS image called, say,
\texttt{sample-image.eps}, you can convert it to a PDF image called
\texttt{sample-image.pdf}, by using the following command in a
\glslink{terminal}{terminal or command prompt}:
\begin{verbatim}
epstopdf sample-image.eps
\end{verbatim}
or (full path name may be required)
\begin{verbatim}
perl epstopdf sample-image.eps
\end{verbatim}

To insert an image file into your document, you first need to specify that 
you want to use the \isty{graphicx} package. So the following must 
go in the \gls{preamble}:
\begin{codeS}
\glsni{usepackage}\marg{graphicx}
\end{codeS}
The image can then be included in your document using the command
\begin{definition}
\gls{includegraphics}\oarg{\meta{key-val options}}\marg{\meta{filename}}
\end{definition}
where \meta{filename} is the name of your image file \emph{without the file
extension}\faq{\dq{Modern} graphics file names}{grffilenames}, and
\meta{key-val options} is a comma-separated list of options that can be
used to change the way the image is displayed. Note that if you have an 
image where the file name contain spaces or multiple dots, you need to
use the \isty{grffile} package:
\begin{codeS}
\glsni{usepackage}\marg{graphicx,grffile}
\end{codeS}

\xminisec{Example:} 

Suppose you had a file called \texttt{shapes.pdf}, then to
include it in your document you would do:
\begin{codeS}
\glsni{includegraphics}\marg{shapes}
\end{codeS}%
which would produce:
\begin{resultS*}
\incGraphics[alt=Image of some shapes]{exercises/shapes}
\end{resultS*}

You can specify a full or relative pathname, but you must use a
forward slash \gls{dir.slash} as the directory divider, even if you are using
Windows. For example:
\begin{codeS}
\glsni{includegraphics}\marg{pictures\glsni{dir.slash}shapes}
\end{codeS}
means the file \texttt{pictures/shapes.pdf} on Unix-type systems,
and it means the file \texttt{pictures\backslashsym shapes.pdf} on
Windows.\footnote{Or \texttt{shapes.png} or \texttt{shapes.jpg} or
\texttt{shapes.eps}. The example assumes a PDF image file.} This is
mainly because the backslash character is a \LaTeX\ special
character indicating a command, but it also helps portability
between platforms.

You can specify the order of the file types to look for with the command
\begin{definition}
\gls{DeclareGraphicsExtensions}\marg{\meta{ext-list}}
\end{definition}%
where \meta{ext-list} is a comma-separated list of extensions.  
For example, 
you might want to search first for PDF files, then for PNG files,
then for JPG files and finally for EPS files:
\begin{codeS}
\glsni{DeclareGraphicsExtensions}\marg{.pdf,.png,.jpg,.eps}
\end{codeS}%
The default for \iPDFLaTeX\ is:
\begin{flushleft}\ttfamily
.png\comma .pdf\comma .jpg\comma .mps\comma .jpeg\comma .jbig2\comma .jb2\comma .PNG\comma .PDF\comma .JPG\comma .JPEG\comma .JBIG2\comma .JB2\comma .eps
\end{flushleft}

The \gls{optional} \meta{key-val options}
should be a comma-separated list of
\meta{key}\texttt{=}\meta{value} pairs. Common options are:

\begin{fwlist}{\ikeyvalopt{includegraphics}{trim}=\meta{l} \meta{b} \meta{r} \meta{t}}
\fwitem{\ikeyvalopt{includegraphics}{angle}=\meta{\mathorit{x}}} rotate the image
by \degrees{x} anticlockwise.

\fwitem{\ikeyvalopt{includegraphics}{width}=\meta{length}} scale the image so that the
width is \meta{length}.  (Remember to specify the
\htmlref{units}{tab:units}.)

\fwitem{\ikeyvalopt{includegraphics}{height}=\meta{length}}  scale the image so that the
height is \meta{length}.  (Remember to specify the
\htmlref{units}{tab:units}.)

\fwitem{\ikeyvalopt{includegraphics}{scale}=\meta{value}} Scale the image by \meta{value}

\fwitem{\ikeyvalopt{includegraphics}{trim}=\meta{l} \meta{b} \meta{r} \meta{t}}  Specifies
the amount to remove from each side.  For example
\begin{codeS}
 \glsni{includegraphics}\oarg{trim=1 2 3 4}\marg{shapes}
\end{codeS}
crops the image by 1bp from the left, 2bp from the bottom, 3bp from
the right and 4bp from the top. (Recall the \iunit{bp}{big point} unit
from \tableref{tab:units}.)\reportpagebreak

\fwitem{\ikeyvalopt{includegraphics}{draft}}  Don't actually print the image, just draw a
box of the same size and print the filename inside it. 

\end{fwlist}

\xminisec{Example:}

This example first rotates the image by \degrees{45} anticlockwise,
then scales it so that the width is 1~inch.

\begin{codeS}
\glsni{includegraphics}\oarg{angle=45,width=1in}\marg{shapes}
\end{codeS}

\begin{resultS}[Image of shapes rotated then scaled]
\includegraphics[angle=45,width=1in]{exercises/shapes}
\end{resultS}

Note that this isn't the same as scaling and then rotating:
\begin{codeS}
\glsni{includegraphics}\oarg{width=1in,angle=45}\marg{shapes}
\end{codeS}

\begin{resultS}[Image of shapes scaled then rotated]
\includegraphics[width=1in,angle=45]{exercises/shapes}
\end{resultS}

You can also scale an image relative to the text area using the
\gls{length} registers \gls{textwidth} and \gls{textheight}. For
example, to scale a portrait image so that its height is
three-quarters of the text area height, you can do:
\begin{codeS}
\glsni{includegraphics}\oarg{height=0.75\glsni{textheight}}\marg{shapes}
\end{codeS}
or to scale a landscape image so that its width is half the text
area width, you can do:
\begin{codeS}
\glsni{includegraphics}\oarg{height=0.5\glsni{textwidth}}\marg{shapes}
\end{codeS}

\xminisec{Note:}

The \gls{includegraphics} command is another form of box (see
\sectionref{sec:boxes}), and can be used in the middle of a line of
text, just like the \gls{env-tabular} environment. See
\sectionref{sec:figures} to find out how to put the image in a
figure with a caption.

\xminisec{Example:}

\refstepcounter{object}\label{example:inlinegraphic}%
Recall the \iunit{ex}{relative unit} unit of measure from
\tableref{tab:units}. This can be used to scale
an image relative to the font size:
\begin{code}
An image can be inserted into a line
of text like this:
\glsni{includegraphics}\oarg{height=2ex}\marg{shapes}
\end{code}
\begin{resultS}[Image of typeset output: the included image has been
scaled so that it fits into the line of text]
An image can be inserted into a line
of text like this:
\includegraphics[height=2ex]{exercises/shapes}
\end{resultS}

\setnode{graphictrans}
\section{Graphical Transformations}
\label{sec:graphtrans}

The \isty{graphicx} package also provides commands to rotate, resize, 
reflect and scale text.  They are as follows:
\begin{definition}
\gls{rotatebox}\oarg{\meta{option list}}\marg{\meta{angle}}\marg{\meta{text}}
\end{definition}
Rotates \meta{text} by \meta{angle} (degrees anti-clockwise by
default). The optional argument \meta{option list} is a
comma-separated list of any of the following options:

\begin{itemize}
  \item \ikeyvalopt{rotatebox}{units}=\meta{number}

    The number of units in one full anti-clockwise rotation. So
    \texttt{units=-360} means that \meta{angle} specifies degrees
    clockwise whereas \texttt{units=6.283185} means that
    \meta{angle} is in radians.

  \item \ikeyvalopt{rotatebox}{origin}=\meta{label}

    The point of rotation.
    The value \meta{label} may contain one from either or both of
    the two lists: \texttt{lrc} (left, right, centre) and \texttt{tbB}
    (top, bottom, baseline). Alternatively the origin may be
    specified using the following two keys:

  \item \ikeyvalopt{rotatebox}{x}=\meta{dimen}

  \item \ikeyvalopt{rotatebox}{y}=\meta{dimen}
\end{itemize}

\xminisec{Example:}
\begin{code}
\begin{alltt}
base line
\glsni{rotatebox}\marg{45}\marg{Some text}
\glsni{rotatebox}\oarg{units=-360}\marg{45}\marg{Some text}
\glsni{rotatebox}\oarg{units=-360,origin=rB}\marg{45}\marg{Some text}
\glsni{rotatebox}\oarg{x=3em,y=3em}\marg{45}\marg{Some text}
base line
\end{alltt}
\end{code}\screenpagebreak
\begin{result}[Image of typeset output: each occurrence of 'Some text' has been rotated]
base line
\rotatebox{45}{Some text}
\rotatebox[units=-360]{45}{Some text}
\rotatebox[units=-360,origin=rB]{45}{Some text}
\rotatebox[x=3em,y=3em]{45}{Some text}
base line
\end{result}

\begin{definition}
\gls{scalebox}\marg{\meta{h scale}}\oarg{\meta{v scale}}\marg{\meta{text}}
\end{definition}
Scales \meta{text} by \meta{h scale} in both directions if \meta{v
scale} omitted, otherwise scales \meta{text} by \meta{h scale}
horizontally and \meta{y scale} vertically.

\xminisec{Example:}
\begin{codeS}
\glsni{scalebox}\marg{0.8}\marg{Some text}
\end{codeS}
\begin{resultS}[Some text (scaled by a factor of 0.8)]
\scalebox{0.8}{Some text}
\end{resultS}
Compare with:
\begin{codeS}
\glsni{scalebox}\marg{0.8}\oarg{1.2}\marg{Some text}
\end{codeS}
\begin{resultS}[Some text (aspect ratio changed)]
\scalebox{0.8}[1.2]{Some text}
\end{resultS}

\begin{definition}
\gls{reflectbox}\marg{\meta{text}}
\end{definition}
Reflects \meta{text} (equivalent to \glsni{scalebox}\marg{-1}\oarg{1}\marg{\meta{text}}).

\xminisec{Example:}
\begin{codeS}
\glsni{reflectbox}\marg{Some text}
\end{codeS}
\begin{resultS}[Some text (mirror image)]
\reflectbox{Some text}
\end{resultS}

\begin{definition}
\gls{resizebox}\marg{\meta{h length}}\marg{\meta{v length}}\marg{\meta{text}}
\end{definition}
Scales \meta{text} so that it is \meta{h length} wide and \meta{v
length} high. To preserve the aspect ratio, use \gls{resizebox.exclamchar}
instead of one of the dimensions.

\xminisec{Example:}
\begin{code}
\begin{alltt}
\glsni{resizebox}\marg{12mm}\marg{1cm}\marg{Some text}
\glsni{resizebox}\marg{\glsni{resizebox.exclamchar}}\marg{1cm}\marg{Some text}
\end{alltt}
\end{code}%
\begin{resultS}[Image of typeset output: two instances of 'Some text' scaled with and without changing the aspect ratio.]
\dohtmlcolorbox{%
\resizebox{12mm}{1cm}{Some text}
\resizebox{!}{1cm}{Some text}
}
\end{resultS}

\setnode{graphicxoptions}
\section{Package Options}

The \isty{graphicx} package can have the following options passed to 
it:
\begin{fwlist}{\optfmt{hiderotate}}
\fwitem{\mdseries \istyopt{graphicx}{draft}} Don't actually display the
images, just print the filename in a box of the correct size.  This
is useful if you want to print out a draft copy of a document to
check the text rather than the images.

\fwitem{\mdseries \istyopt{graphicx}{final}} Opposite of 
\istyopt{graphicx}{draft} (default).

\fwitem{\mdseries \istyopt{graphicx}{hiderotate}} Don't show rotated
text.

\fwitem{\mdseries \istyopt{graphicx}{hidescale}} Don't show scaled text.

\end{fwlist}

Remember that relevant options passed to the class file\indexCLS\ also 
affect packages.

\xminisec{Example (Draft Mode):}

Draft mode helps to speed up compilation of a large document when
you are editing the text. In the preamble:
\begin{codeS}
\glsni{usepackage}\oarg{draft}\marg{graphicx}
\end{codeS}
or
\begin{code}
\begin{alltt}
\glsni{documentclass}\oarg{draft}\marg{scrbook}

\glsni{usepackage}\marg{graphicx}
\end{alltt}
\end{code}
Later in the document:
\begin{codeS}
\glsni{includegraphics}\oarg{width=1in,angle=45}\marg{pictures/shapes}
\end{codeS}
\begin{resultS}[Only bounding box and file name drawn]
 \incGraphics{pictures/draftimage}
\end{resultS}

\begin{exercise}{Using the \sty{graphicx} Package}{ex:graphicx}

Download the image file
\htmladdnormallink{\texttt{shapes.pdf}}{exercises/shapes.pdf}%
\latex{ from \url{\exerciseurl/}} (or create your own image),
and include it into your document. (You
can \downloadorview{graphic} an example solution.)
\end{exercise}

\bookpagebreak\screenpagebreak
For more information on the \isty{graphicx} package see
\latexgraphic\ or the \isty{graphicx} documentation.

\xminisec{Related \gls{ukfaq} topics:}
\begin{itemize}
\item \faqlink{How to import graphics into (La)TeX documents}{impgraph}
\item \faqlink{Imported graphics in PDFLaTeX}{pdftexgraphics}
\item \faqlink{Imported graphics in dvips}{dvipsgraphics}
\item \faqlink{Imported graphics in dvipdfm}{dvipdfmgraphics}
\item \faqlink{Importing graphics from \dq{somewhere else}}{graphicspath}
\item \faqlink{Portable imported graphics}{graph-pspdf}
\item \faqlink{Repeated graphics in a document}{repeatgrf}
\item \faqlink{Limit the width of imported graphics}{grmaxwidth}
\item \faqlink{Top-aligning imported graphics}{topgraph}
\item \faqlink{Labelling graphics}{labelfig}
\item \faqlink{Graphics division by zero}{divzero}
\end{itemize}

%%%%%%%%%%%%%%%%%%%%%%%%%%% FLOATS %%%%%%%%%%%%%%%%%%%%%%%%%%%%%%%%%

\setnode{floats}
\chapter{Floats}
\label{ch:floats}

Figures and tables are referred to as \dq{floats} because they are
\emph{floated} to the nearest location. This prevents ugly large
spaces appearing on the page if there isn't enough room for the
figure or table before the page break. Floats have a caption and
associated number\faq{The style of captions}{captsty}. It is
customary for captions to appear at the bottom of figures but at the
top of tables~\cite{turabian96,oxford}.

For both figures and tables, the caption is generated using the 
command:
\begin{definition}
\gls{caption}\oarg{\meta{short caption}}\marg{\meta{text}}
\end{definition}%
Note that the \cmdname{caption} command has a
moving argument\index{argument!moving|hyperpage}, so
\glspl{fragile}\faq{Footnotes in captions}{ftncapt} will need to be
protected using \gls{protect}.  The \glsterm{optional}
\meta{short caption} is used to provide an alternative shorter
caption for the list of figures or list of tables, akin to the
optional argument to the \htmlref{sectioning
commands}{sec:section}\latex{\ifscreenorother{}{\ described in
\sectionref{sec:section}}}.

\xminisec{Note:}

Although the \glsni{caption} command can have an optional
short title, in general, captions should be brief. They should not
contain lots of description or background detail~\cite{turabian96}.
That type of information should be placed in the main text not the
caption.

\xminisec{Positioning:}

\refstepcounter{object}\label{obj:positioning}%
\faq{Wide figures in two-column documents}{widefigs}Both the 
\glsni{env-figure} and \glsni{env-table} environments have
an optional argument \meta{placement specifiers}, which indicates
permissible locations for the float. This may be a combination of
\texttt{h} (\dq{here}), \texttt{t} (top), \texttt{b} (bottom) and
\texttt{p} (page of floats.) Note that this only gives a general
guideline as to where the float will end up. The final location is
governed by other factors, such as space left on the page and the
proportion of text to floats on the page. If you omit one or more of the
placement specifiers, then you are prohibiting the float from being
placed in that location. A common mistake is to do
\begin{alltt}
\wrong\glsni{begin}\marg{figure}\oarg{h}
\end{alltt}
which says \dq{I~want the figure here and it can't go anywhere
else!} If the figure \emph{can't} be placed exactly here (for
example, there may not be enough room on the page), then you have
given it no alternative location, which can result in this and all
subsequent figures being dumped at the end of the chapter or
document, or can result in a fatal error when running
\LaTeX\faq{\dq{Too many unprocessed floats}}{tmupfl}. You may be
able to manage with only one of the other options, for example,
\begin{alltt}
\glsni{begin}\marg{figure}\oarg{t}
\end{alltt}
(In fact, modern \TeX\ distributions now replace \oarg{h} with
\oarg{t} if the float can't be placed.)
However, if you have a large number of floats it is advisable to
provide as many options as possible:
\begin{alltt}
\glsni{begin}\marg{figure}\oarg{htbp}
\end{alltt}
Similarly for tables.

If you are absolutely adamant that an image must go \dq{right here},
then it's not a float, and you shouldn't be using the 
\glsni{env-figure} environment. It's just a horizontal box, like the
example \xpageref{earlier}{example:inlinegraphic}. Similarly for
tabulated material.

It's worth bearing in mind what the Oxford Style Manual~\cite{oxford} has to say:
\begin{quote}
\dq{Text must not be read into it so as to give [the figure] an
explicit and fixed introduction, for example \sq{in the following
figure}: the final placement is determined by page breaks, which
cannot be anticipated before setting, and this makes rewording the
text necessary if the illustration does not fit the make-up of the
page.}
\end{quote}
Turabian~\cite{turabian96} gives the same advice (and reiterates it for 
figures):
\begin{quote}
\dq{All text references to a table should be by a number, not by an
introductory phrase such as \sq{in the following table}.}
\end{quote}

\setnode{figures}
\section{Figures}
\label{sec:figures}

Figures are created using the \gls{env-figure} environment.
\begin{definition}
\glsni{begin}\marg{figure}\oarg{\meta{placement specifiers}}
\end{definition}%
This environment may contain one or more captions (specified, as
described \htmlref{above}{ch:floats}, with
the \glsni{caption} command) but page breaks are not allowed in the
contents of a \glsni{env-figure} environment. The optional argument
\meta{placement specifiers} is as described
\htmlref{above}{obj:positioning}.

Recall from \chapterref{ch:graphicx} that we
can include an image in our document with
the command \gls{includegraphics} defined in the \isty{graphicx}
package.  We can put our
\htmladdnormallink{\texttt{shapes.pdf}}{exercises/shapes.pdf} image
into a figure as follows:
\begin{code}
\glsni{begin}\marg{figure}\oarg{htbp}\newline
\strut~~\glsni{includegraphics}\marg{shapes}\newline
\strut~~\glsni{caption}\marg{Some Shapes}\newline
\glsni{end}\marg{figure}
\end{code}%
So far so good, but our picture needs to be centred.  This can
be done using the \gls{centering} declaration mentioned in
\sectionref{sec:declaration}:
\begin{code}
\glsni{begin}\marg{figure}\oarg{htbp}\newline
\strut~~\glsni{centering}\newline
\strut~~\glsni{includegraphics}\marg{shapes}\newline
\strut~~\glsni{caption}\marg{Some Shapes}\newline
\glsni{end}\marg{figure}
\end{code}
The \gls{caption} command generates a number, just like
\gls{section}, so we can
\htmlref{cross-reference}{sec:crossref} it with \gls{ref} and
\gls{label}.  First, let's label the figure:
\begin{code}
\glsni{begin}\marg{figure}\oarg{htbp}\newline
\strut~~\glsni{centering}\newline
\strut~~\glsni{includegraphics}\marg{shapes}\newline
\strut~~\glsni{caption}\marg{Some Shapes}\newline
\strut~~\glsni{label}\marg{fig:shapes}\newline
\glsni{end}\marg{figure}
\end{code}%
Now we can reference it:
\begin{codeS}
Figure\glsni{tildechar}\glsni{ref}\marg{fig:shapes} shows some shapes.
\end{codeS}%
(As \htmlref{before}{cmd:tilde} we use \gls{tildechar}
to make an unbreakable space.)  This produces the following output
in the text:
\begin{resultS}[Figure 7.1 shows some shapes.]
Figure~\ref*{fig:shapes} shows some shapes.
\end{resultS}%
and produces \figureref{fig:shapes}.

\begin{figure}[hbtp]
\figconts
  {exercises/shapes}
  {%
    \caption{Some Shapes}%
  }%
  {fig:shapes}
\end{figure}

\xminisec{Important Note:}

If you want to change the caption font, \textbf{\em don't} do, e.g.:
\begin{alltt}
\wrong\glsni{caption}\marg{\gls{bfseries} Some Shapes}
\end{alltt}
Recall \gls{addtokomafont} from \sectionref{sec:section}. This can
also be used to change the fonts used by the caption.\bookpagebreak
\begin{alltt}
\correct\glsni{addtokomafont}\marg{caption}\marg{\glsni{bfseries}}
\end{alltt}
Similarly for the caption label. For example:
\begin{codeS}
\glsni{addtokomafont}\marg{captionlabel}\marg{\glsni{scshape}}
\end{codeS}

\xminisec{List of Figures}

Just as we were able to generate a \htmlref{table of
contents}{sec:toc} using \gls{tableofcontents}, we can also
generate a list of figures using the command
\begin{definition}
\gls{listoffigures}
\end{definition}
This creates a file with the extension \texttt{.lof}\indexLOF\ (see
\sectionref{sec:auxiliary}).
As with \reportlinebreak\glsni{tableofcontents} you will need to \LaTeX\ your document twice to
get the list of figures up-to-date, unless
you're using \iappname{latexmk} (as described in
\sectionref{sec:crossref}) in which case it will be done
automatically.

\begin{exercise}{Creating Figures}{ex:figures}

If you did \exerciseref{ex:graphicx}, you should have a document
with an image in it. You now need to put this image into a
\gls{env-figure} environment. Remember to centre the image, and give
the figure a caption. Next, try labelling the figure and referencing
it in the text. You could also put in a list of figures after the
table of contents.
You can \downloadorview{figures} an example.

\end{exercise}

\setnode{sidebysidefigs}
\subsection{Side-By-Side Figures}
\label{sec:sidebysidefigs}

Recall at the start of \sectionref{sec:figures}, I mentioned that
the \gls{env-figure} environment may contain one or more captions.
In most cases, you'll just have a single caption per
\glsni{env-figure} environment, but sometimes you may want to have
two figures side-by-side, in which case you'll need two captions
within the same \gls{env-figure} environment in order to keep the
figures together.

To do this, we can use the 
\gls{env-minipage} environment, which was covered 
in \sectionref{sec:boxes}. Recall that 
the \gls*{env-minipage} environment creates a 
\Index{horizontal box}, which means that two mini-pages can
be placed side-by-side on the same line. All you need to do 
now, is place one image and caption in one mini-page, and the
other image and caption in the neighbouring mini-page. (Do you remember what effect is obtained
by placing a \htmlref{percent symbol}{obj:suppresseol} at the end of a
line?)\doifbook{\footnote{See page~\pageref{obj:suppresseol}.}}\bookpagebreak
\begin{code}
\glsni{begin}\marg{figure}\oarg{htbp}\newline
\strut~\glsni{begin}\marg{minipage}\marg{0.5\glsni{linewidth}}\newline
\strut~~\glsni{centering}\newline
\strut~~\glsni{includegraphics}\marg{circle}\newline
\strut~~\glsni{caption}\marg{A Circle}\newline
\strut~~\glsni{label}\marg{fig:circle}\newline
\strut~\glsni{end}\marg{minipage}\glsni{percentchar}\newline
\strut~\glsni{begin}\marg{minipage}\marg{0.5\glsni{linewidth}}\newline
\strut~~\glsni{centering}\newline
\strut~~\glsni{includegraphics}\marg{rectangle}\newline
\strut~~\glsni{caption}\marg{A Rectangle}\newline
\strut~~\glsni{label}\marg{fig:rectangle}\newline
\strut~\glsni{end}\marg{minipage}\newline
\glsni{end}\marg{figure}
\end{code}%
\latex{The above code produces 
Figures~\ref{fig:circle-sbs} and~\vref{fig:rectangle-sbs}. }Note that each
mini-page uses \gls{centering} to centre its contents, and the
label is also placed in the same mini-page, after the
\gls{caption} command. If the \gls{label} was not in the same
\htmlref{scope}{sec:group} as the \gls{caption}, the
cross-reference would be incorrect.

A common mistake when trying to create side-by-side figures is
to do:
\begin{alltt}
\glsni{begin}\marg{figure}[htbp]
 \glsni{begin}\marg{minipage}\marg{0.5\glsni{linewidth}}
  \glsni{centering}
  \glsni{includegraphics}\marg{circle}
  \glsni{caption}\marg{A Circle}
  \glsni{label}\marg{fig:circle}
 \glsni{end}\marg{minipage}
\wrong
 \glsni{begin}\marg{minipage}\marg{0.5\glsni{linewidth}}
  \glsni{centering}
  \glsni{includegraphics}\marg{rectangle}
  \glsni{caption}\marg{A Rectangle}
  \glsni{label}\marg{fig:rectangle}
 \glsni{end}\marg{minipage}
\glsni{end}\marg{figure}
\end{alltt}
This produces one figure on top of the other, instead of
side-by-side\html{, as illustrated in \figureref{fig:circle-sbs} and
\figureref{fig:rectangle-sbs}}. This is because the blank
line indicates a paragraph break, so each minipage is in a separate
paragraph, so it's not possible for them to be on the same line.

If you want a bit of spacing in your code to make it more readable,
use \gls{percentchar} to comment out the paragraph break. For
example:
\begin{alltt}
\glsni{end}\marg{minipage}%
\correct\glsni{percentchar}
\glsni{begin}\marg{minipage}\marg{0.5\glsni{linewidth}}
\end{alltt}


\begin{figure}[htbp]
  \begin{makeimage}\end{makeimage}
  \begin{minipage}{0.5\linewidth}
    \figuretop{fig:circle-sbs}
    \centering
    \incGraphics[alt=Image of a circle]{exercises/circle}
    \caption{A Circle}
    \label{fig:circle-sbs}
  \end{minipage}%
  \begin{minipage}{0.5\linewidth}
    \figuretop{fig:rectangle-sbs}
    \centering
    \incGraphics[alt=Image of a rectangle]{exercises/rectangle}
    \caption{A Rectangle}
    \label{fig:rectangle-sbs}
  \end{minipage}
\end{figure}

%%%%%%%%%%%%%%%%%%%%%% Tables %%%%%%%%%%%%%%%%%%%%%%%%%

\setnode{tables}
\section{Tables}
\label{sec:tables}

Tables are produced in much the same way as figures, except that the
\gls{env-table} environment is used instead.
\begin{definition}
\glsni{begin}\marg{table}\oarg{\meta{placement specifiers}}
\end{definition}
Where the optional argument \meta{placement specifiers} is as described
\xpageref{earlier}{obj:positioning}.

As mentioned at the \htmlref{start of this chapter}{ch:floats},
tables typically have the caption at the top of the
table~\cite{turabian96}. With the \koma\ classes, such as
\icls{scrartcl}, \icls{scrreprt} and \icls{scrbook}, use the class
option \scrclsopt[tableabove]{captions} to ensure that the vertical
spacing appears correctly between the caption and the table content
and put \gls{caption} at the start of the \envname{table}
environment. \faq{Tables longer than a single page}{longtab}Page
breaks are not permitted in the \glsni{env-table}
environment. (The \isty{longtable} package can be used for that
instead. See the \sty{longtable} \htmlref{documentation}{sec:texdoc} for 
further details.)

\xminisec{Example:}

In the \glsterm{preamble}:
\begin{codeS}
\glsni{documentclass}\oarg{captions=tableabove}\marg{scrbook}
\end{codeS}
Later in the document:
\begin{code}
\begin{alltt}
\glsni{begin}\marg{table}\oarg{htbp}
 \glsni{caption}\marg{A Sample Table} 
 \glsni{label}\marg{tab:sample}
 \glsni{centering}
 \glsni{begin}\marg{tabular}\marg{lr}
 Item  \glsni{ampchar} Cost\glsni{tab.dbbackslashchar}
 Video \glsni{ampchar} 8.99\glsni{tab.dbbackslashchar}
 CD    \glsni{ampchar} 9.99\glsni{tab.dbbackslashchar}
 DVD   \glsni{ampchar} 15.00
 \glsni{end}\marg{tabular}
\glsni{end}\marg{table}
\end{alltt}
\end{code}%
This produces \tableref{tab:sample}.

\begin{table}[htbp]
\caption{A Sample Table}
\label{tab:sample}
\centering
\begin{makeimage}
\begin{tabular}{lr}
Item & Cost\\
Video & 8.99\\
CD    & 9.99\\
DVD   & 15.00
\end{tabular}
\end{makeimage}
\end{table}

Again, the \gls{centering} declaration is used to centre the
\gls{env-tabular} environment.
As with figures, you can create a list of tables using the command
\begin{definition}
\gls{listoftables}
\end{definition}
This creates a file with the extension \texttt{.lot}\indexLOT\ (see
\sectionref{sec:auxiliary}).
As with the table of contents and list of figures, you will need to
\LaTeX\ your document twice to get the list of tables up-to-date, unless
you're using \iappname{latexmk} (as described in
\sectionref{sec:crossref}) in which case it will be done
automatically.

\begin{exercise}{Creating Tables}{ex:tables}
If you did \exerciseref{ex:tabular}, you should have 
a \gls{env-tabular} environment in your document.  Try
turning this into a table, and add \tableref{tab:sample}.  
You could also try adding a list of tables.
As before, you can \downloadorview{tables} the solution.
\end{exercise}

\setnode{sidebysidetab}
\subsection{Side-by-Side Tables}
\label{sec:sidebysidetab}

You can create side-by-side tables using an analogous method to 
the side-by-side figures approach described
\htmlref{above}{sec:sidebysidefigs}.

\xminisec{Example:}

This example is similar to the one in \sectionref{sec:sidebysidefigs}.
Again, take care to ensure that there is no paragraph break between the
two \gls{env-minipage} environments.
\begin{code}
\begin{alltt}
\glsni{begin}\marg{table}
 \glsni{begin}\marg{minipage}\marg{0.5\glsni{linewidth}}
  \glsni{caption}\marg{Prices for 2011}
  \glsni{label}\marg{tab:prices2011}
  \glsni{centering}
  \glsni{begin}\marg{tabular}\marg{lr}
   Item \glsni{ampchar} Price (\pounds)\glsni{tab.dbbackslashchar}
   Widgets \glsni{ampchar} 10.99\glsni{tab.dbbackslashchar}
   Whatsits \glsni{ampchar} 5.99
  \glsni{end}\marg{tabular}
 \glsni{end}\marg{minipage}\glsni{percentchar}
\glsni{percentchar}
 \glsni{begin}\marg{minipage}\marg{0.5\glsni{linewidth}}
  \glsni{caption}\marg{Prices for 2012}
  \glsni{label}\marg{tab:prices2012}
  \glsni{centering}
  \glsni{begin}\marg{tabular}\marg{lr}
   Item \glsni{ampchar} Price (\glsni{pounds})\glsni{tab.dbbackslashchar}
   Widgets \glsni{ampchar} 11.99\glsni{tab.dbbackslashchar}
   Whatsits \glsni{ampchar} 6.99
  \glsni{end}\marg{tabular}
 \glsni{end}\marg{minipage}\glsni{percentchar}
\glsni{end}\marg{table}
\end{alltt}
\end{code}
This produces Tables~\ref{tab:prices2011} and~\vref{tab:prices2012}.

\begin{table}[htbp]
 \begin{minipage}{0.5\linewidth}
  \caption{Prices for 2011}
  \label{tab:prices2011}
  \centering
  \begin{makeimage}
  \begin{tabular}{lr}
   Item & Price (\pounds)\\
   Widgets & 10.99\\
   Whatsits & 5.99
  \end{tabular}
  \end{makeimage}
 \end{minipage}%
%
 \begin{minipage}{0.5\linewidth}
  \caption{Prices for 2012}
  \label{tab:prices2012}
  \centering
  \begin{makeimage}
  \begin{tabular}{lr}
   Item & Price (\pounds)\\
   Widgets & 11.99\\
   Whatsits & 6.99
  \end{tabular}
  \end{makeimage}
 \end{minipage}%
\end{table}

\setnode{sidewaysfloats}
\section{Sideways Floats}
\label{sec:sidewaysfloats}

The \isty{rotating} package provides the \gls{env-sidewaysfigure}
environment:
\begin{definition}
\glsni{begin}\marg{sidewaysfigure}
\end{definition}
and the \gls{env-sidewaystable} environment:
\begin{definition}
\glsni{begin}\marg{sidewaystable}
\end{definition}
which are like \gls{env-figure} and \gls{env-table}, respectively, but rotate the
entire float (including caption) sideways. This sideways float is
always placed on a page of its own.

If you have used the \clsopt{twoside} class option (or you are
using a class like \icls{scrbook}, which defaults to that option)
then the sideways floats will be rotated clockwise or
anti-clockwise, depending on whether they fall on an even (verso) or
odd (recto) numbered page. (Requires a second \LaTeX\ run to get it
correct.) 

\xminisec{Example:}
\begin{code}
\begin{alltt}
\glsni{begin}\marg{sidewaysfigure}
 \glsni{centering}
 \glsni{includegraphics}\oarg{width=0.75\glsni{textheight}}\marg{shapes}
 \glsni{caption}\marg{A Sideways Figure}
 \glsni{label}\marg{fig:sideways}
\glsni{end}\marg{sidewaysfigure}
\end{alltt}
\end{code}
The above code produces \figureref{fig:sideways}.

\dosidewaysfigure
 {fig:sideways}
 {\includegraphics[width=0.75\textheight]{exercises/shapes}}
 {A Sideways Figure}

\begin{htmlonly}
  \begin{figure}
    \figuretop{fig:sideways}
    \centering
\makeimg{A sideways figure}{%
\colorbox{white}{%
    \rotatebox[br]{90}{%
      \begin{tabular}{c}
      \includegraphics{exercises/shapes}\\
      {\usekomafont{captionlabel}Figure~\ref*{fig:sideways}}\captionformat A Sideways Figure
      \end{tabular}
    }%
}%
}
    \caption{A Sideways Figure}
    \label{fig:sideways}
  \end{figure}
\end{htmlonly}

\setnode{subfloats}
\section{Sub-Floats}
\label{sec:subfloats}

Some floats have sub-floats within them. For example, a figure may
contain several sub-figures, each of which requires a caption.
The simplest way to do this is to use the \isty{subcaption} package
that provides the \gls{env-subfigure} and \gls{env-subtable} environments:
\begin{definition}
\glsni{begin}\marg{subfigure}\oarg{\meta{pos}}\marg{\meta{width}}
\end{definition}
\begin{definition}
\glsni{begin}\marg{subtable}\oarg{\meta{pos}}\marg{\meta{width}}
\end{definition}
Within these environments, you can use \gls{caption} to create a
sub-caption. (In addition to the main \gls{caption} for the
containing \gls{env-figure} or \gls{env-table}
environment.)

\xminisec{Note:}
The \isty{subcaption} package requires the \isty{caption} package,
but doesn't automatically load it, so you'll need to load
both:\reportpagebreak
\begin{codeS}
\glsni{usepackage}\marg{caption,subcaption}
\end{codeS}

\xminisec{Example:}

This is very similar to the side-by-side figures example from
\sectionref{sec:sidebysidefigs}.
 
\begin{code}[0.9\linewidth]
\glsni{begin}\marg{figure}\oarg{htbp}\newline
\strut~\glsni{begin}\marg{subfigure}\oarg{b}\marg{0.5\glsni{linewidth}}\newline
\strut~~\glsni{centering}\newline
\strut~~\glsni{includegraphics}\marg{rectangle}\newline
\strut~~\glsni{caption}\marg{Rectangle}\glsni{label}\marg{fig:rectangle}\newline
\strut~\glsni{end}\marg{subfigure}\glsni{percentchar}\newline
\glsni{percentchar}\newline
\strut~\glsni{begin}\marg{subfigure}\oarg{b}\marg{0.5\glsni{linewidth}}\newline
\strut~~\glsni{centering}\newline
\strut~~\glsni{includegraphics}\marg{circle}\newline
\strut~~\glsni{caption}\marg{Circle}\glsni{label}\marg{fig:circle}\newline
\strut~\glsni{end}\marg{subfigure}\glsni{percentchar}\newline
\glsni{caption}\marg{Two Shapes}\newline
\glsni{label}\marg{fig:shape}\newline
\glsni{end}\marg{figure}
\end{code}

This produces \figureref{fig:shapes2}.
Elsewhere in the document, the figure and its components can be
referenced:
\begin{code}
Figure\gls{tildechar}\gls{ref}\marg{fig:shapes2} shows some shapes.  
Figure\gls{tildechar}\gls{ref}\marg{fig:rectangle} shows a rectangle and 
Figure\gls{tildechar}\gls{ref}\marg{fig:circle} shows a circle.
\end{code}
which produces the following text:
\begin{result}[Figure 7.5 shows some shapes. Figure 7.5a shows a
rectangle and Figure 7.5b shows a circle.]
Figure~\ref*{fig:shapes2} shows some shapes.
Figure~\ref*{fig:rectangle} shows a rectangle and
Figure~\ref*{fig:circle} shows a circle.\relax
\end{result}

You can also reference just the sub-float using
\begin{definition}
\gls{subref}\marg{\meta{label}}
\end{definition}
which is analogous to \gls{ref}, but only displays the sub-float
number without the number associated with its containing float.

\xminisec{Example:}
\begin{code}
Figure\glsni{tildechar}\cmdname{ref}\marg{fig:shapes2}
shows: (\glsni{subref}\marg{fig:rectangle}) a rectangle and
(\glsni{subref}\marg{fig:circle}) a circle.
\end{code}
produces
\begin{resultS}[Figure 7.5 shows: (a) a rectangle and (b) a circle.]
Figure~\ref*{fig:shapes2} shows: (\subref*{fig:rectangle}) a rectangle and (\subref*{fig:circle}) a circle.
\end{resultS}

\begin{figure}[htbp]
\begin{makeimage}\end{makeimage}
\figuretop{fig:shapes2}
\centering
\begin{latexonly}
 \begin{subfigure}[b]{0.5\linewidth}
   \centering
   \includegraphics{exercises/rectangle}
   \caption{Rectangle}\label{fig:rectangle}
 \end{subfigure}%
 \begin{subfigure}[b]{0.5\linewidth}
   \centering
   \includegraphics{exercises/circle}
   \caption{Circle}\label{fig:circle}
 \end{subfigure}%
\end{latexonly}
\begin{htmlonly}
\begin{tabular}{cc}
\htmladdimg{exercises/rectangle.png} \hspace{0.8in}&
\htmladdimg{exercises/circle.png}\\
(a) Rectangle & (b) Circle
\end{tabular}
\end{htmlonly}
\caption{Two Shapes}
\label{fig:shapes2}
\end{figure}

\xminisec{Note:}

The subfigure labels (a, b, etc) should typically be in
italic~\cite{turabian96}. This can be achieved with the
\isty{caption} package using:
\begin{definition}
\gls{DeclareCaptionLabelFormat}\marg{\meta{format-name}}\marg{\meta{code}}
\end{definition}
where \meta{format-name} is the name for this new format and
\meta{code} is the code used to format the label
where \texttt{\glsni{hashchar}2} gets replaced by the reference number.

Once you have defined a new format, you can then use
\begin{definition}
\gls{captionsetup}\oarg{\meta{type}}\marg{\meta{options}}
\end{definition}
to switch to that new format. For subfloats, \meta{type} needs to be
set to \texttt{sub}. The second argument \meta{options} is a
\meta{key}=\meta{value} comma-separated list. The key that sets the format is
\ikeyvalopt{captionsetup}{labelformat}. (For further details about
both \glsni{DeclareCaptionLabelFormat} and \glsni{captionsetup}, see
the \isty{caption} package \htmlref{documentation}{sec:texdoc}.)

For example, to create a format called \texttt{em-noparens} that
displays the number in an emphasized font without parentheses:
\begin{codeS}
\glsni{DeclareCaptionLabelFormat}\marg{em-noparens}\marg{\glsni{emph}\marg{\glsni{hashchar}2}}
\end{codeS}
Now switch to that new format:
\begin{codeS}
\glsni{captionsetup}\oarg{sub}\marg{labelformat=em-noparens}
\end{codeS}

Note that this only changes the caption label format. It doesn't affect the font 
used by \gls{ref} or \glsni{subref}.  For \gls{ref}, you can use the 
\isty{fncylab} package, which provides the command:
\begin{definition}
\gls{labelformat}\marg{\meta{ctr}}\marg{\meta{defn}}
\end{definition}
Within \meta{defn}, use \glsni{hashchar}1 to represent the subfigure value
and use \glsi{thefigure} for the encapsulating figure number. For
example:
\begin{codeS}
\glsni{labelformat}\marg{subfigure}\marg{\glsni{thefigure}\glsni{emph}\marg{\glsni{hashchar}1}}
\end{codeS}
Now
\begin{codeS}
\glsni{ref}\marg{fig:circle}
\end{codeS}
will produce
\begin{resultS}[7.5a (the letter a is emphasized)]
\ref*{fig:shapes2}\emph{a}
\end{resultS}
Unfortunately, this doesn't work for \glsni{subref}. Instead you will have 
to do, for example, the following in the text:
\begin{codeS}
\glsni{emph}\marg{\glsni{subref}\marg{fig:circle}}
\end{codeS}

If you want to add parentheses, the above can be modified
to:
\begin{code}
\glsni{DeclareCaptionLabelFormat}\marg{em-parens}\marg{(\glsni{emph}\marg{\glsni{hashchar}2})}\newline
\glsni{captionsetup}\oarg{sub}\marg{labelformat=em-parens}\newline
\glsni{labelformat}\marg{subfigure}\marg{\glsni{thefigure}(\glsni{emph}\marg{\glsni{hashchar}1})}
\end{code}
For \glsni{subref}, you will have to do, for example, the following
in the text:
\begin{codeS}
(\glsni{emph}\marg{\glsni{subref}\marg{fig:circle}})
\end{codeS}

\begin{exercise}{Creating Sub-Figures}{ex:subfloats}

Download the image files \htmladdnormallink{\texttt{rectangle.pdf}}{exercises/rectangle.pdf}
and \htmladdnormallink{\texttt{circle.pdf}}{exercises/circle.pdf}
\latex{from \url{\exerciseurl/}} (or create your own images) and 
add \figureref{fig:shapes2} to your document.  You can
\downloadorview{subfloats} the solution.
\end{exercise}

%%%%%%%%%%%%%%%%%%% DEFINING COMMANDS %%%%%%%%%%%%%%%%%%%

\setnode{newcom}
\chapter{Defining Commands}
\label{ch:newcom}

It's possible to define your own \glspl{command} or
redefine existing ones.  Be very careful about redefining existing
commands; don't redefine a command simply because you want to use the
name, only redefine it if you are making a modification.  For
example, if you want to change the format of the current date, you
would redefine \gls{today}, but if you want to define a command
to display a specific date, you should define a new command with a
different name.

There are several reasons why you might want to define
a new command:
\begin{enumerate}

\item Reduce typing:

Suppose you have a series of commands or text that
you find yourself frequently using, then you could define
a command to do all these other commands for you.

\xminisec{Example:}  

Suppose you want a lot of large bold slanted sans-serif
portions of text within your document.  Every time you type those
portions of text, you will have to do something like:
\begin{codeS}
\glsni{textsf}\marg{\glsni{large}\glsni{bfseries}\glsni{slshape} Some text}
\end{codeS}%
It would be much easier if you could use just one command
to do all that, called, say, \cmdname{largeboldsfsl}:
\begin{codeS}
\begin{verbatim}
\largeboldsfsl{Some text}
\end{verbatim}
\end{codeS}%
or you could call it, say, \cmdname{lbsfsl} which is shorter,
but slightly less memorable:
\begin{codeS}
\begin{verbatim}
\lbsfsl{Some text}
\end{verbatim}
\end{codeS}

\item Ensure consistency:

You may find that you want to format an object a certain
way.

\xminisec{Example:}

Recall near the end of \sectionref{sec:subfloats}, I~suggested the
following to reference a subfigure (when using \gls{subref} instead
of \gls{ref}):
\begin{codeS}
(\glsni{emph}\marg{\glsni{subref}\marg{fig:circle}})
\end{codeS}
For consistency, you might want to define a command, say,
\begin{alltt}
\cmdname{formattedsubref}\marg{\meta{label}}
\end{alltt}
that was the same as 
(\glsni{emph}\marg{\glsni{subref}\marg{\meta{label}}}).

\xminisec{Another Example}

\refstepcounter{object}\label{example:keywords}%
Suppose your document has a lot of keywords
in it, and you want to format these keywords in a different font,
say sans-serif, so that they stand out.  You could just
do:
\begin{codeS}
A \glsni{textsf}\marg{command} usually begins with a backslash.
\end{codeS}%
however, it is better to define a new command called, say,
\cmdname{keyword} that will typeset its argument in a
sans-serif font.  That way it becomes a lot easier
to change the format at some later date.  For example,
you may decide to splash out and have your keywords
typed in a particular colour.  In which case, all you
need to do is simply change the definition of the
command \cmdname{keyword}, otherwise you'll have to 
go through your entire document looking for keywords
and changing each one which could be very time consuming if
you have a large document.  You might also decide at some
later date to make an index for your document.  Indexing
all the keywords then becomes much simpler, as again all you'll
need to do is modify the \cmdname{keyword} command.

\end{enumerate}

New commands are defined using the command:
\begin{definition}
\gls{newcommand}\marg{\meta{cmd}}\oarg{\meta{n-args}}\oarg{\meta{default}}\marg{\meta{text}}
\end{definition}%
The first \gls{mandatory} \meta{cmd} is
the name of your new command, which must start with a backslash.
The \gls{optional} \meta{n-args} specifies
how many arguments your new command must take.  The next
optional argument \meta{default} will be discussed later.
The final mandatory argument \meta{text} specifies what \LaTeX\ 
should do every time it encounters this command.

\xminisec{Example (No Parameters):}

Let's begin with a trivial example.  Suppose I~wanted to write a
document about a particular course, say \dq{Programming \textemdash\  Languages
and Software Construction}, and I~had to keep writing the course
title, then I~might decide to define a command that prints the course
title rather than having to laboriously type it out every time.
Let's call our new command \cmdname{coursetitle}.  We want the
following code:
\begin{codeS}
The course \glsni{emph}\marg{\cmdname{coursetitle}} is an undergraduate course.
\end{codeS}%
to produce the following output:
\begin{result}[coursetitle.html]
The course \emph{Programming \textemdash\ Languages and Software Construction}
is an undergraduate course.
\end{result}%
Clearly this command doesn't need any arguments, so we don't need to
worry about the optional argument \meta{n-args} to
\gls{newcommand}, and the only thing our new command needs to do
is print:
\begin{alltt}
Programming \glsni{emdash} Languages and Software Construction
\end{alltt}
so we would define our new command as follows:
\begin{code}
\glsni{newcommand}\marg{\cmdname{coursetitle}}\marg{Programming \glsni{emdash} Languages and Software Construction}
\end{code}%
Commands must always be defined before they are used.  The
best place to define commands is in the 
\gls{preamble}:
\begin{code}
\begin{alltt}
\glsni{documentclass}\marg{scrartcl}

\glsni{newcommand}\marg{\cmdname{coursetitle}}\marg{Programming \glsni{emdash} Languages 
and Software Construction}

\glsni{begin}\marg{document}

\glsni{section}\marg{\cmdname{coursetitle}}

The course \glsni{emph}\marg{\cmdname{coursetitle}} is an undergraduate course.

\glsni{end}\marg{document}
\end{alltt}
\end{code}

\xminisec{Example (One Mandatory Argument):}

\refstepcounter{object}\label{sec:newcomarg}%
Now let's try defining a command that takes an \gls{argument} (or
parameter).  Let's go back to our \htmlref{\cmdname{keyword}
example}{example:keywords}\latex{ \vpageref{example:keywords}}.  This
command needs to take one argument that is the keyword.  Let's
suppose we want keywords to come out in
\htmlref{sans-serif}{sec:fontstyle}, then we could do:
\begin{codeS}
\glsni{newcommand}\marg{\cmdname{keyword}}\oarg{1}\marg{\glsni{textsf}\marg{\glsni{hashchar}1}}
\end{codeS}%
In this case we have used the optional argument \meta{n-args} to
\gls{newcommand}.  We want our command \cmdname{keyword} to have
one argument, so we have \texttt{[1]}.  In \texttt{\glsni{textsf}\marg{\glsni{hashchar}1}} the
\gls{hashchar}\texttt{1} represents the first argument.  (If we had
more than one argument, \glsni{hashchar}\texttt{2} would represent the second
argument, \glsni{hashchar}\texttt{3} would represent the third argument etc.\ up to a
maximum of~9\faq{How to break the 9-argument limit}{moren9}.) So
\begin{verbatim}
\keyword{commands}
\end{verbatim}
will be equivalent to 
\begin{alltt}
\glsni{textsf}\marg{commands}
\end{alltt}
and
\begin{verbatim}
\keyword{environment}
\end{verbatim}
will be equivalent to
\begin{alltt}
\glsni{textsf}\marg{environment}
\end{alltt}
and so on.
Again, it's best to put the command definition in the preamble to
ensure the command won't be used before it's defined.

\begin{code}
\begin{alltt}
\glsni{documentclass}\marg{scrartcl}

\glsni{newcommand}\marg{\cmdname{keyword}}\oarg{1}\marg{\glsni{textsf}\marg{\glsni{hashchar}1}}

\glsni{begin}\marg{document}

A \cmdname{keyword}\marg{command} usually begins with a backslash.

\glsni{end}\marg{document}
\end{alltt}
\end{code}

Now if we want to change the way the keywords are formatted,
we can simply change the definition of \cmdname{keyword}.
Let's modify our code so that the keyword is now in a 
slanted sans-serif font:
\begin{code}
\begin{alltt}
\glsni{documentclass}\marg{scrartcl}

\glsni{newcommand}\marg{\cmdname{keyword}}\oarg{1}\marg{\glsni{textsf}\marg{\glsni{slshape} \glsni{hashchar}1}}

\glsni{begin}\marg{document}

A \cmdname{keyword}\marg{command} usually begins with a backslash.

\glsni{end}\marg{document}
\end{alltt}
\end{code}%
Let's go one stage further.  The \isty{color} 
\htmlref{package}{sec:packages} provides the \gls{declaration}:
\begin{definition}
\gls{color}\marg{\meta{col-name}}
\end{definition}
which switches the foreground colour to \meta{col-name}. It also
provides the text-block command:
\begin{definition}
\gls{textcolor}\marg{\meta{col-name}}\marg{\meta{text}}
\end{definition}
which sets \meta{text} in the colour given by \meta{col-name}.

So let's use the \isty{color} package to make our keywords
blue:\screenpagebreak
\begin{code}
\begin{alltt}
\glsni{documentclass}\marg{scrartcl}

\glsni{usepackage}\marg{color}

\glsni{newcommand}\marg{\cmdname{keyword}}\oarg{1}\marg{\glsni{textsf}\marg{\glsni{slshape}\glsni{color}\marg{blue}\glsni{hashchar}1}}

\glsni{begin}\marg{document}

A \cmdname{keyword}\marg{command} usually begins with a backslash.

\glsni{end}\marg{document}
\end{alltt}
\end{code}
Or we could index the keywords.  To do this we need the
\isty{makeidx} \htmlref{package}{sec:packages} and the commands 
\gls{makeindex}, 
\begin{inlinedef}\gls{index}\marg{\meta{text}}\end{inlinedef}
and \gls{printindex}:
\begin{code}
\begin{alltt}
\glsni{documentclass}\marg{scrartcl}

\glsni{usepackage}\marg{makeidx}

\glsni{makeindex}

\glsni{newcommand}\marg{\cmdname{keyword}}\oarg{1}\marg{\glsni{textsf}\marg{\glsni{slshape} \glsni{hashchar}1}\glsni{index}\marg{\glsni{hashchar}1}}

\glsni{begin}\marg{document}

A \cmdname{keyword}\marg{command} usually begins with a backslash.

\glsni{printindex}

\glsni{end}\marg{document}
\end{alltt}
\end{code}
For further information about how to create an index, 
see \latexguide\ or \latexcomp.
Alternatively, if you want a brief overview, try \latexthesis.

Since it is unlikely that the keyword will contain a paragraph break,
we should indicate that this is a 
\gls{short} using the \htmlref{starred form}{itm:starredcommand}:
\begin{codeS}
\glsni{newcommand}*\marg{\cmdname{keyword}}\oarg{1}\marg{\glsni{textsf}\marg{\glsni{slshape} \glsni{hashchar}1}\glsni{index}\marg{\glsni{hashchar}1}}
\end{codeS}
Now if you forget to add the closing brace, for example,
\verb|\keyword{command|, then \gls{tex}['s] error checking mechanism will 
pick up the error sooner. This will give an error message that
looks like:
\begin{verbatim}
! Paragraph ended before \keyword was complete.
<to be read again>
                   \par
l.604
\end{verbatim}
This at least gives you the line number (604 in this example) of the
end of the paragraph where the error has occurred.

If you don't used the starred form of \gls{newcommand}, then
you will get the somewhat less than helpful error:
\begin{verbatim}
! File ended while scanning use of \keyword.
\end{verbatim}
If you have a very large document, it may take a while to track
down where exactly you have missed a brace.

\begin{exercise}{Defining a New Command}{ex:newcom}
Try typing up the following code:
\begin{bcode}
\begin{alltt}
\glsni{documentclass}\marg{scrartcl}

\glsni{newcommand}*\marg{\cmdname{keyword}}\oarg{1}\marg{\glsni{textsf}\marg{\glsni{hashchar}1}}

\glsni{begin}\marg{document}

A \cmdname{keyword}\marg{command} usually begins with a backslash.

Segments of code may be \cmdname{keyword}\marg{grouped}.

Some \cmdname{keyword}\marg{commands} take one or more \cmdname{keyword}\marg{arguments}.
\glsni{end}\marg{document}
\end{alltt}
\end{bcode}

Then modify your code so that the keywords are in a slanted
sans-serif font or modify your code so that the keywords come out in
blue (using the \isty{color} package as in the example earlier).
Again you can \downloadorview{newcom} the result.

\xminisec{For the more adventurous:}

If you want to create an index as in the previous example, you will
need to use the application \iappname{makeindex}. If you used
\iappname{latexmk} back in \sectionref{sec:crossref}, you can just
carry on using that as before. If not you need to do the following
in TeXworks:
\begin{enumerate}
  \item Create the PDF as described in \sectionref{sec:texworks}.

  \item Select \menu{MakeIndex} from the drop-down list next to the
   build (typeset) button (see \figureref{fig:texworks9}).

  \item Click on the build button. If all goes well, you won't see
    anything different. If you see something like the following:
\begin{alltt}
Couldn't find input index file exercise\theexercise\space nor exercise\theexercise.idx.
\end{alltt}
    then you probably forgot to add the command \gls{makeindex} to
    the preamble. Add it in and go back to Step~1.

  \item Select \menu{pdfLaTeX} from the drop-down list and build the
    PDF file again. Move to the last page of the PDF, and you should
    see the index.
\end{enumerate}

\begin{figure}[htbp]
\figconts
  {pictures/texworks9}
  {%
    \caption{Selecting MakeIndex in TeXWorks}
  }%
  {fig:texworks9}
\end{figure}

\end{exercise}

%%%%%%%%%%%%% Defining Commands with an Optional Argument %%%%%%%%%%%%%

\setnode{newcomopt}
\section{Defining Commands with an Optional Argument}
\label{sec:newcomopt}

As mentioned earlier, the \gls{newcommand} command has a second
optional argument \meta{default}.  This allows you to define a
command with an optional argument\faq{More than one optional
argument}{twooptarg}.  For example, suppose we want a command called,
say, \cmdname{price}.  Suppose we want the following code:
\begin{codeS}
\begin{verbatim}
\price{100}
\end{verbatim}
\end{codeS}%
to produce the following output:
\begin{resultS}[price.html]
\pounds100 excl VAT @ 17.5\%\relax
\end{resultS}%
and let's suppose we want an optional argument
so that we can change the VAT.  That is, we
would want the following code:
\begin{codeS}
\begin{verbatim}
\price[20]{30}
\end{verbatim}
\end{codeS}%
to produce the following output:
\begin{resultS}[price2.html]
\pounds30 excl VAT @ 20\%
\end{resultS}%
Therefore we want to define a command
such that if the optional argument is absent 
we will have \texttt{17.5}, and if it is present
the optional argument will be substituted instead.
This command can be defined as follows:
\begin{codeS}
\gls{newcommand}\marg{\cmdname{price}}\oarg{2}\oarg{17.5}\marg{\glsni{pounds} \gls{hashchar}2 excl VAT @ \glsni{hashchar}1\glsni{percent}}
\end{codeS}%
Here, \texttt{\glsni{hashchar}1} represents the optional argument (by default
\texttt{17.5}) and \texttt{\glsni{hashchar}2} represents the mandatory argument (the
second argument if the optional argument is present, or the only
argument if the optional argument is absent.)

As before, since the argument is unlikely to contain a paragraph
break, we should indicate that it is a \gls{short} using the \htmlref{starred form}{itm:starredcommand}:
\begin{codeS}
\gls{newcommand}*\marg{\cmdname{price}}\oarg{2}\oarg{17.5}\marg{\glsni{pounds} \glsni{hashchar}2 excl VAT @ \glsni{hashchar}1\gls{percent}}
\end{codeS}%

\begin{exercise}{Defining Commands with an Optional Argument}{ex:newcomopt}
In this exercise, you will need to define a slightly modified
version of the above example.  Try defining a command called,
say, \cmdname{cost}.  It should take one optional argument and
one mandatory argument.  Without the optional argument, it
behaves in the same way as the \cmdname{price} example above,
so that, say,
\begin{codeS}
\begin{verbatim}
\cost{50}
\end{verbatim}
\end{codeS}%
will produce
\begin{resultS}[price3.html]
\pounds50 excl VAT @ 17.5\%\null
\end{resultS}%
but with the optional argument, you can change the
\verb|excl VAT @ 17.5\%| bit.  So that, say,
\begin{codeS}
\begin{verbatim}
\cost[inc VAT]{50}
\end{verbatim}
\end{codeS}%
will produce
\begin{resultS}[price4.html]
\pounds50 inc VAT\mbox{}% LaTeX2HTML is getting confused!
\end{resultS}

You can \downloadorview{newcomopt} the solution.

\xminisec{For the more adventurous:}
If you did \exerciseref{ex:newcom} and you modified \cmdname{keyword}
so that it indexed the keyword, you may have noticed that
\begin{codeS}
\verb|\keyword{command}|
\end{codeS}
and
\begin{codeS}
\verb|\keyword{commands}|
\end{codeS}
produced separate entries in the index. It would be better to have an
optional argument to override the indexing mechanism. For example,
\begin{codeS}
\verb|\keyword{command}|
\end{codeS}
should print and index the word \dq{command},
whereas
\begin{codeS}
\verb|\keyword[command]{commands}|
\end{codeS}
should print \dq{commands}
and index \dq{command}. In other words, we need an optional
argument that defaults to the mandatory argument if it is not
present.
This is how to achieve that type of effect:\footnote{Recall from \chapterref{ch:def} the percent symbol discards the space resulting from the end
of line character.}
\begin{code}
\glsni{newcommand}*\marg{\cmdname{keyword}}\oarg{2}\oarg{\cmdname{keywordentry}}\marg{\glsni{percentchar}\newline
\strut~~\icmdname{def}\cmdname{keywordentry}\marg{\glsni{hashchar}2}\glsni{percentchar}\newline
\strut~~\glsni{textsf}\marg{\glsni{hashchar}2}\glsni{percentchar}\newline
\strut~~\gls{index}\marg{\glsni{hashchar}1}\glsni{percentchar}\newline
}
\end{code}
In this example, the default value for the optional argument
is set to the command \cmdname{keywordentry}. At the start of 
\cmdname{keyword} this is defined to be the mandatory argument 
(as specified by \verb|#2|) using \gls{tex}['s] \cmdname{def}
command:\footnote{\cmdname{def} is too complicated for an
introductory \LaTeX\ guide but, if you're interested, read \emph{The
\TeX{}book}~\cite{texbook}.}
\begin{alltt}
\cmdname{def}\cmdname{keywordentry}\marg{\glsni{hashchar}2}
\end{alltt}
Then typeset the keyword (given in the mandatory
argument \texttt{\glsni{hashchar}2}) in a sans-serif font:
\begin{alltt}
\glsni{textsf}\marg{\glsni{hashchar}2}
\end{alltt}
Now index the term using the optional argument (\texttt{\glsni{hashchar}1}):
\begin{alltt}
\glsni{index}\marg{\glsni{hashchar}1}
\end{alltt}
If an optional argument is specified, \texttt{\glsni{hashchar}1} will be the given
argument, but if the optional argument is missing, \texttt{\glsni{hashchar}1} will
be \cmdname{keywordentry}, which has earlier been set to
the mandatory argument \texttt{\glsni{hashchar}2}.
\end{exercise}

%%%%%%%%%%%%%%%%%%%%%%%%% Redefining Commands %%%%%%%%%%%%%%%%%%%%%

\setnode{renewcom}
\section{Redefining Commands}
\label{sec:renewcom}

Commands can be redefined using the command:
\begin{definition}
\gls{renewcommand}\marg{\meta{cmd}}\oarg{\meta{n-args}}\oarg{\meta{default}}\marg{\meta{text}}
\end{definition}\reportpagebreak\noindent
This has exactly the same format as \gls{newcommand} but is
used for redefining existing commands.
Again there is a \htmlref{starred version}{itm:starredcommand} to indicate that the command
is a \gls{short}.

\xminisec{Caveat:}

\warning Never redefine a command whose existing function is unknown
to you or just because you want to use a particular command name,
regardless of its previous function. By way of illustration: as a
production editor, I~have to combine articles by different authors
into a single book. Each author supplies the \LaTeX\ code for their
own article. Every so often, I~get code that redefines a command for
the convenience of the author. Later on another author tries to use
the same command, on the assumption that the command behaves
according to its original definition. This tends to involve the
accent commands as they are short and that saves the author typing.
It usually goes along these lines: author~A redefines \gls{c} (the
cedilla accent command) to display a maths bold
\dq{\boldc} to indicate a
vector. Later, author~B, uses the cedilla accent, say, in the name
Fran\c{c}ois:
\begin{alltt}
Fran\gls{c}\marg{c}ois
\end{alltt}
Author~A's hack turns this into
Fran\boldc cois.

\xminisec{Example (Redefining List Labels):}

Recall the \gls{env-itemize} environment discussed in
\sectionref{sec:itemize}.  You may have up to four nested
\gls{env-itemize} environments, the labels for the outer environment
are specified by the command \gls{labelitemi}, the labels for
the second level are specified by \gls{labelitemii}, the third
by \gls{labelitemiii} and the fourth by \gls{labelitemiv}.
By default, \glsni{labelitemi} is a bullet point (\makeimg{bullet
symbol}{\labelitemi}),
\glsni{labelitemii} is an en dash (\makeimg{en dash}{\labelitemii}), 
\glsni{labelitemiii} is an
asterisk (\makeimg{asterisk}{\labelitemiii}) and 
\glsni{labelitemiv} is a centred dot (\makeimg{centered
dot}{\labelitemiv}).  These can be changed by redefining
\glsni{labelitemi} etc.

Recall from \tableref{tab:symbols} that the command
\gls{dag} produces a dagger symbol, we can use this symbol
instead of a bullet point:
\begin{code}
\glsni{renewcommand}*\marg{\glsni{labelitemi}}\marg{\gls{dag}}\newline
\glsni{begin}\marg{itemize}\newline
\mbox{}\newline
\gls{item} Animal\newline
\mbox{}\newline
\glsni{item} Mineral\newline
\mbox{}\newline
\glsni{item} Vegetable\newline
\mbox{}\newline
\glsni{end}\marg{itemize}
\end{code}%
Output:
\begin{result}[List uses dagger symbol instead of a bullet point for
the marker]
\renewcommand*{\labelitemi}{\dag}

\begin{itemize}

\item Animal

\item Mineral

\item Vegetable

\end{itemize}
\end{result}

Here's another example, it uses the \Index{PostScript}
font ZapfDingbats via the \isty{pifont} package:
\begin{code}
\glsni{renewcommand}*\marg{\glsni{labelitemi}}\marg{\gls{ding}\marg{43}}\newline
\glsni{begin}\marg{itemize}\newline
\mbox{}\newline
\glsni{item} Animal\newline
\mbox{}\newline
\glsni{item} Mineral\newline
\mbox{}\newline
\glsni{item} Vegetable\newline
\mbox{}\newline
\glsni{end}\marg{itemize}
\end{code}\bookpagebreak\noindent
Output:
\begin{result}[List uses a pointing finger instead of a bullet point
for the marker]
\renewcommand*{\labelitemi}{\ding{43}}

\begin{itemize}

\item Animal

\item Mineral

\item Vegetable

\end{itemize}
\end{result}%
In the above example, it would actually be easier to use the
\gls{env-dinglist} environment defined in the \isty{pifont} package:
\begin{code}
\begin{alltt}
\glsni{begin}\marg{dinglist}\marg{43}

\glsni{item} Animal

\glsni{item} Mineral

\glsni{item} Vegetable\screenpagebreak

\glsni{end}\marg{dinglist}
\end{alltt}
\end{code}


\xminisec{Example (Redefining the Default Font):}

Recall from \sectionref{sec:changingfonts} that the default font
family is usually the serif (Roman) family. So what happens if you
need to write your entire document in, say, Helvetica?  The default
font family name is stored in:
\begin{definition}
\gls{familydefault}
\end{definition}
This command is usually
defined to be just \gls{rmdefault}, which in turn stores the name of the
default serif font (initially \texttt{cmr}, Computer Modern Roman). If 
you want the default font to be sans-serif, all you
need do is add the following line to the
\glsterm{preamble}:
\begin{codeS}
\glsni{renewcommand}\marg{\glsni{familydefault}}\marg{\gls{sfdefault}}
\end{codeS}
\glsni{sfdefault} stores the name of the default sans-serif font (initially \texttt{cmss},
Computer Modern Sans-Serif) and
the \isty{helvet} package redefines \glsni{sfdefault} to \texttt{phv},
which is the identifier for the Helvetica font.
So the following document will be in Helvetica:
\begin{alltt}
\glsni{documentclass}\marg{scrartcl}

\glsni{usepackage}\marg{helvet}
\glsni{renewcommand}\marg{\glsni{familydefault}}\marg{\glsni{sfdefault}}

\glsni{begin}\marg{document}
This is a sample document.
\glsni{end}\marg{document}
\end{alltt}
Similarly, if you want the default font to be monospaced (typewriter) then you'd need to do:
\begin{codeS}
\glsni{renewcommand}\marg{\glsni{familydefault}}\marg{\gls{ttdefault}}
\end{codeS}

Incidentally, you may have noticed in \sectionref{sec:changingfonts}
that although I~said I'd used the \isty{anttor} and
\isty{libris} packages to set the serif and sans-serif families for
\latexhtml{this book}{the PDF version of this book}, I
didn't mention anything about the typewriter (monospaced) font. I~used
the TXTT font, but that doesn't have a corresponding package. You
just redefine \glsni{ttdefault} to \texttt{txtt}:
\begin{codeS}
\glsni{renewcommand}*\marg{\glsni{ttdefault}}\marg{txtt}
\end{codeS}

\xminisec{Example (Redefining Fixed Names):}

You may have noticed that \LaTeX\ automatically generates pieces of
text such as \dq{Chapter}, \dq{Figure}, \dq{Bibliography}.  These are
generated by the commands shown in \tableref{tab:textlab}.

\begin{table}[hbtp]
\caption[Predefined Names]{Predefined Names 
(\supdag Book and report style classes (such as \icls{scrreprt} and \icls{scrbook}), 
\supddag article-style classes (such as \icls{scrartcl}), 
remainder book, report and article-style classes)}
\label{tab:textlab}
\centering
\begin{tabular}{ll}
\toprule
\bfseries Command & \bfseries Default Text\\
\midrule
\gls{contentsname} & \contentsname\\
\gls{listfigurename} & \listfigurename\\
\gls{listtablename} & \listtablename\\
\gls{bibname}\supdag & \bibname\\
\gls{refname}\supddag & \refname\\
\gls{indexname} & \indexname\\
\gls{figurename} & \figurename\\
\gls{tablename} & \tablename\\
\gls{partname} & \partname\\
\gls{chaptername}\supdag & \chaptername\\
\gls{appendixname} & \appendixname\\
\gls{abstractname} & \abstractname
\\\bottomrule
\end{tabular}
\end{table}

You can change the defaults using \gls{renewcommand}.  For
example, suppose you want the table of contents to be labelled
\dq{Table of Contents}, instead of the default \dq{Contents}, you
would need to do:
\begin{codeS}
\glsni{renewcommand}*\marg{\glsni{contentsname}}\marg{Table of Contents}
\end{codeS}
\faq{Changing the words babel uses}{latexwords}In fact, the
\isty{babel} package (see \sectionref{sec:babel}) uses this method
to redefine the commands in \tableref{tab:textlab} whenever you
switch language using \reportlinebreak\screenlinebreak
\gls{selectlanguage} or within the contents of
the \gls{env-otherlanguage} environment. This unfortunately has the
side-effect that means if you try to redefine these commands,
\sty{babel} will automatically overwrite your definition whenever
there's a language change, which includes at the beginning of the
\gls{env-document} environment. Instead you need to use
\sty{babel}'s \glsni{addto} mechanism.
\begin{definition}
\gls{addto}\marg{\meta{command}}\marg{\meta{code}}
\end{definition}
This patches the definition of a command (specified in the first
\gls{argument}) adding \meta{code} to the end of the command
definition. Whenever \sty{babel} switches the current language, it
uses the command \cmdname{captions}\meta{language}, which performs
all the redefinitions of commands like those listed in
\tableref{tab:textlab}. For example, if you are using \sty{babel}
with the \istyopt{babel}{english} option and you want to change
\glsni{contentsname} so that it does \dq{Table of Contents} instead
of \dq{Contents}, you need to do:
\begin{code}
\begin{alltt}
\glsni{addto}\marg{\cmdname{captionsenglish}}\marg{\glsni{percentchar}
  \glsni{renewcommand}\marg{\glsni{contentsname}}\marg{Table of Contents}\glsni{percentchar}
}
\end{alltt}
\end{code}

\xminisec{Notes:}

\warning Take care if you want to patch an existing command. For
example, suppose you want to append something to the action of a
command. You might be tempted to do
\begin{alltt}
\wrong\glsni{renewcommand}*\marg{\cmdname{foo}}\marg{\cmdname{foo} Something else}
\end{alltt}
This will cause an infinite loop where \cmdname{foo} recursively
calls itself. Instead you should use one of the commands provided by
the \isty{etoolbox} package (such as \gls{appto}, which has the same
syntax as \sty{babel}'s \glsni{addto} described above). For further details,
read the \sty{etoolbox} \htmlref{documentation}{sec:texdoc}.

\begin{exercise}{Renewing Commands}{ex:renewcom}
If you did Exercises~\ref{ex:figures} and \ref{ex:tables}, go back to
that document and changed the figures and tables so that they
are labelled \dq{Fig} and \dq{Tab} instead of \dq{Figure} and \dq{Table}.
Hint: you need to redefine \glsni{tablename} and \glsni{figurename}.

You can \downloadorview{renewcom} the solution.
\end{exercise}

%%%%%%%%%%%%%%%%%%%%%%%% MATHEMATICS %%%%%%%%%%%%%%%%%%%%%%%%%%%%%%%

\setnode{mathematics}
\chapter{Mathematics}
\label{ch:maths}

As mentioned in the \htmlref{introduction}{ch:intro}, \LaTeX\ is
particularly good at typesetting mathematics.  In order to
use any of the maths commands we need to be in one of the
mathematics \glspl{environment}.  There are two
basic types of mathematics: \keyword{in-line maths} and
\keyword{displayed maths}. In-line maths is mathematics that
occurs within a line of text, for example:
\begin{resultS}[inlinemaths.html]
The variable $x$ is transformed by the function $f(x)$.
\end{resultS}%
Displayed maths is mathematics that occurs on a line
of its own.  For example:
\begin{result}[displaymath.html]
A polynomial is a function of the form
\[
f(x) = \sum_{i=0}^na_ix^i
\]
\end{result}
The maths environments switch to \LaTeX's \dq{math mode}, which uses
specialist maths fonts and spacing rather than just using an italic
font.

If you want to typeset any mathematics, I~strongly advise using the
\istystart{amsmath} package:
\begin{alltt}
\glsni{usepackage}\marg{amsmath}
\end{alltt}
This patches some existing \LaTeX\ commands and environments and
also provides many useful additions.

This chapter is just an introduction to typesetting mathematics in
\LaTeX. If you want a comprehensive guide, I~recommend you read \emph{Math Mode}
by Herbert~Vo\ss~\cite{voss}, which can be access via
\appname{texdoc} (see \sectionref{sec:texdoc}):
\begin{verbatim}
texdoc mathmode
\end{verbatim}

%%%%%%%%%%%%%%%% In-Line Mathematics %%%%%%%%%%%%%%%%%%%

\setnode{inlinemaths}
\section{In-Line Mathematics}
\label{sec:inline}

In-line mathematics is created using the \gls{env-math}
environment. (Note U.S. spelling --- \dq{math} not \dq{maths}).

\xminisec{Example:}
\begin{code}
The variable \glsni{begin}\marg{math}x\glsni{end}\marg{math} is transformed
by the function \glsni{begin}\marg{math}f(x)\glsni{end}\marg{math}.
\end{code}
It's somewhat cumbersome having to type
\glsni{begin}\marg{math} and \glsni{end}\marg{math} and it also
makes the \gls{source} a little difficult to read so there
are shorthand notations that can be used instead:
\gls{beginmath} is equivalent to \glsni{begin}\marg{math} and
\gls{endmath} is equivalent to \glsni{end}\marg{math}.  So
the example above can be rewritten:
\begin{codeS}
The variable \glsni{beginmath}x\glsni{endmath} is transformed by the function 
\glsni{beginmath}f(x)\glsni{endmath}.
\end{codeS}%
There is an even shorter notation:
The \htmlref{special character}{sec:chars} \gls{dollarchar} is equivalent
to both \glsni{begin}\marg{math} and \glsni{end}\marg{math}:
\begin{codeS}
The variable \glsni{dollarchar}x\glsni{dollarchar} is transformed by the 
function \glsni{dollarchar}f(x)\glsni{dollarchar}.
\end{codeS}%
This is considerably easier to type and
to read, but you need to make sure that all your \glsni{dollarchar}
symbols have matching pairs.  The above code will look like:
\begin{resultS}[inlinemaths2.html]
The variable $x$ is transformed
by the function $f(x)$.
\end{resultS}%
The other advantage in using \glsni{dollarchar} over \glsni{beginmath} and
\glsni{endmath} is that \glsni{dollarchar} is a \gls{robust}, 
whereas \glsni{beginmath} and \glsni{endmath} are 
\glspl{fragile} and
will need to be protected if they occur in a \glslink{fragile}{moving
argument}.\casemedia{}{}{ (See \sectionref{sec:fragile}.)}{}

Note: you should always make sure you are in maths mode to
typeset any variables (such as $x$, $y$, $z$), as this will
ensure that the correct maths fonts are used, as well as the
appropriate spacing. Similarly, don't use \glsni{dollarchar}
as a short cut for an italic font.
\begin{codeS}[\linewidth]
Notice the \glsni{dollarchar}difference\glsni{dollarchar} between 
\glsni{dollarchar}(x', y', z')\glsni{dollarchar} and
\glsni{textit}\marg{(x', y', z')}.
\end{codeS}%
\begin{resultS}[Image of result: the brackets are upright in the
maths font but slanted in the italic font. The apostrophe becomes a
prime in the maths font. The spacing is also
different.]
Notice the $difference$ between $(x', y', z')$ and
\textit{(x', y', z')}.
\end{resultS}

%%%%%%%%%%%%%%%% Displayed Mathematics %%%%%%%%%%%%%%%%%%%%%%

\setnode{displayedmaths}
\section{Displayed Mathematics}
\label{sec:displayed}

One-line unnumbered displayed mathematics can be created using:
\begin{definition}
\gls{begindispmath}\meta{maths}\gls{enddispmath}
\end{definition}
where \meta{maths} is the mathematics to be displayed.

\xminisec{Example:}
\begin{code}
\begin{alltt}
A linear function is a function of the form
\glsni{begindispmath} y = mx + c \glsni{enddispmath}
\end{alltt}
\end{code}%
Output:
\begin{result}[displaymath2.html]
A linear function is a function of the form
\[ y = mx + c \]
\end{result}%

\warning Don't use the \envname{displaymath} environment or \verb|$$|\ldots\verb|$$|~\cite{l2tabu}.
Use \glsni{begindispmath} and \glsni{enddispmath} with
the \isty{amsmath} package.

The \glsni{env-equation}\refstepcounter{object}\label{env:equation} 
environment provides something similar to
\glsni{begindispmath} \glsni{enddispmath}, except
that the equation is numbered.  Modifying the above example:
\begin{code}
\begin{alltt}
A linear function is a function of the form
\glsni{begin}\marg{equation}
y = mx + c
\glsni{end}\marg{equation}
\end{alltt}
\end{code}%
results in the following output:
\begin{result}[As before but the equation now has a number in
brackets on the right.]
A linear function is a function of the form
\begin{equation}
y = mx + c
\end{equation}
\end{result}
Normal text can be inserted into the equation using
\begin{definition}
\gls{text}\marg{\meta{text}}
\end{definition}
which is provided by the \isty{amsmath} package.

\xminisec{Example:}
\begin{codeS}
\glsni{begindispmath} x = 2 \glsni{text}\marg{ and } y = -1 \glsni{enddispmath}
\end{codeS}%
results in the following output:
\begin{result}[x equals 2 and y equals minus 1 (the word 'and' is in normal text font)]
\[ x = 2 \text{ and } y = -1 \]
\end{result}

\faq{Re-using an equation}{reuseq}%
Recall from \sectionref{sec:crossref} that we can
\htmlref{cross-reference}{sec:crossref} most things that \LaTeX\
automatically numbers using \gls{ref} and \gls{label}.
Equations can be cross-referenced in the same way:\reportpagebreak
\begin{code}
\begin{alltt}
Equation\glsni{tildechar}\glsni{ref}\marg{eqn:linear} is a linear function.
\glsni{begin}\marg{equation}
\glsni{label}\marg{eqn:linear}
f(x) = mx + c
\glsni{end}\marg{equation}\screenpagebreak
\end{alltt}
\end{code}
\begin{result}[eqref.html]
Equation~\ref*{eqn:linear} is a linear function.
\begin{equation}
\label{eqn:linear}
f(x) = mx + c
\end{equation}
\end{result}
Equation numbers are usually given in parentheses, which can be done
using:
\begin{codeS}
Equation\glsni{tildechar}(\glsni{ref}\marg{eqn:linear})
\end{codeS}
The \isty{amsmath} package provides a convenient short cut:
\begin{definition}
\gls{eqref}\marg{\meta{label}}
\end{definition}
So the above can be written as:
\begin{codeS}
Equation\glsni{tildechar}\glsni{eqref}\marg{eqn:linear}
\end{codeS}
\begin{resultS}[Equation (9.2)]
% fake it to suppress hyperlink
Equation~(\ref*{eqn:linear})
\end{resultS}

\xminisec{Note:}

Both the \gls{env-equation} environment and
\gls{begindispmath}\ldots\gls{enddispmath} are only designed for one
line of maths. Therefore you must not have any line breaks or
paragraph breaks within them.
The following will cause an error:
\begin{alltt}
\glsni{begin}\marg{equation}
\wrong
f(x) = mx + c
\wrong
\glsni{end}\marg{equation}
\end{alltt}
Either remove the blank lines or comment them out:
\begin{alltt}
\glsni{begin}\marg{equation}
\glsni{percentchar}\correct
f(x) = mx + c
\glsni{percentchar}\correct
\glsni{end}\marg{equation}
\end{alltt}

\setnode{align}
\section{Multiple Lines of Displayed Maths}
\label{sec:align}

The \isty{amsmath} package provides the \gls{env-align}
and \gls{env-align*} environments for aligned equations.
The starred version doesn't number the equations. These environments 
provide pairs of left- and right-aligned columns.
As with the \gls{env-tabular} environment, use \gls{ampchar} to
separate columns and \gls{tab.dbbackslashchar} to separate rows.
Unlike the \gls{env-tabular} environment, there is no argument as
the column specifiers are predefined. Another difference is that no
page breaks can occur in the \gls{env-tabular} environment, but it's
possible to allow a page break in \glsni{env-align} or
\glsni{env-align*} using
\begin{definition}
\gls{displaybreak}\oarg{\meta{n}}
\end{definition}
immediately before the \gls{tab.dbbackslashchar} where it is to take
effect. The optional argument is a number from~0 to~4 indicating the
desirability to break the page (from~0 the least to~4 the most).

If you want to mix numbered and unnumbered rows, you can use
\begin{definition}
\gls{notag}
\end{definition}
to suppress the numbering for a particular row in the \glsni{env-align} environment.
This command must go before \gls{tab.dbbackslashchar} at the end of the
row. The default equation numbering can be overridden for a
particular row using:
\begin{definition}
\gls{tag}\marg{\meta{tag}}
\end{definition}
where \meta{tag} is the replacement for the equation number.

\screenpagebreak\warning Don't use the \envname{eqnarray} or \envname{eqnarray*}
environments. They're obsolete~\cite{l2tabu}.


\xminisec{Example (Unnumbered):}

\begin{code}
\glsni{begin}\marg{align*}\newline
y \glsni{ampchar}= 2x + 2\glsni{tab.dbbackslashchar}\newline
\strut~~\glsni{ampchar}= 2(x+1)\newline
\glsni{end}\marg{align*}
\end{code}
\begin{result}[align.html]
\begin{align*}
y &= 2x + 2\\
  &= 2(x+1)
\end{align*}
\end{result}
Note that the equals sign is placed at the start of the second
column, \emph{after} the ampersand \gls{ampchar}. This ensures the
correct amount of spacing on either side. If the first line of the
above equation was changed to:\bookpagebreak
\begin{alltt}
\wrong y =\gls{ampchar} 2x + 2\gls{tab.dbbackslashchar}
\end{alltt}
there wouldn't be enough space on the right of the equal sign:
\begin{htmlresult}{y equals 2x + 2}
\begin{makeimage}
\begin{align*}
y =& 2x + 2
\end{align*}
\end{makeimage}
\end{htmlresult}

\xminisec{Example (One Row Numbered):}

\begin{code}
\glsni{begin}\marg{align}\newline
y \glsni{ampchar}= 2x + 2\glsni{notag}\glsni{tab.dbbackslashchar}\newline
\strut~~\glsni{ampchar}= 2(x+1)\newline
\glsni{end}\marg{align}
\end{code}
\begin{result}[Image of result: as the earlier example but the second equation has been numbered]
\begin{align}
y &= 2x + 2\notag\\
  &= 2(x+1)
\end{align}
\end{result}\screenpagebreak

\xminisec{Example (Four Columns):}

\begin{code}
\glsni{begin}\marg{align*}\newline
y \glsni{ampchar}= 2x + 2 \glsni{ampchar} z 
 \glsni{ampchar}= 6x + 3\glsni{tab.dbbackslashchar}\newline
\strut~~\glsni{ampchar}= 2(x+1) \glsni{ampchar}~~~\glsni{ampchar}= 3(2x+1)\newline
\glsni{end}\marg{align*}
\end{code}
\begin{result}[4colalign.html]
\begin{align*}
y &= 2x + 2 & z &= 6x + 3\\
  &= 2(x+1) &  &= 3(2x+1)
\end{align*}
\end{result}

As with \gls{env-equation}, you can cross-reference individual rows
of an \glsni{env-align} environment, but you must remember to
put \gls{label} \emph{before} the end of row \gls{tab.dbbackslashchar}
separator. You can reference a row in the \glsni{env-align*}
environment if you have assigned it a tag with \gls{tag}, but don't
try labelling a row in the \glsni{env-align} environment where the
numbering has been suppressed with \gls{notag}.

\xminisec{Example (Cross-Referenced):}

This example has two numbered equations in an \gls{env-align} environment,
both of which are labelled and referenced:
\begin{code}
The function \glsni{dollarchar}f(x)\glsni{dollarchar} is given in 
Equation\glsni{tildechar}\glsni{eqref}\marg{eq:fx},
and its derivative \glsni{dollarchar}f'(x)\glsni{dollarchar} is given in 
Equation\glsni{tildechar}\glsni{eqref}\marg{eq:dfx}.
\newline
\glsni{begin}\marg{align}\newline
f(x) \glsni{ampchar}= 2x + 1
\glsni{label}\marg{eq:fx}\glsni{tab.dbbackslashchar}\newline
f'(x) \glsni{ampchar}= 2 \glsni{label}\marg{eq:dfx}\newline
\glsni{end}\marg{align}
\end{code}
\begin{result}[alignref.html]
The function $f(x)$ is given in Equation~(\ref*{eq:fx}), and its
derivative $f'(x)$ is given in Equation~(\ref*{eq:dfx}).
\begin{align}
f(x) &= 2x + 1\label{eq:fx}\\
f'(x) &= 2 \label{eq:dfx}
\end{align}
\end{result}

Recall the command \gls{text}\marg{\meta{text}} from the
\htmlref{previous section}{sec:displayed}. This can be used
within cells of the \glsni{env-align} and \glsni{env-align*}
environments, but the \isty{amsmath} package also provides
\begin{definition}
\gls{intertext}\marg{\meta{text}}
\end{definition}
which can be used for a line of interjection between the rows. This
command may only go right after \gls{tab.dbbackslashchar}.

\xminisec{Example}

\begin{code}
\glsni{begin}\marg{align*}\newline
y \glsni{ampchar}= 2x +
2\glsni{tab.dbbackslashchar}\newline
\glsni{intertext}\marg{Using the distributive law:}\newline
  \glsni{ampchar}= 2(x+1)\newline
\glsni{end}\marg{align*}
\end{code}
\begin{result}[intertext.html]
\begin{align*}
y &= 2x + 2\\
\intertext{Using the distributive law:}
  &= 2(x+1)
\end{align*}
\end{result}\bookpagebreak

There are other environments for multiple-line displayed maths, but
they are beyond the scope of this book. See
the \isty{amsmath} documentation for further details.

%%%%%%%%%%%%%%%% Mathematical Commands %%%%%%%%%%%%%%%%%%%%

\setnode{mathscom}
\section{Mathematical Commands}
\label{sec:mathscom}

Most of the \glspl{command} described in this
section may only be used in one of the mathematics environments.  If
you try to use a mathematics command outside a maths environment you
will get a \dq{\texttt{Missing \$ inserted}} error message.

%------------------ Maths Fonts ------------------

\setnode{mathsfonts}
\subsection{Maths Fonts}
\label{sec:mathsfonts}

Just as we are able to \htmlref{change text fonts}{sec:fonts} using
the commands \gls{textrm}, \gls{textbf} etc, we can also
use commands to change the maths font.  Basic maths font changing
commands are shown in \tableref{tab:mathsfonts}.

\begin{table}[tbhp]
\caption{Maths Font Changing Commands}
\label{tab:mathsfonts}
\centering
\begin{tabular}{lll}
\toprule
\bfseries Command & \bfseries Example Input & \bfseries Corresponding Output\\
 & & (Computer Modern)\\
\midrule
\begin{inlinedef}\gls{mathrm}\marg{\meta{maths}}\end{inlinedef} &
\verb|$\mathrm{xyz}$| & $\cmmathrm{xyz}$\\
\begin{inlinedef}\gls{mathsf}\marg{\meta{maths}}\end{inlinedef} &
\verb|$\mathsf{xyz}$| & $\cmmathsf{xyz}$\\
\begin{inlinedef}\gls{mathtt}\marg{\meta{maths}}\end{inlinedef} &
\verb|$\mathtt{xyz}$| & $\cmmathtt{xyz}$\\
\begin{inlinedef}\gls{mathit}\marg{\meta{maths}}\end{inlinedef} &
\verb|$\mathit{xyz}$| & $\cmmathit{xyz}$\\
\begin{inlinedef}\gls{mathbf}\marg{\meta{maths}}\end{inlinedef} &
\verb|$\mathbf{xyz}$| & $\cmmathbf{xyz}$\\
\begin{inlinedef}\gls{mathcal}\marg{\meta{maths}}\end{inlinedef} &
\verb|$\mathcal{XYZ}$| & $\cmmathcal{XYZ}$
\\\bottomrule
\end{tabular}
\end{table}

The calligraphic fonts via \glsni{mathcal} are only available for upper-case 
characters\faq{Better script fonts for maths}{scriptfonts}.
\tableref{tab:amsmathfont} lists additional font commands supplied 
with the \isty{amsmath} and \isty{amsfonts} 
\htmlref{packages}{sec:packages}.

\begin{table}[hbtp]
\caption[The amsfonts and amsmath Font Commands]{The 
\isty{amsfonts}\supddag\ and \isty{amsmath}\supdag\ Font Commands}
\label{tab:amsmathfont}
\centering
\begin{tabular}{lll}
\toprule
\bfseries Command & \bfseries Example Input & \bfseries Example Output\\
\midrule
\begin{inlinedef}\supddag\gls{mathbb}\marg{\meta{maths}}\end{inlinedef} &
\verb|$\mathbb{A+B=C}$| & $\mathbb{A+B=C}$\\
\begin{inlinedef}\supddag\gls{mathfrak}\marg{\meta{maths}}\end{inlinedef} &
\verb|$\mathfrak{A+B=C}$| & $\mathfrak{A+B=C}$\\
\begin{inlinedef}\supdag\gls{boldsymbol}\marg{\meta{maths}}\end{inlinedef} &
\verb|$\boldsymbol{A+B=C}$| & $\boldsymbol{A+B=C}$\\
\begin{inlinedef}\supdag\gls{pmb}\marg{\meta{symbol}}\end{inlinedef} &
\verb|$\pmb{+-=}$| & $\pmb{+-=}$
\\\bottomrule
\end{tabular}
\end{table}

%----------------- Greek Letters -----------------------

\setnode{greek}
\subsection{Greek Letters}
\label{sec:greek}

Greek letters that differ from the corresponding Roman letters are
obtained by placing a backslash in front of the name.\footnote{So,
for example, there is no omicron since it looks the same as a Roman
o.} Lower case and upper case Greek letters are shown in
\tableref{tab:greek} and \tableref{tab:Greek}, respectively. There are
also some variants of certain symbols, such as \glsni{vartheta} as
opposed to \glsni{theta}.

\begin{table}[htbp]
\caption{Lower Case Greek Letters}
\label{tab:greek}
\centering
\begin{tabular}{llllll}
\gls{alpha} & $\alpha$ &
\gls{beta} & $\beta$ &
\gls{gamma} & $\gamma$ \\
\gls{delta} & $\delta$  &
\gls{epsilon} & $\epsilon$ &
\gls{varepsilon} & $\varepsilon$\\
\gls{zeta} & $\zeta$ &
\gls{eta} & $\eta$ &
\gls{theta} & $\theta$ \\
\gls{vartheta} & $\vartheta$ &
\gls{iota} & $\iota$ & 
\gls{kappa} & $\kappa$ \\
\gls{lambda} & $\lambda$ &
\gls{mu} & $\mu$ &
\gls{nu} & $\nu$ \\
\gls{xi} & $\xi$ &
\gls{pi} & $\pi$ &
\gls{varpi} & $\varpi$\\
\gls{rho} & $\rho$ &
\gls{varrho} & $\varrho$ &
\gls{sigma} & $\sigma$ \\
\gls{varsigma} & $\varsigma$ &
\gls{tau} & $\tau$ &
\gls{upsilon} & $\upsilon$ \\
\gls{phi} & $\phi$ &
\gls{varphi} & $\varphi$ &
\gls{chi} & $\chi$ \\
\gls{psi} & $\psi$ &
\gls{omega} & $\omega$ 
\end{tabular}
\end{table}

\begin{table}[htbp]
\caption{Upper Case Greek Letters}
\label{tab:Greek}
\centering
\begin{tabular}{llllll}
\gls{Gamma} & $\Gamma$ &
\gls{Delta} & $\Delta$ &
\gls{Theta} & $\Theta$ \\
\gls{Lambda} & $\Lambda$ &
\gls{Xi} & $\Xi$ &
\gls{Pi} & $\Pi$ \\
\gls{Sigma} & $\Sigma$ &
\gls{Upsilon} & $\Upsilon$ &
\gls{Phi} & $\Phi$ \\
\gls{Psi} & $\Psi$ &
\gls{Omega} & $\Omega$
\end{tabular}
\end{table}

\xminisec{Example:}

The following code
\begin{codeS}
\glsni{begindispmath} x' = x + \glsni{Delta} x \glsni{enddispmath}
\end{codeS}
produces:
\begin{resultS}[x prime equals x plus capital Delta x]
\[ x' = x + \Delta x \]
\end{resultS}

%------------- Subscripts and Superscripts ------------

\setnode{subsupscripts}
\subsection{Subscripts and Superscripts}
\label{sec:scripts}

Subscripts are obtained either by the command
\begin{definition}
\gls{sb}\marg{\meta{maths}}
\end{definition}%
or by the \htmlref{special character}{sec:chars}:
\begin{definition}
\gls{underscorechar}\marg{\meta{maths}}
\end{definition}%
Superscripts are obtained either by the command
\begin{definition}
\gls{sp}\marg{\meta{maths}}
\end{definition}%
or by the special character:
\begin{definition}
\gls{circumchar}\marg{\meta{maths}}
\end{definition}%

\xminisec{Examples:}
\begin{enumerate}
\item  This example uses \glsni{sb} and \glsni{sp}:
\begin{codeS}
\glsni{begindispmath}y = x\glsni{sb}\marg{1}\glsni{sp}\marg{2} +
x\glsni{sb}\marg{2}\glsni{sp}\marg{2}\glsni{enddispmath}
\end{codeS}%

\item This example uses \glsni{underscorechar} and \glsni{circumchar}
\begin{codeS}
\glsni{begindispmath}y = x\glsni{underscorechar}\marg{1}\glsni{circumchar}\marg{2}
 + x\glsni{underscorechar}\marg{2}\glsni{circumchar}\marg{2}\glsni{enddispmath}
\end{codeS}%

\item Recall from 
\latexhtml{page~\pageref{itm:argnogrp}}{\htmlref{earlier}{itm:argnogrp}} that
\glspl{mandatory} only consisting of one
character don't need to be grouped, so the above code can also be
written as:
\begin{codeS}
\glsni{begindispmath}y = x\glsni{underscorechar}1\glsni{circumchar}2 
 + x\glsni{underscorechar}2\glsni{circumchar}2\glsni{enddispmath}
\end{codeS}%
This is simpler than the first two examples. However it's a good idea to be in the
habit of always using braces in case you forgot them when they're
needed.
\end{enumerate}
All three of the above examples produce the same output:
\begin{resultS}[subsup.html]
\[y = x_1^2 + x_2^2\]
\end{resultS}%

Notice how the subscript gets tucked under the slope of the $Y$ in:
\begin{codeS}
\glsni{begindispmath}
Y\glsni{underscorechar}\marg{1}\glsni{circumchar}\marg{2}
\glsni{enddispmath}
\end{codeS}
\begin{resultS}[ysubsup1.html]
\[ Y_{1}^{2} \]
\end{resultS}
Compare with
\begin{codeS}
\glsni{begindispmath}
Y\marg{}\glsni{underscorechar}\marg{1}\glsni{circumchar}\marg{2}
\glsni{enddispmath}
\end{codeS}
\begin{resultS}[ysubsup2.html]
\[ Y{}_{1}^{2} \]

\end{resultS}

\xminisec{Example (Nested)}

Subscripts and superscripts can also be nested (note that
it is now necessary to group the argument to the superscript command):
\begin{codeS}
\glsni{begindispmath}
 f(x) = e\glsni{circumchar}\marg{x\glsni{underscorechar}1}
\glsni{enddispmath}
\end{codeS}%
which produces
\begin{resultS}[nestedsup.html]
\[
f(x) = e^{x_1}
\]
\end{resultS}

This example isn't quite right as e isn't actually a variable
and shouldn't be typeset in italic.  The correct way to do this
is:\bookpagebreak
\begin{codeS}
\glsni{begindispmath}
 f(x) =
\gls{mathrm}\marg{e}\glsni{circumchar}\marg{x\glsni{underscorechar}1}
\glsni{enddispmath}
\end{codeS}%
which results in:
\begin{resultS}[Image: as above but the letter e is upright]
\[ f(x) = \mathrm{e}^{x_1} \]
\end{resultS}

\refstepcounter{object}\label{obj:e}If you are going to use e a lot,
it will be simpler to \htmlref{define a new command}{ch:newcom} to do
this. The definition should go in the
\glsterm{preamble}:
\begin{codeS}
\gls{newcommand}\marg{\cmdname{e}}\marg{\glsni{mathrm}\marg{e}}
\end{codeS}%
Then in the document:
\begin{codeS}
\glsni{begindispmath} f(x\glsni{underscorechar}1, x\glsni{underscorechar}2) 
 = \cmdname{e}\glsni{circumchar}\marg{x\glsni{underscorechar}1\glsni{circumchar}2} 
  +
\cmdname{e}\glsni{circumchar}\marg{x\glsni{underscorechar}2\glsni{circumchar}2}
\glsni{enddispmath}
\end{codeS}%
\begin{resultS}[nestedsup2.html]
\[
f(x_1, x_2) = \e^{x_1^2} + \e^{x_2^2}
\]
\end{resultS}

Take care when nesting subscripts or superscripts. The following
\begin{alltt}\wrong
x\glsni{underscorechar}1\glsni{underscorechar}2
\end{alltt}
will give a \texttt{!\ Double subscript} error.

%-------------- Functional Names --------------

\setnode{functionnames}
\subsection{Functional Names}
\label{sec:funcnames}

Functions such as log and tan can't simply be typed in as
\texttt{log} or \texttt{tan} otherwise they will come out looking
like the variables \mathorit{l} times
\mathorit{o} times \mathorit{g} (\makeimg{l o g}{$ l o g$})
or \mathorit{t} times \mathorit{a} times \mathorit{n} 
(\makeimg{t a n}{$ t a n$}).  Instead you should use
one of the commands listed in \tableref{tab:fnnames}.
The functions denoted with \supdag\ can have limits by using
the subscript command \gls{underscorechar} or the superscript command \gls{circumchar}\faq{Sub- and
superscript positioning for operators}{limits}.
In addition, the modulo commands listed in \tableref{tab:mod} are also available.

\begin{table}[bhtp]
\caption[Function Names]{Function Names (\supdag
indicates command may have limits, \supddag defined
by \protect\isty{amsmath}).}
\label{tab:fnnames}
\centering
\begin{tabular}{llllll}
\gls{arccos} & $\arccos$ &
\gls{arcsin} & $\arcsin$ &
\gls{arctan} & $\arctan$ \\
\gls{arg} & $\arg$ &
\gls{cos} & $\cos$ &
\gls{cosh} & $\cosh$ \\
\gls{cot} & $\cot$ &
\gls{coth} & $\coth$ &
\gls{csc} & $\csc$ \\
\gls{deg} & $\deg$ &
\gls{det}\supdag & $\det$ &
\gls{dim} & $\dim$ \\
\gls{exp} & $\exp$ &
\gls{gcd}\supdag & $\gcd$ &
\gls{hom} & $\hom$ \\
\gls{inf}\supdag & $\inf$ &
\gls{injlim}\supdagddag & $\injlim$ &
\gls{ker} & $\ker$ \\
\gls{lg} &  $\lg$ &
\gls{lim}\supdag & $\lim$ &
\gls{liminf}\supdag & $\liminf$ \\
\gls{limsup}\supdag & $\limsup$ &
\gls{ln} & $\ln$ &
\gls{log} & $\log$ \\
\gls{max}\supdag &  $\max$ &
\gls{min}\supdag & $\min$ &
\gls{Pr}\supdag & $\Pr$ \\
\gls{projlim}\supdagddag & $\projlim$ &
\gls{sec} & $\sec$ &
\gls{sin} & $\sin$ \\
\gls{sinh} & $\sinh$ &
\gls{sup}\supdag & $\sup$ &
\gls{tan} & $\tan$ \\
\gls{tanh} & $\tanh$ &
\gls{varinjlim}\supdagddag & $\varinjlim$ &
\gls{varliminf}\supdagddag & $\varliminf$ \\
\gls{varlimsup}\supdagddag & $\varlimsup$ &
\gls{varprojlim}\supdagddag &  $\varprojlim$ &
\end{tabular}
\end{table}

\begin{table}[htbp]
\caption[Modulo Commands]{Modulo Commands (\supddag
defined by \protect\isty{amsmath} package)}
\label{tab:mod}
\centering
\begin{tabular}{lll}
\toprule
\bfseries Command & \bfseries Example Input & \bfseries Example Output\\
\midrule
\gls{bmod} & \verb|$m \bmod n$| & $m \bmod n$\\
\gls{pmod}\marg{\meta{maths}} & \verb|$m \pmod{n}$| &$m \pmod{n}$\\
\gls{mod}\marg{\meta{maths}}\supddag & \verb|$m \mod{n}$| &$m \mod{n}$\\
\gls{pod}\marg{\meta{maths}}\supddag & \verb|$m \pod{n}$| &$m \pod{n}$
\\\bottomrule
\end{tabular}
\end{table}

\xminisec{Example (Trigonometric Functions):}

This example uses the cos and sin functions and also the 
\htmlref{Greek letter}{sec:greek} theta.
\begin{codeS}
\glsni{begindispmath} z = r(\glsni{cos}\glsni{theta} + i\glsni{sin}\glsni{theta})
\glsni{enddispmath}
\end{codeS}%
\begin{resultS}[Image: z equals r (cos theta plus i sign theta)]
\[
z = r(\cos\theta + i\sin\theta)
\]
\end{resultS}

\xminisec{Example (Limit):}

The command \gls{infty} is the 
infinity symbol \makeimg{image of infinity symbol}{$\infty$}, and the command \gls{to} displays an 
arrow pointing to the right.  Note the use of \gls{underscorechar} since the
limit is a subscript.
\begin{codeS}
\glsni{begindispmath}
  \glsni{lim}\glsni{underscorechar}\marg{x\glsni{to}\glsni{infty}} f(x) 
\glsni{enddispmath}
\end{codeS}%
\begin{resultS}[Image: limit as x tends to infinity of f of x]
\[
\lim_{x\to\infty} f(x)
\]
\end{resultS}%
The operators with limits behave differently depending on whether
they are in \htmlref{displayed}{sec:displayed} or 
\htmlref{in-line}{sec:inline} maths.
Notice the difference when the same code appears in-line:
\begin{codeS}
In a line of text 
\glsni{dollarchar}\glsni{lim}\glsni{underscorechar}\marg{x\glsni{to}\glsni{infty}}
f(x)\glsni{dollarchar}
\end{codeS}%
which now displays as:
\begin{resultS}[Image: the limit appears to the side of lim instead
of below it.]
In a line of text $\lim_{x\to\infty} f(x)$
\end{resultS}

\xminisec{Example (With Subscript):}

This is another example of a functional name using a subscript:
\begin{codeS}
\glsni{begindispmath} \glsni{min}\glsni{underscorechar}x f(x) \glsni{enddispmath}
\end{codeS}%
\begin{resultS}[Image: the x appears below the word min.]
\[
\min_x f(x)
\]
\end{resultS}%
Again, notice the difference when it is used in-line:
\begin{codeS}
In a line of text \glsni{dollarchar}\glsni{min}\glsni{underscorechar}x 
f(x)\glsni{dollarchar}
\end{codeS}%
\begin{resultS}[Image: the x appears to the side of the word min.]
In a line of text $\min_x f(x)$
\end{resultS}
\bookpagebreak

\setnode{declaremathop}
\subsubsection{Defining New Functional Operators}
\label{sec:declaremathop}

It may be that you want a function that isn't specified in
\tableref{tab:fnnames}. In this case, the \isty{amsmath} provides the
\textbf{\em \gls{preamble} only} command
\begin{definition}
\gls{DeclareMathOperator}\marg{\meta{cmd}}\marg{\meta{operator
name}}
\end{definition}
or its \htmlref{starred variant}{itm:starredcommand}
\begin{definition}
\glsni{DeclareMathOperator}*\marg{\meta{cmd}}\marg{\meta{operator
name}}
\end{definition}
\faq{Defining a new log-like function in LaTeX}{newfunction}%
Both versions define a command called \meta{cmd}, which must start
with a backslash, that typesets \meta{operator name} as a function
name.  The starred version is for function names that can take
limits (like \gls{lim} and \gls{min} described above).

\xminisec{Example (Operator Without Limits):}

Suppose I want a function called card, which represents the
cardinality of a set $\mathcal{S}$. First I~need to define the new
operator command (which I'm going to call \cmdname{card}) \emph{in
the \gls{preamble}}:
\begin{codeS}
\glsni{DeclareMathOperator}\marg{\cmdname{card}}\marg{card}
\end{codeS}
This operator doesn't take any limits, so I~have used the unstarred
version.

Later in the document, I~can use this new operator command:
\begin{codeS}
\glsni{begindispmath} n = \cmdname{card}(\glsni{mathcal}\marg{S}) \glsni{enddispmath}
\end{codeS}%
\begin{resultS}[Image: n equals card (S)]
\[
n = \card(\mathcal{S})
\]
\end{resultS}%
In this example \gls{mathcal} is used as sets are
typically represented in a calligraphic font.

\xminisec{Example (Operator With Limits):}

Suppose I~now want a function called mode, which represents the
mode of a set of numbers. First, I~define the operator command in the
preamble:
\begin{codeS}
\glsni{DeclareMathOperator}*\marg{\cmdname{mode}}\marg{mode}
\end{codeS}
This operator needs to be able to have a subscript, so I~have used
the starred version.

Later in the document, I~can use this new operator command:
\begin{codeS}
\glsni{begindispmath} x\glsni{underscorechar}m = 
\cmdname{mode}\glsni{underscorechar}\marg{x \glsni{in} \glsni{mathcal}\marg{S}}(x)
\glsni{enddispmath}
\end{codeS}
\begin{resultS}[Image: x subscript m equals mode of x in S of x]
\[
x_m = \mode_{x \in \mathcal{S}}(x)
\]
\end{resultS}

%----------- Fractions ------------

\bookpagebreak
\setnode{fractions}
\subsection{Fractions}
\label{sec:fractions}

Fractions are created using the command
\begin{definition}
\gls{frac}\marg{\meta{numerator}}\marg{\meta{denominator}}
\end{definition}
The \isty{amsmath} package also provides the command
\begin{definition}
\gls{cfrac}\oarg{\meta{pos}}\marg{\meta{numerator}}\marg{\meta{denominator}}
\end{definition}
which is designed for continued fractions. The optional argument
\texttt{pos} can be used for left (\texttt{l}) or right (\texttt{r})
placement of any of the numerators. (The default is centred.)

\xminisec{Example:}

A simple fraction:
\begin{codeS}
\glsni{begindispmath} 
 \glsni{frac}\marg{1}\marg{1+x} 
\glsni{enddispmath}
\end{codeS}
Produces:
\begin{resultS}[Image of fraction with 1 as the numerator and 1+x as
the denominator.]
\[ \frac{1}{1+x} \]
\end{resultS}
Compare with:
\begin{codeS}
In-line:
\glsni{dollarchar} 
 \glsni{frac}\marg{1}\marg{1+x} 
\glsni{dollarchar}
\end{codeS}
which produces:
\begin{resultS}[Image: as before but fraction is smaller.]
In-line: $\frac{1}{1+x}$
\end{resultS}

\xminisec{Example (Nested):}
\begin{codeS}
\glsni{begindispmath} 
\glsni{frac}\marg{1+\glsni{frac}\marg{1}\marg{x}}\marg{1+x+x\glsni{circumchar}2}
\glsni{enddispmath}
\end{codeS}%
\begin{resultS}[Image of fraction with (1+1/x) as the numerator and
(1 + x + x squared as the denominator.]
\[
\frac{1+\frac{1}{x}}{1+x+x^2}
\]
\end{resultS}

\xminisec{Example (Continued Fraction);}
A continued fraction (example taken from \isty{amsmath}
documentation and uses \glsni{sqrt}, described in \sectionref{sec:roots},
and \glsni{dotsb}, described in
\sectionref{sec:mathssym}):
\begin{code}
\begin{alltt}
\glsni{begindispmath}
\glsni{cfrac}\marg{1}\marg{\glsni{sqrt}\marg{2}+
\glsni{cfrac}\marg{1}\marg{\glsni{sqrt}\marg{2}+
\glsni{cfrac}\marg{1}\marg{\glsni{sqrt}\marg{2}+\glsni{dotsb}
}}}
\glsni{enddispmath}
\end{alltt}
\end{code}\screenpagebreak
\begin{result}[Image of a continued fraction (the denominator has a
fraction with a fraction in its denominator etc).]
\[
\cfrac{1}{\sqrt{2}+
\cfrac{1}{\sqrt{2}+
\cfrac{1}{\sqrt{2}+\dotsb
}}}
\]
\end{result}

\xminisec{Example (A Derivative):}
\begin{codeS}
\glsni{begindispmath} f'(x) = \glsni{frac}\marg{df}\marg{dx} \glsni{enddispmath}
\end{codeS}%
\begin{resultS}[Image: f'(x) = d f by d x]
\[
f'(x) = \frac{df}{dx}
\]
\end{resultS}%
As with \dq{e}, the differential operator \dq{d} should
be in an upright font as it is not a variable:
\begin{code}
\begin{alltt}
\glsni{begindispmath}
  f'(x) = \glsni{frac}\marg{\glsni{mathrm}\marg{d}f}\marg{\glsni{mathrm}\marg{d}x}
\glsni{enddispmath}
\end{alltt}
\end{code}%
\begin{resultS}[Image: as above but the letter d is upright]
\[
f'(x) = \frac{\mathrm{d}f}{\mathrm{d}x}
\]
\end{resultS}

The above example is rather cumbersome, particularly
if you have a lot of derivatives, so it might be easier to
\htmlref{define a new command}{ch:newcom}\latex{ (see \chapterref{ch:newcom})}. In the 
\glsterm{preamble} define:
\begin{codeS}
\gls{newcommand}\marg{\cmdname{deriv}}\oarg{2}\marg{\glsni{frac}\marg{\glsni{mathrm}\marg{d}\glsni{hashchar}1}\marg{\glsni{mathrm}\marg{d}\glsni{hashchar}2}}
\end{codeS}%
Then in the document:
\begin{codeS}
\glsni{begindispmath} f'(x) = \cmdname{deriv}\marg{f}\marg{x} \glsni{enddispmath}
\end{codeS}%
\begin{resultS}[Image: same as before]
\newcommand{\deriv}[2]{\frac{\mathrm{d}#1}{\mathrm{d}#2}}

\[ f'(x) = \deriv{f}{x} \]
\end{resultS}

\xminisec{Example (Partial Derivative):}
\label{itm:pderiv}%

Partial derivatives can be obtained similarly
using the command \gls{partial} to display the partial
derivative symbol. As in the previous example, first define a new
command to format a partial derivative in the
\glsterm{preamble}:
\begin{codeS}
\gls{newcommand}\marg{\cmdname{pderiv}}\oarg{2}\marg{\glsni{frac}\marg{\glsni{partial} \glsni{hashchar}1}{\glsni{partial} \glsni{hashchar}2}}
\end{codeS}%
Then in the document:
\begin{codeS}
\glsni{begindispmath} f\glsni{underscorechar}x = \cmdname{pderiv}\marg{f}\marg{x} \glsni{enddispmath}
\end{codeS}%
\begin{resultS}[Image: f subscript x = partial d f by partial d x]
\[
f_x = \pderiv{f}{x}
\]
\end{resultS}

\xminisec{Example (Double Partial Derivative):}

\begin{code}
\begin{alltt}
\glsni{begindispmath}
  f\glsni{underscorechar}\marg{xy} = \glsni{frac}\marg{\glsni{partial}\glsni{circumchar}2 f}\marg{\glsni{partial} x \glsni{partial} y}
\glsni{enddispmath}
\end{alltt}
\end{code}%
\begin{resultS}[Image: f subscript xy = partial d 2 f by partial d x partial d y.]
\[
f_{xy} = \frac{\partial^2 f}{\partial x \partial y}
\]
\end{resultS}

\xminisec{Example (First principles):}
\begin{code}
\glsni{begindispmath}\newline
\mbox{}\qquad f'(x) = \gls{lim}\gls{underscorechar}\marg{\gls{Delta} x \gls{to}
0}\newline
\mbox{}\qquad \glsni{frac}\marg{f(x + \gls{Delta} x)-f(x)}\marg{\gls{Delta}
x}\newline
\glsni{enddispmath}
\end{code}
\begin{resultS}
\[
f'(x) = \lim_{\Delta x\to0}
  \frac{f(x+\Delta x) - f(x)}{\Delta x}
\]
\end{resultS}

%------------ Roots --------------

\setnode{roots}
\subsection{Roots}
\label{sec:roots}

Roots are obtained using the command
\begin{definition}
\gls{sqrt}\oarg{\meta{order}}\marg{\meta{maths}}
\end{definition}%
without the optional argument \meta{order}
it will produce a simple square root. Cubic roots etc
can be obtained using the optional argument.

\xminisec{Examples:}

\begin{enumerate}
\item A square root:
\begin{codeS}
\glsni{begindispmath} \glsni{sqrt}\marg{a+b} \glsni{enddispmath}
\end{codeS}%
\begin{resultS}[Image: square root of (a + b)]
\[
\sqrt{a+b}
\]
\end{resultS}

\item A cubic root:
\begin{codeS}
\glsni{begindispmath} \glsni{sqrt}\oarg{3}\marg{a+b} \glsni{enddispmath}
\end{codeS}%
\begin{resultS}[Image: cubic root of (a + b)]
\[
\sqrt[3]{a+b}
\]
\end{resultS}

\item An \ensuremath{n}th root:
\begin{codeS}
\glsni{begindispmath} \glsni{sqrt}\oarg{n}\marg{a+b} \glsni{enddispmath}
\end{codeS}%
\begin{resultS}[Image: the nth root of (a + b)]
\[
\sqrt[n]{a+b}
\]
\end{resultS}
\end{enumerate}

%-------- Mathematical Symbols -------------

\setnode{mathssym}
\subsection{Mathematical Symbols}
\label{sec:mathssym}

Relational symbols\faq{Where can I find the symbol for
\ldots}{symbols} are shown in \tableref{tab:relsym}.  If you want a
negation that is not shown, you can obtain it by preceding the
symbol with the command \gls{not}.  For example:
\glsni{not}\glsni{subset} produces the symbol 
\makeimg{image of subset symbol with slash through it}{$\not\subset$}.

\begin{table}[hbtp]
\caption{Relational Symbols}
\label{tab:relsym}
\centering
\begin{tabular}{llllll}
\gls{approx} & $\approx$ &
\gls{asymp} & $\asymp$ &
\gls{bowtie} & $\bowtie$ \\
\gls{cong} & $\cong$ &
\gls{dashv} & $\dashv$ &
\gls{doteq} & $\doteq$ \\
\gls{equiv} & $\equiv$ &
\gls{frown} & $\frown$ &
\gls{ge} or \gls{geq} & $\geq$ \\
\gls{gg} & $\gg$ &
\gls{in} & $\in$ &
\gls{le} or \gls{leq} & $\leq$ \\
\gls{ll} & $\ll$ &
\gls{mid} or \verb"|" & $\mid$ &
\gls{models} & $\models$ \\ 
\gls{neq} & $\neq$ &
\gls{ni} & $\ni$ &
\gls{notin} & $\notin$ \\
\gls{parallel} & $\parallel$ &
\gls{prec}& $\prec$ &
\gls{preceq}& $\preceq$ \\
\gls{perp} & $\perp$ &
\gls{propto} & $\propto$ &
\gls{sim} & $\sim$ \\
\gls{simeq} & $\simeq$ &
\gls{smile} & $\smile$ &
\gls{sqsubseteq} & $\sqsubseteq$ \\
\gls{sqsupseteq} & $\sqsupseteq$ &
\gls{subset} & $\subset$ &
\gls{subseteq} & $\subseteq$ \\
\gls{succ} & $\succ$ &
\gls{succeq} & $\succeq$ &
\gls{supset} & $\supset$ \\
\gls{supseteq} & $\supseteq$ &
\gls{vdash} & $\vdash$ &
\end{tabular}
\end{table}

Binary operator symbols are shown in \tableref{tab:binop}, and arrow
symbols are shown in \xtableref{tab:arrows}. There are also over and
under arrows (\xtableref{tab:overunderarrows}) that have an argument.
The over arrows put an extendible arrow over their argument, and the
under arrows put an extendible arrow under their argument. In
addition, the \isty{amsmath} package provides extensible arrows that
take a superscript and, optionally, a subscript:
\begin{definition}
\gls{xleftarrow}\oarg{\meta{subscript}}\marg{\meta{superscript}}
\end{definition}
\begin{definition}
\gls{xrightarrow}\oarg{\meta{subscript}}\marg{\meta{superscript}}
\end{definition}

\xminisec{Example:}
\begin{code}
\glsni{begindispmath}\newline
\mbox{}\qquad  A \glsni{xleftarrow}\marg{n+m-p} B  
 \glsni{xrightarrow}\oarg{X}\marg{n+p} C\newline
\glsni{enddispmath}
\end{code}
\begin{result}[Image of result: the arrow subscripts and superscript are centred over the arrow.]
\[
  A \xleftarrow{n+m-p} B \xrightarrow[X]{n+p} C
\]
\end{result}

\begin{table}[tbhp]
\caption{Binary Operator Symbols}
\label{tab:binop}
\centering
\begin{tabular}{llllll}
\gls{amalg} & $\amalg$ &
\gls{ast} & $\ast$ &
\gls{bullet} & $\bullet$ \\
\gls{bigcirc} & $\bigcirc$ &
\gls{bigtriangledown} & $\bigtriangledown$ &
\gls{bigtriangleup} & $\bigtriangleup$ \\
\gls{cap} & $\cap$ &
\gls{cdot} & $\cdot$ &
\gls{circ} & $\circ$ \\
\gls{cup} & $\cup$ &
\gls{dagger} & $\dagger$ &
\gls{ddagger} & $\ddagger$ \\
\gls{diamond} & $\diamond$ &
\gls{div} & $\div$ &
\gls{mp} & $\mp$ \\
\gls{odot} & $\odot$ &
\gls{ominus} & $\ominus$ &
\gls{oplus} & $\oplus$ \\
\gls{oslash} & $\oslash$ &
\gls{otimes} & $\otimes$ &
\gls{pm} & $\pm$ \\
\gls{setminus} & $\setminus$ &
\gls{sqcap} & $\sqcap$ &
\gls{sqcup} & $\sqcup$ \\
\gls{star} & $\star$ &
\gls{times} & $\times$ &
\gls{triangleleft} & $\triangleleft$ \\
\gls{triangleright} & $\triangleright$ &
\gls{uplus} & $\uplus$ &
\gls{vee} & $\vee$ \\
\gls{wedge} & $\wedge$ &
\gls{wr} & $\wr$
\end{tabular}
\end{table}

\begin{table}[tbhp]
\caption{Arrow Symbols}
\label{tab:arrows}
\centering
\begin{tabular}{llll}
\gls{downarrow} & $\downarrow$ &
\gls{Downarrow} & $\Downarrow$ \\
\gls{hookleftarrow} & $\hookleftarrow$ &
\gls{hookrightarrow} & $\hookrightarrow$ \\
\gls{leftarrow} or \gls{gets} & $\leftarrow$ &
\gls{Leftarrow} & $\Leftarrow$ \\
\gls{leftharpoondown} & $\leftharpoondown$ &
\gls{leftharpoonup} & $\leftharpoonup$ \\
\gls{leftrightarrow} & $\leftrightarrow$ &
\gls{Leftrightarrow} & $\Leftrightarrow$ \\
\gls{longleftarrow} & $\longleftarrow$ &
\gls{Longleftarrow} & $\Longleftarrow$ \\
\gls{longleftrightarrow} & $\longleftrightarrow$ &
\gls{Longleftrightarrow} & $\Longleftrightarrow$ \\
\gls{longmapsto} & $\longmapsto$ &
\gls{longrightarrow} & $\longrightarrow$ \\
\gls{Longrightarrow} & $\Longrightarrow$ &
\gls{mapsto} & $\mapsto$ \\
\gls{nearrow} & $\nearrow$ &
\gls{nwarrow} & $\nwarrow$ \\
\gls{rightarrow} or \gls{to} & $\to$ &
\gls{Rightarrow} & $\Rightarrow$ \\
\gls{rightharpoondown} & $\rightharpoondown$ &
\gls{rightharpoonup} & $\rightharpoonup$ \\
\gls{rightleftharpoons} & $\rightleftharpoons$ &
\gls{searrow} & $\searrow$ \\
\gls{swarrow} & $\swarrow$ &
\gls{uparrow} & $\uparrow$ \\
\gls{Uparrow} & $\Uparrow$ &
\gls{updownarrow} & $\updownarrow$ \\
\gls{Updownarrow} & $\Updownarrow$ &
\end{tabular}
\end{table}

\begin{table}[htbp]
\caption[Over and Under Arrows]{Over and Under Arrows
(\supdag defined by \protect\isty{amsmath})}
\label{tab:overunderarrows}
\centering
\begin{tabular}{@{}lll@{}}
\toprule
\bfseries Definition & \multicolumn{2}{c}{\bfseries Example}\\
\midrule
\begin{inlinedef}\gls{overleftarrow}\marg{\meta{maths}}\end{inlinedef}
  & \verb|\overleftarrow{ABC}| & $\overleftarrow{ABC}$\\
\begin{inlinedef}\gls{overrightarrow}\marg{\meta{maths}}\end{inlinedef}
  & \verb|\overrightarrow{ABC}| & $\overrightarrow{ABC}$\\
\begin{inlinedef}\gls{overleftrightarrow}\marg{\meta{maths}}\end{inlinedef}\supdag
  & \verb|\overleftrightarrow{ABC}| & $\overleftrightarrow{ABC}$\\
\begin{inlinedef}\gls{underleftarrow}\marg{\meta{maths}}\end{inlinedef}\supdag
  & \verb|\underleftarrow{ABC}| & $\underleftarrow{ABC}$\\
\begin{inlinedef}\gls{underrightarrow}\marg{\meta{maths}}\end{inlinedef}\supdag
  & \verb|\underrightarrow{ABC}| & $\underrightarrow{ABC}$\\
\begin{inlinedef}\gls{underleftrightarrow}\marg{\meta{maths}}\end{inlinedef}\supdag
  & \verb|\underleftrightarrow{ABC}| & $\underleftrightarrow{ABC}$
\\\bottomrule
\end{tabular}
\end{table}

\bookpagebreak
\begin{table}[htbp]
\caption{Symbols with Limits}
\label{tab:symlim}
\centering
\begin{tabular}{llllll}
\gls{sum} & $\sum$ &
\gls{int} & $\int$ &
\gls{oint} & $\oint$ \\[5pt]
\gls{prod} & $\prod$ &
\gls{coprod} & $\coprod$ &
\gls{bigcap} & $\bigcap$ \\[5pt]
\gls{bigcup} & $\bigcup$ &
\gls{bigsqcup} & $\bigsqcup$ &
\gls{bigvee} & $\bigvee$ \\[5pt]
\gls{bigwedge} & $\bigwedge$ &
\gls{bigodot} & $\bigodot$ &
\gls{bigotimes} & $\bigotimes$ \\[5pt]
\gls{bigoplus} & $\bigoplus$ &
\gls{biguplus} & $\biguplus$
\end{tabular}
\end{table}

Symbols that can have limits are shown in \tableref{tab:symlim}.
The size of these symbols depends on whether they are in displayed 
maths or in-line maths. 

\xminisec{Example (Displayed Summation and Product):}

The limits of summations and products are placed above and below the
symbol in displayed maths:
\begin{code}
\begin{alltt}
\glsni{begindispmath}
  f(x) = \glsni{sum}\glsni{underscorechar}\marg{i=1}\glsni{circumchar}\marg{n} x\glsni{underscorechar}i + \glsni{prod}\glsni{underscorechar}\marg{i=1}\glsni{circumchar}\marg{n} x\glsni{underscorechar}i
\glsni{enddispmath}
\end{alltt}
\end{code}
\begin{resultS}[Image: f(x) = sum from i equals 1 to n x subscript i
+ product from i equals 1 to n x subscript i.]
\[
f(x) = \sum_{i=1}^{n} x_i + \prod_{i=1}^{n} x_i
\]
\end{resultS}

\xminisec{Example (In-line Summation and Product):}

The limits of summations and products are placed to the right of the
symbol in in-line maths:
\begin{code}
\begin{alltt}
In a line of text:
\glsni{begin}\marg{math}
  f(x) = \glsni{sum}\glsni{underscorechar}\marg{i=1}\glsni{circumchar}\marg{n} x\glsni{underscorechar}i + \glsni{prod}\glsni{underscorechar}\marg{i=1}\glsni{circumchar}\marg{n} x\glsni{underscorechar}i
\glsni{end}\marg{math}
\end{alltt}
\end{code}
\begin{resultS}[Image: as before but summation and product symbols
are smaller and limits are to the side instead of above and below]
In a line of text:
\begin{math}
f(x) = \sum_{i=1}^{n} x_i + \prod_{i=1}^{n} x_i
\end{math}
\end{resultS}

\xminisec{Multiline Sub- or Superscripts}

The \isty{amsmath} package provides the command:
\begin{definition}
\gls{substack}\marg{\meta{maths}}
\end{definition}
which can be used for multiline sub- or superscripts. Within the
argument \meta{maths} use \gls{tab.dbbackslashchar} to separate
rows. For example:
\begin{code}
\glsni{begindispmath}\newline
\strut~\glsni{sum}\glsni{underscorechar}\marg{\newline
\strut~~\glsni{substack}\newline
\strut~~\marg{\newline
\strut~~~~i \glsni{in}
\glsni{mathcal}\marg{I}\glsni{tab.dbbackslashchar}\newline
\strut~~~~i \glsni{neq} 0\newline
\strut~~}\newline
\strut~}\newline
\strut~x\glsni{underscorechar}i\newline
\glsni{enddispmath}
\end{code}
\begin{result}[substack.html]
\[
 \sum_{
   \substack
   {
     i \in \mathcal{I}\\
     i \neq 0
   }
 }
 x_i
\]
\end{result}

\setnode{ellipses}
\subsection{Ellipses}
\label{sec:ellipses}

Ellipsis (omission mark)\index{ellipses (omission marks)|hyperpage} commands are 
shown in \tableref{tab:ellipsis}. The \isty{amsmath} package also
provides: \glsni{dotsc} for dots with commas, \glsni{dotsb} for dots
with binary operators/relations, \glsni{dotsm} for multiplication
dots, \glsni{dotsi} for dots with integrals and \glsni{dotso} for other
dots, which can be used as replacements for \glsni{ldots} and
\glsni{cdots}.

\begin{table}[hbtp]
\caption[Ellipses]{Ellipses (\supdag\ provided by \sty{amsmath}
package)}
\label{tab:ellipsis}
\centering
\begin{tabular}{ll@{\hspace{4\tabcolsep}}ll@{\hspace{4\tabcolsep}}llllll}
\gls{vdots} & $\vdots$ &
\gls{cdots} & $\cdots$ &
\gls{dotsb}\supdag & $\dotsb$ &
\gls{dotsi}\supdag & $\dotsi$ &
\gls{dotsm}\supdag & $\dotsm$\\
\gls{ddots} & $\ddots$ &
\gls{ldots} & $\ldots$ &
\gls{dotsc}\supdag & $\dotsc$ &
\gls{dotso}\supdag & $\dotso$ 
\end{tabular}
\end{table}

\xminisec{Example (Low Ellipsis):}

This example uses the command 
\refstepcounter{object}\label{cmd:forall}\gls{forall} to produce
the \dq{for all} symbol $\forall$, and it also uses 
\glsni{spacesym} (backslash space) to make a space before the for all symbol. The
\isty{amsmath} \dq{dots with commas} ellipsis \glsni{dotsc} is used rather
than the standard \glsni{ldots}:
\begin{code}
\begin{alltt}
\glsni{begindispmath}
a\glsni{underscorechar}ix\glsni{underscorechar}i = b\glsni{underscorechar}i\glsni{spacesym}\glsni{forall} i = 1,\glsni{dotsc}, n
\glsni{enddispmath}
\end{alltt}
\end{code}%
\begin{resultS}[forall.html]
\[
a_ix_i = b_i\ \forall i = 1,\dotsc, n
\]
\end{resultS}

\xminisec{Example (Centred ellipsis):}

This example uses the \isty{amsmath} \dq{dots with binary operators\slash relations}
\glsni{dotsb} instead of the standard \glsni{cdots}:
\begin{code}
\begin{alltt}
\glsni{begindispmath}
y = a\glsni{underscorechar}1 + a\glsni{underscorechar}2 + \glsni{dotsb} + a\glsni{underscorechar}n 
\glsni{enddispmath}
\end{alltt}
\end{code}%
\begin{resultS}[cdots.html]
\[
y = a_1 + a_2 + \dotsb + a_n 
\]
\end{resultS}

\begin{exercise}{Maths: Fractions and Symbols}{ex:mathssym}
This exercise uses a fraction, a square root, subscripts, 
superscripts and symbols.
Try to reproduce the following output:
\begin{result}[maths1ex.html]
The quadratic equation
\[
\sum_{i=0}^{2} a_i x^i = 0
\]
has solutions given by
\[
x = \frac{-a_1 \pm \sqrt{a_1^2 - 4 a_2 a_0}}{2a_2}
\]
\end{result}%
Again you can \downloadorview{mathssym} the solution.
\end{exercise}

%---------- Delimiters ------------

\setnode{delimiters}
\subsection{Delimiters}
\label{sec:delimiters}

Placing brackets around a tall object in maths mode, such as
fractions, does not look right if you use normal sized brackets.
For example:
\begin{code}
\glsni{begindispmath}\newline
\strut~~(\glsni{frac}\marg{1}\marg{1+x})\newline
\glsni{enddispmath}
\end{code}%
results in:
\begin{resultS}[Image: normal sized brackets around fraction.]
\[
(\frac{1}{1+x})
\]
\end{resultS}\bookpagebreak\noindent
Instead, you can automatically resize the delimiters using the
commands:
\begin{definition}
\gls{left}\meta{delimiter}
\end{definition}%
and
\begin{definition}
\gls{right}\meta{delimiter}
\end{definition}%
Rewriting the above example:
\begin{code}
\begin{alltt}
\glsni{begindispmath}
  \glsni{left}( \glsni{frac}\marg{1}\marg{1+x} \glsni{right})
\glsni{enddispmath}
\end{alltt}
\end{code}%
produces:
\begin{resultS}[Image: brackets now the same size as the fraction.]
\[
  \left( \frac{1}{1+x} \right)
\]
\end{resultS}%

Note that you must always have matching \cmdname{left} and
\cmdname{right} commands, although the delimiters used may be
different. \refstepcounter{object}\label{obj:invisdelim}If you want
one of the delimiters to be invisible, use a
\gls{delimiter.periodchar}\ (full stop) as the delimiter. Available
delimiters are shown in
\tableref{tab:delimiters}. (Note for a vertical bar
delimiter it's best to use \isty{amsmath}'s \glsni{lvert} command instead of
\gls{delim.barchar} and \glsni{lVert} instead of \gls{csbar}.) Sometimes using \glsni{left} and
\glsni{right} doesn't produce the optimal sized delimiters. In which
case you can use additional commands provided by the \isty{amsmath}
package shown in \tableref{tab:big}.

\begin{table}[hbtp]
\caption[Delimiters]{Delimiters (\supdag defined by
\protect\isty{amsmath})}
\label{tab:delimiters}
\centering
\begin{tabular}{llllllll}
\begin{inlinedef}\gls{openparen}\end{inlinedef} &
\makeimg{open round bracket}{$($} &
\begin{inlinedef}\gls{closeparen}\end{inlinedef} &
\makeimg{close round bracket}{$)$} &
\begin{inlinedef}\gls{delimiter.opensq}\end{inlinedef} & 
\makeimg{open square bracket}{$[$} &
\begin{inlinedef}\gls{delimiter.closesq}\end{inlinedef} & 
\makeimg{close square bracket}{$]$} \\
\begin{inlinedef}\gls{leftbrace}\end{inlinedef} & 
\makeimg{open curly brace}{$\{$} &
\begin{inlinedef}\gls{rightbrace}\end{inlinedef} & 
\makeimg{close curly brace}{$\}$} &
\begin{inlinedef}\gls{lvert}\supdag\end{inlinedef} & 
\makeimg{left vertical bar}{$\lvert$} &
\begin{inlinedef}\gls{rvert}\supdag\end{inlinedef} & 
\makeimg{right vertical bar}{$\rvert$} \\
\begin{inlinedef}\gls{lVert}\supdag\end{inlinedef} & 
\makeimg{left double vertical bar}{$\lVert$} &
\begin{inlinedef}\gls{rVert}\supdag\end{inlinedef} & 
\makeimg{right double vertical bar}{$\rVert$} &
\begin{inlinedef}\gls{langle}\end{inlinedef} & 
\makeimg{left angle bracket}{$\langle$} &
\begin{inlinedef}\gls{rangle}\end{inlinedef} & 
\makeimg{right angle bracket}{$\rangle$} \\
\begin{inlinedef}\gls{lfloor}\end{inlinedef} & 
\makeimg{left floor}{$\lfloor$} &
\begin{inlinedef}\gls{rfloor}\end{inlinedef} & 
\makeimg{right floor}{$\rfloor$} &
\begin{inlinedef}\gls{lceil}\end{inlinedef} & 
\makeimg{left ceiling}{$\lceil$} &
\begin{inlinedef}\gls{rceil}\end{inlinedef} & 
\makeimg{right ceiling}{$\rceil$}\\
\begin{inlinedef}\gls{uparrow}\end{inlinedef} & 
\makeimg{up arrow}{$\uparrow$} &
\begin{inlinedef}\gls{downarrow}\end{inlinedef} &
\makeimg{down arrow}{$\downarrow$} &
\begin{inlinedef}\gls{Uparrow}\end{inlinedef} & 
\makeimg{double lined up arrow}{$\Uparrow$} &
\begin{inlinedef}\gls{Downarrow}\end{inlinedef} & 
\makeimg{double lined down arrow}{$\Downarrow$}\\
\begin{inlinedef}\gls{updownarrow}\end{inlinedef} & 
\makeimg{up down arrow}{$\updownarrow$} &
\begin{inlinedef}\gls{Updownarrow}\end{inlinedef} & 
\makeimg{double lined up down arrow}{$\Updownarrow$} &
\begin{inlinedef}\gls{delim.slash}\end{inlinedef} & 
\makeimg{forward slash}{$/$} &
\begin{inlinedef}\gls{backslash}\end{inlinedef} & 
\makeimg{backslash}{$\backslash$} 
\end{tabular}
\end{table}

\begin{table}
\caption[Delimiter Sizing]{Additional Commands Provided by
\protect\isty{amsmath} for Delimiter Sizing}
\label{tab:big}
\centering
\renewcommand*{\arraystretch}{2.2}
\begin{tabular}{@{}ll@{\hspace{4\tabcolsep}}ll@{}}
\toprule
\multicolumn{2}{c@{\hspace{4\tabcolsep}}}{\bfseries Definitions} &
\multicolumn{2}{c}{\bfseries Example}\\
\midrule
\multicolumn{2}{c@{\hspace{4\tabcolsep}}}{\em Default Size} & \verb|$( X )$| & $(X)$\\
\begin{inlinedef}\gls{bigl}\meta{delim}\end{inlinedef} &
\begin{inlinedef}\gls{bigr}\meta{delim}\end{inlinedef} &
\verb|$\bigl( X \bigr)$| & 
\makeimg{(X)}{$\bigl(X\bigr)$} \\
\begin{inlinedef}\gls{Bigl}\meta{delim}\end{inlinedef} &
\begin{inlinedef}\gls{Bigr}\meta{delim}\end{inlinedef} &
\verb|$\Bigl( X \Bigr)$| & 
\makeimg{(X)}{$\Bigl(X\Bigr)$} \\
\begin{inlinedef}\gls{biggl}\meta{delim}\end{inlinedef} &
\begin{inlinedef}\gls{biggr}\meta{delim}\end{inlinedef} &
\verb|$\bigl( X \biggr)$| & 
\makeimg{(X)}{$\biggl(X\biggr)$} \\
\begin{inlinedef}\gls{Biggl}\meta{delim}\end{inlinedef} &
\begin{inlinedef}\gls{Biggr}\meta{delim}\end{inlinedef} &
\verb|$\Biggl( X \Biggr)$| & 
\makeimg{(X)}{$\Biggl(X\Biggr)$} \\
\bottomrule
\end{tabular}
\renewcommand*{\arraystretch}{1}% Keep LaTeX2HTML happy
\end{table}

\xminisec{Example (Vertical Bar Delimiters):}

\begin{code}
\glsni{begindispmath}\newline
\glsni{left}\glsni{lvert}\newline
\glsni{frac}\marg{1}\marg{1+x}\newline
\glsni{right}\glsni{rvert}\newline
\glsni{enddispmath}
\end{code}%
\begin{resultS}[Image: vertical bars same height as fraction.]
\[
\left\lvert\frac{1}{1+x}\right\rvert
\]
\end{resultS}

\xminisec{Example (Delimiter with Subscript):}

Delimiters can take limits:
\begin{code}
\glsni{begindispmath}\newline
\glsni{left}\glsni{lvert}\newline
\glsni{frac}\marg{1}\marg{1+x}\newline
\glsni{right}\glsni{rvert}\glsni{underscorechar}\marg{x=0}\newline
\glsni{enddispmath}
\end{code}\reportpagebreak
\begin{resultS}[Image: as before but the right vertical bar has a subscript.]
\[
\left\lvert\frac{1}{1+x}\right\rvert_{x=0}
\]
\end{resultS}

\xminisec{Example (Mismatch):}
The left and right delimiters don't have to match:
\begin{code}
\begin{alltt}
\glsni{begindispmath}
\glsni{left}\glsni{delimiter.opensq}\glsni{frac}\marg{1}\marg{1+x}\glsni{right}\glsni{rangle}
\glsni{enddispmath}
\end{alltt}
\end{code}%
\begin{resultS}[Image: large open square bracket followed by
fraction followed by large right angle bracket.]
\[
\left[\frac{1}{1+x}\right\rangle
\]
\end{resultS}

\xminisec{Example (An invisible delimiter):}

Every \gls{right} must have a matching \gls{left} (and vice versa), so 
use a \gls{delimiter.periodchar}\ (full stop) for an invisible
delimiter.
\begin{code}
\begin{alltt}
\glsni{begindispmath}
 \glsni{left}\glsni{delimiter.periodchar}
  \glsni{frac}\marg{\glsni{partial} f}\marg{\glsni{partial} x}
 \glsni{right}\glsni{rvert}\glsni{underscorechar}\marg{x=0}
\glsni{enddispmath}
\end{alltt}
\end{code}
\begin{resultS}[Image: partial fraction followed by large vertical
line with a subscript.]
\[
\left.
\frac{\partial f}{\partial x}
\right\rvert_{x=0}
\]
\end{resultS}

\refstepcounter{object}\label{introeq}
We have now covered enough to reproduce the equation shown in 
\chapterref{ch:intro}:
\begin{code}
\begin{alltt}
\glsni{newcommand}*\marg{\cmdname{pderiv}}\oarg{2}\marg{\glsni{frac}\marg{\glsni{partial} \glsni{hashchar}1}\marg{\glsni{partial} \glsni{hashchar}2}}
\glsni{newcommand}*\marg{\cmdname{e}}\marg{\glsni{mathrm}\marg{e}}

\glsni{begindispmath}
\cmdname{pderiv}\marg{\glsni{circumchar}2\glsni{mathcal}\marg{L}}\marg{\marg{z\glsni{underscorechar}i\glsni{circumchar}\glsni{rho}}\glsni{circumchar}2} =
-\cmdname{pderiv}\marg{\glsni{rho}\glsni{underscorechar}i}\marg{z\glsni{underscorechar}i\glsni{circumchar}\glsni{rho}}
\glsni{left}(
  \cmdname{pderiv}\marg{v\glsni{underscorechar}i}\marg{\glsni{rho}\glsni{underscorechar}i} \glsni{frac}\marg{\cmdname{e}\glsni{circumchar}\marg{v\glsni{underscorechar}i}}\marg{1-\cmdname{e}\glsni{circumchar}\marg{v\glsni{underscorechar}i}}
  + v\glsni{underscorechar}i \glsni{frac}\marg{\cmdname{e}\glsni{circumchar}\marg{v\glsni{underscorechar}i}\cmdname{pderiv}\marg{v\glsni{underscorechar}i}\marg{\glsni{rho}\glsni{underscorechar}i}(1-\cmdname{e}\glsni{circumchar}\marg{v\glsni{underscorechar}i})
             +\cmdname{e}\glsni{circumchar}\marg{2v\glsni{underscorechar}i}\cmdname{pderiv}\marg{v\glsni{underscorechar}i}\marg{\glsni{rho}\glsni{underscorechar}i}}\marg{(1-\cmdname{e}\glsni{circumchar}\marg{v\glsni{underscorechar}i})\glsni{circumchar}2}
\glsni{right})
\glsni{enddispmath}
\end{alltt}
\end{code}%
\begin{result}[Image showing a second order partial derivative.]
\[
\pderiv{^2\mathcal{L}}{{z_i^\rho}^2} =
-\pderiv{\rho_i}{z_i^\rho}
\left(
  \pderiv{v_i}{\rho_i} \frac{\e^{v_i}}{1-\e^{v_i}}
  + v_i \frac{\e^{v_i}\pderiv{v_i}{\rho_i}(1-\e^{v_i})+\e^{2v_i}\pderiv{v_i}{\rho_i}}{(1-\e^{v_i})^2}
\right)
\]
\end{result}

\xminisec{Note:}
The above code looks a bit complicated, and there are
so many braces that it can be easy to lose track, so here are some
ways of making it a little easier to type:
\begin{enumerate}
\item Whenever you start a new environment type in the
\gls{begin} and \gls{end} bits first, and then insert
whatever goes inside the environment.  This ensures that you always
have a matching \gls{begin} and \gls{end}. The same goes for 
\glsni{begindispmath} and \glsni{enddispmath}.

\item Whenever you type any braces, always type the opening and
closing braces first, and then insert whatever goes in between.  This
will ensure that your braces always match up.
\end{enumerate}

So keeping these notes in mind, let's try typing in the code in a
methodical manner:\screenpagebreak
\begin{enumerate}
\item Start and end the displayed maths mode:
\begin{code}
\glsni{begindispmath}\newline
\glsni{enddispmath}
\end{code}

\item We now need a partial derivative.
(The command \cmdname{pderiv} is defined as described
\htmlref{earlier}{itm:pderiv}\latex{\ifscreenorother{}{\ on 
page~\pageref{itm:pderiv}}}. Make sure you remember to define it,
preferably in the \gls{preamble}.)
\begin{code}
{\color{midgray}
\glsnl{begindispmath}}\newline
\cmdname{pderiv}\marg{}\marg{}\newline
{\color{midgray}
\glsnl{enddispmath}}
\end{code}

\item Let's do the first argument. This partial derivative is
actually a double derivative, which means we need a squared bit on
the top along with a calligraphic L:
\begin{code}
\begin{alltt}
{\color{midgray}\glsnl{begindispmath}
\cmdname{pderiv}\marg*{\textcolor{black}{\glsni{circumchar}2 \glsni{mathcal}\marg{L}}}\marg*{}
\glsnl{enddispmath}}
\end{alltt}
\end{code}

\item The second argument is the $z_i^\rho$ squared bit.  This is a 
nested superscript
\verb|{z_i^\rho}^2|:
\begin{code}
\begin{alltt}
\color{midgray}\glsnl{begindispmath}
\cmdname{pderiv}\marg*{\textasciicircum2 \cmdname{mathcal}\marg{L}}\marg*{\textcolor{black}{\marg{z\glsni{underscorechar}i\glsni{circumchar}\glsni{rho}}\glsni{circumchar}2}}
\glsnl{enddispmath}
\end{alltt}
\end{code}

\item We can do the next partial derivative in the same way.  This 
one is slightly easier to do:
\begin{code}
{\color{midgray}\glsnl{begindispmath}\newline
\cmdname{pderiv}\marg*{\glsnl{circumchar}2 \glsnl{mathcal}{L}}\marg*{\marg*{z\glsnl{underscorechar}i\glsnl{circumchar}\glsnl{rho}}\glsnl{circumchar}2}} =\newline
-\cmdname{pderiv}\marg{\glsni{rho}\glsni{underscorechar}i}\marg{z\glsni{underscorechar}i\glsni{circumchar}\glsni{rho}}\newline
{\color{midgray}\glsnl{enddispmath}}
\end{code}

\item Delimiters also need to occur in pairs, like curly braces and 
\glsni{begin} and \glsni{end}, so let's do them next:
\begin{code}
{\color{midgray}\verb|\[|\\
\verb|\pderiv{^2 \mathcal{L}}{{z_i^\rho}^2} =|\\
\verb|-\pderiv{\rho_i}{z_i^\rho}|}\\
\strut~~\glsni{left}(\\
\strut~~\glsni{right})\\
{\color{midgray}\verb|\]|}
\end{code}\screenpagebreak

\item Now we need to do the bits inside the brackets.  First of all 
we have yet another partial derivative:
\begin{code}
{\color{midgray}\verb|\[|\\
\verb|\pderiv{^2 \mathcal{L}}{{z_i^\rho}^2} =|\\
\verb|-\pderiv{\rho_i}{z_i^\rho}|\\
\strut~~\glsnl{left}\glsnl{openparen}}\\
\strut~~~~\cmdname{pderiv}\marg{v\glsni{underscorechar}i}\marg{\glsni{rho}\glsni{underscorechar}i}\\
{\color{midgray}~~\glsnl{right}\glsnl{closeparen}\\
\verb|\]|}
\end{code}

\item Now we have a fraction following the partial derivative from
the previous step.
(Make sure you use braces for the exponential bit: \verb|\e^{v_i}|
($\e^{v_i}$) is not the same as \verb|\e^v_i| ($\e^v_i$). The
command \cmdname{e} is defined as described
\htmlref{earlier}{sec:scripts}\latex{ in \sectionref{sec:scripts}}.
Make sure you define it, preferably in the \gls{preamble}.)
\begin{code}
{\color{midgray}\verb|\[|\\
\verb|\pderiv{^2 \mathcal{L}}{{z_i^\rho}^2} =|\\
\verb|-\pderiv{\rho_i}{z_i^\rho}|\\
\strut~~\glsnl{left}\glsnl{openparen}\\
\strut~~~~\cmdname{pderiv}\marg*{v\glsnl{underscorechar}i}\marg*{\glsnl{rho}\glsnl{underscorechar}i}} \glsni{frac}\marg{\cmdname{e}\glsni{circumchar}\marg{v\glsni{underscorechar}i}}\marg{1-\cmdname{e}\glsni{circumchar}\marg{v\glsni{underscorechar}i}}\\
{\color{midgray}~~\glsnl{right}\glsnl{closeparen}\\
\verb|\]|}
\end{code}

\item This is followed by $v_i$ times another fraction:
\begin{code}\par\noindent
{\color{midgray}\verb|\[|\\
\verb|\pderiv{^2 \mathcal{L}}{{z_i^\rho}^2} =|\\
\verb|-\pderiv{\rho_i}{z_i^\rho}|\\
\strut~~\glsnl{left}\glsnl{openparen}\\
\strut~~~~\cmdname{pderiv}\marg*{v\glsnl{underscorechar}i}\marg*{\glsnl{rho}\glsnl{underscorechar}i} \glsnl{frac}\marg*{\cmdname{e}\glsnl{circumchar}\marg*{v\glsnl{underscorechar}i}}\marg*{1-\cmdname{e}\glsnl{circumchar}\marg*{v\glsnl{underscorechar}i}}}\\
\strut~~~~+ v\glsni{underscorechar}i \glsni{frac}\marg{}\marg{}\\
{\color{midgray}~~\glsnl{right}\glsnl{closeparen}\\
\color{midgray}\verb|\]|}
\end{code}

\item The bottom part of the
fraction (the denominator) is easier than the top, so let's do that first:
\begin{code}
\color{midgray}\glsnl{begindispmath}\newline
\cmdname{pderiv}\marg*{\textasciicircum2 \cmdname{mathcal}\marg*{L}}\marg*{\marg*{z\_i\textasciicircum\cmdname{rho}}\textasciicircum2} =\newline
-\cmdname{pderiv}\marg*{\cmdname{rho}\_i}\marg*{z\_i\textasciicircum\cmdname{rho}}\newline
\strut~~\cmdname{left}(\newline
\strut~~~~\cmdname{pderiv}\marg*{v\_i}\marg*{\cmdname{rho}\_i} \cmdname{frac}\marg*{\cmdname{e}\textasciicircum\marg*{v\_i}}\marg*{1-\cmdname{e}\textasciicircum\marg*{v\_i}}\newline
\strut~~~~+ v\_i \cmdname{frac}\marg*{}\marg*{\textcolor{black}{(1-\cmdname{e}\glsni{circumchar}\marg{v\glsni{underscorechar}i})\glsni{circumchar}2}}\newline
\strut~~\cmdname{right})\newline
\glsnl{enddispmath}
\color{black}
\end{code}

% Something weird is causing a shift in the 
% frame here.
% Compensating by ending and restarting enumerate
% environment
\breakenumi

\item Now for the top part of the fraction (the numerator).
To refresh your memory, it should look like:
\[
\e^{v_i}\pderiv{v_i}{\rho_i}(1-\e^{v_i})
+ \e^{2v_i}\pderiv{v_i}{\rho_i}
\]
That's a bit
complicated, so let's break it down:

\begin{enumerate}

\item The first term is:
\begin{alltt}
\cmdname{e}\glsni{circumchar}\marg{v\glsni{underscorechar}i}
\end{alltt}

\item The next term is another partial derivative:
\begin{alltt}
\cmdname{pderiv}\marg{v\glsni{underscorechar}i}\marg{\glsni{rho}\glsni{underscorechar}i}
\end{alltt}

\item Then we have:
\begin{alltt}
(1-\cmdname{e}\glsni{circumchar}\marg{v\glsni{underscorechar}i})
\end{alltt}

\item Next we have to add on:
\begin{alltt}
+\cmdname{e}\glsni{circumchar}\marg{2v\glsni{underscorechar}i}
\end{alltt}

\item And finally we have:
\begin{alltt}
\cmdname{pderiv}\marg{v\glsni{underscorechar}i}\marg{\glsni{rho}\glsni{underscorechar}i}
\end{alltt}
\end{enumerate}

So the numerator is:
\begin{code}
\begin{verbatim}
\e^{v_i}\pderiv{v_i}{\rho_i}(1-\e^{v_i})
  + \e^{2v_i}\pderiv{v_i}{\rho_i}
\end{verbatim}
\end{code}

Inserting this into our code:\reportpagebreak
\begin{code}
\begin{alltt}
\color{midgray}\glsnl{begindispmath}
\cmdname{pderiv}\marg*{\textasciicircum2\cmdname{mathcal}\marg*{L}}\marg*{\marg*{z\_i\textasciicircum\cmdname{rho}}\textasciicircum2} =
-\cmdname{pderiv}\marg*{\cmdname{rho}\_i}\marg*{z\_i\textasciicircum\cmdname{rho}}
\cmdname{left}(
  \cmdname{pderiv}\marg*{v\_i}\marg*{\cmdname{rho}\_i} \cmdname{frac}\marg*{\cmdname{e}\textasciicircum\marg*{v\_i}}\marg*{1-\cmdname{e}\textasciicircum\marg*{v\_i}}
  \color{midgray}+ v\_i \cmdname{frac}\marg*{\textcolor{black}{\cmdname{e}\textasciicircum\marg{v\_i}\cmdname{pderiv}\marg{v\_i}\marg{\cmdname{rho}\_i}(1-\cmdname{e}\textasciicircum\marg{v\_i})
             +\cmdname{e}\textasciicircum\marg{2v\_i}\cmdname{pderiv}\marg{v\_i}\marg{\cmdname{rho}\_i}}}\marg*{(1-\cmdname{e}\textasciicircum\marg*{v\_i})\textasciicircum2}
\cmdname{right})
\glsnl{enddispmath}
\end{alltt}
\color{black}
\end{code}

\end{enumerate}


%----------- Arrays -------------

\setnode{arrays}
\subsection{Arrays}
\label{sec:arrays}

Mathematical structures such as matrices and vectors require elements
to be arranged in rows and columns.  Just as we can align material in
rows and columns in text mode using the
\gls{env-tabular} environment (\sectionref{sec:tabular}), we can do the
same in maths mode using the \gls{env-array} environment.  The
\glsni{env-array} environment has the same format as the
\gls{env-tabular} environment, however it must be in maths mode.
The column half-gaps are given by the \gls{length} register
\gls{arraycolsep} (analogous to \gls{tabcolsep}). 

\xminisec{Example:}

\begin{code}
\glsni{begindispmath}\newline
\glsni{begin}\marg{array}\marg{rrr}\newline
0 \gls{ampchar} 1 \gls{ampchar} 19\gls{tab.dbbackslashchar}\newline
-6 \gls{ampchar} 10 \gls{ampchar} 200\newline
\glsni{end}\marg{array}\newline
\glsni{enddispmath}
\end{code}

\begin{result}[Image shows 3 columns and 2 rows.]
\[
\begin{array}{rrr}
0 & 1 & 19\\
-6 & 10 & 200
\end{array}
\]
\end{result}

\xminisec{Example (Adding Delimiters):}

\begin{code}
\begin{alltt}
\glsni{begindispmath}
\glsni{left}\glsni{openparen}
 \glsni{begin}\marg{array}\marg{rrr}
   0 \glsni{ampchar} 1 \glsni{ampchar} 19\glsni{tab.dbbackslashchar}
  -6 \glsni{ampchar} 10 \glsni{ampchar} 200
 \glsni{end}\marg{array}
\glsni{right}\glsni{closeparen}
\glsni{enddispmath}
\end{alltt}
\end{code}%
\begin{resultS}[Image: as before but has large round brackets on either
side.]
\[
\left(
\begin{array}{rrr}
0 & 1 & 19\\
-6 & 10 & 200
\end{array}
\right)
\]
\end{resultS}

\xminisec{Adding a Vertical Rule:}

A vertical rule can be added using \gls{array.barchar} in the column
specifier. For example:
\begin{code}
\begin{alltt}
\glsni{begindispmath}
\glsni{left}\glsni{openparen}
 \glsni{begin}\marg{array}\marg{rr\glsni{array.barchar}r}
   0 \glsni{ampchar} 1 \glsni{ampchar} 19\glsni{tab.dbbackslashchar}
  -6 \glsni{ampchar} 10 \glsni{ampchar} 200
 \glsni{end}\marg{array}
\glsni{right}\glsni{closeparen}
\glsni{enddispmath}
\end{alltt}
\end{code}%
\begin{resultS}[Image: 3 columns and 2 rows with a vertical line
between the second and third column.]
\[
\left(
\begin{array}{rr|r}
0 & 1 & 19\\
-6 & 10 & 200
\end{array}
\right)
\]
\end{resultS}

\xminisec{Example (Cases):}
This example uses an 
\htmlref{invisible delimiter}{obj:invisdelim}:
\begin{code}
\begin{alltt}
\glsni{begindispmath}
f(x) = 
 \glsni{left}\glsni{leftbrace}
  \glsni{begin}\marg{array}\marg{rl}
   -1 \glsni{ampchar} x \glsni{ltchar} 0\glsni{tab.dbbackslashchar}
   0  \glsni{ampchar} x = 0\glsni{tab.dbbackslashchar}
   +1 \glsni{ampchar} x \glsni{gtchar} 0
  \glsni{end}\marg{array}
 \glsni{right}\glsni{delimiter.periodchar}
\glsni{enddispmath}
\end{alltt}
\end{code}%
\begin{resultS}[Image: array has large open brace on the left and
nothing on the right of it.]
\[
f(x) = 
\left\{
\begin{array}{rl}
-1 & x < 0\\
0 & x = 0\\
+1 & x > 0
\end{array}
\right.
\]
\end{resultS}
This can be rewritten more compactly using the \isty{amsmath}
\gls{env-cases} environment:
\begin{code}
\gls{begindispmath}\newline
f(x) = \newline
\strut~\glsni{begin}\marg{cases}\newline
\strut~~-1 \glsni{ampchar} x < 0\glsni{tab.dbbackslashchar}\newline
\strut~~0  \glsni{ampchar} x = 0\glsni{tab.dbbackslashchar}\newline
\strut~~+1 \glsni{ampchar} x > 0\newline
\strut~\glsni{end}\marg{cases}\newline
\gls{enddispmath}
\end{code}
\begin{resultS}[Image: as previous example but first column is left
aligned.]
\[
f(x) = 
\begin{cases}
-1 & x < 0\\
0 & x = 0\\
+1 & x > 0
\end{cases}
\]
\end{resultS}

The \isty{amsmath} package provides some convenient environments to
typeset matrices: \gls{env-pmatrix}, \gls{env-bmatrix},
\gls{env-Bmatrix}, \gls{env-vmatrix} and \gls{env-Vmatrix}. These
are similar to the \gls{env-array} environment except there is no
argument, and they add (respectively) $\bigl(~\bigr)$,
$\bigl[~\bigr]$, $\bigl\{~\bigr\}$, $\bigl\lvert~\bigr\rvert$ and
$\bigl\lVert~\bigr\rVert$ delimiters. There is also the
\gls{env-matrix} environment that doesn't have any delimiters.

\xminisec{Example:}
\begin{code}
\begin{alltt}
\glsni{begin}\marg{equation}
 \glsni{begin}\marg{pmatrix}
  a \glsni{ampchar} b\glsni{tab.dbbackslashchar}
  c \glsni{ampchar} d
 \glsni{end}\marg{pmatrix}
\glsni{end}\marg{equation}
\end{alltt}
\end{code}
\begin{result}[Image: 2 by 2 matrix with round delimiters and an
equation number.]
\begin{equation}
\begin{pmatrix}
a & b\\
c & d
\end{pmatrix}
\end{equation}
\end{result}

The \isty{amsmath} package also provides the environment
\gls{env-smallmatrix} designed for in-line use. You need to add any
delimiters explicitly.

\xminisec{Example:}
\begin{code}
\begin{alltt}
Here is a small matrix
\glsni{begin}\marg{math}
 \glsni{left}\glsni{openparen}
  \glsni{begin}\marg{smallmatrix}
  a \glsni{ampchar} b\glsni{tab.dbbackslashchar}
  c \glsni{ampchar} d
  \glsni{end}\marg{smallmatrix}
 \glsni{right}\glsni{closeparen}
\glsni{end}\marg{math}
in a line of text.
\end{alltt}
\end{code}
\begin{resultS}[Image: matrix is small enough to fit into the line
of text.]
Here is a small matrix
\begin{math}
\left(
\begin{smallmatrix}
a & b\\
c & d
\end{smallmatrix}
\right)
\end{math}
in a line of text.
\end{resultS}


%----------- Vectors ------------

\setnode{vectors}
\subsection{Vectors}
\label{sec:vec}

A variable representing a vector can be typeset using the command:
\begin{definition}
\gls{vec}\marg{\meta{variable}}
\end{definition}\screenpagebreak

\xminisec{Example:}
\begin{codeS}
\glsni{begindispmath}
\glsni{vec}\marg{x}
\glsni{enddispmath}
\end{codeS}
\begin{resultS}[Image: x with a small right arrow above it.]
\[ \vec{x} \]
\end{resultS}

Vectors are often typeset in bold.  This can be 
done by \htmlref{redefining}{sec:renewcom} the \glsni{vec}
command\doifbook{ (see \sectionref{sec:renewcom})}.
You could use \gls{mathbf}, for example:
\begin{code}
\begin{alltt}
\gls{renewcommand}\marg{\glsni{vec}}\oarg{1}\marg{\glsni{mathbf}\marg{\glsni{hashchar}1}}

\glsni{begindispmath}
  \glsni{vec}\marg{x}\glsni{cdot}\glsni{vec}\marg{\glsni{xi}} = z
\glsni{enddispmath}
\end{alltt}
\end{code}\bookpagebreak\noindent
\begin{resultS}[Image: upright bold x centred dot xi equals z]
\[
\mathbf{x}\cdot\mathbf{\xi} = z
\]
\end{resultS}
\refstepcounter{object}\label{cmd:vec}However, as you may have
noticed, the 
Greek letter $\xi$ has not come out in bold.  Here's an alternative 
(using \gls{boldsymbol} defined in the \isty{amsfonts} \htmlref{package}{sec:packages}):
\begin{code}
\gls{renewcommand}\marg{\glsni{vec}}\oarg{1}\marg{\glsni{boldsymbol}\marg{\glsni{hashchar}1}}\newline
\glsni{begindispmath}\newline
\strut~~\glsni{vec}\marg{x}\glsni{cdot}\glsni{vec}\marg{\glsni{xi}} = z\newline
\glsni{enddispmath}
\end{code}%
\begin{resultS}[As before but x and xi in italic bold]
\[
\boldsymbol{x}\cdot\boldsymbol{\xi} = z
\]
\end{resultS}

Located (or position) vectors, on the other hand, are usually typeset with a right
arrow, but the default definition of \glsni{vec} produces an
arrow that is too small:
\begin{codeS}
\glsni{begindispmath}
\glsni{vec}\marg{OP}
\glsni{enddispmath}
\end{codeS}
\begin{resultS}[Image: the letters O P with a small right arrow
above them.]
\[
\vec{OP}
\]
\end{resultS}
Instead, use \gls{overrightarrow} (\tableref{tab:overunderarrows}):
\begin{codeS}
\glsni{begindispmath}
\glsni{overrightarrow}\marg{OP}
\glsni{enddispmath}
\end{codeS}
\begin{resultS}[Image: as above but with a longer arrow that spans
both letters.]
\[
\overrightarrow{OP}
\]
\end{resultS}
You might prefer to define separate commands for a located vector and a vector
variable.

\xminisec{Example:}
In the \glsterm{preamble}, define \cmdname{lvec} for a
located vector and \cmdname{bvec} for a vector variable:
\begin{code}
\begin{alltt}
\gls{newcommand}*\marg{\cmdname{lvec}}\oarg{1}\marg{\glsni{overrightarrow}\marg{\glsni{hashchar}1}}
\glsni{newcommand}*\marg{\cmdname{bvec}}\oarg{1}\marg{\glsni{boldsymbol}\marg{\glsni{hashchar}1}}
\end{alltt}
\end{code}
Later in the document:
\begin{code}
Let \glsni{dollarchar}\cmdname{bvec}\marg{u}=(x, y)\glsni{dollarchar}
represent \glsni{dollarchar}\cmdname{lvec}\marg{OP}\glsni{dollarchar}, then\newline
\glsni{begindispmath}\newline
\strut~~\glsni{lVert} \cmdname{bvec}\marg{u} \glsni{rVert}
    = \glsni{sqrt}\marg{x\glsni{circumchar}2 + y\glsni{circumchar}2}\newline
\glsni{enddispmath}
\end{code}
\begin{result}[Image: the vector u appeas in italic bold and the
letters O P are as above.]
\newcommand*{\lvec}[1]{\overrightarrow{#1}}%
\newcommand*{\bvec}[1]{\boldsymbol{#1}}%
Let $\bvec{u}=(x, y)$ represent $\lvec{OP}$, then
\[
  \lVert \bvec{u} \rVert = \sqrt{x^2 + y^2}
\]
\end{result}
\bookpagebreak

\begin{exercise}{Maths: Vectors and Arrays}{ex:vectors}
Try to produce the following:
\renewcommand{\vec}[1]{\boldsymbol{#1}}
\begin{result}[maths2ex.html]
\[
\boldsymbol{A}\vec{x} =
\begin{pmatrix}
0 & 1\\
2 & 3
\end{pmatrix}
\begin{pmatrix}
1\\
2
\end{pmatrix}
=
\begin{pmatrix}
2\\
8
\end{pmatrix}
= \vec{y}
\]
\end{result}%
As before, you can \downloadorview{vectors} the solution.
\end{exercise}

%--------------- Mathematical Spacing -------------------

\setnode{mathssp}
\subsection{Mathematical Spacing}
\label{sec:mathsspacing}

\LaTeX\ deals with mathematical spacing fairly well, but sometimes
you may find you want to adjust the spacing yourself.  Available
spacing commands are listed in \tableref{tab:spacing}.

\begin{table}[hbtp]
\caption[Mathematical Spacing Commands]{Mathematical Spacing Commands
(\supdag provided by \protect\istyend{amsmath})}
\label{tab:spacing}
\centering
\begin{tabular}{lll}
\toprule
\bfseries Command & \bfseries Example Input & \bfseries Example Output\\
\midrule
 & \verb|$AB$| & $AB$ \\
\gls{thinspace} or \gls{comma} & \verb|$A\,B$| & $A\,B$\\
\gls{medspace}\supdag\ or \gls{colon} & \verb|$A\:B$| & $A\:B$\\
\gls{thickspace}\supdag\ or \gls{semicolon} & \verb|$A\;B$| & $A\;B$\\
\gls{quad} & \verb|$A\quad B$| & $A\quad B$\\
\gls{qquad} & \verb|$A\qquad B$| & $A\qquad B$\\
\gls{negthinspace} or \gls{exclam} & \verb|$A\!B$| & $A\!B$ \\
\gls{negmedspace}\supdag & \verb|$A\negmedspace B$| & $A\negmedspace B$ \\
\gls{negthickspace}\supdag & \verb|$A\negthickspace B$| & $A\negthickspace B$ \\
\bottomrule
\end{tabular}
\end{table}

\begin{exercise}{More Mathematics}{ex:moremaths}
This exercise uses the spacing command \glsni{qquad}. In addition,
it has a
\htmlref{function name}{sec:declaremathop}, diag, and it uses the
\gls{forall} and \htmlref{ellipses}{sec:ellipses} symbols. It also
\htmlref{redefines the \cmdname{vec}
command}{cmd:vec}\indexCom{vec}, as was done in the previous
section, uses the
\gls{env-bmatrix} environment (see \sectionref{sec:arrays}), and has \htmlref{subscripts and superscripts}{sec:scripts}.

Try to reproduce the following output:
\begin{result}[maths3ex.html]
\renewcommand{\vec}[1]{\boldsymbol{#1}}
The set of linear equations:
\[
a_ix_i = b_i  \qquad \forall i=1, \ldots, n
\]
can be written as a matrix equation:
\[
\diag(\vec{a})\vec{x} = \vec{b}
\]
where $\vec{x} = (x_1, \ldots, x_n)^T$, $\vec{b} = (b_1, \ldots, b_n)^T$
and
\[
\diag(\vec{a}) = 
\begin{bmatrix}
a_1 & 0 & \cdots & 0\\
0 & a_2 & \ddots & \vdots\\
\vdots & \ddots & \ddots & 0\\
0 & \cdots & 0 & a_n
\end{bmatrix}
\]
\end{result}

Again, you can \downloadorview{moremaths} the solution.
\end{exercise}

%%%%%%%%%%%%%%%%%% DEFINING ENVIRONMENTS %%%%%%%%%%%%%%%%

\setnode{newenv}
\chapter{Defining Environments}
\label{ch:newenv}

Just as you can \htmlref{define new commands}{ch:newcom}, you can 
also define new \glspl{environment}.  The command
\begin{definition}
\gls{newenvironment}\marg{\meta{env-name}}\oarg{\meta{n-args}}\oarg{\meta{default}}\marg{\meta{begin-code}}\marg{\meta{end-code}}
\end{definition}%
is used to define a new environment. As with new commands, you can
use the optional argument \meta{n-args} to define an environment with
arguments, and \meta{default} to define an environment with an
optional argument.

The first argument \meta{env-name} is the name of your new
environment.  Remember that the environment name must not have a
backslash.  The mandatory arguments \meta{begin-code} and
\meta{end-code} indicate what \LaTeX\ \reportlinebreak should do at the beginning and
end of the environment. Note that although \meta{begin-code} can
reference the arguments using \gls{hashchar}\texttt{1} etc, the
\meta{end-code} part can't.

\xminisec{Example (An Exercise Environment):}

Let's first consider an example of an
environment without any arguments.  Let's make an environment called,
say, \envname{exercise} that prints \textbf{Exercise} in bold and
typesets the contents of the environment in italic, with a gap
between the title and the contents.  In other words,
we want the following code:
\begin{code}
\glsni{begin}\marg{exercise}\newline
This is a sample.\newline
\glsni{end}\marg{exercise}
\end{code}
to produce the following output:
\begin{result}[newenv1.html]
\par\noindent
\textbf{Exercise}
\begin{itshape}
\par
\vspace{\baselineskip}%
\noindent
This is a sample.
\end{itshape}
\end{result}
(In the \htmlref{next chapter}{ch:counters} we will add numbering.)

Let's first consider what we want this environment to do: we can get
the word \dq{Exercise} in bold using \gls{textbf},
and the italic font can be
obtained by using the \gls{env-itshape} environment (recall
\sectionref{sec:fonts}).  
So, at the start of our
new environment we need
\begin{codeS}
\glsni{textbf}\marg{Exercise}\glsni{begin}\marg{itshape}
\end{codeS}
and at the end of our new
environment we need to end the \glsni{env-itshape} environment:
\begin{codeS}
\glsni{end}\marg{itshape}
\end{codeS}
Putting the above together into the new environment definition:
\begin{code}\obeyspaces
\glsni{newenvironment}\marg{exercise}\glsni{percentchar} environment name\newline
\marg{\glsni{percentchar} begin code\newline
\strut~~\glsni{textbf}\marg{Exercise}\glsni{begin}\marg{itshape}\glsni{percentchar}\newline
}\glsni{percentchar}\newline
\marg{\glsni{end}\marg{itshape}}\glsni{percentchar} end code
\end{code}
Let's try it out:
\begin{code}
\begin{verbatim}
\begin{exercise}
This is a sample.
\end{exercise}
\end{verbatim}
\end{code}%
\begin{resultS}[newenv2.html]
\textbf{Exercise}
\begin{itshape}
This is a sample.
\end{itshape}
\end{resultS}
Not quite right.  Let's put a paragraph break after
\textbf{Exercise}, and put one before it as well. The command
\gls{par} can be used to make a paragraph break and the extra bit of
vertical spacing can be produced using \glsni{vspace}. The
\gls{length} \gls{baselineskip} is the interline spacing.
Modifications are shown in bold \modification{like this}.
\begin{code}\obeyspaces
\glsnl{newenvironment}\marg*{exercise}\glsnl{percentchar} environment name\newline
\marg*{\glsnl{percentchar} begin code\newline
\strut~~\modification{\glsni{par}\glsni{vspace}\marg{\glsni{baselineskip}}\glsni{percentchar}}\newline
\strut~~\glsnl{textbf}\marg*{Exercise}\glsnl{begin}\marg*{itshape}\glsnl{percentchar}\newline
\strut~~\modification{\glsni{par}\glsni{vspace}\marg{\glsni{baselineskip}}\glsni{percentchar}}\newline
}\glsnl{percentchar}\newline
\marg*{\glsnl{end}\marg*{itshape}}\glsnl{percentchar} end code
\end{code}%
Let's have a look at the output now:
\begin{result}
\dohtmlcolorbox{%
\par\vspace{\baselineskip}%
\textbf{Exercise}\begin{itshape}%
\par\vspace{\baselineskip}
This is a sample.
\end{itshape}
}
\end{result}
\faq{There's a space added after my environment}{spaftend}%
The indent at the start of each line is caused by the normal
paragraph indentation. This can be suppressed using \gls{noindent}.
It's also a good idea to suppress
any spaces immediately following \verb|\begin{exercise}| and \verb|\end{exercise}|, which can be done
using \gls{ignorespaces} and \reportlinebreak\screenlinebreak\gls{ignorespacesafterend}.
Modifications are again shown in bold \modification{like this}.
\begin{code}\obeyspaces
\glsnl{newenvironment}\marg*{exercise}\glsnl{percentchar} environment name\newline
\marg*{\glsnl{percentchar} begin code\newline
\strut~~\glsnl{par}\glsnl{vspace}\marg*{\glsnl{baselineskip}}\modification{\glsni{noindent}}\newline
\strut~~\glsnl{textbf}\marg*{Exercise}\glsnl{begin}\marg*{itshape}\glsnl{percentchar}\newline
\strut~~\glsnl{par}\glsnl{vspace}\marg*{\glsnl{baselineskip}}\modification{\glsni{noindent}\glsni{ignorespaces}}\newline
}\glsnl{percentchar}\newline
\marg*{\glsnl{percentchar} end code\newline
\strut~~\glsnl{end}\marg{itshape}\modification{\glsni{ignorespacesafterend}}\newline 
}
\end{code}%
The \envname{exercise} environment now appears as:
\begin{result}
\par\vspace{\baselineskip}\noindent
\textbf{Exercise}\begin{itshape}%
\par\vspace{\baselineskip}\noindent
This is a sample.
\end{itshape}
\end{result}

Now let's modify our code so that the environment takes an
argument.  The argument should indicate the exercise topic.
For example, the following code:
\begin{code}
\begin{verbatim}
\begin{exercise}{An Example}
This is a sample.
\end{exercise}
\end{verbatim}
\end{code}%
should produce the following result:
\begin{result}[newenv3.html]
\par\vspace{\baselineskip}\noindent\textbf{Exercise (An Example)}
\begin{itshape}\par\vspace{\baselineskip}\noindent
This is a sample.
\end{itshape}
\end{result}%
As with \gls{newcommand}, \gls{hashchar}\texttt{1} is
used to indicate the first argument.  We can now modify the code as
follows:
\begin{code}\obeyspaces
\glsnl{newenvironment}\marg*{exercise}\modification{\oarg{1}}\glsnl{percentchar} environment name\newline
\marg*{\glsnl{percentchar} begin code\newline
\strut~~\glsnl{par}\glsnl{vspace}\marg*{\glsnl{baselineskip}}\glsnl{noindent}\newline
\strut~~\glsnl{textbf}\marg*{Exercise\modification{ (\glsni{hashchar}1)}}\glsnl{begin}\marg*{itshape}\glsnl{percentchar}\newline
\strut~~\glsnl{par}\glsnl{vspace}\marg*{\glsnl{baselineskip}}\glsnl{noindent}\glsnl{ignorespaces}\newline
}\glsnl{percentchar}\newline
\marg*{\glsnl{percentchar} end code\newline
\strut~~\glsnl{end}\marg*{itshape}\glsnl{ignorespacesafterend}\newline
}
\end{code}

\setnode{renewenv}
\section{Redefining Environments}
\label{sec:renewenv}

It is also possible to redefine an environment using:
\begin{definition}[\linewidth]
\gls{renewenvironment}\marg{\meta{env-name}}\oarg{\meta{n-args}}\oarg{\meta{default}}\marg{\meta{begin-code}}\marg{\meta{end-code}}
\end{definition}%
As with \gls{renewcommand}, only redefine an existing environment if
you want a modified version of that environment rather than because
you like the environment name.

\begin{exercise}{Defining a New Environment}{ex:newenv}
If you did any of the exercises from \exerciseref{ex:section}
to \exerciseref{ex:tables}, go back to the document you
created and define the \envname{exercise} environment as in
the example above.  Then try creating some exercises using this
environment. You could, maybe, put an exercise in the first
chapter, and then another one in the second chapter.
Again you can \downloadorview{newenv} an example.
\end{exercise}

%%%%%%%%%%%%%%%%% Counters %%%%%%%%%%%%%%%%%%%%%%%%%

\setnode{counters}
\chapter{Counters}
\label{ch:counters}

As we have seen, \LaTeX\ automatically generates numbers for
chapters, sections, equations etc. These numbers are stored in
\keyword{counters}. The names of these counters are usually the same
as the name of the object with which it is associated but without
any backslash. For example, the \gls{chapter} command has an associated
counter called \icounter{chapter}, the \gls{footnote} command has an
associated counter called \icounter{footnote}, the
\gls{env-equation} environment has an associated counter called
\icounter{equation}, the \gls{env-figure} environment has an associated counter
called \icounter{figure} and the \gls{env-table} environment has an associated
counter called \icounter{table}. There is also a counter called
\icounter{page} that keeps track of the current page number.

The value of a counter can be displayed using the command
\begin{definition}
\icmdname{the}\meta{counter}
\end{definition}%
where \meta{counter} is the name of the associated counter.  Note
that \meta{counter} does not go in curly braces and adjoins
\cmdname{the} (for example, \gls{thepage}\faq{Page number is wrong at
start of page}{wrongpn}, \reportlinebreak\gls{thesection} or
\gls{thechapter}).   In fact, we have already encountered
\reportlinebreak\screenlinebreak\gls{thefigure} in \sectionref{sec:subfloats}.

\xminisec{Example:}
\begin{code}
\begin{alltt}
This page is Page\glsni{tildechar}\glsni{thepage}.
The current chapter is Chapter\glsni{tildechar}\glsni{thechapter}.
\end{alltt}
\end{code}\label{example:thepage}%
\begin{resultS}[This page is Page 226. The current chapter is Chapter 11.]
% LaTeX2HTML doesn't store the page number in the image file
% so use \pageref to fake \thepage instead
This page is Page~\pageref*{example:thepage}.
The current chapter is Chapter~\thechapter.
\end{resultS}%
New counters can be created using the command:
\begin{definition}
\gls{newcounter}\marg{\meta{counter-name}}\oarg{\meta{outer-counter}}
\end{definition}%
The \gls{mandatory} \meta{counter-name}
is the name of your new counter (no backslash in the name).  For
example, let's define a counter called \texttt{exercise} to keep
track of each exercise. (Recall the exercise example from
\chapterref{ch:newenv}.)
\begin{codeS}
\glsni{newcounter}\marg{exercise}
\end{codeS}%
We can now display the value of the counter using
the command \cmdname{theexercise}.  At the moment the counter
has the value zero, the value can be changed using one
of the following commands:

\begin{fwlist}{\cmdname{addtocounter}\marg{\meta{counter}}\marg{\meta{num}}}
\fwitem{\gls{stepcounter}\marg{\meta{counter}}} Increments 
\meta{counter} by 1

\fwitem{\gls{refstepcounter}\marg{\meta{counter}}} As above, but 
allows you to cross-\reportlinebreak reference the counter using \gls{label} and 
\gls{ref}

\fwitem{\gls{setcounter}\marg{\meta{counter}}\marg{\meta{num}}} 
Sets the counter to \meta{num}

\fwitem{\gls{addtocounter}\marg{\meta{counter}}\marg{\meta{num}}} 
Adds \meta{num} to \meta{counter}
\end{fwlist}

A couple of the commands above take a number \meta{num} as one of 
the arguments.  If you want to use another counter for this argument, 
you need to use
\begin{definition}
\gls{value}\marg{\meta{counter}}
\end{definition}%
For example, if you want to set our new \texttt{exercise} counter to 
the same value as the \texttt{page} counter, you would do
\begin{codeS}
\glsni{setcounter}\marg{exercise}\marg{\glsni{value}\marg{page}}
\end{codeS}%
Let's go back to the \envname{exercise} environment you
created in \exerciseref{ex:newenv}.  The exercises really ought to have
an associated number, and this number should be incremented each time 
we use the exercise environment.  So let's modify our code to do this.
Modifications are illustrated in bold \modification{like this}:
\begin{code}\obeyspaces
\modification{\cmdname{newcounter}\marg{exercise}}\newline
\strut\newline
\glsni{newenvironment}\marg{exercise}\oarg{1}\glsni{percentchar} environment name\newline
\marg{\glsni{percentchar} begin code\newline
\strut~~\glsni{par}\glsni{vspace}\marg{\glsni{baselineskip}}\glsni{noindent}\newline
\strut~~\modification{\glsni{refstepcounter}\marg{exercise}}\glsni{percentchar}\newline
\strut~~\glsni{textbf}\marg{Exercise \modification{\cmdname{theexercise}\glsni{spacesym}}(\glsni{hashchar}1)}\glsni{percentchar}\newline
\strut~~\glsni{begin}\marg{itshape}\glsni{percentchar}\newline
\strut~~\glsni{par}\glsni{vspace}\marg{\glsni{baselineskip}}\glsni{percentchar}\newline
\strut~~\glsni{noindent}\glsni{ignorespaces}\newline
}\glsni{percentchar}\newline
\marg{\glsni{percentchar} end code\newline
\strut~~\glsni{end}\marg{itshape}\glsni{percentchar}\newline
\strut~~\modification{\glsni{par}\glsni{vspace}\marg{\glsni{baselineskip}}\glsni{percentchar}}\newline
\strut~~\modification{\glsni{noindent}}\glsni{ignorespacesafterend}\newline
}
\end{code}%
Note that the counter needs to be incremented before it is used.  
I've also added an extra \gls{vspace} at the end of the
environment and a paragraph break.
Since we've used \glsni{refstepcounter} instead of 
\glsni{stepcounter} we can cross-reference our
\envname{exercise} environment:\reportpagebreak
\begin{code}
\begin{alltt}
Exercise\glsni{tildechar}\glsni{ref}\marg{ex:simple} is a simple exercise.

\glsni{begin}\marg{exercise}\marg{Simple Exercise}
\glsni{label}\marg{ex:simple}\glsni{percentchar}
This is a simple exercise.
\glsni{end}\marg{exercise}
\end{alltt}
\end{code}%
This produces 
the following output:
\begin{result}[newcounter.html]
Exercise~1 is a simple exercise.
\par\vspace{\baselineskip}\noindent\textbf{Exercise 1 (Simple Exercise)}%
\begin{itshape}\par\vspace{\baselineskip}\noindent
This is a simple exercise.
\end{itshape}
\end{result}

The counter representation can be changed by redefining
\cmdname{theexercise}\faq{Redefining counters'
\cmdname{the}-commands}{the-commands} using the \gls{renewcommand}
command described in \sectionref{sec:renewcom}.
The following commands can be used to display the counter:

\begin{fwlist}{\cmdname{fnsymbol}\marg{\meta{counter}}}
\fwitem{\gls{arabic}\marg{\meta{counter}}} Arabic numeral (1, 2, 3, \ldots)

\fwitem{\gls{Roman}\marg{\meta{counter}}} Upper case Roman numeral 
(I, II, III, \ldots)

\fwitem{\gls{roman}\marg{\meta{counter}}} Lower case Roman numeral 
(i, ii, iii, \ldots)

\fwitem{\gls{alph}\marg{\meta{counter}}} Lower case letter 
(a, b, c, \ldots, z)

\fwitem{\gls{Alph}\marg{\meta{counter}}} Upper case letter 
(A, B, C, \ldots, Z)

\fwitem{\gls{fnsymbol}\marg{\meta{counter}}}
A footnote symbol 
(\makeimg{asterisk dagger double dagger section mark
paragraph mark double bar double asterisk two single daggers
two double daggers}{\footnotesymbols})
\end{fwlist}

\xminisec{Example:}

To make the chapter numbers appear as upper case Roman 
numerals you would do:
\begin{codeS}
\gls{renewcommand}\marg{\gls{thechapter}}\marg{\glsni{Roman}\marg{chapter}}
\end{codeS}

You may have noticed that \gls{newcounter} has an optional
argument \meta{outer-counter}.  This is for use if you require the
new counter to be reset every time \meta{outer-counter} is
incremented\faq{Master and slave counters}{addtoreset}.  For example,
the section numbers in the \icls{scrbook} class are dependent on the
chapter numbers.  Each time a new chapter is started, the section
numbers are reset.  Suppose we want our \texttt{exercise} counter to
be dependent on the \icounter{chapter} counter, we would do
\begin{codeS}
\gls{newcounter}\marg{exercise}\oarg{chapter}
\end{codeS}

Note that if you make a counter dependent on another counter like
this, the default action of \cmdname{the}\meta{counter} remains the
same, so \cmdname{theexercise} won't print the chapter number.
To make the chapter number appear as well, we need to redefine
\cmdname{theexercise} (recall \sectionref{sec:renewcom}):
\begin{codeS}
\gls{renewcommand}\marg{\cmdname{theexercise}}\marg{\glsni{thechapter}.\glsni{arabic}\marg{exercise}}
\end{codeS}%
Notice the use of \gls{thechapter} instead of,
say, \verb|\arabic{chapter}|.  This way
we don't need to keep track of the \icounter{chapter} counter format.

\xminisec{Example (Footnote Markers):}

The \icounter{footnote} counter is reset at the start of each
chapter but by default the chapter number isn't displayed in
\gls{thefootnote}. In \latexhtml{this book}{the PDF version of
this document} \glsni{thefootnote} was redefined so that it
displays the chapter number:
\begin{codeS}
\gls{renewcommand}\marg{\glsni{thefootnote}}\marg{\glsni{thechapter}.\glsni{arabic}\marg{footnote}}
\end{codeS}

\begin{exercise}{Using Counters}{ex:counters}
Modify the document from \exerciseref{ex:newenv} so that
the \envname{exercise} environment has a counter.
Make the counter dependent on the chapter.
You can \downloadorview{counters} an example.
\end{exercise}

%%%%%%%%%%%%%%%%%%%%%% APPENDICES %%%%%%%%%%%%%%%%%%

\appendix

%%%%%%%%%%%%%%%%%%% Installing Packages %%%%%%%%%%%%%%

\setnode{installsty}
\chapter{Downloading and Installing Packages}
\label{ch:installsty}

New \LaTeX\ packages are being created all the time\faq{Installing things
on a (La)TeX system}{instpackages}, so you may find that there are some
packages that you don't have on your installation. In this case, if
you don't have the package you want, you can download
it\faq{Installation using MiKTeX package manager}{inst-miktex*} from
\gls{ctan}. Before discussing installing new packages, it is
a good idea for you to understand the \gls{tds}\faq{What is the TDS?}{tds}.

All the files that make up the \TeX\ distribution are stored in a
standard hierarchical structure. The root directory of the main
distribution is usually called \texttt{texmf} or \texttt{texmf-dist}. Its location
depends on your system. For example, if you are using
\itexdistro{TeX~Live}~2012 on UNIX/Linux, it will probably be located in
\texttt{/usr\slash local\slash texlive\slash 2012\slash texmf-dist} or if you are using \itexdistro{MiKTeX} it
may be located in \verb|c:\texmf| or \texttt{c:\backslashsym Program
Files\backslashsym texmf}. Whichever system you are using, I~shall
refer to this directory as \meta{TEXMF}. So, if you are using
\texdistro{TeX~Live}~2012, \meta{TEXMF}\texttt{\slash doc} refers to the directory
\texttt{/usr\slash local\slash texlive\slash 2012\slash texmf-dist\slash doc}, or if you are using
\texdistro{MiKTeX}, \meta{TEXMF}\verb|\doc| refers to the folder
\verb|c:\texmf\doc| or \texttt{c:\backslashsym Program
Files\backslashsym texmf\backslashsym doc}.
In general, you should not make any modifications to the
\meta{TEXMF} directory tree as it will get overridden whenever you
update your \TeX\ distribution.

You should also have a local texmf tree. Again, the location of the
local texmf tree depends on your system. If you are using
\texdistro{TeX~Live}, it may be 
\texttt{/usr\slash local\slash texlive\slash texmf-local}. 
If you are using \texdistro{MiKTeX}, it
may be \texttt{c:\backslashsym localtexmf} or
\texttt{c:\backslashsym Program Files\backslashsym localtexmf}.
Whichever system you are using, I~shall refer to this directory as
\meta{TEXMF-LOCAL}. There is also the \meta{TEXMF-HOME} directory.
On UNIX-like systems this is usually \verb|~/texmf|. On Windows it's
usually in your user folder. This is the one where you typically
install any new classes or packages.

These directories must all have the same structure. The principle
sub-directories relating to \LaTeX\ are illustrated in
\figureref{fig:tds}. It may be that your \meta{TEXMF-HOME} directory
doesn't exist or doesn't contain some of these sub-directories, if
so, you will need to create them.

You can use the \iappname{kpsewhich} application to find out the 
locations of \meta{TEXMF-LOCAL} and
\meta{TEXMF-HOME}. Since \appname{kpsewhich} is a command-line
application, you will need a \htmlref{command prompt or terminal}{sec:terminal} 
open\latex{ (see \sectionref{sec:terminal})}.
At the command prompt, type
\begin{verbatim}
kpsewhich -var-value=TEXMFHOME
\end{verbatim}
to display the location of \meta{TEXMF-HOME} or
\begin{verbatim}
kpsewhich -var-value=TEXMFLOCAL
\end{verbatim}
to display the location of \meta{TEXMF-LOCAL}.
(Remember to press the enter \enter\ key at the end of the line.)

\refstepcounter{object}\label{obj:docdir}%
The documentation for \LaTeX\ classes and packages can be found
in the \texttt{doc\slash latex} sub-directories: \meta{TEXMF}\texttt{\slash doc\slash latex},
\meta{TEXMF-LOCAL}\texttt{\slash doc\slash latex} and \meta{TEXMF-HOME}\texttt{\slash doc\slash latex}.

\begin{figure}[htbp]
\figconts
  {pictures/tds}
  {%
    \caption{The \protect\TeX\protect\space Directory Structure (TDS) Showing 
    the Main \protect\LaTeX-Related Sub-Directories.}%
  }%
  {fig:tds}
\end{figure}

Some packages are supplied in this format.\footnote{Complete list at \url{http://mirror.ctan.org/install/macros/latex/contrib/}.} For example, the package
\sty{sample-package} may be distributed in a compressed file \texttt{sample-package.tds.zip},
which contains the files
\begin{verbatim}
doc/latex/sample-package/sample-package.pdf
tex/latex/sample-package/sample-package.sty
tex/latex/sample-package/sample-foo.sty
tex/latex/sample-package/sample-bar.sty
\end{verbatim}
In this case all you need to do is decompress the contents of 
the archive into the
\meta{TEXMF-LOCAL} or \meta{TEXMF-HOME} directory.

\refstepcounter{object}\label{obj:runtexhash}%
On older \TeX-distributions, you would then need to \htmlref{refresh the
\TeX\ database}{sec:refresh}\latex{ (described in
\sectionref{sec:refresh}\ifscreenorother{}{\ on
page~\pageref{sec:refresh}})}. With new distributions, you don't
need to do this if you are installing a new package into your
\meta{TEXMF-HOME} directory.

\xminisec{Example (Unix-Like):}
To install \texttt{sample-package.tds.zip} (assuming you're in the
same directory as that file):
\begin{verbatim}
unzip -d ~/texmf sample-package.tds.zip
\end{verbatim}

\setnode{dtxins}
\section{DTX and INS Files}
\label{sec:dtx}
\faq{Documented LaTeX sources (.dtx files)}{dtx}Not all packages are
provided in the \gls{tds} format. Instead (or additionally) many are supplied
with the code and documentation all bundled together in one file.
This file usually has the extension \texttt{.dtx}, and it usually
comes with an installation script that has the extension
\texttt{.ins}. Once you have downloaded the \texttt{.dtx} and
\texttt{.ins} files, you will then have to extract the code before
you can use it. Let's go back to the previous example. The package
\sty{sample-package} is now distributed in a DTX file, so the
\texttt{sample-package.zip} archive contains the files
\begin{verbatim}
sample-package.dtx sample-package.ins
\end{verbatim}
(with hopefully a \texttt{README} or \texttt{INSTALL} file). Note
that this archive, unlike the TDS one, doesn't contain any
\texttt{.sty} files. The documentation source and the package code
(\texttt{sample-package.sty}, \texttt{sample-foo.sty} and
\texttt{sample-bar.sty}) are all contained in the file
\texttt{sample-package.dtx}. This is how to extract them:
\begin{enumerate}
\item Extract the contents of \texttt{sample-package.zip} to a
temporary directory.

\item\label{itm:ins} Run \LaTeX\ on the file \texttt{sample-package.ins}. If
you are using a \latexhtml{terminal}{\htmlref{terminal}{sec:terminal}}, you can type the
following at the command prompt:
\begin{verbatim}
latex sample-package.ins
\end{verbatim}
If you are using a front-end, such as TeXWorks, open the 
\texttt{.ins} file
(for example \texttt{sample-package.ins}), and click on the build/typeset button.

This will create the files containing the package code. In this
example it will create the main package file \texttt{sample-package.sty}
and supplementary packages \texttt{sample-foo.sty} and \texttt{sample-bar.sty}.

\item Make a sub-directory of \meta{TEXMF-LOCAL}\texttt{/tex/latex}%
\footnote{or \meta{TEXMF-LOCAL}\texttt{\backslashsym tex\backslashsym latex} on 
Windows} in which to place these files. In this example, the package is
called \dq{sample-package}, so make a sub-directory called \texttt{sample-package}.

\item Move the files created in \objectref{Step}{itm:ins} into
the new sub-directory you created in the previous step.

\item\label{itm:dtx} Run \iPDFLaTeX\ on the file \texttt{sample-package.dtx}.
(The same as in \objectref{Step}{itm:ins}, but use the file
\texttt{sample-package.dtx} instead of \texttt{sample-package.ins}.) This will 
create a file called \texttt{sample-package.pdf}. You may need to repeat
this step to ensure that the cross references are up-to-date.
Check the \texttt{README} file or \texttt{INSTALL} file to see
if there is anything else you need to do. (If you have downloaded
the package from CTAN, it's possible that the documentation has 
already been supplied, as package authors are encouraged to supply
a PDF version of the documentation for on-line viewing. If so,
you can omit this step.)

\item  Make a sub-directory of \meta{TEXMF-LOCAL}\texttt{/doc/latex}%
\footnote{or \meta{TEXMF-LOCAL}\texttt{\backslashsym doc\backslashsym latex} on 
Windows} in which to place the documentation. In this example, the package is
called \dq{sample-package}, so make a sub-directory called
\texttt{sample-package}.

\item Move the files created in \objectref{Step}{itm:dtx} into
the new sub-directory you created in the previous step.

\end{enumerate}

As mentioned \htmlref{above}{obj:runtexhash}, on older
\TeX-distributions, you would then need to \htmlref{refresh the
\TeX\ database}{sec:refresh}, but this isn't required for
\meta{TEXMF-HOME} installs on new distributions.

\setnode{texhash}
\section{Refreshing the \texorpdfstring{\protect\TeX}{TeX}\protect\space Database}
\label{sec:refresh}

On older \gls{tex} distributions you had to refresh the \TeX\
database whenever you installed new classes or packages. With newer
installations you don't need to do this if you install them in your
\meta{TEXMF-HOME} directory, except under certain circumstances (for
example, you're using using a networked drive). If it turns out that
\TeX\ can't find a new class or package you have installed in
\meta{TEXMF-HOME} you will need to update the database using the
\iappname{texhash} (or \iappname{mktexlsr}) application. This is a
command-line application, so you need a \htmlref{terminal or command
prompt}{sec:terminal}\latex{ (see \sectionref{sec:terminal})}.

For example, on UNIX/Linux, to update \meta{TEXMF-HOME} (the
directory \verb|~/texmf|) you need to type the following at the
command prompt:
\begin{verbatim}
texhash ~/texmf
\end{verbatim}

If you are using a modern \TeX\ distribution, such as
\itexdistro{MiKTeX}, \itexdistro{TeX~Live} or \itexdistro{MacTeX} there 
should be a package manager that has a package installation
and refresh facility. For example, \texdistro{TeX~Live} comes with
the TeX~Live Manager (\iappname{tlmgr} or \iappname{mactlmgr}) and
recent versions of \texdistro{MiKTeX} have an application called
\appname{MiKTeX Update Wizard} which can automatically download and
install known packages.

If you experience any problems, contact your system administrator for help
or try one of the resources listed in \appendixref{ch:help}.

\xminisec{Related \gls{ukfaq} topics:}
\begin{itemize}
\item \faqlink{Installing things on a (La)TeX system}{instpackages}
\item \faqlink{Installing files \dq{where (La)TeX can find them}}{wherefiles}
\item \faqlink{Installation using MiKTeX package manager}{inst-miktex*}
\item \faqlink{\dq{Temporary} installation of (La)TeX files}{tempinst}
\item \faqlink{\dq{Private} installations of files}{privinst}
\end{itemize}


%%%%%%%%%%%%%%%%%%%%%% ERRORS %%%%%%%%%%%%%%%%%%%%%%%%%%%

\setnode{commonerrors}
\chapter{Common Errors}
\label{ch:errors}

\begin{itemize}
\item
If you're running \LaTeX\ from a \gls{terminal} and the only message that gets displayed is:
\begin{verbatim}
latex: Command not found.
\end{verbatim}
or
\begin{verbatim}
Bad command or file name
\end{verbatim}
then you have either mistyped the command name, or you
don't have \LaTeX\ installed on your computer, or your path hasn't
been set up correctly.  First check that you have typed the command
correctly, then check to see if you have \TeX\ installed.  Failing
that, contact your system administrator for help or try one of the
resources listed in \appendixref{ch:help}.

\item
If you're running \LaTeX\ from a \gls{terminal} and you get the message (or something similar):
\begin{verbatim}
This is TeX, Version 3.14159 (Web2C 7.3.1)
! I can't find file `sample'.
<*> sample
          
Please type another input file name: 
\end{verbatim}
then you have either misspelt the filename or you are in the wrong
directory.  If you have misspelt the filename, simply type in the
correct name at the prompt.  If you are in the wrong directory or you
want to quit, type \texttt{X} followed by the return character \enter.
To check you are in
the right directory, on a Unix-like system you can type:
\begin{verbatim}
ls
\end{verbatim}
This will list the contents of
the directory.  If you are certain that you have spelt the filename
correctly and that you are in the right directory, there may be
something wrong with your path, in which case contact your system
administrator.

\item
Error messages will usually look something like:
\begin{verbatim}
! Undefined control sequence.
l.1 \docmentclass
                 [12pt]{scrartcl}
? 
\end{verbatim}
The first line is the error message.  In this example I~have misspelt
the command \cmdname{documentclass}.  The next line begins with
\texttt{l.}\ followed by a number.  This is the line number in the
source code where the error occurred.  In this case the error
occurred on line~1.  Following the line number is the input line
\LaTeX\ has processed so far, and staggered on the next line is the
remainder of the input line.

Here's another example.  Suppose line~8 of my 
\gls{source}
looks like:
\begin{verbatim}
The date today is: \toady, which is nice to know.
\end{verbatim}
The error in this case is the misspelling of the command 
\cmdname{today}.  The error message will appear as follows:

\begin{latexonly}
\noindent
\begin{tabular}{l@{}l}
\multicolumn{2}{l}{%
\texttt{!\ Undefined control sequence.}\rlap{\color{blue}%
\tikz[baseline=-1ex]{\draw[latex-] (0em,0ex) -- (4em,2ex);
\pgftext[left,center,at={\pgfpoint{4em}{2ex}}]{Error Message}
}}}\\
\texttt{l.}\smash{\llap{\color{blue}%
\tikz[baseline]{\draw[latex-] (0em,0ex) -- (-3em,-4ex);
\pgftext[right,top,at={\pgfpoint{-3em}{-3ex}}]{\shortstack{Line\\Number}}}}}%
\texttt{8 The date today is }%
\smash{\llap{\color{blue}%
\tikz[baseline]{\draw[latex-] (0em,0ex) -- (-3em,-3ex);
\pgftext[right,top,at={\pgfpoint{-3em}{-3ex}}]{Error}}}}%
\cmdname{toady}%
\smash{\rlap{\color{blue}%
\tikz[baseline]{\draw[latex-] (0em,0ex) -- (0em,-6ex);
\pgftext[left,top,at={\pgfpoint{0em}{-6.5ex}}]{This is how far \LaTeX\ has got}}}}\\
 &\texttt{, which is nice to know.}\smash{\rlap{\color{blue}%
\tikz[baseline]{\draw[latex-] (0em,0.5ex) -- (3em,0.5ex);
\pgftext[left,center,at={\pgfpoint{3em}{0.5ex}}]{\shortstack{Rest of\\the line}}}}}\\
\texttt{?}\color{blue}%
\tikz[baseline]{\draw[latex-] (0em,0ex) -- (2em,-4ex);
\pgftext[left,top,at={\pgfpoint{2em}{-4ex}}]{\LaTeX\ prompt}}
\end{tabular}

\vspace{\baselineskip}
\end{latexonly}
\html{\htmladdimg[alt="annotated error message"]{pictures/errormessage.png}}

At the \LaTeX\ prompt, you can either type \texttt{h} for a help 
message, or type \texttt{x} to exit \LaTeX\ and go back to your 
source code and fix the problem.
\end{itemize}

There follows below a list of common error messages. If your
problem isn't listed there, try the \gls{ukfaq}.

%%%%%%%%%%%%%%%%%%%% * (No message) %%%%%%%%%%%%%%%%%%%%%%%%%%%

\setnode{asteriskprompt}
\section{* (No message, just an asterisk prompt!)}

You've gone into \gls{tex}!  This is probably because 
you've forgotten the \texttt{\gls{end}\marg{document}}.  The asterisk is the 
\TeX\ prompt.  At this point the best thing to do is to abort the
\TeX\ run.

%%%%%%%%%%%%% Argument of \cline has an extra } %%%%%%%%%%%%%%%%

\setnode{clineextrabrace}
\section{Argument of \textbackslash cline has an extra \}}

If this error occurred on the first line in your \gls{env-tabular}
\glsterm{environment}, you may have forgotten the
\glsterm{argument} to the \gls*{env-tabular} environment.

%%%%%%%%%%%%% Argument of \multicolumn has an extra } %%%%%%%%%%%%%%%%

\setnode{multicolumnextrabrace}
\section{Argument of \textbackslash multicolumn has an extra \}}

If this error occurred on the first line in your \gls{env-tabular}
\glsterm{environment}, you may have forgotten the
\glsterm{argument} to the \gls*{env-tabular} environment.

%%%%%%%%%%%%%%%% \begin{...} ended by \end{...} %%%%%%%%%%%%%%%%%%%%

\setnode{mismatchedend}
\section{\textbackslash begin\{\texorpdfstring{\ldots}{...}\} ended by 
\textbackslash end\{\texorpdfstring{\ldots}{...}\}}

The beginning of your environment doesn't have a matching end.
\begin{itemize}
\item Check to make sure you have spelt the name of the environment 
correctly.

You will get this error message if you do, say,
\begin{alltt}\wrong
\glsni{end}\marg{docment}
\end{alltt}
instead of
\begin{alltt}\correct
\glsni{end}\marg{document}
\end{alltt}

\item Check that for every \gls{begin} you have a corresponding 
\gls{end} with the same name.
\end{itemize}

%%%%%%%%%%%%%%%%%%% Bad math environment delimiter %%%%%%%%%%%%%%%%

\setnode{badmathdelim}
\section{Bad math environment delimiter}

Only a certain type of character may be used as a 
\htmlref{delimiter}{sec:delimiters} (for example,
\gls{openparen} \gls{closeparen} \gls{delimiter.opensq}
\gls{delimiter.closesq}), check which one you have specified.  This error may 
also occur if you have forgotten a \gls{right} or not used it in the
same \glslink{group}{scope}.
(Remember to use a \gls{delimiter.periodchar} if you want an invisible delimiter) or 
you may have forgotten to end your array environment with 
\verb|\end{array}|.

%%%%%%%%%%%%%%%%%%%% Can only be used in preamble %%%%%%%%%%%%%%%%%%

\setnode{onlypreamble}
\section{Can only be used in preamble.}

Some commands, such as \gls{usepackage} may only appear in the 
\gls{preamble}.  Check to see where you have put 
it.  For example, this error will be caused by doing:
\begin{alltt}
\glsni{documentclass}\marg{scrartcl}

\glsni{begin}\marg{document}\wrong

\glsni{usepackage}\marg{graphicx}
\end{alltt}
instead of
\begin{alltt}
\glsni{documentclass}\marg{scrartcl}

\glsni{usepackage}\marg{graphicx}\correct

\glsni{begin}\marg{document}
\end{alltt}

%%%%%%%%%%%%%%%%%% Command ... already defined %%%%%%%%%%%%%%%%%%%%%

\setnode{alreadydef}
\section{Command \texorpdfstring{\protect\ldots}{...}\protect\space already defined}

You have tried to define a \gls{command} which 
already exists.  Try giving it a different name.  Remember never to
redefine a command if you don't know what the command originally
does.

Alternatively, you have tried to define an
\gls{environment} which already exists.  Give the new
environment a different name.  Again, never redefine an environment
where you don't know what the original environment does.

%%%%%%%%%%%%% Display math should end with $$ %%%%%%%%%%%%%%%%%%%%%

\setnode{missingenddispmath}
\section{Display math should end with \$\$}

You may have a dollar sign (\gls{dollarchar}) in a displayed maths
environment, such as the \gls{env-equation} environment.  
Remember that \verb|$| is short hand for \verb|\begin{math}| or 
\verb|\end{math}|, so you can't end one of the other environments 
with a \verb|$|. (This error message is in fact a bit confusing,
as it seems to be suggesting that you end a displayed maths
environment with \verb|$$| which you also shouldn't do\faq{Why use
\glsnl{begindispmath}\ldots\glsnl{enddispmath} in place of \$\$\ldots\$\$}{dolldoll}.)

%%%%%%%%%%%%%% Environment ... undefined %%%%%%%%%%%%%%%%%%%%%%%%%

\setnode{envundefined}
\section{Environment \texorpdfstring{\protect\ldots}{...}\protect\space undefined}

\LaTeX\ doesn't recognise the environment you have specified.

\begin{itemize}
\item Check you have spelt the environment name correctly.

You will get this error if you do, say,
\begin{alltt}\wrong
\glsni{begin}\marg{docment}
\end{alltt}
instead of
\begin{alltt}\correct
\glsni{begin}\marg{document}
\end{alltt}

\item If it's your own environment, check you have defined the
environment before using it.

\item If the environment is defined in a package, check you 
have included the package using the 
\htmlref{\cmdname{usepackage}}{sec:packages} command.

\end{itemize}

%%%%%%%%%% Extra alignment tab has been changed to \cr %%%%%%%%%%%%

\setnode{tab2cr}
\section{Extra alignment tab has been changed to \textbackslash cr}

You have too many ampersands (\gls{ampchar}) in one row.  The most 
probable cause is that you have forgotten the end of row command
\gls{tab.dbbackslashchar} on the previous row.  Remember also that if you have a
\gls{multicolumn} command to span more than one column, you
should have fewer \glsni{ampchar}s in that row. This error can also occur
from a confusion over the two meanings of \gls{dbbackslashchar}: a
line break within a paragraph cell and a row break.  In which case,
you need to use \gls{tabularnewline} instead of
\gls{tab.dbbackslashchar}.

%%%%%%%%%%%%%%%%%%%%%%%%% Extra \right %%%%%%%%%%%%%%%%%%%%%%%%%%

\setnode{extraright}
\section{Extra \textbackslash right}

There are a number of possible causes.  The most probable is that you
have a \gls{right} that doesn't have a matching \gls{left}.
(Remember left comes before right.)  Another possible cause is that
you have missed out \verb|\end{array}|.  (Remember that
\glspl{environment} provide implicit
\glsing{group}, and \glsni{left} and its matching
\glsni{right} must appear within the same group level.)

%%%%%%%%%%%%%% File ended while scanning use of ... %%%%%%%%%%%%%

\setnode{fileendedwhilescanning}
\section{File ended while scanning use of \texorpdfstring{\ldots}{...}}

The most usual cause of this error is a missing closing brace.

You will get this error if you do, say,
\begin{alltt}\wrong
\glsni{end}\glsni{leftbracechar}document
\end{alltt}
instead of
\begin{alltt}\correct
\glsni{end}\marg{document}
\end{alltt}

%%%%%%%%%%%%%%%%%%%%%% File not found %%%%%%%%%%%%%%%%%%%%%%%%%%

\setnode{filenotfound}
\section{File not found}

\LaTeX\ can't find the file you have specified.  You will be given
the opportunity to type in the correct filename at the prompt.  If
you want to quit, simply type \texttt{X} followed by the return
character \enter.

\begin{itemize}

\item Make sure that you have spelt the filename correctly.

This error will be caused by, say,
\begin{alltt}\wrong
\gls{documentclass}\marg{scrarticle}
\end{alltt}
instead of
\begin{alltt}\correct
\gls{documentclass}\marg{scrartcl}
\end{alltt}

If this is the case, simply type in the correct name at the prompt 
(followed by the return character \enter) then go back and fix the
spelling in the \gls{source}.

\item Make sure that the file is in the same directory as your
document or in the \LaTeX\ path.  If the file is in another directory
(not in the \LaTeX\ path), you will need to specify the pathname, but
remember that when using \LaTeX\ under Windows, you need to use a
forward slash (\glsi{dir.slash}) as the directory divider, as a backslash would be
interpreted as a command.  For example, if you have a file called
\texttt{shapes.pdf} in the subdirectory \texttt{pictures} then you
would get a \dq{file not found} error message if you did
\begin{alltt}\wrong
\gls{includegraphics}\marg{shapes}
\end{alltt}
instead of
\begin{alltt}\correct
\gls{includegraphics}\marg{pictures/shapes}
\end{alltt}

\item If the file is a \htmlref{package}{sec:packages} or
\htmlref{class file}{sec:cls}, it's possible that you don't have that
file installed on your computer.  If this is the case you will need
to download and install it as described in
\appendixref{ch:installsty}.  Remember that you need to
refresh the database after installing a new package or class file.

\end{itemize}

%%%%%%%%%%%%%%%%%%%% Illegal character in array arg %%%%%%%%%%%%%%%%%%

\setnode{illegalarraychar}
\section{Illegal character in array arg}

You have used a character in the \gls{argument} of a
\gls{env-tabular} or \gls{env-array} environment that is not
allowed. The standard available characters are: \texttt{r} (right
justified), \texttt{l} (left justified), \texttt{c} (centred) and
\texttt{p}, as well as \gls{atchar}\marg{\meta{inter-col text}}.
Remember that if you want to use the \gls{gtcol}\marg{\meta{decl}} or 
\gls{ltcol}\marg{\meta{decl}} specifiers, you must include the
\isty{array} package.

This error will also occur if you have forgotten the argument to the
\gls*{env-tabular} or \gls*{env-array} environment.

%%%%%%%%%%%%%%%%%%% Illegal parameter number in definition %%%%%%%%%%

\setnode{illegalparamnum}
\section{Illegal parameter number in definition}

You have referred to a \htmlref{parameter (argument)
number}{sec:newcomarg} that is greater than the number of parameters
you have specified. For example, suppose you defined the command to
have only one parameter, then you can't use \gls{hashchar}\texttt{2}
which refers to the second, non-existent, parameter. Remember that
you need to specify how many parameters you want in the
\gls{optional} to \glsni{newcommand}, otherwise it will be assumed
that the command has no arguments.

%%%%%%%%%%%%%%%%%%%%% Illegal unit of measure %%%%%%%%%%%%%%%%%%%%%%%

\setnode{illegalunit}
\section{Illegal unit of measure (pt inserted)}

You have either not specified a unit when giving a \gls{length} (even zero
lengths must have a unit) or you have specified an invalid unit or
you have misspelt the unit.  Available units are listed in
\tableref{tab:units}.

%%%%%%%%%%%%%%%%% Lonely \item %%%%%%%%%%%%%%%%%%%%%%%%%

\setnode{lonelyitem}
\section{Lonely \textbackslash item}

The command \gls{item} may only appear in one of the list making
environments (such as \gls{env-itemize}).  Make sure you haven't
forgotten your environment.

%%%%%%%%%%%%%%% Misplaced alignment tab character & %%%%%%%%%%%%%%%%%

\setnode{misplacedtab}
\section{Misplaced alignment tab character \&}

You have used the special character \gls{ampchar} where you shouldn't
have.  Recall from \sectionref{sec:chars} that if you want an \& sign 
to appear you need to do \gls{amp} not just \glsni{ampchar}.

You would have got this error message if you had done, say,
\begin{alltt}\wrong
\glsni{ampchar} our equipment
\end{alltt}
instead of
\begin{alltt}\correct
\glsni{amp} our equipment
\end{alltt}

%%%%%%%%%%%%%%%%%%%%% Missing } inserted %%%%%%%%%%%%%%%%%%%%%%%%%%%

\setnode{missingendbrace}
\section{Missing \} inserted}

You have missed a closing curly brace, or you may have missed out an
argument.

\xminisec{Example:} 
If the following line occurs in a tabular environment:
\begin{alltt}\wrong
 \gls{ampchar} \gls{multicolumn}\marg{2}\marg{c}\gls{tab.dbbackslashchar}
\end{alltt}
this will produce the error. (The third argument to 
\gls{multicolumn} has been omitted.)

%%%%%%%%%%%%%%%%%%%%% Missing $ inserted %%%%%%%%%%%%%%%%%%%%%%%%%%%

\setnode{missingdollar}
\section{Missing \$ inserted}

This message can be caused by a number of errors:
\begin{itemize}
\item You might have missed the beginning of one of the mathematics
\glspl{environment} (that is, you've used a
\gls{command} that must only appear in maths mode).

\item You may have typed \gls{dollarchar} instead of \gls{dollar} (you
actually want a dollar symbol to appear).  Recall from
\sectionref{sec:chars} that if you
want a \$ sign to appear you need to do \glsni{dollar} not just \glsni{dollarchar}.

You would have got this error message if you had done, say,
\begin{alltt}\wrong
expenditure came to \glsni{dollarchar}2000.00
\end{alltt}
instead of
\begin{alltt}\correct
expenditure came to \glsni{dollar}2000.00
\end{alltt}

\item You may have missed the end of a mathematics environment, or
you may have a paragraph break within an in-line or
displayed maths environment, where it is
not permitted.  Make sure you don't have any blank lines within the
environment. If you want a blank line in your code to make it easier
to edit, try having a percent sign at the start of an empty line to
ensure that the line is ignored by \LaTeX. For example:
\begin{alltt}
\glsni{begin}\marg{equation}
\glsni{percentchar}
E = mc\glsni{circumchar}2
\glsni{percentchar}
\glsni{end}\marg{equation}
\end{alltt}
\end{itemize}

%%%%%%%%%%%%%%% Missing \begin{document} %%%%%%%%%%%%%%%%%

\setnode{missingbegindoc}
\section{Missing \textbackslash begin\{document\}}

You have put some text outside of the document 
\glsterm{environment}.  Check the following:
\begin{itemize}
\item You have remembered \verb|\begin{document}|

This error would be caused by, say,
\begin{alltt}
\glsni{documentclass}\marg{scrartcl}
\wrong
This is a simple document
\end{alltt}
instead of
\begin{alltt}
\glsni{documentclass}\marg{scrartcl}

\glsni{begin}\marg{document}\correct
This is a simple document.
\end{alltt}

\item You haven't placed any text before \verb|\begin{document}|. 
For example:
\begin{alltt}
\glsni{documentclass}\marg{scrartcl}

This is a simple document\wrong 
\glsni{begin}\marg{document}
\end{alltt}
instead of
\begin{alltt}
\glsni{documentclass}\marg{scrartcl}

\glsni{begin}\marg{document}\correct
This is a simple document
\end{alltt}

\item You haven't missed out the backslash at the start of either 
\gls{documentclass} or \glsni{begin}\marg{document}

This error would be caused by, say,
\begin{alltt}\wrong
documentclass\marg{scrartcl}
\end{alltt}
instead of
\begin{alltt}\correct
\glsni{documentclass}\marg{scrartcl}
\end{alltt}
\end{itemize}

%%%%%%%%%%%%%%%%%%% Missing delimiter %%%%%%%%%%%%%%%%%%%%%%%

\setnode{missingdelim}
\section{Missing delimiter}

You have forgotten to specify the type of delimiter you want (for
example, 
\gls{openparen} \gls{closeparen} \gls{delimiter.opensq} \gls{delimiter.closesq}
\gls{leftbrace} \gls{rightbrace}). (Remember to use a \gls{delimiter.periodchar} if you want 
an invisible delimiter, and remember that if you want a 
curly brace, you must have a backslash followed by the curly brace.)

\xminisec{Example:}

This error will occur if you do, say,
\begin{alltt}\wrong
f(x) = \glsni{left}\glsni{leftbracechar}
\glsni{begin}\marg{array}\marg{ll}
0 \glsni{ampchar} x \cmdname{leq} 0\gls{tab.dbbackslashchar}
1 \glsni{ampchar} x > 1
\glsni{end}\marg{array}
\glsni{right}.
\end{alltt}
instead of
\begin{alltt}\correct
f(x) = \glsni{left}\glsni{leftbrace}
\glsni{begin}\marg{array}\marg{ll}
0 \glsni{ampchar} x \cmdname{leq} 0\gls{tab.dbbackslashchar}
1 \glsni{ampchar} x > 1
\glsni{end}\marg{array}
\glsni{right}.
\end{alltt}

%%%%%%%%%%%%%%%% Missing \endcsname inserted %%%%%%%%%%%%%%%%

\setnode{missingendcsname}
\section{Missing \textbackslash endcsname inserted}

This is a \gls{tex} error rather than a \LaTeX\ error
which makes it harder to determine the cause, however it can be
caused by placing a backslash in front of the name of an
\gls{environment}.  (Remember that
\gls{environment} names do not contain a backslash.)

This error will be caused by, say,
\begin{alltt}\wrong
\glsni{begin}\marg{\glsni{sffamily}}
\end{alltt}
instead of
\begin{alltt}\correct
\glsni{begin}\marg{sffamily}
\end{alltt}

%%%%%%%%%%%%%% Missing \endgroup inserted %%%%%%%%%%%%%%%%%%

\setnode{missingendgroup}
\section{Missing \textbackslash endgroup inserted}

A number of things could have caused this.  You may have missed out
the end of an \gls{environment}, or you may have an
environment inside of another environment it's not allowed to be in.
For example, this error can be caused by placing an
\envname{eqnarray} environment inside a \envname{displaymath}
environment, which is not allowed. (But, of course, you haven't
used either of those obsolete environments~\cite{l2tabu}, have you?) 

%%%%%%%%%%%%%%% Missing number %%%%%%%%%%%%%%%%%%

\setnode{missingnum}
\section{Missing number, treated as zero}

\LaTeX\ is expecting a number. If your command takes more than one
\gls{argument}, check to make sure the arguments are in
the correct order.  For example, if you are using a
\gls{env-minipage} environment, you might have
omitted the \gls{mandatory} which
specifies the width of the minipage, or you may have the
\glspl{optional} the wrong way round.  The placement specifier
should come first, followed by the height.

If you are using \gls{addtocounter} or \gls{setcounter} 
remember that the second argument must be a number, so if you want 
the value of a counter as the argument you must use \gls{value}.
This error can be caused by, say,
\begin{alltt}\wrong
\glsni{setcounter}\marg{exercise}\marg{chapter}
\end{alltt}
instead of
\begin{alltt}\correct
\glsni{setcounter}\marg{exercise}\marg{\glsni{value}\marg{chapter}}
\end{alltt}

%%%%%%%%%%%%%% Paragraph ended %%%%%%%%%%

\setnode{parended}
\section{Paragraph ended before \textbackslash begin was complete}

You've probably missed a closing brace at the end of the argument to
\gls{begin}.  This error will be caused by, say,
\begin{alltt}\wrong
\glsni{begin}\glsni{leftbracechar}document
\end{alltt}
instead of
\begin{alltt}\correct
\glsni{begin}\marg{document}
\end{alltt}

%%%%%%%%%%%%%%%%%%% Runaway argument %%%%%%%%%%%%%%%%%%%%%%%%%%%%

\setnode{runawayarg}
\section{Runaway argument}

There are a number of possible causes of this error:

\begin{itemize}

\item Paragraph breaks are not permitted in the
\glspl{argument} of \glspl{short}.
If there is a corresponding \gls{environment} then you
should use that instead.  For example, this error message will be
generated by doing, say,
\begin{alltt}
\glsni{textbf}\glsni{leftbracechar}This is a simple document.
Here is the first paragraph.
\wrong
Here is the second paragraph.\glsni{rightbracechar}
\end{alltt}
instead of\bookpagebreak
\begin{alltt}
\glsni{begin}\marg{bfseries}
This is a simple document.
Here is the first paragraph.
\correct
Here is the second paragraph.
\glsni{end}\marg{bfseries}
\end{alltt}

\item The closing brace of a 
\gls{mandatory} is missing:
This error will be caused by, say,
\begin{alltt}\wrong
\glsni{title}\glsni{leftbracechar}A Simple Document
\end{alltt}
instead of
\begin{alltt}\correct
\glsni{title}\marg{A Simple Document}
\end{alltt}

\item This error can also be caused by omitting the 
\gls{mandatory} of an \gls{environment}.  For example, 
this error will occur
if you do, say,
\begin{alltt}\wrong
\glsni{begin}\marg{thebibliography}
\glsni{bibitem}\marg{kopka95} A Guide to \glsni{LaTeXe}
\end{alltt}
instead of
\begin{alltt}\correct
\glsni{begin}\marg{thebibliography}\marg{1}
\glsni{bibitem}\marg{kopka95} A Guide to \glsni{LaTeXe}
\end{alltt}

\end{itemize}

%%%%%%%%%%%%%%% Something's wrong--perhaps a missing \item %%%%%%%%%%%%%%%%

\setnode{missingitem}
\section{Something's wrong--perhaps a missing \textbackslash item}

You may have missed an \gls{item} command.  The first object in
a list environment must either be an \gls{item} command, or another
list environment.

This error will be caused by, say,
\begin{alltt}\wrong
\glsni{begin}\marg{itemize}
Animal
\glsni{item} Vegetable
\glsni{item} Mineral
\glsni{end}\marg{itemize}
\end{alltt}
instead of
\begin{alltt}\correct
\glsni{begin}\marg{itemize}
\glsni{item} Animal
\glsni{item} Vegetable
\glsni{item} Mineral
\glsni{end}\marg{itemize}
\end{alltt}

\bookpagebreak
This error can also be caused by a missing \gls{bibitem} in the
\htmlref{bibliography}{sec:bib}.  For example, the error will occur
if you do, say,
\begin{alltt}\wrong
\glsni{begin}\marg{thebibliography}\marg{1}
A Guide to \glsni{LaTeXe}
\end{alltt}
instead of
\begin{alltt}\correct
\glsni{begin}\marg{thebibliography}\marg{1}
\glsni{bibitem}\marg{kopka95} A Guide to \glsni{LaTeXe}
\end{alltt}

See also \gls{ukfaq} entry: \faqlink{Perhaps a missing \cmdname{item}?}{errmissitem}.

%%%%%%%%%%%%%%%%% There's no line here to end %%%%%%%%%%%%%%%%%%%%%

\setnode{nolinetoend}
\section{There's no line here to end}

You have placed a line breaking command (such as \gls{newline.dbbackslashchar},
\gls{newline} or \gls{linebreak}) where it doesn't make
sense to have one.

%%%%%%%%%%%%%%%%% Undefined control sequence %%%%%%%%%%%%%%%%%%%%%%

\setnode{csundefined}
\section{Undefined control sequence}

\LaTeX\ doesn't understand the \gls{command} you
have used.  

\begin{itemize}
\item Check to see if you have misspelt the command name (remember
that all \LaTeX\ command names are case-sensitive.)

You will get this error if you do, say,
\begin{alltt}\wrong
This is a simple \cmdname{Latex}\glsni{spacesym}document
\end{alltt}
instead of
\begin{alltt}\correct
This is a simple \gls{LaTeX}\glsni{spacesym}document
\end{alltt}

\item Check that you have remembered the space when typing
\gls{spacesym} (backslash space).
For example, this error will occur if you do, say,
\begin{alltt}\wrong
This is a \cmdname{LaTeX}\cmdname{sample} document.
\end{alltt}
instead of
\begin{alltt}\correct
This is a \glsni{LaTeX}\glsni{spacesym}sample document
\end{alltt}

\item If you are using a command that is defined in a
\htmlref{package}{sec:packages} make sure you have included the
package using \gls{usepackage}.

\item Check that your command name hasn't run into the next piece
of text.  For example,
you can do
\begin{alltt}
man\glsni{oe}\marg{}uvre
\end{alltt}
or
\begin{alltt}
man\glsni{oe} uvre
\end{alltt}
or (not recommended)
\begin{alltt}
man\marg{\glsni{oe}}uvre
\end{alltt}
but not
\begin{alltt}\wrong
man\cmdname{oeuvre}
\end{alltt}

\item Check if you have used a backslash instead of a forward slash as a
directory divider.  (Remember that when using \LaTeX\ under Windows,
you need to use a forward slash (\glsi{dir.slash}) as the directory divider, as a
backslash would be interpreted as a command.)

For example, suppose you have a file called \texttt{shapes.pdf} in a 
subdirectory called \texttt{pictures}, then you would get an error 
if you did
\begin{alltt}\wrong
\glsni{includegraphics}\marg{pictures\cmdname{shapes}}
\end{alltt}
instead of
\begin{alltt}\correct
\glsni{includegraphics}\marg{pictures/shapes}
\end{alltt}

\end{itemize}

%%%%%%%%%% You can't use `macro parameter character #' in horizontal mode %%%%%%%%%%%%%%%

\setnode{hashinhmode}
\section{You can't use `macro parameter character \#' in horizontal mode}

You have used the special character \gls{hashchar} where you shouldn't
have.  Recall from \sectionref{sec:chars} that if you want a \# sign to appear you need
to do \gls{hash} not just \glsni{hashchar}.

This error message will be caused by doing, say,
\begin{alltt}\wrong
Item \glsni{hashchar}1
\end{alltt}
instead of
\begin{alltt}\correct
Item \glsni{hash}1
\end{alltt}

\setnode{help}
\chapter{Need More Help?}
\label{ch:help}

First, try to find your query in the \gls{ukfaq}. \Gls{tug}
also has a list of useful resources at
\url{http://tug.org/interest.html}. If you're still stuck, you can
post your question on a (La)TeX forum, newsgroup or mailing list,
such as those listed below. If you do post a question, remember
you're asking people who only have an altruistic interest in
helping. No one is paying them to help you. Most of the class files
and packages were written for free by people who had a need to solve
a particular problem and decided to make their work publicly
available. So no matter how frustrated you're feeling, stick to
being polite. If you can't work out how to use a particular class or
package, don't start by heaping offensive, unconstructive criticism on it as there's a chance
the author will read the message. There's no sense in alienating the
person most qualified to answer your question. In your message,
stick to the following guidelines:
\begin{enumerate}
\item Cut to the chase. In other words, be concise about the nature
of the problem. Don't write lots of long-winded paragraphs.

\item Provide a minimal
example\footnote{see \url{http://www.dickimaw-books.com/latex/minexample/}} that 
illustrates the problem.

\end{enumerate}

\xminisec{Example:}
\begin{verbatim}
I'm trying to use the \foo command in the "bar" package, 
but I'm getting the following error message:
! Undefined control sequence.
l.4 \foo

Here's a minimal example:
\documentclass{scrartcl}
\usepackage{bar}
\begin{document}
\foo{Blah}
\end{document}

I'm using bar version 1.0 (2012/06/30).
\end{verbatim}

\xminisec{Another example:}
\begin{verbatim}
I'm using the \foo command in the "bar" package. According
to the documentation, this command should display its argument
in a bold font, but it's coming out in italic instead. 
Anyone know why?

Here's a minimal example:
\documentclass{scrartcl}
\usepackage{bar}
\begin{document}
\foo{Blah}
\end{document}

I'm using bar version 1.1 (2012/07/30).
\end{verbatim}

There's no guarantee that you will get an answer, but if you follow
the above guidelines, you will increase your chances.

\xminisec{Resources}
\begin{itemize}

\item The \LaTeX\ Community (\url{http://www.latex-community.org/}).

\item \TeX/\LaTeX\ on StackExchange
(\url{http://tex.stackexchange.com/}).

\item \htmladdnormallink{\texttt{comp.text.tex}
newsgroup}{http://groups.google.com/group/comp.text.tex} (use a
newsreader rather than the Google interface if you want to avoid the
spam).

\item \htmladdnormallink{\texttt{texhax} archives}{http://tug.org/pipermail/texhax/}\doifbook{~\cite{texhax}}.

\end{itemize}
I~strongly recommend that you also have a look at the On-Line
Catalogue~\cite{texcat}.
It's also a good idea to look at the documentation that was
installed with your \TeX/\LaTeX\ distribution (see
\sectionref{sec:texdoc}). If you are using \itexdistro{MiKTeX} you can access the
on-line help via the Start Menu: \startmenu{MiKTeX \menuto\ Help}

(Please don't send your problems to me, unless you want to hire a
consultant. I~read both the \LaTeX\ Community Forum and
\texttt{comp.text.tex} and answer relevant questions if I~have time,
but it clogs up my inbox if people keep sending attachments that are in
the order of several megabytes in size.) Besides, you'll reach a
wide group of experts if you post to a newsgroup, forum or mailing
list, rather than a single busy individual.

%%%%%%%%%%%%%%%%%%%%%%%%%% Supplemental Material %%%%%%%%%%%%%%%%%%%

\ifpdf\else\setnode{supplemental}
\chapter{Supplementary Material}
\label{ch:supp}

This appendix chapter contains three sections from earlier versions of this
book. The information here is a little dated, but is provided in
case it's still of any use.

%%%%%%%%%%%%%%%%%%% Terminal + Text Editor %%%%%%%%%%%%%%%%%%%

\setnode{editorandterminal}
\section{Text Editor and Terminal Approach}
\label{sec:editorandterminal}

Creating a \LaTeX\ document using a text editor and a terminal
is an approach often favoured by UNIX-type users. If you have never
used a terminal (i.e.\ you have only ever used point-and-click menu
driven applications) then you will be better off using a
front-end, in which case I suggest you
\begin{latexonly}%
turn to Sections~\ref{sec:TeXnicCenter} and~\ref{sec:winedt} which 
describe \appname{TeXnicCenter} and \appname{WinEdt}, respectively.
\end{latexonly}
\begin{htmlonly}
look at the \htmlref{\appname{TeXnicCenter}}{sec:TeXnicCenter} and 
\htmlref{\appname{WinEdt}}{sec:winedt} sections instead.
\end{htmlonly}
\par
To begin, you will first need a text editor\faq{TeX-friendly editors
and shells}{editors}. There are a number available that are suited 
to using with \LaTeX, some people advocate
\iappnamelink{Emacs}{http://www.gnu.org/software/emacs/emacs.html},
others advocate \iappnamelink{vim}{http://vim.sourceforge.net/}, and
there are various others such as
\iappnamelink{NEdit}{http://www.nedit.org/}.  I prefer to use
\iappname{vim}---I'm not overly keen on using the mouse, and I prefer
being able to issue all commands via the keyboard (although 
there is a GUI version of \iappname{vim}). As with some other
editors, it comes with syntax highlighting, regular expression search
and replace, auto-insertion, and a brace matching mechanism which I
find useful. If you are using version~7 of \iappname{vim}, there is an
integrated spell checker, otherwise there is a spell checker plug-in
called
\htmladdnormallink{vimspell}{http://www.vim.org/scripts/script.php?script_id=465},
so you can check your spelling as you type.  If there is already a
text editor that you are comfortable with, then stick with that,
otherwise try out available editors, and decide which one you prefer.

When using the terminal and text editor approach, most people usually
have at least two terminals open: one to run the editor, the other to
run \LaTeX. This means that you don't have to keep quitting the
editor every time you want to \htmlref{\LaTeX\ your
document}{itm:step2}. Some editors allow you to run commands, but
personally I don't like to use this approach. If your editor has a
GUI interface, then you'll probably only need one terminal open.

Let's get started: start up your text editor. This is usually done by
entering the name of the editor at the command prompt in your
terminal. With most editors you can also specify the filename as
well. If the file doesn't exist, a new one will be created when you
save your document. \figureref{fig:startvim} shows my terminal.  The
command prompt looks like \verb|[nlct@nlctltpc examples]$|.  It will
be different for your system. I have typed \texttt{vim sample1.tex}
at the command prompt. This will start \iappname{vim} with a new file
called \texttt{sample1.tex}.  (In this section, I will be using
\iappname{vim} as the text editor, if you are not using
\iappname{vim}, then substitute the editor of your choice.)

\begin{figure}[htbp]
\figconts[alt=Image of a terminal]
 {pictures/terminal1}
 {\caption{Starting vim from a terminal}}
 {fig:startvim}
\end{figure}

Once I have pressed the return key, my terminal looks like
\figureref{fig:terminal2}. Normally \iappname{vim} starts in visual
command mode, which means that when you start typing text, it will be
interpreted as part of a command. In order to type text into your
file, you will need to enter input mode. There are a number of ways
of doing this, but pressing \texttt{i} will do for now\footnote{For a
complete set of available commands, see the \appname{vim} manual at
\url{http://vimdoc.sourceforge.net/}}. \figureref{fig:terminal3}
shows how my terminal looks when I am in input mode. I can now go
ahead and type in my text (\figureref{fig:terminal4}). To go back to
the visual command mode, press the escape key (Esc). Now that you are
back in the visual command mode, you can save your document, either
using the command \texttt{:w} if you have already given your file a
name, or \texttt{:w }\meta{filename} (e.g.\ \texttt{:w sample1.tex})
if you started \iappname{vim} without specifying a file, see
\figureref{fig:terminal5}. When you want to quit \iappname{vim} you
can do \texttt{:wq} to save and quit or \texttt{:q!} to quit without
saving, but I suggest you don't do this just yet if you have another
terminal available.

\begin{figure}[htbp]
 \figconts[alt=Image of a terminal with vim running]
 {pictures/terminal2}
 {\caption{Starting a new file in vim}}
 {fig:terminal2}
\end{figure}

\begin{figure}[htbp]
  \figconts[alt=Image of a terminal with vim running in input mode]
  {pictures/terminal3}
  {\caption{Input mode in vim}}
  {fig:terminal3}
\end{figure}

\begin{figure}[htbp]
  \figconts[alt=Image of a terminal with LaTeX code typed in vim]
  {pictures/terminal4}
  {\caption{Creating a sample document in vim}}
  {fig:terminal4}
\end{figure}

\begin{figure}[htbp]
  \figconts[alt=Image of new file being saved in vim]
  {pictures/terminal5}
  {\caption{Saving your document in vim (the file name should be 
    omitted if the file already has a name)}}
  {fig:terminal5}
\end{figure}

\objectref{Step}{itm:step1} is now complete, and you are now ready to
move on to \objectref{Step}{itm:step2}: using \LaTeX. Go to your
other terminal (or quit your editor if you only have access to one
terminal) and make sure that you are in the same directory as the
file you just created.  Typing \texttt{ls} at the command prompt will
list the contents of your current directory. If you do this, you
should see the file that you have just created (see
\figureref{fig:terminal6}).  At the command prompt type\faq{Makefiles
for LaTeX documents}{make}:
\refstepcounter{object}\label{obj:termlatex}
\begin{verbatim}
latex sample1.tex
\end{verbatim}
as shown in \figureref{fig:terminal7}. You can omit the \texttt{.tex}
extension if you like, \LaTeX\ will automatically add this if
it has been omitted. If you prefer to use \iPDFLaTeX, type
\begin{verbatim}
pdflatex sample1.tex
\end{verbatim}
instead (again the \texttt{.tex} extension may be omitted.)
You should now see something like \figureref{fig:terminal8}.

\begin{figure}[htbp]
  \figconts[alt=Image of directory listing in a terminal]
  {pictures/terminal6}
  {\caption{Listing the contents of the current directory}}
  {fig:terminal6}
\end{figure}

\begin{figure}[htbp]
  \figconts[alt=Image of a terminal where user has typed the command
   to run LaTeX]
  {pictures/terminal7}
  {\caption{Running \LaTeX}}
  {fig:terminal7}
\end{figure}

\begin{figure}[htbp]
  \figconts[alt=Image of a terminal showing normal LaTeX messages]
  {pictures/terminal8}
  {\caption{Running \LaTeX}}
  {fig:terminal8}
\end{figure}

\refstepcounter{object}\label{obj:transcript}
Numbers appearing in square brackets, e.g.\ \texttt{[1]}, indicate
which page \LaTeX\ is currently processing.  In this case, there is
only one page.  The last line to appear on screen
indicates that information about this \LaTeX\ run has been written
to the log file \texttt{sample1.log}, which you can look at using 
your text editor.

The most important thing to note is the penultimate line\footnote{if
you are using \PDFLaTeX, it will have \texttt{sample1.pdf} instead of
\texttt{sample1.dvi}}:
\begin{verbatim}
Output written on sample1.dvi (1 page, 248 bytes).
\end{verbatim}
This means that the document has been successfully created, and
is one page long.

If you have made a mistake in your \htmlref{source code}{sec:source},
for example suppose you have missed the starting backslash in
\cmdname{documentclass}, then the output will look something like:
\begin{verbatim}
! LaTeX Error: Missing \begin{document}.

See the LaTeX manual or LaTeX Companion for explanation.
Type  H <return>  for immediate help.
 ...

l.1 d
     ocumentclass[a4paper]{article}
?
\end{verbatim}
There are several things you can do at this point, but the easiest
thing to do is to exit \LaTeX\ by typing \texttt{X} followed by the
return key. Go back to your editor, fix the mistake, save the
document, and then try again. (If you do get an error message, check
\chapterref{ch:errors}.) Note that it is important to always save your
document before running \LaTeX.

\refstepcounter{object}\label{obj:xdvi}
You can view the typeset output by loading the file\indexDVI\  
\texttt{sample1.dvi} into a DVI viewer, such as \iappname{xdvi}
or \iappname{kdvi}. To do this type
\begin{verbatim}
xdvi sample1.dvi
\end{verbatim}
or
\begin{verbatim}
kdvi sample1.dvi
\end{verbatim}
at the command prompt (see
\figureref{fig:terminal9}). You will then see the final output,
as shown in \figureref{fig:terminal10}.

\begin{figure}[htbp]
 \figconts[alt=Image of a terminal where user has typed command to
  load DVI file into kdvi]
  {pictures/terminal9}
  {\caption{Load a DVI file into a DVI viewer}}
  {fig:terminal9}
\end{figure}

\begin{figure}[htbp]
  \figconts[alt=Image of DVI file being viewed using kdvi]
  {pictures/terminal10}
  {\caption{Viewing a DVI file in kdvi}}
  {fig:terminal10}
\end{figure}

If you have used \iPDFLaTeX\ instead of \LaTeX, you should have a
file called \texttt{sample1.pdf} instead of \texttt{sample1.dvi}.
You can view this using a PDF viewer, such as \iappname{acroread}
or \iappname{kpdf}.

Some viewers, such as \iappname{kdvi} and \iappname{kpdf} will
automatically reload the file whenever it is modified, in which case
you may like to keep the viewer open, and as you keep editing and
\LaTeX{}ing your document, the viewer will automatically reload the
new versions. Some viewers, such as \iappname{xpdf} don't
automatically reload, but have a reload facility, which you can use
whenever you \htmlref{\LaTeX\ your document}{itm:step2}.

\refstepcounter{object}\label{obj:dvips}
If you like, you can convert your DVI file\indexDVI\ to PostScript 
using \iappname{dvips}. To do this, type the following at the command
prompt in your terminal:
\begin{verbatim}
dvips -o sample1.ps sample1.dvi
\end{verbatim}
(The \texttt{.dvi} extension may be omitted.)
You can then view the PostScript file using \iappname{ghostscript}
or one of its associated applications, such as \iappname{ghostview}
or \iappname{GSview}. I have \iappname{kghostview} installed on
my laptop, so to view the PostScript file, \texttt{sample1.ps},
I would need to type:
\begin{verbatim}
kghostview sample1.ps
\end{verbatim}
(See \figureref{fig:terminal11}.)

\begin{figure}[htbp]
  \figconts[alt=Image of terminal where user has typed command to
   load PostScript file into kghostview]
   {pictures/terminal11}
  {\caption{Loading a PostScript file}}
  {fig:terminal11}
\end{figure}

%%%%%%%%%%%%%%%%% TeXnicCenter %%%%%%%%%%%%%%%%%%%%%%%%%%%%

\setnode{TeXnicCenter}
\section{TeXnicCenter}
\label{sec:TeXnicCenter}

\iappname{TeXnicCenter} is an application that enables you to edit
\LaTeX\ source code, and simply click on a button to pass the source
code to \LaTeX, and then click on another button to view the
resulting typeset document.  Many people prefer this approach to the
\latexhtml{text editor and terminal approach described in the
previous section}{\htmlref{text editor and terminal}{sec:editorandterminal}
approach}. This section gives a brief overview of 
\appname{TeXnicCenter}, however it has been several years since
I last used it\footnote{I used to use TeXnicCenter and
MiKTeX when I was teaching \LaTeX, but that was my limit of using
\TeX\ under Windows}, so this information may be dated.

\appname{TeXnicCenter} is free and can be downloaded from \gls{ctan} in the
\htmladdnormallink{\texttt{systems/win32/TeXnicCenter/}}{http://www.tex.ac.uk/tex-archive/systems/win32/TeXnicCenter/}
directory or from \url{http://www.toolscenter.org/}.  Note that you
must have a \htmlref{\TeX/\LaTeX\ distribution}{obj:MiKTeX} installed
before you install \appname{TeXnicCenter}. If you installed
\itexdistrolink{proTeXt}{http://www.tug.org/protext}, you should
already have \appname{TeXnicCenter} installed. If you have any
problems with installing or running \appname{TeXnicCenter}, go to
their \latexhtml{help page at
\url{http://www.texniccenter.org/help.html}}{\htmladdnormallink{help
page}{http://www.texniccenter.org/help.html}}.

Once the installation is complete, you can then run
\appname{TeXnicCenter} from the Start Menu: \startmenu{TeXnicCenter
\menuto\ TeXnicCenter} Firstly you should see the tip of the day
window (\figureref{fig:TXCtip}.)

\begin{figure}[hbtp]
  \figconts[alt=Image of TeXnicCenter's tip of the day window]
  {pictures/texniccenter/tip}
  {\caption{TeXnicCenter Tip of the Day Window}}
  {fig:TXCtip}
\end{figure}

You can close this window, and then, if this is the first time you
are using \appname{TeXnicCenter} you will have to use the
configuration wizard to set up \appname{TeXnicCenter} correctly.  I
would recommend that you choose the default settings. (Select
\menu{\underline{N}ext}, \menu{\underline{N}ext} and then
\menu{Finish}.)

\begin{figure}[hbtp]
  \figconts[alt=Image of TeXnicCenter's configuration wizard: welcome
   screen]
  {pictures/texniccenter/setup01}
  {\caption{TeXnicCenter Configuration Wizard}}
  {fig:TXCsetup1}
\end{figure}

\begin{figure}[hbtp]
  \figconts[alt=Image of TeXnicCenter's configuration wizard: use
MiKTeX option selected]
  {pictures/texniccenter/setup02}
  {\caption{TeXnicCenter Configuration Wizard}}
  {fig:TXCsetup2}
\end{figure}

\begin{figure}[hbtp]
\figconts[alt=Image of TeXnicCenter's configuration wizard: finish
configuration]
  {pictures/texniccenter/setup03}
  {\caption{TeXnicCenter Configuration Wizard}}
  {fig:TXCsetup3}
\end{figure}

Now you are ready to use \appname{TeXnicCenter}.  It should look like
\figureref{fig:TXCuse1}.

\begin{figure}[hbtp]
\figconts[alt=Image of TeXnicCenter: no files loaded]
  {pictures/texniccenter/use01}
  {\caption{TeXnicCenter}}
  {fig:TXCuse1}
\end{figure}

\mbox{}\refstepcounter{object}\label{TXCnewproject}%
To start a new project select \menu{File \menuto\ New Project}.
This will open the window shown in \figureref{fig:TXCuse2}.

\begin{figure}[hbtp]
\figconts[alt=Image of TeXnicCenter's New Project dialog box]
  {pictures/texniccenter/use02}
  {\caption{New Project Dialog Box}}
  {fig:TXCuse2}
\end{figure}

Enter a name for your project, and specify the directory where you
want to save your work.  For example, I shall call my project
\dq{example} and I want to save it in 
\verb|c:\My Documents\Nicky\example| (see \figureref{fig:TXCuse3}.)

\begin{figure}[hbtp]
\figconts[alt=Image of TeXnicCenter's New Project dialog box:
project name has been entered]
  {pictures/texniccenter/use03}
  {\caption{New Project Dialog Box}}
  {fig:TXCuse3}
\end{figure}

Select the \dq{Empty Project} icon, and click on \dq{Okay}.  You
should now see something like \figureref{fig:TXCuse4}.

\begin{figure}[hbtp]
\figconts[alt=Image of TeXnicCenter: new project on view]
  {pictures/texniccenter/use04}
  {\caption{TeXnicCenter --- New Project Started}}
  {fig:TXCuse4}
\end{figure}

You can now start typing the source code (we'll cover this
\htmlref{later}{ch:simpledoc}).  See \figureref{fig:TXCuse5}.

\begin{figure}[hbtp]
\figconts[alt=Image of TeXnicCenter: user has typed source code in
the edit panel]
  {pictures/texniccenter/use05}
  {\caption{TeXnicCenter --- Typing in Source Code}}
  {fig:TXCuse5}
\end{figure}

Save it by either clicking on the save icon 
\incGraphics[alt=Image of save icon]{pictures/texniccenter/saveicon} or 
select \menu{File \menuto\ Save}

Now select what type of output you want (DVI, PDF or
\Index{PostScript}\refstepcounter{object}\label{obj:txcdvips}) see
\figureref{fig:TXCuse6}.  If this box is blank, then it's possible
that you didn't complete all the steps in the configuration wizard
described above.

\begin{figure}[hbtp]
\figconts[alt=Image of TeXnicCenter: drop-down menu showing]
  {pictures/texniccenter/use06}
  {\caption{TeXnicCenter --- Selecting Output Type}}
  {fig:TXCuse6}
\end{figure}

\refstepcounter{object}\label{TXC:latex}
Now click on the build output icon 
\incGraphics[alt=Image of build icon]{pictures/texniccenter/build} or 
select \menu{Build \menuto\ Build Output}.  The transcript will be
written in the window at the bottom (see \xfigureref{fig:TXCuse7})
This transcript should be the same as described \latexhtml{on
page~\pageref{obj:transcript}
onwards}{\htmlref{earlier}{obj:transcript}}.  If you have selected
\verb|LaTeX => PDF|, then \appname{TeXnicCenter} will use \iPDFLaTeX\
instead of \LaTeX.  If you have selected \verb|LaTeX => PS|, then
\appname{TeXnicCenter} will use \LaTeX\ followed by \iappname{dvips}
(as in \xfigureref{fig:TXCuse7}).  The \appname{dvips} messages will
follow on from the \LaTeX\ messages.  (If you selected the BibTeX or
MakeIndex features when you initialised the project,
\figureref{fig:TXCuse3}, then \iappname{TeXnicCenter} will also use
the \iBiBTeX\ and \appname{MakeIndex} applications.)

\begin{figure}[hbtp]
\figconts[alt=Image of TeXnicCenter: messages have been printed on
the build panel]
  {pictures/texniccenter/use07}
  {\caption[TeXnicCenter (using LaTeX and dvips)]{TeXnicCenter (using \LaTeX\ and \appname{dvips})}}
  {fig:TXCuse7}
\end{figure}

\refstepcounter{object}\label{TXC:yap}
To view the document, click the View Output button 
\incGraphics[alt=Image of view icon]{pictures/texniccenter/view}.
(Note that if you have selected \verb|LaTeX => PDF| or
\verb|LaTeX => PS| you will need \iappname{Adobe Reader}
or \iappname{GSview}, respectively, to view the output file.)

\refstepcounter{object}\label{TXCerrors}
If there are any errors, you can select \menu{Build \menuto\ Next
Error} and it will show you where the error has occured (See
\figureref{fig:TXCuse8}).  If you do have any errors, check
\chapterref{ch:errors}.

\begin{figure}[hbtp]
\figconts[alt=Image of TeXnicCenter: error message displayed in
build panel. A red arrow points to line containing the error in the
edit panel.]
  {pictures/texniccenter/use08}
  {\caption{TeXnicCenter --- Showing Error}}
  {fig:TXCuse8}
\end{figure}


%%%%%%%%%%%%%%%%%%%%%%%%%%%%% WinEdt %%%%%%%%%%%%%%%%%%%%%%%%%%%%%

\setnode{winedt}
\section{WinEdt}
\label{sec:winedt}

\iappname{WinEdt} (not to be confused with \appname{WinEdit} which is
a completely different application) is an application that enables
you to edit \LaTeX\ source code, and simply click on a button to pass
the source code to \LaTeX, and then click on another button to view
the resulting typeset document.  This section gives a brief overview
of \iappname{WinEdt}, however it has been several years since I
last used it, so this information may be dated.

\iappname{WinEdt} is shareware: it can be downloaded from \gls{ctan} in the
\htmladdnormallink{\texttt{systems/win32/winedt}}{http://www.tex.ac.uk/tex-archive/systems/win32/winedt/}
directory or from \url{http://www.winedt.com/} and evaluated for a
trial period of 31 days, after which, if you want to continue to use
it, you must pay the registration fee.  Details of prices and types
of licence available can be found at \url{http://www.winedt.com/}.
Note that members of the \gls{uktug} can benefit from their 
\iappname{WinEdt} licence
scheme (see \url{http://uk.tug.org/membership} for further details).

Again, you must have a \htmlref{\TeX/\LaTeX\
distribution}{obj:MiKTeX} installed before you start.
\appname{WinEdt} is fairly easy to install.  First unpack all the
files, and then run the \appname{setup.exe} application.  I recommend
that you use the default settings.  If you have any problems
installing or using \appname{WinEdt}, go to
\url{http://www.winedt.com/support.html}.

To run \iappname{WinEdt}, select \appname{WinEdt} from the start menu:
\startmenu{WinEdt \menuto\ WinEdt}
It should look like \figureref{fig:winedt1}.

\begin{figure}[hbtp]
\figconts[alt=Image of WinEdt: no files loaded]
  {pictures/winedt/01}%
  {\caption{WinEdt}}
  {fig:winedt1}
\end{figure}

Click on the \dq{New Document} button or select \menu{File \menuto\
New}.  You can now start typing your \htmlref{source
code}{sec:source} into the \appname{WinEdt} window, as shown in
\figureref{fig:winedt3}.

\begin{figure}[htbp]
\figconts[alt=Image of WinEdt: user has typed source code]
  {pictures/winedt/03}
  {\caption{WinEdt}}
  {fig:winedt3}
\end{figure}

You can now save your document using the \menu{File \menuto\ Save as}
menu.  Select the file type to be \texttt{TeX}, and type in the name
of your file, e.g.\ \texttt{sample1.tex}.  See
\figureref{fig:winedt4}.

\begin{figure}[htbp]
  \figconts[alt=Image save file dialog]
  {pictures/winedt/04}
  {\caption{WinEdt --- Saving the File}}
  {fig:winedt4}
\end{figure}

\refstepcounter{object}\label{winedt:latex}
To \LaTeX\ your document, simply click on the \LaTeX\ button
\incGraphics[alt=Image of LaTeX icon]{pictures/winedt/latex}.
The output will appear in an MSDOS Prompt window (see
\figureref{fig:winedt5}).

\begin{figure}[htbp]
\figconts[alt=Image of MSDOS prompt showing standard LaTeX messages]
  {pictures/winedt/05}%
  {\caption[WinEdt --- LaTeX Output]{WinEdt --- \LaTeX\ Output}}
  {fig:winedt5}
\end{figure}

\refstepcounter{object}\label{winedt:yap}
To view your typeset document, click on the \dq{view DVI}\indexDVI\ 
button \incGraphics[alt=Image of previewer icon]{pictures/winedt/yap}.

\refstepcounter{object}\label{obj:winedtdvips}
You can convert your DVI file\indexDVI\ to \Index{PostScript} by clicking 
the \incGraphics[alt=Image of dvips icon]{pictures/winedt/dvips} button.  
If you have \iappname{GSview} installed, you can then view
the PostScript file by clicking on the 
\incGraphics[alt=Image of GSView icon]{pictures/winedt/gsview2} button.  

Depending on which version of \iappname{WinEdt} you have installed,
there may also be a \iPDFLaTeX\ button which you can click on to
create a Portable Document Format (\texttt{.pdf}) document.  If not,
you can click on the 
\incGraphics[alt=Image of MSDOS icon]{pictures/winedt/msdos} button to 
open up an MS-DOS Prompt window, and type \texttt{pdflatex} followed by 
the filename. For example:
\begin{verbatim}
pdflatex sample1.tex
\end{verbatim}
Note that if the filename contains a
space, you will need to use double quotes:
\begin{verbatim}
pdflatex "my file.tex"
\end{verbatim}

\fi

\backmatter

%%%%%%%%%%%%%%%%%%%%%%%%%%% Bibliography %%%%%%%%%%%%%%%%%%%%%%%%%%%%%%%%

\clearpage\phantomsection
\setnode{bibliography}
\bibliographystyle{plain}
\bibliography{novices}

%%%%%%%%%%%%%%%%%%%%%%%%%%%% GLOSSARIES %%%%%%%%%%%%%%%%%%%%%%%%%%%%

\setnode{acronyms}
\printglossary[type=acronym,nonumberlist,style=acronyms]

\setnode{summary}
\setdoublecolumnglo
\printglossary[style=summary]

\setnode{bookindex}
\printindex

%%%%%%%%%%%%%%%%%%%%%%%%% LICENCE %%%%%%%%%%%%%%%%%%%%%%%%%%%%%%%%%%

\setnode{licence}
\input{fdl}

%%%%%%%%%%%%%%%%%%%%%%%%%% HISTORY %%%%%%%%%%%%%%%%%%%%%%%%%%%%

\setnode{history}
\chapter{History}\markright{History}

\historyitem{25th Sept 2012 (Version 1.4)}

\begin{itemize}
\item Paperback edition~1 produced.
\item Added TeXWorks section.
\item Moved \dq{Text editor and Terminal approach}, \dq{TeXnicCenter}
and \dq{WinEdt} sections to new supplementary material.
\item Added hardcopy-related code.
\item Change to \koma\ classes (both for examples and for pdf
versions of this document).
\item Changed from using \sty{subfloat} to \sty{subcaption} package.
\item Added section on inter-sentence spacing.
\item Moved \dq{Errors} and \dq{Where to get Help?} to appendices.
\item Moved section \dq{Downloading and Installing Packages} to
new appendix chapter.
\item Moved introduction to packages to \dq{Creating a Simple
Document}.
\item Moved datetime section to \dq{Creating a Simple Document}
chapter. (Removed reference to ukdate package.)
\item Moved babel section to \dq{Structuring Your Document} chapter.
\item Moved graphicx section to its own chapter.
\item Added section on align.
\item Added \cmdname{cfrac}, \cmdname{substack} and \sty{amsmath} ellipses to maths chapter.
\item Added extensible arrows and \cmdname{bigl} etc to maths
chapter.
\item Added booktabs.
\item Moved lengths chapter to section in definitions.
\item Added summary chapter with commands hyperlinked to their
definitions in the summary.
\item Changed definitions chapter to use a glossary structure.
\item Moved bibliography into bib file.
\item Added \sty{varioref}.
\item Removed dependency on \sty{html} package (for pdf versions)
to avoid conflict between \sty{html} and \sty{varioref}
(\sty{html} package functions not defined by \sty{hyperref}
now emulated; \sty{comment} package loaded to provide
\envname{htmlonly} environment).
\item Removed image of equation written in Word (Microsoft have
improved their equation rendering) and added link to
Murray Sargent~III blog~\cite{sargentiii}.
\item Added section on what a terminal/command prompt is.
\item Added section on auxiliary files.
\item Added section on Perl.
\item Added information about latexmk
\item Mentioned grffile package.
\item Mentioned on-the-fly EPS conversion.
\item Mentioned \sty{etoolbox}'s \cmdname{appto} and \sty{babel}'s
\cmdname{addto}.
\item Changed to UTF-8 and mostly changed to using code points
instead of named entities in HTML files.
\item Moved the document's home page from
\url{http://theoval.cmp.uea.ac.uk/~nlct/latex/novices/} to
\url{http://www.dickimaw-books.com/latex/novices/}.
\end{itemize}

\historyitem{15th Jan 2008 (Version 1.3)}

The main reason behind this change was to increase accessibility and
conform to W3C guidelines. If you are experiencing problems relating
to accessibility, please let me know (clearly stating the problem).
\begin{itemize}
\item Corrected error in the university's post code on the
title page
\item Added alternative text tags to more of the images, and made
some of the images hyperlinks to a more detailed description of the
image.
\item Added information on how to break ligatures.
\item Moved information on TeX to the introduction, and removed
section on TeX that was in the "Some Definitions" chapter.
\item Document nodes now have permanent names instead of the generic
node\meta{n}.html which \LaTeX2HTML generates by default.
\item Went back to using straight double quotes in the HTML document
as the fancy typographic double quotes are nonstandard.

\end{itemize}

\historyitem{8th May 2007 (Version 1.2)}

\begin{itemize}
\item Links to \gls{ukfaq} added.
\item Overview made into a separate section, and tidied up a bit.
\item Added some extra definitions: moving arguments and fragile
commands, robust commands, short and long commands.
\item Changed \dq{Text editor and Terminal approach} to deal with
Unix-type systems rather than MS-DOS.
\item Moved section on \envname{tabular} environment.
\item Added section on boxes and mini-pages.
\item Segmented section on font changing commands.
\item Segmented section describing \sty{graphicx}.
\item Added section on the \sty{babel} package.
\item Updated and segmented section on downloading and installing
new packages.
\item Added section on side-by-side figures.
\item Updated section on sub-figures to use the new \sty{subfloat}
package instead of the obsolete \sty{subfigure} package.
\item Added \dq{Need More Help?} chapter.
\end{itemize}

\backcoverheading
(See \url{http://www.gnu.org/licenses/fdl-howto-opt.html#SEC2}.)

\backcovertext
\end{document}
